%% ---- begin preamble.tex
\documentclass[11pt]{amsart}
\usepackage[margin=1.15in]{geometry}
\usepackage{amssymb,amsmath,amsthm,mathtools}
\usepackage{needspace,microtype,etoolbox}
\usepackage{tikz-cd}
\usepackage{listings,xcolor,textcomp}
\definecolor{LeanKeyword}{HTML}{254E91}
\definecolor{LeanType}{HTML}{006B65}
\definecolor{LeanComment}{HTML}{526B46}
\definecolor{LeanString}{HTML}{8A4727}
\lstdefinelanguage{Lean}{
  sensitive=true,
  morekeywords={abbrev,by,class,def,end,exact,fun,import,in,instance,let,
    namespace,noncomputable,open,scoped,section,theorem,variable,where},
  morekeywords=[2]{Prop,Type,Set,NontriviallyNormedField,NormedAddCommGroup,
    NormedSpace,TopologicalSpace,ChartedSpace,ModelWithCorners,IsOpen,
    ContMDiffOn,ContMDiffVectorBundle},
  morecomment=[l]{--},morecomment=[s]{/-}{-/},morestring=[b]"
}
% Render Lean's Unicode notation with both the editor and PDF export engines.
\lstnewenvironment{leancode}{\lstset{
  language=Lean,
  basicstyle=\small\ttfamily,columns=fullflexible,keepspaces=true,
  keywordstyle=\color{LeanKeyword}\bfseries,
  keywordstyle=[2]\color{LeanType},
  commentstyle=\color{LeanComment}\itshape,
  stringstyle=\color{LeanString},
  showstringspaces=false,upquote=true,
  literate={∀}{{\ensuremath{\forall}}}1
           {γ}{{\ensuremath{\gamma}}}1
           {→}{{\ensuremath{\to}}}1
           {↦}{{\ensuremath{\mapsto}}}1
           {ℕ}{{\ensuremath{\mathbb N}}}1
           {∞}{{\ensuremath{\infty}}}1
           {ω}{{\ensuremath{\omega}}}1
           {×}{{\ensuremath{\times}}}1
           {⋀}{{\ensuremath{\bigwedge}}}1
           {𝓘}{{\ensuremath{\mathcal I}}}1
           {₁}{{\ensuremath{{}_1}}}1
           {₂}{{\ensuremath{{}_2}}}1
}}{}
\usepackage[colorlinks=true,linkcolor=blue,urlcolor=blue,citecolor=blue]{hyperref}
\hypersetup{pdftitle={Divided-power analytic maps between normed spaces},
 pdfauthor={Jack McCarthy},pdfsubject={Divided-power analytic maps between normed spaces},
 pdfkeywords={divided-power analyticity, polynomial laws, nonarchimedean analysis}}
\newtheorem{theorem}{Theorem}[section]
\newtheorem{lemma}[theorem]{Lemma}
\newtheorem{proposition}[theorem]{Proposition}
\newtheorem{corollary}[theorem]{Corollary}
\theoremstyle{definition}
\newtheorem{definition}[theorem]{Definition}
\theoremstyle{remark}
\newtheorem{remark}[theorem]{Remark}
\numberwithin{equation}{section}
\newcommand{\N}{\mathbb N}
\newcommand{\Z}{\mathbb Z}
\newcommand{\F}{\mathbb F}
\newcommand{\cL}{\mathcal L}
\newcommand{\Alt}{\operatorname{Alt}}
\newcommand{\Mult}{\operatorname{Mult}}
\newcommand{\End}{\operatorname{End}}
\newcommand{\linf}{\ell^\infty}
\newcommand{\sgn}{\operatorname{sgn}}
\newcommand{\spn}{\operatorname{span}}
\newcommand{\id}{\mathrm{id}}
\newcommand{\abs}[1]{\lvert #1\rvert}
\newcommand{\norm}[1]{\lVert #1\rVert}
\newcommand{\cC}{\mathcal{C}}
\newcommand{\cD}{\mathcal{D}}
\newcommand{\coeff}{\operatorname{coeff}}
\newcommand{\ord}{\operatorname{ord}}
\newcommand{\rank}{\operatorname{rank}}
\newcommand{\type}{\operatorname{type}}
\newcommand{\sdim}{\operatorname{sdim}}
\newcommand{\incl}{\mathrm{incl}}
\newcommand{\Cop}{C^{\mathrm{op}}}
\newcommand{\Cvec}{C^{\mathrm{vec}}}
\newcommand{\Kh}{\widehat K}
\theoremstyle{plain}
\newtheorem{introtheorem}{Theorem}
\renewcommand{\theintrotheorem}{\Alph{introtheorem}}
\theoremstyle{definition}
\newtheorem{example}[theorem]{Example}
\newtheorem{conjecture}[theorem]{Conjecture}
\newtheorem{question}[theorem]{Question}
\theoremstyle{plain}
\newcommand{\alt}{\operatorname{alt}}
\newcommand{\sh}{\operatorname{sh}}
\newcommand{\ev}{\operatorname{ev}}
\newcommand{\ob}{\operatorname{ob}}
\newcommand{\Ext}{\operatorname{Ext}}
\newcommand{\Poly}{\operatorname{Poly}}
\newcommand{\hot}{\mathbin{\widehat\otimes}}
\newcommand{\hLambda}{\widehat\Lambda}
\newcommand{\hGamma}{\widehat\Gamma}
\newcommand{\cN}{\mathcal N}
\newcommand{\cS}{\mathcal S}
\newcommand{\cR}{\mathcal R}
\newcommand{\cG}{\mathcal G}
\newcommand{\cM}{\mathcal M}
\newcommand{\cP}{\mathcal P}
\newcommand{\fD}{\mathfrak D}
\newcommand{\fX}{\mathfrak X}
\newcommand{\Lsym}{\Lambda_{\mathrm{sym}}}
%% ---- end preamble.tex
%% ---- begin front.tex
\setlength{\emergencystretch}{1em}
\title{Divided-power analytic maps between normed spaces}
\author{Jack McCarthy}
\thanks{This paper was developed and written with Anthropic's Claude and is not formalized.}
\makeatletter
\g@addto@macro\@setauthors{\vskip6pt{\centering\footnotesize October 2026\par}}
\makeatother
\begin{document}
\begin{abstract}
A bounded polynomial law is a map between normed spaces over a nonarchimedean field that is a homogeneous polynomial on
finite-dimensional subspaces and is bounded by a power of the norm. A germ is divided-power analytic, or DP, if it is a
sum of such laws with geometric bounds. The algebra of composition and of the formal inverse is Roby's. Here we give the
estimates that make a calculus over normed spaces: re-expansion on the whole ball, composition, inverse functions,
bundles, forms, jets and base change. Among germs with Banach values, the DP germs are the smallest class that contains
the analytic germs and has target descent and source descent. Between ultrametric spaces, the ambient analytic germs, the
germs analytic after a cover, the strongly DP germs and the DP germs coincide. Over a complete field, wherever the
analyticity of Bourbaki's map $\alpha_k$ between ultrametric Banach spaces is decided here, the agreement of DP germs with
analytic germs there is decided with the same answer. Both fail over complete, not spherically complete fields of
characteristic $p>0$; one case of the field is undecided. For general norms we list the disagreements between DP and
analytic germs and what forces each one. An example on $\ell^1$ shows that minimality on general sources depends on target
descent into non-ultrametric spaces: the germs analytic after a cover exclude this example and, over a complete field, also have a calculus. In
characteristic $p$, composition closes up strictly below DP: the tree analytic germs are closed under composition, have
inverse functions between Banach spaces, and are strictly fewer than the DP germs.
\end{abstract}
\maketitle
\setcounter{tocdepth}{1}
\begingroup
\small
\linespacing=8pt
\tableofcontents\removelastskip
\endgroup
\clearpage
%% ---- end front.tex
%% ---- begin sec-01-intro.tex
\section{Introduction}\label{dp:sec:intro}

Let $K$ be a nonarchimedean nontrivially valued field, fix $\varpi\in K$ with $0<\abs\varpi<1$, and let $P,W$ be normed $K$-spaces, not necessarily complete and not necessarily ultrametric. A \emph{bounded polynomial law} of degree $n$, or \emph{law}, is a map $a\colon P\to W$ that is a homogeneous polynomial of degree $n$ on every finite-dimensional subspace and satisfies $\norm{a(h)}\le C\norm h^n$ for some real $C$; the least such $C$ is $\norm a$. A map $f$ from an open subset of $P$ to $W$ is \emph{divided-power analytic}, or \emph{DP}, at $x_0$ if $f(x_0+h)=\sum_{n\ge0}a_n(h)$ for $\norm h<R$, with laws $a_n$ of degree $n$ and $\norm{a_n}\le MR^{-n}$ (Definition~\ref{dp:def:laws}). The name does not mean a series $\sum_na_nx^n/n!$. No division occurs, and if $K$ is complete and $P$ is finite-dimensional, a DP germ is a convergent power series (Theorem~\ref{dp:thm:conservativity}(5)). The name refers to the target of the universal law: a law of degree $n$ is a bounded linear map on the span $\Gamma^n(P)$ of the powers $h^{\otimes n}$, for the largest seminorm $\norm\cdot_\Gamma$ with $\norm{h^{\otimes n}}_\Gamma\le\norm h^n$ (Theorem~\ref{dp:thm:normed-roby}). Without the bound, this is Roby's description of laws by divided powers~\cite[Th.~IV.1, Prop.~IV.5]{Ro}.

A map is \emph{analytic} at $x_0$ if each $a_n$ above is the diagonal of an $n$-linear map $\beta_n\colon P^n\to W$ with $\norm{\beta_n}\le MR^{-n}$. Analytic germs are DP, with the same radius and bound. Roby rejected multilinear diagonals as a definition of polynomial maps: factorials get in the way, even for vector spaces over a finite field, and the multilinear map is not unique~\cite[p.~214]{Ro}. Section~\ref{dp:sec:laws} identifies laws with Roby's homogeneous polynomial laws. So, without the bounds, the laws of positive degree of a DP germ form a formal law in Roby's sense~\cite[Th.~VI.1, Th.~VI.2]{Ro}. So a DP germ is a Roby formal law plus a constant, with geometric bounds. The algebra of composition and of the formal inverse is Roby's~\cite[VI~\S8, Prop.~VI.6, Cor.; Th.~VI.3, Prop.~VI.8]{Ro}. We add only the estimates that make the series converge (Theorems~\ref{dp:thm:calculus} and~\ref{dp:thm:inverse}).

\subsection*{Laws that are not ambient}

Paper~\cite{I} enlarges the target. A map $f$ is \emph{ambient analytic}, or \emph{AA}, at $x_0$ if $j\circ f$ is analytic at $x_0$ for a linear isometry $j$ of $W$ into a normed space, with the diagonals of the coefficients in $j(W)$. A law $a$ of degree $n$ is \emph{ambient} if $j\circ a$ is the diagonal of a bounded $n$-linear map for such a $j$. An AA germ is a DP germ whose laws are ambient, with geometric bounds on the norms of these multilinear maps~\cite[Proposition~14.9]{I}. Over every $K$ of characteristic $p>0$, the class AA is not closed under composition~\cite[Theorem~14.17]{I}. If $\operatorname{char}K=p>0$, then in degrees below $p$ every law is the diagonal of a bounded multilinear map (Section~\ref{dp:sec:nonambient}). From degree $p$ on this can fail, even after an enlargement of the target.

\begin{introtheorem}[laws that are not ambient]\label{dp:intro:A}
Let $\operatorname{char}K=p>0$.
\begin{enumerate}
\item Over every $K$, in every degree $n\ge p$, there is a law of degree $n$ and norm at most $1$ that is not ambient: the map $(\lambda,h)\mapsto\lambda^{n-p}h^{\otimes p}$, with $\lambda^0=1$, from $K\oplus\mathcal E$, with the maximum norm, to $\Gamma^p(\mathcal E)$ with the norm $\norm\cdot_\Gamma$, for a suitable $\ell^1$-sum $\mathcal E$ of finite-dimensional spaces.
\item If $K$ is discretely valued or $p=2$, composition fails on $H_0=c_0\oplus G$, with the maximum norm, where $c_0=c_0(\N,K)$ and $G$ is a suitable $\ell^1$-sum of finite-dimensional spaces. Let $X=c_0$, and give $\Gamma^p(X)$ the projective norm. The map $f_X(x,w)=(x^{\otimes p},w)$ into $\Gamma^p(X)\oplus G$ is AA at every point, the map $g_X(s,w)=s\otimes w$ into the projective tensor product of $\Gamma^p(X)$ and $G$ is the diagonal of a bounded bilinear map, and $g_X\circ f_X$ is a law of degree $p+1$, DP everywhere and AA at no point.
\end{enumerate}
\end{introtheorem}

This is Proposition~\ref{dp:prop:degree-p}(2) and Theorem~\ref{dp:thm:c0-failure}. Over densely valued $K$ of odd characteristic, whether this failure occurs on $c_0\oplus G$ is open~\cite[Remark~14.20]{I}.

\subsection*{The two descents and minimality}

Consider classes of germs with Banach values and three properties of such a class. (C1): it contains the analytic germs. (C2), \emph{target descent}: if $j\colon W\to Z$ is a linear isometry of Banach spaces and $j\circ f$ is in the class, then $f$ is. (C3), \emph{source descent}: if $Q\colon P'\to P$ is a bounded linear surjection with a preimage constant and $f\circ Q$ is in the class, then $f$ is. Here $Q$ has \emph{preimage constant} $\mathsf c$ if every $h\in P$ has a preimage of norm at most $\mathsf c\norm h$.

\begin{introtheorem}[the smallest class with these descents]\label{dp:intro:B}
\leavevmode
\begin{enumerate}
\item The DP germs with Banach values are the smallest class with \textup{(C1), (C2), (C3)}.
\item Among germs with ultrametric sources and ultrametric targets, the DP germs are the smallest class with \textup{(C1)} and \textup{(C3)}.
\end{enumerate}
\end{introtheorem}

This is Theorem~\ref{dp:thm:minimal}(3) and Theorem~\ref{dp:thm:ultrametric-minimal}(2). Item~(2) is where ultrametric norms simplify. There the targets need not be complete. No target descent is needed: $P$ is the image of $c_{00}(\mathsf I)$, the finitely supported functions on a set $\mathsf I$ with the supremum norm, under a contractive surjection with a preimage constant, and a law on $c_{00}(\mathsf I)$ with values in an ultrametric space is already the diagonal of a multilinear map, with constant $\abs\varpi^{-n}$. Analytic germs have (C1) and, in every characteristic, fail (C2). In characteristic $p$ they fail (C3), and AA, the smallest class with (C1) and (C2), fails it too (Proposition~\ref{dp:prop:other-classes}).

Minimality is a statement about these axioms, and on general sources it depends on target descent into non-ultrametric spaces. Fix $\mathsf e_n\in K$ with $n^n\le\abs{\mathsf e_n}<\abs\varpi^{-1}n^n$, and put
\[
 \mathsf f(h)=\sum_{n\ge1}\mathsf e_nh_1h_2\cdots h_n
\]
on the open unit ball of $\ell^1(\N,K)$, with values in $\Kh$. By the arithmetic--geometric mean inequality, the laws have norm below $\abs\varpi^{-1}$. So $\mathsf f$ is DP at $0$. It is AA, so every class with (C1), (C2) contains it. In every characteristic it is not analytic, since a multilinear map whose diagonal is the law of degree $n$ has norm at least $n^n$ (Example~\ref{dp:ex:l1}). The series converges because of an inequality between real numbers on the source, and no divided power is involved.

Let CA be the class of germs \emph{analytic after a cover}: germs $f$ with $f\circ Q$ analytic at a preimage of the base point, for some bounded linear surjection $Q$ with a preimage constant. Among germs with Banach values it is the smallest class with (C1) and (C3) (Lemma~\ref{dp:lem:cover-flat}(3)). For $K$ complete it also has composition, pairing, inverse and implicit functions, it contains every law, and it does not contain $\mathsf f$ (Theorem~\ref{dp:thm:descents-independent}(1)). So on general sources (C1) and (C3) reach only CA. A class with (C1) and (C3) whose germs may have non-ultrametric sources contains every DP germ on an ultrametric source, with every target (Corollary~\ref{dp:cor:um-cover}); the cover used there is not ultrametric. In characteristic $p$, whether, among germs with ultrametric sources and Banach targets, the DP germs are the smallest class with (C1), (C2), (C3) is open (Question~\ref{dp:q:c0}(2)). The other way up to DP on general sources is an axiom for classes closed under translation in the source, the axiom (S): if laws $b_n$ of degree $n$ have germs at $0$ in the class and geometric bounds on their norms, then the germ of $\sum_nb_n$ at $0$ is in the class. Such a class with (C1), (C3), (S) contains every DP germ with Banach values, and in characteristic $0$ (C1) and (S) suffice (Propositions~\ref{dp:prop:s-axiom} and~\ref{dp:prop:s-char0}). The axiom (S) is close to the definition of DP.

\subsection*{Composition is not enough}

The failure of composition in Theorem~\ref{dp:intro:A} does not lead to DP. In characteristic $p$, repairing composition closes up strictly below DP. Let $\mathrm{AA}^\circ$ be the smallest class of germs with Banach values that contains the analytic germs and is closed under composition, pairing and target descent. Let $\pi$ be the projective seminorm on $\Gamma^N(P)\subseteq P^{\otimes N}$. A germ is AA if and only if the linear maps on $\Gamma^N(P)$ given by its laws have geometric bounds for $\pi$, and DP if and only if they have them for $\norm\cdot_\Gamma$. Between the two lie the \emph{tree norms} $\mathfrak t_N$, iterated projective seminorms along the compositions of the integer $N$, with $\pi\le\mathfrak t_N\le\norm\cdot_\Gamma$. A germ is \emph{tree analytic}, or TA, if its linear maps have geometric bounds for $\mathfrak t_N$ (Definitions~\ref{dp:def:tree} and~\ref{dp:def:ta}). Let $\cC_p$ be the class of DP germs whose laws of degree at most $p$ are ambient. Section~\ref{dp:sec:chain} also has two variants of $\cC_p$: a local one $\cC_p^{\rm loc}$, the germs in $\cC_p$ at every point near the base point, and a uniform one $\cC_p^{\rm u}$.

\begin{introtheorem}[the chain]\label{dp:intro:C}
Let $\operatorname{char}K=p>0$, with $K$ complete or not.
\begin{enumerate}
\item Among germs with Banach values,
\[
 \mathrm{AA}\subsetneq\mathrm{AA}^\circ\subseteq\mathrm{TA}\subseteq\cC_p^{\rm u}\subseteq\cC_p^{\rm loc}\subsetneq\cC_p\subsetneq\mathrm{DP}.
\]
\item TA and $\cC_p$ are closed under composition and pairing, and between Banach spaces they have inverse functions.
\item The map $\gamma^\sharp_p(\lambda,h)=h^{\otimes p}$ on $K\oplus\mathcal E$, for the space $\mathcal E$ of Theorem~\ref{dp:intro:A} and with the norm $\norm\cdot_\Gamma$ on values, is DP at every point, and its germ at no point is in $\cC_p$.
\item TA and $\cC_p$ fail \textup{(C3)}. A class with \textup{(C1), (C2)} has \textup{(C3)} only if it contains every DP germ with Banach values.
\end{enumerate}
\end{introtheorem}

This is Theorem~\ref{dp:thm:chain}, with Theorems~\ref{dp:thm:ta-composition}, \ref{dp:thm:ta-inverse}, \ref{dp:thm:cp}, \ref{dp:thm:cp-inverse} and~\ref{dp:thm:singles-out}. By Theorem~\ref{dp:thm:singles-out}, TA has (C1), (C2), composition, pairing, inverse and implicit functions, re-expansion on the whole ball, derivatives and bundle lifting. Of the properties listed there, only (C3) is not shared by TA, and with (C1), (C2) it forces DP. The separations in Theorem~\ref{dp:intro:C} and the failure of (C3) are all in characteristic $p$. In characteristic $0$, whether $\mathrm{AA}\ne\mathrm{DP}$ and whether $\mathrm{TA}\ne\mathrm{DP}$ are open. Between ultrametric spaces the two papers study one class, from outside in~\cite{I} and from inside here: there $\mathrm{AA}=\mathrm{TA}=\mathrm{DP}$.

\subsection*{Jets as a second test}

Jets are not on that list, and they test a class in another way. For DP manifolds $\mathcal M,\mathcal N$ with normed models, in every characteristic, the jet space $J^r(\mathcal M,\mathcal N)$ is a DP manifold, with DP forgetful maps, prolongations and functoriality (Theorem~\ref{dp:thm:dp-jets}). For analytic manifolds the question is governed by a section. We say that $(\mathrm{Sec}_k)$ holds for $(P,W)$ if the diagonal map from bounded $k$-linear maps $P^k\to W$ to laws $P\to W$ of degree $k$ has a bounded linear section. The \emph{operator family} of a chart change sends a point to the linear map on jets with values in $W$ that the chart change induces at that point. If $(\mathrm{Sec}_k)$ holds for $1\le k\le r$, the jet bundles of analytic manifolds with models $P,W$ have analytic transitions and operator families. The \emph{shear manifold} of a normed space $V$ is $V\times V\times K$ with the two global charts $\id$ and $(x,z,\lambda)\mapsto(x,z+\lambda x,\lambda)$. If on the shear manifold of $V$ the operator family of the jet bundle with values in $W$ is analytic at one point, then $(\mathrm{Sec}_k)$ holds for $(V,W)$ and $1\le k\le r$ (Theorem~\ref{dp:thm:analytic-jets}).

Let $K$ be complete of characteristic $p>0$ and $r\ge p$. The scalar jet bundles $J^r(\mathcal M;K)$ of all analytic manifolds $\mathcal M$ with ultrametric models have analytic transitions and operator families if and only if $K$ is spherically complete (Corollary~\ref{dp:cor:jet-dichotomy}). This separates analyticity from $\mathrm{AA}=\mathrm{DP}$. It does not separate AA from DP, as the family that fails is a map between ultrametric spaces, where AA and DP coincide. The separation from AA is Proposition~\ref{dp:prop:jets-not-ambient}(4), over every $K$ of characteristic $p$: for $r\ge p$, on the shear manifold of $\mathcal E$, with target the completion of $\Gamma^p(\mathcal E)$ for $\norm\cdot_\Gamma$, the operator family of the jet bundle with values in that target is DP at every point and misses $\cC_p$ at every point. So it is neither AA nor TA at any point.

\subsection*{The whole picture}

Between ultrametric spaces the classes above, apart from the analytic germs, are one class. The AA germs, the TA germs, the strongly DP germs of Section~\ref{dp:sec:strong}, the CA germs and the DP germs coincide there (\cite[Theorem~14.39(1)]{I}, Theorem~\ref{dp:thm:singles-out}(4), Lemmas~\ref{dp:lem:strong-formal}(3) and~\ref{dp:lem:cover-flat}(5)). The comparison with analyticity then turns on the field. For $k\ge1$, normed spaces $E,D,F,F'$, $u\in\cL(E,D)$, $v\in\cL(F,F')$ and $k$-linear maps $\mathsf m\colon D^k\to F$, let $B(u,v)(\mathsf m)=v\circ\mathsf m\circ(u,\dots,u)$ be the joint action, and let $A^k_{E,D;F}(u)=B(u,\id_F)$ on alternating $\mathsf m$. For Banach spaces this is the formula of~\cite[n\textsuperscript{o}~7.8.1]{Bou} for $\alpha_k(\id_F,u)$~\cite[Section~10.2]{I}. Bourbaki defines the functor $\alpha_k$ when $K$ has characteristic $0$ or the bundles have finite rank (Section~\ref{dp:sec:functors}); we use the formula for all normed spaces. Let $K$ be complete, and consider two statements.
\begin{itemize}
\item[$(\mathrm D_K)$] Every map between open subsets of ultrametric Banach spaces that is DP at a point is analytic there.
\item[$(\mathrm B_K)$] For all ultrametric Banach spaces $E,D,F,F'$ and every $k\ge1$, the maps $A^k_{E,D;F}$ and $B$, restricted to alternating maps, are analytic everywhere.
\end{itemize}

\begin{introtheorem}[agreement with analyticity]\label{dp:intro:D}
Let $K$ be complete.
\begin{enumerate}
\item $(\mathrm D_K)$ implies $(\mathrm B_K)$. Both hold if $\operatorname{char}K=0$ or if $K$ is discretely valued. Both fail if $\operatorname{char}K=p>0$ and $K$ is not spherically complete.
\item Let $P$ be ultrametric and $W$ an ultrametric Banach space. If $f$ is DP at $x_0$ and not analytic there, then $\operatorname{char}K=p>0$, $K$ is densely valued, the completion of $P$ is not linearly homeomorphic to any $c_0(\mathsf I,K)$, and $W$ has no equivalent ultrametric spherically complete norm.
\item Let $P$ be finite-dimensional and $W$ any normed space. A map from an open subset of $P$ to $W$ is DP at a point if and only if it is analytic there.
\end{enumerate}
\end{introtheorem}

This is Theorem~\ref{dp:thm:field-level} and Theorem~\ref{dp:thm:conservativity}(2),(5). So, wherever $(\mathrm B_K)$ is decided here, $(\mathrm D_K)$ is decided with the same answer. Between ultrametric Banach spaces, the disagreements listed here occur over the complete, not spherically complete $K$ of characteristic $p$, for some spaces, and there, for $k\ge p$ and some ultrametric Banach spaces $E,D$, the map $A^k_{E,D;K}$ is analytic at no point (Proposition~\ref{dp:prop:strict}(1)). One case is undecided. If $K$ is densely valued and spherically complete of characteristic $p$, we decide neither $(\mathrm D_K)$ nor $(\mathrm B_K)$. Whether $(\mathrm B_K)$ implies $(\mathrm D_K)$ is open (Section~\ref{dp:sec:questions}).

Section~\ref{dp:sec:ledger} lists the disagreements between DP and analyticity stated in this paper and in~\cite{I}, by field and by norms (Table~\ref{dp:tab:ledger}), and says what forces each one (Theorem~\ref{dp:thm:forcing}). We do not claim that the list is complete: in characteristic $p$, on ultrametric sources with non-ultrametric targets, whether DP and AA differ is open.

The strong norm appears both as a norm and as a competing class. For a law $a$ and vectors $y_1,\dots,y_l\in P$, expand $a(\sum_it_iy_i)$ as a polynomial in scalars $t_1,\dots,t_l$. The \emph{strong norm} of $a$ is the least constant $C$ such that the coefficient of each monomial $\prod_it_i^{d_i}$ has norm at most $C\prod_i\norm{y_i}^{d_i}$, for all $l$ and all $y_1,\dots,y_l$, and a germ is \emph{strongly DP} if its laws have geometric bounds in this norm. As a norm it loses no radius on ultrametric spaces (Theorems~\ref{dp:thm:strong-inverse} and~\ref{dp:thm:strong-recenter}). For $W$ an ultrametric Banach space and every $P$, a germ is strongly DP if and only if it is analytic after a cover, if and only if its laws descend to the ultrametric hull of $P$, the quotient of $P$ by the kernel of the seminorm $\mathsf u$ defined below, with geometric bounds there (Theorem~\ref{dp:thm:strong-um}). As a class it keeps (C2) and (C3), and it can lose (C1) (Theorem~\ref{dp:thm:strong-descent}). On $\ell^1(\N,K)$, with a non-ultrametric target, an analytic germ need not be strongly DP (Example~\ref{dp:ex:an-not-sdp}). On the same source, over every $K$ and in every characteristic, strongly DP germs have neither composition nor inverse functions (Theorem~\ref{dp:thm:strong-failure}).

We also have two results for all norms. If $W$ is complete, a map that is DP at $x_0$ with radius $R$ is DP at every point of the open ball of radius $R$ around $x_0$, with no hypothesis on $K$ or on the norms (Theorem~\ref{dp:thm:whole-ball}), though over some discretely valued $K$ its largest radius at a point $x_1$ of that ball can be smaller than $R-\norm{x_1-x_0}$ (Example~\ref{dp:ex:radius-drop}). Every law $a$ of degree $n$ extends to a law $a_{K'}$ on the base change to every valued extension $K'$ of $K$, with $\norm{a_{K'}}\le C_0^n\norm a$ for $C_0=8\mathrm e/\abs\varpi^2$ (Theorem~\ref{dp:thm:base-change-law}). Base change does not separate DP from analyticity by stability: both classes are stable, analytic germs with no loss of radius (Theorem~\ref{dp:thm:an-base-change}(1)), and DP germs with a possible loss of radius (Theorem~\ref{dp:thm:germ-base-change}(1)), which occurs with the maximum norm on the base change (Example~\ref{dp:ex:base-radius}). A DP germ whose base change is analytic or AA over $K'$ is already AA over $K$, since base change is a particular enlargement of the target (Theorems~\ref{dp:thm:an-base-change}(3) and~\ref{dp:thm:aa-base-change}(3)).

\subsection*{Organization and norms}

Section~\ref{dp:sec:conv} fixes conventions and recalls what we use from~\cite{I}. Sections~\ref{dp:sec:laws} and~\ref{dp:sec:germs} give the calculus of laws and germs. Section~\ref{dp:sec:nonambient} gives laws that are not ambient, with proofs in Appendix~\ref{dp:app:proofs}. Section~\ref{dp:sec:descent} gives descent and minimality, and Section~\ref{dp:sec:chain} the classes between AA and DP. Sections~\ref{dp:sec:forms}, \ref{dp:sec:quotients} and~\ref{dp:sec:base-change} treat forms, difference quotients and base change. Section~\ref{dp:sec:agreement} compares DP with analyticity, Section~\ref{dp:sec:functors} treats vector functors and n\textsuperscript{o}~7.8 of~\cite{Bou}, and Section~\ref{dp:sec:jets} treats jets. Sections~\ref{dp:sec:strong} and~\ref{dp:sec:cover} study the strong norm and CA, Section~\ref{dp:sec:ledger} is the list of disagreements, and Section~\ref{dp:sec:questions} collects questions.

We work with general norms as far as they extend naturally. The mixed-term bound, the whole-ball theorem, composition, inverse functions, descents, bundles, forms, jets and base change need no ultrametric hypothesis. Ultrametric norms enter where they are used, and we say so there. The main places and reasons are these.
\begin{itemize}
\item Theorem~\ref{dp:thm:mixed}(4). On an ultrametric source the mixed-term bound loses no factor $n^n$. With an ultrametric target as well, Lemma~\ref{dp:lem:strong-coefficients} gives the constants.
\item Theorem~\ref{dp:thm:ultrametric-minimal}. Between ultrametric spaces, source descent reaches DP without target descent, since laws on the cover are diagonals with geometric constants.
\item Sections~\ref{dp:sec:agreement} and~\ref{dp:sec:strong}, on targets. In an ultrametric target a sum is bounded by its maximum. So multilinear maps into it see only the hull seminorm $\mathsf u(x)=\inf\max_i\norm{x_i}$ of the source, the infimum over the decompositions $x=\sum_ix_i$, and a coefficient, a sum with one term for each ordering, is no larger than its largest term.
\end{itemize}
Every theorem is proved here or cited from~\cite{I} or from the bibliography.
%% ---- end sec-01-intro.tex
%% ---- begin sec-02-conventions.tex
\section{Conventions and inputs from \texorpdfstring{\cite{I}}{[I]}}\label{dp:sec:conv}

This section fixes notation and recalls the inputs from~\cite{I} used throughout. Other results of~\cite{I} are cited where they are used.

\emph{Standing assumptions.} Throughout, $K$ is a field with a nontrivial nonarchimedean absolute value $\abs\cdot$, with completion $\Kh$, and $\varpi\in K$ is fixed with $0<\abs\varpi<1$. The powers of $\varpi$ have distinct absolute values, so $K$ is infinite. If $K$ is discretely valued, $\vartheta\in K$ is fixed with $\abs{K^\times}=\abs\vartheta^{\Z}$ and $\abs\vartheta<1$. Put $\mathfrak c_K=\abs\vartheta$ in that case, and $\mathfrak c_K=1$ otherwise.

The letters $P,P'$ denote source spaces and $W,Y,Z$ target spaces. All are normed $K$-spaces, not necessarily complete or ultrametric. A normed space is \emph{ultrametric} if $\norm{x+y}\le\max(\norm x,\norm y)$ for all its vectors $x,y$. The completion of $W$ is $\widehat W$. Products and finite direct sums carry the maximum norm. The open ball of radius $R$ around $x$ is $B(x,R)$. The \emph{$\ell^1$-sum} of normed spaces $Y_i$ is the space of the families $(y_i)_i$, $y_i\in Y_i$, with $\sum_i\norm{y_i}<\infty$, and this sum as norm~\cite[Definition~14.15]{I}. For a set $J$ let $\ell^1(J)=\ell^1(J,K)$ be the $\ell^1$-sum of copies of $K$ indexed by $J$, and let $\ell^1_{00}(J)$ be its subspace of finitely supported vectors.

\emph{Tensors.} For $n\ge1$, $P^{\otimes n}$ is the algebraic tensor power. For $t\in P^{\otimes n}$ put
\[
 \pi(t)=\inf\sum_l\prod_{i=1}^n\norm{h_{l,i}},\qquad
 \pi_{\max}(t)=\inf\max_l\prod_{i=1}^n\norm{h_{l,i}},
\]
both infima over the representations $t=\sum_lh_{l,1}\otimes\dots\otimes h_{l,n}$. Here $\pi$ is the projective seminorm and $\pi_{\max}$ is an ultrametric seminorm with $\pi_{\max}\le\pi$. For $n=0$ we read $P^{\otimes0}=K$, $\pi=\pi_{\max}=\abs\cdot$, $h^{\otimes0}=1$, and $\norm h^0=1$ also at $h=0$.

The space $\Gamma^n(P)\subseteq P^{\otimes n}$ is the span of the powers $h^{\otimes n}$, $h\in P$. As $K$ is infinite, it is the space of symmetric tensors~\cite[Section~L.3, opening paragraph]{I}. We write $\Gamma^n_\pi(P)$ for $\Gamma^n(P)$ with the restriction of $\pi$. On $\Gamma^n(P)$ we also use the seminorm
\begin{equation}\label{dp:eq:gamma-norm}
 \norm t_\Gamma=\inf\Bigl\{\sum_i\abs{c_i}\norm{h_i}^n:\ t=\sum_ic_i\,h_i^{\otimes n}\Bigr\},
\end{equation}
so that $\pi\le\norm\cdot_\Gamma$. If the bounded linear forms on $P$ separate points, then $\pi$ is a norm on $P^{\otimes n}$ and $\norm\cdot_\Gamma$ is a norm on $\Gamma^n(P)$. Indeed, the products of such forms are bounded for $\pi$, and they separate the points of $P^{\otimes n}$, as one checks on a finite-dimensional subspace containing the factors of a given tensor.

For $h_1,\dots,h_r\in P$ and integers $\alpha_i\ge0$ with $\abs\alpha=\alpha_1+\dots+\alpha_r=n$, the \emph{orbit sum} $O_\alpha(h_1,\dots,h_r)$ is the sum of the tensors $h_{i_1}\otimes\dots\otimes h_{i_n}$ over the tuples $(i_1,\dots,i_n)$ in which each $i$ occurs $\alpha_i$ times. With $\lambda^\alpha=\prod_i\lambda_i^{\alpha_i}$, expanding the power gives~\cite[(14.1)]{I}
\begin{equation}\label{dp:eq:orbit}
 \Bigl(\sum_{i=1}^r\lambda_ih_i\Bigr)^{\otimes n}=\sum_{\abs\alpha=n}\lambda^\alpha\,O_\alpha(h_1,\dots,h_r)\qquad(\lambda\in K^r),
\end{equation}
and $O_\alpha(h_1,\dots,h_r)\in\Gamma^n(P)$, since each coefficient of a polynomial with values in a subspace lies in that subspace, by interpolation.

For integers $n,n'\ge0$, $s\in P^{\otimes n}$ and $t\in P^{\otimes n'}$, the \emph{shuffle product} $s\cdot t\in P^{\otimes(n+n')}$ is the sum of the $\binom{n+n'}n$ tensors obtained from $s\otimes t$ by the permutations of the $n+n'$ factors that keep the order of the first $n$ factors and the order of the last $n'$. For $n=0$ or $n'=0$ the tensor $s\otimes t$ is the scalar multiple. The product is bilinear. Permutations of the factors preserve $\pi$, and $\pi(s\otimes t)\le\pi(s)\pi(t)$, so $\pi(s\cdot t)\le\binom{n+n'}n\pi(s)\pi(t)$. For $h,h'\in P$ we have $h^{\otimes n}\cdot h'^{\otimes n'}=O_{(n,n')}(h,h')$. So, by bilinearity and since orbit sums lie in $\Gamma^{n+n'}(P)$, the product maps $\Gamma^n(P)\times\Gamma^{n'}(P)$ into $\Gamma^{n+n'}(P)$. In the completed projective tensor powers it induces the shuffle product $*$ of~\cite[Section~L.3]{I}.

\emph{Enlarging the target.} Bounded polynomial laws, or \emph{laws}, their norm $\norm a$, maps ambient of degree $n$, and DP germs with radius, bound and coefficient laws are as in Definition~\ref{dp:def:laws}, which repeats~\cite[Definition~14.8]{I}. A law $a\colon P\to W$ of degree $n$ has a unique linear map $L_a\colon\Gamma^n(P)\to W$ with $L_a(h^{\otimes n})=a(h)$ for $h\in P$~\cite[Lemma~Q.1]{I}. Its \emph{ambient constant} is $\operatorname{amb}(a)=\sup\{\norm{L_a(t)}/\pi(t):t\in\Gamma^n(P)\}\in[0,\infty]$, where a quotient with denominator $0$ is read as $\infty$ if its numerator is positive and as $0$ if its numerator is $0$. The law $a$ is ambient of degree $n$ if and only if $\operatorname{amb}(a)<\infty$, and $\norm a\le\operatorname{amb}(a)$~\cite[text after Definition~14.8]{I}.

Let $U\subseteq P$ be open, $f\colon U\to W$, $x_0\in U$, $R>0$ with $B(x_0,R)\subseteq U$, and $M\ge0$. By~\cite[Definition~14.7]{I}, $f$ is \emph{ambient analytic}, or \emph{AA}, at $x_0$ with radius $R$ and bound $M$ if there are a linear isometry $j\colon W\to Z$ into a normed space and bounded $n$-linear maps $\beta_n\colon P^n\to Z$, $n\ge0$, with $\norm{\beta_n}\le MR^{-n}$, with $\beta_n(h,\dots,h)\in j(W)$ for all $h\in P$, and with $jf(x_0+h)=\sum_n\beta_n(h,\dots,h)$ for $\norm h<R$. It is AA at $x_0$ if this holds for some $R,M$, and AA on $U$ if it is AA at every point of $U$. For $Z=W$ and $j=\id_W$ we say that $f$ is \emph{analytic} at $x_0$ with radius $R$ and bound $M$, analytic at $x_0$ with radius $R$ if this holds for some $M$, and analytic at $x_0$ if it holds for some $R,M$. In~\cite{I} the bound $M$ and the maps $\beta_n$ are written $\mathfrak b$ and $b_n$.

By~\cite[Proposition~14.9]{I}: (1) $f$ is AA at $x_0$ with radius $R$ and bound $M$ if and only if it is DP at $x_0$ with radius $R$ and its coefficient laws $a_n$ satisfy $\operatorname{amb}(a_n)\le MR^{-n}$ for all $n$; it is then DP with bound $M$, and $ja_n(h)=\beta_n(h,\dots,h)$ for every $j$ and $(\beta_n)_n$ as above with this radius and any bound. (2) A law $a$ of degree $n$ such that $j\circ a$ is analytic at one point, for a linear isometry $j$ of $W$ into a normed space, is ambient of degree $n$.

Composition fails~\cite[Theorem~14.17]{I}. Let $\operatorname{char}K=p>0$, with $K$ complete or not, and let $f_X\colon H^{\rm at}\to W^{\rm at}$ and $g_X\colon W^{\rm at}\to Y^{\rm at}$ be the maps of~\cite[Definition~14.16]{I}, with $X=X^{\rm at}$ there. Then $f_X$ is AA at every point, $g_X$ is the diagonal of a bounded bilinear map, and $g_X\circ f_X$ is a law of degree $p+1$, DP everywhere and AA at no point. Whether composition fails in characteristic $0$ is open~\cite[Remark~14.20]{I}.

\emph{Tame pairs.} For $n\ge1$ let $\mathfrak a_n(P,W)$ be the supremum of $\operatorname{amb}(a)/\norm a$ over the nonzero laws $a\colon P\to W$ of degree $n$, and $0$ if there are none. The pair $(P,W)$ is \emph{tame} if $\sup_{n\ge1}\mathfrak a_n(P,W)^{1/n}<\infty$, and $P$ is \emph{tame} if $(P,W)$ is tame for every normed $W$. Put $\mathfrak a_n(P)=\sup\{\norm t_\Gamma/\pi(t):t\in\Gamma^n(P)\}$, with the same reading of quotients~\cite[Definition~14.36]{I}. Then $\mathfrak a_n(P,W)\le\mathfrak a_n(P)$, and $P$ is tame if and only if $\sup_{n\ge1}\mathfrak a_n(P)^{1/n}<\infty$~\cite[text after Definition~14.36]{I}.

For $W$ Banach, $(P,W)$ is tame if and only if every map from an open subset of $P$ into $W$ that is DP at a point is AA there~\cite[Proposition~14.37]{I}. By~\cite[Theorem~14.38]{I}: (1) if $\mathfrak a_n(P,W)\le\mathfrak q_0\mathfrak q^n$ for $n\ge1$, with $\mathfrak q_0,\mathfrak q\ge1$, then a germ from $P$ to $W$ that is DP at $x_0$ with radius $R$ and bound $M$ is AA there with radius $R/\mathfrak q$ and bound $\mathfrak q_0M$; (2) if $(P,W)$ is tame, $f$ from an open subset of $P$ into a normed space $Y$ is DP at $x_0$, and $g$ from an open subset of $Y$ into $W$ is DP at $f(x_0)$, then $g\circ f$ is AA at $x_0$; (3) if $P$ is Banach, $(P,P)$ is tame, and $f$ from an open subset of $P$ into $P$ is DP at $x_0$, with coefficient law $a_1$ bijective with bounded inverse, then $f$ has a local inverse germ, which is AA at $f(x_0)$ and at every point of a ball around it.

The examples are in~\cite[Theorem~14.39]{I}. (1) If $P,W$ are ultrametric, then $\mathfrak a_n(P,W)\le(\abs\varpi\mathfrak c_K)^{-n}$. (2) If $P=\ell^1(J)$ or $P=\ell^1_{00}(J)$, then $\mathfrak a_n(P)\le(8\mathrm e/\abs\varpi^2)^n$, with $\mathrm e$ Euler's number. (3) If $\operatorname{char}K=0$ and $P$ is ultrametric, then $\sup_{n\ge1}\mathfrak a_n(P)^{1/n}<\infty$, so $P$ is tame. (4) If $W$ is ultrametric and $P$ is any normed space, then every law from $P$ to $W$ is ambient, and $\mathfrak a_n(P,W)<\infty$ in each degree; whether $(P,W)$ is always tame is open. Between ultrametric spaces the AA germs and the DP germs coincide~\cite[Theorem~14.39(1)]{I}.

\emph{Alternating maps.} For normed spaces $E,D,F,F'$ and $k\ge1$, $\cL(E,D)$ is the normed space of bounded linear maps $E\to D$ with the operator norm, $\Mult^k(E;F)$ is the normed space of bounded $k$-linear maps $E^k\to F$, and $\Alt^k(E;F)$ is its subspace of alternating maps, those vanishing whenever two arguments agree. For $k=0$ both spaces are $F$. The inclusion $j_{E,F}\colon\Alt^k(E;F)\to\Mult^k(E;F)$ is a linear isometry with closed range~\cite[Section~2]{I}. The alternating degree $k$ is Bourbaki's $d$ in~\cite[n\textsuperscript{o}~7.8]{Bou}. Following~\cite[Section~2]{I}, for $u\in\cL(E,D)$, $v\in\cL(F,F')$ and $\beta\in\Mult^k(D;F)$ the \emph{joint action} at $z_1,\dots,z_k\in E$ is
\[
 B^k_{E,D;F,F'}(u,v)(\beta)(z_1,\dots,z_k)=v\bigl(\beta(uz_1,\dots,uz_k)\bigr),
\]
written $B$ in~\cite{I}. In the later sections $B(u,v)$, with operators as arguments, is this action, and $B(x,R)$, with a point and a radius, is the ball. The precomposition is the map $A^k_{E,D;F}$ from $\cL(E,D)$ to $\cL(\Alt^k(D;F),\Alt^k(E;F))$ that sends $u$ to the restriction of $B^k_{E,D;F,F}(u,\id_F)$ to alternating maps. By~\cite[Proposition~2.2]{I}, the joint action is the diagonal of a bounded $(k+1)$-linear map, and $A^k_{E,D;F}$ followed by $j_{E,F}\circ{-}$ is the diagonal of a bounded $k$-linear map of norm at most $1$. The projective exterior power $\Lambda^k_\pi E$ is the normed quotient of the algebraic exterior power by the null space of the seminorm built from wedges as $\pi$ is built from tensors, $\widehat\Lambda^k_\pi E$ is its completion, and $\Alt^k(E;F)\cong\cL(\Lambda^k_\pi E,F)$ isometrically, with precomposition by $u$ becoming precomposition by $\Lambda^k_\pi u$~\cite[Section~3.1, Proposition~3.1]{I}.

\emph{Bundles.} Let $\mathsf{Vec}_K$ be the category of normed $K$-spaces and bounded linear maps, with hom spaces $\cL(E,D)$. Let $m\in\N\cup\{\infty,\omega\}$; this is the index written $n$ in~\cite{I}. The class $C^m$ is as in the conventions of~\cite[Section~1]{I}, and $C^0$ means continuous. Fix a $C^m$ manifold $\mathcal M$ modeled on $P$. Then $\mathsf{VB}^m_K(\mathcal M)$ is the category of $C^m$ normed vector bundles over $\mathcal M$, with morphisms over $\id_{\mathcal M}$ whose local operator-valued expressions $x\mapsto T(x)\in\cL(E,F)$ are $C^m$. Let $\epsilon=(\epsilon_1,\dots,\epsilon_q)$ specify a variance in each variable. Put $\mathbf C^\epsilon=\mathsf C_1\times\dots\times\mathsf C_q$, with $\mathsf C_i=\mathsf{Vec}_K$ or $\mathsf{Vec}_K^{\mathrm{op}}$, and let $\mathbf B^\epsilon_{\mathcal M}$ be the corresponding product of bundle categories, with the same variances. A functor $\mathcal F\colon\mathbf C^\epsilon\to\mathsf{Vec}_K$ is $C^m$ if all its maps on products of hom spaces are jointly $C^m$. Let $\operatorname{Fun}^{C^m}(\mathbf C^\epsilon,\mathsf{Vec}_K)$ be the full subcategory of the functor category on these functors, and write $\eta$ for a natural transformation between two of them. The statement of~\cite[Theorem~2.1]{I} is this.
\begin{quote}
For a fixed $C^m$ manifold $\mathcal M$, fiberwise application defines, up to canonical natural isomorphism, a functor
\[
\begin{aligned}
 \mathcal B^m_{\mathcal M}\colon
 \operatorname{Fun}^{C^m}(\mathbf C^\epsilon,\mathsf{Vec}_K)
 &\longrightarrow
 \operatorname{Fun}(\mathbf B^\epsilon_{\mathcal M},\mathsf{VB}^m_K(\mathcal M)),\\
 \mathcal F&\longmapsto\mathcal F_{\mathcal M},\qquad
 \eta\longmapsto\eta_{\mathcal M}.
\end{aligned}
\]
The resulting bundles, bundle morphisms and natural transformations are given fiberwise by $\mathcal F$, its action on maps, and the components of $\eta$, respectively. The same conclusion holds if regularity on whole hom spaces is replaced by preservation of $C^m$ families parametrized by open subsets of the model space $P$.
\end{quote}
Preservation of $C^m$ families means that $\mathcal F\circ\gamma$ is $C^m$ for each $C^m$ map $\gamma$ from an open subset of $P$ into a product of hom spaces of $\mathbf C^\epsilon$. The theorem also holds for restricted categories of fibers containing the bundle morphisms~\cite[text after Theorem~2.1]{I}.

The proof in~\cite{I} uses of the class $C^m$ only these four properties. (i) Maps of the class on open sets compose, tuples of them into products are in the class, and the evaluation of operators is in the class. (ii) If $x\mapsto S(x)$ and $x\mapsto T(x)$ are local operator families of the class, so is $x\mapsto S(x)T(x)$. (iii) Constants are in the class. (iv) The class passes to restrictions to open subsets of $P$, and membership depends only on the map near each point. Step~1 there composes the transitions with the maps of the functor on hom spaces and evaluates operators, by (i), and glues chart by chart, by (iv). Step~2 applies the maps of the functor to local operator families, by (i), multiplies by operator families for changes of trivialization, by (ii), and works chart by chart, by (iv). Step~3 uses (iii), and Step~4 uses (iv). The category $\mathsf{VB}^m_K(\mathcal M)$ itself uses (ii) for the composition of morphisms and (iii) for identities. Corollary~\ref{dp:cor:bundles} applies this to DP maps.
%% ---- end sec-02-conventions.tex
%% ---- begin sec-03-laws.tex
\section{Bounded polynomial laws}\label{dp:sec:laws}

This section defines laws and divided-power germs and collects the estimates on coefficients used in the rest of the paper. The algebra of laws is Roby's~\cite{Ro}; the bounds below are the estimates we need. We keep the conventions of Section~\ref{dp:sec:conv}: $P$, $P'$, $W$ and $Y$ are normed spaces, not necessarily complete or ultrametric. We fix the constants
\[
 C_1=8/\abs\varpi^2,\qquad C_2=16/\abs\varpi^2,\qquad C_0=\mathrm eC_1=8\mathrm e/\abs\varpi^2.
\]

\subsection{Laws and germs}

\begin{definition}\label{dp:def:laws}
Let $n\ge0$.
\begin{enumerate}
\item A map $a\colon P\to W$ is a \emph{bounded polynomial law of degree $n$}, or a \emph{law}, if, for all $h_1,\dots,h_r\in P$, the map $s\mapsto a(s_1h_1+\dots+s_rh_r)$ on $K^r$ is a homogeneous polynomial of degree $n$ with coefficients in $W$, and $\norm a=\sup_{h\ne0}\norm{a(h)}/\norm h^n$ is finite; a law of degree $0$ is constant, and then $\norm a=\norm{a(0)}$. For $\alpha=(\alpha_1,\dots,\alpha_r)$ with integers $\alpha_c\ge0$ and $\abs\alpha=\alpha_1+\dots+\alpha_r=n$, $a^{[\alpha]}(h_1,\dots,h_r)\in W$ is the coefficient of $s_1^{\alpha_1}\cdots s_r^{\alpha_r}$, unique since $K$ is infinite.
\item A map $a\colon P\to W$ is \emph{ambient of degree $n$} if $ja(h)=\beta(h,\dots,h)$ for all $h\in P$, for a linear isometry $j$ of $W$ into a normed space $Z$ and a bounded $n$\nobreakdash-linear map $\beta\colon P^n\to Z$.
\item Let $U\subseteq P$ be open, $f\colon U\to W$, $x_0\in U$, $R>0$ with $B(x_0,R)\subseteq U$, and $M\ge0$. The map $f$ is \emph{divided-power analytic}, or \emph{DP}, at $x_0$ with radius $R$ and bound $M$ if $f(x_0+h)=\sum_na_n(h)$ for $\norm h<R$, with laws $a_n\colon P\to W$ of degree $n$ and $\norm{a_n}\le MR^{-n}$. It is DP at $x_0$ with radius $R$ if this holds for some $M$, DP at $x_0$ if it holds for some $R,M$, and DP on $U$ if it is DP at every point of $U$. The $a_n$ are the \emph{coefficient laws} of $f$ at $x_0$. By \cite[Lemma~A.1]{I} on the lines $s\mapsto x_0+sh$, they depend only on $f$.
\end{enumerate}
\end{definition}

The constancy in Definition~\ref{dp:def:laws}(1) follows from $r=1$, comparing $s=0$ with $s=1$. The space $\cP_n(P;W)$ is the vector space of laws $P\to W$ of degree $n$ with the norm $\norm a$. Sums of laws are laws, as the polynomials add. Here $\N=\{0,1,2,\dots\}$ and $[n]=\{1,\dots,n\}$. We write $\alpha!=\prod_c\alpha_c!$, $s^\gamma=\prod_cs_c^{\gamma_c}$, and $1^k$ for a run of $k$ ones. So $(n-k,1^k)=(n-k,1,\dots,1)$. The space $\Mult^k(P;W)$ consists of the bounded $k$\nobreakdash-linear maps, with $\Mult^0(P;W)=W$, and $\cL(P,\Mult^{k-1}(P;W))=\Mult^k(P;W)$ isometrically for $k\ge1$.

Without the bound on $\norm a$, the maps in Definition~\ref{dp:def:laws}(1) are the maps $P\to W$ induced by the polynomial laws on $(P,W)$ in the sense of Roby~\cite{Ro} that are homogeneous of degree $n$, and each of these maps determines its law, since $K$ is an infinite field and the linear forms on the vector space $W$ separate its points~\cite[Th.~I.1, Props.~I.1, I.8, IV.6, IV.7]{Ro}. The $a^{[\alpha]}$ are the values on $P$ of Roby's polarizations~\cite[II~\S1]{Ro}.

Recall the seminorm $\norm\cdot_\Gamma$ of~\eqref{dp:eq:gamma-norm}. Let $\Gamma^n_{\mathrm s}(P)$ be the quotient of $\Gamma^n(P)$ by the kernel of $\norm\cdot_\Gamma$, normed by $\norm\cdot_\Gamma$, let $\hGamma^n(P)$ be its completion, and let $\mathsf v_n\colon P\to\Gamma^n_{\mathrm s}(P)$ be the map $\mathsf v_n(h)=[h^{\otimes n}]$.

\begin{theorem}[laws as linear maps]\label{dp:thm:normed-roby}
Let $n\ge0$.
\begin{enumerate}
\item For every $a\in\cP_n(P;W)$ there is a unique linear map $L_a\colon\Gamma^n(P)\to W$ with $L_a(h^{\otimes n})=a(h)$ for $h\in P$, and $\norm{L_a(t)}\le\norm a\,\norm t_\Gamma$.
\item If $\mathsf L\colon\Gamma^n(P)\to W$ is linear with $\norm{\mathsf L(t)}\le C\norm t_\Gamma$, then $h\mapsto\mathsf L(h^{\otimes n})$ is a law of degree $n$ and norm at most $C$.
\item By item (1), $L_a$ vanishes on the kernel of $\norm\cdot_\Gamma$, and we also write $L_a$ for the induced map on $\Gamma^n_{\mathrm s}(P)$. The map $a\mapsto L_a$ is an isometric linear isomorphism from $\cP_n(P;W)$ onto $\cL(\Gamma^n_{\mathrm s}(P),W)$. If $W$ is Banach, this space is $\cL(\hGamma^n(P),W)$. Equivalently, $\mathsf v_n$ is a universal law of degree $n$ and norm at most $1$: every law factors uniquely as $a=L_a\circ\mathsf v_n$, with $\norm{L_a}=\norm a$.
\item The seminorm $\norm\cdot_\Gamma$ is the largest seminorm $\sigma$ on $\Gamma^n(P)$ with $\sigma(h^{\otimes n})\le\norm h^n$ for all $h\in P$.
\end{enumerate}
\end{theorem}

Without the bounds, items (1)--(3) are~\cite[Th.~IV.1, Prop.~IV.5]{Ro}. The bound is the estimate we need.

\begin{proof}
\emph{Step 1 (items (1) and (2): citation).} These are \cite[Lemma~Q.1(1),(2)]{I}.

\emph{Step 2 (item (3): the isometry).} A linear map bounded for $\norm\cdot_\Gamma$ vanishes on the kernel of $\norm\cdot_\Gamma$, so it is the same as a bounded linear map on $\Gamma^n_{\mathrm s}(P)$. We use the same letter for it. By item (2), applied to the quotient map, $\mathsf v_n$ is a law of norm at most $1$. If $W$ is Banach, a bounded linear map on $\Gamma^n_{\mathrm s}(P)$ extends uniquely to the completion with the same norm. The extension is linear over $K$, and over $\Kh$ for the $\Kh$\nobreakdash-structure that the complete space $W$ carries by continuity. By items (1) and (2), $a\mapsto L_a$ and $\mathsf L\mapsto\mathsf L\circ\mathsf v_n$ are mutually inverse, and item (1) gives $\norm{L_a}\le\norm a$. Since $\norm{h^{\otimes n}}_\Gamma\le\norm h^n$ by~\eqref{dp:eq:gamma-norm}, also $\norm{a(h)}=\norm{L_a(h^{\otimes n})}\le\norm{L_a}\norm h^n$, and for $n=0$ this reads $\norm{a(0)}\le\norm{L_a}$. So $\norm a\le\norm{L_a}$.

\emph{Step 3 (item (4): the largest seminorm).} We have $\norm{h^{\otimes n}}_\Gamma\le\norm h^n$. If $\sigma$ is a seminorm with $\sigma(h^{\otimes n})\le\norm h^n$, then $\sigma(t)\le\sum_l\abs{\lambda_l}\norm{h_l}^n$ for every representation $t=\sum_l\lambda_lh_l^{\otimes n}$. So $\sigma\le\norm\cdot_\Gamma$.
\end{proof}

\begin{lemma}[elementary facts]\label{dp:lem:elementary}
Let $a\in\cP_n(P;W)$.
\begin{enumerate}
\item A map $P\to W$ is a law of degree $1$ if and only if it is bounded linear, with the same norm. So $\cP_1(P;W)=\cL(P,W)$ isometrically.
\item If $T_2\in\cL(P',P)$ and $T_1\in\cL(W,Y)$, then $T_1\circ a\circ T_2\in\cP_n(P';Y)$, with $\norm{T_1\circ a\circ T_2}\le\norm{T_1}\,\norm a\,\norm{T_2}^n$.
\item For $w_1,\dots,w_r\in P$, scalars $\lambda_c$ and $\abs\alpha=n$, we have
\[
 a^{[\alpha]}(\lambda_1w_1,\dots,\lambda_rw_r)=\prod_c\lambda_c^{\alpha_c}\,a^{[\alpha]}(w_1,\dots,w_r).
\]
Arguments $w_c$ with $\alpha_c=0$ may be deleted. Also $a^{[\alpha]}(w_1,\dots,w_r)=0$ if $w_c=0$ for some $c$ with $\alpha_c\ge1$, and $a(w_1+\dots+w_r)=\sum_{\abs\alpha=n}a^{[\alpha]}(w_1,\dots,w_r)$.
\item (Partial coefficients are laws.) Let $h_1,\dots,h_r\in P$, let $i,k\ge0$, and let $\gamma\in\N^r$ with $\abs\gamma=n-i$, respectively $\abs\gamma=n-k$. The map $y\mapsto a^{[(\gamma,i)]}(h_1,\dots,h_r,y)$ lies in $\cP_i(P;W)$, and for $y_1,\dots,y_{r'}\in P$ and $\delta\in\N^{r'}$ with $\abs\delta=i$ its coefficients are
\begin{equation}\label{dp:laws:partial}
 \bigl(a^{[(\gamma,i)]}(h_1,\dots,h_r,\cdot\,)\bigr)^{[\delta]}(y_1,\dots,y_{r'})=a^{[(\gamma,\delta)]}(h_1,\dots,h_r,y_1,\dots,y_{r'}).
\end{equation}
The map sending $(y_1,\dots,y_k)$ to $a^{[(\gamma,1^k)]}(h_1,\dots,h_r,y_1,\dots,y_k)$ is bounded, symmetric and $k$\nobreakdash-linear.
\item (Uniqueness.) If $\sum_na_n(h)=\sum_na_n'(h)$ for $\norm h<R$, with $a_n,a_n'\in\cP_n(P;W)$ and $\sup_n(\norm{a_n}+\norm{a_n'})R^n<\infty$, then $a_n=a_n'$ for all $n$.
\end{enumerate}
\end{lemma}

\begin{proof}
\emph{Step 1 (items (1)--(3): substitutions).} For (1), let $a$ have degree $1$. Then $a(s_1h_1+s_2h_2)=s_1\mathsf w_1+s_2\mathsf w_2$ with $\mathsf w_1,\mathsf w_2\in W$, and $s=(1,0),(0,1)$ give $\mathsf w_1=a(h_1)$, $\mathsf w_2=a(h_2)$. So $a$ is additive and homogeneous. The converse is clear. For (2), the polynomial $s\mapsto T_1a(\sum_cs_cT_2h_c)$ has as coefficients the images under $T_1$ of those of $a$ at the vectors $T_2h_c$. For (3), substitute $s_c\mapsto\lambda_cs_c$, take $\lambda_c=0$ for the vanishing statement, set $s_c=0$ for the variables that do not occur to delete arguments, and set $s=(1,\dots,1)$ for the last formula.

\emph{Step 2 (item (4): the coefficients of a partial law).} Let $y=\sum_{l=1}^{r'}\tau_ly_l$ with $\tau_l\in K$. Expand $a(\sum_cs_ch_c+z\sum_l\tau_ly_l)$ in the variables $s_c,z$ in two ways. In the vectors $h_1,\dots,h_r,y$, the coefficient of $s^\gamma z^i$ is $a^{[(\gamma,i)]}(h_1,\dots,h_r,y)$. In the vectors $h_1,\dots,h_r,y_1,\dots,y_{r'}$, with the scalars $s_c$ and $z\tau_l$, the same coefficient is $\sum_{\abs\delta=i}a^{[(\gamma,\delta)]}(h_1,\dots,h_r,y_1,\dots,y_{r'})\,\tau^\delta$. This is a homogeneous polynomial of degree $i$ in $\tau$ with the coefficients~\eqref{dp:laws:partial}. Its value at $y$ has norm at most $C_0^n\norm a\,\frac{n!}{\gamma!\,i!}\prod_c\norm{h_c}^{\gamma_c}\,\norm y^i$, by~\cite[(Q.4)]{I}. For the second assertion, apply the first in each slot $y_l$ separately with $i=1$, and use item (1). Symmetry holds because permuting $y_1,\dots,y_k$ permutes the variables attached to them. Boundedness is~\cite[(Q.4)]{I} again.

\emph{Step 3 (item (5): one line at a time).} Let $h\ne0$ and $\mathsf w_n=a_n(h)-a_n'(h)$. Then $\sum_n\lambda^n\mathsf w_n$ converges to $0$ for $\abs\lambda<R/\norm h$, and $\norm{\mathsf w_n}(R/\norm h)^n\le(\norm{a_n}+\norm{a_n'})R^n$ is bounded. So $\mathsf w_n=0$ by \cite[Lemma~A.1]{I}. The case $h=0$ is clear.
\end{proof}

\subsection{Mixed terms}

The next theorem is the main estimate. Items (1)--(3) hold on every normed space and are proved in~\cite{I}. On ultrametric sources, item (4) removes the factor $m^m\prod_c\alpha_c^{-\alpha_c}$ of item (3). The reason is that the points at which the polynomial is interpolated stay in the ball of radius $\max_c\norm{w_c}$.

\begin{theorem}[mixed terms]\label{dp:thm:mixed}
Let $a\colon P\to W$ be a law of degree $m$, let $w_1,\dots,w_r\in P$, and let $\alpha\in\N^r$ with $\abs\alpha=m$. We read $0^0=1$.
\begin{enumerate}
\item We have $\norm{a^{[\alpha]}(w_1,\dots,w_r)}\le(8/\abs\varpi)^m\norm a\,(\norm{w_1}+\dots+\norm{w_r})^m$.
\item For $x,y\in P$ and $0\le i\le m$, we have $\norm{a^{[(m-i,i)]}(x,y)}\le C_2^m\norm a\,\norm x^{m-i}\norm y^i$. So $a(x+y)=\sum_{i\le m}a^{[(m-i,i)]}(x,y)$ is a finite expansion at $x$ by laws in $y$, and $a$ is continuous and DP on $P$.
\item We have
\begin{equation}\label{dp:laws:scaled}
 \norm{a^{[\alpha]}(w_1,\dots,w_r)}\le C_1^m\norm a\,\frac{m^m}{\prod_c\alpha_c^{\alpha_c}}\prod_c\norm{w_c}^{\alpha_c}\le C_0^m\norm a\,\frac{m!}{\alpha!}\prod_c\norm{w_c}^{\alpha_c}.
\end{equation}
\item If $P$ is ultrametric, then, for every normed $W$,
\begin{equation}\label{dp:laws:ultrametric}
 \norm{a^{[\alpha]}(w_1,\dots,w_r)}\le C_1^m\norm a\prod_c\norm{w_c}^{\alpha_c}.
\end{equation}
\end{enumerate}
\end{theorem}

\begin{proof}
\emph{Step 1 (items (1)--(3): citation and the scaled bound).} Items (1) and (2) and the outer inequality of~\eqref{dp:laws:scaled} are~\cite[Lemma~Q.2(1),(2),(3)]{I}. The first inequality of~\eqref{dp:laws:scaled} is Step~4 of its proof, as follows. If $m=0$, then $a$ is constant. Otherwise delete the $w_c$ with $\alpha_c=0$; if a remaining $w_c$ vanishes, both sides vanish. Choose $\lambda_c\in K$ with $\abs\varpi\alpha_c/m<\abs{\lambda_c}\norm{w_c}\le\alpha_c/m$. Then $\sum_c\abs{\lambda_c}\norm{w_c}\le1$, so $\norm{a^{[\alpha]}(\lambda_1w_1,\dots,\lambda_rw_r)}\le(8/\abs\varpi)^m\norm a$ by item (1). Undoing the scaling with Lemma~\ref{dp:lem:elementary}(3) gives the first inequality. The second follows from $m^m\le\mathrm e^mm!$ and $\alpha_c^{\alpha_c}\ge\alpha_c!$.

\emph{Step 2 (nodes).} Let $\varpi'\in K$ with $0<\abs{\varpi'}<1$. For an integer $\ell\ge0$ with binary digits $\ell_u$, put $\xi_\ell=\sum_u\ell_u\varpi'^u\in K$. Then $\abs{\xi_\ell}\le1$. For $\ell\ne\ell'$ the least $u$ with $\ell_u\ne\ell'_u$ is $\ord_2(\ell-\ell')$, the largest $u$ with $2^u\mid\ell-\ell'$, so $\abs{\xi_\ell-\xi_{\ell'}}=\abs{\varpi'}^{\ord_2(\ell-\ell')}$. For $0\le\ell\le d$ the numbers $\abs{\ell-\ell'}$, with $\ell'\le d$ and $\ell'\ne\ell$, are $1,\dots,\ell$ and $1,\dots,d-\ell$, and $\ord_2(\ell!)\le\ell$. Hence
\begin{equation}\label{dp:laws:nodes}
 \prod_{\ell'\le d,\ \ell'\ne\ell}\abs{\xi_\ell-\xi_{\ell'}}^{-1}\le\abs{\varpi'}^{-d}\qquad(0\le\ell\le d).
\end{equation}
So a divided difference of order $d$ on $\xi_0,\dots,\xi_d$ of a map into a normed space is a combination of its $d+1$ values at these nodes, with coefficients of absolute value at most $\abs{\varpi'}^{-d}$.

\emph{Step 3 (item (4) on the unit ball).} Let $P$ be ultrametric and use Step~2 with $\varpi'=\varpi$. If $m=0$, then $a$ is constant and~\eqref{dp:laws:ultrametric} is clear. Let $m\ge1$. By Lemma~\ref{dp:lem:elementary}(3) we may assume that $\alpha_c\ge1$ for all $c$, so $r\le m$. First let $\norm{w_c}\le1$ for all $c$. We show that
\begin{equation}\label{dp:laws:unit}
 \norm{a^{[\alpha]}(w_1,\dots,w_r)}\le(8/\abs\varpi)^m\norm a.
\end{equation}
If $r=1$, then $a^{[\alpha]}(w_1)=a(w_1)$. Let $r\ge2$ and put $\Phi(s)=a(s_1w_1+\dots+s_{r-1}w_{r-1}+w_r)$ on $K^{r-1}$. This is a polynomial of total degree at most $m$ with coefficients in $W$, and $a^{[\alpha]}(w_1,\dots,w_r)$ is its coefficient of $s_1^{\alpha_1}\cdots s_{r-1}^{\alpha_{r-1}}$. Newton's interpolation formula on $\xi_0,\dots,\xi_m$ in each variable gives $\Phi(s)=\sum_g\mathsf d_g\prod_{c<r}\mathsf N_{g_c}(s_c)$, over $g\in\{0,\dots,m\}^{r-1}$. Here $\mathsf N_d(s_c)=\prod_{\ell<d}(s_c-\xi_\ell)$ has coefficients of absolute value at most $1$, and $\mathsf d_g\in W$ is the iterated divided difference of $\Phi$, of order $g_c$ in $s_c$. A divided difference of order $d$ kills polynomials of degree below $d$, so $\mathsf d_g=0$ unless $\abs g\le m$. By Step~2, $\mathsf d_g$ is a combination of at most $\prod_c(g_c+1)\le2^{\abs g}$ values of $\Phi$ at points with coordinates $\xi_\ell$, with coefficients of absolute value at most $\abs\varpi^{-\abs g}$. Each such value is $a$ at a point $\sum_{c<r}\xi_{\ell_c}w_c+w_r$. This point has norm at most $\max_c\norm{w_c}\le1$, because $P$ is ultrametric and $\abs{\xi_\ell}\le1$. So each value has norm at most $\norm a$, and $\norm{\mathsf d_g}\le(2/\abs\varpi)^m\norm a$ by the triangle inequality in $W$. The coefficient of $s_1^{\alpha_1}\cdots s_{r-1}^{\alpha_{r-1}}$ is a sum over at most $\binom{m+r-1}{r-1}\le2^{m+r-1}\le4^m$ indices $g$ of $\mathsf d_g$ times a scalar of absolute value at most $1$. This gives~\eqref{dp:laws:unit}.

\emph{Step 4 (item (4) by scaling).} If some $w_c=0$, both sides of~\eqref{dp:laws:ultrametric} vanish. Otherwise choose $\lambda_c\in K$ with $\abs\varpi<\abs{\lambda_c}\norm{w_c}\le1$. By Lemma~\ref{dp:lem:elementary}(3) and~\eqref{dp:laws:unit},
\[
 \norm{a^{[\alpha]}(w_1,\dots,w_r)}=\prod_c\abs{\lambda_c}^{-\alpha_c}\norm{a^{[\alpha]}(\lambda_1w_1,\dots,\lambda_rw_r)}\le(8/\abs\varpi)^m\norm a\prod_c\bigl(\norm{w_c}/\abs\varpi\bigr)^{\alpha_c},
\]
which is~\eqref{dp:laws:ultrametric}.
\end{proof}

For $a\in\cP_n(P;W)$, $h\in P$ and $0\le i,k\le n$, define the \emph{Hasse coefficients}
\begin{equation}\label{dp:laws:hasse}
 D^{(i)}_a(h)=a^{[(n-i,i)]}(h,\cdot\,)\in\cP_i(P;W),\qquad a^{(k)}(h)=a^{[(n-k,1^k)]}(h,\cdot,\dots,\cdot)\in\Mult^k(P;W),
\end{equation}
which lie in these spaces by Lemma~\ref{dp:lem:elementary}(4). For $k=1$ the two agree, by Lemma~\ref{dp:lem:elementary}(1).

\begin{lemma}[Hasse coefficients are laws]\label{dp:lem:hasse-law}
Let $a\in\cP_n(P;W)$.
\begin{enumerate}
\item The map $h\mapsto D^{(i)}_a(h)$ is a law of degree $n-i$ with values in $\cP_i(P;W)$, of norm at most $C_2^n\norm a$.
\item The map $h\mapsto a^{(k)}(h)$ is a law of degree $n-k$ with values in $\Mult^k(P;W)$, of norm at most $\mathrm e^kn^kC_1^n\norm a$, with $0^0=1$, and its values are symmetric. For $k=1$ this is a law into $\cL(P,W)$ of norm at most $\mathrm e\,n\,C_1^n\norm a$.
\item We have $a(h+y)=\sum_{i=0}^nD^{(i)}_a(h)(y)$.
\item For $0\le k<n$, we have $(a^{(k)})^{(1)}(h)(v)(v_1,\dots,v_k)=a^{(k+1)}(h)(v,v_1,\dots,v_k)$, where the left side is the Hasse coefficient~\eqref{dp:laws:hasse} of the law $a^{(k)}$ of degree $n-k$.
\item We have $a^{(k)}(h)(h,\dots,h)=\frac{n!}{(n-k)!}\,a(h)$. In particular, for $n\ge1$, $a^{(n)}$ is the constant $a^{[1^n]}$, and $a^{[1^n]}(h,\dots,h)=n!\,a(h)$.
\end{enumerate}
\end{lemma}

The algebra of items (1) and (3) is~\cite[Prop.~IV.10; II~\S4]{Ro}, and item (5) is~\cite[Th.~II.1]{Ro}. The bounds are the estimates we need.

\begin{proof}
\emph{Step 1 (item (1): the expansion on a span).} Let $h_1,\dots,h_r\in P$. Expanding $a(z\sum_cs_ch_c+z'y)$ in the vectors $h_1,\dots,h_r,y$ and reading off the coefficient of $z^{n-i}z'^i$ gives
\begin{equation}\label{dp:laws:hasse-expansion}
 D^{(i)}_a\Bigl(\sum_cs_ch_c\Bigr)(y)=\sum_{\abs\gamma=n-i}s^\gamma a^{[(\gamma,i)]}(h_1,\dots,h_r,y).
\end{equation}
This is a homogeneous polynomial of degree $n-i$ in $s$, with coefficients in $\cP_i(P;W)$ by Lemma~\ref{dp:lem:elementary}(4). The norm bound is Theorem~\ref{dp:thm:mixed}(2).

\emph{Step 2 (item (2): the same expansion with linear slots).} The same computation with $a(z\sum_cs_ch_c+\sum_l\tau_lv_l)$ and the coefficient of $z^{n-k}\tau_1\cdots\tau_k$ gives a homogeneous polynomial of degree $n-k$ in $s$, with coefficients in $\Mult^k(P;W)$ by Lemma~\ref{dp:lem:elementary}(4). By~\eqref{dp:laws:scaled} with $\alpha=(n-k,1^k)$, the norm is at most $C_1^n\norm a\,n^n(n-k)^{-(n-k)}$. For $k<n$ we have $n^n(n-k)^{-(n-k)}=n^k(1+k/(n-k))^{n-k}\le\mathrm e^kn^k$, and for $k=n$ the factor is $n^n=n^k$.

\emph{Step 3 (items (3) and (4): coefficients).} Item (3) is the case $z=z'=1$ of $a(z h+z'y)=\sum_ia^{[(n-i,i)]}(h,y)z^{n-i}z'^i$. For (4), the value $a^{(k)}(s_0h+s_1v)(v_1,\dots,v_k)$ is the coefficient of $z^{n-k}\tau_1\cdots\tau_k$ in
\[
 a\Bigl((zs_0)h+(zs_1)v+\sum_l\tau_lv_l\Bigr)=\sum_{\gamma=(\gamma_0,\dots,\gamma_{k+1})}a^{[\gamma]}(h,v,v_1,\dots,v_k)(zs_0)^{\gamma_0}(zs_1)^{\gamma_1}\prod_l\tau_l^{\gamma_{l+1}},
\]
and its coefficient of $s_0^{n-k-1}s_1$ is $a^{[(n-k-1,1,1^k)]}(h,v,v_1,\dots,v_k)=a^{(k+1)}(h)(v,v_1,\dots,v_k)$.

\emph{Step 4 (item (5): the diagonal).} The value $a^{(k)}(h)(h,\dots,h)$ is the coefficient of $z^{n-k}\tau_1\cdots\tau_k$ in $a((z+\sum_l\tau_l)h)=(z+\sum_l\tau_l)^na(h)$, which is $\frac{n!}{(n-k)!}a(h)$.
\end{proof}

\begin{lemma}[polarization]\label{dp:lem:polarization}
Let $a\colon P\to W$ be a law of degree $n\ge1$ and $h_1,\dots,h_n\in P$. Then
\begin{equation}\label{dp:laws:polarization}
 a^{[1^n]}(h_1,\dots,h_n)=\sum_{S\subseteq[n]}(-1)^{n-\abs S}\,a\Bigl(\sum_{c\in S}h_c\Bigr).
\end{equation}
If $P$ is ultrametric, then $\norm{a^{[1^n]}(h_1,\dots,h_n)}\le2^n\norm a\max_c\norm{h_c}^n$, and $\le\norm a\max_c\norm{h_c}^n$ if $W$ is also ultrametric.
\end{lemma}

The coefficient $a^{[1^n]}$ is the value on $P$ of Roby's complete polarization~\cite[II~\S1]{Ro}.

\begin{proof}
\emph{Step 1 (the formula).} For $S\subseteq[n]$ let $1_S\in K^n$ be its indicator vector. By definition, $a(\sum_{c\in S}h_c)=\sum_{\abs\gamma=n}(1_S)^\gamma a^{[\gamma]}(h_1,\dots,h_n)$, and $(1_S)^\gamma$ is $1$ if the support of $\gamma$ lies in $S$ and $0$ otherwise. Hence $\sum_S(-1)^{n-\abs S}(1_S)^\gamma=\sum_{S\supseteq\operatorname{supp}\gamma}(-1)^{n-\abs S}=(1-1)^{n-\abs{\operatorname{supp}\gamma}}$. This is $1$ if $\gamma$ has full support and $0$ otherwise. The only $\gamma$ with $\abs\gamma=n$ and full support is $1^n$. Summing over $\gamma$ gives~\eqref{dp:laws:polarization}.

\emph{Step 2 (the bounds).} If $P$ is ultrametric, then $\norm{\sum_{c\in S}h_c}\le\max_c\norm{h_c}$, so each of the $2^n$ terms of~\eqref{dp:laws:polarization} has norm at most $\norm a\max_c\norm{h_c}^n$. In an ultrametric $W$ the norm of the sum is at most the maximum.
\end{proof}

\subsection{The lifting norm and the strong norm}

\begin{definition}\label{dp:def:strong-norm}
Let $a\in\cP_n(P;W)$.
\begin{enumerate}
\item The \emph{lifting norm} of $a$ is $\operatorname{lift}(a)=\inf\{\norm\beta:\ \beta\colon P^n\to W$ bounded $n$\nobreakdash-linear with $\beta(h,\dots,h)=a(h)$ for all $h\in P\}$, with $\inf\emptyset=\infty$. Always $\norm a\le\operatorname{lift}(a)$.
\item The \emph{strong norm} of $a$ is
\[
 \lvert\!\lvert\!\lvert a\rvert\!\rvert\!\rvert=\sup\frac{\norm{a^{[\alpha]}(h_1,\dots,h_r)}}{\norm{h_1}^{\alpha_1}\cdots\norm{h_r}^{\alpha_r}},
\]
the supremum over $r\ge1$, nonzero $h_1,\dots,h_r\in P$ and $\abs\alpha=n$. For $n=0$ read $\lvert\!\lvert\!\lvert a\rvert\!\rvert\!\rvert=\norm{a(0)}$.
\item For $U,x_0,R,M$ as in Definition~\ref{dp:def:laws}(3), the map $f\colon U\to W$ is \emph{strongly DP}, or \emph{sDP}, at $x_0$ with radius $R$ and bound $M$ if $f(x_0+h)=\sum_na_n(h)$ for $\norm h<R$, with laws $a_n$ of degree $n$ and $\lvert\!\lvert\!\lvert a_n\rvert\!\rvert\!\rvert\le MR^{-n}$. It is sDP at $x_0$ with radius $R$ if this holds for some $M$, sDP at $x_0$ if it holds for some $R,M$, and sDP on $U$ if it is sDP at every point of $U$. By Lemma~\ref{dp:lem:strong-formal}(3), such an $f$ is DP at $x_0$ with the same radius and bound, and the $a_n$ are its coefficient laws.
\end{enumerate}
\end{definition}

By Lemma~\ref{dp:lem:elementary}(3), with $0^0=1$,
\begin{equation}\label{dp:laws:strong}
 \norm{a^{[\alpha]}(h_1,\dots,h_r)}\le\lvert\!\lvert\!\lvert a\rvert\!\rvert\!\rvert\prod_c\norm{h_c}^{\alpha_c}\qquad(h_1,\dots,h_r\in P).
\end{equation}
By Lemma~\ref{dp:lem:strong-formal}(3) below, $\lvert\!\lvert\!\lvert a\rvert\!\rvert\!\rvert$ is finite and at least $\norm a$, so $\lvert\!\lvert\!\lvert\cdot\rvert\!\rvert\!\rvert$ is a norm on $\cP_n(P;W)$.

We use one case of~\eqref{dp:laws:partial} often. For $a\in\cP_n(P;W)$, $x\in P$ and $0\le i\le n$, the map $b(y)=a^{[(n-i,i)]}(x,y)$ is a law of degree $i$, and for $h_1,\dots,h_r\in P$ and $\abs\delta=i$,
\begin{equation}\label{dp:eq:coefficient-identity}
 b^{[\delta]}(h_1,\dots,h_r)=a^{[(n-i,\delta)]}(x,h_1,\dots,h_r),\qquad a(x+y)=\sum_{i\le n}a^{[(n-i,i)]}(x,y).
\end{equation}
The first formula is~\eqref{dp:laws:partial} with one argument $x$, and the second is Lemma~\ref{dp:lem:hasse-law}(3).

\begin{lemma}[formal properties of the strong norm]\label{dp:lem:strong-formal}
Let $a\in\cP_n(P;W)$, with $P$ and $W$ arbitrary normed spaces.
\begin{enumerate}
\item If $T_2\in\cL(P',P)$ and $T_1\in\cL(W,Y)$, then $\lvert\!\lvert\!\lvert T_1\circ a\circ T_2\rvert\!\rvert\!\rvert\le\norm{T_1}\,\lvert\!\lvert\!\lvert a\rvert\!\rvert\!\rvert\,\norm{T_2}^n$. For a linear isometry $j$ of $W$ into a normed space, $\lvert\!\lvert\!\lvert j\circ a\rvert\!\rvert\!\rvert=\lvert\!\lvert\!\lvert a\rvert\!\rvert\!\rvert$.
\item Let $Q\colon P'\to P$ be a bounded linear surjection with preimage constant $\mathsf c$, that is, every $h\in P$ has a preimage $\tilde h$ with $\norm{\tilde h}\le\mathsf c\norm h$. Then $\lvert\!\lvert\!\lvert a\rvert\!\rvert\!\rvert\le\mathsf c^n\lvert\!\lvert\!\lvert a\circ Q\rvert\!\rvert\!\rvert$.
\item We have $\norm a\le\lvert\!\lvert\!\lvert a\rvert\!\rvert\!\rvert\le(C_1n)^n\norm a$, with $0^0=1$. If $P$ is ultrametric, then $\lvert\!\lvert\!\lvert a\rvert\!\rvert\!\rvert\le C_1^n\norm a$, for every $W$. So every law is sDP at every point. On an ultrametric $P$, and for every $W$, a map that is DP at $x_0$ with radius $R$ and bound $M$ is sDP at $x_0$ with radius $R/C_1$ and bound $M$. Conversely, for all $P,W$, a map that is sDP at $x_0$ with radius $R$ and bound $M$ is DP at $x_0$ with radius $R$ and bound $M$.
\item (a) Let $W$ be any normed space. For $0\le i\le n$ and $x\in P$, we have $\lvert\!\lvert\!\lvert a^{[(n-i,i)]}(x,\cdot\,)\rvert\!\rvert\!\rvert\le\lvert\!\lvert\!\lvert a\rvert\!\rvert\!\rvert\,\norm x^{n-i}$, and the law $D^{(i)}_a$ of Lemma~\ref{dp:lem:hasse-law}(1) has strong norm at most $\lvert\!\lvert\!\lvert a\rvert\!\rvert\!\rvert$ as a law into $(\cP_i(P;W),\lvert\!\lvert\!\lvert\cdot\rvert\!\rvert\!\rvert)$.

(b) Let $W$ be Banach, and let $f$ be sDP at $x_0$ with radius $R$, bound $M$ and laws $a_n$. For $\norm x<R$ and $i\ge0$, the series $b_i(x)=\sum_{n\ge i}a_n^{[(n-i,i)]}(x,\cdot\,)$ converges in $(\cP_i(P;W),\lvert\!\lvert\!\lvert\cdot\rvert\!\rvert\!\rvert)$, and $f(x_0+x+y)=\sum_ib_i(x)(y)$ for $\norm y<R-\norm x$. The map $x_0+x\mapsto b_i(x)$ into $(\cP_i(P;W),\lvert\!\lvert\!\lvert\cdot\rvert\!\rvert\!\rvert)$ is sDP at $x_0$ with radius $R$ and bound $MR^{-i}$. If $W$ is not complete, part (b) holds with $f$ regarded as a map into $\widehat W$.
\item If $W$ is ultrametric, then $\lvert\!\lvert\!\lvert a\rvert\!\rvert\!\rvert\le\operatorname{lift}(a)$. For every $W$, if $\operatorname{lift}(a)<\infty$, then
\[
 \norm{a^{[\alpha]}(w_1,\dots,w_r)}\le\frac{n!}{\alpha!}\operatorname{lift}(a)\prod_c\norm{w_c}^{\alpha_c}.
\]
\end{enumerate}
\end{lemma}

\begin{proof}
\emph{Step 1 (items (1) and (2): linear maps).} We have $(T_1\circ a\circ T_2)^{[\alpha]}(h_1,\dots,h_r)=T_1a^{[\alpha]}(T_2h_1,\dots,T_2h_r)$, so (1) follows from~\eqref{dp:laws:strong}. An isometry $j$ gives equality of the norms of the coefficients. For (2), let $Q\tilde h_c=h_c$ with $\norm{\tilde h_c}\le\mathsf c\norm{h_c}$. Then $a^{[\alpha]}(h_1,\dots,h_r)=(a\circ Q)^{[\alpha]}(\tilde h_1,\dots,\tilde h_r)$, and~\eqref{dp:laws:strong} for $a\circ Q$ gives the claim.

\emph{Step 2 (item (5): expanding a lift).} Let $\beta$ be a lift of $a$. Expanding $\beta(\sum_cs_cw_c,\dots,\sum_cs_cw_c)$ shows that $a^{[\alpha]}(w_1,\dots,w_r)$ is the sum of $\beta(w_{c_1},\dots,w_{c_n})$ over the $n!/\alpha!$ tuples $(c_1,\dots,c_n)$ in which each $c$ occurs $\alpha_c$ times. Each term has norm at most $\norm\beta\prod_c\norm{w_c}^{\alpha_c}$. In an ultrametric $W$ the norm of the sum is at most the maximum. Take the infimum over $\beta$.

\emph{Step 3 (item (4), part (a): the bounds).} By~\eqref{dp:eq:coefficient-identity} and~\eqref{dp:laws:strong}, the coefficient of $a^{[(n-i,i)]}(x,\cdot\,)$ at $h_1,\dots,h_r$ of index $\delta$ has norm at most $\lvert\!\lvert\!\lvert a\rvert\!\rvert\!\rvert\,\norm x^{n-i}\prod_c\norm{h_c}^{\delta_c}$. By~\eqref{dp:laws:hasse-expansion}, the coefficient of index $\gamma$ of the law $D^{(i)}_a$ at $w_1,\dots,w_r$ is $a^{[(\gamma,i)]}(w_1,\dots,w_r,\cdot\,)$. By~\eqref{dp:laws:partial} and~\eqref{dp:laws:strong}, it has strong norm at most $\lvert\!\lvert\!\lvert a\rvert\!\rvert\!\rvert\prod_c\norm{w_c}^{\gamma_c}$.

\emph{Step 4 (item (3): comparing the norms).} The case $r=1$ of the supremum gives $\norm a\le\lvert\!\lvert\!\lvert a\rvert\!\rvert\!\rvert$. The first inequality of~\eqref{dp:laws:scaled} and $\alpha_c^{\alpha_c}\ge1$ give $\lvert\!\lvert\!\lvert a\rvert\!\rvert\!\rvert\le(C_1n)^n\norm a$, and Theorem~\ref{dp:thm:mixed}(4) gives the ultrametric bound. By~\eqref{dp:eq:coefficient-identity} and Step~3, a law is at each point $x_0$ a finite sum of laws of finite strong norm, so it is sDP at $x_0$ with every radius. If $P$ is ultrametric and $f$ is DP at $x_0$ with radius $R$, bound $M$ and laws $a_n$, then $\lvert\!\lvert\!\lvert a_n\rvert\!\rvert\!\rvert\le C_1^nMR^{-n}=M(R/C_1)^{-n}$, and the expansion holds for $\norm h<R/C_1$, as $C_1>1$. The converse follows from $\norm{a_n}\le\lvert\!\lvert\!\lvert a_n\rvert\!\rvert\!\rvert$.

\emph{Step 5 (completeness).} Let $W$ be Banach. Then $(\cP_i(P;W),\lvert\!\lvert\!\lvert\cdot\rvert\!\rvert\!\rvert)$ is complete. Indeed, let $(\mathsf b_N)$ be a Cauchy sequence, and let $\mathsf b$ be its pointwise limit, which exists by~\eqref{dp:laws:strong} with $r=1$. By~\eqref{dp:laws:strong}, all the coefficients $\mathsf b_N^{[\delta]}(h_1,\dots,h_r)$ are Cauchy in $W$. Their limits are the coefficients of the polynomials $s\mapsto\mathsf b(\sum_cs_ch_c)$. Also $\norm{\mathsf b}\le\sup_N\lvert\!\lvert\!\lvert\mathsf b_N\rvert\!\rvert\!\rvert<\infty$, so $\mathsf b$ is a law. If $\lvert\!\lvert\!\lvert\mathsf b_N-\mathsf b_{N'}\rvert\!\rvert\!\rvert\le\varepsilon$ for $N,N'\ge N_0$, then letting $N'\to\infty$ in~\eqref{dp:laws:strong} gives $\lvert\!\lvert\!\lvert\mathsf b_N-\mathsf b\rvert\!\rvert\!\rvert\le\varepsilon$ for $N\ge N_0$.

\emph{Step 6 (item (4), part (b): derivatives).} Let $\norm x<R$. By Step~3, the term of degree $n$ of $b_i(x)$ has strong norm at most $MR^{-n}\norm x^{n-i}=MR^{-i}(\norm x/R)^{n-i}$. So the series converges, by Step~5. The double series
\[
 \sum_n\sum_{i\le n}\norm{a_n^{[(n-i,i)]}(x,y)}\le\sum_nMR^{-n}\sum_{i\le n}\norm x^{n-i}\norm y^i\le\sum_nM\Bigl(\frac{\norm x+\norm y}R\Bigr)^n
\]
is finite for $\norm y<R-\norm x$. Summing over $i$ first gives $\sum_na_n(x+y)=f(x_0+x+y)$ by~\eqref{dp:eq:coefficient-identity}. Summing over $n$ first gives $\sum_ib_i(x)(y)$. The coefficient law of degree $m$ of $x_0+x\mapsto b_i(x)$ is $D^{(i)}_{a_{m+i}}$. By Step~3 its strong norm is at most $\lvert\!\lvert\!\lvert a_{m+i}\rvert\!\rvert\!\rvert\le(MR^{-i})R^{-m}$. If $W$ is not complete, regard the $a_n$ as laws into $\widehat W$ and use item (1).
\end{proof}

\subsection{Ultrametric spaces}

Let $\mathcal O=\{\lambda\in K:\abs\lambda\le1\}$, with residue field $\kappa$. A subgroup of $\mathbb R_{>0}$ is discrete or dense, so exactly one of the following holds: (a) $\kappa$ is infinite; (b) $\kappa$ is finite and $\abs{K^\times}$ is dense; (c) $K$ is discretely valued and $\kappa$ has $\mathsf q<\infty$ elements. Put $\theta_n(K)=1$ and $\theta_K=1$ in cases (a) and (b). In case (c), with $\vartheta$ as in Section~\ref{dp:sec:conv}, put
\[
 \theta_n(K)=\abs\vartheta^{-V_{\mathsf q}(n)},\qquad V_{\mathsf q}(n)=\sum_{u\ge1}\lfloor n/\mathsf q^u\rfloor,\qquad\theta_K=\abs\vartheta^{1/(\mathsf q-1)}.
\]
Summing over the base-$\mathsf q$ digits of $n$ gives $V_{\mathsf q}(n)=(n-\sum_un_u)/(\mathsf q-1)\le(n-1)/(\mathsf q-1)$ for $n\ge1$, where $n_u$ are the digits. So $\theta_n(K)\le\theta_K^{1-n}$ for $n\ge1$. The \emph{scaling loss} $\mathfrak s(P)$ is the supremum of the $\eta\in(0,1]$ such that every nonzero $h\in P$ has a multiple $\lambda h$, $\lambda\in K$, with $\eta\le\norm{\lambda h}\le1$. Then $\abs\varpi\le\mathfrak s(P)\le1$, as powers of $\varpi$ rescale each $h$ into norms in $(\abs\varpi,1]$. Also $\mathfrak s(P)=1$ if $\abs{K^\times}$ is dense, and also if $\norm{P\setminus0}\subseteq\abs{K^\times}$, for instance for the space $c_0(J,K)$ of families in $K$ indexed by a set $J$ and tending to $0$, with the supremum norm.

\begin{lemma}[coefficients on ultrametric spaces]\label{dp:lem:strong-coefficients}
Let $P$ and $W$ be ultrametric, $a\colon P\to W$ a law of degree $n$, $w_1,\dots,w_r\in P$ and $\abs\alpha=n$.
\begin{enumerate}
\item We have
\[
 \norm{a^{[\alpha]}(w_1,\dots,w_r)}\le\theta_n(K)\norm a\max_c\norm{w_c}^n.
\]
For every $\varpi'\in K$ with $0<\abs{\varpi'}<1$, we also have
\[
 \norm{a^{[\alpha]}(w_1,\dots,w_r)}\le\abs{\varpi'}^{-n}\norm a\max_c\norm{w_c}^n.
\]
\item We have $\lvert\!\lvert\!\lvert a\rvert\!\rvert\!\rvert\le\theta_n(K)\mathfrak s(P)^{-n}\norm a$. So $\lvert\!\lvert\!\lvert a\rvert\!\rvert\!\rvert=\norm a$ if $\abs{K^\times}$ is dense, and also if $\kappa$ is infinite and $\mathfrak s(P)=1$. In case (c), for $n\ge1$, $\lvert\!\lvert\!\lvert a\rvert\!\rvert\!\rvert\le\abs\vartheta^{-V_{\mathsf q}(n)}\mathfrak s(P)^{-n}\norm a\le\theta_K^{1-n}\mathfrak s(P)^{-n}\norm a$.
\end{enumerate}
\end{lemma}

\begin{proof}
\emph{Step 1 (nodes).} For pairwise distinct nodes $\xi_0,\xi_1,\dots$ in $\mathcal O$ put
\[
 \mathsf D_d=\max_{\ell\le d}\prod_{\ell'\le d,\,\ell'\ne\ell}\abs{\xi_\ell-\xi_{\ell'}}^{-1}.
\]
A divided difference of order $d$ on $\xi_0,\dots,\xi_d$ is a combination of values with coefficients of absolute value at most $\mathsf D_d$. We use three systems of nodes. If $\kappa$ is infinite, take nodes with pairwise distinct residues, so that $\mathsf D_d=1$. For any $\varpi'$, the binary nodes of Step~2 of the proof of Theorem~\ref{dp:thm:mixed} have $\mathsf D_d\le\abs{\varpi'}^{-d}$, by~\eqref{dp:laws:nodes}. In case (c), let $\rho_0=0,\rho_1,\dots,\rho_{\mathsf q-1}\in\mathcal O$ represent the residues and put $\xi_\ell=\sum_u\rho_{\ell_u}\vartheta^u$, with $\ell_u$ the base-$\mathsf q$ digits of $\ell$. For a nonzero integer $e$, let $\ord_{\mathsf q}(e)$ be the largest $u$ with $\mathsf q^u\mid e$. For $\ell\ne\ell'$ the digits agree below position $\ord_{\mathsf q}(\ell-\ell')$ and differ there, so $\abs{\xi_\ell-\xi_{\ell'}}=\abs\vartheta^{\ord_{\mathsf q}(\ell-\ell')}$. The numbers $\abs{\ell-\ell'}$, with $\ell'\le d$ and $\ell'\ne\ell$, are $1,\dots,\ell$ and $1,\dots,d-\ell$. We have $\sum_{e=1}^d\ord_{\mathsf q}(e)=V_{\mathsf q}(d)$ for every $d$. The function $V_{\mathsf q}$ is increasing, and it is superadditive, since $\lfloor e/\mathsf q^u\rfloor+\lfloor e'/\mathsf q^u\rfloor\le\lfloor(e+e')/\mathsf q^u\rfloor$ for integers $e,e'\ge0$. So we get $\mathsf D_d=\max_{\ell\le d}\abs\vartheta^{-V_{\mathsf q}(\ell)-V_{\mathsf q}(d-\ell)}\le\abs\vartheta^{-V_{\mathsf q}(d)}$.

\emph{Step 2 (item (1): Newton interpolation).} If $n=0$, then $a$ is constant. Otherwise we may assume that every $\alpha_c\ge1$ and $r\ge2$, since for $r=1$ the coefficient is $a(w_1)$. As in Step~3 of the proof of Theorem~\ref{dp:thm:mixed}, $a^{[\alpha]}(w_1,\dots,w_r)$ is a coefficient of $\Phi(s)=a(s_1w_1+\dots+s_{r-1}w_{r-1}+w_r)$, and Newton's formula on nodes as in Step~1 gives $\Phi(s)=\sum_{\abs g\le n}\mathsf d_g\prod_{c<r}\mathsf N_{g_c}(s_c)$ with $\mathsf N_d\in\mathcal O[s_c]$. Here $\mathsf d_g$ is a combination of values of $\Phi$ at node points, with coefficients of absolute value at most $\prod_c\mathsf D_{g_c}$. With residue nodes this product is $1$. With the nodes of case (c) it is at most $\abs\vartheta^{-\sum_cV_{\mathsf q}(g_c)}\le\abs\vartheta^{-V_{\mathsf q}(n)}$, as $V_{\mathsf q}$ is superadditive and increasing and $\abs g\le n$. With binary nodes it is at most $\abs{\varpi'}^{-n}$. Each value has norm at most $\norm a\max_c\norm{w_c}^n$, as $P$ is ultrametric and $\abs{\xi_\ell}\le1$. As $W$ is ultrametric, $\mathsf d_g$ obeys the same bound with the product as factor, and so does the coefficient of $\Phi$, which is a combination of the $\mathsf d_g$ with coefficients in $\mathcal O$. This gives the second bound of item (1), and the first in cases (a) and (c). In case (b), choose $\varpi'$ with $\abs{\varpi'}$ arbitrarily close to $1$.

\emph{Step 3 (item (2): scaling).} Let $w_1,\dots,w_r$ be nonzero and let $\eta\in(0,1]$ be as in the definition of $\mathfrak s(P)$. Choose $\lambda_c\in K$ with $\eta\le\norm{\lambda_cw_c}\le1$ for each $c$. By item (1) and Lemma~\ref{dp:lem:elementary}(3),
\[
 \norm{a^{[\alpha]}(w_1,\dots,w_r)}=\prod_c\abs{\lambda_c}^{-\alpha_c}\norm{a^{[\alpha]}(\lambda_1w_1,\dots,\lambda_rw_r)}\le\theta_n(K)\norm a\prod_c\bigl(\norm{w_c}/\eta\bigr)^{\alpha_c}.
\]
Let $\eta$ tend to $\mathfrak s(P)$ and divide by $\prod_c\norm{w_c}^{\alpha_c}$. The equality cases follow from $\norm a\le\lvert\!\lvert\!\lvert a\rvert\!\rvert\!\rvert$, and the last bound from $\theta_n(K)\le\theta_K^{1-n}$.
\end{proof}

\begin{remark}
In case (c) the constant $\theta_{\mathsf q+1}(K)=\abs\vartheta^{-1}$ of Lemma~\ref{dp:lem:strong-coefficients}(2) is attained. Let $P=K^2$ with the maximum norm and $\omega(s_1,s_2)=(s_1^{\mathsf q}s_2-s_1s_2^{\mathsf q})/\vartheta$, with values in $K$. For $s_1,s_2\in\mathcal O$ we have $\abs{s_1^{\mathsf q}-s_1},\abs{s_2-s_2^{\mathsf q}}\le\abs\vartheta$, so $\omega(s_1,s_2)=(s_1^{\mathsf q}-s_1)s_2/\vartheta+s_1(s_2-s_2^{\mathsf q})/\vartheta$ has absolute value at most $1$. Every nonzero $h\in P$ is $\mu h'$ with $\abs\mu=\norm h$ and $\norm{h'}=1$, and $\omega$ is homogeneous of degree $\mathsf q+1$. So $\norm\omega\le1$. The coefficient of $\omega$ of index $(\mathsf q,1)$ at the unit vectors is $1/\vartheta$. So $\omega\ne0$ and $\lvert\!\lvert\!\lvert\omega\rvert\!\rvert\!\rvert\ge\abs\vartheta^{-1}\norm\omega$, while here $\mathfrak s(P)=1$ and $V_{\mathsf q}(\mathsf q+1)=1$.
\end{remark}
%% ---- end sec-03-laws.tex
%% ---- begin sec-04-a-germs.tex
\section{Divided-power germs}\label{dp:sec:germs}

A DP germ is a sum of laws with geometric bounds. This section gives the calculus of such sums, shows that the laws are determined by the germ, and shows that, for complete $W$, the germ stays DP on its whole ball of convergence.

\subsection{Expansions}\label{dp:ssec:expansions}

A normed space $P$ has \emph{finite coordinates} if, for some $r\ge0$, there is a bounded linear bijection $T\colon K^r\to P$ with bounded inverse, where $K^r$ carries the maximum norm.

\begin{theorem}[divided-power calculus]\label{dp:thm:calculus}
\leavevmode
\begin{enumerate}
\item \textup{(Mixed terms.)} If $b\colon W\to Y$ is a bounded polynomial law of degree $m$, then
\[
 \norm{b^{[\alpha]}(w_1,\dots,w_r)}\le(8/\abs\varpi)^m\norm b\,(\norm{w_1}+\dots+\norm{w_r})^m
 \qquad(w_1,\dots,w_r\in W,\ \abs\alpha=m).
\]
Every bounded polynomial law is continuous and DP on its whole domain.
\item \textup{(Composition.)} Let $f$ map an open subset of $P$ to $W$ and $g$ an open subset of $W$ to $Y$. If $f$ is
DP at $x_0$ and $g$ is DP at $f(x_0)$, then $g\circ f$, which is defined near $x_0$, is DP at $x_0$.
\item \textup{(Comparison.)} A map ambient of degree $n$, for instance the diagonal of a bounded $n$\nobreakdash-linear
map, is a bounded polynomial law of degree $n$, and a map analytic at a point is DP there. If $P$ has finite
coordinates, a map DP at a point is analytic there.
\item \textup{(Alternating maps.)} Let $E$, $D$, $F$ and $F'$ be normed spaces and $k\ge1$. The maps
$A^k_{E,D;F}\colon\cL(E,D)\to\cL(\Alt^k(D;F),\Alt^k(E;F))$, $\cL(E,D)\to\cL(\Lambda^k_\pi E,\Lambda^k_\pi D)$,
$u\mapsto\Lambda^k_\pi u$, and $\cL(E,D)\to\cL(\widehat\Lambda^k_\pi E,\widehat\Lambda^k_\pi D)$,
$u\mapsto\widehat\Lambda^k_\pi u$, are ambient of degree $k$, and the joint action $B^k_{E,D;F,F'}$ of
\cite[Proposition~2.2]{I}, restricted to alternating maps, is ambient of degree $k+1$ as a map from
$\cL(E,D)\times\cL(F,F')$ to $\cL(\Alt^k(D;F),\Alt^k(E;F'))$. As bounded polynomial laws, all four have norm at most $1$.
\end{enumerate}
\end{theorem}

\begin{proof}
\emph{Step 1 (items (1) and (2)).} Item~(1) is Theorem~\ref{dp:thm:mixed}(1),(2). Item~(2) is \cite[Lemma~Q.3]{I}. The algebra of composition follows from \cite[VI~\S8, Prop.~VI.6, Cor.; Prop.~I.2; II~\S1]{Ro}; the bound on the coefficient laws of $g\circ f$ is the estimate we need.

\emph{Step 2 (item (3)).} Let $ja(h)=\beta(h,\dots,h)$ as in Definition~\ref{dp:def:laws}(2). By \cite[Lemma~Q.1(3)]{I}, $a$ is a bounded polynomial law of degree $n$, with $\norm a\le\norm\beta$. In particular
analytic maps are DP. Now let $P$ have finite coordinates, with $T\colon K^r\to P$ as in the definition above. If $r=0$, then $P=0$ and there is nothing to prove. So let $r\ge1$, let $e_1,\dots,e_r$ be the unit vectors of $K^r$, and let $a\colon P\to W$ be a bounded polynomial law of degree $n$. Then
$a(Tz)=\sum_{\abs\gamma=n}v_\gamma z^\gamma$ with $v_\gamma=(a\circ T)^{[\gamma]}(e_1,\dots,e_r)$, and
$\norm{v_\gamma}\le(8r/\abs\varpi)^n\norm T^n\norm a$ by~(1). The $n$\nobreakdash-linear map
\[
 \beta_a(h^1,\dots,h^n)=\sum_{t_1\le\dots\le t_n}v_{\gamma(t)}\prod_{c=1}^n(T^{-1}h^c)_{t_c},
\]
summed over the nondecreasing words $t$ in $\{1,\dots,r\}$, where $\gamma(t)$ counts the letters of $t$, has diagonal
$a$, since each monomial $z^\gamma$ comes from exactly one such word, and
$\norm{\beta_a}\le\binom{n+r-1}n(8r\norm T\norm{T^{-1}}/\abs\varpi)^n\norm a\le2^{r-1}(16r\norm T\norm{T^{-1}}/\abs\varpi)^n\norm a$.
Applied to the coefficient laws of a DP expansion, this bound, geometric in $n$, gives an analytic expansion with a
smaller positive radius.

\emph{Step 3 (item (4)).} By \cite[Proposition~2.2]{I} and its proof, the restriction of $B^k_{E,D;F,F'}$ to
alternating maps is the diagonal of a bounded $(k+1)$\nobreakdash-linear map of norm at most $1$ with values in
$\cL(\Alt^k(D;F),\Mult^k(E;F'))$, and its values lie in $\cL(\Alt^k(D;F),\Alt^k(E;F'))$, which postcomposition by
$j_{E,F'}$ embeds isometrically. The same holds for $A^k_{E,D;F}(u)=B^k_{E,D;F,F}(u,\id_F)$ in degree $k$.

Let $\otimes^k_\pi E$ be the normed quotient of $E^{\otimes k}$ by the null space of $\pi$, as in \cite[Section~3.1]{I}. For $u_1,\dots,u_k\in\cL(E,D)$, the operator $u_1\otimes\dots\otimes u_k\colon\otimes^k_\pi E\to\otimes^k_\pi D$ has norm at most $\prod_c\norm{u_c}$ and is multilinear in $(u_1,\dots,u_k)$. Let $\wedge^k_E\colon\otimes^k_\pi E\to\Lambda^k_\pi E$ and $\wedge^k_D\colon\otimes^k_\pi D\to\Lambda^k_\pi D$ be the canonical maps. Postcomposition with the inclusion $\Lambda^k_\pi D\subseteq\widehat\Lambda^k_\pi D$ is a linear isometry of $\cL(\Lambda^k_\pi E,\Lambda^k_\pi D)$ into $\cL(\Lambda^k_\pi E,\widehat\Lambda^k_\pi D)$. By \cite[Lemma~3.7]{I}, $\Theta\mapsto\Theta\circ\wedge^k_E$ is a linear isometry of
$\cL(\Lambda^k_\pi E,\widehat\Lambda^k_\pi D)$ into $\cL(\otimes^k_\pi E,\widehat\Lambda^k_\pi D)$. Restriction to the dense
subspace $\Lambda^k_\pi E$ is an isometric isomorphism of $\cL(\widehat\Lambda^k_\pi E,\widehat\Lambda^k_\pi D)$ onto
$\cL(\Lambda^k_\pi E,\widehat\Lambda^k_\pi D)$, as $\widehat\Lambda^k_\pi D$ is complete. These isometries map
$\Lambda^k_\pi u$ and $\widehat\Lambda^k_\pi u$ to $\wedge^k_D\circ(u\otimes\dots\otimes u)$, the diagonal of the
$k$\nobreakdash-linear map $(u_1,\dots,u_k)\mapsto\wedge^k_D\circ(u_1\otimes\dots\otimes u_k)$ of norm at most $1$. Item~(3)
bounds the four norms by $1$.
\end{proof}

DP and analyticity differ. Let $\operatorname{char}K=p>0$ and $k\ge p$. For all normed $E,D,F$, the map $A^k_{E,D;F}$ is DP on $\cL(E,D)$ by
Theorem~\ref{dp:thm:calculus}(4) and~(1). For the spaces $E$, $D=E$, $F$ of \cite[Theorem~7.2(1)]{I} it is analytic at
no point.

\begin{lemma}[coefficient laws are determined by the germ]\label{dp:lem:coefficients-unique}
Let $f$ be DP at $x_0$ with radius $R$ and coefficient laws $a_n$, and let $h\in P$.
\begin{enumerate}
\item If $\mathsf r>0$, $w_n\in W$, $\sup_n\norm{w_n}\mathsf r^n<\infty$ and $f(x_0+sh)=\sum_ns^nw_n$ for all $s\in K$ with $\abs s<\mathsf r$, then $w_n=a_n(h)$ for all $n$. In particular, the $a_n$ depend only on the germ of $f$ at $x_0$: a second DP expansion at $x_0$ has the same laws.
\item Let $j\colon W\to Z$ be a linear isometry into a normed space, and let $\beta_n\colon P^n\to Z$ be bounded $n$\nobreakdash-linear maps with $\sup_n\norm{\beta_n}R'^n<\infty$ for some $R'>0$ and $jf(x_0+h')=\sum_n\beta_n(h',\dots,h')$ for $\norm{h'}<R'$. Then $ja_n(h)=\beta_n(h,\dots,h)$ for all $n$. No condition on the diagonals of the $\beta_n$ is assumed; the identity shows that they lie in $j(W)$.
\end{enumerate}
\end{lemma}

\begin{proof}
\emph{Step 1 (item (1)).} Let $\mathsf r'>0$ satisfy $\mathsf r'\le\mathsf r$ and $\mathsf r'\norm h<R$. For $\abs s<\mathsf r'$ we have $f(x_0+sh)=\sum_ns^na_n(h)$, and $\norm{a_n(h)}\mathsf r'^n$ is bounded. So $\sum_ns^n(a_n(h)-w_n)$ converges to $0$ for $\abs s<\mathsf r'$, and \cite[Lemma~A.1]{I} gives $a_n(h)=w_n$. A second expansion gives a second sequence $w_n$, with the second laws evaluated at $h$.

\emph{Step 2 (item (2)).} Let $\mathsf r'>0$ satisfy $\mathsf r'\norm h<\min(R,R')$. For $\abs s<\mathsf r'$ both $\sum_ns^nja_n(h)$ and $\sum_ns^n\beta_n(h,\dots,h)$ converge to $jf(x_0+sh)$ in $Z$, with geometric bounds. Apply \cite[Lemma~A.1]{I} in $Z$ to their difference.
\end{proof}

\begin{definition}\label{dp:def:hasse}
Let $f$ be DP at $x_0$ with coefficient laws $a_n$. The \emph{Hasse derivative of order $n$} of $f$ at $x_0$ is $D^{(n)}f(x_0)=a_n\in\cP_n(P;W)$. The \emph{Hasse jet of order $r$} of $f$ at $x_0$, for $r\in\N\cup\{\infty\}$, is $(D^{(n)}f(x_0))_{0\le n\le r}$. The \emph{Hasse jet} is the one of order $\infty$. For $P=K$, $D^{(n)}f(x_0)(h)=w_nh^n$, with $w_n$ the $n$-th coefficient of the power series of $f$ at $x_0$.
\end{definition}

\begin{theorem}[quasi-analyticity]\label{dp:thm:quasi-analytic}
Let $U\subseteq P$ be open, $x_0\in U$, and let $f,g\colon U\to W$ be DP at $x_0$.
\begin{enumerate}
\item If $D^{(n)}f(x_0)=D^{(n)}g(x_0)$ for all $n$, then $f=g$ on $B(x_0,R)$ for every $R$ such that $f$ and $g$ are both DP at $x_0$ with radius $R$. Conversely, if $f=g$ near $x_0$, their Hasse jets at $x_0$ agree.
\item \textup{(Identity on the ball.)} If $f$ is DP at $x_0$ with radius $R$ and vanishes on a nonempty open subset $U'$ of $B(x_0,R)$, then $D^{(n)}f(x_0)=0$ for all $n$ and $f=0$ on $B(x_0,R)$.
\item \textup{(No identity on larger sets.)} If $\phi$ is a nonzero bounded linear functional on $P$ and $w\in W$ is nonzero, then the map equal to $w$ on $\{y\in P:\abs{\phi(y)}\le1\}$ and to $0$ elsewhere is DP on $P$ and not constant. If $P$ is ultrametric and nonzero, the same holds with an open ball of $P$ in place of this set.
\item If $W$ is complete and $(a_n)$ is a sequence with $a_n\in\cP_n(P;W)$ and $\sup_n\norm{a_n}R^n<\infty$ for some $R>0$, then $h\mapsto\sum_na_n(h)$, $\norm h<R$, is DP at $0$ with Hasse jet $(a_n)$. So for $W$ complete, taking the Hasse jet is a bijection from DP germs at $x_0$ onto sequences of laws with geometric bounds.
\end{enumerate}
\end{theorem}

\begin{proof}
\emph{Step 1 (item (1)).} Both maps equal $\sum_na_n(h)$ at $x_0+h$ for $\norm h<R$, with the common laws $a_n$. The converse is Lemma~\ref{dp:lem:coefficients-unique}(1).

\emph{Step 2 (item (2), on a line through $x_0$).} Let $M=\sup_n\norm{a_n}R^n$. Let $y\in P$ with $\norm y<R$, and choose $\mathsf r_0>1$ with $\mathsf r_0\norm y<R$. Put $\mathsf c_n=a_n(y)$ and $\mathsf g(s)=f(x_0+sy)$. Then $\norm{\mathsf c_n}\le M\mathsf r_0^{-n}$ and $\mathsf g(s)=\sum_n\mathsf c_ns^n$ for $\abs s<\mathsf r_0$. Suppose that $\mathsf g(s)=0$ for $\abs{s-1}<\delta$, where $0<\delta\le1$. For $\abs\lambda\le1$,
\[
 \mathsf g(1+\lambda)=\sum_n\sum_{i\le n}\binom ni \mathsf c_n\lambda^i.
\]
This double series converges absolutely in $\widehat W$, since $\abs{\binom ni}\le1$ in $K$ and $\sum_n(n+1)M\mathsf r_0^{-n}<\infty$. So $\mathsf g(1+\lambda)=\sum_i\mathsf b_i\lambda^i$ with $\mathsf b_i=\sum_{n\ge i}\binom ni\mathsf c_n$ and $\norm{\mathsf b_i}\le\sum_n(n+1)M\mathsf r_0^{-n}$. By \cite[Lemma~A.1]{I} in $\widehat W$ with radius $\delta$, all $\mathsf b_i$ vanish. So $\mathsf g(1+\lambda)=0$ for all $\abs\lambda\le1$, hence $\mathsf g(s)=0$ for $\abs s<1$. By the same lemma with radius $1$, all $\mathsf c_n$ vanish.

\emph{Step 3 (item (2), the ball).} Let $x'\in U'$ and $y_0=x'-x_0$. For $y$ near $y_0$ and $s$ near $1$, $x_0+sy\in U'$. So Step~2 gives $a_n(y)=0$ for all $n$ and all $y$ in a neighborhood of $y_0$. Let $n\ge0$ and $h\in P$. The polynomial $s\mapsto a_n(y_0+sh)$ vanishes for all small $\abs s$, hence identically, and its coefficient of $s^n$ is $a_n(h)$. So $a_n=0$, and $f=0$ on $B(x_0,R)$.

\emph{Step 4 (item (3)).} The set $V_\phi=\{y:\abs{\phi(y)}\le1\}$ is open and closed: $\phi$ is continuous, and if $\abs\lambda\le1$ and $\abs{\lambda'-\lambda}<1$ then $\abs{\lambda'}\le1$, while if $\abs\lambda>1$ and $\abs{\lambda'-\lambda}<\abs\lambda$ then $\abs{\lambda'}=\abs\lambda$. So the map is locally constant. A locally constant map is DP at every point, with $a_0$ its value and $a_n=0$ for $n\ge1$. It is not constant: $\phi(y)\ne0$ for some $y$, and suitable multiples of $y$ lie inside and outside $V_\phi$. On an ultrametric space an open ball is open and closed, and multiples of a nonzero vector lie outside it.

\emph{Step 5 (item (4)).} For $\norm h\le R_1<R$, $\norm{a_n(h)}\le\sup_m(\norm{a_m}R^m)(R_1/R)^n$, so the series converges in $W$. So the map is DP at $0$ with these laws. The map from germs to jets is injective by item~(1), and it is surjective by the first assertion, after a translation to $x_0$.
\end{proof}

\subsection{The whole ball}\label{dp:ssec:ball}

\begin{lemma}[node lemma]\label{dp:lem:node}
Let $\varpi'\in K$ with $0<\abs{\varpi'}<1$ and $\mathsf a_{\varpi'}=\log_2(1/\abs{\varpi'})$. Let $W$ be a normed space, $n\ge1$, $\tau\in K^\times$, and $\mathsf p(s)=\sum_{i=0}^n\mathsf c_is^i$ with $\mathsf c_i\in W$. Let $\mathsf k$ be least with $2^{\mathsf k}\ge n+1$, and let $S=\{\sum_{l<\mathsf k}\mathsf d_l\varpi'^l:\mathsf d_l\in\{0,1\}\}$. Then $S$ has $N=2^{\mathsf k}\le2n$ elements, $\abs\xi\le1$ for $\xi\in S$, and with $\mathsf V=\max_{\xi\in S}\norm{\mathsf p(\tau\xi)}$,
\begin{equation}\label{dp:4a:eq:node}
 \norm{\mathsf c_i}\le2n\,\mathsf V\,\abs\tau^{-i}\Bigl(\frac{2\mathrm en}i\Bigr)^{\mathsf a_{\varpi'} i}\qquad(1\le i\le n).
\end{equation}
If $K$ has $n+1$ elements of absolute value at most $1$ at pairwise distance $1$, for instance if its residue field is infinite, then $\norm{\mathsf c_i}\le(n+1)\mathsf V'\abs\tau^{-i}$ for $0\le i\le n$, with $\mathsf V'$ the maximum of $\norm{\mathsf p(\tau\xi)}$ over these elements $\xi$.
\end{lemma}

\begin{proof}
\emph{Step 1 (spacing).} Let $\xi\ne\xi'$ in $S$ have first differing digit at $l$. Then $\abs{\xi-\xi'}=\abs{\varpi'}^l$, since the digit difference is $\pm1$ and the later terms are smaller. So $S$ has $2^{\mathsf k}$ elements. As $2^{\mathsf k-1}\le n$, $N\le2n$. For fixed $\xi$, the number of $\xi'$ at distance $\abs{\varpi'}^l$ is $2^{\mathsf k-l-1}$. Hence $\prod_{\xi'\ne\xi}\abs{\xi-\xi'}=\abs{\varpi'}^{\mathsf D}$ with $\mathsf D=\sum_{l<\mathsf k}l2^{\mathsf k-l-1}$, the same for all $\xi$. Taking $\xi=0$: the set $S\setminus\{0\}$ has $2^{\mathsf k-l-1}$ elements of absolute value $\abs{\varpi'}^l$, and $\prod_{\eta\in S\setminus\{0\}}\abs\eta=\abs{\varpi'}^{\mathsf D}$.

\emph{Step 2 (Lagrange coefficients).} Let $\ell_\xi(s)=\prod_{\eta\ne\xi}(s-\eta)/\prod_{\eta\ne\xi}(\xi-\eta)$, with $\eta$ in $S$. The coefficient of $s^i$ in the numerator is, up to sign, the elementary symmetric function of degree $N-1-i$ of $S\setminus\{\xi\}$. A nonzero product in it runs over a subset $\mathsf T\subseteq S\setminus\{0,\xi\}$ of size $N-1-i$. Let $\mathsf I=(S\setminus\{0\})\setminus\mathsf T$, of size $i$. Then $\prod_{\mathsf T}\abs\eta=\abs{\varpi'}^{\mathsf D}/\prod_{\mathsf I}\abs\eta\le\abs{\varpi'}^{\mathsf D}/\mathsf Q_i$, with $\mathsf Q_i$ the product of the $i$ smallest absolute values on $S\setminus\{0\}$. So the coefficient of $s^i$ in $\ell_\xi$ has absolute value at most $1/\mathsf Q_i$.

\emph{Step 3 (size of $\mathsf Q_i$).} The $c$-th smallest absolute value on $S\setminus\{0\}$ is $\abs{\varpi'}^{\mathsf k-1-\lfloor\log_2c\rfloor}$: one element at level $\mathsf k-1$, two at level $\mathsf k-2$, and so on. So $\mathsf Q_i=\abs{\varpi'}^{\mathsf E}$ with
\[
 \mathsf E=\sum_{c\le i}(\mathsf k-1-\lfloor\log_2c\rfloor)\le\sum_{c\le i}(\mathsf k-\log_2c)=i\mathsf k-\log_2i!\le i\log_2(\mathrm eN/i),
\]
using $i!\ge(i/\mathrm e)^i$. As $\abs{\varpi'}^{-\log_2\mathsf y}=\mathsf y^{\mathsf a_{\varpi'}}$ for $\mathsf y>0$, we get $1/\mathsf Q_i\le(\mathrm eN/i)^{\mathsf a_{\varpi'} i}\le(2\mathrm en/i)^{\mathsf a_{\varpi'} i}$.

\emph{Step 4 (interpolation).} The polynomial $\mathsf p(\tau s)$ has degree at most $n\le N-1$. So $\mathsf p(\tau s)=\sum_\xi\mathsf p(\tau\xi)\ell_\xi(s)$, and $\mathsf c_i\tau^i$ is the sum over $\xi$ of $\mathsf p(\tau\xi)$ times the coefficient of $s^i$ in $\ell_\xi$. There are $N\le2n$ terms, which gives~\eqref{dp:4a:eq:node}. In the second case, interpolate on the $n+1$ given elements: the numerators have coefficients of absolute value at most $1$, and the denominators have absolute value $1$.
\end{proof}

\begin{definition}\label{dp:def:radius}
Let $f$ be DP at $x_0$ with coefficient laws $a_n$. The series $\sum_na_n(h)$ is the \emph{law series} of $f$ at $x_0$. The \emph{DP radius} of $f$ at $x_0$ is $\rho_{\mathrm{DP}}(f,x_0)=(\limsup_n\norm{a_n}^{1/n})^{-1}\in(0,\infty]$. By Lemma~\ref{dp:lem:coefficients-unique}(1) it depends only on the germ of $f$ at $x_0$. If $f$ is DP at $x_0$ with radius $R$, then $\rho_{\mathrm{DP}}(f,x_0)\ge R$. A map DP at $x_0$ can differ from its law series away from $x_0$ (Theorem~\ref{dp:thm:quasi-analytic}(3)). So in statements about radii at another point we assume that $f$ is the sum of its law series on the ball of the radius in question.
\end{definition}

\begin{theorem}[a law series is DP on its whole ball]\label{dp:thm:whole-ball}
Let $W$ be complete, and let $f\colon U\to W$ be DP at $x_0$ with radius $R$ and bound $M$. Let $x_1=x_0+x$ with $\varrho=\norm x<R$. Let $\varpi'\in K$ with $0<\abs{\varpi'}<1$, $\mathsf a_{\varpi'}=\log_2(1/\abs{\varpi'})$, $0<\sigma<R-\varrho$ and $\mathsf t=(\varrho+\sigma)/R<1$. Then $f$ is DP at $x_1$ with radius
\begin{equation}\label{dp:4a:eq:radius}
 r_1=\min\Bigl(R-\varrho,\ \abs{\varpi'}\sigma\,\bigl(\tfrac{\ln(1/\mathsf t)}{4\mathsf a_{\varpi'}}\bigr)^{\mathsf a_{\varpi'}}\Bigr)
\end{equation}
and bound $M\max\bigl(R/(R-\varrho),\,2\sqrt{\mathsf t}/(1-\sqrt{\mathsf t})^2\bigr)$. There is no further hypothesis on $K$, and the norms need not be ultrametric. Consequently:
\begin{enumerate}
\item $f$ is DP at every point of $B(x_0,R)$.
\item Let $K$ be densely valued, and let $f$ be the sum of its law series on $B(x_0,\rho_{\mathrm{DP}}(f,x_0))$. Then $\rho_{\mathrm{DP}}(f,x_2)\ge\rho_{\mathrm{DP}}(f,x_0)-\norm{x_2-x_0}$ for every $x_2\in B(x_0,\rho_{\mathrm{DP}}(f,x_0))$.
\item Let the residue field of $K$ be infinite, and let $f$ be the sum of its law series on $B(x_0,\rho_{\mathrm{DP}}(f,x_0))$. Then, for every $x_2$ in this ball and every $\varpi'$ as above,
\[
 \rho_{\mathrm{DP}}(f,x_2)\ge\abs{\varpi'}\bigl(\rho_{\mathrm{DP}}(f,x_0)-\norm{x_2-x_0}\bigr).
\]
\end{enumerate}
\end{theorem}

\begin{proof}
\emph{Step 1 (one law).} Let $a_n$ be the coefficient laws, so that $\norm{a_n}\le MR^{-n}$. Put $a_{n,i}(h)=a_n^{[(n-i,i)]}(x,h)$ for $0\le i\le n$. By~\eqref{dp:eq:coefficient-identity}, $a_{n,i}$ is a bounded polynomial law of degree $i$ in $h$, and $a_n(x+h)=\sum_{i\le n}a_{n,i}(h)$. Let $h\ne0$ and $n\ge1$. The polynomial $\mathsf p(s)=a_n(x+sh)$ has the coefficients $a_{n,i}(h)$. Choose $\tau\in K^\times$ with $\abs{\varpi'}\sigma<\abs\tau\norm h\le\sigma$. For $\xi$ in the set $S$ of Lemma~\ref{dp:lem:node}, $\norm{\mathsf p(\tau\xi)}\le\norm{a_n}(\varrho+\sigma)^n\le M\mathsf t^n$. By~\eqref{dp:4a:eq:node},
\begin{equation}\label{dp:4a:eq:one-law}
 \norm{a_{n,i}(h)}\le2nM\mathsf t^n\Bigl(\frac{\norm h}{\abs{\varpi'}\sigma}\Bigr)^i\Bigl(\frac{2\mathrm en}i\Bigr)^{\mathsf a_{\varpi'} i}\qquad(1\le i\le n).
\end{equation}

\emph{Step 2 (summing in $n$).} Put $\mathsf w=\sqrt{\mathsf t}$ and $\mathsf G=(4\mathsf a_{\varpi'}/\ln(1/\mathsf t))^{\mathsf a_{\varpi'}}$. For real $\mathsf y>0$, the maximum of $(2\mathrm e\mathsf y)^{\mathsf a_{\varpi'}}\mathsf w^{\mathsf y}$ is at $\mathsf y=\mathsf a_{\varpi'}/\ln(1/\mathsf w)$, and it equals $(2\mathsf a_{\varpi'}/\ln(1/\mathsf w))^{\mathsf a_{\varpi'}}=\mathsf G$. With $\mathsf y=n/i$ this gives $(2\mathrm en/i)^{\mathsf a_{\varpi'} i}\mathsf w^n\le\mathsf G^i$. As $\mathsf t^n=\mathsf w^n\mathsf w^n$, we may divide~\eqref{dp:4a:eq:one-law} by $\norm h^i$ and take the supremum over $h\ne0$. This gives $\norm{a_{n,i}}\le2n\mathsf w^nM(\mathsf G/(\abs{\varpi'}\sigma))^i$ for $1\le i\le n$. Hence
\begin{equation}\label{dp:4a:eq:sum-n}
 \sum_{n\ge i}\norm{a_{n,i}}\le M\Bigl(\sum_{n\ge1}2n\mathsf w^n\Bigr)\Bigl(\frac{\mathsf G}{\abs{\varpi'}\sigma}\Bigr)^i=\frac{2\mathsf wM}{(1-\mathsf w)^2}\Bigl(\frac{\mathsf G}{\abs{\varpi'}\sigma}\Bigr)^i\qquad(i\ge1).
\end{equation}

\emph{Step 3 (completeness of the space of laws).} The space $\cP_i(P;W)$ is complete, as $W$ is. Indeed, on the span of $h_1,\dots,h_r$ the coefficients of a law are fixed linear combinations of its values at finitely many points, by interpolation, as $K$ is infinite. So a Cauchy sequence of laws converges pointwise to a law, and in norm.

\emph{Step 4 (the expansion).} By~\eqref{dp:4a:eq:sum-n} and Step~3, $a'_i=\sum_{n\ge i}a_{n,i}$ converges in $\cP_i(P;W)$ for $i\ge1$, with
\begin{equation}\label{dp:4a:eq:recentered}
 \norm{a'_i}\le\frac{2\mathsf wM}{(1-\mathsf w)^2}\Bigl(\frac{\mathsf G}{\abs{\varpi'}\sigma}\Bigr)^i.
\end{equation}
Also $\sum_n\norm{a_n(x)}\le MR/(R-\varrho)$, so $a'_0=\sum_na_n(x)$ converges in $W$, with $\norm{a'_0}\le MR/(R-\varrho)$. Since $r_1\le\abs{\varpi'}\sigma/\mathsf G$, the number $\norm{a'_i}r_1^i$ is at most the stated bound for all $i$. Let $\norm h<r_1$. Then $\norm{x+h}<R$, and the family $(a_{n,i}(h))_{n,i}$ is absolutely summable, by $\norm{a_{n,i}(h)}\le\norm{a_{n,i}}\norm h^i$ and~\eqref{dp:4a:eq:sum-n} for $i\ge1$, and by the bound on $\sum_n\norm{a_n(x)}$ for $i=0$. Grouping it, and using $a_n(x+h)=\sum_{i\le n}a_{n,i}(h)$,
\[
 f(x_1+h)=\sum_na_n(x+h)=\sum_n\sum_{i\le n}a_{n,i}(h)=\sum_ia'_i(h).
\]
So $f$ is DP at $x_1$ with radius $r_1$ and the stated bound, and the $a'_i$ are its coefficient laws.

\emph{Step 5 (items).} Item~(1) follows with $\sigma=(R-\varrho)/2$. For items~(2) and~(3), we apply Steps~1--4 to other data. Let $x_2\in B(x_0,\rho_{\mathrm{DP}}(f,x_0))$ and $\varrho_2=\norm{x_2-x_0}$. Let $R_2$ satisfy $\varrho_2<R_2<\rho_{\mathrm{DP}}(f,x_0)$, and let $0<\sigma<R_2-\varrho_2$. Then $f$ is DP at $x_0$ with radius $R_2$, for some bound. So Steps~1--4 hold with $x_2-x_0$, $\varrho_2$, $R_2$ in place of $x$, $\varrho$, $R$, and with this bound as $M$. The laws $a'_i$ are then the coefficient laws of $f$ at $x_2$, and they do not depend on $\varpi'$.

For item~(2), \eqref{dp:4a:eq:recentered} gives $\rho_{\mathrm{DP}}(f,x_2)\ge\abs{\varpi'}\sigma/\mathsf G$ for every $\varpi'$. Since $K$ is densely valued, we may let $\abs{\varpi'}\to1$. Then $\mathsf a_{\varpi'}\to0$ and $\mathsf G\to1$, so $\rho_{\mathrm{DP}}(f,x_2)\ge\sigma$. Let $\sigma\uparrow R_2-\varrho_2$, and then $R_2\uparrow\rho_{\mathrm{DP}}(f,x_0)$.

For item~(3), use the second case of Lemma~\ref{dp:lem:node} in Step~1. It gives $\norm{a_{n,i}(h)}\le(n+1)M\mathsf t^n(\norm h/(\abs{\varpi'}\sigma))^i$. So $\norm{a'_i}\le M(1-\mathsf t)^{-2}(\abs{\varpi'}\sigma)^{-i}$, and $\rho_{\mathrm{DP}}(f,x_2)\ge\abs{\varpi'}\sigma$. Let $\sigma$ and $R_2$ tend to their limits as before.
\end{proof}

\begin{remark}\label{dp:rem:whole-ball}
The bound of Theorem~\ref{dp:thm:mixed}(2) loses $C_2^n$ for every $i$, and by itself it reaches only $\norm x<R/C_2$. By~\eqref{dp:4a:eq:one-law} the mixed coefficient $a_n^{[(n-i,i)]}(x,h)$ has no loss geometric in $n$ when $i$ is small against $n$. Take $\sigma=(R-\varrho)/2$ in~\eqref{dp:4a:eq:radius}. As $\varrho\to R$, the radius $r_1$ divided by $(R-\varrho)^{1+\mathsf a_{\varpi'}}R^{-\mathsf a_{\varpi'}}$ tends to $\abs{\varpi'}/\bigl(2(8\mathsf a_{\varpi'})^{\mathsf a_{\varpi'}}\bigr)$.
\end{remark}

The next two examples show that the inequality of Theorem~\ref{dp:thm:whole-ball}(2) does not hold over every field, and that the DP radius can drop below $\rho_{\mathrm{DP}}(f,x_0)$ inside the ball. In both, $P$ is ultrametric and $W=K$. For $\mathsf u>0$ one may replace the absolute value of $K$ and the norm of $P$ by their $\mathsf u$-th powers. They stay ultrametric, and every DP radius is raised to the power $\mathsf u$. The inequality of Theorem~\ref{dp:thm:whole-ball}(2) is not invariant under this. The strict inequality $\rho_{\mathrm{DP}}(f,x_1)<\rho_{\mathrm{DP}}(f,x_0)$ is invariant.

\begin{example}[the drop over every complete discretely valued field]\label{dp:ex:radius-drop}
Let $K$ be complete and discretely valued, with arbitrary residue field, and let $W=K$. Let $P=K^2$ with the ultrametric norm $\norm{(h_1,h_2)}=\max(\abs{h_1},\abs\vartheta^{1/2}\abs{h_2})$, so that the scaling loss of Section~\ref{dp:sec:laws} is $\mathfrak s(P)=\abs\vartheta^{1/2}<1$. Put $R=\abs\vartheta^{-1/4}$ and, on $B(0,R)$,
\[
 f(h)=\frac1{1-\vartheta(h_1h_2)^2}=\sum_{m\ge0}\vartheta^m(h_1h_2)^{2m}.
\]
\begin{enumerate}
\item $\rho_{\mathrm{DP}}(f,0)=R$.
\item At $x_1=(1,0)$ we have $\norm{x_1}=1<R$ and $\rho_{\mathrm{DP}}(f,x_1)=1<R$, for every value of $\abs\vartheta$.
\item $\rho_{\mathrm{DP}}(f,x_1)<R-\norm{x_1}$ if and only if $\abs\vartheta<1/16$. So the inequality of Theorem~\ref{dp:thm:whole-ball}(2) fails for every local field with $\mathsf q\ge17$ residue elements at its module normalization $\abs\vartheta=1/\mathsf q$, for instance for $\mathbb Q_{17}$.
\end{enumerate}
\end{example}

\begin{proof}[Verification]
\emph{Step 1 (item (1)).} All norms on $P$ lie in $\abs\vartheta^{\Z}\cup\abs\vartheta^{\Z+1/2}$, so every nonzero vector has a multiple of norm in $[\abs\vartheta^{1/2},1]$, and the multiples of $(0,1)$ have norms in $\abs\vartheta^{\Z+1/2}$ only. This gives $\mathfrak s(P)=\abs\vartheta^{1/2}$. If $h_1h_2=0$, the ratio below vanishes. Let $h_1,h_2\ne0$ with $\abs{h_1}=\abs\vartheta^{\mathsf m_1}$ and $\abs{h_2}=\abs\vartheta^{\mathsf m_2}$, $\mathsf m_1,\mathsf m_2\in\Z$. If $\mathsf m_1\le\mathsf m_2$, then $\norm h=\abs\vartheta^{\mathsf m_1}$ and $\abs{h_1h_2}/\norm h^2=\abs\vartheta^{\mathsf m_2-\mathsf m_1}\le1$. If $\mathsf m_1\ge\mathsf m_2+1$, then $\norm h=\abs\vartheta^{\mathsf m_2+1/2}$ and the ratio is $\abs\vartheta^{\mathsf m_1-\mathsf m_2-1}\le1$. There is equality at $(1,1)$. So $\abs{\vartheta(h_1h_2)^2}<1$ on $B(0,R)$, the coefficient law of degree $4m$ is $a_{4m}=\vartheta^m(h_1h_2)^{2m}$, of norm exactly $\abs\vartheta^m$, and $\limsup_n\norm{a_n}^{1/n}=\abs\vartheta^{1/4}$.

\emph{Step 2 (the laws at $x_1$).} Put $a'_i(h)=\sum_{c+2m=i}\vartheta^m\binom{2m}ch_1^ch_2^{2m}$. Each monomial satisfies $\abs{\vartheta^mh_1^ch_2^{2m}}\le\abs\vartheta^m\norm h^c(\abs\vartheta^{-1/2}\norm h)^{2m}=\norm h^i$, and $\abs{\binom{2m}c}\le1$. So $\norm{a'_i}\le1$. For $\norm h<1$ we have $\norm{x_1+h}=1<R$, the monomials tend to $0$ with $i$, and regrouping $\sum_m\vartheta^m(1+h_1)^{2m}h_2^{2m}$ gives $f(x_1+h)=\sum_ia'_i(h)$. So $f$ is DP at $x_1$ with radius $1$, and $\rho_{\mathrm{DP}}(f,x_1)\ge1$. Conversely $a'_{2m}(0,1)=\vartheta^m$ and $\norm{(0,1)}^{2m}=\abs\vartheta^m$, so $\norm{a'_{2m}}\ge1$, and $\rho_{\mathrm{DP}}(f,x_1)\le1$.

\emph{Step 3 (item (3)).} $R-1>1$ if and only if $\abs\vartheta^{-1/4}>2$, that is, $\abs\vartheta<1/16$.
\end{proof}

\begin{example}[the drop with $R\in\abs{K^\times}$ and $\mathfrak s(P)=1$]\label{dp:ex:radius-drop-b}
Let $K$ be complete and discretely valued, with residue field of $\mathsf q<\infty$ elements, let $W=K$, let $\mathcal O$ be the valuation ring, and let $P=K^2$ with the maximum norm, so that $\mathfrak s(P)=1$. For $l\ge1$ fix a set $\mathsf A_l$ of representatives of $\mathcal O/\vartheta^l\mathcal O$ and a set $\mathsf B_l$ of representatives of $\vartheta\mathcal O/\vartheta^l\mathcal O$, with $0\in\mathsf B_l$ representing $0$. Put $n_l=\mathsf q^l+\mathsf q^{l-1}$, $\mathsf v_l=n_l-(\mathsf q^l-1)/(\mathsf q-1)$,
\[
 \mathsf H_l(h)=\prod_{\zeta\in\mathsf A_l}(h_2-\zeta h_1)\prod_{\zeta'\in\mathsf B_l}(h_1-\zeta'h_2),\qquad f=\sum_{l\ge1}\vartheta^{\mathsf v_l}\mathsf H_l\quad\text{on }B(0,R),\qquad R=\abs\vartheta^{-1}.
\]
\begin{enumerate}
\item $\norm{\mathsf H_l}\le\abs\vartheta^{(\mathsf q^l-1)/(\mathsf q-1)}$, and $f$ is DP at $0$ with radius $R\in\abs{K^\times}$ and bound $1$.
\item At $x_1=(1,0)\in B(0,R)$, $\rho_{\mathrm{DP}}(f,x_1)\le\abs\vartheta^{-1+1/(\mathsf q(\mathsf q-1))}<R$.
\end{enumerate}
\end{example}

\begin{proof}[Verification]
\emph{Step 1 (item (1)).} The factors of $\mathsf H_l$ are, up to sign, the determinants of $h$ with the $n_l$ vectors $(1,\zeta)$ and $(\zeta',1)$, and these vectors represent the points of $\mathbb P^1(\mathcal O/\vartheta^l\mathcal O)$ once each. Let $h\in\mathcal O^2$ with $\norm h=1$, and let $1\le c\le l$. Exactly $\mathsf q^{l-c}$ of these points reduce modulo $\vartheta^c$ to the class of $h$, and the factor of such a point is divisible by $\vartheta^c$. A point that reduces to the class of $h$ exactly for the levels $c\le c_0$ has its factor divisible by $\vartheta^{c_0}$, and it occurs in $c_0$ pairs (point, level). So $\abs{\mathsf H_l(h)}$ is at most $\abs\vartheta$ to the power of the number of such pairs, that is, $\abs{\mathsf H_l(h)}\le\abs\vartheta^{\sum_{c=1}^l\mathsf q^{l-c}}$. All norms on $P$ lie in $\abs{K^\times}$, so homogeneity gives the bound on $\norm{\mathsf H_l}$. Hence $\norm{\vartheta^{\mathsf v_l}\mathsf H_l}\le\abs\vartheta^{n_l}=R^{-n_l}$. The $n_l$ are strictly increasing, so these are the coefficient laws of $f$ at $0$.

\emph{Step 2 (one coefficient).} The polynomial $\mathsf H_l(1,s)=\prod_\zeta(s-\zeta)\prod_{\zeta'}(1-\zeta's)$ has degree at most $n_l-1$ in $s$, as $0\in\mathsf B_l$. Its coefficients lie in $\mathcal O$. The first product is monic of degree $\mathsf q^l$, and the second is $1$ plus terms in $s^c$, $c\ge1$, with coefficients of absolute value at most $\abs\vartheta$. So the coefficient of $s^{\mathsf q^l}$ has absolute value $1$.

\emph{Step 3 (item (2)).} For $\abs s\le1$ the terms of $\sum_l\vartheta^{\mathsf v_l}\mathsf H_l(1,s)$ tend to $0$, so $f(1,s)=\sum_iw_is^i$ with $w_i=\sum_l\vartheta^{\mathsf v_l}[s^i]\mathsf H_l(1,s)$ and $\abs{w_i}\le1$. By Theorem~\ref{dp:thm:whole-ball}(1), $f$ is DP at $x_1$. Let $a'_i$ be its coefficient laws. By Lemma~\ref{dp:lem:coefficients-unique}(1), $a'_i(0,1)=w_i$. The sequence $\mathsf v_l$ is strictly increasing, since $\mathsf v_{l+1}-\mathsf v_l=\mathsf q^{l-1}(\mathsf q^2-\mathsf q-1)>0$. Fix $l$ and $i=\mathsf q^l$. For $l'<l$, $\mathsf H_{l'}(1,s)$ has degree at most $n_{l'}-1<\mathsf q^l$. For $l'>l$ the term of $l'$ contributes at most $\abs\vartheta^{\mathsf v_{l'}}<\abs\vartheta^{\mathsf v_l}$. So $\abs{w_{\mathsf q^l}}=\abs\vartheta^{\mathsf v_l}$ by Step~2. As $\norm{(0,1)}=1$, this gives $\norm{a'_{\mathsf q^l}}\ge\abs\vartheta^{\mathsf v_l}$. Finally $\mathsf v_l/\mathsf q^l\to1+\frac1{\mathsf q}-\frac1{\mathsf q-1}=1-\frac1{\mathsf q(\mathsf q-1)}$, which gives the bound.
\end{proof}
%% ---- end sec-04-a-germs.tex
%% ---- begin sec-04-b-maps.tex
\subsection{Derivatives}\label{dp:ssec:derivatives}

In this subsection $U\subseteq P$ is open, $f\colon U\to W$ is a map and $x_0\in U$. We use the Hasse derivatives $D^{(n)}f(x_0)=a_n$ of Subsection~\ref{dp:ssec:expansions}, the Hasse coefficients $D^{(j)}_a$ and $a^{(k)}$ of~\eqref{dp:laws:hasse}, the isometry $\cP_1(P;W)=\cL(P,W)$, and the isometric identification of $\cL(P,\Mult^k(P;W))$ with $\Mult^{k+1}(P;W)$. Here $C^\infty$ means iterated Fr\'echet differentiability.

\begin{theorem}[differentiability]\label{dp:thm:differentiable}
Let $f$ be DP at $x_0$ with radius $R$, bound $M$ and coefficient laws $a_n$.
\begin{enumerate}
\item \textup{(Strict differentiability.)} For $0<\rho<R/(2C_2)$ and $\norm h,\norm{h'}\le\rho$,
\begin{equation}\label{dp:4b:eq:strict}
 \norm{f(x_0+h)-f(x_0+h')-a_1(h-h')}\le\mathsf L(\rho)\norm{h-h'},\qquad
 \mathsf L(\rho)=\sum_{n\ge2}nC_2^nMR^{-n}(2\rho)^{n-1},
\end{equation}
and $\mathsf L(\rho)\to0$ as $\rho\to0$. So $f$ is strictly differentiable at $x_0$, with $Df(x_0)=a_1$.
\item \textup{(Hasse derivatives.)} Let $f$ be DP on $U$. For every $j\ge0$ the map $x\mapsto D^{(j)}f(x)$ from $U$ to $\cP_j(P;W)$ is DP on $U$. For $\norm h<R/C_2$,
\begin{equation}\label{dp:4b:eq:hasse-expansion}
 D^{(j)}f(x_0+h)=\sum_{n\ge j}D^{(j)}_{a_n}(h),
\end{equation}
the series converging in $\cP_j(P;W)$, with $n$th term a law of degree $n-j$ in $h$ of norm at most $C_2^nMR^{-n}$. So $D^{(j)}f$ is DP at $x_0$ with radius $R/C_2$ and bound $M(C_2/R)^j$, and
\begin{equation}\label{dp:4b:eq:hasse-iterated-law}
 D^{(i)}\bigl(D^{(j)}f\bigr)(x_0)=\bigl(h\mapsto a_{i+j}^{[(i,j)]}(h,\cdot)\bigr)\in\cP_i\bigl(P;\cP_j(P;W)\bigr).
\end{equation}
\item \textup{(Derivative.)} Let $f$ be DP on $U$. Then $Df\colon U\to\cL(P,W)$ exists and is DP on $U$, and $Df(x_0+h)=\sum_{n\ge1}a_n^{(1)}(h)$ for $\norm h<R/C_2$, with $\norm{a_n^{(1)}}\le\mathrm e\,n\,C_1^nMR^{-n}$.
\item \textup{(Higher derivatives.)} Let $f$ be DP on $U$. Then $f$ is $C^\infty$ on $U$. For every $k\ge1$, $D^kf\colon U\to\Mult^k(P;W)$ is DP on $U$, with $D^kf(x_0+h)=\sum_{n\ge k}a_n^{(k)}(h)$ on a ball around $0$, and
\begin{equation}\label{dp:4b:eq:frechet-law}
 D^kf(x_0)(v_1,\dots,v_k)=a_k^{[1^k]}(v_1,\dots,v_k),\qquad D^kf(x_0)(h,\dots,h)=k!\,a_k(h).
\end{equation}
In particular $D^kf(x_0)$ is symmetric, in every characteristic.
\end{enumerate}
If $W$ is complete, $f$ is DP on $B(x_0,R)$ by Theorem~\ref{dp:thm:whole-ball}, so items \textup{(2)--(4)} apply with $U=B(x_0,R)$.
\end{theorem}

The radius $R/C_2$ in~\eqref{dp:4b:eq:hasse-expansion} is the DP radius that the bound on the terms gives at $x_0$. For $f$ DP on $U$, the identity itself holds for all $\norm h<R$. By Step~4 of the proof of Theorem~\ref{dp:thm:whole-ball}, applied to $f$ as a map into $\widehat W$, the coefficient law of degree $j$ at $x_0+h$ is $\sum_{n\ge j}D^{(j)}_{a_n}(h)$, and the argument of Step~3 below places this series in $\cP_j(P;W)$. So the formula of item~(3) for $Df$ also holds for all $\norm h<R$. The Hasse derivatives are Roby's divided partial derivatives~\cite[II~\S5]{Ro}; the bounds are the estimates we need.

\begin{proof}
\emph{Step 1 (item~(1)).} For $n\ge2$, expanding $a_n(h)=a_n(h'+(h-h'))$ and using Theorem~\ref{dp:thm:mixed}(2),
\begin{multline*}
 \norm{a_n(h)-a_n(h')}=\Bigl\lVert\sum_{i=1}^na_n^{[(n-i,i)]}(h',h-h')\Bigr\rVert\\
 \le C_2^n\norm{a_n}\sum_{i=1}^n\norm{h'}^{n-i}\norm{h-h'}^i\le nC_2^n\norm{a_n}(2\rho)^{n-1}\norm{h-h'},
\end{multline*}
as $\norm{h'}\le2\rho$ and $\norm{h-h'}\le2\rho$. Summing over $n\ge2$ gives~\eqref{dp:4b:eq:strict}, since $a_0$ cancels and $a_1(h)-a_1(h')=a_1(h-h')$ by Lemma~\ref{dp:lem:elementary}(1). The series $\mathsf L(\rho)$ converges for $2C_2\rho<R$. For $\rho\le\rho_0<R/(2C_2)$ every term has $n-1\ge1$, so $\mathsf L(\rho)\le(\rho/\rho_0)\mathsf L(\rho_0)$, and $\mathsf L(\rho)\to0$. With $h'=0$, \eqref{dp:4b:eq:strict} says that $a_1$ is the Fr\'echet derivative of $f$ at $x_0$, and the uniform form is strict differentiability.

\emph{Step 2 (item~(2): the double series).} Let $\norm h+\norm y<R/C_2$. By Lemma~\ref{dp:lem:hasse-law}(3) and~(1), $a_n(h+y)=\sum_{j\le n}D^{(j)}_{a_n}(h)(y)$, and the $j$th term has norm at most $C_2^nMR^{-n}\norm h^{n-j}\norm y^j$. So the double family is absolutely summable in $\widehat W$. As $C_2>1$, also $\norm{h+y}<R$, so
\[
 f(x_0+h+y)=\sum_j\mathsf m_j(h)(y),\qquad \mathsf m_j(h)=\sum_{n\ge j}D^{(j)}_{a_n}(h),
\]
with $\mathsf m_j(h)$ converging in the complete space $\cP_j(P;\widehat W)$ (Step~3 of the proof of Theorem~\ref{dp:thm:whole-ball}) and $\norm{\mathsf m_j(h)}\le M(C_2/R)^j(1-C_2\norm h/R)^{-1}$.

\emph{Step 3 (item~(2): identification).} As $f$ is DP at $x_0+h$, uniqueness of law expansions (Lemma~\ref{dp:lem:elementary}(5), in $\widehat W$) gives $\mathsf m_j(h)=D^{(j)}f(x_0+h)\in\cP_j(P;W)$. The partial sums of $\mathsf m_j(h)$ lie in $\cP_j(P;W)$, which sits isometrically in $\cP_j(P;\widehat W)$, and their limit lies there too. So~\eqref{dp:4b:eq:hasse-expansion} converges in $\cP_j(P;W)$. By Lemma~\ref{dp:lem:hasse-law}(1) its $n$th term is a law of degree $n-j$ in $h$ of norm at most $C_2^nMR^{-n}=M(C_2/R)^j(R/C_2)^{-(n-j)}$. This is a DP expansion at $x_0$, and $x_0\in U$ was arbitrary. Uniqueness in the normed space $\cP_j(P;W)$ says that the law of degree $i$ is the term $n=i+j$, which is~\eqref{dp:4b:eq:hasse-iterated-law}.

\emph{Step 4 (item~(3)).} By item~(1) at $x_0+h$, $Df(x_0+h)=D^{(1)}f(x_0+h)$. Apply item~(2) with $j=1$, and the bound of Lemma~\ref{dp:lem:hasse-law}(2).

\emph{Step 5 (item~(4)).} We argue by induction on $k$. Suppose that $D^{k-1}f\colon U\to\Mult^{k-1}(P;W)$ is DP on $U$ with coefficient laws $a_n^{(k-1)}$ at $x_0$; for $k=1$ this is the hypothesis, with $\Mult^0(P;W)=W$ and $a_n^{(0)}=a_n$. Item~(3), for this map, shows that $D^kf$ exists and is DP on $U$, and Lemma~\ref{dp:lem:hasse-law}(4) identifies its coefficient laws as $a_n^{(k)}$. DP maps are continuous by item~(1), so $f$ is $C^k$ for every $k$. At $h=0$ only $n=k$ contributes, and $a_k^{(k)}(0)=a_k^{[(0,1^k)]}(0,\cdot)=a_k^{[1^k]}$ as the argument with exponent $0$ may be deleted (Lemma~\ref{dp:lem:elementary}(3)). This is symmetric by Lemma~\ref{dp:lem:elementary}(4), and the second formula of~\eqref{dp:4b:eq:frechet-law} is Lemma~\ref{dp:lem:hasse-law}(5).
\end{proof}

In the following statement, $h^{\times m}$ denotes $m$ consecutive arguments equal to $h$, and $\alpha!=\prod_c\alpha_c!$.

\begin{proposition}[Hasse and Fr\'echet derivatives]\label{dp:prop:frechet-hasse}
Let $f$ and $g$ be DP on $U$, and let $x\in U$.
\begin{enumerate}
\item For $i,j\ge0$ and $h\in P$, $D^{(i)}\bigl(D^{(j)}f\bigr)(x)(h)(h)=\binom{i+j}iD^{(i+j)}f(x)(h)$.
\item For a law $a$ of degree $k$, $h_1,\dots,h_r\in P$ and $\abs\alpha=k$,
\begin{equation}\label{dp:4b:eq:full-polarization}
 a^{[1^k]}\bigl(h_1^{\times\alpha_1},\dots,h_r^{\times\alpha_r}\bigr)=\alpha!\,a^{[\alpha]}(h_1,\dots,h_r).
\end{equation}
So $a^{[1^k]}=0$ if and only if $a^{[\alpha]}=0$ for every $r$ and every $\alpha\in\N^r$ with $\abs\alpha=k$ and $\alpha!\ne0$ in $K$. If $k!\ne0$ in $K$, then $a^{[1^k]}=0$ if and only if $a=0$.
\item $D^kf(x)=D^kg(x)$ for all $k\ge0$ if and only if $(D^{(k)}f(x)-D^{(k)}g(x))^{[\alpha]}=0$ for all $k$ and all $\alpha$ with $\abs\alpha=k$ and $\alpha!\ne0$ in $K$. If $\operatorname{char}K=0$, this holds if and only if $f=g$ near $x$, and $D^{(k)}f(x)(h)=D^kf(x)(h,\dots,h)/k!$ for every $f$ that is DP on $U$.
\item Let $p=\operatorname{char}K>0$, let $\phi$ be a nonzero bounded linear functional on $P$ and $w\in W\setminus\{0\}$. The map $f_1(x)=\phi(x)^pw$ is a law, $D^kf_1\equiv0$ for all $k\ge1$, $D^{(1)}f_1\equiv0$, and $D^{(p)}f_1(x)(h)=\phi(h)^pw\ne0$ when $\phi(h)\ne0$. Also $f_1(0)=0$, and $f_1$ is not identically $0$ on any neighborhood of $0$. So neither the Fr\'echet derivatives nor the iterates of first-order Hasse derivatives determine a DP germ.
\end{enumerate}
\end{proposition}

\begin{proof}
\emph{Step 1 (item~(1)).} By~\eqref{dp:4b:eq:hasse-iterated-law} at $x$, the left side is $a^{[(i,j)]}(h,h)$ with $a=D^{(i+j)}f(x)$. This is the coefficient of $s^is'^j$ in $a(sh+s'h)=(s+s')^{i+j}a(h)$.

\emph{Step 2 (item~(2)).} Let $v_1,\dots,v_k$ be the arguments on the left of~\eqref{dp:4b:eq:full-polarization}, in consecutive blocks of sizes $\alpha_1,\dots,\alpha_r$, and let $\mathsf S_c$ be the sum of the variables $t_l$ of block $c$. Then $a(\sum_lt_lv_l)=a(\sum_c\mathsf S_ch_c)=\sum_{\abs\gamma=k}\mathsf S^\gamma a^{[\gamma]}(h_1,\dots,h_r)$. The coefficient of $t_1\cdots t_k$ on the left is $a^{[1^k]}(v_1,\dots,v_k)$. On the right, the blocks are disjoint, so $\mathsf S^\gamma$ contains $t_1\cdots t_k$ only for $\gamma=\alpha$, and then with coefficient $\alpha!$. For the kernel: if $a^{[1^k]}=0$, then~\eqref{dp:4b:eq:full-polarization} gives $a^{[\alpha]}=0$ when $\alpha!\ne0$; conversely $1^k$ is one of these $\alpha$. If $k!\ne0$, take $r=1$ and $\alpha=(k)$, for which $a^{[(k)]}=a$.

\emph{Step 3 (item~(3)).} By~\eqref{dp:4b:eq:frechet-law}, $D^kf(x)=D^kg(x)$ means $(D^{(k)}f(x)-D^{(k)}g(x))^{[1^k]}=0$. For $k=0$ both sides are $f(x)$ and $g(x)$, as $D^0f(x)=f(x)=D^{(0)}f(x)$. Apply item~(2). In characteristic $0$ all $k!\ne0$, so the Hasse derivatives agree, and Theorem~\ref{dp:thm:quasi-analytic} gives $f=g$ near $x$. The last formula is~\eqref{dp:4b:eq:frechet-law}.

\emph{Step 4 (item~(4)).} The map $f_1$ is a law of degree $p$ with $\norm{f_1}\le\norm\phi^p\norm w$, and it is DP on $P$ by Theorem~\ref{dp:thm:calculus}(1). By additivity of Frobenius, $f_1(x+h)=\phi(x)^pw+\phi(h)^pw$. So $D^{(p)}f_1(x)=\phi(\cdot)^pw$ and $D^{(n)}f_1(x)=0$ for $n\ne0,p$. Also $(\sum_lt_l\phi(v_l))^p=\sum_lt_l^p\phi(v_l)^p$ has no monomial $t_1\cdots t_p$, so the coefficient $a^{[1^p]}$ of $a=\phi(\cdot)^pw$ vanishes, and~\eqref{dp:4b:eq:frechet-law} gives $D^kf_1\equiv0$ for $k\ge1$. Scaling a vector $h$ with $\phi(h)\ne0$ gives the last claim.
\end{proof}

\begin{proposition}[chain rule]\label{dp:prop:chain-rule}
Let $f$ be DP at $x_0$ with coefficient laws $a_n$, and let $g$ map an open subset of $W$ to $Y$, DP at $f(x_0)$ with coefficient laws $b_m$. Then $g\circ f$ is DP at $x_0$, $D^{(0)}(g\circ f)(x_0)=b_0$, and for $d\ge1$,
\begin{equation}\label{dp:4b:eq:chain-rule}
 D^{(d)}(g\circ f)(x_0)(h)=\sum_{\deg\alpha=d}b_{\abs\alpha}^{[\alpha]}\bigl(a_1(h),\dots,a_d(h)\bigr),
\end{equation}
the finite sum over $\alpha\in\N^d$ with $\sum_nn\alpha_n=d$. The law of degree $d$ of $g\circ f$ has linearization $\sum_{r=1}^dL_{b_r}\circ\mathsf Q_{r,d}$ on $\Gamma^d(P)$, with $\mathsf Q_{r,d}$ as in~\cite[Lemma~Q.4]{I}. It involves only the $a_n$ with $1\le n\le d$ and the $b_m$ with $m\le d$.
\end{proposition}

\begin{proof}
Formula~\eqref{dp:4b:eq:chain-rule} is~\cite[Lemma~Q.3]{I}, and the linearization is~\cite[Lemma~Q.4]{I}. The condition $\deg\alpha=d$ forces $\alpha_n=0$ for $n>d$ and $\abs\alpha\le d$. The algebra is~\cite[VI~\S8, Prop.~VI.6, Cor.; Prop.~I.2; II~\S1]{Ro}; the bound behind~\cite[Lemma~Q.3]{I} is the estimate we need.
\end{proof}

\subsection{Products and inverse functions}\label{dp:ssec:inverse}

Let $\mathcal A$ be a normed $K$\nobreakdash-algebra, so $\norm{ww'}\le\norm w\norm{w'}$. A map into a product with the maximum norm is DP at a point when its components are, since a pair of laws of degree $n$ is a law of degree $n$.

\begin{proposition}[products and inversion]\label{dp:prop:products}
\leavevmode
\begin{enumerate}
\item If $a,b\colon P\to\mathcal A$ are laws of degrees $m$ and $n$, then $ab\colon h\mapsto a(h)b(h)$ is a law of degree $m+n$ with $\norm{ab}\le\norm a\norm b$.
\item Let $f,g\colon U\to\mathcal A$ be DP at $x_0$ with radii $R,R'$, bounds $M,M'$ and coefficient laws $a_n,b_n$, and let $R_1=\min(R,R')$. Then $(fg)(x_0+h)=\sum_d\mathsf m_d(h)$ for $\norm h<R_1$, with $\mathsf m_d=\sum_{m+n=d}a_mb_n$, $\norm{\mathsf m_d}\le(d+1)MM'R_1^{-d}$, and $\norm{\mathsf m_d}\le MM'R_1^{-d}$ if $\mathcal A$ is ultrametric. So $fg$ is DP at $x_0$.
\item For $U\subseteq P$ open, the DP maps $U\to\mathcal A$ form a $K$\nobreakdash-algebra, unital if $\mathcal A$ is, and they form a sheaf of $K$\nobreakdash-algebras on $P$. For $\mathcal A=K$ we write $\mathcal O_{\rm DP}$ for this sheaf.
\item Let $\mathcal A$ be a unital Banach algebra or $\mathcal A=K$. Let $g\colon U\to\mathcal A$ be DP at $x_0$ with radius $R$ and coefficient laws $a_n$, with $\sup_{n\ge1}\norm{a_n}R^n\le M$ and $M>0$, and let $w_0=g(x_0)$ be invertible. Then $x\mapsto g(x)^{-1}$ is defined near $x_0$ and DP at $x_0$, with every radius below $R/(1+\norm{w_0^{-1}}M)$, and with radius $R\min(1,1/(\norm{w_0^{-1}}M))$ if $\mathcal A$ is ultrametric.
\item The stalks of $\mathcal O_{\rm DP}$ are local rings: a germ is invertible if and only if its value at the point is nonzero.
\end{enumerate}
\end{proposition}

\begin{proof}
\emph{Step 1 (item~(1)).} On the span of finitely many vectors, $a$ and $b$ are homogeneous polynomials with coefficients in $\mathcal A$, and their pointwise product is a homogeneous polynomial of degree $m+n$. The norm bound is immediate.

\emph{Step 2 (item~(2)).} For $\norm h<R_1$ both series converge absolutely, so their Cauchy product converges in $\widehat{\mathcal A}$ to $f(x_0+h)g(x_0+h)\in\mathcal A$. Its partial sums lie in $\mathcal A$, so it converges in $\mathcal A$. The bounds on $\mathsf m_d$ are item~(1) with the triangle or the ultrametric inequality. For $\theta<R_1$, $\sup_d\norm{\mathsf m_d}\theta^d<\infty$.

\emph{Step 3 (item~(3)).} Sums and scalar multiples are clear, constants are DP, and products are item~(2). Being DP on $U$ is the condition of being DP at every point, and this depends only on the germ. So restriction and gluing hold.

\emph{Step 4 (item~(4): the majorant).} Write $g(x_0+h)=w_0(1+\zeta(h))$ with $\zeta=\sum_{n\ge1}\zeta_n$, $\zeta_n=w_0^{-1}a_n$, so $\norm{\zeta_n}\le\norm{w_0^{-1}}MR^{-n}$. For real $\theta\ge0$ let $\bar\zeta(\theta)=\sum_{n\ge1}\norm{\zeta_n}\theta^n$, which is at most $\norm{w_0^{-1}}M(\theta/R)(1-\theta/R)^{-1}$ for $\theta<R$. Let $\mathsf i_0=1$ and
\[
 \mathsf i_d=\sum_{l\ge1}(-1)^l\sum_{n_1+\dots+n_l=d,\ n_i\ge1}\zeta_{n_1}\cdots\zeta_{n_l}\qquad(d\ge1).
\]
By item~(1), $\mathsf i_d$ is a law of degree $d$, and for $d\ge1$, $\norm{\mathsf i_d}$ is at most the coefficient of $\theta^d$ in $(1-\bar\zeta(\theta))^{-1}$. This power series has nonnegative coefficients and converges when $\bar\zeta(\theta)<1$, which holds for $\theta<R/(1+\norm{w_0^{-1}}M)$. So $\sup_d\norm{\mathsf i_d}\theta^d<\infty$ for such $\theta$.

\emph{Step 5 (item~(4): the inverse).} Let $\norm h\le\theta$ with $\theta$ as in Step~4. Then $\sum_l\norm{\zeta(h)}^l\le\sum_l\bar\zeta(\theta)^l<\infty$. So $\sum_l(-\zeta(h))^l$ converges absolutely to $(1+\zeta(h))^{-1}$, and regrouping the absolutely summable family of products by degree gives $\sum_d\mathsf i_d(h)$. Hence $g(x_0+h)^{-1}=\sum_d\mathsf i_d(h)w_0^{-1}$, a DP expansion. If $\mathcal A=K$, the same computation runs in $\Kh$. There $(1+\zeta(h))^{-1}$ lies in $K$, as $\abs{\zeta(h)}<1$, and so do the partial sums, so the series converges in $K$.

\emph{Step 6 (item~(4): ultrametric $\mathcal A$).} Put $\theta_0=R\min(1,1/(\norm{w_0^{-1}}M))$. Let $\norm h<\theta_0$ and $\delta=\max(\norm{w_0^{-1}}M\norm h/R,\norm h/R)<1$. Then $\norm{\zeta(h)}<1$, so $\sum_l(-\zeta(h))^l$ converges to $(1+\zeta(h))^{-1}$, in $\Kh$ if $\mathcal A=K$. A product $\zeta_{n_1}(h)\cdots\zeta_{n_l}(h)$ with $\sum_in_i=d$ has norm at most $(\norm{w_0^{-1}}M\norm h/R)^l(\norm h/R)^{d-l}\le\delta^d$, and for each $d$ there are finitely many. So the family is summable and may be regrouped by degree, which gives $\sum_d\mathsf i_d(h)$, in $K$ if $\mathcal A=K$ as in Step~5. Next, $\norm{\zeta_n}\theta_0^n\le(\norm{w_0^{-1}}M\theta_0/R)(\theta_0/R)^{n-1}\le1$ for $n\ge1$, so item~(1) and the ultrametric inequality give $\norm{\mathsf i_d}\theta_0^d\le1$ for $d\ge1$. Hence $g(x_0+h)^{-1}=\sum_d\mathsf i_d(h)w_0^{-1}$ is a DP expansion with radius $\theta_0$ and bound $\norm{w_0^{-1}}$.

\emph{Step 7 (item~(5)).} A germ with nonzero value at $x$ is nonzero near $x$, as DP maps are continuous, and its inverse is DP at each nearby point by item~(4). So the noninvertible germs are those that vanish at $x$, and they form an ideal.
\end{proof}

\begin{theorem}[inverse functions]\label{dp:thm:inverse}
Let $P$ and $W$ be Banach spaces, where $K$ need not be complete. Let $f\colon U\to W$ be DP at $x_0\in U$, and suppose that $\mathsf T_0=Df(x_0)$ is bijective with bounded inverse. Then there are open sets $V\ni x_0$ and $V'\ni f(x_0)$ such that $f$ restricts to a bijection $V\to V'$, DP on $V$, whose inverse is DP on $V'$. In positive degree $d$, the coefficient law of the inverse at $f(x_0)$ is $\mathsf g_d\circ\mathsf T_0^{-1}$, with $\mathsf g_d$ given by the recursion~\cite[(Q.14)]{I} for the map $\mathsf T_0^{-1}\circ f$.
\end{theorem}

The formal inverse is Roby's~\cite[Th.~VI.3, Prop.~VI.8]{Ro}; only the majorant in the proof of~\cite[Lemma~Q.14]{I} is the estimate we need.

\begin{proof}
\emph{Step 1 (normal form).} Let $R$ be a radius of $f$ at $x_0$, with coefficient laws $a_n$. The map $\tilde f=\mathsf T_0^{-1}\circ f\colon U\to P$ is DP at $x_0$, with coefficient laws $\mathsf T_0^{-1}\circ a_n$, and its law of degree $1$ is $\id_P$ by Theorem~\ref{dp:thm:differentiable}(1). Put $z_0=\tilde f(x_0)$. By~\cite[Lemma~Q.14]{I} there are open sets $U_0\ni x_0$ and $U_1\ni z_0$ and the map $\tilde g(z)=x_0+\sum_{d\ge1}\mathsf g_d(z-z_0)$, DP at $z_0$ with some radius $R_0$, such that $\tilde f$ restricts to a bijection $U_0\to U_1$ whose inverse is $\tilde g$ on $U_1$.

\emph{Step 2 (open sets).} As $P$ and $W$ are complete, Theorem~\ref{dp:thm:whole-ball} shows that $\tilde g$ is DP on $B(z_0,R_0)$ and that $f$ and $\tilde f$ are DP on $B(x_0,R)$. DP maps are continuous. So $B_1=U_1\cap B(z_0,R_0)$ and $V_0'=\{y\in B_1:\tilde g(y)\in B(x_0,R)\}$ are open and contain $z_0$. Put $V=\tilde g(V_0')$. Then $V=U_0\cap B(x_0,R)\cap\tilde f^{-1}(B_1)$, which is open. Put $V'=\mathsf T_0(V_0')$, which is open because $\mathsf T_0$ is a homeomorphism. Then $f=\mathsf T_0\circ\tilde f$ restricts to a bijection $V\to V'$ with inverse $\tilde g\circ\mathsf T_0^{-1}$. The map $f$ is DP on $V\subseteq B(x_0,R)$, and $\tilde g\circ\mathsf T_0^{-1}$ is DP on $V'$ by Theorem~\ref{dp:thm:calculus}(2).

\emph{Step 3 (the coefficient laws).} As $\mathsf T_0^{-1}f(x_0)=z_0$, for small $\norm h$ we have $(\tilde g\circ\mathsf T_0^{-1})(f(x_0)+h)=x_0+\sum_{d\ge1}\mathsf g_d(\mathsf T_0^{-1}h)$. Each $\mathsf g_d\circ\mathsf T_0^{-1}$ is a law of degree $d$ by Lemma~\ref{dp:lem:elementary}(2). For fixed $h$ the series along the line $sh$ converges for small $\abs s$, so its terms have a geometric bound. So Lemma~\ref{dp:lem:coefficients-unique}(1) says that the coefficient law of degree $d\ge1$ of the inverse at $f(x_0)$ is $\mathsf g_d\circ\mathsf T_0^{-1}$.
\end{proof}

\begin{corollary}[implicit functions]\label{dp:cor:implicit}
Let $P$, $Y$ and $Y'$ be Banach spaces, let $U\subseteq P\times Y$ be open, and let $\Phi\colon U\to Y'$ be DP at $(x_0,y_0)$ with $\Phi(x_0,y_0)=0$. Suppose that $\partial_Y\Phi(x_0,y_0)=D\Phi(x_0,y_0)(0,\cdot)$ is a bijection $Y\to Y'$ with bounded inverse. Then there are open sets $V_1\ni x_0$ and $V_2\ni y_0$ and a map $\psi\colon V_1\to V_2$, DP on $V_1$, such that for $(x,y)\in V_1\times V_2$: $\Phi(x,y)=0$ if and only if $y=\psi(x)$.
\end{corollary}

\begin{proof}
\emph{Step 1 (a map to invert).} The map $\Theta(x,y)=(x,\Phi(x,y))$ into $P\times Y'$ is DP at $(x_0,y_0)$. Put $\partial_P\Phi=D\Phi(x_0,y_0)(\cdot,0)$, and write $\partial_Y\Phi$ for $\partial_Y\Phi(x_0,y_0)$. The derivative of $\Theta$ at $(x_0,y_0)$, which is $(v,v')\mapsto(v,\partial_P\Phi v+\partial_Y\Phi v')$, has the bounded inverse $(v,v'')\mapsto(v,(\partial_Y\Phi)^{-1}(v''-\partial_P\Phi v))$.

\emph{Step 2 (shrinking).} Theorem~\ref{dp:thm:inverse} gives open sets $V\ni(x_0,y_0)$ and $V'\ni(x_0,0)$ and a DP inverse $\Theta^{-1}\colon V'\to V$, of the form $(x,v'')\mapsto(x,\Theta_2(x,v''))$. Choose open sets with $(x_0,y_0)\in V_0\times V_2\subseteq V$. The image $\Theta(V_0\times V_2)$ is open, as $\Theta|_V$ is a homeomorphism onto $V'$, and it contains $(x_0,0)$. So $V_1=\{x\in V_0:(x,0)\in\Theta(V_0\times V_2)\}$ is an open neighborhood of $x_0$. As $\Theta$ preserves the first coordinate, $(x,0)=\Theta(x',y)$ with $(x',y)\in V_0\times V_2$ forces $x'=x$. So $\Theta_2(x,0)\in V_2$ for $x\in V_1$.

\emph{Step 3 (the function).} Put $\psi(x)=\Theta_2(x,0)$, which is DP on $V_1$ by composition with the bounded linear map $x\mapsto(x,0)$ and the projection to $Y$. For $(x,y)\in V_1\times V_2$: $\Phi(x,y)=0$ if and only if $\Theta(x,y)=(x,0)$, if and only if $(x,y)=\Theta^{-1}(x,0)=(x,\psi(x))$, by injectivity of $\Theta$ on $V$.
\end{proof}

\begin{remark}[completeness is needed]\label{dp:rem:completeness}
Let $P=c_{00}(\N,K)$ be the space of finitely supported sequences $x=(\mathsf x_1,\mathsf x_2,\dots)$ with the supremum norm, $\sh$ the right shift, $e_1$ the first unit vector, and $f(x)=x-\mathsf x_1\sh(x)$. Then $f$ is the identity minus the diagonal of the bounded bilinear map $(x,y)\mapsto\mathsf x_1\sh(y)$, so it is analytic and DP on $P$, with $Df(0)=\id_P$. But $te_1\notin f(P)$ for every $t\ne0$. If $f(x)=te_1$, then $\mathsf x_1=t$, as $\sh(x)$ has first coordinate $0$, and so $x=te_1+t\sh(x)$. If $l$ is the largest index with $\mathsf x_l\ne0$, the right side has the nonzero coordinate $t\mathsf x_l$ at $l+1$, a contradiction. So $f$ is not open at $0$, and Theorem~\ref{dp:thm:inverse} fails for this incomplete $P$. In $c_0(\N,K)$ with $0<\abs t<1$, the equation has the solution $\mathsf x_l=t^l$.
\end{remark}

\subsection{Manifolds and bundles}\label{dp:ssec:bundles}

\begin{definition}\label{dp:def:manifold}
Let $P$ be a normed space.
\begin{enumerate}
\item A \emph{DP atlas} on a Hausdorff space $\mathcal M$, modeled on $P$, is a family of homeomorphisms $\chi_i\colon U_i\to\chi_i(U_i)$ from an open cover of $\mathcal M$ onto open subsets of $P$ whose transitions $\chi_i\circ\chi_j^{-1}$ are DP on $\chi_j(U_i\cap U_j)$. A \emph{DP manifold} is $\mathcal M$ with a maximal DP atlas. A \emph{DP map} is a continuous map between DP manifolds that is DP in charts.
\item A \emph{DP vector bundle} over $\mathcal M$ with fiber a normed space $E$ is a vector bundle with local trivializations whose transition families $x\mapsto\mathsf o_{ij}(x)\in\cL(E,E)$ are DP in charts. A \emph{morphism} is a bundle map over $\id_{\mathcal M}$ whose local operator-valued expressions are DP.
\item The \emph{tangent bundle} $T\mathcal M$ has fiber $P$ and transition families $x\mapsto D(\chi_i\circ\chi_j^{-1})(\chi_j(x))$.
\item Let $P_0\subseteq P$ be a closed subspace with a bounded linear projection $\mathsf r_0\colon P\to P_0$. A subset $\cS\subseteq\mathcal M$ is a \emph{DP submanifold} modeled on $P_0$ if every point of $\cS$ lies in a chart $\chi\colon U\to P$ with $\chi(U\cap\cS)=\chi(U)\cap P_0$.
\end{enumerate}
\end{definition}

\begin{corollary}[divided-power bundles]\label{dp:cor:bundles}
Theorem~2.1 of~\cite{I}, displayed in Section~\ref{dp:sec:conv}, holds with $C^m$ replaced by DP on open sets throughout, for the DP manifolds and DP vector bundles of Definition~\ref{dp:def:manifold}(1),(2), whose chart changes and operator-valued transition families are DP on open sets. The replacement is also made in the definitions that the theorem uses: bundle morphisms are those of Definition~\ref{dp:def:manifold}(2), and a functor is DP if all its maps on products of hom spaces are DP at every point of these products. For $k\ge1$, the alternating bifunctor $(E,F)\mapsto\Alt^k(E;F)$ and the functors $E\mapsto\Lambda^k_\pi E$ and $E\mapsto\widehat\Lambda^k_\pi E$ satisfy its hypothesis. So in every characteristic and for all normed fibers, the alternating and the projective exterior power bundles of a DP vector bundle are DP vector bundles, with DP induced morphisms.
\end{corollary}

\begin{proof}
\emph{Step 1 (the calculus).} By Theorem~\ref{dp:thm:calculus}(2), maps that are DP on open sets compose. By Theorem~\ref{dp:thm:calculus}(3) and~(1), bounded multilinear maps, in particular bounded linear maps and evaluation of operators, and constants are DP on their whole domains: a bounded $r$\nobreakdash-linear map $\beta$ on $E_1\times\dots\times E_r$ is the diagonal of the bounded $r$\nobreakdash-linear map $(v^1,\dots,v^r)\mapsto\beta(v^1_1,\dots,v^r_r)$ on the $r$th power of this product. A map into a product is DP at a point when its components are. Composition of operators is bounded bilinear, so local operator families compose. Being DP at a point depends only on the map near that point.

\emph{Step 2 (the theorem).} The proof of~\cite[Theorem~2.1]{I} uses of the class $C^m$ only the four properties (i)--(iv) listed in Section~\ref{dp:sec:conv}, and Step~1 gives them for DP maps. Independence of the trivialization multiplies by the DP families $\mathsf o_{ij}$ and $\mathsf o_{ji}$, so no inversion is needed.

\emph{Step 3 (the functors).} By Theorem~\ref{dp:thm:calculus}(4) and~(1), the maps of the three functors on hom spaces, which for the bifunctor is the joint action $B$, are bounded polynomial laws and DP on their whole domains.
\end{proof}

\begin{proposition}[categories]\label{dp:prop:categories}
\leavevmode
\begin{enumerate}
\item A homeomorphism from an open subset of $\mathcal M$ onto an open subset of $P$ is a \emph{chart}. A chart is \emph{compatible} with a DP atlas if its transitions, in both directions, with every chart of the atlas are DP. Compatibility is transitive, in the sense that two charts compatible with a DP atlas have DP transitions, so every DP atlas lies in a unique maximal DP atlas. DP manifolds and DP maps form a category. For each DP manifold $\mathcal M$, DP vector bundles over $\mathcal M$ form a category, to which Corollary~\ref{dp:cor:bundles} applies.
\item Let $f\colon\mathcal M\to\mathcal N$ be a DP map, with models $P$ and $P'$. Then $f$ is $C^\infty$ in charts. $T\mathcal M$ is a DP vector bundle and a DP manifold modeled on $P\times P$. The map $Tf$, in charts $(x,v)\mapsto(f(x),Df(x)v)$, is DP, and it is a bundle map over $f$ with DP local operator family. And $\mathcal M\mapsto T\mathcal M$ is a functor on DP manifolds.
\item For every normed space $F$ and $k\ge1$, the bundles $\Alt^k(T\mathcal M;F)$, $\Lambda^k_\pi T\mathcal M$, $\widehat\Lambda^k_\pi T\mathcal M$ and $\Alt^k(T\mathcal M;T\mathcal M)$ are DP vector bundles, in every characteristic.
\item A DP submanifold $\cS$ modeled on $P_0$ is a DP manifold modeled on $P_0$ with the restricted charts, and the inclusion $\cS\to\mathcal M$ is DP.
\end{enumerate}
\end{proposition}

\begin{proof}
\emph{Step 1 (item~(1)).} Let $\chi,\chi'$ be compatible with every chart of a DP atlas, and let $\mathsf u$ lie in both domains. For a chart $\chi_i$ of the atlas at $\mathsf u$, near $\chi(\mathsf u)$ we have $\chi'\circ\chi^{-1}=(\chi'\circ\chi_i^{-1})\circ(\chi_i\circ\chi^{-1})$, which is DP at $\chi(\mathsf u)$ by Theorem~\ref{dp:thm:calculus}(2). So the charts compatible with the atlas form a DP atlas, and it contains every DP atlas that contains the given one. The two categories exist by Step~1 of the proof of Corollary~\ref{dp:cor:bundles}: DP maps compose, identities are DP, being DP at a point depends only on the germ, and local operator families compose.

\emph{Step 2 (item~(2)).} In charts, $f$ is DP on open sets, hence $C^\infty$ with DP derivative $Df$ on open sets, by Theorem~\ref{dp:thm:differentiable}(3),(4). The transition families of $T\mathcal M$ are therefore DP into $\cL(P,P)$, after composition with the DP transitions, and they satisfy the cocycle identity by the chain rule for Fr\'echet derivatives. The map $(x,v)\mapsto(f(x),Df(x)v)$ is the composite of the DP map $(x,v)\mapsto(f(x),Df(x),v)$ with the bounded bilinear evaluation, hence DP, and its local operator family is $x\mapsto Df(x)$. Applied to the transitions of $\mathcal M$, this gives the DP atlas of $T\mathcal M$. The total space is Hausdorff: points over different base points are separated by preimages of neighborhoods in $\mathcal M$, and points of one fiber by a trivialization. The chain rule shows that $Tf$ does not depend on the charts, and that $\mathcal M\mapsto T\mathcal M$ is a functor.

\emph{Step 3 (item~(3)).} Apply Corollary~\ref{dp:cor:bundles} to the DP bundle $T\mathcal M$, and, for the bifunctor $\Alt^k$, to the pair of $T\mathcal M$ with the constant bundle of fiber $F$ and to the pair $(T\mathcal M,T\mathcal M)$, whose map on hom spaces is the joint action $B$ of Theorem~\ref{dp:thm:calculus}(4).

\emph{Step 4 (item~(4)).} Let $\chi\colon U\to P$ and $\chi'\colon U'\to P$ be charts as in Definition~\ref{dp:def:manifold}(4). As $\chi(U\cap\cS)=\chi(U)\cap P_0$ and $\chi$ is injective, $P_0\cap\chi(U\cap U')=\chi(U\cap U'\cap\cS)$, which is open in $P_0$. On it $\chi'\circ\chi^{-1}$ has values in $P_0$, so it equals $\mathsf r_0\circ\chi'\circ\chi^{-1}$. So the transition of the restricted charts is $x\mapsto\mathsf r_0(\chi'\circ\chi^{-1}(x))$, the composite of the bounded linear inclusion $P_0\to P$, a DP map and the bounded linear map $\mathsf r_0$, and it is DP. The space $\cS$ is Hausdorff, and the restricted charts are homeomorphisms onto open subsets of $P_0$. In the chart $\chi$, the inclusion of $\cS$ is the inclusion $P_0\to P$.
\end{proof}
%% ---- end sec-04-b-maps.tex
%% ---- begin sec-05-nonambient.tex
\section{Laws that are not ambient}\label{dp:sec:nonambient}\label{sec:dp}

Let $p=\operatorname{char}K>0$. In degrees below $p$ every bounded polynomial law is ambient. A law of degree $0$ is constant. Let $a\colon P\to W$ be a law of degree $n$ with $1\le n<p$, and put
\[
 \beta(h_1,\dots,h_n)=a^{[(1,\dots,1)]}(h_1,\dots,h_n)/n!.
\]
This map is $n$\nobreakdash-linear: compare coefficients in $a(s_1h_1+s_1'h_1'+s_2h_2+\dots+s_nh_n)$ after the substitutions $s_1'=s_1$, $s_1'=0$ and $s_1=0$, scale $s_1$, and do the same in every slot. It has diagonal $a$, since $a((s_1+\dots+s_n)h)=(s_1+\dots+s_n)^na(h)$ gives $a^{[(1,\dots,1)]}(h,\dots,h)=n!\,a(h)$. It is bounded by Theorem~\ref{dp:thm:mixed}(1). The algebra is \cite[Th.~II.1]{Ro}; the bound is the estimate we need.

We also use the one-point lemma \cite[Proposition~14.9(2)]{I}: a law $a\colon P\to W$ of degree $n$ such that $j\circ a$ is analytic at one point, for a linear isometry $j$ of $W$ into a normed space, is ambient of degree $n$.

In degrees at least $p$ bounded polynomial laws need not be ambient. Let $c_0=c_0(\N,K)$, and identify $K^m$ with the span of the first $m$ unit vectors, $m\ge1$. Recall that $\Gamma^p(K^m)\subseteq(K^m)^{\otimes p}$ is the space of symmetric tensors. Put $Q_m=(K^m)^{\otimes p}/\Gamma^p(K^m)$, and let $\bar v\in Q_m$ be the image of $v\in(K^m)^{\otimes p}$. Let $I_m$ be the set of the $p^{pm}$ tuples $l=(y_l^1,\dots,y_l^p)$ of vectors in $K^m$ with entries in the prime field $\F_p$, and put $y_l=y_l^1\otimes\dots\otimes y_l^p$. Let $N_m\subseteq\ell^1(I_m,K)$ consist of the vectors $(\phi(\bar y_l))_{l\in I_m}$, $\phi$ a linear form on $Q_m$. Let $G_m=\ell^1(I_m,K)/N_m$ with the quotient norm ($N_m$ is closed, being the common kernel of finitely many linear forms on $K^{I_m}$, each a combination of the bounded coordinate forms), and let $\varepsilon_l\in G_m$ be the image of the unit vector at $l$.

Let $G$ be the $\ell^1$-sum of the spaces $G_m$, $m\ge1$. We read $G_m$ as the $m$-th summand of $G$. Let $H=K\oplus c_0\oplus G$, with the maximum norm; it is a Banach space if $K$ is complete. Let $V$ be the span of the tensors $u^{\otimes p}\otimes g$, $u\in c_0$, $g\in G$, in $c_0^{\otimes p}\otimes G$, with the norm
\[
 \nu(t)=\inf\Bigl\{\sum_i\abs{c_i}\norm{u_i}^p\norm{g_i}:t=\sum_ic_i\,u_i^{\otimes p}\otimes g_i\Bigr\},
\]
which dominates the projective norm $\pi$ of $c_0^{\otimes p}\otimes G$. Both are norms, since the coordinate forms of $c_0$, and the linear forms on the spaces $G_m$ composed with the projections of $G$ onto $G_m$, are bounded and separate points.

\begin{proposition}[laws that are not ambient]\label{dp:prop:nonambient}
Assume that either (a) $K$ is discretely valued, that is, $\abs{K^\times}=\abs\vartheta^{\Z}$ for some $\vartheta\in K$ with $0<\abs\vartheta<1$, or (b) $p=2$. Put $t_m=\sum_{l\in I_m}y_l\otimes\varepsilon_l$.
\begin{enumerate}
\item $t_m\in V$ and $\pi(t_m)\le p^{pm}$; under~(a), $\nu(t_m)\ge mp^{pm}/\bigl(2^pM_{\abs\vartheta}(1+\lfloor\log_pm\rfloor)\bigr)$ with $M_{\abs\vartheta}=\max_{i\ge0}(i+1)\abs\vartheta^{\,i}$, and under~(b), $\nu(t_m)\ge m4^{m-1}$.
\item For $n\ge p+1$, $R_n(\lambda,x,g)=\lambda^{n-p-1}\,x^{\otimes p}\otimes g$, with $\lambda^0=1$, defines a bounded polynomial law $R_n\colon H\to V$ of degree $n$ and norm at most $1$ that is not ambient of degree $n$, also as a map into the completion of $V$.
\end{enumerate}
\end{proposition}
The proof is in Appendix~\ref{dp:app:proofs}.

Under~(a) or~(b) of Proposition~\ref{dp:prop:nonambient}, the same lower bounds hold for the smaller norm $\mu$ below, and this shows that composition fails on $c_0\oplus G$. From here on elements of $G$ are written $w$, since $g_X$ will be a map. As $\Gamma^p(c_0)$ is spanned by the powers $u^{\otimes p}$, we have $V=\Gamma^p(c_0)\otimes G$. Let $\mu$ be the projective norm of $\Gamma^p_\pi(c_0)\otimes G$ on $V$, as in \cite[(14.4)]{I} with $X=c_0$ and $n=p$:
\[
 \mu(t)=\inf\Bigl\{\sum_q\pi(\mathsf s_q)\norm{w_q}:t=\sum_q\mathsf s_q\otimes w_q,\ \mathsf s_q\in\Gamma^p(c_0),\ w_q\in G\Bigr\}.
\]
Here $\pi(\mathsf s_q)$ is the projective norm of $\mathsf s_q$ in $c_0^{\otimes p}$. Then $\pi\le\mu\le\nu$ on $V$, so $\mu$ is a norm, as $\pi$ is. For the first inequality expand each $\mathsf s_q$ nearly optimally in $c_0^{\otimes p}$. For the second, a representation for $\nu$ is one for $\mu$, as $\pi(c\,u^{\otimes p})\le\abs c\norm u^p$. Let $\pi_m$ be the projective norm of $(K^m)^{\otimes p}$, with the maximum norm on $K^m$. Let $y^{(m)}$ be the family $(y_l)_{l\in I_m}$ in $c_0^{\otimes p}$, in the sense of \cite[Definition~14.15]{I}, with its factors $y_l^1,\dots,y_l^p$. Let $\mathfrak l(y^{(m)})$ and $\mathfrak l_p(c_0)$ be the constants of \cite[(14.3)]{I}. A \emph{projection onto the symmetric tensors} is a linear map of $(K^m)^{\otimes p}$ into $\Gamma^p(K^m)$ that fixes $\Gamma^p(K^m)$.

\begin{lemma}[lower bounds on $c_0$]\label{dp:lem:c0-lower}
Assume~\textup{(a)} or~\textup{(b)} of Proposition~\ref{dp:prop:nonambient}, and let $m\ge1$. Put $\mathsf b_m=mp^{pm}/\bigl(2^pM_{\abs\vartheta}(1+\lfloor\log_pm\rfloor)\bigr)$ under~\textup{(a)} and $\mathsf b_m=m4^{m-1}$ under~\textup{(b)}. If both hold, the statements below hold with each value.
\begin{enumerate}
\item For every projection $\Psi_K$ of $(K^m)^{\otimes p}$ onto the symmetric tensors, $\sum_{l\in I_m}\pi_m(\Psi_Ky_l)\ge\mathsf b_m$.
\item In $c_0$, $\mathfrak l(y^{(m)})\ge\mathsf b_m/p^{pm}$. So $\mathfrak l(y^{(m)})\ge m/\bigl(2^pM_{\abs\vartheta}(1+\lfloor\log_pm\rfloor)\bigr)$ under~\textup{(a)} and $\mathfrak l(y^{(m)})\ge m/4$ under~\textup{(b)}, and $\mathfrak l_p(c_0)=\infty$.
\item $\mu(t_m)\ge\mathsf b_m$.
\end{enumerate}
\end{lemma}
The proof is in Appendix~\ref{dp:app:proofs}.

The universal space $G(y^{(m)})$ of \cite[Definition~14.15]{I} is $G_m$, since $I_m$ contains the tuples of unit vectors, so the span of the $y_l$ is $(K^m)^{\otimes p}$, and this span meets $\Gamma^p(c_0)$ in $\Gamma^p(K^m)$. So, under~(a) or~(b) of Proposition~\ref{dp:prop:nonambient}, $G$ is a bad space for $X=c_0$, $n=p$, by Lemma~\ref{dp:lem:c0-lower}(2). Let $H_0=c_0\oplus G$ and $W_X=\Gamma^p_\pi(c_0)\oplus G$, with the maximum norms, let $Y_X=(V,\mu)$, and let
\[
 f_X\colon H_0\to W_X,\quad f_X(x,w)=(x^{\otimes p},w);\qquad
 g_X\colon W_X\to Y_X,\quad g_X(\mathsf s,w)=\mathsf s\otimes w,
\]
as in \cite[(14.5)]{I}.

\begin{theorem}[composition fails on $c_0\oplus G$]\label{dp:thm:c0-failure}
Assume~\textup{(a)} or~\textup{(b)} of Proposition~\ref{dp:prop:nonambient}.
\begin{enumerate}
\item $f_X$ is AA at every point, with every radius, the bound depending on the radius and the point.
\item $g_X$ is the diagonal of a bounded bilinear map, so it is analytic everywhere.
\item $g_X\circ f_X$ is a bounded law of degree $p+1$, DP everywhere, and AA at no point, also as a map into the completion of $Y_X$.
\item $\mathfrak k_{(p,1)}(H_0)=\infty$, for the constant of \cite[Definition~14.21]{I}.
\end{enumerate}
\end{theorem}
The proof is in Appendix~\ref{dp:app:proofs}; it applies \cite[Proposition~14.18]{I}. Over densely valued $K$ of odd characteristic it is open whether $\mathfrak l_p(c_0)=\infty$ \cite[Question~14.42(2)]{I}, and it is open whether this failure occurs on $c_0\oplus G$ \cite[Remark~14.20]{I}. A failure on $c_0\oplus G$ would give $\mathfrak l_p(c_0)=\infty$, by \cite[Theorem~14.24]{I}. Over every $K$ of characteristic $p$, composition fails on another space, by \cite[Theorem~14.17]{I}.

Let $\kappa$ be the residue field of $K$. A projection of $(\kappa^m)^{\otimes p}$ onto the symmetric tensors is defined as for $K$, with the space of all symmetric tensors. The bounds under~(a) rest on Cohen's structure theorem~\cite{Co}, on \cite[Lemma~K.2(2)]{I}, and on \cite[Lemma~Q.6(4)]{I}, by which every such projection takes, on average over the tensors $y_l$, values whose $m\times m^{p-1}$ matrices of entries have large rank. The bounds under~(b) rest on \cite[Proposition~N.12]{I}, from which, for $m\ge2$, every projection $\Psi_K$ of $(K^m)^{\otimes2}$ onto the symmetric tensors has some $\Psi_K(y_l)$ whose determinant has absolute value at least $1$, and on \cite[Lemma~N.15(2)]{I}, which turns this into a lower bound for its projective norm. Condition~(a) holds under condition~\textup{(L)} of \cite[Section~7.1]{I}, with the element $\vartheta$ and the number $\varrho$ there, and then $M_{\abs\vartheta}\le M_\varrho$.

By Theorem~\ref{dp:thm:calculus}(1), $R_n$ is DP on $H$. Under~(a) or~(b) of Proposition~\ref{dp:prop:nonambient}, by item~(2) there and the one-point lemma, $jR_n$ is analytic at no point for every linear isometry $j$ of $V$, or of its completion, into a normed space.

In degree $p$ itself bounded polynomial laws need not be ambient either, over every $K$ of characteristic $p$. Recall the space $\Gamma^p(P)$ and the seminorm $\norm\cdot_\Gamma$ of~\eqref{dp:eq:gamma-norm} with $n=p$, which dominates the projective seminorm $\pi$ of $P^{\otimes p}$; both are norms if the bounded linear forms on $P$ separate points. For integers $m\ge1$ and $d\ge0$ let $\N^m_d$ be the set of $\alpha\in\N^m$ with $\abs\alpha=d$, and let $e_1,\dots,e_m$ be the unit vectors of $K^m$. For a symmetric tensor $T\in(\F_p^m)^{\otimes p}$ and $\alpha\in\N^m_p$ let $T_\alpha$ be its coefficient at $e_{i_1}\otimes\dots\otimes e_{i_p}$ for any $(i_1,\dots,i_p)$ in which each $i$ occurs $\alpha_i$ times. For $\gamma\in\N^m_{p-1}$ let the \emph{slice} $\partial^\gamma T\in\F_p^m$ be the vector whose $i$-th entry is $T_{\gamma^{+i}}$, where $\gamma^{+i}\in\N^m_p$ is $\gamma$ with its $i$-th entry raised by $1$. For a real number $\tau\ge1$ let $\mathcal E_{T,\tau}$ be $K^m$ with the norm
\begin{equation}\label{dp:eq:weighted-norm}
 \norm h_{T,\tau}=\inf\Bigl\{\sum_{i=1}^m\abs{z_i}+\tau\sum_{\gamma\in\N^m_{p-1}}\abs{z_\gamma}:
 h=\sum_{i=1}^mz_ie_i+\sum_{\gamma\in\N^m_{p-1}}z_\gamma\,\partial^\gamma T\Bigr\},
\end{equation}
the infimum over the scalars $z_i,z_\gamma\in K$. It is a norm, and the coordinate forms of $\mathcal E_{T,\tau}$ are bounded: as the entries of $T$ lie in $\F_p$ and $\tau\ge1$, no entry of $h$ has absolute value above $\norm h_{T,\tau}$. The tensor $T$ lies in $\Gamma^p(\mathcal E_{T,\tau})$.

\begin{proposition}[laws that are not ambient, over every $K$]\label{dp:prop:degree-p}
Let $p=\operatorname{char}K>0$.
\begin{enumerate}
\item For every $m\ge1$ there are a nonzero symmetric tensor $T\in(\F_p^m)^{\otimes p}$ and a real number $\tau\ge1$ such that, in $\Gamma^p(\mathcal E_{T,\tau})$,
\begin{equation}\label{dp:eq:degree-p-ratio}
 \norm T_\Gamma\ge\frac{(p-1)!}{3p^{3/2}(2p-1)!}\Bigl(\frac m{1+\log_pm}\Bigr)^{1/2}\pi(T).
\end{equation}
\item Choose such $T^{(m)}$ and $\tau^{(m)}$ for every $m\ge1$, and let $\mathcal E$ be the $\ell^1$-sum of the spaces $\mathcal E_m=\mathcal E_{T^{(m)},\tau^{(m)}}$, $m\ge1$; it is a Banach space if $K$ is complete. For $n\ge p$, the map $(\lambda,h)\mapsto\lambda^{n-p}h^{\otimes p}$, with $\lambda^0=1$, from $K\oplus\mathcal E$, with the maximum norm, to $\Gamma^p(\mathcal E)$, with the norm $\norm\cdot_\Gamma$ of~\eqref{dp:eq:gamma-norm}, is a bounded polynomial law of degree $n$ and norm at most $1$ that is not ambient of degree $n$, also as a map into the completion of $\Gamma^p(\mathcal E)$.
\end{enumerate}
\end{proposition}
The proof is in Appendix~\ref{dp:app:proofs}.

A representation of $T$ in $\Gamma^p(\mathcal E_{T,\tau})$ gives a symmetric tensor over $\F_p$ in $m+\binom{m+p-2}{p-1}$ symbols, which stand for the vectors $e_i$ and $\partial^\gamma T$. It maps to $T$, and its nonzero coefficients have total weight at most the cost $\sum_i\abs{c_i}\norm{h_i}^p$ of the representation; a coefficient at two or more slices weighs at least $\tau^2$, the others at least $1$. The coefficients at exactly one slice act on $T$ through a linear map $\mathsf L^\sharp$ of Lemma~\ref{dp:lem:linear}, an identity for symmetric tensors in characteristic $p$ by which $\id-\mathsf L^\sharp$ has rank at least $(p-1)!\,m^p/(2p-1)!$. So for most $T$ every such tensor has many nonzero coefficients, or many at two or more slices. By Theorem~\ref{dp:thm:calculus}(1) the law of~(2) is DP on $K\oplus\mathcal E$, while, as for $R_n$, its composite with a linear isometry of $\Gamma^p(\mathcal E)$, or of its completion, into a normed space is analytic at no point. For $n\ge p+1$ this gives laws with the properties of Proposition~\ref{dp:prop:nonambient}(2) over every $K$, on another space.

For $h_1,\dots,h_p\in P$ the tensor $\sum_{\sigma\in S_p}h_{\sigma(1)}\otimes\dots\otimes h_{\sigma(p)}$ lies in $\Gamma^p(P)$, by \cite[(4.2)]{I}; the elements of the span of these tensors are \emph{hollow}. For $t\in\Gamma^p(P)$ let
\[
 \norm t_{\mathrm F}=\inf\{\norm{t'}_\Gamma:t'\in\Gamma^p(P),\ t-t'\text{ hollow}\},
\]
the quotient seminorm of $\norm\cdot_\Gamma$ modulo the hollow tensors. For $C\ge0$, $P$ satisfies \emph{$(\mathrm N_P)$ with constant $C$} if $\norm t_\Gamma\le C\pi(t)$ for every hollow $t$, and \emph{$(\mathrm F_P)$ with constant $C$} if $\norm t_{\mathrm F}\le C\pi(t)$ for every $t\in\Gamma^p(P)$. A \emph{Frobenius law} is a map $a\colon P\to W$ with $a(h+h')=a(h)+a(h')$ and $a(\lambda h)=\lambda^pa(h)$ for $h,h'\in P$ and $\lambda\in K$, and with $\norm a=\sup_{h\ne0}\norm{a(h)}/\norm h^p$ finite; it is a bounded polynomial law of degree $p$, as $a(s_1h_1+\dots+s_rh_r)=s_1^pa(h_1)+\dots+s_r^pa(h_r)$.

\begin{theorem}[laws of degree $p$]\label{dp:thm:frobenius}
Let $p=\operatorname{char}K>0$.
\begin{enumerate}
\item $P$ satisfies $(\mathrm F_P)$ with constant $p$, and with constant $\dim P$ if $\dim P<p$.
\item Every Frobenius law $a\colon P\to W$ is ambient of degree $p$, with a $p$\nobreakdash-linear map $\beta$ of norm at most $p\norm a$ in Definition~\ref{dp:def:laws}(2).
\item If $P$ satisfies $(\mathrm N_P)$ with constant $C$, then $\norm t_\Gamma\le(p+(p+1)C)\pi(t)$ for every $t\in\Gamma^p(P)$, and every bounded polynomial law $a\colon P\to W$ of degree $p$ is ambient of degree $p$, with a $p$\nobreakdash-linear map $\beta$ of norm at most $(p+(p+1)C)\norm a$. Conversely, if every bounded polynomial law of degree $p$ on $P$, with values in any normed space, is ambient of degree $p$, then $P$ satisfies $(\mathrm N_P)$ with some constant.
\end{enumerate}
\end{theorem}
The proof is in Appendix~\ref{dp:app:proofs}. It rests on Lemma~\ref{dp:lem:bases}: if $t\in\Gamma^p(P)$ is a sum of elementary tensors with nonzero factors, some of these factors form a basis of the span of all of them for which $t$ is congruent, modulo the hollow tensors, to a combination of the $p$-th powers of the basis vectors in which the power of each basis vector $h$, taken as a factor of a term $h_1\otimes\dots\otimes h_p$, has a coefficient of absolute value at most $\norm{h_1}\cdots\norm{h_p}/\norm h^p$. By~(3), the space $K\oplus\mathcal E$ of Proposition~\ref{dp:prop:degree-p}(2) does not satisfy $(\mathrm N_{K\oplus\mathcal E})$ with any constant, as the law there with $n=p$ is not ambient of degree $p$.

From here through Proposition~\ref{dp:prop:all-norms} the characteristic of $K$ is arbitrary. For $n\ge1$ recall that $\mathfrak a_n(P)\in[0,\infty]$ is the infimum of the $C\ge0$ with $\norm t_\Gamma\le C\pi(t)$ for all $t\in\Gamma^n(P)$, as in \cite[Definition~14.36]{I}; if it is finite, it is the least such $C$. For $m\ge1$ let $[m]=\{1,\dots,m\}$, and let $c_0([m],K)$ be $K^m$ with the maximum norm, read as the span of the first $m$ unit vectors of $c_0(\N,K)$. Recall that $\mathfrak c_K=\abs\vartheta$ if $K$ is discretely valued, with $\vartheta$ as in Proposition~\ref{dp:prop:nonambient}(a), and $\mathfrak c_K=1$ otherwise.

\begin{proposition}[ambience constants]\label{dp:prop:ultrametric}
Let $n\ge1$.
\begin{enumerate}
\item If $\mathfrak a_n(P)<\infty$, then every bounded polynomial law $a\colon P\to W$ of degree $n$ is ambient of degree $n$, with an $n$\nobreakdash-linear map $\beta$ of norm at most $\mathfrak a_n(P)\norm a$ in Definition~\ref{dp:def:laws}(2). Conversely, if every bounded polynomial law of degree $n$ on $P$, with values in any normed space, is ambient of degree $n$, then $\mathfrak a_n(P)<\infty$.
\item $\mathfrak a_n(P)\le(n+1)\,\mathfrak a_{n+1}(K\oplus P)$, with the maximum norm on $K\oplus P$.
\item If $P$ is ultrametric, then $\mathfrak a_n(P)\le\mathfrak c_K^{-n}\sup_{m\ge1}\mathfrak a_n(c_0([m],K))\le\mathfrak c_K^{-n}\,\mathfrak a_n(c_0(\N,K))$.
\item If $P$ is ultrametric, then $\mathfrak a_2(P)\le\frac32\,\mathfrak c_K^{-2}$.
\end{enumerate}
\end{proposition}
The proof is in Appendix~\ref{dp:app:proofs}. By~(1) and~(3), if every bounded polynomial law of degree $n$ on $c_0(\N,K)$ is ambient of degree $n$, then so is every bounded polynomial law of degree $n$ on every ultrametric normed space. By~(1) and~(4), every bounded polynomial law of degree $2$ on an ultrametric normed space is ambient of degree $2$.

\begin{proposition}[a bound for all norms]\label{dp:prop:all-norms}
Let $n\ge1$.
\begin{enumerate}
\item If $n!\ne0$ in $K$, then $\mathfrak a_n(P)\le(2n/\abs\varpi)^n/\abs{n!}$ for every normed space $P$, and $\mathfrak a_n(P)\le(2/\abs\varpi)^n/\abs{n!}$ if $P$ is ultrametric.
\item If $n\ge p=\operatorname{char}K>0$, then $\mathfrak a_n(K\oplus\mathcal E)=\infty$, for the space $\mathcal E$ of Proposition~\ref{dp:prop:degree-p}(2).
\end{enumerate}
So $\sup_P\mathfrak a_n(P)<\infty$ if and only if $n!\ne0$ in $K$.
\end{proposition}
The proof is in Appendix~\ref{dp:app:proofs}.

\begin{remark}[the norms $\mu$ and $\nu$]\label{dp:rem:nu-mu}
Let $p=\operatorname{char}K=2$. Then $\nu\le\mathfrak a_2(c_0)\mu\le\frac32\mathfrak c_K^{-2}\mu$ on $V$. Indeed, let $t=\sum_q\mathsf s_q\otimes w_q$ with $\mathsf s_q\in\Gamma^2(c_0)$ and $w_q\in G$. Representing each $\mathsf s_q$ with cost close to $\norm{\mathsf s_q}_\Gamma\le\mathfrak a_2(c_0)\pi(\mathsf s_q)$ gives a representation for $\nu$, and $\mathfrak a_2(c_0)\le\frac32\mathfrak c_K^{-2}$ by Proposition~\ref{dp:prop:ultrametric}(4). So $\mu$ and $\nu$ are equivalent. Let $\operatorname{pr}\colon H\to H_0$ be the coordinate projection. By \cite[Proposition~14.12(1)(b)]{I}, $f_X\circ\operatorname{pr}$ is AA everywhere, and $R_3=g_X\circ f_X\circ\operatorname{pr}$ as maps into the vector space $V$. So for $p=2$ the law $R_3$ is, up to an equivalent target norm, an analytic map after an AA map.
\end{remark}

Let $\mathrm{An}$, $\mathrm{AA}$ and $\mathrm{DP}$ be the classes of germs that are analytic, AA and DP. Then $\mathrm{An}\subseteq\mathrm{AA}$ by definition, and $\mathrm{AA}\subseteq\mathrm{DP}$ by \cite[Proposition~14.9(1)]{I}. Both inclusions are strict in characteristic $p$. The first is also strict between ultrametric spaces, for $K$ complete and not spherically complete.

\begin{proposition}[characteristic $p$: both inclusions are strict]\label{dp:prop:strict}
Let $p=\operatorname{char}K>0$ and $k\ge p$.
\begin{enumerate}
\item Each of the following maps is AA at every point and analytic at no point.
\begin{enumerate}
\item $A^k_{E,E;F}$ on $\cL(E,E)$, for the Banach spaces $E,F$ of \cite[Theorem~7.2(1)]{I}.
\item If $p=k=2$, or if $K$ satisfies condition~\textup{(L)} of \cite[Section~7.1]{I}: the power map $c_0\to\hGamma^k_\pi c_0$, $h\mapsto h^{\otimes k}$.
\item If $K$ is complete and not spherically complete: $A^k_{E,D;K}$ on $\cL(E,D)$, for the ultrametric Banach spaces $E,D$ of \cite[Theorem~7.2(2)]{I}. Its source $\cL(E,D)$ and its target $\cL(\Alt^k(D;K),\Alt^k(E;K))$ are ultrametric Banach spaces.
\item If $K$ is discretely valued or $p=2$: the power map $\operatorname{pw}\colon c_0\to\Gamma^p_\pi(c_0)$, $\operatorname{pw}(h)=h^{\otimes p}$, also as a map into the completion of $\Gamma^p_\pi(c_0)$.
\end{enumerate}
\item If $K$ is discretely valued or $p=2$, the laws of Proposition~\ref{dp:prop:nonambient}(2), and over every $K$ of characteristic $p$ the laws of Proposition~\ref{dp:prop:degree-p}(2), are DP at every point and AA at no point.
\end{enumerate}
\end{proposition}

\begin{proof}
\emph{Step 1 (items (1)(a) and (1)(c)).} By \cite[Proposition~2.2]{I}, $A^k_{E,D;F}$ followed by $j_{E,F}\circ{-}$ is the diagonal of a bounded $k$-linear map, read on alternating inputs, and $j_{E,F}\circ{-}$ is a linear isometry whose range, the operators with values in $\Alt^k(E;F)$, is closed. Expanding by multilinearity gives a finite expansion at every point. The condition on diagonals holds as the range is closed. So $A^k_{E,D;F}$ is AA at every point. It is analytic at no point in~(1)(a) and~(1)(c), by \cite[Theorem~7.2]{I}. In~(1)(c), an operator norm with ultrametric target is ultrametric, and $\Alt^k(D;K)$ and $\Alt^k(E;K)$ are ultrametric, as their values lie in $K$. These spaces are Banach, as $K$ and $D$ are complete.

\emph{Step 2 (item (1)(b)).} The power map is the diagonal of the bounded $k$-linear map $(x_1,\dots,x_k)\mapsto x_1\otimes\dots\otimes x_k$ into the completion of $(c_0^{\otimes k},\pi)$. The subspace $\hGamma^k_\pi c_0$ of this completion is closed, by \cite[Definition~3.8]{I}, so the condition on diagonals holds as in Step~1. Expanding by multilinearity shows that the power map is AA at every point. By \cite[Proposition~2.5]{I}, it is analytic at some point if and only if it is the diagonal of a bounded $k$-linear map into $\hGamma^k_\pi c_0$. By \cite[Lemma~3.11]{I}, this holds if and only if $c_0$ is polarizable in degree $k$. That fails for $p=k=2$ by \cite[Theorem~7.20(4)]{I}, and under~\textup{(L)} by \cite[Corollary~8.3]{I}.

\emph{Step 3 (item (1)(d)).} The map $\operatorname{pw}$ is $f_X$ on the summand $c_0$ of $H_0$, followed by the projection of $W_X$ onto $\Gamma^p_\pi(c_0)$. So it is AA at every point, by Theorem~\ref{dp:thm:c0-failure}(1) and \cite[Proposition~14.12(1)]{I}, also into the completion. The completion of $\Gamma^p_\pi(c_0)$ is $\hGamma^p_\pi c_0$ \cite[Section~Q.4]{I}. Suppose $\operatorname{pw}$ is analytic at $u_0$ as a map into $\hGamma^p_\pi c_0$; this covers the map into $\Gamma^p_\pi(c_0)$. Write $\operatorname{pw}(u_0+u)=\sum_{i\ge0}\beta_i(u,\dots,u)$ for $\norm u<R$, with $\norm{\beta_i}\le MR^{-i}$. The bilinear map $(\mathsf s,w)\mapsto\mathsf s\otimes w$ has norm at most $1$ for $\mu$, so it extends to a bilinear map $\boxtimes$ of norm at most $1$ from $\hGamma^p_\pi c_0\times G$ to the completion of $Y_X$. Fix $w_0\in G$. For $h=(u,w)$ in $H_0$ with $\norm h<R$,
\[
 (g_X\circ f_X)\bigl((u_0,w_0)+h\bigr)=\sum_{i\ge0}\bigl(\beta_i(u,\dots,u)\boxtimes w_0+\beta_i(u,\dots,u)\boxtimes w\bigr).
\]
This series converges absolutely, so it may be grouped by degree. The term of degree $n$ is the diagonal of the $n$-linear map $(u_r,w_r)_{1\le r\le n}\mapsto\beta_n(u_1,\dots,u_n)\boxtimes w_0+\beta_{n-1}(u_1,\dots,u_{n-1})\boxtimes w_n$, of norm at most $M(\norm{w_0}+R)R^{-n}$; for $n=0$ only the first term occurs. So $g_X\circ f_X$ is analytic at $(u_0,w_0)$ as a map into the completion of $Y_X$, against Theorem~\ref{dp:thm:c0-failure}(3).

\emph{Step 4 (item (2)).} Let $a$ be one of these laws, of degree $n$. It is DP at every point by Theorem~\ref{dp:thm:calculus}(1). If $a$ were AA at a point, then $j\circ a$ would be analytic there for a linear isometry $j$ of its target. By the one-point lemma, $a$ would be ambient of degree $n$, against Propositions~\ref{dp:prop:nonambient}(2) and~\ref{dp:prop:degree-p}(2).
\end{proof}
%% ---- end sec-05-nonambient.tex
%% ---- begin sec-06-descent.tex
\section{Descent and minimality}\label{dp:sec:descent}

We prove that divided-power analyticity descends along isometries of the target and along surjections of the source. Then we write each divided-power germ as an analytic germ after one enlargement of the target and one cover of the source. It follows that the divided-power germs are the smallest class containing the analytic germs with these two descents (Theorem~\ref{dp:thm:minimal}). Between ultrametric spaces source descent alone suffices (Theorem~\ref{dp:thm:ultrametric-minimal}). The section ends with what this minimality depends on (Example~\ref{dp:ex:l1}).

The analytic and the AA germs are as in \cite[Definition~14.7]{I}, with $M$ for the bound. Recall $C_1=8/\abs\varpi^2$ and $C_0=\mathrm eC_1$.

\subsection{The two descents}\label{dp:ssec:two-descents}

\begin{theorem}[target descent]\label{dp:thm:target-descent}
Let $j\colon W\to Z$ be a linear isometry of normed spaces, $U\subseteq P$ open, $f\colon U\to W$ and $x_0\in U$. Suppose that $j\circ f$ is DP at $x_0$ with radius $R$, bound $M$ and coefficient laws $a'_n\colon P\to Z$.
\begin{enumerate}
\item If every $a'_n$ takes values in $j(W)$, then $f$ is DP at $x_0$ with radius $R$ and bound $M$, with coefficient laws $a_n=j^{-1}\circ a'_n$ and $\norm{a_n}=\norm{a'_n}$.
\item If $j(W)$ is closed in $Z$, for instance if $W$ is complete, the hypothesis of~(1) holds. So divided-power analyticity, at a point or on an open set, descends along linear isometries with closed range, with the same radius and bound.
\end{enumerate}
\end{theorem}

\begin{proof}
\emph{Step 1 (item (1)).} Let $h_1,\dots,h_r\in P$ and $\abs\alpha=n$. The vector $(a'_n)^{[\alpha]}(h_1,\dots,h_r)$ is a coefficient of the polynomial $\lambda\mapsto a'_n(\sum_i\lambda_ih_i)$ on $K^r$, whose values lie in $j(W)$. As $K$ is infinite, interpolation shows that this coefficient lies in $j(W)$. So $a_n(\sum_i\lambda_ih_i)=\sum_\alpha\lambda^\alpha j^{-1}(a'_n)^{[\alpha]}(h_1,\dots,h_r)$ is a homogeneous polynomial with coefficients in $W$, and $\norm{a_n}=\norm{a'_n}\le MR^{-n}$ because $j$ is isometric. For $\norm h<R$ the partial sums $\sum_{n\le N}a_n(h)\in W$ have images converging to $jf(x_0+h)$, so they converge to $f(x_0+h)$.

\emph{Step 2 (item (2)).} For $h=0$ the values $a'_n(0)$ are $jf(x_0)$ and $0$. Fix $h\ne0$ and put $\rho=R/\norm h$. For $\abs s<\rho$ we have $\psi(s)=jf(x_0+sh)=\sum_ns^na'_n(h)$, with $\norm{a'_n(h)}\rho^n\le M$. We show $a'_n(h)\in j(W)$ by induction on $n$. First, $a'_0(h)=jf(x_0)$. If $a'_m(h)\in j(W)$ for $m<n$, then for $0<\abs s<\rho$,
\begin{gather*}
 \delta_n(s)=s^{-n}\Bigl(\psi(s)-\sum_{m<n}s^ma'_m(h)\Bigr)\in j(W),\\
 \norm{\delta_n(s)-a'_n(h)}\le\sum_{m>n}M\rho^{-m}\abs s^{m-n}=M\rho^{-n}\frac{\abs s/\rho}{1-\abs s/\rho}.
\end{gather*}
As $K$ is nontrivially valued, $s$ can tend to $0$, so $a'_n(h)$ lies in the closure of $j(W)$, which is $j(W)$. If $W$ is complete, then $j(W)$ is complete, hence closed.
\end{proof}

Without closed range, target descent can fail at every point.

\begin{example}[a dense isometry]\label{dp:ex:dense-isometry}
Let $K$ be complete. Let $Z=c_0(\N,K)$ and $\Delta=\{y\in K:\abs y<1\}$. Put $\mathsf F(y)=(y^i)_{i\ge0}\in Z$, with $0^0=1$. Let $W$ be the linear span of $\mathsf F(\Delta)$, with the norm of $Z$, and let $j\colon W\to Z$ be the inclusion. Then:
\begin{enumerate}
\item $j\circ\mathsf F$ is analytic at every point of $\Delta$, with radius $1$ and bound $1$. So is $\mathsf F$ as a map into the closure $\overline W$ of $W$ in $Z$, which is a completion of $W$.
\item $\mathsf F\colon\Delta\to W$ is DP at no point of $\Delta$, and it is differentiable at no point as a $W$-valued map.
\end{enumerate}
So in \cite[Definition~14.7]{I} the condition on diagonals cannot be dropped: the maps $\beta_n$ of the proof expand $j\circ\mathsf F$ with no condition on their diagonals, while $\mathsf F$ is AA at no point, as AA germs are DP \cite[Proposition~14.9(1)]{I}.
\end{example}

\begin{proof}
\emph{Step 1 (item (1)).} Let $y_0\in\Delta$. For $\abs s<1$ we have $\abs{y_0+s}<1$ and, coordinatewise, $(y_0+s)^i=\sum_{n\le i}\binom iny_0^{i-n}s^n$. Put $z_n=(\binom iny_0^{i-n})_{i\ge0}$, with entry $0$ for $i<n$. Then $z_n\in Z$ and $\norm{z_n}\le1$, as $\abs{\binom in}\le1$ and $\abs{y_0}^{i-n}\to0$ as $i\to\infty$. Since $Z$ is complete, $\sum_ns^nz_n$ converges in $Z$, and coordinatewise it converges to $\mathsf F(y_0+s)$. So $j\mathsf F(y_0+s)=\sum_n\beta_n(s,\dots,s)$ with the $n$-linear maps $\beta_n(s_1,\dots,s_n)=s_1\cdots s_nz_n$, of norm at most $1$. Now apply Theorem~\ref{dp:thm:target-descent}(2) to the inclusion of the closed subspace $\overline W$ in $Z$, with Lemma~\ref{dp:lem:coefficients-unique}(2). This gives $z_n=\beta_n(1,\dots,1)\in\overline W$, so the $\beta_n$ expand $\mathsf F$ as a map into $\overline W$.

\emph{Step 2 (the vector that leaves $W$).} Suppose $\mathsf F$ is DP into $W$ at $y_0$, with laws $a_n$. By Lemma~\ref{dp:lem:coefficients-unique}(2), $ja_n(s)=\beta_n(s,\dots,s)$. So $z_1=(\mathsf d_i)_{i\ge0}$, with $\mathsf d_i=iy_0^{i-1}$ and $\mathsf d_0=0$, lies in $W$. The vector $z_1$ is also the limit in $Z$ of the difference quotients $s^{-1}(\mathsf F(y_0+s)-\mathsf F(y_0))$, so a derivative of $\mathsf F$ in $W$ at $y_0$ would also force $z_1\in W$.

\emph{Step 3 ($z_1\notin W$).} Let $\sh$ be the shift, $(\sh\mathsf t)_i=\mathsf t_{i+1}$, on sequences $\mathsf t$. The identity $(i+2)y^{i+1}-2y(i+1)y^i+y^2iy^{i-1}=0$ gives $(\sh-y_0)^2z_1=0$, while $((\sh-y_0)z_1)_0=\mathsf d_1-y_0\mathsf d_0=1$. Suppose $z_1=\sum_{l\le r}\lambda_l\mathsf F(y_l)$ with distinct $y_l\in\Delta$. Since $(\sh-y_l)\mathsf F(y_l)=0$, also for $y_l=0$, we get $\prod_l(\sh-y_l)z_1=0$. The polynomials $\mathsf g\in K[X]$ with $\mathsf g(\sh)z_1=0$ form an ideal, and its monic generator divides both $(X-y_0)^2$ and the squarefree $\prod_l(X-y_l)$. So it divides $X-y_0$, and $(\sh-y_0)z_1=0$. This is a contradiction.
\end{proof}

\begin{corollary}[values in the completion]\label{dp:cor:completion}
Let $U\subseteq P$ be open, $f\colon U\to W$ and $x_0\in U$. Then $f$ is DP at $x_0$ as a $W$-valued map if and only if it is DP at $x_0$ as a $\widehat W$-valued map and, for every $h\in P$ and $n\ge0$, the coefficient of $s^n$ in $s\mapsto f(x_0+sh)$, computed in $\widehat W$, lies in $W$. Radii and bounds agree. The second condition can fail at every point of an open set on which $f$ is analytic into $\widehat W$ (Example~\ref{dp:ex:dense-isometry}).
\end{corollary}

\begin{proof}
Apply Theorem~\ref{dp:thm:target-descent}(1) with $j\colon W\to\widehat W$, and Lemma~\ref{dp:lem:coefficients-unique}(1). The converse is immediate.
\end{proof}

Recall from Lemma~\ref{dp:lem:strong-formal}(2) that a bounded linear surjection $Q\colon P'\to P$ has \emph{preimage constant} $\mathsf c>0$ if every $h\in P$ has a preimage $\tilde h$ with $\norm{\tilde h}\le\mathsf c\norm h$.

\begin{proposition}[source descent]\label{dp:prop:source-descent}
Let $Q\colon P'\to P$ be a bounded linear surjection with preimage constant $\mathsf c$. Let $U\subseteq P$ be open, $f\colon U\to W$, $\tilde x_0\in Q^{-1}(U)$ and $x_0=Q(\tilde x_0)$.
\begin{enumerate}
\item If $f\circ Q$ is DP at $\tilde x_0$ with radius $R'$, bound $M$ and coefficient laws $b_n$, then $b_n=a_n\circ Q$ for bounded laws $a_n\colon P\to W$ with $\norm{a_n}\le\mathsf c^n\norm{b_n}$, and $f$ is DP at $x_0$ with radius $R'/\mathsf c$, bound $M$ and coefficient laws $a_n$.
\item If $f$ is DP at $x_0$ with radius $R$, bound $M$ and coefficient laws $a_n$, then $f\circ Q$ is DP at $\tilde x_0$ with radius $R/\norm Q$ (any radius if $Q=0$), bound $M$ and coefficient laws $a_n\circ Q$.
\end{enumerate}
\end{proposition}

\begin{proof}
\emph{Step 1 ($b_n$ is constant on fibers).} Let $\tilde h\in P'$ and $\mathsf k\in\ker Q$. For small $s$ we have $f\circ Q(\tilde x_0+s(\tilde h+\mathsf k))=f\circ Q(\tilde x_0+s\tilde h)$. Both sides are power series in $s$ with geometric bounds, with coefficients $b_n(\tilde h+\mathsf k)$ and $b_n(\tilde h)$. By Lemma~\ref{dp:lem:coefficients-unique}(1), $b_n(\tilde h+\mathsf k)=b_n(\tilde h)$. So $a_n(h)=b_n(\tilde h)$, for any preimage $\tilde h$ of $h$, is well defined, and $b_n=a_n\circ Q$.

\emph{Step 2 ($a_n$ is a bounded law).} For $h_1,\dots,h_r\in P$ with preimages $\tilde h_i$, the vector $\sum_i\lambda_i\tilde h_i$ is a preimage of $\sum_i\lambda_ih_i$. So $a_n(\sum_i\lambda_ih_i)=b_n(\sum_i\lambda_i\tilde h_i)$ is a homogeneous polynomial of degree $n$ with coefficients in $W$. Taking $\norm{\tilde h}\le\mathsf c\norm h$ gives $\norm{a_n(h)}\le\norm{b_n}\mathsf c^n\norm h^n$.

\emph{Step 3 (expansion).} For $\norm h<R'/\mathsf c$ take $\tilde h$ with $\norm{\tilde h}\le\mathsf c\norm h<R'$. Then $x_0+h\in U$ and $f(x_0+h)=f\circ Q(\tilde x_0+\tilde h)=\sum_nb_n(\tilde h)=\sum_na_n(h)$. Also $\norm{a_n}\le\mathsf c^nMR'^{-n}=M(R'/\mathsf c)^{-n}$.

\emph{Step 4 (item (2)).} The map $a_n\circ Q$ is a bounded law with $\norm{a_n\circ Q}\le\norm Q^n\norm{a_n}$, and $f\circ Q(\tilde x_0+\tilde h)=\sum_na_n(Q\tilde h)$ for $\norm Q\norm{\tilde h}<R$.
\end{proof}

\subsection{The canonical cover}\label{dp:ssec:cover}

\begin{definition}[canonical cover]\label{dp:def:cover}
Recall $\ell^1(J)$ and $\ell^1_{00}(J)$ from Section~\ref{dp:sec:conv}, and write $e_\theta$ for the unit vectors. For a normed space $P$ let $I_P=\{h\in P:\abs\varpi<\norm h\le1\}$, and let $q_P\colon\ell^1_{00}(I_P)\to P$ be the linear map with $q_P(e_h)=h$. If $K$ and $P$ are complete, let $\hat q_P\colon\ell^1(I_P)\to P$, $\mathsf v\mapsto\sum_{h\in I_P}\mathsf v_hh$.
\end{definition}

The maps $q_P$, $\hat q_P$ are contractive, as $\norm{\sum_h\mathsf v_hh}\le\sum_h\abs{\mathsf v_h}\norm h\le\norm{\mathsf v}$. Each has preimage constant $\abs\varpi^{-1}$. Indeed, for $h\ne0$ choose $\lambda\in K$ with $\abs\varpi<\abs\lambda\norm h\le1$. Then $\lambda h\in I_P$, $h=q_P(\lambda^{-1}e_{\lambda h})$ and $\abs\lambda^{-1}<\abs\varpi^{-1}\norm h$.

For a set $J$ with a fixed total order, and a law $b\colon\ell^1_{00}(J)\to W$ of degree $n$, we use the following construction. For a nondecreasing word $\mathsf w=(\mathsf w_1\le\dots\le\mathsf w_n)$ in $J$, with distinct letters $\theta_1<\dots<\theta_r$ occurring $\alpha_1,\dots,\alpha_r$ times, put $v_{\mathsf w}=b^{[\alpha]}(e_{\theta_1},\dots,e_{\theta_r})$. For finitely supported $h^1,\dots,h^n$ put
\begin{equation}\label{dp:eq:word-lift}
 \beta_b(h^1,\dots,h^n)=\sum_{\mathsf w}h^1_{\mathsf w_1}\cdots h^n_{\mathsf w_n}\,v_{\mathsf w},
\end{equation}
a finite sum over nondecreasing words. This map is $n$-linear. Let $h$ be supported in a finite set $J_0\subseteq J$. By Lemma~\ref{dp:lem:elementary}(3), $b(h)$ is the sum over $\alpha\in\N^{J_0}$ with $\abs\alpha=n$ of $h^\alpha b^{[\alpha]}((e_\theta)_{\theta\in J_0})$, and arguments with $\alpha_\theta=0$ can be deleted. Each monomial $h^\alpha$ comes from exactly one nondecreasing word. So $\beta_b(h,\dots,h)=b(h)$. For $n=0$, $\beta_b$ is the constant $b$.

\begin{lemma}[laws on $\ell^1_{00}$ are diagonals]\label{dp:lem:laws-cover}
Let $W$ be normed, and fix a total order on $J$. Every bounded law $b\colon\ell^1_{00}(J)\to W$ of degree $n$ is the diagonal of the $n$-linear map $\beta_b\colon\ell^1_{00}(J)^n\to W$, with $\norm{\beta_b}\le C_0^n\,n!\,\norm b$. So $b$ is analytic everywhere.
\end{lemma}

\begin{proof}
By Theorem~\ref{dp:thm:mixed}(3), $\norm{v_{\mathsf w}}\le C_0^n\norm b\,n!/\prod_c\alpha_c!\le C_0^nn!\norm b$. Also $\sum_{\mathsf w}\prod_c\abs{h^c_{\mathsf w_c}}\le\prod_c\norm{h^c}$. The last assertion is \cite[Proposition~2.5]{I}, with $j=\id_W$.
\end{proof}

The factor $n!$ is not geometric. Removing it is what the enlargement in Theorem~\ref{dp:thm:normal-form} is for. Recall $\operatorname{lift}(a)$ from Definition~\ref{dp:def:strong-norm}(1).

\begin{lemma}[lifts after the cover]\label{dp:lem:lift-cover}
Let $W$ be normed and let $a\colon P\to W$ be a law of degree $n$. Fix a total order on $I_P$. The map $\beta_{a\circ q_P}$ of~\eqref{dp:eq:word-lift}, for $J=I_P$, has norm at most $\lvert\!\lvert\!\lvert a\rvert\!\rvert\!\rvert$. So $\operatorname{lift}(a\circ q_P)\le\lvert\!\lvert\!\lvert a\rvert\!\rvert\!\rvert$.
\end{lemma}

\begin{proof}
Let $b=a\circ q_P$, and take $\beta_b$ of~\eqref{dp:eq:word-lift} for $J=I_P$. For a word with letters $\xi_1<\dots<\xi_r$ in $I_P$ and multiplicities $\alpha$ we have $v_{\mathsf w}=a^{[\alpha]}(\xi_1,\dots,\xi_r)$, as $q_P$ is linear. As $\norm{\xi_c}\le1$, Definition~\ref{dp:def:strong-norm} gives $\norm{v_{\mathsf w}}\le\lvert\!\lvert\!\lvert a\rvert\!\rvert\!\rvert$. Conclude as in Lemma~\ref{dp:lem:laws-cover}.
\end{proof}

\begin{remark}[failures of the two descents for analyticity]\label{dp:rem:an-no-descent}
\emph{Target.} Let $\operatorname{char}K=p$, $k\ge p$, and let $E,F$ be as in \cite[Theorem~7.2(1)]{I}. Then $(j_{E,F}\circ{-})\circ A^k_{E,E;F}$ is analytic everywhere, $j_{E,F}\circ{-}$ has closed range, and $A^k_{E,E;F}$ is analytic at no point \cite[Remark~14.13]{I}. In every characteristic, the germ $\mathsf f$ of Example~\ref{dp:ex:l1} shows the same.

\emph{Source.} Let $\operatorname{char}K=p$ and let $a$ be the law of Proposition~\ref{dp:prop:degree-p}(2) with $n=p$, into the completion of $\Gamma^p(\mathcal E)$. By Lemma~\ref{dp:lem:laws-cover}, $a\circ q_{K\oplus\mathcal E}$ is analytic everywhere. By \cite[Proposition~14.9(2)]{I} with $j=\id$, $a$ is analytic at no point, since it is not ambient.
\end{remark}

\begin{theorem}[normal form]\label{dp:thm:normal-form}
Let $W$ be a Banach space, $U\subseteq P$ open, $f\colon U\to W$, $x_0\in U$, and $\tilde x_0\in\ell^1_{00}(I_P)$ with $q_P(\tilde x_0)=x_0$. Let $\mathfrak U$, $u_0$ be as in \cite[(14.8)]{I}, and let $j_W\colon W\to W\hot_\pi\mathfrak U$, $w\mapsto w\otimes u_0$, which is a linear isometry into a Banach space \cite[Remark~14.41]{I}. The following are equivalent:
\begin{enumerate}
\item $f$ is DP at $x_0$;
\item $f\circ q_P$ is AA at $\tilde x_0$;
\item there are a Banach space $Z$ and a linear isometry $j\colon W\to Z$ such that $j\circ f\circ q_P$ is analytic at $\tilde x_0$;
\item $j_W\circ f\circ q_P$ is analytic at $\tilde x_0$.
\end{enumerate}
Radius $R$ and bound $M$ in~(1) give radius $R/C_0$ and bound $M$ in~(2), (3), (4). Radius $R'$ and bound $M$ in~(3) or~(4) give radius $\abs\varpi R'$ and bound $M$ in~(1).
\end{theorem}

\begin{proof}
\emph{Step 1 ((1) implies (2)).} By Proposition~\ref{dp:prop:source-descent}(2) and $\norm{q_P}\le1$, $f\circ q_P$ is DP at $\tilde x_0$ with radius $R$, bound $M$ and laws $b_n=a_n\circ q_P$. By \cite[Theorem~14.39(2)]{I} and $\norm{L_{b_n}(t)}\le\norm{b_n}\norm t_\Gamma$ (Theorem~\ref{dp:thm:normed-roby}(1)), we get
\begin{equation}\label{dp:eq:amb-cover}
 \operatorname{amb}(b_n)\le C_0^nMR^{-n}=M(R/C_0)^{-n},
\end{equation}
and for $n=0$ the law is a constant of norm at most $M$. As $C_0>1$, \cite[Proposition~14.9(1)]{I} gives that $f\circ q_P$ is AA with radius $R/C_0$ and bound $M$.

\emph{Step 2 ((2) implies (3)).} This is \cite[Remark~14.13]{I}, with the same radius and bound.

\emph{Step 3 ((3) implies (1)).} The map $j\circ f\circ q_P$ is DP at $\tilde x_0$ with radius $R'$ and bound $M$ \cite[Proposition~14.9(1)]{I}. As $W$ is complete, Theorem~\ref{dp:thm:target-descent}(2) gives the same for $f\circ q_P$. Proposition~\ref{dp:prop:source-descent}(1) with $\mathsf c=\abs\varpi^{-1}$ gives that $f$ is DP at $x_0$ with radius $\abs\varpi R'$ and bound $M$.

\emph{Step 4 ((1) implies (4)).} This is the construction of \cite[Remark~14.41]{I}, stated there on $\ell^1(J)$. On $\ell^1_{00}(J)$, with $J=I_P$, every sum below is finite. Recall that $u_0=\sum_{i\le N}\mathsf u_{N,i}$ with $\norm{\mathsf u_{N,i}}\le1/N$. Fix $n\ge1$. For finitely supported $\alpha\in\N^J$ with $\abs\alpha=n$ let $O_\alpha(e)$ be the orbit sum of the unit vectors in the support of $\alpha$, and $w_\alpha=L_{b_n}(O_\alpha(e))$. There are $N_\alpha=n!/\prod_\theta\alpha_\theta!$ words $\mathsf w\in J^n$ of content $\alpha$, that is, in which each $\theta$ occurs $\alpha_\theta$ times. So $\pi(O_\alpha(e))\le N_\alpha$ and $\norm{w_\alpha}\le\operatorname{amb}(b_n)N_\alpha$. Number these words, and put $\mathsf z_{\mathsf w}=w_\alpha\otimes\mathsf u_{N_\alpha,i}$ if $\mathsf w$ is the $i$-th. Then $\norm{\mathsf z_{\mathsf w}}\le\operatorname{amb}(b_n)$, and the $\mathsf z_{\mathsf w}$ of content $\alpha$ sum to $w_\alpha\otimes u_0$. Put
\[
 \beta_n(h^1,\dots,h^n)=\sum_{\mathsf w\in J^n}h^1_{\mathsf w_1}\cdots h^n_{\mathsf w_n}\,\mathsf z_{\mathsf w}.
\]
Then $\norm{\beta_n}\le\operatorname{amb}(b_n)\le M(R/C_0)^{-n}$ by~\eqref{dp:eq:amb-cover}, and by~\eqref{dp:eq:orbit},
\[
 \beta_n(h,\dots,h)=\sum_\alpha h^\alpha w_\alpha\otimes u_0=L_{b_n}(h^{\otimes n})\otimes u_0=j_W(b_n(h)).
\]
With $\beta_0=j_W(f(x_0))$ this expands $j_W\circ f\circ q_P$ at $\tilde x_0$ with radius $R/C_0$ and bound $M$.

\emph{Step 5 ((4) implies (3)).} Take $j=j_W$.
\end{proof}

\subsection{Minimality}\label{dp:ssec:minimal}

Here a \emph{germ} is a quadruple $(P,U,f,x_0)$: $P$ a normed space, $U\subseteq P$ open, $x_0\in U$ and $f\colon U\to W$ with $W$ a Banach space. Consider three properties of a class $\cC$ of germs.
\begin{itemize}
\item[(C1)] Every germ $(P,U,f,x_0)$ with $f$ analytic at $x_0$ is in $\cC$.
\item[(C2)] \emph{(Target descent.)} If $j\colon W\to Z$ is a linear isometry of Banach spaces and $(P,U,j\circ f,x_0)\in\cC$, then $(P,U,f,x_0)\in\cC$.
\item[(C3)] \emph{(Source descent.)} If $Q\colon P'\to P$ is a bounded linear surjection with a preimage constant, $\tilde x_0\in Q^{-1}(U)$ and $(P',Q^{-1}(U),f\circ Q,\tilde x_0)\in\cC$, then $(P,U,f,Q(\tilde x_0))\in\cC$.
\end{itemize}

\begin{theorem}[the smallest class with these descents]\label{dp:thm:minimal}
\leavevmode
\begin{enumerate}
\item The DP germs satisfy (C1), (C2), (C3).
\item Every class of germs with (C1), (C2), (C3) contains every DP germ. It suffices to require (C1) for the sources $\ell^1_{00}(J)$, (C2) for the isometries $j_W$ of Theorem~\ref{dp:thm:normal-form} and (C3) for the canonical covers $q_P$.
\item So the DP germs are the smallest class of germs with (C1), (C2), (C3).
\item Let $K$ be complete. The DP germs with Banach sources are the smallest class of germs with Banach sources that has (C1) and (C2) and has (C3) for surjections between Banach spaces.
\end{enumerate}
\end{theorem}

\begin{proof}
\emph{Step 1 (item (1)).} (C1) is Theorem~\ref{dp:thm:calculus}(3). For (C2), an isometry from a Banach space has closed range, and Theorem~\ref{dp:thm:target-descent}(2) applies. (C3) is Proposition~\ref{dp:prop:source-descent}(1).

\emph{Step 2 (item (2)).} Let $(P,U,f,x_0)$ be a DP germ, and choose $\tilde x_0$ with $q_P(\tilde x_0)=x_0$. By Theorem~\ref{dp:thm:normal-form}(4), $(\ell^1_{00}(I_P),q_P^{-1}(U),j_W\circ f\circ q_P,\tilde x_0)\in\cC$ by (C1). So $(\ell^1_{00}(I_P),q_P^{-1}(U),f\circ q_P,\tilde x_0)\in\cC$ by (C2), and $(P,U,f,x_0)\in\cC$ by (C3). Item~(3) follows from items~(1) and~(2).

\emph{Step 3 (item (4)).} Let $P$ be a Banach space and choose $\tilde x_0\in\ell^1(I_P)$ with $\hat q_P(\tilde x_0)=x_0$. As $K$ is complete, $\ell^1(I_P)$ is a Banach space, and $\hat q_P$ is contractive with preimage constant $\abs\varpi^{-1}$. By \cite[Theorem~14.39(2)]{I} for $\ell^1(I_P)$, Steps~1 and~2 of the proof of Theorem~\ref{dp:thm:normal-form} apply with $\hat q_P$ in place of $q_P$. They give a Banach space $Z$ and a linear isometry $j\colon W\to Z$ with $j\circ f\circ\hat q_P$ analytic at $\tilde x_0$. Then (C1), (C2), (C3) apply as in Step~2, and all sources involved are Banach spaces.
\end{proof}

\begin{proposition}[the two other classes]\label{dp:prop:other-classes}
\leavevmode
\begin{enumerate}
\item The analytic germs satisfy (C1). Over every $K$ they fail (C2). If $\operatorname{char}K=p$ they fail (C3). If $\operatorname{char}K=0$, then (C2) and (C3) hold for germs whose source $P$ is ultrametric, the source $P'$ of $Q$ in (C3) being arbitrary.
\item The AA germs are the smallest class of germs with (C1) and (C2). If $\operatorname{char}K=p$ they fail (C3).
\end{enumerate}
\end{proposition}

\begin{proof}
\emph{Step 1 (item (1)).} The failures are Example~\ref{dp:ex:l1} and Remark~\ref{dp:rem:an-no-descent}. Let $\operatorname{char}K=0$ and $P$ be ultrametric. If $j\circ f$ is analytic, it is DP, so $f$ is DP (Theorem~\ref{dp:thm:target-descent}(2)), so $f$ is analytic \cite[Proposition~14.40(2)]{I}. If $f\circ Q$ is analytic, it is DP, so $f$ is DP (Proposition~\ref{dp:prop:source-descent}(1)), and again $f$ is analytic.

\emph{Step 2 (item (2): minimality).} Analytic germs are AA, with $j=\id$. If $\cC$ has (C1) and (C2), and $f$ is AA at $x_0$, then \cite[Remark~14.13]{I} gives a Banach space $Z$ and an isometry $j$ with $j\circ f$ analytic. So $f\in\cC$.

\emph{Step 3 (item (2): the AA germs have (C2)).} Let $j\colon W\to Z$ be an isometry of Banach spaces with $j\circ f$ AA. By \cite[Remark~14.13]{I} there is an isometry $j'$ of $Z$ into a Banach space with $j'\circ j\circ f$ analytic. So $f$ is AA, through $j'\circ j$, by the same remark.

\emph{Step 4 (item (2): failure of (C3)).} Let $a$ be the law of Remark~\ref{dp:rem:an-no-descent}. It has Banach values and is not ambient. If it were AA at a point, \cite[Remark~14.13]{I} and \cite[Proposition~14.9(2)]{I} would make it ambient. And $a\circ q_{K\oplus\mathcal E}$ is analytic, hence AA, everywhere.
\end{proof}

In characteristic $0$, on non-ultrametric sources, the analytic germs with ultrametric Banach targets have (C3) (Corollary~\ref{dp:cor:two-gaps}). For non-ultrametric sources and targets in characteristic $0$, whether they have (C3) is open.

The descents (C2), (C3) are properties of $C^\infty$ maps. For (C2) this is \cite[Lemma~2.4]{I}. For (C3):

\begin{proposition}[$C^m$ descends along source quotients]\label{dp:prop:cn-descends}
Let $Q\colon P'\to P$ be a bounded linear surjection with preimage constant $\mathsf c$, $U\subseteq P$ open, $W$ normed, $m\in\N\cup\{\infty\}$ and $f\colon U\to W$. If $f\circ Q$ is $C^m$ on $Q^{-1}(U)$, then $f$ is $C^m$ on $U$.
\end{proposition}

\begin{proof}
\emph{Step 1 ($m=0$).} Let $x_0=Q(\tilde x_0)\in U$. For $x$ near $x_0$ choose $\tilde x$ with $Q\tilde x=x$ and $\norm{\tilde x-\tilde x_0}\le\mathsf c\norm{x-x_0}$. Then $f(x)-f(x_0)=f\circ Q(\tilde x)-f\circ Q(\tilde x_0)\to0$.

\emph{Step 2 ($m=1$).} Let $x=Q\tilde x$ and $T=D(f\circ Q)(\tilde x)$. For $\mathsf k\in\ker Q$ and small $s$ we have $f\circ Q(\tilde x+s\mathsf k)=f(x)$, so $T\mathsf k=0$. Hence $T=T_0\circ Q$ for a unique linear $T_0\colon P\to W$, and $\norm{T_0h}=\norm{T\tilde h}\le\mathsf c\norm T\norm h$. Also $\norm{f(x+h)-f(x)-T_0h}=\norm{f\circ Q(\tilde x+\tilde h)-f\circ Q(\tilde x)-T\tilde h}=o(\norm{\tilde h})=o(\norm h)$, so $Df(x)=T_0$. For $x'$ near $x$ choose $\tilde x'$ with $Q\tilde x'=x'$ and $\norm{\tilde x'-\tilde x}\le\mathsf c\norm{x'-x}$. Then $\norm{Df(x')-Df(x)}\le\mathsf c\norm{D(f\circ Q)(\tilde x')-D(f\circ Q)(\tilde x)}\to0$.

\emph{Step 3 (induction).} We induct on $m$, for all normed $W$. The map $Q^*\colon\cL(P,W)\to\cL(P',W)$, $T_0\mapsto T_0\circ Q$, is linear with $\norm{Q^*T_0}\ge\norm{T_0}/\mathsf c$. By Step~2 its range is the set of operators that vanish on $\ker Q$, which is closed. By the chain rule, $D(f\circ Q)=Q^*\circ(Df)\circ Q$. If $f\circ Q$ is $C^{m+1}$, then $D(f\circ Q)$ is $C^m$ with values in the closed range of $Q^*$. By \cite[Lemma~2.4]{I} it is $C^m$ as a map into that range. Composing with the bounded inverse of $Q^*$ on its range, $(Df)\circ Q$ is $C^m$. By induction, applied to $Df\colon U\to\cL(P,W)$, $Df$ is $C^m$, so $f$ is $C^{m+1}$. For $m=\infty$ apply this for every finite $m$.
\end{proof}

So Theorem~\ref{dp:thm:minimal} reads: the DP germs are the smallest class that contains the analytic germs and has the two descents of $C^\infty$ maps.

\subsection{Ultrametric spaces}\label{dp:ssec:um-minimal}

Between ultrametric spaces the DP germs are already the AA germs, through an ultrametric enlargement \cite[Theorem~14.39(1)]{I}. Here ultrametric norms simplify further: no target descent is needed. The reason is in the bound of Lemma~\ref{dp:lem:laws-cover}. Its factor $n!$ counts the words that share a monomial, and removing it needs an enlargement of $W$. If $W$ is ultrametric and the cover has the supremum norm, the sum~\eqref{dp:eq:word-lift} has no count of terms and each $v_{\mathsf w}$ is bounded geometrically. So the cover alone makes each law a $W$-valued diagonal. In this subsection, unless said otherwise, the target of a germ is any normed space. Conditions (C1) and (C3) are read among the germs of each statement, so in (C3) the source $P'$ of $Q$ is also one of the sources of that statement.

\begin{theorem}[between ultrametric spaces]\label{dp:thm:ultrametric-minimal}
Let $P$ and $W$ be ultrametric normed spaces. Let $q^{\rm u}_P\colon c_{00}(I_P)\to P$, $e_h\mapsto h$, where $c_{00}(I_P)$ has the supremum norm. Then $q^{\rm u}_P$ is contractive with preimage constant $\abs\varpi^{-1}$. For $U\subseteq P$ open, $f\colon U\to W$ and $x_0=q^{\rm u}_P(\tilde x_0)\in U$:
\begin{enumerate}
\item $f$ is DP at $x_0$ if and only if $f\circ q^{\rm u}_P$ is analytic at $\tilde x_0$. Radius $R$ and bound $M$ give radius $\abs\varpi R$ and bound $M$, in both directions.
\item Among germs with ultrametric sources and ultrametric targets, the DP germs are the smallest class with (C1) and (C3).
\item The analytic germs between ultrametric spaces have (C3) if and only if every DP germ between ultrametric spaces is analytic. This fails for $K$ complete, not spherically complete, of characteristic $p$.
\item Let $K$ be complete and let $P$, $W$ be Banach. Items~(1) and~(2) hold, within Banach spaces, with the cover $c_0(I_P)\to P$, $\mathsf v\mapsto\sum_h\mathsf v_hh$.
\end{enumerate}
\end{theorem}

\begin{proof}
\emph{Step 1 (the cover).} As $P$ is ultrametric, $\norm{\sum_h\mathsf v_hh}\le\max_h\abs{\mathsf v_h}\norm h\le\norm{\mathsf v}_\infty$. The preimages are those of the paragraph after Definition~\ref{dp:def:cover}, of supremum norm $\abs\lambda^{-1}<\abs\varpi^{-1}\norm h$.

\emph{Step 2 (laws on the cover).} Let $J$ be a totally ordered set and $b\colon c_{00}(J)\to W$ a bounded law of degree $n$. Define $v_{\mathsf w}$ and $\beta_b$ as before~\eqref{dp:eq:word-lift}; the same computation gives $\beta_b(h,\dots,h)=b(h)$. On the span of the letters of $\mathsf w$, which carries the maximum norm, \cite[Lemma~Q.2(4)]{I} gives $\norm{v_{\mathsf w}}\le\abs\varpi^{-n}\norm b$. As $W$ is ultrametric, $\norm{\beta_b(h^1,\dots,h^n)}\le\max_{\mathsf w}\prod_c\abs{h^c_{\mathsf w_c}}\norm{v_{\mathsf w}}\le\abs\varpi^{-n}\norm b\prod_c\norm{h^c}_\infty$. So $\norm{\beta_b}\le\abs\varpi^{-n}\norm b$.

\emph{Step 3 (item (1)).} Let $f$ be DP at $x_0$ with radius $R$, bound $M$ and laws $a_n$. By Proposition~\ref{dp:prop:source-descent}(2), $f\circ q^{\rm u}_P$ is DP at $\tilde x_0$ with laws $a_n\circ q^{\rm u}_P$ of norm at most $MR^{-n}$. By Step~2 each is the diagonal of an $n$-linear map into $W$ of norm at most $M(\abs\varpi R)^{-n}$. So $f\circ q^{\rm u}_P$ is analytic at $\tilde x_0$ with radius $\abs\varpi R$ and bound $M$. Conversely, an analytic germ with radius $R$ and bound $M$ is DP with the same data \cite[Proposition~14.9(1)]{I}, and Proposition~\ref{dp:prop:source-descent}(1) with $\mathsf c=\abs\varpi^{-1}$ gives radius $\abs\varpi R$ and bound $M$.

\emph{Step 4 (item (2)).} The DP germs contain the analytic germs and have (C3). Let $\cC$ have (C1) and (C3), and let $f$ be DP at $x_0$. By item~(1), $f\circ q^{\rm u}_P$ is analytic, so it is in $\cC$. As $c_{00}(I_P)$ is ultrametric, the germ stays in the class of germs considered, and (C3) puts $f$ in $\cC$.

\emph{Step 5 (item (3)).} If the analytic germs have (C3) and $f$ is DP, then $f\circ q^{\rm u}_P$ is analytic by item~(1), so $f$ is analytic. Conversely, if the analytic germs are the DP germs, they have (C3) by Proposition~\ref{dp:prop:source-descent}(1). For the failure, \cite[Proposition~14.40(3)]{I} gives a function on an ultrametric Banach space that is AA, hence DP, at every point and analytic at no point.

\emph{Step 6 (item (4)).} The map $\mathsf v\mapsto\sum_h\mathsf v_hh$ is defined: the sum converges, as $\mathsf v_hh\to0$ in the complete ultrametric $P$. The map has the properties of Step~1. Let $b\colon c_0(J)\to W$ be a bounded law of degree $n$, with $v_{\mathsf w}$ as in Step~2. For $h^1,\dots,h^n\in c_0(J)$ and $\varepsilon>0$, only finitely many words have $\prod_c\abs{h^c_{\mathsf w_c}}>\varepsilon$. So~\eqref{dp:eq:word-lift} converges in the Banach space $W$, with $\norm{\beta_b}\le\abs\varpi^{-n}\norm b$. The maps $h\mapsto\beta_b(h,\dots,h)$ and $b$ are continuous (Theorem~\ref{dp:thm:calculus}(1)) and agree on the dense subspace $c_{00}(J)$. Steps~3 and~4 then apply.
\end{proof}

\begin{corollary}[ultrametric sources, every target]\label{dp:cor:um-cover}
Let $P$ be ultrametric and $W$ normed, and let $f$ be DP at $x_0$ with radius $R$ and bound $M$. Then $f\circ q_P$ is analytic at every preimage of $x_0$, with radius $R/C_1$ and bound $M$. So every class of germs that has (C1) for the sources $\ell^1_{00}(J)$ and (C3) for the canonical covers contains every DP germ on an ultrametric source, with no target descent.
\end{corollary}

\begin{proof}
Let $a_n$ be the laws of $f$ at $x_0$. By Lemma~\ref{dp:lem:lift-cover} and Lemma~\ref{dp:lem:strong-formal}(3), $a_n\circ q_P$ is the diagonal of a bounded $n$-linear map into $W$ of norm at most $\lvert\!\lvert\!\lvert a_n\rvert\!\rvert\!\rvert\le C_1^n\norm{a_n}\le M(R/C_1)^{-n}$. By Proposition~\ref{dp:prop:source-descent}(2), these laws expand $f\circ q_P$.
\end{proof}

For $P\ne0$ the source $\ell^1_{00}(I_P)$ of this cover is not ultrametric. In characteristic $0$ the DP germs on ultrametric sources are the analytic germs \cite[Proposition~14.40(2)]{I}, so (C1) alone gives them. In characteristic $p$, whether, among germs with ultrametric sources and Banach targets, the DP germs are the smallest class with (C1), (C2), (C3) is open; see Question~\ref{dp:q:c0}(2).

\subsection{What the descents reach}\label{dp:ssec:reach}

Let $(Z;W_1,W_2)$ and $(Z';W_1',W_2')$ be filtered spaces as in \cite[Definition~14.1]{I}, with $Z,Z'$ Banach, and with quotient maps $\operatorname{qt}\colon W_1\to W_1/W_2$ and $\operatorname{qt}'\colon W_1'\to W_1'/W_2'$. Let $w_0\in W_1$. A \emph{morphism germ} at $w_0$ is the germ of a map $\Phi$, from an open neighborhood $\mathsf O$ of $w_0$ in $Z$ into $Z'$, analytic at $w_0$, with $\Phi(\mathsf O\cap W_1)\subseteq W_1'$ and with $\Phi(w+w')-\Phi(w)\in W_2'$ whenever $w,w+w'\in\mathsf O\cap W_1$ and $w'\in W_2$. It induces a germ $\varphi$ at $\operatorname{qt}(w_0)$ from $W_1/W_2$ to $W_1'/W_2'$, by $\varphi(\operatorname{qt}(w))=\operatorname{qt}'(\Phi(w))$. As $Z'$ is complete, rearranging the series shows that a map analytic at $w_0$ is analytic on a ball around $w_0$, so this agrees with \cite[Definition~14.25]{I} when $W_2=W_2'=0$. If $W_2=0$, the induced germs are AA: $\Phi|_{W_1}$ is AA into the Banach space $W_1'$ \cite[Remark~14.13]{I}, and so is $\operatorname{qt}'\circ\Phi|_{W_1}$ \cite[Proposition~14.12(1)]{I}. With a quotient source, every DP germ between Banach spaces is induced, after the identifications of Proposition~\ref{dp:prop:subquotients}(2). For $K$ complete of characteristic $p$ these include germs that are AA at no point: the law $a$ of Remark~\ref{dp:rem:an-no-descent}, by Step~4 of the proof of Proposition~\ref{dp:prop:other-classes}, as AA at a point does not change under an equivalent norm on the source. In characteristic $0$ we do not decide whether they go beyond the AA germs: this is the case of Banach sources in \cite[Question~14.42(1)]{I}.

\begin{proposition}[subquotient sources]\label{dp:prop:subquotients}
\leavevmode
\begin{enumerate}
\item Every germ induced by a morphism germ between filtered spaces as above is DP at $\operatorname{qt}(w_0)$.
\item Let $K$ be complete, let $P,W$ be Banach spaces and let $f\colon U\to W$ be DP at $x_0\in U$. Then $f$ is induced by a morphism germ, after two identifications: of $P$, with an equivalent norm, with $\ell^1(I_P)/\ker\hat q_P$, and of $W$ with a closed subspace of a Banach space.
\end{enumerate}
\end{proposition}

\begin{proof}
\emph{Step 1 (item (1)).} As a map from $W_1$ into $Z'$, $\Phi|_{W_1}$ is analytic at $w_0$, hence DP at $w_0$ (Theorem~\ref{dp:thm:calculus}(3)). Its values near $w_0$ lie in the closed subspace $W_1'$, so it is DP into $W_1'$ (Theorem~\ref{dp:thm:target-descent}(2)). Composing with the bounded linear $\operatorname{qt}'$ keeps this. Near $w_0$ we have $\varphi\circ\operatorname{qt}=\operatorname{qt}'\circ\Phi|_{W_1}$, so $\varphi\circ\operatorname{qt}$ is DP at $w_0$. The map $\operatorname{qt}$ is open, so $\varphi$ is defined on the open set $\operatorname{qt}(\mathsf O\cap W_1)$, and $\operatorname{qt}$ has preimage constant $\abs\varpi^{-1}$ for the quotient norm. Proposition~\ref{dp:prop:source-descent}(1) gives that $\varphi$ is DP at $\operatorname{qt}(w_0)$.

\emph{Step 2 (item (2)).} By the proof of Theorem~\ref{dp:thm:minimal}(4), there are $\tilde x_0\in\ell^1(I_P)$ with $\hat q_P(\tilde x_0)=x_0$ and a linear isometry $j$ of $W$ into a Banach space $Z'$ with $\Phi=j\circ f\circ\hat q_P$ analytic at $\tilde x_0$. The map $\Phi$ has values in the closed subspace $j(W)$ and is constant on the cosets of $\ker\hat q_P$. So it is a morphism germ from $(\ell^1(I_P);\ell^1(I_P),\ker\hat q_P)$ to $(Z';j(W),0)$. As $\hat q_P$ is contractive with preimage constant $\abs\varpi^{-1}$, it induces a bijection of $\ell^1(I_P)/\ker\hat q_P$ onto $P$ with $\norm{\hat q_P(\mathsf v)}\le\norm{[\mathsf v]}\le\abs\varpi^{-1}\norm{\hat q_P(\mathsf v)}$. Under this bijection and $j$, the germ induced by $\Phi$ is $f$.
\end{proof}

An equivalent norm on the source changes the radius and keeps the class of DP germs. So, for $K$ complete and Banach spaces, the germs induced with quotient sources allowed are exactly the DP germs. This is Theorem~\ref{dp:thm:minimal} again: sub-pairs are target descent, and quotient sources are source descent.

Now the example that shows what this minimality rests on.

\begin{example}[the $\ell^1$ germ]\label{dp:ex:l1}
Let $P=\ell^1(\N,K)$ and $W=\Kh$. Choose $\mathsf e_n\in K$ with $n^n\le\abs{\mathsf e_n}<\abs\varpi^{-1}n^n$, and put
\[
 a_n(h)=\mathsf e_nh_1h_2\cdots h_n,\qquad\mathsf f(h)=\sum_{n\ge1}a_n(h)\qquad(\norm h<1).
\]
For $K$ complete this is the germ of \cite[Theorem~14.33(1)]{I}. Over every $K$:
\begin{enumerate}
\item $\norm{a_n}<\abs\varpi^{-1}$, so $\mathsf f$ is DP at $0$ with radius $1$ and bound $\abs\varpi^{-1}$.
\item $\mathsf f$ is AA at $0$ with radius $1/C_0$ and bound $\abs\varpi^{-1}$. So $j\circ\mathsf f$ is analytic at $0$ for a linear isometry $j$ of $\Kh$ into a Banach space.
\item $\mathsf f$ is not analytic at $0$.
\end{enumerate}
So the analytic germs fail (C2) in every characteristic. By items~(1) and~(3), finite coordinates cannot be dropped from Theorem~\ref{dp:thm:calculus}(3).
\end{example}

\begin{proof}
\emph{Step 1 (item (1)).} By the arithmetic--geometric mean inequality, $\prod_{i\le n}\abs{h_i}\le(\norm h/n)^n$. So $\norm{a_n}\le\abs{\mathsf e_n}n^{-n}<\abs\varpi^{-1}$.

\emph{Step 2 (item (2)).} By \cite[Theorem~14.39(2)]{I}, $\operatorname{amb}(a_n)\le C_0^n\norm{a_n}<\abs\varpi^{-1}C_0^n$. Then \cite[Proposition~14.9(1)]{I} and \cite[Remark~14.13]{I} apply.

\emph{Step 3 (item (3)).} Suppose $\mathsf f(h)=\sum_n\beta_n(h,\dots,h)$ near $0$, with $n$-linear $\beta_n\colon P^n\to\Kh$ and $\norm{\beta_n}\le MR^{-n}$. By Lemma~\ref{dp:lem:coefficients-unique}(2) with $j=\id$, $\beta_n(h,\dots,h)=a_n(h)$. As $K$ is infinite, comparing coefficients of $s_1\cdots s_n$ at $h=\sum_{i\le n}s_ie_i$ gives $\sum_{\sigma\in S_n}\beta_n(e_{\sigma(1)},\dots,e_{\sigma(n)})=\mathsf e_n$. The absolute value is nonarchimedean and $\norm{e_i}=1$, so $\norm{\beta_n}\ge\abs{\mathsf e_n}\ge n^n$. This is not at most $MR^{-n}$ for large $n$.
\end{proof}

\emph{The mechanism.} The norm of $\ell^1$ adds real numbers. So the arithmetic--geometric mean inequality applies, with a gain of $n^{-n}$, and $\mathsf e_nh_1\cdots h_n$ stays bounded with $\abs{\mathsf e_n}$ about $n^n$. A multilinear map $\beta$ with this diagonal has $\sum_\sigma\beta(e_{\sigma(1)},\dots,e_{\sigma(n)})=\mathsf e_n$. Over $\mathbb C$ the $n!$ orderings could share the load: $\beta$ could take values $\mathsf e_n/n!$, of size about $\mathrm e^n$. In $\Kh$ they cannot, as a sum equal to $\mathsf e_n$ has a term of absolute value at least $\abs{\mathsf e_n}$. Multilinear maps into an ultrametric space see only the supremum norm on $\ell^1$ (Proposition~\ref{dp:prop:hull}(1) and the paragraph before it), and there $h_1\cdots h_n$ has size $1$. So $\sum_na_n$ converges on the ball of $\ell^1$ because of an inequality between real numbers on the source. No divided power occurs: the example has the same form for $K=\mathbb C((t))$, in residue characteristic $0$. A non-ultrametric enlargement $Z$ can hold the shares in separate coordinates, which is why $j\circ\mathsf f$ is analytic for some $j$.

\emph{What this says about minimality.} By Example~\ref{dp:ex:l1}(2), every class with (C1) and (C2) contains $\mathsf f$. Among germs with Banach values, the class CA of germs analytic after a cover (Section~\ref{dp:sec:cover}) is the smallest class with (C1) and (C3), and it does not contain $\mathsf f$ (Lemma~\ref{dp:lem:cover-flat}(3), Theorem~\ref{dp:thm:cah}(5)). Within ultrametric targets, the strongly DP germs have (C1), (C3) and target descent along isometries between ultrametric spaces (Theorem~\ref{dp:thm:strong-descent}), and do not contain $\mathsf f$, as the coefficient of $a_n$ at $e_1,\dots,e_n$ has absolute value at least $n^n$. So the minimality of Theorem~\ref{dp:thm:minimal} on general sources rests on target descent into non-ultrametric spaces. On general sources, (C1) and (C3) alone do not reach the DP germs. Two axioms that reach them when added are (C2), by Theorem~\ref{dp:thm:minimal}, and the axiom (S) of Section~\ref{dp:sec:ledger}, by Proposition~\ref{dp:prop:s-axiom}(4). On ultrametric sources, (C1) and (C3) alone reach them, through covers that are not ultrametric (Corollary~\ref{dp:cor:um-cover}). Section~\ref{dp:sec:ledger} lists what each disagreement with analyticity is forced by.
%% ---- end sec-06-descent.tex
%% ---- begin sec-07-a-tree.tex
\section{Classes between ambient and divided-power analyticity}\label{dp:sec:chain}

Paper~\cite{I} shows that ambient analyticity is not closed under composition. Over every $K$ of characteristic $p$, on the space $H^{\rm at}$ of~\cite[Definition~14.16]{I}, the map $f_X$ is AA at every point, the map $g_X$ on $W^{\rm at}$ is the diagonal of a bounded bilinear map, and $g_X\circ f_X$ is a bounded law of degree $p+1$ that is DP everywhere and AA at no point~\cite[Theorem~14.17]{I}. In general the question depends on the source. Let $\mathcal C(N,r)$ be the set of compositions of $N$ into $r$ positive parts, and for $\vec n\in\mathcal C(N,r)$ let $\pi_{\vec n}$ be the projective seminorm of $\Gamma^{n_1}_\pi(P)\otimes\dots\otimes\Gamma^{n_r}_\pi(P)$, which contains $\Gamma^N(P)$. If $\mathfrak k(P)=\sup_{N,\vec n}\bigl(\sup_{t\in\Gamma^N(P)}\pi_{\vec n}(t)/\pi(t)\bigr)^{1/N}$ is finite, then composites of AA germs on $P$ are AA; for $K$ complete the converse holds, already for an analytic germ after an AA germ, with values in Banach spaces~\cite[Definition~14.21, Theorem~14.22]{I}. This section describes what composition produces from analytic germs, and compares the result with DP.

\subsection{Tree norms}\label{dp:ssec:tree}

We begin with the constants $\mathfrak k_{\vec n}(P)=\sup_t\pi_{\vec n}(t)/\pi(t)$ of~\cite[Definition~14.21]{I}, with $0/0=0$ and $c/0=\infty$ for $c>0$, and with $\mathfrak a_N(P)$ as in~\cite[Definition~14.36]{I}.

\begin{lemma}[bounds for the composition constants]\label{dp:lem:kappa-bounds}
Let $P$ be normed and $\vec n\in\mathcal C(N,r)$.
\begin{enumerate}
\item On $\Gamma^N(P)$ we have $\pi\le\pi_{\vec n}\le\norm\cdot_\Gamma$, so $\mathfrak k_{\vec n}(P)\le\mathfrak a_N(P)$. If $r=1$ or $r=N$, then $\pi_{\vec n}=\pi$.
\item If $n_i!\ne0$ in $K$ for all $i$, then $\mathfrak k_{\vec n}(P)\le\prod_in_i!/\abs{n_i!}$, where the numerator is a real number and $\abs{n_i!}$ is the absolute value of the image of $n_i!$ in $K$. If $\operatorname{char}K=p$ and all $n_i<p$, then $\mathfrak k_{\vec n}(P)\le\prod_in_i!$.
\item $\mathfrak k_{\vec n}(\ell^1(J,K))\le1$ for every set $J$.
\end{enumerate}
\end{lemma}

\begin{proof}
\emph{Step 1 (item (1)).} The inequalities are in~\cite[Definition~14.21]{I}. For $r=1$ the definition of $\pi_{\vec n}$ reads $\pi_{(N)}=\pi$. For $r=N$ all $n_i$ are $1$ and $\Gamma^1(P)=P$.

\emph{Step 2 (item (2)).} Let $n!\ne0$ in $K$ and $\operatorname{Sym}_n=\frac1{n!}\sum_{\sigma\in S_n}\sigma$ on $P^{\otimes n}$. For $x_1,\dots,x_n\in P$, the tensor $\operatorname{Sym}_n(x_1\otimes\dots\otimes x_n)=\frac1{n!}O_{(1,\dots,1)}(x_1,\dots,x_n)$ lies in $\Gamma^n(P)$, by~\eqref{dp:eq:orbit}, and it has $\pi$ at most $(n!/\abs{n!})\norm{x_1}\cdots\norm{x_n}$. Let $t\in\Gamma^N(P)$. The operator $\operatorname{Sym}_{n_1}\otimes\dots\otimes\operatorname{Sym}_{n_r}$ fixes the powers, so it fixes $t$. Apply it to a representation of $t$ in $P^{\otimes N}$ of cost close to $\pi(t)$. This gives a representation of $t$ in $\bigotimes_i\Gamma^{n_i}(P)$ of cost at most $\prod_in_i!/\abs{n_i!}$ times that cost. In characteristic $p$, for $n<p$ the element $n!$ is a nonzero element of $\F_p$, so $\abs{n!}=1$.

\emph{Step 3 (the norm on $\ell^1$).} Let $P=\ell^1(J,K)$, with coordinates $t(\vec u)$ on $P^{\otimes N}$, $\vec u\in J^N$. Every representation $t=\sum_qx_{q,1}\otimes\dots\otimes x_{q,N}$ satisfies $\sum_{\vec u}\abs{t(\vec u)}\le\sum_q\prod_i\norm{x_{q,i}}$, so $\pi(t)\ge\sum_{\vec u}\abs{t(\vec u)}$. In particular $\pi$ is a norm. For $t\in(K^{J_1})^{\otimes N}$ with $J_1\subseteq J$ finite, the expansion of $t$ in the unit vectors $\mathbf e_u$ gives equality.

\emph{Step 4 (item (3), finitely supported $t$).} Let $J_1\subseteq J$ be finite, $\mathbf e=(\mathbf e_u)_{u\in J_1}$, and $t\in\Gamma^N(K^{J_1})$. Call the elements of $J_1^N$ words; a word has content $\bar\alpha\in\N^{J_1}$ if each $u\in J_1$ occurs $\bar\alpha_u$ times. As $t$ is symmetric, $t=\sum_{\bar\alpha}t_{\bar\alpha}O_{\bar\alpha}(\mathbf e)$ over the $\bar\alpha\in\N^{J_1}$ with $\abs{\bar\alpha}=N$, where $t_{\bar\alpha}$ is the coordinate of $t$ at the words of content $\bar\alpha$. By Step~3, $\pi(t)=\sum_{\bar\alpha}(N!/\bar\alpha!)\abs{t_{\bar\alpha}}$. A word of content $\bar\alpha$ splits uniquely into $r$ blocks of lengths $n_i$, of contents $\alpha_i$, so
\[
 t=\sum_{\bar\alpha}t_{\bar\alpha}\sum_{\substack{\alpha_1+\dots+\alpha_r=\bar\alpha\\\abs{\alpha_i}=n_i}}O_{\alpha_1}(\mathbf e)\otimes\dots\otimes O_{\alpha_r}(\mathbf e).
\]
Here $O_{\alpha_i}(\mathbf e)\in\Gamma^{n_i}(P)$ and $\pi(O_{\alpha_i}(\mathbf e))\le n_i!/\alpha_i!$. Counting words by their blocks gives $\sum\prod_in_i!/\alpha_i!=N!/\bar\alpha!$, the sum over the same index set. So $\pi_{\vec n}(t)\le\pi(t)$.

\emph{Step 5 (item (3), general $t$).} Let $t=\sum_qc_qh_q^{\otimes N}\in\Gamma^N(P)$, let $\mathsf p_{J_1}$ be the coordinate projection onto $K^{J_1}$, of norm at most $1$, and $t_{J_1}=\mathsf p_{J_1}^{\otimes N}t$. Then $\pi_{\vec n}(t_{J_1})\le\pi(t_{J_1})\le\pi(t)$ by Step~4, and by item~(1),
\[
 \pi_{\vec n}(t-t_{J_1})\le\norm{t-t_{J_1}}_\Gamma\le\sum_q\abs{c_q}\norm{h_q^{\otimes N}-(\mathsf p_{J_1}h_q)^{\otimes N}}_\Gamma.
\]
By Step~3, $\pi$ is a norm on $\Gamma^N(P)$, and so is $\norm\cdot_\Gamma\ge\pi$. The map $h\mapsto h^{\otimes N}$ is a law of norm at most $1$ into the normed space $(\Gamma^N(P),\norm\cdot_\Gamma)$ by~\cite[Lemma~Q.1(2)]{I}, so by~\cite[(Q.3)]{I} the right side tends to $0$ as $J_1$ increases, because $\norm{h_q-\mathsf p_{J_1}h_q}\to0$. So $\pi_{\vec n}(t)\le\pi_{\vec n}(t_{J_1})+\pi_{\vec n}(t-t_{J_1})\le\pi(t)+\pi_{\vec n}(t-t_{J_1})$, and letting $J_1$ increase gives $\pi_{\vec n}(t)\le\pi(t)$.
\end{proof}

\begin{definition}[tree norms]\label{dp:def:tree}
Let $P$ be normed. For $\ell\ge1$ and $N\ge1$ we define seminorms $\mathfrak t^{(\ell)}_N$ on $\Gamma^N(P)$, and for $\vec n\in\mathcal C(N,r)$ seminorms $\mathfrak t^{(\ell)}_{\vec n}$ on $\bigotimes_i\Gamma^{n_i}(P)$, as follows. Put $\mathfrak t^{(1)}_N=\pi$. Given the $\mathfrak t^{(\ell)}_n$, let $\mathfrak t^{(\ell)}_{\vec n}$ be the projective seminorm of the spaces $(\Gamma^{n_i}(P),\mathfrak t^{(\ell)}_{n_i})$:
\begin{equation}\label{dp:eq:tree-root}
 \mathfrak t^{(\ell)}_{\vec n}(t)=\inf\Bigl\{\sum_q\prod_{i=1}^r\mathfrak t^{(\ell)}_{n_i}(s_{q,i}):\ t=\sum_qs_{q,1}\otimes\dots\otimes s_{q,r},\ s_{q,i}\in\Gamma^{n_i}(P)\Bigr\}.
\end{equation}
Put $\mathfrak t^{(\ell+1)}_1=\norm\cdot$, and $\mathfrak t^{(\ell+1)}_N=\max_{r\ge2}\max_{\vec n\in\mathcal C(N,r)}\mathfrak t^{(\ell)}_{\vec n}$ for $N\ge2$. Also $\mathfrak t^{(\ell)}_0=\abs\cdot$ on $\Gamma^0(P)=K$. We write $\mathfrak t^{(\ell)}_N(P)$ when the space must be named.
\end{definition}

Note that $\mathfrak t^{(1)}_{\vec n}=\pi_{\vec n}$ and $\mathfrak t^{(\ell)}_{(N)}=\mathfrak t^{(\ell)}_N$. The recursion runs over rooted trees with $N$ leaves: a composition is the root, and the compositions of its parts are the next level.

\begin{lemma}[stabilization]\label{dp:lem:tree-stable}
Let $P$ be normed. We have $\mathfrak t^{(\ell)}_1=\norm\cdot$ for all $\ell\ge1$. For $N\ge2$ and $\ell\ge N-1$, the seminorm $\mathfrak t^{(\ell)}_N$ does not depend on $\ell$.
\end{lemma}

\begin{proof}
For the first claim, $\pi$ on $\Gamma^1(P)=P$ is $\norm\cdot$. For the second use induction on $N$. For $N=2$ the only composition with $r\ge2$ is $(1,1)$, and $\mathfrak t^{(\ell)}_{(1,1)}=\pi$, so $\mathfrak t^{(\ell)}_2=\pi$ for all $\ell$. For $N\ge3$ and $\ell\ge N-1$, the seminorm $\mathfrak t^{(\ell)}_N$ is built from the $\mathfrak t^{(\ell-1)}_{n_i}$ with $n_i\le N-1$ and $\ell-1\ge n_i-1$, and these do not depend on $\ell$ by induction.
\end{proof}

Put $\mathfrak t_N=\mathfrak t^{(N)}_N$ for $N\ge1$ and $\mathfrak t_0=\abs\cdot$ on $\Gamma^0(P)=K$, so that $\mathfrak t^{(\ell)}_n=\mathfrak t_n$ for $\ell\ge\max(n,1)$. Let $\mathfrak t_{\vec n}$ be~\eqref{dp:eq:tree-root} with $\mathfrak t_{n_i}$ in place of $\mathfrak t^{(\ell)}_{n_i}$. Then $\mathfrak t^{(\ell)}_{\vec n}=\mathfrak t_{\vec n}$ for $\ell\ge N$, so
\begin{equation}\label{dp:eq:tree-fixed}
 \mathfrak t_N=\max_{r\ge2}\max_{\vec n\in\mathcal C(N,r)}\mathfrak t_{\vec n}\qquad(N\ge2).
\end{equation}

We use the shuffle product $s\cdot t$ of Section~\ref{dp:sec:conv}; it maps $\Gamma^m(P)\times\Gamma^{m'}(P)$ into $\Gamma^{m+m'}(P)$. We read $1\cdot t=t$.

\begin{lemma}[properties of the tree norms]\label{dp:lem:tree}
Let $P$ be normed and $\ell,N\ge1$.
\begin{enumerate}
\item $\pi\le\mathfrak t^{(\ell)}_N\le\mathfrak t^{(\ell+1)}_N\le\mathfrak t_N\le\norm\cdot_\Gamma$, and $\mathfrak t^{(\ell)}_N(h^{\otimes N})\le\norm h^N$.
\item For a bounded linear map $T\colon P\to P'$ we have $\mathfrak t^{(\ell)}_N(P')(T^{\otimes N}t)\le\norm T^N\mathfrak t^{(\ell)}_N(P)(t)$. The same holds with $\mathfrak t_N$ in place of $\mathfrak t^{(\ell)}_N$.
\item $\mathfrak t^{(2)}_N=\max_{\vec n}\pi_{\vec n}$ over the $\vec n$ with $r\ge2$, for $N\ge2$. On $\ell^1(J,K)$ we have $\mathfrak t^{(\ell)}_N=\mathfrak t_N=\pi$.
\item If $\operatorname{char}K=p$ and $N\le p$, then $\mathfrak t_N\le\mathsf c_N\pi$, with $\mathsf c_N$ depending only on $N$.
\item \textup{(Shuffle bound.)} For $s\in\Gamma^m(P)$, $t\in\Gamma^{m'}(P)$, $m,m'\ge0$,
\begin{equation}\label{dp:eq:shuffle}
 \mathfrak t^{(\ell)}_{m+m'}(s\cdot t)\le\binom{m+m'}m\mathfrak t^{(\ell)}_m(s)\,\mathfrak t^{(\ell)}_{m'}(t).
\end{equation}
So, for $x\in P$, $\mathfrak t^{(\ell)}_{m+m'}(x^{\otimes m}\cdot t)\le\binom{m+m'}m\norm x^m\mathfrak t^{(\ell)}_{m'}(t)$. Both bounds hold with $\mathfrak t_{m+m'},\mathfrak t_m,\mathfrak t_{m'}$ in place of the seminorms of level $\ell$.
\end{enumerate}
\end{lemma}

\begin{proof}
\emph{Step 1 (item (1)).} First $\mathfrak t^{(\ell)}_N\le\norm\cdot_\Gamma$, by induction on $\ell$. For $\ell=1$ this is $\pi\le\norm\cdot_\Gamma$. For the step let $N\ge2$, $\vec n\in\mathcal C(N,r)$, $r\ge2$. A representation $t=\sum_qc_qh_q^{\otimes N}$ gives $t=\sum_qc_qh_q^{\otimes n_1}\otimes\dots\otimes h_q^{\otimes n_r}$, of cost at most $\sum_q\abs{c_q}\norm{h_q}^N$ in~\eqref{dp:eq:tree-root}, since $\mathfrak t^{(\ell)}_{n_i}\le\norm\cdot_\Gamma$. Taking $t=h^{\otimes N}$ gives the bound on powers. Next $\mathfrak t^{(\ell)}_N\le\mathfrak t^{(\ell+1)}_N$, by induction on $\ell$. For $\ell=1$ and $N\ge2$, Lemma~\ref{dp:lem:kappa-bounds}(1) gives $\pi\le\pi_{\vec n}=\mathfrak t^{(1)}_{\vec n}\le\mathfrak t^{(2)}_N$. For the step, $\mathfrak t^{(\ell-1)}_{n_i}\le\mathfrak t^{(\ell)}_{n_i}$ gives $\mathfrak t^{(\ell-1)}_{\vec n}\le\mathfrak t^{(\ell)}_{\vec n}$, hence $\mathfrak t^{(\ell)}_N\le\mathfrak t^{(\ell+1)}_N$. Finally $\mathfrak t^{(\ell)}_N\le\mathfrak t_N$ by monotonicity and Lemma~\ref{dp:lem:tree-stable}.

\emph{Step 2 (item (2)).} Induction on $\ell$. The case $\ell=1$ is the same statement for $\pi$. For the step, $T^{\otimes N}$ carries a representation in~\eqref{dp:eq:tree-root} to the representation with factors $T^{\otimes n_i}s_{q,i}$, and $\prod_i\norm T^{n_i}=\norm T^N$. The statement for $\mathfrak t_N$ is the case $\ell=N$.

\emph{Step 3 (item (3)).} The first assertion is the definition, as $\mathfrak t^{(1)}_{\vec n}=\pi_{\vec n}$. On $\ell^1(J,K)$ use induction on $N$. If $\mathfrak t_n=\pi$ for $n<N$, then $\mathfrak t_{\vec n}=\pi_{\vec n}$ for $r\ge2$, and $\pi_{\vec n}\le\pi$ on $\Gamma^N(P)$ by Lemma~\ref{dp:lem:kappa-bounds}(3). So $\mathfrak t_N\le\pi$ by~\eqref{dp:eq:tree-fixed}, and item~(1) gives the rest.

\emph{Step 4 (item (4)).} Let $\mathsf c_1=1$. Let $2\le N\le p$ and $r\ge2$. All parts of $\vec n$ are below $p$, so $\pi_{\vec n}\le\prod_in_i!\,\pi$ by Lemma~\ref{dp:lem:kappa-bounds}(2). By induction $\mathfrak t_{n_i}\le\mathsf c_{n_i}\pi$, so $\mathfrak t_{\vec n}\le\prod_i\mathsf c_{n_i}\,\pi_{\vec n}$. By~\eqref{dp:eq:tree-fixed} we may take $\mathsf c_N=\max_{\vec n}\prod_i\mathsf c_{n_i}n_i!$.

\emph{Step 5 (item (5), set-up).} The cases $m=0$ or $m'=0$ are trivial. Let $m,m'\ge1$ and use induction on $\ell$. For $\ell=1$: each shuffle is an isometry for $\pi$, and $\pi(s\otimes t)\le\pi(s)\pi(t)$. Assume~\eqref{dp:eq:shuffle} at level $\ell$. Put $N=m+m'$ and let $\vec n\in\mathcal C(N,r)$, $r\ge2$. It is enough to show $\mathfrak t^{(\ell)}_{\vec n}(s\cdot t)\le\binom Nm\mathfrak t^{(\ell+1)}_m(s)\mathfrak t^{(\ell+1)}_{m'}(t)$.

\emph{Step 6 (item (5), the blocks).} A shuffle of $m$ and $m'$ letters into $N$ slots, divided into blocks of sizes $n_i$, determines the numbers $m_i,m'_i$ of letters of each kind in block $i$, with $m_i+m'_i=n_i$, and a shuffle in each block. The letters of each kind enter in order. Fix a vector $(m_i)_i$ with $0\le m_i\le n_i$ and $\sum_im_i=m$. Read $s$ in $\bigotimes_{m_i>0}\Gamma^{m_i}(P)$ and $t$ in $\bigotimes_{m'_i>0}\Gamma^{m'_i}(P)$, with representations $s=\sum_q\bigotimes_is_{q,i}$ and $t=\sum_{q'}\bigotimes_it_{q',i}$; where $m_i=0$ or $m'_i=0$, read the missing factor as $1$. The sum of the shuffles with this vector is then
\begin{equation}\label{dp:eq:shuffle-blocks}
 \sum_{q,q'}\bigotimes_{i=1}^r\bigl(s_{q,i}\cdot t_{q',i}\bigr).
\end{equation}
By~\eqref{dp:eq:shuffle} at level $\ell$, the cost of~\eqref{dp:eq:shuffle-blocks} in~\eqref{dp:eq:tree-root} is at most
\[
 \prod_i\binom{n_i}{m_i}\Bigl(\sum_q\prod_i\mathfrak t^{(\ell)}_{m_i}(s_{q,i})\Bigr)\Bigl(\sum_{q'}\prod_i\mathfrak t^{(\ell)}_{m'_i}(t_{q',i})\Bigr).
\]
If $(m_i)_i$ has at least two nonzero parts, the infimum of the first factor over representations is $\mathfrak t^{(\ell)}_{\vec m}(s)$, for the composition $\vec m$ of $m$ formed by the nonzero $m_i$, and this is at most $\mathfrak t^{(\ell+1)}_m(s)$. If it has one nonzero part, use $s$ itself: the factor is $\mathfrak t^{(\ell)}_m(s)\le\mathfrak t^{(\ell+1)}_m(s)$. The same holds for $t$. Summing over the vectors $(m_i)_i$, and using $\sum\prod_i\binom{n_i}{m_i}=\binom Nm$, we get $\mathfrak t^{(\ell)}_{\vec n}(s\cdot t)\le\binom Nm\mathfrak t^{(\ell+1)}_m(s)\mathfrak t^{(\ell+1)}_{m'}(t)$. Take the maximum over $\vec n$. The second bound follows with item~(1). For $\mathfrak t$, take $\ell=\max(m+m',1)$: then the seminorms of level $\ell$ in degrees $m+m'$, $m$, $m'$ are $\mathfrak t_{m+m'},\mathfrak t_m,\mathfrak t_{m'}$.
\end{proof}

\subsection{Tree analytic germs}\label{dp:ssec:ta}

\begin{definition}[tree analytic germs, the composition closure]\label{dp:def:ta}
Let $U\subseteq P$ be open, $x_0\in U$, $R>0$ with $B(x_0,R)\subseteq U$, $M\ge0$, $\ell\ge1$.
\begin{enumerate}
\item A map $f\colon U\to W$ is $\mathrm{TA}_\ell$ at $x_0$ with radius $R$ and bound $M$ if $f(x_0+h)=\sum_{N\ge0}L_N(h^{\otimes N})$ for $\norm h<R$, with linear maps $L_N\colon\Gamma^N(P)\to W$ such that
\begin{equation}\label{dp:eq:ta-bound}
 \norm{L_Nt}\le MR^{-N}\mathfrak t^{(\ell)}_N(t)\qquad(N\ge0,\ t\in\Gamma^N(P)).
\end{equation}
It is \emph{tree analytic}, or TA, at $x_0$ with radius $R$ and bound $M$ if this holds with $\mathfrak t_N$ in place of $\mathfrak t^{(\ell)}_N$. It is $\mathrm{TA}_\ell$ or TA at $x_0$ if this holds for some $R,M$, and on $U$ if it holds at every point of $U$.
\item $\mathrm{AA}^\circ$ is the smallest class of germs of maps from open subsets of normed spaces into Banach spaces that contains the analytic germs and is closed under composition, under pairing $(f_1,f_2)$ into the sum with the maximum norm, and under target descent: if $j\colon W\to Z$ is a linear isometry of Banach spaces and $j\circ f$ is in the class, then so is $f$.
\end{enumerate}
\end{definition}

\begin{remark}[first properties]\label{dp:rem:ta-basic}
Let $f\colon U\to W$ be $\mathrm{TA}_\ell$ or TA at $x_0$ with radius $R$ and bound $M$, with linear maps $L_N$.
\begin{enumerate}
\item The same holds with every smaller radius and every larger bound. If $f$ is $\mathrm{TA}_\ell$, it is $\mathrm{TA}_{\ell'}$ for $\ell'\ge\ell$, and it is TA, with the same radius and bound.
\item By Lemma~\ref{dp:lem:tree}(1) and~\cite[Lemma~Q.1(2)]{I}, $a_N(h)=L_N(h^{\otimes N})$ is a law with $\norm{a_N}\le MR^{-N}$. So $f$ is DP at $x_0$ with radius $R$ and bound $M$, and the $L_N=L_{a_N}$ are the linearizations of its coefficient laws. Hence $\mathrm{TA}_\ell\subseteq\mathrm{TA}_{\ell+1}\subseteq\mathrm{TA}\subseteq\mathrm{DP}$, with the same radius and bound.
\item By~\cite[Proposition~14.9(1)]{I}, $\mathrm{TA}_1=\mathrm{AA}$, with the same radius and bound.
\item An AA germ into a Banach space $W$ is in $\mathrm{AA}^\circ$: the enlargement $Z$ may be taken Banach~\cite[text after Definition~14.7]{I}, then $j\circ f$ is analytic, and target descent applies.
\end{enumerate}
\end{remark}

The next lemma is the estimate for the map $\mathsf Q_{r,N}$ of~\cite[(Q.8)]{I}. The algebra of~\cite[(Q.8)]{I} follows from~\cite[VI~\S8, Prop.~VI.6, Cor.; Prop.~I.2; II~\S1]{Ro}, see~\cite[Section~Q.4]{I}; the bound below is the estimate we need.

\begin{lemma}[grafting]\label{dp:lem:grafting}
Let $P,W$ be normed. For $n\ge1$ let $L_n\colon\Gamma^n(P)\to W$ be linear and $\mathsf m_n\ge0$, with $\norm{L_nt}\le\mathsf m_n\mathfrak t_n(t)$ for $t\in\Gamma^n(P)$. For $1\le r\le N$ put
\[
 \mathsf Q_{r,N}=\sum_{\vec n\in\mathcal C(N,r)}L_{n_1}\otimes\dots\otimes L_{n_r},\qquad
 \mathsf B_{r,N}=\sum_{\vec n\in\mathcal C(N,r)}\prod_i\mathsf m_{n_i}.
\]
Then $\mathsf Q_{r,N}$ maps $\Gamma^N(P)$ to $\Gamma^r(W)$, and:
\begin{enumerate}
\item if $\ell\ge1$ and $\norm{L_nt}\le\mathsf m_n\mathfrak t^{(\ell)}_n(t)$ for all $n$ and $t$, then for all $e\ge1$, $\mathfrak t^{(e)}_r(W)(\mathsf Q_{r,N}t)\le\mathsf B_{r,N}\mathfrak t^{(\ell+e)}_N(P)(t)$;
\item $\mathfrak t_r(W)(\mathsf Q_{r,N}t)\le\mathsf B_{r,N}\mathfrak t_N(P)(t)$;
\item if each $L_n$ is the restriction to $\Gamma^n(P)$ of a linear map $\tilde L_n\colon(P^{\otimes n},\pi)\to W$ of norm at most $\mathsf m_n$, then for all $e\ge1$, $\mathfrak t^{(e)}_r(W)(\mathsf Q_{r,N}t)\le\mathsf B_{r,N}\mathfrak t^{(e)}_N(P)(t)$.
\end{enumerate}
\end{lemma}

\begin{proof}
\emph{Step 1 (values).} By Lemma~\ref{dp:lem:tree}(1) and~\cite[Lemma~Q.1(2)]{I}, $h\mapsto L_n(h^{\otimes n})$ is a law. The hypotheses of items~(1) and~(3) imply the standing bound $\norm{L_nt}\le\mathsf m_n\mathfrak t_n(t)$, by Lemma~\ref{dp:lem:tree}(1). So Steps~2 and~3 of the proof of~\cite[Lemma~Q.4]{I} show that $\mathsf Q_{r,N}$ has values in $\Gamma^r(W)$.

\emph{Step 2 (item (1), base cases).} Use induction on $e$, for all $r,N$. If $r=1$, then $\mathsf Q_{1,N}=L_N$, $\mathsf B_{1,N}=\mathsf m_N$ and $\mathfrak t^{(e)}_1=\norm\cdot$, and $\mathfrak t^{(\ell)}_N\le\mathfrak t^{(\ell+e)}_N$. Let $e=1$, $r\ge2$. Applying $\bigotimes_iL_{n_i}$ to a representation of $t$ in~\eqref{dp:eq:tree-root} gives
\[
 \pi\Bigl(\bigl(\bigotimes_iL_{n_i}\bigr)t\Bigr)\le\prod_i\mathsf m_{n_i}\,\mathfrak t^{(\ell)}_{\vec n}(t)\le\prod_i\mathsf m_{n_i}\,\mathfrak t^{(\ell+1)}_N(t).
\]
Sum over $\vec n$.

\emph{Step 3 (item (1), the step).} Assume item~(1) at level $e$, and let $r\ge2$. Let $\vec r\in\mathcal C(r,r')$, $r'\ge2$. A composition $\vec n\in\mathcal C(N,r)$ cut into consecutive blocks of lengths $r_1,\dots,r_{r'}$ has block sums $\vec N\in\mathcal C(N,r')$, and it is determined by $\vec N$ and the blocks, which are elements of the $\mathcal C(N_i,r_i)$. So, reading $t$ in $\bigotimes_i\Gamma^{N_i}(P)$,
\begin{equation}\label{dp:eq:grafting-cut}
 \mathsf Q_{r,N}(t)=\sum_{\vec N\in\mathcal C(N,r')}\Bigl(\bigotimes_{i=1}^{r'}\mathsf Q_{r_i,N_i}\Bigr)(t),\qquad
 \mathsf B_{r,N}=\sum_{\vec N\in\mathcal C(N,r')}\prod_i\mathsf B_{r_i,N_i},
\end{equation}
with empty sums for $N_i<r_i$. Let $\delta>0$. For each $\vec N$ take a representation $t=\sum_q\bigotimes_is_{q,i}$, $s_{q,i}\in\Gamma^{N_i}(P)$, with
\[
 \sum_q\prod_i\mathfrak t^{(\ell+e)}_{N_i}(s_{q,i})\le\mathfrak t^{(\ell+e)}_{\vec N}(t)+\delta\le\mathfrak t^{(\ell+e+1)}_N(t)+\delta.
\]
Then $\sum_q\bigotimes_i\mathsf Q_{r_i,N_i}(s_{q,i})$ represents $(\bigotimes_i\mathsf Q_{r_i,N_i})t$ in $\bigotimes_i\Gamma^{r_i}(W)$. By item~(1) at level $e$, its cost in~\eqref{dp:eq:tree-root} for $W$ at level $e$ is at most $\prod_i\mathsf B_{r_i,N_i}(\mathfrak t^{(\ell+e+1)}_N(t)+\delta)$. Sum over $\vec N$ with~\eqref{dp:eq:grafting-cut}, let $\delta\to0$, and take the maximum over $\vec r$.

\emph{Step 4 (item (2)).} Fix $N$. Replacing $L_n$ by $0$ for $n>N$ changes neither $\mathsf Q_{r,N}$ nor $\mathsf B_{r,N}$. By Lemma~\ref{dp:lem:tree-stable}, $\mathfrak t_n=\mathfrak t^{(N)}_n$ for $n\le N$, so the hypothesis of item~(1) holds with $\ell=N$. Item~(1) with $\ell=e=N$ gives item~(2), as $\mathfrak t^{(N)}_r=\mathfrak t_r$ for $r\le N$ and $\mathfrak t^{(2N)}_N=\mathfrak t_N$.

\emph{Step 5 (item (3)).} The induction is that of Steps~2 and~3, with $\mathfrak t^{(e)}$ in place of $\mathfrak t^{(\ell+e)}$. If $r=1$: $\norm{L_Nt}\le\mathsf m_N\pi(t)\le\mathsf m_N\mathfrak t^{(e)}_N(t)$. If $e=1$, $r\ge2$: $\mathsf Q_{r,N}$ is the restriction to $\Gamma^N(P)$ of $\sum_{\vec n}\bigotimes_i\tilde L_{n_i}$ on $P^{\otimes N}$; applied to a representation of $t$ in $P^{\otimes N}$, this gives $\pi(\mathsf Q_{r,N}t)\le\mathsf B_{r,N}\pi(t)$. For the step, choose the representations with $\sum_q\prod_i\mathfrak t^{(e)}_{N_i}(s_{q,i})\le\mathfrak t^{(e)}_{\vec N}(t)+\delta\le\mathfrak t^{(e+1)}_N(t)+\delta$.
\end{proof}

\begin{theorem}[composition]\label{dp:thm:ta-composition}
Let $\ell,e\ge1$. Let $f\colon U\to W$ be $\mathrm{TA}_\ell$ at $x_0$ with radius $R$ and bound $M>0$, and let $g$ map an open subset of $W$ into a normed space $Y$, $\mathrm{TA}_e$ at $f(x_0)$ with radius $R'$ and bound $M'$. Let $\Xi_N$ be the laws of $g\circ f$ at $x_0$, as in~\cite[Lemma~Q.3]{I}.
\begin{enumerate}
\item For $N\ge1$ and $t\in\Gamma^N(P)$,
\begin{equation}\label{dp:eq:ta-composition}
 \norm{L_{\Xi_N}t}\le M'\frac M{R'}\Bigl(1+\frac M{R'}\Bigr)^{N-1}R^{-N}\mathfrak t^{(\ell+e)}_N(t),
\end{equation}
and $g\circ f$ is $\mathrm{TA}_{\ell+e}$ at $x_0$ with radius $R_1=\min(R/2,RR'/(8C_0M))$ and bound $M'$.
\item The same holds with TA in place of $\mathrm{TA}_\ell$, $\mathrm{TA}_e$, $\mathrm{TA}_{\ell+e}$, and $\mathfrak t_N$ in~\eqref{dp:eq:ta-composition}.
\item If $f$ is analytic at $x_0$ with radius $R$ and bound $M$, then~\eqref{dp:eq:ta-composition} holds with $\mathfrak t^{(e)}_N$, and $g\circ f$ is $\mathrm{TA}_e$ at $x_0$ with radius $R_1$ and bound $M'$.
\item $\mathrm{TA}_\ell$ and TA are closed under pairing, sums, bounded linear postcomposition, precomposition with analytic maps, restriction to open subsets, and target descent along linear isometries $j\colon W\to Z$ for which the laws of $j\circ f$ lie in $j(W)$.
\end{enumerate}
\end{theorem}

\begin{proof}
\emph{Step 1 (the estimate).} Let $a_n$, $b_r$ be the coefficient laws of $f$, $g$. By Remark~\ref{dp:rem:ta-basic}(2), both germs are DP, so~\cite[Lemmas~Q.3, Q.4]{I} apply, with $b_r$ for the law called $d_r$ there: for $\norm h<R_1$, $g(f(x_0+h))=\sum_N\Xi_N(h)$, and $L_{\Xi_N}=\sum_{r=1}^NL_{b_r}\circ\mathsf Q_{r,N}$. Apply Lemma~\ref{dp:lem:grafting}(1) with $L_n=L_{a_n}$ and $\mathsf m_n=MR^{-n}$. Then $\mathsf B_{r,N}=\binom{N-1}{r-1}M^rR^{-N}$, since $\abs{\mathcal C(N,r)}=\binom{N-1}{r-1}$. By~\eqref{dp:eq:ta-bound} for $g$,
\[
 \norm{L_{b_r}\mathsf Q_{r,N}t}\le M'R'^{-r}\mathfrak t^{(e)}_r(W)(\mathsf Q_{r,N}t)\le M'R'^{-r}\binom{N-1}{r-1}M^rR^{-N}\mathfrak t^{(\ell+e)}_N(t).
\]
Summing over $r$ gives~\eqref{dp:eq:ta-composition}.

\emph{Step 2 (radius and bound).} As $C_0\ge1$, we have $1+M/R'\le\max(2,8C_0M/R')=R/R_1$. So~\eqref{dp:eq:ta-composition} gives $\norm{L_{\Xi_N}t}\le M'(1+M/R')^NR^{-N}\mathfrak t^{(\ell+e)}_N(t)\le M'R_1^{-N}\mathfrak t^{(\ell+e)}_N(t)$ for $N\ge1$. For $N=0$, $\Xi_0=g(f(x_0))$ has norm at most $M'$. This proves item~(1).

\emph{Step 3 (items (2), (3)).} For item~(2) use Lemma~\ref{dp:lem:grafting}(2) in Step~1. For item~(3), by~\cite[Proposition~14.9(1)]{I} with $j=\id_W$, the law $a_n$ is the diagonal of a bounded $n$-linear map of norm at most $MR^{-n}$. Let $\tilde L_n$ be its linearization on $(P^{\otimes n},\pi)$; then $L_{a_n}=\tilde L_n|_{\Gamma^n(P)}$. Use Lemma~\ref{dp:lem:grafting}(3) in Step~1.

\emph{Step 4 (item (4)).} By Remark~\ref{dp:rem:ta-basic}(1) we may take a common radius and level for two germs, and raise a bound to make it positive. For a pairing, the linearizations are pairs, and the norm of the sum is the maximum. For sums and bounded linear postcomposition the bound~\eqref{dp:eq:ta-bound} is immediate. Precomposition with analytic maps is item~(3), and for TA item~(2), as analytic germs are TA. Restriction is clear. For target descent, let $L^Z_N$ be the linearizations of the laws of $j\circ f$. By hypothesis $L^Z_N$ has values in $j(W)$, since $\Gamma^N(P)$ is spanned by the powers. Put $L_N=j^{-1}L^Z_N$. Then $\norm{L_Nt}=\norm{L^Z_Nt}$, and the series $\sum_NL_N(h^{\otimes N})$ converges to $f(x_0+h)$, since $j$ is an isometry.
\end{proof}

\begin{theorem}[the composition closure]\label{dp:thm:composition-closure}
Over every $K$, every germ in $\mathrm{AA}^\circ$ is in $\bigcup_\ell\mathrm{TA}_\ell$. If $K$ is complete, then $\mathrm{AA}^\circ$ is the class of germs in $\bigcup_\ell\mathrm{TA}_\ell$ with Banach values.
\end{theorem}

\begin{proof}
\emph{Step 1 (the inclusion $\subseteq$).} The class $\bigcup_\ell\mathrm{TA}_\ell$ contains the analytic germs, by Remark~\ref{dp:rem:ta-basic}(3), and Theorem~\ref{dp:thm:ta-composition} gives closure under composition and pairing. For target descent, let $j\colon W\to Z$ be a linear isometry of Banach spaces, and let $j\circ f$ be $\mathrm{TA}_\ell$. It is DP, by Remark~\ref{dp:rem:ta-basic}(2), and $j(W)$ is closed. So Theorem~\ref{dp:thm:target-descent}(2) puts the laws of $j\circ f$ in $j(W)$, and Theorem~\ref{dp:thm:ta-composition}(4) applies.

From here on $K$ is complete. Fix $P$ and $x_0\in P$.

\emph{Step 2 (universal germs).} Let $\mathbf W_\ell$ be the completed $\ell^1$-sum of the separated completions of $(\Gamma^N(P),\mathfrak t^{(\ell)}_N)$, $N\ge1$, with $\iota^{(\ell)}_N\colon\Gamma^N(P)\to\mathbf W_\ell$ the natural maps. Put $\mathsf v_\ell(h)=\sum_{N\ge1}\iota^{(\ell)}_N(h^{\otimes N})$ for $\norm h<1$. By Lemma~\ref{dp:lem:tree}(1), $\mathsf v_\ell$ is $\mathrm{TA}_\ell$ at $0$ with radius $1$ and bound $1$, and $\norm{\mathsf v_\ell(h)}<1$ for $\norm h<\frac12$. Its coefficient laws are $a_n(h)=\iota^{(\ell)}_n(h^{\otimes n})$.

\emph{Step 3 (one more level).} Let $\mathsf s\in K^\times$. We show by induction that the germ at $x_0$ of $\mathsf v'_\ell(x)=\mathsf v_\ell(\mathsf s(x-x_0))$ is in $\mathrm{AA}^\circ$. By Theorem~\ref{dp:thm:ta-composition}(3), $\mathsf v'_1$ is AA, so it is in $\mathrm{AA}^\circ$, by Remark~\ref{dp:rem:ta-basic}(4). Let $\mathsf v'_\ell$ be in $\mathrm{AA}^\circ$. Let $Y$ be the completed $\ell^1$-sum of the spaces $\widehat\otimes^r_\pi\mathbf W_\ell$, $r\ge1$, with inclusions $\iota'_r$, and $g(w)=\sum_{r\ge1}\iota'_r(w^{\otimes r})$ for $\norm w<1$. This $g$ is analytic at $0$ with radius $1$ and bound $1$, and $L_{b_r}=\iota'_r$ on $\Gamma^r(\mathbf W_\ell)$. By~\cite[Lemma~Q.4]{I}, the composite $g\circ\mathsf v_\ell$ has laws $\Xi_N$ with $L_{\Xi_N}=\sum_r\iota'_r\mathsf Q_{r,N}$, where $(\bigotimes_iL_{a_{n_i}})t$ is the image of $t$ in the block of $\mathbf W_\ell^{\otimes r}$ indexed by $\vec n$. Blocks with different $\vec n$ are disjoint, also for different $N$. By~\cite[Lemma~Q.5(2),(1)]{I}, for finitely many $t_N\in\Gamma^N(P)$,
\begin{equation}\label{dp:eq:closure-norm}
 \Bigl\lVert\sum_NL_{\Xi_N}t_N\Bigr\rVert_Y=\sum_N\mathfrak x_N(t_N),\qquad
 \mathfrak x_N(t)=\sum_{r=1}^N\sum_{\vec n\in\mathcal C(N,r)}\mathfrak t^{(\ell)}_{\vec n}(t)\ge\mathfrak t^{(\ell+1)}_N(t).
\end{equation}
Let $Y_0$ be the closed span in $Y$ of the ranges of the $L_{\Xi_N}$. The germ $g\circ\mathsf v'_\ell$ is in $\mathrm{AA}^\circ$. By~\cite[Lemma~Q.3]{I}, near $x_0$ it equals $\sum_{N\ge1}L_{\Xi_N}((\mathsf s(x-x_0))^{\otimes N})$, so it has values in $Y_0$, and by target descent its germ into $Y_0$ is in $\mathrm{AA}^\circ$. By~\eqref{dp:eq:closure-norm} the assignment $L_{\Xi_N}t\mapsto\iota^{(\ell+1)}_Nt$ is well defined on the span of the ranges and has norm at most $1$. It extends to a bounded linear map $\mathbf c\colon Y_0\to\mathbf W_{\ell+1}$, and, applying the continuous map $\mathbf c$ term by term to this series, $\mathbf c\circ g\circ\mathsf v'_\ell=\mathsf v'_{\ell+1}$ near $x_0$. A bounded linear map is analytic, so $\mathsf v'_{\ell+1}$ is in $\mathrm{AA}^\circ$.

\emph{Step 4 (general germs).} Let $f\colon U\to W$ be $\mathrm{TA}_\ell$ at $x_0$ with radius $R$ and bound $M$, with $W$ Banach and linear maps $L_N$. Choose $\mathsf s\in K$ with $\abs{\mathsf s}\ge1/R$. By~\eqref{dp:eq:ta-bound}, $\norm{\mathsf s^{-N}L_Nt}\le M\mathfrak t^{(\ell)}_N(t)$ for $N\ge1$. So $\iota^{(\ell)}_N(t)\mapsto\mathsf s^{-N}L_Nt$ extends to a bounded linear map $\mathbf L\colon\mathbf W_\ell\to W$, and
\[
 f(x)=f(x_0)+\mathbf L\bigl(\mathsf v_\ell(\mathsf s(x-x_0))\bigr)\qquad(\norm{x-x_0}<\abs{\mathsf s}^{-1}).
\]
The map $w\mapsto f(x_0)+\mathbf Lw$ is analytic between Banach spaces. So $f$ is in $\mathrm{AA}^\circ$ by Step~3.
\end{proof}

For $K$ complete, Step~3 shows that $\mathsf v'_\ell$ arises from the AA germ $\mathsf v'_1$ by $\ell-1$ steps, each a composition with an analytic germ, a target descent and a bounded linear map.

\begin{theorem}[re-expansion and derivatives]\label{dp:thm:ta-recenter}
Let $W$ be Banach and let $f$ be $\mathrm{TA}_\ell$ at $x_0$ with radius $R$ and bound $M$. Let $x_1\in B(x_0,R)$ and $\varrho=\norm{x_1-x_0}$. Then $f$ is $\mathrm{TA}_\ell$ at $x_1$ with radius $R-\varrho$ and bound $MR/(R-\varrho)$. Moreover $f$ is $C^\infty$ on $B(x_0,R)$, and for each $\mathsf d\ge1$ and $0<R_2<R$, the map $D^{\mathsf d}f$ is $\mathrm{TA}_\ell$ at $x_0$ with radius $R_2$, into $\Mult^{\mathsf d}(P;W)$. The same holds with TA in place of $\mathrm{TA}_\ell$.
\end{theorem}

\begin{proof}
\emph{Step 1 (re-expansion).} Let $L_N$ be as in Definition~\ref{dp:def:ta}, $\xi=x_1-x_0$. For $\norm h<R-\varrho$, \eqref{dp:eq:orbit} gives $(\xi+h)^{\otimes N}=\sum_{n=0}^N\xi^{\otimes(N-n)}\cdot h^{\otimes n}$. For $t\in\Gamma^n(P)$ put $L'_n(t)=\sum_{N\ge n}L_N(\xi^{\otimes(N-n)}\cdot t)$. By Lemma~\ref{dp:lem:tree}(5),
\[
 \norm{L'_nt}\le\sum_{N\ge n}MR^{-N}\binom Nn\varrho^{N-n}\mathfrak t^{(\ell)}_n(t)=MR^{-n}\Bigl(1-\frac\varrho R\Bigr)^{-n-1}\mathfrak t^{(\ell)}_n(t)=\frac{MR}{R-\varrho}(R-\varrho)^{-n}\mathfrak t^{(\ell)}_n(t).
\]
The sum over $N$ in this display is a sum of real numbers, and the series defining $L'_nt$ converges in the Banach space $W$. The double series is absolutely summable, since $\sum_N\sum_{n\le N}MR^{-N}\binom Nn\varrho^{N-n}\norm h^n=M\sum_N((\varrho+\norm h)/R)^N$ is finite. So $f(x_1+h)=\sum_nL'_n(h^{\otimes n})$.

\emph{Step 2 (the first derivative).} By Step~1, $f$ is DP on $B(x_0,R)$, so it is $C^\infty$ there by Theorem~\ref{dp:thm:differentiable}, and $Df(x_1)$ is the law of degree $1$ at $x_1$. For $\norm h<R$ and $z\in P$, Step~1 with $x_1=x_0+h$ gives $Df(x_0+h)z=\sum_{N\ge1}L_N(h^{\otimes(N-1)}\cdot z)$. For $n\ge0$ and $t\in\Gamma^n(P)$ put $\mathsf D_n(t)z=L_{n+1}(t\cdot z)$. By~\eqref{dp:eq:shuffle}, $\norm{\mathsf D_nt}\le(n+1)MR^{-n-1}\mathfrak t^{(\ell)}_n(t)$ in $\cL(P,W)$. For $0<R_2<R$ the number $M_2=\sup_n(n+1)MR^{-n-1}R_2^n$ is finite, so the $\mathsf D_n$ satisfy~\eqref{dp:eq:ta-bound} with radius $R_2$ and bound $M_2$. For $\norm h<R_2$ the series $\sum_n\mathsf D_n(h^{\otimes n})$ converges absolutely in the Banach space $\cL(P,W)$, and by the formula above its sum is $Df(x_0+h)$. So $Df$ is $\mathrm{TA}_\ell$ at $x_0$ with radius $R_2$ and bound $M_2$.

\emph{Step 3 (higher derivatives).} Use induction on $\mathsf d$. Let $D^{\mathsf d}f$ be $\mathrm{TA}_\ell$ at $x_0$ with every radius below $R$, and let $0<R_2<R$. The space $\Mult^{\mathsf d}(P;W)$ is Banach. Apply Step~2 to $D^{\mathsf d}f$, with a radius between $R_2$ and $R$, and use the isometry $\cL(P,\Mult^{\mathsf d}(P;W))=\Mult^{\mathsf d+1}(P;W)$. The argument for TA is the same with $\mathfrak t_n$, by Lemma~\ref{dp:lem:tree}(5).
\end{proof}

So, among germs with Banach values, TA and each $\mathrm{TA}_\ell$ are open conditions, and so is $\mathrm{AA}^\circ$ if $K$ is complete, by Theorem~\ref{dp:thm:composition-closure}. By Theorem~\ref{dp:thm:whole-ball}, a DP germ with Banach values is also DP on its whole ball. The difference is the radius. Theorem~\ref{dp:thm:ta-recenter} gives $R-\varrho$ over every $K$, for a radius $R$ in the sense of~\eqref{dp:eq:ta-bound}. For a radius in the sense of DP this formula can fail over discretely valued $K$, by Example~\ref{dp:ex:radius-drop}. The two proofs differ in the cost of mixed terms. By~\eqref{dp:eq:shuffle} it is $\binom Nn$ for the tree norms. For laws the bounds we have are~\cite[(Q.2)]{I}, which loses a factor geometric in $N$, and~\eqref{dp:4a:eq:one-law}.
%% ---- end sec-07-a-tree.tex
%% ---- begin sec-07-b-chain.tex
\subsection{Inverse functions and the chain}\label{dp:ssec:chain}

The tree norms, the class TA and the grafting map $\mathsf Q_{r,N}$ are those of Section~\ref{dp:ssec:tree} and Section~\ref{dp:ssec:ta}. For a germ with coefficient laws $a_N$ we write $L_N=L_{a_N}$.

The algebra of the formal inverse is \cite[Th.~VI.3, Prop.~VI.8]{Ro}. The bound below is the estimate we need.

\begin{theorem}[inverse functions for TA]\label{dp:thm:ta-inverse}
Let $P,W$ be Banach spaces, and let $f$ map an open subset of $P$ into $W$, TA at $x_0$, with $\mathsf T_0=Df(x_0)$ bijective with bounded inverse. Then the local inverse germ of $f$ at $f(x_0)$, which is DP by Theorem~\ref{dp:thm:inverse}, is TA at $f(x_0)$.

Let $P,Y,Y'$ be Banach spaces, let $\Phi$ map an open subset of $P\times Y$ into $Y'$, TA at $(x_0,y_0)$, with $\Phi(x_0,y_0)=0$ and with $D\Phi(x_0,y_0)(0,\cdot)\colon Y\to Y'$ bijective with bounded inverse. Then the map $\psi$ of Corollary~\ref{dp:cor:implicit}, with $\Phi(x,y)=0$ if and only if $y=\psi(x)$ near $(x_0,y_0)$, is TA at $x_0$.
\end{theorem}

\begin{proof}
We may assume $P\ne0$. Let $f$ be TA at $x_0$ with radius $R$, bound $M>0$ and coefficient laws $a_N$, and put $M_{\mathsf T}=\norm{\mathsf T_0^{-1}}M$.

\emph{Step 1 (the recursion).} By Theorem~\ref{dp:thm:differentiable}(1), $a_1=\mathsf T_0$. The germ $\mathsf T_0^{-1}\circ f$ is DP at $x_0$ with values in $P$, and its law of degree $1$ is $\id_P$. Apply \cite[Lemma~Q.14]{I} to it, and write $\mathsf R_m=-\mathsf T_0^{-1}\circ a_m$ for the laws called $\mathsf P_m$ there; so the remainder is $\mathsf R=\sum_{m\ge2}\mathsf R_m$. We obtain laws $\mathsf g_N$ of degree $N$ with $\mathsf g_1=\id_P$, a series $\mathsf g=\sum_N\mathsf g_N$ on a ball around $0$, and the local inverse germ
\begin{equation}\label{dp:eq:ta-inverse-germ}
 y\longmapsto x_0+\mathsf g\bigl(\mathsf T_0^{-1}(y-f(x_0))\bigr)
\end{equation}
of $f$ at $f(x_0)$. For $N\ge2$, the law $\mathsf g_N$ is the sum over $2\le r\le N$ and over the $\alpha$ with $\abs\alpha=r$, $\deg\alpha=N$ of the coefficients $\mathsf R_r^{[\alpha]}(\mathsf g_1,\dots,\mathsf g_{N-1})$. By Steps~1 and~2 of the proof of \cite[Lemma~Q.4]{I}, with inner laws $\mathsf g_n$ and outer law $\mathsf R_r$, this reads, on powers and so on all of $\Gamma^N(P)$,
\begin{equation}\label{dp:eq:ta-inverse-recursion}
 L_{\mathsf g_N}=\sum_{r=2}^NL_{\mathsf R_r}\circ\mathsf Q_{r,N},\qquad
 \mathsf Q_{r,N}=\sum_{\vec n\in\mathcal C(N,r)}L_{\mathsf g_{n_1}}\otimes\dots\otimes L_{\mathsf g_{n_r}}.
\end{equation}
As $r\ge2$, all parts $n_i$ are below $N$.

\emph{Step 2 (the tree bound).} Since $L_{\mathsf R_r}=-\mathsf T_0^{-1}L_r$, we have $\norm{L_{\mathsf R_r}s}\le M_{\mathsf T}R^{-r}\mathfrak t_r(s)$ for $s\in\Gamma^r(P)$. Define $\mathsf y_1=1$ and, for $N\ge2$,
\begin{equation}\label{dp:eq:ta-majorant}
 \mathsf y_N=\sum_{r=2}^NM_{\mathsf T}R^{-r}\sum_{\vec n\in\mathcal C(N,r)}\prod_{i=1}^r\mathsf y_{n_i}.
\end{equation}
We show by induction that $\norm{L_{\mathsf g_N}t}\le\mathsf y_N\mathfrak t_N(t)$ for $t\in\Gamma^N(P)$. For $N=1$ this is an equality, as $\mathfrak t_1=\norm\cdot$. Let $N\ge2$. Lemma~\ref{dp:lem:grafting}(2), with the maps $L_{\mathsf g_n}$ and constants $\mathsf y_n$ for $n<N$ and zero maps for $n\ge N$, which do not occur as $r\ge2$, gives
\[
 \mathfrak t_r(\mathsf Q_{r,N}t)\le\Bigl(\sum_{\vec n\in\mathcal C(N,r)}\prod_i\mathsf y_{n_i}\Bigr)\mathfrak t_N(t).
\]
With~\eqref{dp:eq:ta-inverse-recursion} and~\eqref{dp:eq:ta-majorant} this is the claim.

\emph{Step 3 (a positive radius).} The coefficient of $\theta^N$ in $(\sum_{1\le i<N}\mathsf y_i\theta^i)^r$ is $\sum_{\vec n\in\mathcal C(N,r)}\prod_i\mathsf y_{n_i}$. So~\eqref{dp:eq:ta-majorant} is the sequence of Step~3 of the proof of \cite[Lemma~Q.14]{I}, with $M_{\mathsf T}$ for $\mathfrak b_{\mathsf T}$ and $1/R$ for $\varkappa$. Step~4 of that proof solves a quadratic equation and holds for all positive values of these two constants. Hence there is $R_0>0$ with $\mathsf Y(\theta)=\sum_N\mathsf y_N\theta^N<\infty$ for $0\le\theta<R_0$.

\emph{Step 4 (the inverse germ).} Take $0<R_1<R_0$ so small that~\eqref{dp:eq:ta-inverse-germ} is the inverse germ for $\norm{\mathsf T_0^{-1}(y-f(x_0))}<R_1$. Its laws are $x_0$ in degree $0$ and $\mathsf g_N\circ\mathsf T_0^{-1}$ in degree $N\ge1$, with linearization $L_{\mathsf g_N}\circ(\mathsf T_0^{-1})^{\otimes N}$. By Lemma~\ref{dp:lem:tree}(2) and Step~2,
\[
 \norm{L_{\mathsf g_N}((\mathsf T_0^{-1})^{\otimes N}t)}\le\mathsf y_N\norm{\mathsf T_0^{-1}}^N\mathfrak t_N(t)\le\mathsf Y(R_1)\bigl(R_1/\norm{\mathsf T_0^{-1}}\bigr)^{-N}\mathfrak t_N(t)\qquad(t\in\Gamma^N(W)).
\]
So the inverse germ is TA at $f(x_0)$ with radius $R_1/\norm{\mathsf T_0^{-1}}$ and bound $\max(\norm{x_0},\mathsf Y(R_1))$.

\emph{Step 5 (implicit functions).} The map $\Theta(x,y)=(x,\Phi(x,y))$ is TA at $(x_0,y_0)$, by the pairing in Theorem~\ref{dp:thm:ta-composition}. Its derivative there is bijective with bounded inverse, as in the proof of Corollary~\ref{dp:cor:implicit}. By the first part, its local inverse germ at $(x_0,0)$ is TA. This inverse preserves the first coordinate. Let $\tilde\psi(x)$ be the second coordinate of its value at $(x,0)$. Then $\tilde\psi$ is TA at $x_0$, by precomposition with the bounded linear map $x\mapsto(x,0)$ and postcomposition with a projection (Theorem~\ref{dp:thm:ta-composition}). Also $\Phi(x,\tilde\psi(x))=0$, and $\tilde\psi$ is continuous with $\tilde\psi(x_0)=y_0$. By the uniqueness in Corollary~\ref{dp:cor:implicit}, $\tilde\psi=\psi$ near $x_0$.
\end{proof}

\begin{remark}[depth]\label{dp:rem:ta-depth}
The proof does not bound the depth of the trees. It bounds $L_{\mathsf g_N}$ by $\mathfrak t_N$, and not by $\mathfrak t^{(\ell)}_N$ for one $\ell$. So Theorem~\ref{dp:thm:ta-inverse} does not give inverse functions for $\mathrm{AA}^\circ$. See Section~\ref{dp:sec:questions}.
\end{remark}

\begin{corollary}[tree analytic bundles]\label{dp:cor:ta-bundles}
Let all models and fibers be Banach spaces. Then \cite[Theorem~2.1]{I}, displayed in Section~\ref{dp:sec:conv}, holds with $C^m$ replaced by TA on open sets, for manifolds and vector bundles whose transition maps are TA on open sets. The replacement is also made in the definitions that the theorem uses: morphisms are bundle maps over the identity whose local operator-valued expressions are TA on open sets, and a functor is TA if all its maps on products of hom spaces are jointly TA on open sets. The alternating bifunctor $(E,F)\mapsto\Alt^k(E;F)$ and the functor $E\mapsto\hLambda^k_\pi E$, on Banach spaces, satisfy its hypothesis. So, in every characteristic, the alternating bundle of a TA vector bundle is a TA vector bundle, with TA induced morphisms.
\end{corollary}

\begin{proof}
\emph{Step 1 (the calculus).} By Theorem~\ref{dp:thm:ta-composition}, maps that are TA on open sets compose and pair, and they restrict to open subsets. Bounded multilinear maps and constants are analytic at every point, so they are TA. In particular evaluation and composition of operators are TA. Being TA at a point depends only on the map near that point.

\emph{Step 2 (the theorem).} The proof of Corollary~\ref{dp:cor:bundles} uses of its class only the properties of Step~1.

\emph{Step 3 (the functors).} By Theorem~\ref{dp:thm:calculus}(4), the maps of these functors on hom spaces are bounded ambient laws. After an isometry of the target, such a law is the diagonal of a bounded multilinear map. So it is analytic at every point after that isometry, with diagonals in the image of the target, and it is AA at every point. As AA is $\mathrm{TA}_1$, it is TA at every point.
\end{proof}

For $\vec n\in\mathcal C(N,r)$ put $\operatorname{mult}(\vec n)=N!/\prod_in_i!$, a positive integer, and write $\abs{\operatorname{mult}(\vec n)}$ for the absolute value of its image in $K$. Let $p$ be a prime with $\abs p<1$ in $K$, that is, $p=\operatorname{char}K$, or $\operatorname{char}K=0$ and $p$ is the residue characteristic. If $N=\sum_i\mathsf d_ip^i$ in base $p$, the \emph{digit root} of $N$ is the composition that lists each $p^i$ exactly $\mathsf d_i$ times, in increasing order. Let $\operatorname{ds}_p(M)$ be the sum of the base-$p$ digits of $M$.

\begin{proposition}[base-$p$ roots]\label{dp:prop:base-p}
Let $P$ be a normed space and $t\in\Gamma^N(P)$.
\begin{enumerate}
\item \textup{(Block refinement.)} Let $\vec n\in\mathcal C(N,r)$, $r\ge1$. For each $i$ let $\vec m^{(i)}$ be a composition of $n_i$ with $\operatorname{mult}(\vec m^{(i)})\ne0$ in $K$, and let $\vec m$ be the concatenation of the $\vec m^{(i)}$. Then for every $\ell\ge1$
\[
 \mathfrak t^{(\ell)}_{\vec n}(t)\le\prod_{i=1}^r\frac{\operatorname{mult}(\vec m^{(i)})}{\abs{\operatorname{mult}(\vec m^{(i)})}}\ \mathfrak t^{(\ell)}_{\vec m}(t).
\]
The same holds with $\mathfrak t_{\vec n},\mathfrak t_{\vec m}$. In particular, taking $r=1$, so that $\vec n=(N)$, every composition $\vec m$ of $N$ with $\operatorname{mult}(\vec m)\ne0$ in $K$ has $\mathfrak t_N(t)\le\frac{\operatorname{mult}(\vec m)}{\abs{\operatorname{mult}(\vec m)}}\mathfrak t_{\vec m}(t)$.
\item \textup{(Digit roots.)} For a prime $p$ with $\abs p<1$ in $K$, the digit root $\vec n$ of $N$ has $\abs{\operatorname{mult}(\vec n)}=1$.
\item \textup{(Low degrees.)} Let $\operatorname{char}K=p$ and $2\le N\le p^2$. Then there is a constant $\mathsf A(N,p)$, depending only on $N$ and $p$, with
\[
 \mathfrak t^{(2)}_N\le\mathfrak t^{(\ell)}_N\le\mathfrak t_N\le\mathsf A(N,p)\,\mathfrak t^{(2)}_N\qquad(\ell\ge2).
\]
So all $\mathfrak t^{(\ell)}_N$, $\ell\ge2$, are equivalent in these degrees.
\end{enumerate}
\end{proposition}

\begin{proof}
\emph{Step 1 (item (1)).} Let $r_i$ be the number of parts of $\vec m^{(i)}$, and let $t=\sum_q\bigotimes_i\bigotimes_{i'}s_{q,i,i'}$ with $s_{q,i,i'}\in\Gamma^{m^{(i)}_{i'}}(P)$. For each $i$, the $\operatorname{mult}(\vec m^{(i)})$ shuffle permutations of the slots of block $i$ fix the symmetric tensor $t$. Summing over the product of these sets of permutations gives
\[
 \prod_i\operatorname{mult}(\vec m^{(i)})\,t=\sum_q\bigotimes_{i=1}^r\bigl(s_{q,i,1}\cdot\ldots\cdot s_{q,i,r_i}\bigr),
\]
a representation of a multiple of $t$ in $\bigotimes_i\Gamma^{n_i}(P)$. The shuffle bound of Lemma~\ref{dp:lem:tree}(5), iterated at level $\ell$, gives $\mathfrak t^{(\ell)}_{n_i}(s_{q,i,1}\cdot\ldots\cdot s_{q,i,r_i})\le\operatorname{mult}(\vec m^{(i)})\prod_{i'}\mathfrak t^{(\ell)}_{m^{(i)}_{i'}}(s_{q,i,i'})$. Divide by $\prod_i\abs{\operatorname{mult}(\vec m^{(i)})}$ and take the infimum over representations. For $\mathfrak t$, take $\ell\ge N$.

\emph{Step 2 (item (2)).} By Legendre's formula the exponent of $p$ in $M!$ is $(M-\operatorname{ds}_p(M))/(p-1)$. So the exponent of $p$ in $N!$ is $\sum_i\mathsf d_i(p^i-1)/(p-1)$, which is the exponent of $p$ in $\prod_i(p^i!)^{\mathsf d_i}$. Hence $\operatorname{mult}(\vec n)$ is an integer prime to $p$. Take integers $\mathsf a,\mathsf b$ with $\mathsf ap+\mathsf b\operatorname{mult}(\vec n)=1$. As $\abs p<1$ and integers have absolute value at most $1$, the ultrametric inequality forces $\abs{\operatorname{mult}(\vec n)}=1$.

\emph{Step 3 (item (3)).} The first two inequalities are Lemma~\ref{dp:lem:tree}(1). Let $\vec n$ be a composition of $N$ with at least two parts. Each $n_i$ is below $p^2$, so its digit root $\vec m^{(i)}$ has parts in $\{1,p\}$. Let $\vec m$ be the concatenation, a composition of $N$ with at least two parts. By items (1) and (2), $\mathfrak t_{\vec n}\le\prod_i\operatorname{mult}(\vec m^{(i)})\,\mathfrak t_{\vec m}$. By Lemma~\ref{dp:lem:tree}(4), $\mathfrak t_p\le\mathsf c_p\pi$, where we may take $\mathsf c_p\ge1$, and $\mathfrak t_1=\norm\cdot$. So $\mathfrak t_{\vec m}\le\mathsf c_p^{N/p}\pi_{\vec m}$. Moreover $\pi_{\vec m}\le\mathfrak t^{(2)}_N$, by Lemma~\ref{dp:lem:tree}(3). By~\eqref{dp:eq:tree-fixed}, $\mathfrak t_N$ is the maximum of the $\mathfrak t_{\vec n}$ over the compositions with at least two parts, so item (3) follows, with $\mathsf A(N,p)$ the maximum over $\vec n$ of $\prod_i\operatorname{mult}(\vec m^{(i)})\,\mathsf c_p^{N/p}$.
\end{proof}

Item~(3) includes $N=p^2$. The argument stops at $N=p^2+1$. Among the compositions of $p^2+1$ with at least two parts, $(1,p^2)$ and $(p^2,1)$ have the part $p^2$, and in characteristic $p$ this part cannot be refined: for a composition $\vec m$ of $p^2$ with at least two parts, the exponent of $p$ in $\operatorname{mult}(\vec m)$ is $(\sum_j\operatorname{ds}_p(m_j)-1)/(p-1)>0$, so $\operatorname{mult}(\vec m)=0$ in $K$. The first degree this leaves undecided is $p^2+1$: whether $\mathfrak t_{p^2+1}$ is bounded by a constant times $\mathfrak t^{(2)}_{p^2+1}$; see Section~\ref{dp:sec:questions}. For fixed $p$, item~(3) covers finitely many degrees, so it does not decide whether $\mathrm{AA}^\circ=\mathrm{TA}$.

In the rest of this subsection, let $\operatorname{char}K=p$. Let $\mathcal E$ and $\Gamma^p(\mathcal E)$, with the norm $\norm\cdot_\Gamma$, be as in Proposition~\ref{dp:prop:degree-p}, and let $\hGamma^p_\Gamma(\mathcal E)$ be the completion of this normed space, which is the space $\hGamma^p(\mathcal E)$ of Section~\ref{dp:sec:laws}.

\begin{definition}[the class $\cC_p$]\label{dp:def:cp}
$\cC_p$ is the class of germs $(P,U,f,x_0)$, with $P,W$ normed, $U\subseteq P$ open, $f\colon U\to W$, such that $f$ is DP at $x_0$ and its coefficient laws of degree $n\le p$ are ambient of degree $n$.
\end{definition}

\begin{theorem}[composition and target descent do not reach DP]\label{dp:thm:cp}
Let $K$ be of characteristic $p$, complete or not.
\begin{enumerate}
\item $\mathrm{An}\subseteq\mathrm{AA}\subseteq\cC_p\subseteq\mathrm{DP}$.
\item $\cC_p$ is closed under composition, pairing, bounded linear postcomposition, restriction to open subsets and to linear subspaces, and target descent along linear isometries $j$ for which the laws of $j\circ f$ lie in $j(W)$.
\item The map $\gamma^\sharp_p\colon K\oplus\mathcal E\to\Gamma^p(\mathcal E)$, $\gamma^\sharp_p(\lambda,h)=h^{\otimes p}$, is DP at every point, and its germ at no point is in $\cC_p$, also as a map into $\hGamma^p_\Gamma(\mathcal E)$.
\end{enumerate}
So the smallest class of germs that contains the analytic germs and is closed under composition, pairing and target descent along linear isometries with closed range is contained in $\cC_p$ and is not DP. The same holds among germs with Banach values.
\end{theorem}

\begin{proof}
\emph{Step 1 (item (1)).} By \cite[Proposition~14.9(1)]{I}, AA germs are DP with ambient laws, and analytic germs are AA.

\emph{Step 2 (composition).} Let $f$ be in $\cC_p$ at $x_0$ with laws $a_n$, and $g$ in $\cC_p$ at $f(x_0)$ with laws $b_r$. By \cite[Lemma~Q.3]{I}, $g\circ f$ is DP at $x_0$. Let $1\le N\le p$. By \cite[Lemma~Q.4]{I}, its law of degree $N$ has linearization $\sum_{r=1}^NL_{b_r}\circ\mathsf Q_{r,N}$, with $\mathsf Q_{r,N}=\sum_{\vec n\in\mathcal C(N,r)}\bigotimes_iL_{a_{n_i}}$. The term $r=1$ is $L_{b_1}\circ L_{a_N}$ with $b_1$ bounded linear, so it is bounded for $\pi$ by $\norm{b_1}\operatorname{amb}(a_N)$. For $r\ge2$ all parts of $\vec n$ are below $p$, so $\pi_{\vec n}\le\prod_in_i!\,\pi$ on $\Gamma^N(P)$, by Lemma~\ref{dp:lem:kappa-bounds}. Hence $L_{b_r}\circ\mathsf Q_{r,N}$ is bounded for $\pi$ by $\operatorname{amb}(b_r)\sum_{\vec n}\prod_in_i!\operatorname{amb}(a_{n_i})$. This is finite, as $r\le p$ and all $n_i\le p$. So the law of degree $N$ of $g\circ f$ is ambient. The law of degree $0$ is a constant, so it is ambient.

\emph{Step 3 (the other operations).} For a pair $(a,a')$ of laws, $\operatorname{amb}((a,a'))$ is the maximum of $\operatorname{amb}(a)$ and $\operatorname{amb}(a')$. For bounded linear $\Upsilon$, $\operatorname{amb}(\Upsilon a)\le\norm\Upsilon\operatorname{amb}(a)$. Restriction: on a linear subspace $P'\subseteq P$, the projective seminorm of $P'^{\otimes n}$ dominates that of $P^{\otimes n}$. Target descent: by Theorem~\ref{dp:thm:target-descent}, $f$ is DP with laws $j^{-1}\tilde a_n$, for $\tilde a_n$ the laws of $j\circ f$, and $\operatorname{amb}(j^{-1}\tilde a_n)=\operatorname{amb}(\tilde a_n)$.

\emph{Step 4 (item (3)).} $\gamma^\sharp_p$ is a bounded law of degree $p$ and norm at most $1$, so it is DP at every point, by Theorem~\ref{dp:thm:calculus}(1). At every point its law of degree $p$ is $\gamma^\sharp_p$ itself, the top-degree part of the polynomial. By Proposition~\ref{dp:prop:degree-p}(2) with $n=p$, it is not ambient of degree $p$, also into the completion.

\emph{Step 5 (the last sentence).} The class $\cC_p$ contains the analytic germs and has the operations of item (2). Target descent along an isometry with closed range is an instance of item (2), by Theorem~\ref{dp:thm:target-descent}(2). By item (3), with values in the Banach space $\hGamma^p_\Gamma(\mathcal E)$ for the last claim, $\cC_p$ is not DP.
\end{proof}

\begin{theorem}[inverse functions for $\cC_p$]\label{dp:thm:cp-inverse}
Let $P,W$ be Banach spaces and $f$ a germ from $P$ to $W$ at $x_0$ with $Df(x_0)$ bijective with bounded inverse. If $f$ is in $\cC_p$ at $x_0$, then its local inverse germ is in $\cC_p$ at $f(x_0)$. If $f$ is in $\cC_p$ at every point of a neighborhood of $x_0$, then its local inverse germ is in $\cC_p$ at every point of a neighborhood of $f(x_0)$.
\end{theorem}

\begin{proof}
\emph{Step 1 (one point).} Use the notation of Step~1 of the proof of Theorem~\ref{dp:thm:ta-inverse}, a step that uses only that $f$ is DP: the inverse germ is DP with laws $\mathsf g_N\circ\mathsf T_0^{-1}$, and~\eqref{dp:eq:ta-inverse-recursion} holds. We show by induction that $\operatorname{amb}(\mathsf g_N)<\infty$ for $N\le p$. For $N=1$ this is clear. Let $2\le N\le p$ and $r\ge2$. All parts of $\vec n\in\mathcal C(N,r)$ are below $p$, so $\pi_{\vec n}\le\prod_in_i!\,\pi$ by Lemma~\ref{dp:lem:kappa-bounds}, and $\pi(\mathsf Q_{r,N}t)\le\sum_{\vec n}\prod_in_i!\operatorname{amb}(\mathsf g_{n_i})\pi(t)$. As $r\le p$, $\operatorname{amb}(\mathsf R_r)\le\norm{\mathsf T_0^{-1}}\operatorname{amb}(a_r)<\infty$. So $\operatorname{amb}(\mathsf g_N)<\infty$, and $\operatorname{amb}(\mathsf g_N\circ\mathsf T_0^{-1})\le\operatorname{amb}(\mathsf g_N)\norm{\mathsf T_0^{-1}}^N$. The law of degree $0$ of the inverse germ is the constant $x_0$, so it is ambient.

\emph{Step 2 (nearby points).} Let $f$ be in $\cC_p$ on a neighborhood $U'$ of $x_0$. By Theorem~\ref{dp:thm:differentiable}(4), $f$ is $C^\infty$ on $U'$, and the operators bijective with bounded inverse are open in $\cL(P,W)$. So we may take $Df(x)$ bijective with bounded inverse for $x\in U'$. By Theorem~\ref{dp:thm:inverse}, after shrinking, there are open sets $U_0\subseteq U'$ and $U_1$ with $f\colon U_0\to U_1$ a bijection whose inverse $\mathsf h$ is DP on $U_1$. Let $y\in U_1$ and $x=\mathsf h(y)$. Step~1 at $x$ gives a local inverse germ $\mathsf h'$ of $f$ at $y$, in $\cC_p$ there. It is continuous and $\mathsf h'(y)=x\in U_0$, so $\mathsf h'$ has values in $U_0$ near $y$. As $f$ is injective on $U_0$, $\mathsf h'=\mathsf h$ near $y$.
\end{proof}

Two more classes sit between TA and $\cC_p$. Let $\cC_p^{\rm loc}$ be the class of germs that are in $\cC_p$ at every point of a neighborhood of the base point. Let $\cC_p^{\rm u}$ be the class of DP germs at $x_0$ such that for some radius $R$ and bound $M$, with $L_N$ the linearizations of the coefficient laws,
\begin{equation}\label{dp:eq:cpu}
 \norm{L_N(h^{\otimes(N-n)}\cdot t)}\le MR^{-N}\norm h^{N-n}\pi(t)\qquad(1\le n\le\min(N,p),\ h\in P,\ t\in\Gamma^n(P)).
\end{equation}
The class $\cC_p^{\rm loc}$ is closed under composition, as a DP germ is continuous at its base point. By Theorem~\ref{dp:thm:cp}(2) at each nearby point, and Theorem~\ref{dp:thm:target-descent}(2), it is closed under pairing, bounded linear postcomposition, restriction, and target descent along linear isometries with closed range.

\begin{theorem}[the chain]\label{dp:thm:chain}
For every $K$ of characteristic $p$, complete or not, among germs with Banach values,
\[
 \mathrm{AA}\subsetneq\mathrm{AA}^\circ\subseteq\mathrm{TA}\subseteq\cC_p^{\rm u}\subseteq\cC_p^{\rm loc}\subsetneq\cC_p\subsetneq\mathrm{DP}.
\]
\end{theorem}

\begin{proof}
\emph{Step 1 (AA and $\mathrm{AA}^\circ$).} Let $f$ be AA with Banach values, through an isometry $j$ into a Banach space. Then $j$ has closed range and $j\circ f$ is analytic, so $f\in\mathrm{AA}^\circ$ by target descent. Now use the spaces of \cite[Theorem~14.17]{I}. Over every $K$, as stated after that theorem, $g_X$ extends to the diagonal $\hat g_X\colon\widehat W^{\rm at}\to\widehat Y^{\rm at}$ of a bounded bilinear map, and $f_X$ is AA into $\widehat W^{\rm at}$. So $\hat g_X\circ f_X$ is in $\mathrm{AA}^\circ$, and it is AA at no point.

\emph{Step 2 ($\mathrm{AA}^\circ\subseteq\mathrm{TA}\subseteq\cC_p^{\rm u}$).} The first inclusion is Theorem~\ref{dp:thm:composition-closure}. Let $f$ be TA with radius $R$ and bound $M$, and let $n\le\min(N,p)$. By Lemma~\ref{dp:lem:tree}(5) and~(4),
\[
 \norm{L_N(h^{\otimes(N-n)}\cdot t)}\le MR^{-N}\binom Nn\norm h^{N-n}\mathfrak t_n(t)\le M\max_{i\le p}\mathsf c_i\,(R/2)^{-N}\norm h^{N-n}\pi(t).
\]
So~\eqref{dp:eq:cpu} holds with radius $R/2$ and bound $M\max_{i\le p}\mathsf c_i$.

\emph{Step 3 ($\cC_p^{\rm u}\subseteq\cC_p^{\rm loc}$).} Let $f$ satisfy~\eqref{dp:eq:cpu}, with $f$ DP at $x_0$ with radius $R$; shrinking $R$ we may use one radius for both. By Theorem~\ref{dp:thm:whole-ball}, $f$ is DP on $B(x_0,R)$. By Theorem~\ref{dp:thm:differentiable}(2), for small $\norm h$ the law of degree $n$ of $f$ at $x_0+h$ is $v\mapsto\sum_{N\ge n}a_N^{[(N-n,n)]}(h,v)$, and $a_N^{[(N-n,n)]}(h,v)=L_N(h^{\otimes(N-n)}\cdot v^{\otimes n})$ by Step~1 of the proof of \cite[Lemma~Q.4]{I}. For $1\le n\le p$ and $\norm h<R$, the series $L'_n(t)=\sum_{N\ge n}L_N(h^{\otimes(N-n)}\cdot t)$ converges in the Banach target, by~\eqref{dp:eq:cpu}, and $\norm{L'_n(t)}\le MR^{-n}(1-\norm h/R)^{-1}\pi(t)$. As $L'_n$ linearizes that law, the law is ambient. The law of degree $0$ is a constant, so it is ambient. So $f$ is in $\cC_p$ at $x_0+h$ for small $\norm h$.

\emph{Step 4 ($\cC_p^{\rm loc}\ne\cC_p$).} On $K\oplus\mathcal E$ let $\mathsf G_p(\lambda,h)=\lambda h^{\otimes p}$, with values in $\hGamma^p_\Gamma(\mathcal E)$. It is a bounded law of degree $p+1$, so it is DP everywhere. At $0$ its laws of degree at most $p$ vanish, so its germ at $0$ is in $\cC_p$. At $(\lambda_0,0)$ with $\lambda_0\ne0$ we have $\mathsf G_p(\lambda_0+\lambda,h)=\lambda_0h^{\otimes p}+\lambda h^{\otimes p}$. So the law of degree $p$ there is $\lambda_0\gamma^\sharp_p$, which is not ambient, by Theorem~\ref{dp:thm:cp}(3). Such points accumulate at $0$. So the germ of $\mathsf G_p$ at $0$ is in $\cC_p$ and not in $\cC_p^{\rm loc}$.

\emph{Step 5 ($\cC_p\ne\mathrm{DP}$).} This is Theorem~\ref{dp:thm:cp}(3) with values in $\hGamma^p_\Gamma(\mathcal E)$.
\end{proof}

So, in characteristic $p$ and among germs with Banach values, the classes $\mathrm{AA}^\circ$, TA and $\cC_p$ are closed under composition and lie strictly between AA and DP, and $\mathrm{AA}^\circ\ne\cC_p$. Between Banach spaces, TA and $\cC_p$ also have the inverse function theorem. The chain does not say whether $\mathrm{AA}^\circ$, TA, $\cC_p^{\rm u}$ and $\cC_p^{\rm loc}$ are distinct. The conjecture and the questions on this are in Section~\ref{dp:sec:questions}.

\subsection{What singles out DP}\label{dp:ssec:singles}

In this subsection classes consist of germs with Banach values, and (C1), (C2), (C3) are the conditions of Section~\ref{dp:ssec:minimal}: containing the analytic germs, target descent, source descent. The germ $\mathsf f$ on $\ell^1(\N,K)$, with coefficients $\mathsf e_n$ and values in $\Kh$, is that of Example~\ref{dp:ex:l1}. Let $\mathfrak S$ be the class of germs $f$ at $x_0$ such that $f\circ Q$ is analytic at a preimage of $x_0$, for some bounded linear surjection $Q$ with a preimage constant.

\begin{theorem}[what singles out DP]\label{dp:thm:singles-out}
\leavevmode
\begin{enumerate}
\item DP is the smallest class with \textup{(C1), (C2), (C3)} (Theorem~\ref{dp:thm:minimal}).
\item \textup{(Dropping (C3).)} AA is the smallest class with \textup{(C1), (C2)}. Adding composition and pairing gives $\mathrm{AA}^\circ$. The class TA, which contains $\mathrm{AA}^\circ$, has \textup{(C1), (C2)}, composition and pairing. It also has inverse and implicit functions, re-expansion on the whole ball, derivatives and bundle lifting, and for every $K$ of characteristic $p$ it is not DP. So, in characteristic $p$, no list of properties chosen from \textup{(C1), (C2)}, composition, pairing, inverse and implicit functions, re-expansion on the whole ball, derivatives and bundle lifting forces DP. In characteristic $0$, whether $\mathrm{TA}=\mathrm{DP}$ is open.
\item \textup{(Dropping (C2).)} In every characteristic, $\mathfrak S$ is the smallest class with \textup{(C1), (C3)}. The germ of $\mathsf f$ at $0$ is AA, so DP, and it is not in $\mathfrak S$. Neither is its restriction to $\ell^1_{00}(\N,K)$. So \textup{(C2)} cannot be dropped for general norms.
\item \textup{(Ultrametric spaces.)} Between ultrametric spaces, \textup{(C1), (C3)} already give DP (Theorem~\ref{dp:thm:ultrametric-minimal}), and $\mathrm{AA}=\mathrm{TA}=\mathrm{DP}$. In characteristic $0$, whether AA, TA and DP differ on non-ultrametric sources is open.
\item \textup{(Failure of (C3).)} Let $\operatorname{char}K=p$, and read $\gamma^\sharp_p$ with values in $\hGamma^p_\Gamma(\mathcal E)$. Then $\gamma^\sharp_p$ is not in $\cC_p$, while $\gamma^\sharp_p\circ q_{K\oplus\mathcal E}$ is analytic, for the canonical cover $q_{K\oplus\mathcal E}$ of Definition~\ref{dp:def:cover}. So TA, $\mathrm{AA}^\circ$, $\cC_p^{\rm loc}$ and $\cC_p$ fail \textup{(C3)}. In general, a class with \textup{(C1), (C2)} has \textup{(C3)} only if it contains DP.
\end{enumerate}
\end{theorem}

\begin{proof}
\emph{Step 1 (item (2)).} The first sentence is Proposition~\ref{dp:prop:other-classes}. The second is the definition of $\mathrm{AA}^\circ$. TA contains $\mathrm{AA}^\circ$ by Theorem~\ref{dp:thm:composition-closure}, and it has the listed properties by Theorems~\ref{dp:thm:ta-composition}, \ref{dp:thm:ta-recenter}, \ref{dp:thm:ta-inverse} and Corollary~\ref{dp:cor:ta-bundles}. It has (C1), as analytic germs are TA. It has (C2): for a linear isometry $j$ of Banach spaces with $j\circ f$ TA, the germ $j\circ f$ is DP, so its laws lie in $j(W)$ by Theorem~\ref{dp:thm:target-descent}(2), and Theorem~\ref{dp:thm:ta-composition}(4) applies. By Theorem~\ref{dp:thm:chain}, $\mathrm{TA}\ne\mathrm{DP}$.

\emph{Step 2 (item (3), the class).} $\mathfrak S$ contains An: take $Q=\id$. Let $f\circ Q\in\mathfrak S$, with $f\circ Q\circ Q'$ analytic. Then $Q\circ Q'$ is a bounded linear surjection with a preimage constant, the product of the two constants. So $\mathfrak S$ has (C3). A class with (C1), (C3) contains every germ of $\mathfrak S$, by the definition of $\mathfrak S$.

\emph{Step 3 (item (3), the germ).} $\mathsf f$ is AA at $0$ by Example~\ref{dp:ex:l1}(2). Let $P_1$ be $\ell^1(\N,K)$ or $\ell^1_{00}(\N,K)$, and suppose that $\mathsf f|_{P_1}\circ Q$ is analytic at $\tilde x_0$, with $Q\colon P'\to P_1$ a bounded linear surjection with preimage constant $\mathsf c'$, and $Q(\tilde x_0)=0$. Let $e_i$ be the unit vectors, and choose $\tilde e_i\in P'$ with $Q(\tilde e_i)=e_i$ and $\norm{\tilde e_i}\le\mathsf c'$. On $\ell^1_{00}(\N,K)$ the map $\mathsf t(h)=\sum_ih_i\tilde e_i$ is linear, $\norm{\mathsf t(h)}\le\mathsf c'\norm h$, and $Q\circ\mathsf t=\id$. So $\mathsf f(h)=(\mathsf f|_{P_1}\circ Q)(\tilde x_0+\mathsf t(h))$ is analytic at $0$ on $\ell^1_{00}(\N,K)$, with $n$-linear coefficients $\beta_n$ into $\Kh$ and $\norm{\beta_n}\le M'R'^{-n}$. Step~3 of the proof of Example~\ref{dp:ex:l1} uses only the unit vectors, so it gives $n^n\le\abs{\mathsf e_n}\le\norm{\beta_n}$ for all $n$, a contradiction.

\emph{Step 4 (item (4)).} Theorem~\ref{dp:thm:ultrametric-minimal} gives the first claim. By \cite[Theorem~14.39(1)]{I}, $\mathrm{AA}=\mathrm{DP}$ between ultrametric spaces, and $\mathrm{AA}\subseteq\mathrm{TA}\subseteq\mathrm{DP}$.

\emph{Step 5 (item (5)).} $\gamma^\sharp_p$ is not in $\cC_p$ by Theorem~\ref{dp:thm:cp}(3). By Lemma~\ref{dp:lem:laws-cover}, the law $\gamma^\sharp_p\circ q_{K\oplus\mathcal E}$ on $\ell^1_{00}(I_{K\oplus\mathcal E})$ is the diagonal of a bounded $p$-linear map, so it is analytic. A class containing An with (C3) would contain $\gamma^\sharp_p$. By Theorem~\ref{dp:thm:chain} the four named classes lie in $\cC_p$. The last sentence is Theorem~\ref{dp:thm:minimal}.
\end{proof}

\begin{remark}[the three conditions]\label{dp:rem:conditions}
The class of constant germs has (C2), (C3) and not (C1). The class of DP germs together with all germs on infinite-dimensional sources has (C1), (C2) and not (C3): take a discontinuous germ $g$ on $K$, which is not DP, and a bounded linear surjection $Q\colon P'\to K$ with a preimage constant and $\dim P'=\infty$; then $g\circ Q$ is in the class and $g$ is not. The class $\mathfrak S$ has (C1), (C3) and not (C2), by Theorem~\ref{dp:thm:singles-out}(3) and Example~\ref{dp:ex:l1}(2); by Lemma~\ref{dp:lem:cover-flat} it is the class CA of Section~\ref{dp:sec:cover}. In characteristic $p$, AA is a named class with (C1), (C2) and not (C3). In characteristic $0$ we have no such class, as $\mathrm{AA}\ne\mathrm{DP}$ is open there.
\end{remark}

\begin{remark}[the universal law]\label{dp:rem:universal-law}
Theorem~\ref{dp:thm:singles-out} concerns the listed properties. Let $P$ be normed, let $\mathbf U$ be the completed $\ell^1$-sum of the spaces $\hGamma^n(P)$ of Section~\ref{dp:sec:laws}, $n\ge1$, with $\iota_n\colon\Gamma^n(P)\to\mathbf U$ the canonical maps, and put $\mathsf v(h)=\sum_{n\ge1}\iota_n(h^{\otimes n})$ for $\norm h<1$. Let $f$ be DP at $x_0$ with radius $R$, bound $M$ and laws $a_n$, with values in a Banach space $W$. Choose $\mathsf s\in K$ with $\abs{\mathsf s}\ge1/R$. By Theorem~\ref{dp:thm:normed-roby}, $\norm{\mathsf s^{-n}L_{a_n}t}\le M\norm t_\Gamma$. So $\iota_n(t)\mapsto\mathsf s^{-n}L_{a_n}t$ extends to a bounded linear map $\mathbf L\colon\mathbf U\to W$, and
\[
 f(x)=f(x_0)+\mathbf L\bigl(\mathsf v(\mathsf s(x-x_0))\bigr)\qquad(\norm{x-x_0}<1/\abs{\mathsf s}).
\]
So a class that contains the constants and the germs $\mathsf v$ and is closed under sums, bounded linear postcomposition and precomposition with affine maps contains every DP germ with Banach values. So these conditions, with no source descent, also reach every DP germ with Banach values.
\end{remark}

Table~\ref{dp:tab:classes} collects the comparison. Among the properties in it, only source descent and the row on all laws are not shared by TA, and source descent with (C1), (C2) forces DP. For TA, re-expansion keeps the radius $R-\varrho$ over every $K$ (Theorem~\ref{dp:thm:ta-recenter}). For DP, Theorem~\ref{dp:thm:whole-ball}(2) gives this inequality for the DP radius over densely valued $K$, when $f$ is the sum of its law series on the ball, and Example~\ref{dp:ex:radius-drop} shows that it can fail over discretely valued $K$.

\begin{table}[htbp]
\small
\caption{Classes in characteristic $p$, for $K$ complete and Banach values, unless a row says otherwise. A dash means that the entry is not considered here.}\label{dp:tab:classes}
\begin{tabular}{@{}p{0.19\linewidth}p{0.1\linewidth}p{0.085\linewidth}p{0.105\linewidth}p{0.085\linewidth}p{0.105\linewidth}p{0.11\linewidth}@{}}
 & An & AA & $\mathrm{AA}^\circ$ & TA & $\cC_p^{\rm loc}$ & DP\\
\hline
Composition, every $K$ of char.\ $p$ & yes & no & yes & yes & yes & yes\\
Target descent & no & yes & yes & yes & yes & yes\\
Source descent & no & no & no & no & no & yes\\
Inverse functions & --- & no & open & yes & yes & yes\\
Whole ball & yes & yes & yes & yes & --- & yes, radius can drop\\
All bounded laws & no & no & no & no & no & yes\\
Laws bounded for & lifts in $W$ & $\pi$ & $\mathfrak t^{(\ell)}_N$, one $\ell$ & $\mathfrak t_N$ & $\pi$, degree $\le p$ & $\norm\cdot_\Gamma$\\
Between ultrametric spaces & can differ from DP & $=\mathrm{DP}$ & $=\mathrm{DP}$ & $=\mathrm{DP}$ & $=\mathrm{DP}$ & $=\mathrm{DP}$\\
\end{tabular}
\end{table}

\emph{Sources of the entries.} Composition: \cite[Proposition~14.26, proof, Step~1]{I} for An; \cite[Theorem~14.17]{I} for AA; the definition for $\mathrm{AA}^\circ$; Theorem~\ref{dp:thm:ta-composition} for TA; Theorem~\ref{dp:thm:cp}(2) and the paragraph before Theorem~\ref{dp:thm:chain} for $\cC_p^{\rm loc}$; \cite[Lemma~Q.3]{I} for DP. Target descent: Proposition~\ref{dp:prop:other-classes} for An and AA; the definition for $\mathrm{AA}^\circ$; Theorem~\ref{dp:thm:ta-composition}(4) with Theorem~\ref{dp:thm:target-descent}(2) for TA; the paragraph before Theorem~\ref{dp:thm:chain} for $\cC_p^{\rm loc}$; Theorem~\ref{dp:thm:target-descent} for DP. Source descent: Proposition~\ref{dp:prop:other-classes}, Theorem~\ref{dp:thm:singles-out}(5), Proposition~\ref{dp:prop:source-descent}. Inverse functions: \cite[Corollary~14.19(1)]{I} for AA; Theorems~\ref{dp:thm:ta-inverse}, \ref{dp:thm:cp-inverse}, \ref{dp:thm:inverse}; for $\mathrm{AA}^\circ$ see Section~\ref{dp:sec:questions}. Whole ball, meaning membership at every point of the ball of an expansion: Step~5 of the proof of \cite[Proposition~14.12]{I} with $j=\id$, for An; \cite[Proposition~14.12(2)]{I} for AA; Theorems~\ref{dp:thm:composition-closure} and~\ref{dp:thm:ta-recenter} for $\mathrm{AA}^\circ$ and TA; Theorem~\ref{dp:thm:whole-ball} and Example~\ref{dp:ex:radius-drop} for DP. All laws: Theorem~\ref{dp:thm:cp}(3), and Theorem~\ref{dp:thm:calculus}(1) for DP. The row ``laws bounded for'' names the seminorm for which the linearizations are bounded, with geometric constants, by \cite[Proposition~14.9(1)]{I}, Theorem~\ref{dp:thm:composition-closure} and Theorem~\ref{dp:thm:normed-roby}; for An the laws are diagonals of multilinear maps into $W$. Between ultrametric spaces: \cite[Proposition~14.40(3)]{I} for An, which agrees with DP if $K$ is discretely valued or $W$ is spherically complete, by \cite[Proposition~14.40(1)]{I}; the other entries by \cite[Theorem~14.39(1)]{I} and the chain.

\subsection{Composition on all sources}\label{dp:ssec:all-sources}

In this subsection the characteristic of $K$ is arbitrary. We use the data of \cite[Definition~14.15]{I}. For a normed space $X\ne0$, $n\ge1$ and a family $y=(y_l)_{l\in I}$ of elementary tensors of $X^{\otimes n}$, given with their factors, let $E_0$ be the span of the $y_l$ and $\mathcal P(y)$ the set of linear maps $\Psi\colon E_0\to\Gamma^n(X)$ that fix $E_0\cap\Gamma^n(X)$. Then
\[
 \mathfrak l(y)=\inf_{\Psi\in\mathcal P(y)}\frac{\sum_l\pi(\Psi y_l)}{\sum_l\norm{y_l}_{\rm e}},\qquad\mathfrak l_n(X)=\sup_y\mathfrak l(y).
\]
Put $\mathfrak L_n(K)=\sup_P\mathfrak l_n(P)$, over all normed spaces $P\ne0$. As $\Gamma^1(X)=X$, the map $\Psi=\id$ gives $\mathfrak l_1(X)\le1$. If $\operatorname{char}K=0$, symmetrization gives $\Psi$ with $\pi(\Psi y_l)\le(n!/\abs{n!})\norm{y_l}_{\rm e}$, the numerator being a real number; so every $\mathfrak l_n(X)$ is finite, and only the growth in $n$ is in question. If $\operatorname{char}K=p$, then $\mathfrak L_p(K)=\infty$ by \cite[Lemma~Q.9]{I}.

Let $\mathsf A$ be a finite set, $\mathsf w=(\mathsf w_{\mathsf a})_{\mathsf a\in\mathsf A}$ positive real numbers, $m\ge1$, and $\operatorname{qm}\colon K^{\mathsf A}\to K^m$ a linear surjection. Let $X_{\mathsf A,\mathsf w,\operatorname{qm}}$ be $K^m$ with the norm
\[
 \norm h=\inf\Bigl\{\sum_{\mathsf a}\abs{\mathsf z_{\mathsf a}}\mathsf w_{\mathsf a}:\ \operatorname{qm}(\mathsf z)=h\Bigr\}.
\]
This is a quotient norm, and it is a norm since $\ker\operatorname{qm}$, the common kernel of finitely many bounded linear forms of the weighted $\ell^1$-norm on $K^{\mathsf A}$, is closed.

\begin{proposition}[composition on all sources]\label{dp:prop:all-sources}
\leavevmode
\begin{enumerate}
\item Let $P\ne0$ be normed and $\vec n\in\mathcal C(N,r)$, with $\mathfrak l_{n_i}(P)<\infty$ whenever $n_i\ge2$. Then $\mathfrak k_{\vec n}(P)\le\prod_{i:\,n_i\ge2}\abs\varpi^{-1}\mathfrak l_{n_i}(P)$. So $\sup_n\mathfrak l_n(P)^{1/n}<\infty$ implies $\mathfrak k(P)<\infty$, and if $\sup_n\mathfrak L_n(K)^{1/n}<\infty$, then composites of AA germs are AA on every normed source.
\item Let $K$ be complete and $\sup_n\mathfrak L_n(K)^{1/n}=\infty$. Then there are a Banach space $H$ with $\mathfrak k(H)=\infty$, a germ $f$ on $H$ that is AA with values in a Banach space $W$, and an analytic germ $g$ on $W$ with values in a Banach space, such that $g\circ f$ is not AA.
\item For every $n$, $\mathfrak L_n(K)$ is the supremum of $\mathfrak l_n(X)$ over the spaces $X=X_{\mathsf A,\mathsf w,\operatorname{qm}}$.
\end{enumerate}
\end{proposition}

So for $K$ complete, $\mathfrak k(P)<\infty$ holds for every normed $P$ if and only if $\sup_n\mathfrak L_n(K)^{1/n}<\infty$, and by item~(3) this condition concerns finite-dimensional spaces only.

\begin{proof}
\emph{Step 1 (item (3)).} Each $X_{\mathsf A,\mathsf w,\operatorname{qm}}$ is a nonzero normed space, so the supremum in item~(3) is at most $\mathfrak L_n(K)$. Conversely, let $y$ be a family in $P^{\otimes n}$. Let $\mathsf A\subseteq P$ be the finite set of its nonzero factors, $\mathsf w_{\mathsf a}=\norm{\mathsf a}_P$, $Z$ the span of $\mathsf A$, $m=\dim Z$, and identify $Z$ with $K^m$. Let $\operatorname{qm}(\mathsf z)=\sum_{\mathsf a}\mathsf z_{\mathsf a}\mathsf a$, and let $X$ be $Z$ with the resulting norm. The triangle inequality gives $\norm h_P\le\norm h_X$ on $Z$, and $\norm{\mathsf a}_X\le\mathsf w_{\mathsf a}=\norm{\mathsf a}_P$. So the atoms have the same norm in $X$ and in $P$, and $\norm{y_l}_{\rm e}$ is the same in both. As $\Gamma^n$ is the space of symmetric tensors (Section~\ref{dp:sec:conv}), $\Gamma^n(P)\cap Z^{\otimes n}=\Gamma^n(Z)$. Since $E_0\subseteq Z^{\otimes n}$, the set $\mathcal P(y)$ is the same in $X$ and in $P$. The identity from $X$ to $P$ is a contraction, so $\pi_P\le\pi_X$ on $Z^{\otimes n}$. Hence $\mathfrak l(y)$ computed in $P$ is at most $\mathfrak l(y)$ computed in $X$.

\emph{Step 2 (item (1), one block).} Let $t\in\Gamma^N(P)$. For $0\le i\le r$ let $\pi^{(i)}$ be the projective seminorm of $\Gamma^{n_1}_\pi(P)\otimes\dots\otimes\Gamma^{n_i}_\pi(P)\otimes P^{\otimes(n_{i+1}+\dots+n_r)}$, a space that contains $t$. Then $\pi^{(0)}(t)=\pi(t)$ and $\pi^{(r)}(t)=\pi_{\vec n}(t)$. We claim that $\pi^{(i)}(t)\le\abs\varpi^{-1}\mathfrak l_{n_i}(P)\pi^{(i-1)}(t)$ if $n_i\ge2$, and $\pi^{(i)}(t)=\pi^{(i-1)}(t)$ if $n_i=1$. The second holds as $\Gamma^1_\pi(P)=P$. Let $n_i\ge2$. Let $G_0$ be the tensor product of all the factors except block $i$, that is, $\Gamma^{n_1}_\pi(P)\otimes\dots\otimes\Gamma^{n_{i-1}}_\pi(P)\otimes P^{\otimes(n_{i+1}+\dots+n_r)}$, with its projective seminorm $\pi_{G_0}$. Let $G$ be its separated quotient, a normed space, with quotient map $\operatorname{sq}$. Moving block $i$ to the front, we read $t$ in $\Gamma^{n_i}(P)\otimes G_0$, and we put $t'=(1\otimes\operatorname{sq})t\in\Gamma^{n_i}(P)\otimes G$. Moving factors does not change projective seminorms, and the projective seminorm is associative. So $\pi(t')\le\pi^{(i-1)}(t)$, and \cite[Lemma~Q.7(2)]{I} gives $\mu(t')\le\abs\varpi^{-1}\mathfrak l_{n_i}(P)\pi^{(i-1)}(t)$. Let $t'=\sum_qs_q\otimes w'_q$ with $s_q\in\Gamma^{n_i}(P)$, $w'_q\in G$, and choose $w_q\in G_0$ with $\operatorname{sq}(w_q)=w'_q$ and $\pi_{G_0}(w_q)=\norm{w'_q}$. Then $t-\sum_qs_q\otimes w_q$ lies in $\Gamma^{n_i}(P)\otimes\ker\operatorname{sq}$, which has cost $0$. So $\pi^{(i)}(t)\le\sum_q\pi(s_q)\norm{w'_q}$, and taking the infimum, $\pi^{(i)}(t)\le\mu(t')$. This proves the claim.

\emph{Step 3 (item (1), the conclusion).} Multiplying the claims gives the bound on $\mathfrak k_{\vec n}(P)$, for $\pi(t)=0$ as well. Let $\sup_n\mathfrak l_n(P)^{1/n}<\infty$, and let $\mathsf c\ge1$ be at least this supremum. At most $N$ blocks occur and $\sum_in_i=N$, so $\mathfrak k_{\vec n}(P)\le\abs\varpi^{-N}\mathsf c^N$, and $\mathfrak k(P)\le\mathsf c/\abs\varpi$. If $\sup_n\mathfrak L_n(K)^{1/n}<\infty$, this holds for every $P\ne0$, and \cite[Theorem~14.22]{I}, (iii) implies (i), gives the composites.

\emph{Step 4 (item (2), the spaces).} For each $i\ge2$ choose $n_i$ with $\mathfrak L_{n_i}(K)>i^{n_i}$. By item~(3) there are a space $P_i=X_{\mathsf A,\mathsf w,\operatorname{qm}}$ and a family $y^{(i)}$ in $P_i^{\otimes n_i}$ with $\mathfrak l(y^{(i)})\ge i^{n_i}$. As $\mathfrak l_1\le1$, $n_i\ge2$. Let $X$ be the $\ell^1$-sum of the $P_i$, and $G$ the $\ell^1$-sum of the universal spaces $G(y^{(i)})$ of \cite[Definition~14.15]{I}. Over complete $K$ the $P_i$ and $G(y^{(i)})$ are finite-dimensional, so complete. So $X$, $G$ and $H=X\oplus G$, with the maximum norm, are Banach spaces.

Each $P_i$ sits isometrically in $X$, with a contractive coordinate projection $\operatorname{pr}_i$. As $\Gamma^{n_i}(X)\cap P_i^{\otimes n_i}=\Gamma^{n_i}(P_i)$, the space $G(y^{(i)})$ is the same computed in $P_i$ or in $X$. If $\tilde\Psi$ is in $\mathcal P(y^{(i)})$ for $X$, then $\operatorname{pr}_i^{\otimes n_i}\circ\tilde\Psi$ is in $\mathcal P(y^{(i)})$ for $P_i$, with values of no larger $\pi$. So $\mathfrak l(y^{(i)})\ge i^{n_i}$ also holds in $X$.

\emph{Step 5 (item (2), the constant).} Fix $i\ge2$ and write $n=n_i$. Steps~4, 5 and~7 of the proof of \cite[Proposition~14.18]{I} use one family at a time, and neither that $\mathfrak l_n(X)$ is infinite nor that all families have the same degree. For this family they give $\mathfrak k_{(n,1)}(H)\ge\mathfrak l(y^{(i)})/(n+1)\ge i^n/(n+1)$. As $n\ge2$, we have $n/(n+1)\ge2/3$ and $(n+1)^{1/(n+1)}\le2$, so
\[
 \mathfrak k(H)\ge\bigl(i^n/(n+1)\bigr)^{1/(n+1)}\ge i^{2/3}/2.
\]
This holds for all $i\ge2$, so $\mathfrak k(H)=\infty$. As $K$ is complete, \cite[Theorem~14.22]{I} says that \textup{(iii)}, \textup{(iii$'$)} and \textup{(ii)} are equivalent for $H$, so \textup{(ii)} fails, and this is the assertion.
\end{proof}
%% ---- end sec-07-b-chain.tex
%% ---- begin sec-08-forms.tex
\section{Divided-power forms}\label{dp:sec:forms}

Throughout this section $P,P'$ are normed model spaces, $F,F_1,F_2,F_3$ are normed value spaces, and $k,l\ge0$ are integers. The characteristic of $K$ is arbitrary unless this is said. We read $\Mult^0(P;F)=\Alt^0(P;F)=F$. So $\cL(P,\Alt^0(P;F))=\cL(P,F)=\Alt^1(P;F)$.

\subsection{The de Rham complex}\label{dp:ssec:de-rham}

Let $\cM$ be a DP manifold modeled on $P$. A section of the DP bundle $\Alt^k(T\cM;F)$ of Proposition~\ref{dp:prop:categories}(3) is \emph{DP} if its representative in every chart of $\cM$ is DP. A \emph{DP $F$-valued $k$-form} on $\cM$ is a DP section of this bundle; for $k=0$ it is a DP map $\cM\to F$. Write $\Omega^k_{\mathrm{DP}}(\cM;F)$ for the space of these forms. A representative in a chart is a map $\eta\colon U\to\Alt^k(P;F)$ on an open set $U\subseteq P$.

Let $U\subseteq P$ and $U'\subseteq P'$ be open, $f\colon U\to U'$ be DP on $U$, and $\eta\colon U'\to\Alt^k(P';F)$. The \emph{pullback} is
\begin{equation}\label{dp:f:pullback}
 (f^*\eta)(x)=A^k_{P,P';F}\bigl(Df(x)\bigr)\bigl(\eta(f(x))\bigr)=\eta(f(x))\circ\bigl(Df(x),\dots,Df(x)\bigr),
\end{equation}
which for $k=0$ reads $f^*\eta=\eta\circ f$. Let $\delta_k\colon\cL(P,\Alt^k(P;F))\to\Mult^{k+1}(P;F)$ be the linear map, of norm at most $k+1$, given by
\[
 \delta_k(\mathsf T)(v_0,\dots,v_k)=\sum_{i=0}^k(-1)^i\mathsf T(v_i)(v_0,\dots,\widehat{v_i},\dots,v_k).
\]
For a differentiable map $\eta\colon U\to\Alt^k(P;F)$ the \emph{exterior derivative} is $d\eta=\delta_k\circ D\eta$; this is the formula~\cite[(10.3)]{I}. For a bounded bilinear pairing $F_1\times F_2\to F_3$, the \emph{wedge product} of an $F_1$-valued $k$-form and an $F_2$-valued $l$-form is the shuffle formula~\cite[(10.2)]{I}, with values multiplied in the pairing.

\begin{theorem}[DP de Rham complex]\label{dp:thm:de-rham}
In every characteristic:
\begin{enumerate}
\item \textup{(Pullback.)} If $f$ is DP on $U$ and $\eta$ is DP on $U'$, then $f^*\eta$ is DP on $U$, and $(f_1\circ f)^*=f^*f_1^*$ for a further DP map $f_1$ from $U'$ to an open subset of a normed space. Hence a section of $\Alt^k(T\cM;F)$ is DP if and only if its representative in one chart around each point is DP. Representatives in charts $\chi,\chi'$ are related by $\eta_\chi=(\chi'\circ\chi^{-1})^*\eta_{\chi'}$, and DP maps of manifolds pull back DP forms.
\item \textup{(Exterior derivative.)} If $\eta\colon U\to\Alt^k(P;F)$ is DP on $U$, then $d\eta$ has values in $\Alt^{k+1}(P;F)$ and is DP on $U$.
\item \textup{(Naturality.)} For $f,\eta$ as in item~\textup{(1)}, $d(f^*\eta)=f^*(d\eta)$. Hence $d\colon\Omega^k_{\mathrm{DP}}(\cM;F)\to\Omega^{k+1}_{\mathrm{DP}}(\cM;F)$ is well defined, independently of charts.
\item $d\circ d=0$.
\item \textup{(Products.)} The wedge product of DP forms is DP and commutes with pullback, and $d(\eta\wedge\xi)=d\eta\wedge\xi+(-1)^k\eta\wedge d\xi$ for a $k$-form $\eta$. For scalar forms it is associative and graded commutative.
\end{enumerate}
So $\Omega^\bullet_{\mathrm{DP}}(-;K)$ is a sheaf of differential graded algebras on DP manifolds, contravariant for DP maps.
\end{theorem}

\begin{proof}
Steps~2--5 repeat the algebra of Steps~2--6 of the proof of~\cite[Theorem~10.7]{I}. It is written there for scalar values and uses only linearity in the values, bilinearity of the pairing, the chain and product rules, and the symmetry of second derivatives. These hold for DP maps by Theorem~\ref{dp:thm:differentiable}(1), (3) and~(4), in every characteristic.

\emph{Step 1 (pullback).} For $k=0$ this is Theorem~\ref{dp:thm:calculus}(2). Let $k\ge1$. By Theorem~\ref{dp:thm:differentiable}(3) and Theorem~\ref{dp:thm:calculus}(2), the map $x\mapsto(Df(x),\eta(f(x)))$ is DP into $\cL(P,P')\times\Alt^k(P';F)$; a map into a product is DP when its components are, by Step~1 of the proof of Corollary~\ref{dp:cor:bundles}. The map $A=A^k_{P,P';F}$ is DP on $\cL(P,P')$ by Theorem~\ref{dp:thm:calculus}(4) and~(1), and evaluation of operators is bounded bilinear. Composing, $f^*\eta$ is DP. The composition rule is the chain rule. The bundle $\Alt^k(T\cM;F)$ has transitions $A(D(\chi'\circ\chi^{-1}))$, which gives the relation between representatives. So a representative that is DP on a chart domain gives representatives that are DP on the overlaps with every other chart, and being DP at a point depends only on the map near that point. In charts, the pullback along a DP map of manifolds is~\eqref{dp:f:pullback}, so it preserves DP forms.

\emph{Step 2 ($d\eta$).} By Theorem~\ref{dp:thm:differentiable}(3), $D\eta\colon U\to\cL(P,\Alt^k(P;F))$ is DP. The map $\delta_k$ has alternating values. Let $v_i=v_q$ with $i<q$. The summands with index other than $i,q$ contain this vector twice and vanish. The two others differ by moving it across $q-i-1$ entries, so they cancel, with $(-1)^i$ and $(-1)^q(-1)^{q-i-1}=(-1)^{i+1}$; here an alternating map changes sign under a transposition, also in characteristic $2$. Hence $d\eta=\delta_k\circ D\eta$ is a DP map followed by a bounded linear map into $\Alt^{k+1}(P;F)$.

\emph{Step 3 (naturality).} Let $u=Df(x)$ and $u_2=D^2f(x)$. By Step~1 and Theorem~\ref{dp:thm:differentiable}(1), $f^*\eta$ is differentiable into $\Alt^k(P;F)$, and the inclusion into $\Mult^k(P;F)$ is an isometry. In $\Mult^k(P;F)$, the right side of~\eqref{dp:f:pullback} is the bounded $(k+1)$\nobreakdash-linear map $(m,t_1,\dots,t_k)\mapsto m\circ(t_1,\dots,t_k)$ at $(\eta\circ f,Df,\dots,Df)$. So the product rule gives
\begin{multline*}
 D(f^*\eta)(x)(v_0)(v_1,\dots,v_k)=D\eta(f(x))(uv_0)(uv_1,\dots,uv_k)\\
 +\sum_{q=1}^k\eta(f(x))(uv_1,\dots,u_2(v_0,v_q),\dots,uv_k).
\end{multline*}
Apply $\delta_k$. The terms that differentiate $\eta$ give $f^*(d\eta)(x)(v_0,\dots,v_k)$. The remaining terms give the sum over $i\ne q$ of $(-1)^iX_{iq}$, where $X_{iq}$ is $\eta(f(x))$ on the list $(uv_c)_{c\ne i}$ with $uv_q$ replaced by $u_2(v_i,v_q)$. For $i<q$, symmetry of $u_2$ and alternation of $\eta(f(x))$ give $X_{iq}=(-1)^{q-1-i}X_{qi}$, so $(-1)^iX_{iq}+(-1)^qX_{qi}=0$. Applied to transitions, this shows that the local $d\eta_\chi$ satisfy the relation of item~(1), so they define a global DP section.

\emph{Step 4 ($d^2=0$).} Since $\delta_k$ is bounded linear, $D(d\eta)=\delta_k\circ D^2\eta$. So $d(d\eta)(x)(v_0,\dots,v_{k+1})$ is the sum over $i\ne q$ of $\epsilon_{iq}D^2\eta(x)(v_i,v_q)$ on the other arguments in order, with $\epsilon_{iq}=(-1)^{i+q-1}$ for $i<q$ and $\epsilon_{iq}=(-1)^{i+q}$ for $i>q$. The terms $(i,q)$ and $(q,i)$ have opposite signs and, by symmetry of $D^2\eta(x)$, equal values.

\emph{Step 5 (products).} By Step~2 of the proof of~\cite[Theorem~10.7]{I}, the shuffle formula is a bounded bilinear map $\Alt^k(P;F_1)\times\Alt^l(P;F_2)\to\Alt^{k+l}(P;F_3)$ that commutes with simultaneous precomposition, and for scalar values it is associative and graded commutative. So the wedge product of DP forms is DP, by Step~1 of the proof of Corollary~\ref{dp:cor:bundles}, and commutes with~\eqref{dp:f:pullback}. The Leibniz rule is Step~6 of the proof of~\cite[Theorem~10.7]{I}: a derivative index placed in the second block moves past $k$ arguments and contributes $(-1)^k$.

\emph{Step 6 (the sheaf).} Being DP at a point depends only on the map near that point, by Step~1 of the proof of Corollary~\ref{dp:cor:bundles}, and sums of DP maps are DP, since addition is bounded linear. So DP forms restrict to open subsets and glue along open covers. With items~(1)--(5) this gives the last sentence of the theorem.
\end{proof}

Now we compare with the ambient forms of~\cite[Section~10.3]{I}. Let $j\colon\Alt^k(P;K)\to\Mult^k(P;K)$ be the inclusion.

\begin{proposition}[comparison with ambient forms]\label{dp:prop:forms-compare}
Let $K$ be complete, and let $\cM$ be a Hausdorff analytic manifold as in~\cite[Section~10.3]{I}. Its transitions are analytic, hence DP by Theorem~\ref{dp:thm:calculus}(3), so $\cM$ is a DP manifold, whose maximal DP atlas contains the analytic atlas.
\begin{enumerate}
\item Every ambient analytic form $\eta\in\Omega^k_{\mathrm{amb}}(\cM)$ is a DP form. So $\Omega^\bullet_{\mathrm{amb}}(\cM)\subseteq\Omega^\bullet_{\mathrm{DP}}(\cM;K)$, and $d$, $\wedge$ and pullback along analytic maps of~\cite[Theorem~10.7]{I} are the restrictions of those of Theorem~\ref{dp:thm:de-rham}.
\item Let $f\colon U\to U'$ and $\omega\colon U'\to\Alt^k(P';K)$ be analytic on open sets. Then $f^*\omega$ is DP on $U$. Suppose moreover that $K$ has characteristic $p>0$ and is not spherically complete, and that $k\ge p$. Then there are Banach spaces $N,Y$, a polynomial analytic map $f_0\colon N\to Y$ and an analytic map $\omega_0\colon Y\to\Alt^k(Y;K)$ such that the DP form $f_0^*\omega_0$ is analytic at no point as an $\Alt^k(N;K)$-valued map.
\item If the model $P$ has finite coordinates, or if $\operatorname{char}K=0$ and $P$ is ultrametric, then three classes of $k$-forms on $\cM$ coincide: DP forms, ambient analytic forms, and forms whose representatives in the charts of the analytic atlas are analytic into $\Alt^k(P;K)$.
\end{enumerate}
\end{proposition}

\begin{proof}
\emph{Step 1 (item (1)).} Let $\eta_\chi$ be the representative of $\eta$ in a chart $\chi$ of the analytic atlas of $\cM$. The map $j\circ\eta_\chi$ is analytic into the Banach space $\Mult^k(P;K)$. The map $j$ is an isometry, and its range is closed, since vanishing on repeated arguments is a closed condition. So the diagonals of an expansion of $j\circ\eta_\chi$ lie in the range of $j$, by~\cite[Theorem~2.3]{I}, and $\eta_\chi$ is AA at each point of its domain. By~\cite[Proposition~14.9(1)]{I} it is DP there. So $\eta$ is a DP form, by Theorem~\ref{dp:thm:de-rham}(1). As $\eta_\chi$ is differentiable into $\Alt^k(P;K)$, the chain rule gives $D(j\circ\eta_\chi)=j\circ D\eta_\chi$. So~\cite[(10.3)]{I} gives the same form in $\Mult^k$ and in $\Alt^k$. The wedge product and pullback are the same pointwise formulas.

\emph{Step 2 (item (2)).} Analytic maps are DP (Theorem~\ref{dp:thm:calculus}(3)), so Theorem~\ref{dp:thm:de-rham}(1) applies. The maps $f_0,\omega_0$ are those named $h,\omega_0$ in~\cite[Proposition~10.3(2)]{I}, whose hypotheses are the ones stated.

\emph{Step 3 (item (3)).} A form with analytic representatives is ambient analytic, by composing with $j$, and an ambient analytic form is DP by item~(1). Let $\eta$ be a DP form, with representative $\eta_\chi$ in a chart $\chi$ of the analytic atlas. If $P$ has finite coordinates, $\eta_\chi$ is analytic by Theorem~\ref{dp:thm:calculus}(3). If $\operatorname{char}K=0$ and $P$ is ultrametric, it is analytic by~\cite[Proposition~14.40(2)]{I}, with target $\Alt^k(P;K)$. So a DP form has analytic representatives, and the three classes coincide.
\end{proof}

\subsection{Weights}\label{dp:ssec:weights}

\emph{Derivatives of laws.} Let $a\in\cP_n(P;W)$ for a normed space $W$, and let $x\in P$. By Theorem~\ref{dp:thm:mixed}(2), $v\mapsto\sum_{i\le n}a^{[(n-i,i)]}(x,v)$ is the expansion of $a$ at $x$. So by Theorem~\ref{dp:thm:differentiable}(1), with $a^{(1)}$ as in~\eqref{dp:laws:hasse},
\begin{equation}\label{dp:f:law-derivative}
 Da(x)=a^{(1)}(x)\quad(n\ge1),\qquad Da=0\quad(n=0).
\end{equation}
By Lemma~\ref{dp:lem:hasse-law}(1), $x\mapsto Da(x)$ is a bounded polynomial law of degree $n-1$ into $\cL(P,W)$, of norm at most $C_2^n\norm a$ for $n\ge1$. By the case of one derivative in Lemma~\ref{dp:lem:hasse-law}(5), and trivially for $n=0$,
\begin{equation}\label{dp:f:law-euler}
 Da(x)(x)=n\,a(x).
\end{equation}

Let $\eta_n\colon P\to\Alt^k(P;F)$ be a bounded polynomial law of degree $n$. It is a form on $P$, DP everywhere by Theorem~\ref{dp:thm:calculus}(1). Its \emph{total weight} is $n+k$. Its \emph{contraction with the Euler field} $x\mapsto x$ is
\begin{equation}\label{dp:f:contraction}
 (\iota\eta_n)(x)(v_1,\dots,v_{k-1})=\eta_n(x)(x,v_1,\dots,v_{k-1})\quad(k\ge1),\qquad\iota\eta_n=0\quad(k=0).
\end{equation}
For $k\ge1$, $\iota\eta_n\in\cP_{n+1}(P;\Alt^{k-1}(P;F))$ and $\norm{\iota\eta_n}\le\norm{\eta_n}$. Indeed, on $x=\sum_cs_ch_c$ the map~\eqref{dp:f:contraction} is the product of a homogeneous polynomial of degree $n$ in $s$ with a linear one, paired by the bilinear map $(m,x)\mapsto m(x,\cdot)$ on $\Alt^k(P;F)\times P$ of norm at most $1$. Also, by Lemma~\ref{dp:lem:elementary}(2) and Theorem~\ref{dp:thm:de-rham}(2), $d\eta_n=\delta_k\circ D\eta_n$ is a law of degree $n-1$ into $\Alt^{k+1}(P;F)$ for $n\ge1$, of norm at most $(k+1)C_2^n\norm{\eta_n}$, and $d\eta_0=0$. So $\iota d\eta_n$ is defined.

\begin{proposition}[Euler--Cartan identity]\label{dp:prop:euler}
For a law $\eta_n\in\cP_n(P;\Alt^k(P;F))$, we have $(d\iota+\iota d)\eta_n=(n+k)\,\eta_n$, in every characteristic.
\end{proposition}

\begin{proof}
Write $\mathbf v=(v_0,\dots,v_{k-1})$, and $\mathbf v_{\hat\imath}$ for this list without $v_i$.

\emph{Step 1 ($d\iota$, for $k\ge1$).} By the product rule, $D(\iota\eta_n)(x)(v_0)=D\eta_n(x)(v_0)(x,\cdot)+\eta_n(x)(v_0,\cdot)$. Apply $\delta_{k-1}$. Moving $v_i$ to its place costs $(-1)^i$, so the second terms give $k\,\eta_n(x)(\mathbf v)$. Hence
\[
 d\iota\eta_n(x)(\mathbf v)=\sum_{i=0}^{k-1}(-1)^iD\eta_n(x)(v_i)(x,\mathbf v_{\hat\imath})+k\,\eta_n(x)(\mathbf v).
\]

\emph{Step 2 ($\iota d$).} By definition,
\[
 \iota d\eta_n(x)(\mathbf v)=D\eta_n(x)(x)(\mathbf v)+\sum_{i=0}^{k-1}(-1)^{i+1}D\eta_n(x)(v_i)(x,\mathbf v_{\hat\imath}).
\]
The sums cancel against Step~1.

\emph{Step 3 (the sum).} By~\eqref{dp:f:law-euler}, $D\eta_n(x)(x)=n\,\eta_n(x)$. So $(d\iota+\iota d)\eta_n=n\eta_n+k\eta_n$. For $k=0$ only Steps~2 and~3 occur, with empty sums.
\end{proof}

\begin{theorem}[weight splitting]\label{dp:thm:weight-splitting}
Let $F$ be a Banach space and $k\ge1$. Let $\eta\colon U\to\Alt^k(P;F)$ be DP on an open set $U$ containing $0$, with $B(0,R)\subseteq U$, $\eta(x)=\sum_n\eta_n(x)$ for $\norm x<R$, laws $\eta_n\in\cP_n(P;\Alt^k(P;F))$ with $\norm{\eta_n}\le MR^{-n}$, and $d\eta=0$ near $0$.
\begin{enumerate}
\item Every $\eta_n$ is closed.
\item If $\operatorname{char}K=p>0$, the series $\zeta=\sum_{p\nmid n+k}(n+k)^{-1}\iota\eta_n$ converges on $B(0,R)$ and is DP at $0$ and on $B(0,R/C_2)$. The series $\eta_{[p]}=\sum_{p\mid n+k}\eta_n$ is closed on $B(0,R/C_2)$, and $\eta=d\zeta+\eta_{[p]}$ there.
\item If $\operatorname{char}K=0$, the series $\zeta=\sum_n(n+k)^{-1}\iota\eta_n$ converges on $B(0,R)$ and is DP at $0$ and on $B(0,R/C_2)$, and $\eta=d\zeta$ there. So closed $F$-valued DP germs at $0$ of degree $k\ge1$ are exact (Poincar\'e lemma); by translation the same holds at every point.
\end{enumerate}
\end{theorem}

\begin{proof}
\emph{Step 1 (termwise derivative).} Let $r>0$, let $k'\in\{k-1,k\}$, so that $\Alt^{k'}(P;F)$ is complete, and let $g=\sum_na_n$ with laws $a_n\in\cP_n(P;\Alt^{k'}(P;F))$ and $\sup_n\norm{a_n}r^n<\infty$. Then $g$ is DP on $B(0,r)$ by Theorem~\ref{dp:thm:whole-ball}(1). Theorem~\ref{dp:thm:differentiable}(3) gives $Dg(x)=\sum_nDa_n(x)$ for $\norm x<r/C_2$, by~\eqref{dp:f:law-derivative}. Applying $\delta_{k'}$, we get $dg=\sum_nda_n$ on $B(0,r/C_2)$.

\emph{Step 2 (item (1)).} Step~1 with $g=\eta$ and $r=R$ gives $d\eta=\sum_nd\eta_n$ on $B(0,R/C_2)$. For $n\ge1$, the law $d\eta_n$ of degree $n-1$ has norm at most $(k+1)C_2^n\norm{\eta_n}\le(k+1)M(C_2/R)^n$. So this is an expansion as in Definition~\ref{dp:def:laws}(3), and by Lemma~\ref{dp:lem:coefficients-unique}(1) the condition $d\eta=0$ near $0$ forces $d\eta_n=0$ for every $n$.

\emph{Step 3 (inverses of integers).} There is $\mathsf z\ge0$ with $\abs{w^{-1}}\le w^{\mathsf z}$ for every integer $w\ge1$ invertible in $K$. If $\operatorname{char}K=p$, then $w$ prime to $p$ is a root of unity in $\F_p$, so $\abs{w^{-1}}=1$ and $\mathsf z=0$ works. Let $\operatorname{char}K=0$. The set of integers $w$ with $\abs w<1$ is a prime ideal of $\Z$, because $\abs{\cdot}$ is ultrametric and multiplicative. If it is zero, then $\mathsf z=0$ works. Otherwise it is generated by a prime $p_0$; define $\mathsf z>0$ by $\abs{p_0}=p_0^{-\mathsf z}$. If $p_0^e$ is the exact power of $p_0$ dividing $w$, then $\abs w=p_0^{-e\mathsf z}\ge w^{-\mathsf z}$.

\emph{Step 4 (the series, items (2), (3)).} Let $n+k$ be invertible in $K$, and put $\zeta_{n+1}=(n+k)^{-1}\iota\eta_n$; put $\zeta_{n+1}=0$ for the other $n$, and $\zeta_0=0$. By item~(1) and Proposition~\ref{dp:prop:euler}, $d\zeta_{n+1}=\eta_n$ when $n+k$ is invertible. By Step~3, $\norm{\zeta_{n+1}}\le(n+k)^{\mathsf z}MR^{-n}$. So $\sup_n\norm{\zeta_{n+1}}r^{n+1}<\infty$ for every $r<R$, the series $\zeta$ converges on $B(0,R)$, and it is DP at $0$. Step~1 for each $r<R$ shows that $\zeta$ is DP on $B(0,R/C_2)$ and that there $d\zeta$ is the sum of the $\eta_n$ with $n+k$ invertible. In characteristic $0$ this is all of $\eta$. In characteristic $p$, the series $\eta_{[p]}$ has the bounds of $\eta$, so Step~1 and item~(1) give $d\eta_{[p]}=\sum_{p\mid n+k}d\eta_n=0$ on $B(0,R/C_2)$, and $\eta-\eta_{[p]}=d\zeta$ there.

\emph{Step 5 (germs).} Let $\operatorname{char}K=0$, and let $\eta$ be a closed germ at $0$ of degree $k\ge1$ that is DP into the complete space $\Alt^k(P;F)$. It is DP at $0$ with some radius $R$ and bound $M$. By Theorem~\ref{dp:thm:whole-ball}(1) it is DP on $B(0,R)$, so the hypotheses hold with $U=B(0,R)$, and item~(3) applies. Translations are DP with DP inverses, so items~(1) and~(3) of Theorem~\ref{dp:thm:de-rham} carry the statement to every point.
\end{proof}
%% ---- end sec-08-forms.tex
%% ---- begin sec-09-quotients.tex
\section{Slices and difference quotients}\label{dp:sec:quotients}

Throughout this section $K$ is nonarchimedean, $P$ and $W$ are normed spaces, $U$ is an open subset of $P$, $\bar B(x,r)$ is the closed ball, and $\mathcal O=\{\lambda\in K:\abs\lambda\le1\}$. Radius and bound are as in Definition~\ref{dp:def:laws}, and $D^{(j)}f(x)$ is the Hasse derivative of $f$ at $x$, its coefficient law of degree $j$. Proposition~\ref{dp:prop:slices} restates the definition of a DP germ, and Proposition~\ref{dp:prop:div-bgn}(4) restates Theorem~\ref{dp:thm:divided-planes}; the other numbered statements have content of their own. Theorems~\ref{dp:thm:banach-steinhaus}, \ref{dp:thm:planes}, \ref{dp:thm:divided-one} and~\ref{dp:thm:divided-planes} replace the bounds on laws by a Baire hypothesis, by multilinear bounds on planes, or by bounds on divided differences. Theorem~\ref{dp:thm:bgn} and Proposition~\ref{dp:prop:div-bgn}(1)--(3) compare DP germs with difference quotients. Examples~\ref{dp:ex:slice-needed} and~\ref{dp:ex:zero-jets} show what the hypotheses contain.

\subsection{Slices}\label{dp:sq:ssec-slices}

A map on a finite-dimensional subspace $V\subseteq P$ is a \emph{homogeneous polynomial of degree $n$} if it is one, with coefficients in $W$, in the coordinates of some basis of $V$, equivalently of every basis.

\begin{lemma}[affine lines]\label{dp:sq:lines}
Let $n\ge0$ and let $\phi\colon P\to W$ satisfy $\phi(\lambda h)=\lambda^n\phi(h)$ for $\lambda\in K$ and $h\in P$. Suppose that for all $h,y\in P$ the map $s\mapsto\phi(h+sy)$ on $K$ is a polynomial of degree at most $n$ with coefficients in $W$. Then for all $h_1,\dots,h_r\in P$ the map $(s_1,\dots,s_r)\mapsto\phi(s_1h_1+\dots+s_rh_r)$ is a homogeneous polynomial of degree $n$ with coefficients in $W$.
\end{lemma}

\begin{proof}
\emph{Step 1 (separately polynomial).} We show by induction on $r$ that for all $h_0,\dots,h_r\in P$ the map $\Phi(s)=\phi(h_0+s_1h_1+\dots+s_rh_r)$ is a polynomial of degree at most $n$ in each variable. For $r=1$ this is the hypothesis. For $r\ge2$ fix distinct $\mu_0,\dots,\mu_n\in K$, and let $\ell_i$ be the Lagrange polynomials with $\ell_i(\mu_i)=1$ and $\ell_i(\mu_{i'})=0$ for $i'\ne i$. Then $\Phi(s)=\sum_{i\le n}\ell_i(s_1)\phi(h_0+\mu_ih_1+s_2h_2+\dots+s_rh_r)$, and by induction each value of $\phi$ here is a polynomial in $(s_2,\dots,s_r)$ of degree at most $n$ in each variable.

\emph{Step 2 (homogeneity).} Take $h_0=0$ and write $\Phi=\sum_d\Phi_d$ with $\Phi_d$ homogeneous of degree $d$. For fixed $s$ and all $\lambda\in K$, $\sum_d\lambda^d\Phi_d(s)=\Phi(\lambda s)=\lambda^n\Phi(s)$. As $K$ is infinite, $\Phi_d(s)=0$ for $d\ne n$.
\end{proof}

\begin{proposition}[slices]\label{dp:prop:slices}
Let $f\colon U\to W$, $x_0\in U$, $R>0$ with $B(x_0,R)\subseteq U$, and $M\ge0$. The following are equivalent.
\begin{enumerate}
\item $f$ is DP at $x_0$ with radius $R$ and bound $M$.
\item For every finite-dimensional subspace $V\subseteq P$ there are homogeneous polynomials $a_n^V\colon V\to W$ of degree $n$ with $\norm{a_n^V(h)}\le MR^{-n}\norm h^n$ for $h\in V$, and $f(x_0+h)=\sum_na_n^V(h)$ for $h\in V$ with $\norm h<R$.
\item The same as \textup{(2)}, for every subspace $V$ of dimension at most $2$.
\end{enumerate}
\end{proposition}

This is a restatement of the definition: the size of each homogeneous part is measured by its values, with the same $R$ and $M$. Its only content is Lemma~\ref{dp:sq:lines}, which lets planes replace all finite-dimensional subspaces.

\begin{proof}
\emph{Step 1 (restriction).} The restriction of a law of degree $n$ to $V$ is a homogeneous polynomial of degree $n$ with the same bound, so (1) implies (2). And (3) is a special case of (2).

\emph{Step 2 (compatibility).} Assume (3). Let $V\subseteq V'$ have dimension at most $2$ and let $h\in V\setminus\{0\}$. For $\abs s<R/\norm h$ both $\sum_ns^na_n^V(h)$ and $\sum_ns^na_n^{V'}(h)$ converge to $f(x_0+sh)$, and $\Delta_n=a_n^V(h)-a_n^{V'}(h)$ satisfies $\norm{\Delta_n}(R/\norm h)^n\le2M$. By~\cite[Lemma~A.1]{I}, $\Delta_n=0$. At $h=0$ both polynomials give $f(x_0)$ for $n=0$ and $0$ for $n\ge1$. So $a_n^{V'}=a_n^V$ on $V$.

\emph{Step 3 (the laws).} Put $a_n(h)=a_n^{Kh}(h)$. By Step 2, $a_n$ restricts to $a_n^V$ on every $V$ of dimension at most $2$. So $a_n(\lambda h)=\lambda^na_n(h)$, and $s\mapsto a_n(h+sy)$ is a polynomial of degree at most $n$, since $h+sy$ stays in the span of $h$ and $y$. By Lemma~\ref{dp:sq:lines}, $a_n$ is a polynomial law of degree $n$, and $\norm{a_n}\le MR^{-n}$. Finally $f(x_0+h)=\sum_na_n^{Kh}(h)=\sum_na_n(h)$ for $\norm h<R$.
\end{proof}

\begin{example}[what the uniformity contains]\label{dp:ex:slice-needed}
\leavevmode
\begin{enumerate}
\item \emph{Lines through the base point do not suffice.} Let $P=K^2$ with the maximum norm and $U$ its open unit ball. Define $g_1\colon U\to K$ by $g_1(\xi,\eta)=1/(1-\xi)$ if $\abs\xi\ge\abs\eta$ and $g_1(\xi,\eta)=1/(1-\eta)$ otherwise, and put $\psi_n(\xi,\eta)=\xi^n$ if $\abs\xi\ge\abs\eta$ and $\psi_n(\xi,\eta)=\eta^n$ otherwise. Then $g_1$ is continuous and bounded by $1$. For $h\ne0$ and $\abs s<1/\norm h$, $g_1(sh)=\sum_ns^n\psi_n(h)$ with $\abs{\psi_n(h)}=\norm h^n$. So on every line through $0$ it is a power series with the constants $R=M=1$. It is not DP at $0$.
\item \emph{The radius must be uniform in the slice.} Let $P=c_0(\N,K)$ with unit vectors $e_m$, let $B_m=B(\varpi^me_m,\abs\varpi^{m+1})$ for $m\ge1$, and define $g_2\colon P\to K$ by $g_2=\varpi^m$ on $B_m$ and $g_2=0$ off $\bigcup_mB_m$. The $B_m$ are pairwise disjoint, $g_2$ is continuous, $\abs{g_2(h)}\le\norm h$, and $g_2$ is locally constant on $P\setminus\{0\}$. For every finite-dimensional subspace $V$, $g_2$ vanishes on a neighborhood of $0$ in $V$. It is not DP at $0$.
\end{enumerate}
\end{example}

\begin{proof}
\emph{Step 1 (item (1)).} For $s\ne0$, $\abs{s\xi}\ge\abs{s\eta}$ if and only if $\abs\xi\ge\abs\eta$, and $1/(1-s\xi)=\sum_n(s\xi)^n$ for $\abs{s\xi}<1$; this gives the expansions. The sets $\{\abs\xi\ge\abs\eta\}\setminus\{0\}$ and $\{\abs\xi<\abs\eta\}$ are open, and $\abs{g_1(\xi,\eta)-1}\le\norm{(\xi,\eta)}$; this gives continuity at every point. If $g_1(h)=\sum_na_n(h)$ near $0$ with laws $a_n$, then $a_n=\psi_n$ by~\cite[Lemma~A.1]{I} on each line through $0$. But $\psi_2$ is not a law: $s\mapsto\psi_2(s,1)$ equals $1$ for $\abs s<1$ and $s^2$ for $\abs s\ge1$, and no polynomial agrees with both on these infinite sets.

\emph{Step 2 (item (2), disjointness and local constancy).} If $h\in B_m$, then $\abs{h_m}=\abs\varpi^m=\norm h$ and $\abs{h_i}<\abs\varpi^{m+1}$ for $i\ne m$. If also $h\in B_{m'}$ with $m<m'$, then $\abs{h_m}=\abs\varpi^m$ and $\abs{h_m}<\abs\varpi^{m'+1}$, which is impossible. Each $B_m$ is open and closed. Let $x\ne0$. If $\abs\varpi^m<\norm x$, then $B_m\subseteq B(0,\norm x)$, which is disjoint from $B(x,\norm x)$. So $B(x,\norm x)$ meets finitely many $B_m$, and $g_2$ is constant near $x$. At $0$, $\abs{g_2(h)}=\norm h$ on $B_m$.

\emph{Step 3 (item (2), slices).} Let $\widehat V$ be the $\Kh$-span of $V$ in $c_0(\N,\Kh)$, with basis $b_1,\dots,b_r$. By~\cite[Lemma~A.3]{I} there is $C_V$ with $\abs{\lambda_i}\le C_V\norm{\sum_i\lambda_ib_i}$. If $\mathsf y\in V\cap B_m$, then $\mathsf z=\varpi^{-m}\mathsf y\in\widehat V$ has $\norm{\mathsf z}=1$ and $\abs{\mathsf z_m}=1$, while $\abs{\mathsf z_m}\le C_V\max_i\abs{(b_i)_m}$, which tends to $0$ as $m\to\infty$. So $V\cap B_m=\emptyset$ for $m\ge m_V$, and $g_2=0$ on $V\cap B(0,\abs\varpi^{m_V-1})$.

\emph{Step 4 (item (2), not DP).} If $g_2(h)=\sum_na_n(h)$ for $\norm h<\sigma$, then by Step 3 and~\cite[Lemma~A.1]{I} every $a_n$ vanishes, so $g_2=0$ on $B(0,\sigma)$. But $g_2(\varpi^me_m)=\varpi^m\ne0$ for every $m\ge1$ with $\abs\varpi^m<\sigma$.
\end{proof}

\begin{theorem}[Banach--Steinhaus for laws]\label{dp:thm:banach-steinhaus}
Let $P$ be a Banach space, $x_0\in U$, $f\colon U\to W$ and $B(x_0,R)\subseteq U$. Let $a_n\colon P\to W$ be bounded polynomial laws of degree $n$, no bound uniform in $n$ being assumed, and suppose that $f(x_0+h)=\sum_na_n(h)$, convergent in $W$, for $\norm h<R$. Then there are $M\ge0$ and a ball $B(h_1,\delta)\subseteq B(0,R)$ with $\norm{a_n(h)}\le M$ on it for all $n$. For any such $M$, $h_1$ and $\delta$, $\norm{a_n}\le M(2/(\abs\varpi^2\delta))^n$, and $f$ is DP at $x_0$ with radius $\abs\varpi^2\delta/2$ and bound $M$.
\end{theorem}

\begin{proof}
\emph{Step 1 (Baire).} For $\norm h<R$ the terms $a_n(h)$ of a convergent series are bounded in $n$. So the sets $\{h\in B(0,R):\norm{a_n(h)}\le M\text{ for all }n\}$, $M\in\N$, cover $B(0,R)$. They are closed in $B(0,R)$, since laws are continuous (Theorem~\ref{dp:thm:calculus}(1)), and $B(0,R)$ is a Baire space. So one of them contains a ball.

\emph{Step 2 (leading coefficients).} Let $\norm y<\delta$. For $s\in\mathcal O$, $h_1+sy\in B(h_1,\delta)$, so the polynomial $G(s)=a_n(h_1+sy)$ of degree at most $n$ has $\norm{G(s)}\le M$ on $\mathcal O$. Its coefficient of $s^n$ is $a_n(y)$. Let $\mu_0,\dots,\mu_n\in\mathcal O$ be the first $n+1$ nodes of Step~2 of the proof of Theorem~\ref{dp:thm:mixed}, with $\varpi'=\varpi$. By~\eqref{dp:laws:nodes}, $\prod_{i'\ne i}\abs{\mu_i-\mu_{i'}}^{-1}\le\abs\varpi^{-n}$. By Lagrange interpolation $a_n(y)=\sum_iG(\mu_i)\prod_{i'\ne i}(\mu_i-\mu_{i'})^{-1}$. So $\norm{a_n(y)}\le(n+1)\abs\varpi^{-n}M\le(2/\abs\varpi)^nM$.

\emph{Step 3 (scaling).} We may assume $P\ne0$. Let $y\ne0$ and $0<\delta'<\delta$, and choose $\lambda\in K$ with $\abs\varpi\delta'\le\norm{\lambda y}<\delta'$. By Step 2, $\norm{a_n(y)}=\abs\lambda^{-n}\norm{a_n(\lambda y)}\le(\norm y/(\abs\varpi\delta'))^n(2/\abs\varpi)^nM$. Letting $\delta'\to\delta$ gives the bound on $\norm{a_n}$. Both $h_1$ and $h_1+\lambda y$ lie in $B(0,R)$, so $\abs\varpi\delta'\le\norm{\lambda y}<2R$. Hence $\sigma=\abs\varpi^2\delta/2$ satisfies $\sigma\le\abs\varpi R<R$, and $\sup_n\norm{a_n}\sigma^n\le M$.
\end{proof}

\begin{theorem}[analyticity on planes]\label{dp:thm:planes}
Let $K$ be complete.
\begin{enumerate}
\item Every subspace $V\subseteq P$ of dimension $r\le2$ has a linear isomorphism $T\colon K^r\to V$, with the maximum norm on $K^r$, such that $\norm T\le2$ and $\norm{T^{-1}}\le3/\abs\varpi$.
\item Every bounded polynomial law $a\colon V\to W$ of degree $n$ on such a $V$ is the diagonal of a bounded $n$\nobreakdash-linear map $V^n\to W$ of norm at most $(n+1)(96/\abs\varpi^2)^n\norm a\le(192/\abs\varpi^2)^n\norm a$.
\item Let $f\colon U\to W$, $x_0\in U$, $R>0$ with $B(x_0,R)\subseteq U$, and $M\ge0$. Consider the condition: for every subspace $V$ of dimension at most $2$ there are bounded $n$\nobreakdash-linear maps $\beta^V_n\colon V^n\to W$ with $\norm{\beta^V_n}\le MR^{-n}$ and $f(x_0+h)=\sum_n\beta^V_n(h,\dots,h)$ for $h\in V$, $\norm h<R$. If this condition holds with $(R,M)$, then $f$ is DP at $x_0$ with radius $R$ and bound $M$. If $f$ is DP at $x_0$ with radius $R$ and bound $M$, then the condition holds with $(\abs\varpi^2R/192,M)$.
\end{enumerate}
\end{theorem}

So, for $K$ complete, a DP germ is a germ that is analytic on every plane through its base point, with constants independent of the plane.

\begin{proof}
\emph{Step 1 (item (1), $r\le1$).} For $r=0$ there is nothing to prove. For $r=1$ pick $b_1\in V$ with $\abs\varpi<\norm{b_1}\le1$ and put $T(\lambda)=\lambda b_1$. Then $\norm T\le1$ and $\norm{T^{-1}}\le1/\abs\varpi$.

\emph{Step 2 (item (1), $r=2$).} Take $b_1$ as in Step 1. The line $Kb_1$ is complete, as $K$ is, hence closed in $V$. Take $b'\in V\setminus Kb_1$. Then $\mathsf D=\operatorname{dist}(b',Kb_1)>0$, and there is $b''\in Kb_1$ with $\norm{b'-b''}\le2\mathsf D$. Put $b_2=\lambda'(b'-b'')$ with $\abs\varpi<\norm{b_2}\le1$, so that $\operatorname{dist}(b_2,Kb_1)=\abs{\lambda'}\mathsf D\ge\norm{b_2}/2$. For $b=\lambda_1b_1+\lambda_2b_2$ we get $\norm b\ge\abs{\lambda_2}\operatorname{dist}(b_2,Kb_1)\ge\abs{\lambda_2}\norm{b_2}/2$, so $\abs{\lambda_2}\le2\norm b/\abs\varpi$. Then $\abs{\lambda_1}\norm{b_1}\le\norm b+\abs{\lambda_2}\norm{b_2}\le3\norm b$, so $\abs{\lambda_1}\le3\norm b/\abs\varpi$. Also $\norm T\le\norm{b_1}+\norm{b_2}\le2$.

\emph{Step 3 (item (2)).} Let $e_1,\dots,e_r$ be the unit vectors of $K^r$. Then $a(T\lambda)=\sum_{\abs\gamma=n}w_\gamma\lambda^\gamma$ for $\lambda\in K^r$, with $w_\gamma=(a\circ T)^{[\gamma]}(e_1,\dots,e_r)$, and $\norm{w_\gamma}\le(8r\norm T/\abs\varpi)^n\norm a$ by Theorem~\ref{dp:thm:mixed}(1). The $n$\nobreakdash-linear map $(h^1,\dots,h^n)\mapsto\sum_{\mathsf w_1\le\dots\le\mathsf w_n}w_{\gamma(\mathsf w)}\prod_l(T^{-1}h^l)_{\mathsf w_l}$, summed over the nondecreasing words $\mathsf w$ in $\{1,\dots,r\}$ with $\gamma(\mathsf w)$ counting the letters of $\mathsf w$, has diagonal $a$. Its norm is at most $\binom{n+r-1}n(8r\norm T\norm{T^{-1}}/\abs\varpi)^n\norm a$. For $r=2$ this is $(n+1)(96/\abs\varpi^2)^n\norm a$. For $r=1$ it is at most $(8/\abs\varpi^2)^n\norm a$, by the bounds of Step 1, and $r=0$ is clear. Finally $n+1\le2^n$.

\emph{Step 4 (item (3)).} On $V$ the diagonal of $\beta^V_n$ is a homogeneous polynomial of degree $n$ with $\norm{\beta^V_n(h,\dots,h)}\le MR^{-n}\norm h^n$, so the condition gives item (3) of Proposition~\ref{dp:prop:slices}. Conversely, let $f$ be DP with laws $a_n$, $\norm{a_n}\le MR^{-n}$. By item (2) the restriction $a_n|_V$ has a lift $\beta^V_n$ with $\norm{\beta^V_n}\le(192/\abs\varpi^2)^nMR^{-n}$, and the expansion holds for $\norm h<\abs\varpi^2R/192\le R$.
\end{proof}

\subsection{Difference quotients}\label{dp:sq:ssec-quotients}

\begin{definition}\label{dp:def:bgn}
Let $P_0=P$ and $P_{k+1}=P_k\times P_k\times K$, and let $U^{[0]}=U$ and $U^{[k+1]}=\{(x,v,\tau)\in U^{[k]}\times P_k\times K:x+\tau v\in U^{[k]}\}$, an open subset of $P_{k+1}$. A map $f\colon U\to W$ is \emph{$C^0_{\mathrm{BGN}}$} if it is continuous. It is \emph{$C^1_{\mathrm{BGN}}$} if there is a continuous $f^{[1]}\colon U^{[1]}\to W$ with
\[
 f^{[1]}(x,v,\tau)=\tau^{-1}\bigl(f(x+\tau v)-f(x)\bigr)\qquad(\tau\ne0).
\]
It is \emph{$C^{k+1}_{\mathrm{BGN}}$}, $k\ge1$, if it is $C^1_{\mathrm{BGN}}$ and $f^{[1]}$ is $C^k_{\mathrm{BGN}}$ on $U^{[1]}$, and it is \emph{$C^\infty_{\mathrm{BGN}}$} if it is $C^k_{\mathrm{BGN}}$ for all $k$.
\end{definition}

This is the differentiability of~\cite{BGN}, specialized to normed spaces; the statements below concern the definition as written here. The map $f^{[1]}$ is unique, since the points with $\tau\ne0$ are dense in $U^{[1]}$. Put $f^{[0]}=f$ and $f^{[k+1]}=(f^{[k]})^{[1]}$ on $U^{[k+1]}$. By induction $(f^{[1]})^{[k]}=f^{[k+1]}$, and $f$ is $C^k_{\mathrm{BGN}}$ if and only if $f^{[j]}$ exists and is continuous for every $j\le k$.

\begin{theorem}[DP germs and difference quotients]\label{dp:thm:bgn}
Let $f\colon U\to W$ be DP on $U$.
\begin{enumerate}
\item For $x\in U$, $R$ such that $f$ is DP at $x$ with radius $R$, $v\in P$ and $\tau\in K$ with $\abs\tau\norm v<R$, the map $f^{[1]}$ of item \textup{(2)} satisfies
\begin{equation}\label{dp:sq:quotient-series}
 f^{[1]}(x,v,\tau)=\sum_{j\ge1}\tau^{j-1}D^{(j)}f(x)(v),\qquad\text{in particular } f^{[1]}(x,v,0)=D^{(1)}f(x)v=Df(x)v.
\end{equation}
\item $f$ is $C^1_{\mathrm{BGN}}$, and $f^{[1]}$ is DP on $U^{[1]}$.
\item $f$ is $C^\infty_{\mathrm{BGN}}$, and every $f^{[k]}$ is DP on $U^{[k]}$.
\item Let $x\in U$, $v\in P$ and $k\ge0$. If $f(x+\tau v)=\sum_{j\le k}\tau^jw_j+o(\abs\tau^k)$ as $\tau\to0$, with $w_j\in W$, then $w_j=D^{(j)}f(x)(v)$. So the Hasse derivatives are the coefficients of every finite Taylor expansion along lines.
\end{enumerate}
\end{theorem}

\begin{proof}
\emph{Step 1 (the map, item (1)).} Let $\abs\tau\norm v<R$. By homogeneity $\sum_{j\ge0}\tau^jD^{(j)}f(x)(v)$ converges in $W$ to $f(x+\tau v)$. Subtracting $f(x)=D^{(0)}f(x)$ and dividing by $\tau\ne0$ gives~\eqref{dp:sq:quotient-series} for $\tau\ne0$. Define $f^{[1]}$ on $U^{[1]}$ by the quotient for $\tau\ne0$ and by $f^{[1]}(x,v,0)=D^{(1)}f(x)v$, which is $Df(x)v$ by Theorem~\ref{dp:thm:differentiable}(1). Since $\norm{D^{(j)}f(x)(v)}\le MR^{-j}\norm v^j$ for a bound $M$ at $x$,~\eqref{dp:sq:quotient-series} also gives
\begin{equation}\label{dp:sq:limit}
 \lim_{\tau\to0}f^{[1]}(x,v,\tau)=f^{[1]}(x,v,0).
\end{equation}

\emph{Step 2 (DP at points with $\tau_0\ne0$).} Near $(x_0,v_0,\tau_0)\in U^{[1]}$ with $\tau_0\ne0$ we have $f^{[1]}(x,v,\tau)=\tau^{-1}(f(x+\tau v)-f(x))$. The map $(x,v,\tau)\mapsto x+\tau v$ is a bounded linear map plus the diagonal of a bounded bilinear map, hence DP (Theorem~\ref{dp:thm:calculus}(1),(3)). The map $\tau\mapsto\tau^{-1}$ is analytic at $\tau_0$, hence DP there. Scalar multiplication $K\times W\to W$ is bounded bilinear. A pair of DP maps is DP into the product, since pairs of laws are laws for the maximum norm. A finite sum or difference of maps DP at a point is DP there, as sums and differences of laws are laws. By composition (Theorem~\ref{dp:thm:calculus}(2)), $f^{[1]}$ is DP at $(x_0,v_0,\tau_0)$.

\emph{Step 3 (points $(x_0,v_0,0)$: the expansion).} Let $f$ be DP at $x_0$ with radius $R$, bound $M$ and laws $a_n$. Put $\Lambda_0=\max(1,\norm{v_0})$. For $(h,u,\tau)\in P_1$ write $v=v_0+u$. If $\norm h<R$ and $\norm{h+\tau v}<R$, then
\[
 f(x_0+h+\tau v)-f(x_0+h)=\sum_{n\ge1}\sum_{j=1}^n\sum_{i=0}^j\tau^ja_n^{[(n-j,j-i,i)]}(h,v_0,u),
\]
by Lemma~\ref{dp:lem:elementary}(3) and a comparison of the coefficients of $a_n(\lambda_1h+\lambda_2v_0+\lambda_2u)$. Define $\mathsf t_{n,j,i}(h,u,\tau)=\tau^{j-1}a_n^{[(n-j,j-i,i)]}(h,v_0,u)$. Then $\mathsf t_{n,j,i}$ is a bounded polynomial law of degree $n+i-1$ on $P_1$. Indeed, for vectors $(h_c,u_c,\tau_c)$ of $P_1$ the vector $a_n^{[(n-j,j-i,i)]}(\sum_cs_ch_c,v_0,\sum_cs_cu_c)$ is the coefficient of $\lambda_1^{n-j}\lambda_2^{j-i}\lambda_3^i$ in $a_n(\sum_c\lambda_1s_ch_c+\lambda_2v_0+\sum_c\lambda_3s_cu_c)$, so it is a homogeneous polynomial in $s$ of degree $n-j+i$. Also $(\sum_cs_c\tau_c)^{j-1}$ is a homogeneous polynomial in $s$ of degree $j-1$, so the product is homogeneous of degree $n+i-1$ with coefficients in $W$. As $n!/\alpha!\le3^n$ for three indices, Theorem~\ref{dp:thm:mixed}(3) gives
\begin{equation}\label{dp:sq:t-bound}
 \norm{\mathsf t_{n,j,i}(h,u,\tau)}\le(3C_0)^nMR^{-n}\norm{v_0}^{j-i}\,\norm{(h,u,\tau)}^{n+i-1}.
\end{equation}
Fix $0<r_0\le1$ with $3C_0\Lambda_0r_0<R$ and $r_0(1+\Lambda_0+r_0)<R$. For $\norm{(h,u,\tau)}\le r_0$ we have $\norm h<R$ and $\norm{h+\tau v}<R$, and since $\norm{v_0}^{j-i}r_0^i\le\Lambda_0^n$,
\[
 \sum_{n\ge1}\sum_{j=1}^n\sum_{i=0}^j\norm{\mathsf t_{n,j,i}(h,u,\tau)}\le\sum_{n\ge1}(n+1)^2M(3C_0\Lambda_0/R)^nr_0^{n-1}<\infty.
\]
So the family is absolutely summable in $\widehat W$, uniformly on this ball. For $\tau\ne0$ its sum is $f^{[1]}(x_0+h,v,\tau)$. Each term is continuous in $\tau$, so the sum is continuous in $\tau$, and~\eqref{dp:sq:limit} at the point $x_0+h$ shows that its value at $\tau=0$ is $f^{[1]}(x_0+h,v,0)$.

\emph{Step 4 (points $(x_0,v_0,0)$: grouping).} For $N\ge0$ let $\mathsf a_N=\sum_{n+i-1=N}\mathsf t_{n,j,i}$, a finite sum of at most $(N+2)^2$ terms with $n\le N+1$, so a bounded polynomial law of degree $N$ on $P_1$. With $\Lambda=3C_0\Lambda_0\max(1,R^{-1})\ge1$,~\eqref{dp:sq:t-bound} gives $\norm{\mathsf a_N}\le(N+2)^2M\Lambda^{N+1}$. Regrouping by $N$ gives $f^{[1]}(x_0+h,v_0+u,\tau)=\sum_N\mathsf a_N(h,u,\tau)$ for $\norm{(h,u,\tau)}\le r_0$, convergent in $\widehat W$, hence in $W$, as the terms and the sum lie in $W$. So $f^{[1]}$ is DP at $(x_0,v_0,0)$, with radius $\min(r_0,\sigma)$ for any $\sigma<1/\Lambda$.

\emph{Step 5 (items (2), (3)).} By Steps 2--4, $f^{[1]}$ is DP on $U^{[1]}$, hence continuous (Theorem~\ref{dp:thm:differentiable}(1)). So $f$ is $C^1_{\mathrm{BGN}}$. Steps 1--4 apply to $f^{[1]}$ on the open subset $U^{[1]}$ of $P_1$, and item (3) follows by induction.

\emph{Step 6 (item (4)).} By~\eqref{dp:sq:quotient-series}, $f(x+\tau v)=\sum_{j\le k}\tau^jD^{(j)}f(x)(v)+O(\abs\tau^{k+1})$. So $\sum_{j\le k}\tau^j(w_j-D^{(j)}f(x)(v))=o(\abs\tau^k)$. If $j_0$ is the least index with a nonzero difference, dividing by $\tau^{j_0}$ and letting $\tau\to0$ gives a contradiction.
\end{proof}

\subsection{Divided differences}\label{dp:sq:ssec-divided}

For a map $\phi$ from a subset $S\subseteq K$ to $W$ and pairwise distinct $t_0,\dots,t_k\in S$, put
\[
 \operatorname{dd}_k\phi(t_0,\dots,t_k)=\sum_{i=0}^k\frac{\phi(t_i)}{\prod_{l\ne i}(t_i-t_l)}.
\]
Let $\mathsf N_k(t)=\prod_{i<k}(t-t_i)$, and let $\mathsf h_r$ be the complete homogeneous symmetric polynomial of degree $r$, with $\mathsf h_r=0$ for $r<0$.

\begin{lemma}[interpolation]\label{dp:sq:interpolation}
\leavevmode
\begin{enumerate}
\item $\operatorname{dd}_k\phi$ is symmetric and linear in $\phi$, $\operatorname{dd}_0\phi(t_0)=\phi(t_0)$, and for $k\ge1$
\begin{equation}\label{dp:sq:dd-recursion}
 \operatorname{dd}_k\phi(t_0,\dots,t_k)=\frac{\operatorname{dd}_{k-1}\phi(t_0,\dots,t_{k-1})-\operatorname{dd}_{k-1}\phi(t_1,\dots,t_k)}{t_0-t_k}.
\end{equation}
\item For $t\notin\{t_0,\dots,t_{n-1}\}$,
$\phi(t)=\sum_{k<n}\operatorname{dd}_k\phi(t_0,\dots,t_k)\mathsf N_k(t)+\operatorname{dd}_n\phi(t_0,\dots,t_{n-1},t)\mathsf N_n(t)$. For $t=t_l$ with $l<n$, $\phi(t_l)=\sum_{k\le l}\operatorname{dd}_k\phi(t_0,\dots,t_k)\mathsf N_k(t_l)$.
\item For $x_0\in K$ and $m\ge0$, $\operatorname{dd}_k[t\mapsto(t-x_0)^m](t_0,\dots,t_k)=\mathsf h_{m-k}(t_0-x_0,\dots,t_k-x_0)$, of absolute value at most $\max_i\abs{t_i-x_0}^{m-k}$. So $\operatorname{dd}_k$ kills polynomials of degree below $k$. The coefficient of $h^l$ in $\mathsf N_k(x_0+h)$ has absolute value at most $\max_{i<k}\abs{t_i-x_0}^{k-l}$. Here we read $0^0=1$, and the maximum of the empty set as $0$.
\end{enumerate}
\end{lemma}

\begin{proof}
\emph{Step 1 (item (1)).} Compare the coefficient of each $\phi(t_i)$ on both sides of~\eqref{dp:sq:dd-recursion}.

\emph{Step 2 (item (2)).} Induct on $n$ with $\operatorname{dd}_n\phi(t_0,\dots,t_{n-1},t)=\operatorname{dd}_n\phi(t_0,\dots,t_n)+(t-t_n)\operatorname{dd}_{n+1}\phi(t_0,\dots,t_n,t)$, which is~\eqref{dp:sq:dd-recursion} with symmetry; for $t=t_l$ use the formula with $n=l$.

\emph{Step 3 (item (3)).} We may take $x_0=0$, as the differences do not change. Induct on $k$ with~\eqref{dp:sq:dd-recursion} and $\prod_{i<k}(1-t_iX)^{-1}-\prod_{0<i\le k}(1-t_iX)^{-1}=X(t_0-t_k)\prod_{i\le k}(1-t_iX)^{-1}$, the generating function identity of the $\mathsf h_r$. The two bounds hold because these polynomials are sums of monomials with coefficients $\pm1$ and $K$ is nonarchimedean.
\end{proof}

\begin{theorem}[one variable]\label{dp:thm:divided-one}
Let $U\subseteq K$ be open, $x_0\in U$ and $f\colon U\to W$. A germ on $K$ is DP with radius $R$ and bound $M$ if and only if $f(x_0+h)=\sum_kw_kh^k$ for $\abs h<R$ with $\norm{w_k}R^k\le M$, that is, if and only if it is analytic with the same radius and bound.
\begin{enumerate}
\item If $f$ is DP at $x_0$ with radius $R$ and bound $M$, then for $0<\rho<R$, all $k\ge0$ and pairwise distinct $t_0,\dots,t_k\in\bar B(x_0,\rho)$,
\begin{equation}\label{dp:sq:one-variable-bound}
 \norm{\operatorname{dd}_kf(t_0,\dots,t_k)}\le\frac{M}{1-\rho/R}R^{-k},\qquad\text{and}\quad\le MR^{-k}\text{ if $W$ is ultrametric}.
\end{equation}
\item Let $W$ be complete, $\rho,\sigma>0$ with $B(x_0,\rho)\subseteq U$, and $M\ge0$ with $\norm{\operatorname{dd}_kf(t_0,\dots,t_k)}\le M\sigma^{-k}$ for all $k\ge0$ and all pairwise distinct $t_0,\dots,t_k\in B(x_0,\rho)$. Then $f$ is DP at $x_0$ with radius $\min(\rho,\sigma)$ and bound $M$, and its coefficients satisfy $\norm{w_k}\le M\sigma^{-k}$.
\item So for complete $W$, $f$ is DP at $x_0$ if and only if such $\rho,\sigma,M$ exist, and the supremum of the DP radii at $x_0$ is the supremum of $\min(\rho,\sigma)$ over such pairs.
\end{enumerate}
No continuity of $f$ is assumed in item \textup{(2)}.
\end{theorem}

\begin{proof}
\emph{Step 1 (the first sentence).} A law on $K$ of degree $k$ is $h\mapsto w_kh^k$, the diagonal of $(h_1,\dots,h_k)\mapsto w_kh_1\cdots h_k$ of the same norm.

\emph{Step 2 (item (1)).} As $\operatorname{dd}_k$ is a combination of values, Lemma~\ref{dp:sq:interpolation}(3) writes $\operatorname{dd}_kf(t_0,\dots,t_k)$ as the sum over $m\ge k$ of $w_m\mathsf h_{m-k}(t_0-x_0,\dots,t_k-x_0)$, with terms of norm at most $MR^{-m}\rho^{m-k}$. Summing, or taking the maximum if $W$ is ultrametric, gives~\eqref{dp:sq:one-variable-bound}.

\emph{Step 3 (item (2), a Newton series).} Put $R'=\min(\rho,\sigma)$. Fix $0<\rho'<R'$ and pairwise distinct nodes $t_0,t_1,\dots\in\bar B(x_0,\rho')$, which exist as balls of positive radius in $K$ are infinite. For $\abs h\le\rho'$ and $t=x_0+h$ we have $\abs{t-t_i}\le\rho'$. If $t$ is not a node, the remainder in Lemma~\ref{dp:sq:interpolation}(2) has norm at most $M\sigma^{-n}\rho'^n\to0$; if $t$ is a node the expansion terminates. So $f(t)=\sum_k\operatorname{dd}_kf(t_0,\dots,t_k)\mathsf N_k(t)$.

\emph{Step 4 (item (2), power series).} Write $\mathsf N_k(x_0+h)=\sum_{l\le k}\mathsf n_{k,l}h^l$, with $\abs{\mathsf n_{k,l}}\le\rho'^{k-l}$. For $\abs h\le\rho'$ the double family $\operatorname{dd}_kf(t_0,\dots,t_k)\mathsf n_{k,l}h^l$ has norms at most $M(\rho'/\sigma)^k$, so it is absolutely summable. Hence $f(x_0+h)=\sum_lw_lh^l$ in the complete space $W$, with $\norm{w_l}\le M\sigma^{-l}/(1-\rho'/\sigma)$.

\emph{Step 5 (item (2), independence of the choices).} For two choices of $\rho'$ and nodes, both power series represent $f$ on a ball, with geometric bounds, so by~\cite[Lemma~A.1]{I} they have the same coefficients. Letting $\rho'\to0$ gives $\norm{w_l}\le M\sigma^{-l}$, and letting $\rho'\to R'$ gives the expansion on $B(x_0,R')$. Then $\norm{w_l}R'^l\le M$.

\emph{Step 6 (item (3)).} Items (1) and (2), with $\rho<R$ arbitrary in item (1) and $\sigma=R$.
\end{proof}

Let $f\colon U\to W$ and $x_0\in U$. For $\mathbf v=(v_1,v_2)\in P^2$ with $\norm{v_1},\norm{v_2}\le1$ put $F_{\mathbf v}(s_1,s_2)=f(x_0+s_1v_1+s_2v_2)$. For pairwise distinct $t_0,\dots,t_{k_1}$ and pairwise distinct $t'_0,\dots,t'_{k_2}$ in $K$, let $\operatorname{dd}_{k_1,k_2}F_{\mathbf v}(t;t')$ be the divided difference of order $k_1$ in $s_1$ of the divided difference of order $k_2$ in $s_2$. The two operations commute, as both give one finite sum of values.

\begin{theorem}[planes]\label{dp:thm:divided-planes}
\leavevmode
\begin{enumerate}
\item If $f$ is DP at $x_0$ with radius $R$ and bound $M$, put $\rho=R/(2C_2)$ and $\sigma=R/C_2$. For all $\mathbf v$ as above, all $k_1,k_2$ and all pairwise distinct nodes in $B(0,\rho)$, $F_{\mathbf v}$ is defined at the nodes and $\norm{\operatorname{dd}_{k_1,k_2}F_{\mathbf v}(t;t')}\le4M\sigma^{-(k_1+k_2)}$.
\item Let $W$ be complete, let $\rho,\sigma>0$ and $M\ge0$, and suppose that $B(x_0,2\rho)\subseteq U$ and that $\norm{\operatorname{dd}_{k_1,k_2}F_{\mathbf v}(t;t')}\le M\sigma^{-(k_1+k_2)}$ for all $\mathbf v$, all $k_1,k_2$ and all pairwise distinct nodes in $B(0,\rho)$. Then $f$ is DP at $x_0$ with radius $\abs\varpi\min(\rho,\sigma)$ and bound $M$.
\item Lines do not suffice: the map $g_1$ of Example~\textup{\ref{dp:ex:slice-needed}(1)} satisfies, on every line $s\mapsto sh$ with $\norm h=1$ and every $0<\rho<1$, the bound $\abs{\operatorname{dd}_k[s\mapsto g_1(sh)]}\le1$ on $\bar B(0,\rho)$, and it is not DP at $0$.
\end{enumerate}
\end{theorem}

\begin{proof}
\emph{Step 1 (item (1)).} Let $a_n$ be the laws of $f$ at $x_0$. For $\abs{s_1},\abs{s_2}<\rho$ we have $\norm{s_1v_1+s_2v_2}<2\rho<R$, and $F_{\mathbf v}(s)=\sum_n\sum_{m_1+m_2=n}a_n^{[(m_1,m_2)]}(v_1,v_2)s_1^{m_1}s_2^{m_2}$ by Lemma~\ref{dp:lem:elementary}(3), with $\norm{a_n^{[(m_1,m_2)]}(v_1,v_2)}\le C_2^nMR^{-n}$ by Theorem~\ref{dp:thm:mixed}(2). By Lemma~\ref{dp:sq:interpolation}(3) in each variable, the divided differences of $s_1^{m_1}s_2^{m_2}$ have absolute value at most $\rho^{m_1+m_2-k_1-k_2}$. Since $C_2\rho/R=1/2$,
\[
 \norm{\operatorname{dd}_{k_1,k_2}F_{\mathbf v}}\le\sum_{m_1\ge k_1,\,m_2\ge k_2}M(C_2/R)^{m_1+m_2}\rho^{m_1+m_2-k_1-k_2}=4M\sigma^{-(k_1+k_2)}.
\]

\emph{Step 2 (item (2), two-variable Newton series).} Fix $\mathbf v$, put $R'=\min(\rho,\sigma)$, and fix $0<\rho'<R'$ and pairwise distinct nodes $t_0,t_1,\dots$ and $t'_0,t'_1,\dots$ in $\bar B(0,\rho')$. Put $\mathsf N'_k(t)=\prod_{i<k}(t-t'_i)$. For $s\in\bar B(0,\rho')^2$ apply Lemma~\ref{dp:sq:interpolation}(2) in $s_2$, with remainder bounded by the hypothesis for $k_1=0$, and then in $s_1$, with remainder at most $M\sigma^{-n-k_2}\rho'^n$. This gives
\[
 F_{\mathbf v}(s)=\sum_{k_1,k_2}\operatorname{dd}_{k_1,k_2}F_{\mathbf v}(t_0,\dots,t_{k_1};t'_0,\dots,t'_{k_2})\mathsf N_{k_1}(s_1)\mathsf N'_{k_2}(s_2),
\]
with terms of norm at most $M(\rho'/\sigma)^{k_1+k_2}$. Expanding the Newton polynomials as in Step 4 of the proof of Theorem~\ref{dp:thm:divided-one} gives $F_{\mathbf v}(s)=\sum_{m\in\N^2}w^{\mathbf v}_ms^m$ on $\bar B(0,\rho')^2$, with $\norm{w^{\mathbf v}_m}\le M\sigma^{-\abs m}/(1-\rho'/\sigma)^2$. Two-variable power series with geometric bounds that agree on a polydisc have equal coefficients, by~\cite[Lemma~A.1]{I} in each variable. So the $w^{\mathbf v}_m$ do not depend on the choices, $\norm{w^{\mathbf v}_m}\le M\sigma^{-\abs m}$, and the expansion holds for $\abs{s_1},\abs{s_2}<R'$.

\emph{Step 3 (item (2), the laws).} For $h\ne0$ choose $\lambda\in K$ with $\abs\varpi<\norm{\lambda h}\le1$ and put $\mathbf v=(\lambda h,0)$. By Step 2, for $\abs s<\abs\lambda R'$ we have $f(x_0+sh)=\sum_n\lambda^{-n}w^{\mathbf v}_{(n,0)}s^n$, and by~\cite[Lemma~A.1]{I} the coefficients $a_n(h)=\lambda^{-n}w^{\mathbf v}_{(n,0)}$ depend only on $h$. Put $a_n(0)=0$ for $n\ge1$ and $a_0(0)=f(x_0)$. Then $a_n(\theta h)=\theta^na_n(h)$ for $\theta\in K$. For any $\mathbf v=(v_1,v_2)$ with $\norm{v_1},\norm{v_2}\le1$, any $\theta_1,\theta_2\in K$ and small $s$, regrouping gives $f(x_0+s(\theta_1v_1+\theta_2v_2))=\sum_ns^n\sum_{\abs m=n}w^{\mathbf v}_m\theta_1^{m_1}\theta_2^{m_2}$. So $a_n(\theta_1v_1+\theta_2v_2)=\sum_{\abs m=n}w^{\mathbf v}_m\theta_1^{m_1}\theta_2^{m_2}$, by the same lemma. Any two vectors of $P$ are scalar multiples of vectors of norm at most $1$, so $a_n$ is a homogeneous polynomial of degree $n$ on every plane through $0$. Hence it is a polynomial of degree at most $n$ on every affine line, and Lemma~\ref{dp:sq:lines} makes $a_n$ a polynomial law of degree $n$.

\emph{Step 4 (item (2), bounds).} With $\lambda$ and $\mathbf v=(\lambda h,0)$ as in Step 3, $\norm{a_n(h)}=\abs\lambda^{-n}\norm{w^{\mathbf v}_{(n,0)}}\le(\norm h/(\abs\varpi\sigma))^nM$, so $\norm{a_n}\le M(\abs\varpi\sigma)^{-n}$. If $\norm h<\abs\varpi R'$, then $\abs\lambda^{-1}<R'$, and $f(x_0+h)=F_{\mathbf v}(\lambda^{-1},0)=\sum_na_n(h)$. Also $\norm{a_n}(\abs\varpi R')^n\le M$.

\emph{Step 5 (item (3)).} By Example~\ref{dp:ex:slice-needed}(1), $s\mapsto g_1(sh)$ is a power series with coefficients of absolute value $1$ and radius $1$. Apply Theorem~\ref{dp:thm:divided-one}(1) with ultrametric values.
\end{proof}

For $U\subseteq K$ open, a map $\phi\colon U\to W$ is \emph{$C^k_{\mathrm{div}}$} if for every $j\le k$ the map $\operatorname{dd}_j\phi$, on the pairwise distinct tuples of $U$, extends to a continuous map $\overline{\operatorname{dd}}_j\phi$ on $U^{j+1}$; its \emph{diagonal values} are the $\overline{\operatorname{dd}}_j\phi(x,\dots,x)$. It is \emph{$C^\infty_{\mathrm{div}}$} if this holds for all $k$. If a $C^k_{\mathrm{div}}$ map $\phi$ is DP at $x$ and $D^{(k)}\phi(x)(h)=w_kh^k$, then $\overline{\operatorname{dd}}_k\phi(x,\dots,x)=w_k$: recenter Step 2 of the proof of Theorem~\ref{dp:thm:divided-one} at $x$ and let all $t_i\to x$. So the diagonal values are the Hasse derivatives. The next example shows that the bounds of Theorems~\ref{dp:thm:divided-one} and~\ref{dp:thm:divided-planes} cannot be restricted to the diagonal.

\begin{example}[zero jets, not DP]\label{dp:ex:zero-jets}
Let $K$ be complete. Define $\omega\colon\mathcal O\to K$ by $\omega(0)=0$ and $\omega(x)=\varpi^{\mathsf r^2}$ if $\abs\varpi^{\mathsf r+1}<\abs x\le\abs\varpi^{\mathsf r}$, $\mathsf r\ge0$. Put $\bar B_{\mathsf r}=\bar B(0,\abs\varpi^{\mathsf r})$. Say that $\bar B_{\mathsf r}$ \emph{separates} a tuple if some of its entries lie in $\bar B_{\mathsf r}$ and some do not.
\begin{enumerate}
\item $\omega$ is locally constant on $\mathcal O\setminus\{0\}$ and Fr\'echet differentiable on $\mathcal O$ with $\omega'\equiv0$, so it is Fr\'echet $C^\infty$ with all derivatives of order at least $1$ zero. It is $C^\infty_{\mathrm{div}}$, with $\overline{\operatorname{dd}}_k\omega(x,\dots,x)=0$ for all $k\ge1$ and $x\in\mathcal O$.
\item For $k\ge1$ and pairwise distinct $t_0,\dots,t_k\in\mathcal O$,
\begin{equation}\label{dp:sq:omega-bound}
 \abs{\operatorname{dd}_k\omega(t_0,\dots,t_k)}\le\max\bigl\{\abs\varpi^{(\mathsf r-1)^2-\mathsf rk}:\mathsf r\ge1,\ \bar B_{\mathsf r}\text{ separates }(t_0,\dots,t_k)\bigr\},
\end{equation}
where the maximum of the empty set is $0$.
\item $\omega$ is not DP at $0$. For $k\ge1$, $0\le\mathsf r<\mathsf r'$ and the pairwise distinct points $t_0=\varpi^{\mathsf r}$, $t_l=\varpi^{\mathsf r'}(1+\varpi^l)$, $1\le l\le k$, we have $\abs{\operatorname{dd}_k\omega(t_0,\dots,t_k)}=\abs\varpi^{\mathsf r^2-k\mathsf r}$.
\end{enumerate}
\end{example}

\begin{proof}
\emph{Step 1 (decomposition).} Put $\mathsf v_{\mathsf r}=\varpi^{\mathsf r^2}$ and, for $\mathsf r\ge1$, $\mathsf d_{\mathsf r}=\mathsf v_{\mathsf r}-\mathsf v_{\mathsf r-1}$, so that $\abs{\mathsf d_{\mathsf r}}=\abs\varpi^{(\mathsf r-1)^2}$. If $\abs\varpi^{\mathsf r'+1}<\abs x\le\abs\varpi^{\mathsf r'}$, then $x\in\bar B_{\mathsf r}$ if and only if $\mathsf r\le\mathsf r'$. So $\omega=\mathsf v_0+\sum_{\mathsf r\ge1}\mathsf d_{\mathsf r}1_{\bar B_{\mathsf r}}$ pointwise on $\mathcal O$, also at $x=0$.

\emph{Step 2 (one ball).} Let $B=\bar B(0,\rho)$, $\rho>0$, and let $t_0,\dots,t_k$ be pairwise distinct, $k\ge1$. Let $I$ be the set of $i$ with $t_i\in B$ and $J$ the rest. If $I$ or $J$ is empty, $\operatorname{dd}_k1_B(t)=0$. Otherwise, with $\Upsilon_J(x)=\prod_{j\in J}(x-t_j)^{-1}$ and $\mathsf i=\abs I-1$,
\[
 \operatorname{dd}_k1_B(t)=\sum_{i\in I}\frac{\Upsilon_J(t_i)}{\prod_{i'\in I\setminus\{i\}}(t_i-t_{i'})}=\operatorname{dd}_{\mathsf i}\Upsilon_J(t_i:i\in I).
\]
Let $\varrho=\min_{j\in J}\abs{t_j}>\rho$. On $\abs x\le\rho$ we have $(x-t_j)^{-1}=-\sum_mx^mt_j^{-m-1}$, so $\Upsilon_J(x)=\sum_m\mathsf b_mx^m$ with $\abs{\mathsf b_m}\le\varrho^{-m-\abs J}$. By Lemma~\ref{dp:sq:interpolation}(3), $\operatorname{dd}_{\mathsf i}\Upsilon_J(t_i:i\in I)=\sum_{m\ge\mathsf i}\mathsf b_m\mathsf h_{m-\mathsf i}(t_i:i\in I)$, of absolute value at most $\varrho^{-\mathsf i-\abs J}=\varrho^{-k}<\rho^{-k}$.

\emph{Step 3 (item (2)).} By Step 1, $\operatorname{dd}_k\omega(t)=\sum_{\mathsf r\ge1}\mathsf d_{\mathsf r}\operatorname{dd}_k1_{\bar B_{\mathsf r}}(t)$, as $\operatorname{dd}_k$ is a finite combination of values and kills $\mathsf v_0$. By Step 2 only the $\mathsf r$ for which $\bar B_{\mathsf r}$ separates the tuple contribute, each with absolute value at most $\abs\varpi^{(\mathsf r-1)^2-\mathsf rk}$, which tends to $0$ as $\mathsf r\to\infty$. This gives~\eqref{dp:sq:omega-bound}. For $\mathsf T\ge1$ write $\omega=\omega_{\le\mathsf T}+\omega_{>\mathsf T}$ with $\omega_{\le\mathsf T}=\mathsf v_0+\sum_{\mathsf r\le\mathsf T}\mathsf d_{\mathsf r}1_{\bar B_{\mathsf r}}$. The same argument gives $\abs{\operatorname{dd}_k\omega_{>\mathsf T}}\le\varepsilon_{\mathsf T}=\sup_{\mathsf r>\mathsf T}\abs\varpi^{(\mathsf r-1)^2-\mathsf rk}$ on distinct tuples, and $\varepsilon_{\mathsf T}\to0$.

\emph{Step 4 (item (1)).} The sets $\{\abs\varpi^{\mathsf r+1}<\abs x\le\abs\varpi^{\mathsf r}\}$ are open. At $0$, on such a set, $\abs{\omega(x)}/\abs x<\abs\varpi^{\mathsf r^2-\mathsf r-1}$, which tends to $0$ as $\mathsf r\to\infty$; so $\omega'\equiv0$. In particular $\omega$ is continuous, which is the condition for $j=0$. Now let $k\ge1$. Each $\operatorname{dd}_k1_B$ extends continuously to $K^{k+1}$. Indeed, near a tuple of $K^{k+1}$ the partition $I,J$ is constant and the numbers $\abs{t_j}$, $j\in J$, are constant. On that neighborhood the series of Step 2 is a uniformly convergent series of continuous functions, also at tuples with repeated entries. So $\operatorname{dd}_k\omega_{\le\mathsf T}$ extends continuously to a map $\mathsf E_{\mathsf T}$. As distinct tuples are dense, $\abs{\mathsf E_{\mathsf T}-\mathsf E_{\mathsf T'}}\le\max(\varepsilon_{\mathsf T},\varepsilon_{\mathsf T'})$ everywhere. So $\mathsf E_{\mathsf T}$ converges uniformly, $K$ being complete, to a continuous extension of $\operatorname{dd}_k\omega$. For $x\ne0$, $\omega$ is constant near $x$, so the extension vanishes at $(x,\dots,x)$. Tuples in $\bar B_{\mathsf T}^{k+1}$ are separated only by $\bar B_{\mathsf r}$ with $\mathsf r>\mathsf T$, so there $\abs{\operatorname{dd}_k\omega}\le\varepsilon_{\mathsf T}$, and the extension vanishes at $(0,\dots,0)$.

\emph{Step 5 (item (3)).} We have $\abs{t_l}=\abs\varpi^{\mathsf r'}$ for $l\ge1$, and $t_l-t_{l'}=\varpi^{\mathsf r'}(\varpi^l-\varpi^{l'})\ne0$. So $\omega$ takes the value $\mathsf v_{\mathsf r}$ at $t_0$ and $\mathsf v_{\mathsf r'}$ at the others. As $\operatorname{dd}_k$ kills constants, $\operatorname{dd}_k\omega(t)=(\mathsf v_{\mathsf r}-\mathsf v_{\mathsf r'})/\prod_{l\ge1}(t_0-t_l)$, of absolute value $\abs\varpi^{\mathsf r^2-k\mathsf r}$. If $\omega$ were DP at $0$ with radius $R$, then~\eqref{dp:sq:one-variable-bound} would bound $\operatorname{dd}_k\omega$ by a geometric sequence on $\bar B(0,\rho)$, $\rho<R$. But for $\abs\varpi^{\mathsf r_0}\le\rho$, $k\ge\max(1,2\mathsf r_0)$, $\mathsf r=\lfloor k/2\rfloor\ge\mathsf r_0$ and $\mathsf r'=\mathsf r+1$, these points lie in $\bar B(0,\rho)$ and the absolute value is $\abs\varpi^{-\lfloor k^2/4\rfloor}$, which no geometric sequence bounds.
\end{proof}

For a subset $Z$ of a normed space and $g\colon Z\to W$ put $Z^{\{0\}}=Z$, $g^{[0]}=g$ and, recursively,
\begin{gather*}
 Z^{\{j+1\}}=\{(z,v,\tau):\tau\ne0,\ z\in Z^{\{j\}},\ z+\tau v\in Z^{\{j\}}\},\\
 g^{[j+1]}(z,v,\tau)=\tau^{-1}\bigl(g^{[j]}(z+\tau v)-g^{[j]}(z)\bigr).
\end{gather*}
So $g^{[j]}$ is defined on $Z^{\{j\}}$, as a finite combination of values of $g$, and no continuity of $g$ is assumed. By induction on $b$, $(Z^{\{a\}})^{\{b\}}=Z^{\{a+b\}}$ and $(g^{[a]})^{[b]}=g^{[a+b]}$. For $Z$ open and $g$ of class $C^j_{\mathrm{BGN}}$ we have $Z^{\{j\}}\subseteq Z^{[j]}$, and $g^{[j]}$ is the restriction of the map of Definition~\ref{dp:def:bgn}.

\begin{proposition}[divided differences and difference quotients]\label{dp:prop:div-bgn}
\leavevmode
\begin{enumerate}
\item Let $U\subseteq K$ be open, $k\ge0$ and $f\colon U\to W$. Then $f$ is $C^k_{\mathrm{div}}$ if and only if it is $C^k_{\mathrm{BGN}}$, and in that case, on $U^{[k]}$,
\begin{equation}\label{dp:sq:explicit}
 f^{[k]}(z)=\sum_{\mathsf e\in\mathcal I_k}\mathsf p_{\mathsf e}(z)\,\overline{\operatorname{dd}}_{j_{\mathsf e}}f\bigl(\mathsf L_{\mathsf e,0}(z),\dots,\mathsf L_{\mathsf e,j_{\mathsf e}}(z)\bigr),\qquad j_{\mathsf e}\le k,
\end{equation}
a finite sum over a set $\mathcal I_k$, with $\mathsf p_{\mathsf e},\mathsf L_{\mathsf e,i}$ polynomials over $\Z$ in the coordinates of $z$, independent of $f$ and $U$, and $\mathsf L_{\mathsf e,i}(U^{[k]})\subseteq U$. In particular $\omega$ of Example~\textup{\ref{dp:ex:zero-jets}} is $C^\infty_{\mathrm{BGN}}$.
\item Let $Z$ be a subset of a normed space, $g\colon Z\to W$, $\mathsf c(s)=z_0+sv_0$ for $s\in S\subseteq K$ with $\mathsf c(S)\subseteq Z$, and $s_0,\dots,s_j\in S$ pairwise distinct. Put $\zeta^{(0)}(s)=\mathsf c(s)$, $v^{(0)}=v_0$, and for $1\le i\le j$
\[
 \zeta^{(i)}(s)=\bigl(\zeta^{(i-1)}(s_i),\,v^{(i-1)},\,s-s_i\bigr),\qquad v^{(i)}=(0,\dots,0,1).
\]
Then $\zeta^{(j)}(s_0)\in Z^{\{j\}}$ and $\operatorname{dd}_j(g\circ\mathsf c)(s_0,\dots,s_j)=g^{[j]}(\zeta^{(j)}(s_0))$. If $Z$ is open, then $\zeta^{(j)}(s_0)\in Z^{[j]}$ also when the $s_i$ repeat, and its coordinates are polynomial in the $s_i$.
\item Let $f\colon U\to W$, $x_0$, $\mathbf v$ and nodes be as before Theorem~\textup{\ref{dp:thm:divided-planes}}. For $\mathsf s\in K$ let $\zeta'(\mathsf s)$ be the point $\zeta^{(k_2)}(t'_0)$ of item \textup{(2)} for the line $s\mapsto x_0+\mathsf sv_1+sv_2$ and the nodes $t'_0,\dots,t'_{k_2}$; it is affine in $\mathsf s$. Let $\zeta_{\mathbf v}(t;t')$ be the point $\zeta^{(k_1)}(t_0)$ of item \textup{(2)} for the line $\mathsf s\mapsto\zeta'(\mathsf s)$ and the nodes $t_0,\dots,t_{k_1}$. This point is polynomial in $x_0,v_1,v_2$ and the nodes, and it does not depend on $f$. If $x_0+t_iv_1+t'_lv_2\in U$ for all $i,l$, then $\operatorname{dd}_{k_1,k_2}F_{\mathbf v}(t;t')=f^{[k_1+k_2]}(\zeta_{\mathbf v}(t;t'))$. Here $f^{[k]}$ is the map defined before this proposition, with $g=f$ and $Z=U$.
\item Let $W$ be complete and $x_0\in U$. For $\rho>0$ and $k\ge0$ let $\mathcal Z_k(\rho)$ be the set of the points $\zeta_{\mathbf v}(t;t')$ of item \textup{(3)} with $k_1+k_2=k$, $\norm{v_1},\norm{v_2}\le1$ and all nodes in $B(0,\rho)$. Then $f$ is DP at $x_0$ if and only if there are $\rho,\sigma>0$ and $M\ge0$ with $B(x_0,2\rho)\subseteq U$ and $\norm{f^{[k]}(\zeta)}\le M\sigma^{-k}$ for all $k$ and all $\zeta\in\mathcal Z_k(\rho)$. No continuity of $f$ is assumed.
\end{enumerate}
\end{proposition}

\begin{proof}
\emph{Step 1 (item (2)).} Only the last coordinate of $\zeta^{(i)}(s)$ depends on $s$, so $\zeta^{(i)}(s)=\zeta^{(i)}(s')+(s-s')v^{(i)}$, and the same holds for $\zeta^{(0)}$. We show $\zeta^{(i)}(s)\in Z^{\{i\}}$ and $\operatorname{dd}_i(g\circ\mathsf c)(s,s_1,\dots,s_i)=g^{[i]}(\zeta^{(i)}(s))$ for $s\in S\setminus\{s_1,\dots,s_i\}$, by induction on $i$. The case $i=0$ is the definition. Let $i\ge1$. By the induction hypothesis, $\zeta^{(i-1)}(s_i)$ and $\zeta^{(i-1)}(s)$ lie in $Z^{\{i-1\}}$, and $s-s_i\ne0$, so $\zeta^{(i)}(s)\in Z^{\{i\}}$. By~\eqref{dp:sq:dd-recursion}, symmetry and the induction hypothesis at $s$ and at $s_i$, the left side is
\[
 \frac{g^{[i-1]}(\zeta^{(i-1)}(s))-g^{[i-1]}(\zeta^{(i-1)}(s_i))}{s-s_i}=g^{[i]}\bigl(\zeta^{(i-1)}(s_i),v^{(i-1)},s-s_i\bigr).
\]
Take $s=s_0$, $i=j$. If $Z$ is open, the same induction gives $\zeta^{(i)}(s)\in Z^{[i]}$ for all $s\in S$, as $\zeta^{(i-1)}(s_i)$ and $\zeta^{(i-1)}(s)$ lie in $Z^{[i-1]}$. The coordinates are polynomial in the $s_i$ by the recursion.

\emph{Step 2 (item (3)).} In the recursion of item (2), $\mathsf s$ enters only through the coordinate $\zeta^{(0)}(t'_1)=x_0+\mathsf sv_1+t'_1v_2$, or $x_0+\mathsf sv_1+t'_0v_2$ if $k_2=0$. So $\zeta'$ is affine in $\mathsf s$, and the point $\zeta_{\mathbf v}(t;t')$ is polynomial as stated and does not depend on $f$. Now let $x_0+t_iv_1+t'_lv_2\in U$ for all $i,l$, and fix $i$. Item (2) for $g=f$ on the line $s\mapsto x_0+t_iv_1+sv_2$, with $S=\{t'_0,\dots,t'_{k_2}\}$ and nodes $t'$, gives $\zeta'(t_i)\in U^{\{k_2\}}$ and $\operatorname{dd}_{0,k_2}F_{\mathbf v}(t_i;t')=f^{[k_2]}(\zeta'(t_i))$. Item (2) for $g=f^{[k_2]}$ on $Z=U^{\{k_2\}}$, on the line $\zeta'$ with $S=\{t_0,\dots,t_{k_1}\}$ and nodes $t$, gives the claim, since $(f^{[k_2]})^{[k_1]}=f^{[k_1+k_2]}$.

\emph{Step 3 (item (4)).} Item (3) with Theorem~\ref{dp:thm:divided-planes}(1),(2). In item (1) of that theorem, $B(x_0,2\rho)=B(x_0,R/C_2)\subseteq U$, as $C_2\ge1$.

\emph{Step 4 (item (1), the sums).} Assume $f$ is $C^k_{\mathrm{div}}$. By~\eqref{dp:sq:dd-recursion}, symmetry, density of the distinct tuples and continuity, replacing the entries of a tuple one at a time gives, for $j<k$,
\begin{equation}\label{dp:sq:exchange}
 \overline{\operatorname{dd}}_jf(\mathsf x_0,\dots,\mathsf x_j)-\overline{\operatorname{dd}}_jf(\mathsf x'_0,\dots,\mathsf x'_j)=\sum_{i=0}^j(\mathsf x_i-\mathsf x'_i)\,\overline{\operatorname{dd}}_{j+1}f(\mathsf x_0,\dots,\mathsf x_i,\mathsf x'_i,\dots,\mathsf x'_j).
\end{equation}
We prove~\eqref{dp:sq:explicit} by induction on $k$. For $k=0$ take $\mathsf p=1$, $\mathsf L_0(x)=x$. Let $\mathsf g=\mathsf p\cdot\overline{\operatorname{dd}}_jf\circ(\mathsf L_0,\dots,\mathsf L_j)$ be a term of $f^{[k-1]}$ on $Z=U^{[k-1]}$, with $j\le k-1$ and all $\mathsf L_i(Z)\subseteq U$. In $\Z[z,v,\tau]$, $\tau$ divides $\mathsf p(z+\tau v)-\mathsf p(z)$ and $\mathsf L_i(z+\tau v)-\mathsf L_i(z)$; call the quotients $\mathsf p^\flat$ and $\mathsf L_i^\flat$. By~\eqref{dp:sq:exchange}, for $\tau\ne0$,
\begin{multline*}
 \frac{\mathsf g(z+\tau v)-\mathsf g(z)}\tau=\mathsf p^\flat\,\overline{\operatorname{dd}}_jf\bigl(\mathsf L(z+\tau v)\bigr)\\
 +\sum_{i=0}^j\mathsf p(z)\mathsf L_i^\flat\,\overline{\operatorname{dd}}_{j+1}f\bigl(\mathsf L_0(z+\tau v),\dots,\mathsf L_i(z+\tau v),\mathsf L_i(z),\dots,\mathsf L_j(z)\bigr).
\end{multline*}
The right side is defined and continuous on all of $Z^{[1]}$, and it has the form~\eqref{dp:sq:explicit} with $j+1\le k$. So $\mathsf g^{[1]}$ exists and is continuous, and the new $\mathsf L$'s map $Z^{[1]}=U^{[k]}$ into $U$. Summing over the terms gives $f^{[k]}$, so $f$ is $C^k_{\mathrm{BGN}}$.

\emph{Step 5 (item (1), the converse).} Let $f$ be $C^k_{\mathrm{BGN}}$ and $j\le k$. Item (2) with $g=f$ and $\mathsf c(s)=s$ gives a polynomial map $(s_0,\dots,s_j)\mapsto\zeta^{(j)}(s_0)$ from $U^{j+1}$ to $U^{[j]}$. So $f^{[j]}\circ\zeta^{(j)}$ is a continuous extension of $\operatorname{dd}_jf$.
\end{proof}
%% ---- end sec-09-quotients.tex
%% ---- begin sec-10-base-change.tex
\section{Base change}\label{dp:sec:base-change}

Base change does not separate divided-power analyticity from analyticity by stability: both classes are stable, analytic germs with no loss of radius (Theorem~\ref{dp:thm:an-base-change}(1)), DP germs with a possible loss of radius (Theorem~\ref{dp:thm:germ-base-change}(1), Example~\ref{dp:ex:base-radius}). It cannot detect more than AA either: a DP germ whose base change is analytic or AA over $K'$ is AA over $K$, since base change is a particular enlargement of the target (Theorems~\ref{dp:thm:an-base-change}(3) and~\ref{dp:thm:aa-base-change}(3)). Roby's polynomial laws are defined through all base changes~\cite[pp.~219--220]{Ro}, and his $\Gamma$ commutes with base change~\cite[Th.~III.3]{Ro}: the algebra below is his, and the bounds are the estimates we need.

Throughout, $K'$ is a \emph{valued extension} of $K$: a field containing $K$ with an absolute value extending that of $K$. It is nonarchimedean, since integers have absolute value at most $1$, and it is nontrivially valued, so the preceding sections apply over $K'$. A \emph{$K'$-law} is a law over $K'$. A \emph{$K'$-polynomial map} of degree $n$ from $X\otimes_KK'$ to a $K'$-vector space $Q$ is a map that is a homogeneous polynomial of degree $n$ with coefficients in $Q$ on every finite-dimensional $K'$-slice $s\mapsto\sum_is_iz_i$, $s\in K'^r$.

\subsection{Spaces}

\begin{definition}\label{dp:def:base-change}
Let $X$ be a seminormed $K$-space. For $z\in X\otimes_KK'$ put
\[
 \pi(z)=\inf\sum_k\norm{x_k}\abs{\lambda_k},\qquad\pi_{\max}(z)=\inf\max_k\norm{x_k}\abs{\lambda_k},
\]
both infima over the representations $z=\sum_kx_k\otimes\lambda_k$. These are seminorms over $K'$, with scalars acting on the right factor. Also $\pi_{\max}\le\pi$, and $\pi_{\max}$ is ultrametric. The \emph{base change} $X_{K'}$ is the quotient of $X\otimes_KK'$ by the kernel of $\pi$, normed by $\pi$. The space $X^{\max}_{K'}$ is the quotient by the kernel of $\pi_{\max}$, normed by $\pi_{\max}$. In both, $\iota_X$ is the map $x\mapsto[x\otimes1]$. We write $\widehat{X_{K'}}$ and $\widehat{X^{\max}_{K'}}$ for the completions. \emph{With maxima} means that the spaces over $K$ are ultrametric, and that $\pi_{\max}$ and $X^{\max}_{K'}$ replace $\pi$ and $X_{K'}$.
\end{definition}

\begin{lemma}[universal property]\label{dp:bc:universal}
Let $Q$ be a normed $K'$-space and $u\colon X\to Q$ a bounded $K$-linear map. Then $u_{K'}(x\otimes\lambda)=\lambda u(x)$ defines a $K'$-linear map $X_{K'}\to Q$ with $\norm{u_{K'}}=\norm u$, and $u\mapsto u_{K'}$ is a bijection onto the bounded $K'$-linear maps $X_{K'}\to Q$. The same holds with maxima if $Q$ is ultrametric.
\end{lemma}

\begin{proof}
The map $(x,\lambda)\mapsto\lambda u(x)$ is $K$-bilinear. For $z=\sum_kx_k\otimes\lambda_k$ we have $\norm{u_{K'}(z)}\le\norm u\sum_k\norm{x_k}\abs{\lambda_k}$, with a maximum if $Q$ is ultrametric. So $\norm{u_{K'}(z)}\le\norm u\pi(z)$. Conversely $\norm{u(x)}=\norm{u_{K'}(\iota_Xx)}\le\norm{u_{K'}}\norm x$, as $\pi(x\otimes1)\le\norm x$. A $K'$-linear map $S$ is determined by $S\circ\iota_X$, since the vectors $\iota_X(x)$ span. So $S=(S\circ\iota_X)_{K'}$.
\end{proof}

\begin{lemma}[almost orthogonal bases]\label{dp:bc:orthogonal}
Let $K$ be complete, $F\subseteq K'$ a $K$-subspace of finite dimension $d$, and $0<\mathfrak o<1$. Then $F$ has a basis $\omega_1,\dots,\omega_d$ with $\abs{\sum_lc_l\omega_l}\ge\mathfrak o\max_l\abs{c_l}\abs{\omega_l}$ for $c\in K^d$. If $1\in F$, one can take $\omega_1=1$, and then $\psi(\sum_lc_l\omega_l)=c_1$ is $K$-linear with $\psi(1)=1$ and $\abs{\psi(\lambda)}\le\mathfrak o^{-1}\abs\lambda$.
\end{lemma}

\begin{proof}
\emph{Step 1 (one vector).} Let $D\subsetneq F$ be a subspace and $0<\mathfrak o_1<1$. By~\cite[Lemma~A.3]{I}, $D$ is complete, hence closed in $K'$. So $\delta=\operatorname{dist}(\xi,D)>0$ for $\xi\in F\setminus D$. Choose $\mu_0\in D$ with $\abs{\xi-\mu_0}\le\delta/\mathfrak o_1$ and put $\omega=\xi-\mu_0$. For $\mu\in D$ and $c\in K^\times$, $\abs{\mu+c\omega}\ge\abs c\delta\ge\mathfrak o_1\abs{c\omega}$, and then $\abs\mu\le\max(\abs{\mu+c\omega},\abs{c\omega})\le\mathfrak o_1^{-1}\abs{\mu+c\omega}$. So $\abs{\mu+c\omega}\ge\mathfrak o_1\max(\abs\mu,\abs{c\omega})$. This also holds for $c=0$.

\emph{Step 2 (the basis).} Start with $\omega_1=1$ if $1\in F$, and with any nonzero vector otherwise. If $d\ge2$, apply Step~1 with $\mathfrak o_1=\mathfrak o^{1/(d-1)}$ to the span of $\omega_1,\dots,\omega_i$, and induct. The bound on $\psi$ is the case $l=1$ of the inequality.
\end{proof}

\begin{theorem}[$X\to X_{K'}$ is isometric]\label{dp:thm:base-change-isometric}
Let $X$ be a seminormed $K$-space.
\begin{enumerate}
\item For every $K'$, $\pi(x\otimes1)=\norm x$. If $X$ is ultrametric, also $\pi_{\max}(x\otimes1)=\norm x$.
\item \textup{(Completion.)} Let $K'=\Kh$. The map $\mathrm{mult}\colon X\otimes_K\Kh\to\widehat X$, $x\otimes\lambda\mapsto\lambda x$, satisfies $\pi(z)=\norm{\mathrm{mult}(z)}$, and $\pi_{\max}(z)=\norm{\mathrm{mult}(z)}$ if $X$ is ultrametric. So $X_{\Kh}$ is the $\Kh$-span of $X$ in $\widehat X$, and $\widehat{X_{\Kh}}=\widehat X$.
\item \textup{(Tensor powers.)} For $n\ge1$ let $\sigma_n\colon X^{\otimes n}\otimes_KK'\to(X\otimes_KK')^{\otimes_{K'}n}$ be the canonical isomorphism, $\sigma_n((x_1\otimes\dots\otimes x_n)\otimes\lambda)=\lambda(x_1\otimes1)\otimes\dots\otimes(x_n\otimes1)$. It preserves the seminorm $\pi$, where the left side is the base change of $(X^{\otimes n},\pi)$ and the right side is the tensor power of $(X\otimes_KK',\pi)$. The same holds with $\pi_{\max}$ everywhere.
\end{enumerate}
\end{theorem}

\begin{proof}
\emph{Step 1 (item~(1), $K$ complete).} Let $x\otimes1=\sum_kx'_k\otimes\lambda_k$, and let $F$ be the $K$-span of $1$ and the $\lambda_k$. The identity holds in $X\otimes_KF$, which injects into $X\otimes_KK'$. Apply $\id\otimes\psi$, with $\psi$ from Lemma~\ref{dp:bc:orthogonal} for some $0<\mathfrak o<1$. Then $x=\sum_k\psi(\lambda_k)x'_k$, so $\norm x\le\mathfrak o^{-1}\sum_k\abs{\lambda_k}\norm{x'_k}$, with a maximum if $X$ is ultrametric. Take the infimum and let $\mathfrak o\to1$.

\emph{Step 2 (item~(1), general $K$).} The closure of $K$ in $\widehat{K'}$ is a copy of $\Kh$, and $X$ maps isometrically to its separated completion $\widehat X$, a seminormed $\Kh$-space. The map $X\otimes_KK'\to\widehat X\otimes_{\Kh}\widehat{K'}$ carries a representation to one of the same cost. With Step~1 over $\Kh$, this gives $\norm x\le\pi(x\otimes1)\le\norm x$, and the same with maxima.

\emph{Step 3 (item~(2)).} Clearly $\norm{\mathrm{mult}(z)}\le\pi(z)$. Let $z=\sum_{k\le r}x_k\otimes\lambda_k$ and $\varepsilon>0$, and choose $\lambda'_k\in K$ with $\abs{\lambda_k-\lambda'_k}\norm{x_k}\le\varepsilon$. Then $z=(\sum_k\lambda'_kx_k)\otimes1+\sum_kx_k\otimes(\lambda_k-\lambda'_k)$, so $\pi(z)\le\norm{\sum_k\lambda'_kx_k}+r\varepsilon\le\norm{\mathrm{mult}(z)}+2r\varepsilon$, with maxima in the ultrametric case. Let $\varepsilon\to0$.

\emph{Step 4 (item~(3)).} Let $\sum_k\tau_k\otimes\lambda_k$ be a representation, with $\tau_k\in X^{\otimes n}$. Expanding each $\tau_k$ nearly optimally over $K$, and using $\pi(x\otimes1)\le\norm x$, we see that $\sigma_n$ does not increase the seminorm. Conversely, let $z_i=\sum_kx_{i,k}\otimes\lambda_{i,k}$ for $i=1,\dots,n$. Then
\begin{equation}\label{dp:bc:sigma-inverse}
 \sigma_n^{-1}(z_1\otimes\dots\otimes z_n)=\sum_{k_1,\dots,k_n}(x_{1,k_1}\otimes\dots\otimes x_{n,k_n})\otimes\lambda_{1,k_1}\cdots\lambda_{n,k_n},
\end{equation}
of cost at most $\prod_i\sum_k\norm{x_{i,k}}\abs{\lambda_{i,k}}$, with $\prod_i\max_k$ in place of it for $\pi_{\max}$. Take infima over the representations of each $z_i$, and then sum over a representation of a general tensor.
\end{proof}

\begin{theorem}[multilinear maps extend]\label{dp:thm:base-change-multilinear}
Let $\beta\colon P^n\to Q$ be a bounded $K$-multilinear map into a normed $K'$-space $Q$. There is a unique $K'$-multilinear map $\beta^{K'}\colon(P_{K'})^n\to Q$ with $\beta^{K'}(\iota_Px_1,\dots,\iota_Px_n)=\beta(x_1,\dots,x_n)$, and $\norm{\beta^{K'}}=\norm\beta$. The same holds with maxima if $Q$ is ultrametric. For $\beta\colon P^n\to W$ put $\beta_{K'}=(\iota_W\circ\beta)^{K'}\colon(P_{K'})^n\to W_{K'}$; then $\norm{\beta_{K'}}=\norm\beta$.
\end{theorem}

\begin{proof}
For $n=0$, $\beta$ is a vector and $\beta^{K'}=\beta$. Let $n\ge1$. The map $(x_i,\lambda_i)_i\mapsto\lambda_1\cdots\lambda_n\beta(x_1,\dots,x_n)$ is $K$-multilinear, so it defines a $K'$-multilinear map on $(P\otimes_KK')^n$. As in~\eqref{dp:bc:sigma-inverse}, its value at $z_i=\sum_kx_{i,k}\otimes\lambda_{i,k}$ has norm at most $\norm\beta\prod_i\sum_k\norm{x_{i,k}}\abs{\lambda_{i,k}}$, with $\prod_i\max_k$ if $Q$ is ultrametric. So it is at most $\norm\beta\prod_i\pi(z_i)$, and the map passes to the separated quotients. Uniqueness holds since $\iota_P(P)$ spans. The lower bound follows from $\pi(x\otimes1)\le\norm x$. The last claim uses Theorem~\ref{dp:thm:base-change-isometric}(1) for $W$.
\end{proof}

Bounded multilinear maps extend to completions with the same norm; we keep the notation.

\subsection{Laws}

Let $a\colon P\to W$ be a law of degree $n$, with linearization $L_a$. For $n\ge1$ and $z\in P\otimes_KK'$ put $\Theta(z)=\sigma_n^{-1}(z^{\otimes n})$ and $a_{K'}(z)=(L_a\otimes\id)(\Theta(z))$; for $n=0$ put $a_{K'}(z)=a(0)\otimes1$. Let $\norm{a_{K'}}\in[0,\infty]$ be the least $C$ with $\pi(a_{K'}(z))\le C\pi(z)^n$ for all $z$, and for ultrametric $P,W$ let $\norm{a_{K'}}_{\pi_{\max}}$ be the same with $\pi_{\max}$. The \emph{extension constants} are
\[
 \mathfrak e_n(P,K')=\sup\norm{a_{K'}}/\norm a,\qquad\mathfrak e(P,K')=\sup_{n\ge1}\mathfrak e_n(P,K')^{1/n},
\]
the supremum over normed $W$ and nonzero laws $a\colon P\to W$ of degree $n$, and $\mathfrak e_n=1$ if there are none. For ultrametric $P$, the constants $\mathfrak e^{\max}_n(P,K')$ and $\mathfrak e^{\max}(P,K')$ are the same with ultrametric $W$ and $\norm{a_{K'}}_{\pi_{\max}}$.

\begin{proposition}[the extension of a law]\label{dp:prop:law-extension}
\leavevmode
\begin{enumerate}
\item For every representation $z=\sum_{k\le r}x_k\otimes\lambda_k$,
\begin{equation}\label{dp:bc:formula}
 \Theta(z)=\sum_{\abs\alpha=n}O_\alpha(x_1,\dots,x_r)\otimes\lambda^\alpha\in\Gamma^n(P)\otimes_KK',\qquad a_{K'}(z)=\sum_{\abs\alpha=n}a^{[\alpha]}(x_1,\dots,x_r)\otimes\lambda^\alpha.
\end{equation}
The first formula is for $n\ge1$. The map $a_{K'}$ is $K'$-polynomial of degree $n$, and $a_{K'}(x\otimes1)=a(x)\otimes1$.
\item \textup{(Uniqueness.)} Two $K'$-polynomial maps of degree $n$ from $P\otimes_KK'$ to a $K'$-space $Q$ that agree on the vectors $x\otimes1$ are equal.
\item \textup{(Diagonals.)} If $\beta\colon P^n\to Q$ is as in Theorem~\ref{dp:thm:base-change-multilinear} and $\mathrm{diag}(h)=(h,\dots,h)$, then $z\mapsto\beta^{K'}(z,\dots,z)$ is the $K'$-polynomial map that restricts to $x\mapsto\beta(x,\dots,x)$. So $\beta_{K'}\circ\mathrm{diag}=(\beta\circ\mathrm{diag})_{K'}$.
\item $\norm a\le\norm{a_{K'}}$, and $\norm a\le\norm{a_{K'}}_{\pi_{\max}}$ for ultrametric $P,W$. If $P$ has a nonzero law of degree $n$, then $\mathfrak e_n(P,K')=\norm{(\mathsf v_n)_{K'}}$, with $\mathsf v_n$ the universal law of Theorem~\ref{dp:thm:normed-roby}(3).
\end{enumerate}
\end{proposition}

\begin{proof}
\emph{Step 1 (item~(1)).} For $n=0$ the map $a_{K'}$ is constant, and the second formula of~\eqref{dp:bc:formula} is its definition. Let $n\ge1$. Expand $z^{\otimes n}$ over the words in $[r]^n$. By~\eqref{dp:bc:sigma-inverse}, $\sigma_n^{-1}$ sends the term of a word to the tensor of that word in the $x_k$, tensored with the product of the corresponding $\lambda_k$. Grouping the words by content gives the formula for $\Theta(z)$, and $O_\alpha(x)\in\Gamma^n(P)$ by~\eqref{dp:eq:orbit}. Comparing the coefficients of $s^\alpha$ in $L_a$ of~\eqref{dp:eq:orbit} gives $L_a(O_\alpha(x))=a^{[\alpha]}(x)$. On a $K'$-slice, $(\sum_is_iz_i)^{\otimes n}$ is a homogeneous polynomial in $s$ by~\eqref{dp:eq:orbit} over $K'$. Its values lie in $\sigma_n(\Gamma^n(P)\otimes_KK')$ by the formula for $\Theta$, so its coefficients do too, by interpolation over the infinite field $K'$. Apply $(L_a\otimes\id)\circ\sigma_n^{-1}$. The case $r=1$, $\lambda_1=1$ gives the last formula.

\emph{Step 2 (item~(2)).} For fixed $x_1,\dots,x_r\in P$, the difference of the two maps at $\sum_ks_k(x_k\otimes1)$ is a polynomial in $s$ over $K'$ that vanishes on $K^r$. As $K$ is infinite, its coefficients vanish. So it vanishes on $K'^r$. Every $z\in P\otimes_KK'$ has the form $\sum_ks_k(x_k\otimes1)$.

\emph{Step 3 (item~(3)).} The diagonal is homogeneous of degree $n$ on slices, by multilinearity. Apply item~(2), and use $\beta_{K'}=(\iota_W\circ\beta)^{K'}$.

\emph{Step 4 (item~(4)).} By Theorem~\ref{dp:thm:base-change-isometric}(1), $\norm{a(x)}=\pi(a_{K'}(x\otimes1))\le\norm{a_{K'}}\pi(x\otimes1)^n\le\norm{a_{K'}}\norm x^n$, and the same with $\pi_{\max}$. Next, $\mathsf v_n^{[\alpha]}(x)=[O_\alpha(x)]$ by~\eqref{dp:eq:orbit}, so~\eqref{dp:bc:formula} gives $a_{K'}=(L_a\otimes\id)\circ(\mathsf v_n)_{K'}$. A representation shows that $L_a\otimes\id$ has norm at most $\norm{L_a}=\norm a$ for $\pi$. So $\norm{a_{K'}}\le\norm a\norm{(\mathsf v_n)_{K'}}$. Conversely, $\mathsf v_n\ne0$ and $\norm{\mathsf v_n}\le1$ give $\norm{(\mathsf v_n)_{K'}}\le\mathfrak e_n(P,K')$.
\end{proof}

\begin{theorem}[universal bound, all norms]\label{dp:thm:base-change-law}
For every $K'$, all normed $P,W$ and every law $a\colon P\to W$ of degree $n$,
\begin{equation}\label{dp:bc:bound}
 \norm a\le\norm{a_{K'}}\le C_0^n\norm a,\qquad\text{and}\qquad\pi(a_{K'}(z+y)-a_{K'}(z))\le C_0^n\norm a\bigl((\pi(z)+\pi(y))^n-\pi(z)^n\bigr).
\end{equation}
So $a_{K'}$ induces a $K'$-law $P_{K'}\to W_{K'}$ of norm $\norm{a_{K'}}$, which extends to the completions with that norm. Also $\mathfrak e(P,K')\le C_0$, and $\mathfrak e_n(P,\Kh)=1$.
\end{theorem}

\begin{proof}
The case $n=0$ is clear; let $n\ge1$.

\emph{Step 1 (the bounds).} Let $z=\sum_kx_k\otimes\lambda_k$ and $y=\sum_lx'_l\otimes\mu_l$, of costs $\mathsf A$ and $\mathsf B$ for $\pi$. By~\eqref{dp:bc:formula} for the concatenated representation and Lemma~\ref{dp:lem:elementary}(3), $a_{K'}(z+y)-a_{K'}(z)$ is the sum of the $a^{[(\alpha,\alpha')]}(x,x')\otimes\lambda^\alpha\mu^{\alpha'}$ over $\abs\alpha+\abs{\alpha'}=n$ with $\alpha'\ne0$. By Theorem~\ref{dp:thm:mixed}(3) and the multinomial theorem in $\mathbb R$, its cost is at most
\[
 C_0^n\norm a\sum_{\alpha'\ne0}\frac{n!}{\alpha!\alpha'!}\prod_k(\norm{x_k}\abs{\lambda_k})^{\alpha_k}\prod_l(\norm{x'_l}\abs{\mu_l})^{\alpha'_l}=C_0^n\norm a\bigl((\mathsf A+\mathsf B)^n-\mathsf A^n\bigr).
\]
Take infima. With $z=0$ this is the upper bound in~\eqref{dp:bc:bound}. The lower bound is Proposition~\ref{dp:prop:law-extension}(4).

\emph{Step 2 (quotients and completions).} If $\pi(y)=0$, then~\eqref{dp:bc:bound} gives $\pi(a_{K'}(z+y)-a_{K'}(z))=0$. So $a_{K'}$ passes to the separated quotients. Apply~\cite[Lemma~Q.13(1)]{I} over $K'$ for the completions.

\emph{Step 3 (the completion of $K$).} Let $K'=\Kh$. By~\cite[Lemma~Q.13(1)]{I}, $a$ extends to a $\Kh$-law $\hat a\colon\widehat P\to\widehat W$ with $\norm{\hat a}\le\norm a$. Both $\mathrm{mult}\circ a_{\Kh}$ and $\hat a\circ\mathrm{mult}$ are $\Kh$-polynomial and restrict to $a$. So they agree, by Proposition~\ref{dp:prop:law-extension}(2). With Theorem~\ref{dp:thm:base-change-isometric}(2), $\pi(a_{\Kh}(z))=\norm{\hat a(\mathrm{mult}(z))}\le\norm a\pi(z)^n$.
\end{proof}

\begin{theorem}[with maxima]\label{dp:thm:base-change-um}
Let $P,W$ be ultrametric. For every $K'$ and every law $a\colon P\to W$ of degree $n$,
\[
 \norm a\le\norm{a_{K'}}_{\pi_{\max}}\le(\abs\varpi\mathfrak c_K)^{-n}\norm a.
\]
So $a_{K'}$ induces a $K'$-law $P^{\max}_{K'}\to W^{\max}_{K'}$ of norm $\norm{a_{K'}}_{\pi_{\max}}$, which extends to the completions with that norm, and $\mathfrak e^{\max}(P,K')\le(\abs\varpi\mathfrak c_K)^{-1}$. If $K$ is densely valued, then $\norm{a_{K'}}_{\pi_{\max}}=\norm a$ for every $K'$ and every $a$.
\end{theorem}

\begin{proof}
The case $n=0$ is clear. For $n\ge1$, each term of the orbit sum $O_\alpha(x)$ is an elementary tensor of cost $\prod_k\norm{x_k}^{\alpha_k}$. So $\pi_{\max}(O_\alpha(x))\le\prod_k\norm{x_k}^{\alpha_k}$, and~\cite[Theorem~14.39(1), (14.9)]{I} gives $\norm{a^{[\alpha]}(x)}\le(\abs\varpi\mathfrak c_K)^{-n}\norm a\prod_k\norm{x_k}^{\alpha_k}$. As in Step~1 of the proof of Theorem~\ref{dp:thm:base-change-law}, now with maxima, $\pi_{\max}(a_{K'}(z+y)-a_{K'}(z))\le(\abs\varpi\mathfrak c_K)^{-n}\norm a\,\mathsf B\max(\mathsf A,\mathsf B)^{n-1}$, and $\pi_{\max}(a_{K'}(z))\le(\abs\varpi\mathfrak c_K)^{-n}\norm a\,\mathsf A^n$. Take infima. The rest is as in Step~2 there. The lower bound is Proposition~\ref{dp:prop:law-extension}(4). Over a densely valued $K$ we have $\mathfrak c_K=1$, and the bound holds with every $\varpi'\in K$, $0<\abs{\varpi'}<1$, in place of $\varpi$; let $\abs{\varpi'}\to1$.
\end{proof}

\emph{Where ultrametric norms are used.} With sums, the mixed term of index $\alpha$ must be bounded with the weight $n!/\alpha!$, so that the terms sum up, and this is what gives $C_0$. With maxima the multi-indices are not counted, so the constant is $(\abs\varpi\mathfrak c_K)^{-1}$, which is $1$ in the limit over densely valued $K$. Whether the constant for $\pi$ can be taken equal to $1$ over densely valued $K$ is Question~\ref{dp:q:base-change}(1).

For a law $a\colon P\to W$ of degree $n$ put $\operatorname{amb}_{\max}(a)=\sup\{\norm{L_a(t)}/\pi_{\max}(t):t\in\Gamma^n(P)\}$, with the reading of quotients of Section~\ref{dp:sec:conv}. As $\pi_{\max}\le\pi$, we have $\norm a\le\operatorname{amb}(a)\le\operatorname{amb}_{\max}(a)$. If $ja=\beta\circ\mathrm{diag}$, with $j$ a linear isometry into an ultrametric space $Z$ and $\beta$ bounded and $n$-linear, then
\begin{equation}\label{dp:bc:amb-rep}
 \operatorname{amb}_{\max}(a)\le\norm\beta.
\end{equation}
Indeed, the linearization of $\beta$ on $P^{\otimes n}$ is bounded by $\norm\beta$ for $\pi_{\max}$, as $Z$ is ultrametric, and it equals $jL_a$ on $\Gamma^n(P)$. Conversely:

\begin{lemma}[an ultrametric enlargement]\label{dp:bc:envelope}
Let $P,W$ be ultrametric, $\rho>0$, $M\ge0$, and let $a_n\colon P\to W$ be laws of degree $n$ with $\operatorname{amb}_{\max}(a_n)\le M\rho^{-n}$ for all $n\ge0$. There are an ultrametric space $Z$, a linear isometry $j\colon W\to Z$ and $n$-linear maps $\beta_n\colon P^n\to Z$ with $\norm{\beta_n}\le M\rho^{-n}$ and $\beta_n\circ\mathrm{diag}=ja_n$.
\end{lemma}

\begin{proof}
Let $\norm\cdot_Z$ be the seminorm on $W\oplus\bigoplus_nP^{\otimes n}$ given by
\[
 \norm{(v,\tau)}_Z=\inf_t\max\Bigl(\Bigl\lVert v+\sum_nL_{a_n}(t_n)\Bigr\rVert,\ \max_nM\rho^{-n}\pi_{\max}(\tau_n-t_n)\Bigr),
\]
the infimum over $t\in\bigoplus_n\Gamma^n(P)$, and let $Z$ be the quotient by its kernel. Put $j(v)=[v,0]$ and $\beta_n(h_1,\dots,h_n)=[0,h_1\otimes\dots\otimes h_n]$. This is the construction of Step~5 of the proof of~\cite[Theorem~14.39]{I}, with $\rho$ for $\abs\varpi\mathfrak c_KR$ and $M$ for the bound. That step uses only the bound on the $L_{a_n}$ for $\pi_{\max}$.
\end{proof}

\begin{proposition}[non-ambience persists]\label{dp:prop:nonambient-base}
For every $K'$ and every law $a\colon P\to W$ of degree $n$, $\operatorname{amb}(a)\le\operatorname{amb}(a_{K'})$. If $a_{K'}$ is ambient of degree $n$ over $K'$, then $a$ is ambient of degree $n$ over $K$, with a multilinear map of no larger norm. So the laws of Propositions~\ref{dp:prop:nonambient}\textup{(2)} and~\ref{dp:prop:degree-p}\textup{(2)} are not ambient after any base change. With maxima, $\operatorname{amb}_{\max}(a)\le\operatorname{amb}_{\max}(a_{K'})$.
\end{proposition}

\begin{proof}
On $\Gamma^n(P)$ we have $L_{a_{K'}}\circ\iota_P^{\otimes n}=\iota_W\circ L_a$, since both are $K$-linear and agree on powers. A representation of $t$ over $K$ gives one of $\iota_P^{\otimes n}t$ over $K'$ of no larger cost. With Theorem~\ref{dp:thm:base-change-isometric}(1), $\norm{L_a(t)}=\norm{L_{a_{K'}}(\iota_P^{\otimes n}t)}\le\operatorname{amb}(a_{K'})\pi(t)$, and the same with $\pi_{\max}$. If $j'a_{K'}=\beta'\circ\mathrm{diag}$ with a $K'$-linear isometry $j'$, then $j'\circ\iota_W$ is a $K$-linear isometry, and $\beta'\circ\iota_P^{\times n}$, with $\iota_P^{\times n}=(\iota_P,\dots,\iota_P)$, is $K$-multilinear of norm at most $\norm{\beta'}$, with diagonal $j'\iota_Wa$.
\end{proof}

\begin{theorem}[base change computes the ambient constant]\label{dp:thm:base-change-amb}
Let $P,W$ be ultrametric and let $a\colon P\to W$ be a law of degree $n$. Work with maxima.
\begin{enumerate}
\item If $K'$ is densely valued, then $\operatorname{amb}_{\max}(a_{K'})=\norm{a_{K'}}_{\pi_{\max}}$.
\item If $K$ is complete, then $\operatorname{amb}_{\max}(a_{K'})=\operatorname{amb}_{\max}(a)$ for every $K'$.
\item So if $K$ is complete and $K'$ is densely valued, then $\norm{a_{K'}}_{\pi_{\max}}=\operatorname{amb}_{\max}(a)$. If $K$ is densely valued, complete or not, then $\operatorname{amb}_{\max}(a)=\norm a$.
\end{enumerate}
\end{theorem}

\begin{proof}
If $n=0$, then $\operatorname{amb}_{\max}(a)=\operatorname{amb}_{\max}(a_{K'})=\norm{a_{K'}}_{\pi_{\max}}=\norm{a(0)}$, by Theorem~\ref{dp:thm:base-change-isometric}(1). Let $n\ge1$.

\emph{Step 1 (item~(1)).} Apply~\cite[(14.9)]{I} over $K'$ to the law $a_{K'}$ between the ultrametric spaces $P^{\max}_{K'},W^{\max}_{K'}$. Here $\mathfrak c_{K'}=1$, and~\cite[(14.9)]{I} holds for every $\varpi'\in K'$ with $0<\abs{\varpi'}<1$; let $\abs{\varpi'}\to1$. This gives $\operatorname{amb}_{\max}(a_{K'})\le\norm{a_{K'}}_{\pi_{\max}}$. The other inequality holds for every law.

\emph{Step 2 (item~(2), reduction).} One inequality is Proposition~\ref{dp:prop:nonambient-base}. For the other, let $\mathsf t\in\Gamma^n(P^{\max}_{K'})$. Lift it to $\mathsf t_0=\sum_ic_iz_i^{\otimes n}$ with $c_i\in K'$ and $z_i\in P\otimes_KK'$. A representation of $\mathsf t$ lifts at the same cost to a tensor that differs from $\mathsf t_0$ by a sum of elementary tensors with a null factor, and these have $\pi_{\max}=0$. So $\pi_{\max}(\mathsf t_0)\le\pi_{\max}(\mathsf t)$. Put $\mathsf t'=\sum_ic_i\Theta(z_i)\in\Gamma^n(P)\otimes_KK'$. Then $\sigma_n\mathsf t'=\mathsf t_0$, and $L_{a_{K'}}(\mathsf t)$ is the class of $(L_a\otimes\id)(\mathsf t')$.

\emph{Step 3 (item~(2), the expansion).} Here $K$ complete is used. Let $\varepsilon>0$ and $0<\mathfrak o<1$. By Theorem~\ref{dp:thm:base-change-isometric}(3) there is a representation $\mathsf t'=\sum_k\tau_k\otimes\lambda_k$, $\tau_k\in P^{\otimes n}$, with $\max_k\pi_{\max}(\tau_k)\abs{\lambda_k}\le\pi_{\max}(\mathsf t_0)+\varepsilon$. Take a basis $\omega_l$ of the $K$-span of the $\lambda_k$ as in Lemma~\ref{dp:bc:orthogonal}, and write $\lambda_k=\sum_lc_{kl}\omega_l$. Then $\mathsf t'=\sum_lt_l\otimes\omega_l$ with $t_l=\sum_kc_{kl}\tau_k$. As $\mathsf t'\in\Gamma^n(P)\otimes_KK'$ and the $\omega_l$ are independent, $t_l\in\Gamma^n(P)$. Also $\abs{c_{kl}}\abs{\omega_l}\le\mathfrak o^{-1}\abs{\lambda_k}$, so $\pi_{\max}(t_l)\abs{\omega_l}\le\mathfrak o^{-1}(\pi_{\max}(\mathsf t_0)+\varepsilon)$. Therefore
\[
 \pi_{\max}\bigl((L_a\otimes\id)(\mathsf t')\bigr)\le\max_l\norm{L_a(t_l)}\abs{\omega_l}\le\operatorname{amb}_{\max}(a)\,\mathfrak o^{-1}(\pi_{\max}(\mathsf t)+\varepsilon).
\]
Let $\varepsilon\to0$ and $\mathfrak o\to1$.

\emph{Step 4 (item~(3)).} The first claim is items (1) and (2). The second is~\cite[(14.9)]{I} with $\mathfrak c_K=1$ and $\abs{\varpi'}\to1$, for $\varpi'\in K$.
\end{proof}

With maxima the expansion in Step~3 costs nothing. With sums it costs the number of basis vectors, which is unbounded for infinite $K'$. This is the only place where the proof of item~(2) uses ultrametric norms.

\begin{example}[the constant exceeds $1$]\label{dp:ex:base-law}
Let $K=\mathbb Q_p$ with $\abs p=1/p$, $K'=\mathbb Q_{p^2}$, and let $\zeta\in K'$ be a unit whose residue is not in $\F_p$. Let $m\ge1$, $P=c_0([2m],K)$, $W=K$, and
\[
 \phi(s,s')=\frac{s^ps'-ss'^p}p,\qquad b_m(h)=\prod_{i=1}^m\phi(h_{2i-1},h_{2i}),
\]
a law of degree $m(p+1)$. Then $\norm{b_m}=1$ and $\norm{(b_m)_{K'}}_{\pi_{\max}}=\operatorname{amb}_{\max}(b_m)=p^m$. So $\mathfrak e^{\max}_{m(p+1)}(P,K')\ge p^m$.

For $s,s'\in\mathbb Z_p$ we have $s^ps'-ss'^p\equiv0$ modulo $p$. By homogeneity and $\abs{K^\times}=p^{\Z}$, this gives $\abs{\phi(s,s')}\le\max(\abs s,\abs{s'})^{p+1}$. So $\norm{b_m}\le1$. Equality holds at the vector $(1,p,1,p,\dots)$, since $\phi(1,p)=1-p^{p-1}$ is a unit. Next, $\phi$ is the diagonal of a $(p+1)$-linear form of norm at most $p$, so $b_m$ is the diagonal of a form $\Phi_m$ of norm at most $p^m$, and $\operatorname{amb}_{\max}(b_m)\le p^m$ by~\eqref{dp:bc:amb-rep}. Here $W^{\max}_{K'}=K'$, by Lemma~\ref{dp:bc:universal} with $u=\id_K$ and $Q=K'$, and by~\eqref{dp:bc:formula}, $(b_m)_{K'}$ is the same polynomial in the $K'$-coordinates of $z$. With unit vectors $e_i$, the vector $z=\sum_{i\le m}(e_{2i-1}\otimes1+e_{2i}\otimes\zeta)$ has $\pi_{\max}(z)\le1$, and $\abs{\phi(1,\zeta)}=p\abs{\zeta-\zeta^p}=p$, as $\bar\zeta^p\ne\bar\zeta$. So $p^m\le\norm{(b_m)_{K'}}_{\pi_{\max}}\le\operatorname{amb}_{\max}((b_m)_{K'})=\operatorname{amb}_{\max}(b_m)\le p^m$, by Theorem~\ref{dp:thm:base-change-amb}(2).
\end{example}

\subsection{Germs}

Let $U\subseteq P$ be open, and let $f\colon U\to W$ be DP at $x_0$ with coefficient laws $a_n$. We identify the points of $P$ with their images under $\iota_P$. The \emph{base change} $f_{K'}$ is the germ at $x_0$ defined by its series,
\[
 f_{K'}(x_0+z)=\sum_n(a_n)_{K'}(z),\qquad z\in\widehat{P_{K'}},\quad\pi(z)<\bigl(\limsup_n\norm{(a_n)_{K'}}^{1/n}\bigr)^{-1},
\]
with values in $\widehat{W_{K'}}$. By Theorem~\ref{dp:thm:base-change-law} this radius is at least $\rho_{\mathrm{DP}}(f,x_0)/C_0$, and it is $\rho_{\mathrm{DP}}(f_{K'},x_0)$. Statements about $f_{K'}$ concern this series only. A map can stop agreeing with its series outside the ball on which it is represented (Theorem~\ref{dp:thm:quasi-analytic}(3)). So a conclusion about a map $f$ that is DP at $x_0$ with radius $R$ carries the radius $\min(\cdot,R)$.

\begin{theorem}[DP germs]\label{dp:thm:germ-base-change}
Let $f$ be DP at $x_0$ with radius $R$, bound $M$ and coefficient laws $a_n$, and let $\mathfrak e=\mathfrak e(P,K')$.
\begin{enumerate}
\item $f_{K'}$ is DP over $K'$ at $x_0$ with radius $R/\mathfrak e$, bound $M$ and coefficient laws $(a_n)_{K'}$, and $\rho_{\mathrm{DP}}(f,x_0)/\mathfrak e\le\rho_{\mathrm{DP}}(f_{K'},x_0)\le\rho_{\mathrm{DP}}(f,x_0)$.
\item $f_{K'}(x_0+h)=\iota_Wf(x_0+h)$ for $h\in P$, $\norm h<R/\mathfrak e$.
\item \textup{(Uniqueness.)} If $f'$, from an open subset of $\widehat{P_{K'}}$ to $\widehat{W_{K'}}$, is DP over $K'$ at $x_0$ and $f'(x_0+h)=\iota_Wf(x_0+h)$ for small $h\in P$, then $f'=f_{K'}$ near $x_0$, with the same coefficient laws.
\item \textup{(Functoriality.)} If $g$, from an open subset of $W$ to $Y$, is DP at $f(x_0)$, then $(g\circ f)_{K'}=g_{K'}\circ f_{K'}$ near $x_0$.
\item \textup{(Descent.)} Independently of $f$, let $W$ be complete. Let $f'$, from an open subset of $\widehat{P_{K'}}$ to $\widehat{W_{K'}}$, be DP over $K'$ at $x_0\in P$ with radius $R'$ and bound $M'$, with $f'(x_0+h)\in\iota_W(W)$ for small $h\in P$. Then $f'(x_0+h)=\iota_Wf_0(x_0+h)$ on $\norm h<R'$ for a map $f_0$ into $W$ that is DP over $K$ at $x_0$ with radius $R'$ and bound $M'$, and $f'=(f_0)_{K'}$ near $x_0$.
\end{enumerate}
The same holds with maxima, with $\mathfrak e^{\max}(P,K')$.
\end{theorem}

\begin{proof}
\emph{Step 1 (items~(1) and~(2)).} For $n\ge1$, $\norm{(a_n)_{K'}}\le\mathfrak e^n\norm{a_n}\le M(R/\mathfrak e)^{-n}$, and $\norm{(a_0)_{K'}}=\norm{a_0}$. So the series converges for $\pi(z)<R/\mathfrak e$, and this is item~(1) with the lower bound on $\rho_{\mathrm{DP}}$. The upper bound is $\norm{a_n}\le\norm{(a_n)_{K'}}$. Item~(2) is $(a_n)_{K'}(\iota_Ph)=\iota_Wa_n(h)$, since $\mathfrak e\ge1$.

\emph{Step 2 (item~(3)).} Let $a'_n$ be the coefficient laws of $f'$. Let $h\in P$. For $s\in K$ small, $\iota_Wf(x_0+sh)=f'(x_0+sh)=\sum_ns^na'_n(\iota_Ph)$, with geometric bounds. Lemma~\ref{dp:lem:coefficients-unique}(1), for $\iota_W\circ f$, gives $a'_n\circ\iota_P=\iota_W\circ a_n$. By Proposition~\ref{dp:prop:law-extension}(2), $a'_n=(a_n)_{K'}$ on $P_{K'}$, and on the completion by continuity. The two germs agree on the smaller ball.

\emph{Step 3 (item~(4)).} By Theorem~\ref{dp:thm:calculus}(2) over $K'$, $g_{K'}\circ f_{K'}$ is DP over $K'$ at $x_0$. For small $h\in P$, item~(2) gives $g_{K'}(f_{K'}(x_0+h))=g_{K'}(\iota_Wf(x_0+h))=\iota_Yg(f(x_0+h))$. Apply item~(3) to $g\circ f$, which is DP at $x_0$ by the same theorem.

\emph{Step 4 (item~(5)).} Let $a'_n$ be the coefficient laws of $f'$. The maps $a'_n\circ\iota_P$ are $K$-laws into $\widehat{W_{K'}}$ with $\norm{a'_n\circ\iota_P}\le\norm{a'_n}$, as $\pi(\iota_Ph)\le\norm h$. Choose $R''\le R'$ with $f'(x_0+h)\in\iota_W(W)$ for $\norm h<R''$. The map $\iota_W$ is an isometry with closed range, as $W$ is complete. By Theorem~\ref{dp:thm:target-descent}(2) on the ball of radius $R''$, the laws $a'_n\circ\iota_P$ take values in $\iota_W(W)$. So $a_n=\iota_W^{-1}\circ a'_n\circ\iota_P$ are laws with $\norm{a_n}\le M'R'^{-n}$. For $\norm h<R'$ the sum $f'(x_0+h)=\sum_na'_n(\iota_Ph)$ lies in the closed subspace $\iota_W(W)$, so $f_0(x_0+h)=\sum_na_n(h)$ is the required map. Item~(3) gives $f'=(f_0)_{K'}$ near $x_0$.
\end{proof}

\begin{theorem}[analytic germs]\label{dp:thm:an-base-change}
Let $f$ be DP at $x_0$ with radius $R$ and coefficient laws $a_n$.
\begin{enumerate}
\item If $f$ is analytic at $x_0$ with radius $R_1$, bound $M_1$ and multilinear coefficients $\beta_n$, then $f_{K'}$ is analytic over $K'$ at $x_0$ with radius $R_1$, bound $M_1$ and coefficients $(\beta_n)_{K'}$. This holds for every $K$, every $K'$ and all norms.
\item If $f_{K'}$ is analytic over $K'$ at $x_0$ with radius $R_1$ and bound $M_1$, then $\iota_W\circ f\colon U\to\widehat{W_{K'}}$ is analytic over $K$ at $x_0$ with radius $\min(R_1,R)$ and bound $M_1$. Conversely, if $\iota_W\circ f$ is analytic over $K$ at $x_0$ with radius $R_1$ and bound $M_1$, then $f_{K'}$ is analytic over $K'$ at $x_0$ with radius $R_1$ and bound $M_1$.
\item So if $f_{K'}$ is analytic over $K'$ at $x_0$ with radius $R_1$ and bound $M_1$, then $f$ is AA over $K$ at $x_0$ with radius $\min(R_1,R)$ and bound $M_1$, through the isometry $\iota_W$.
\end{enumerate}
The same holds with maxima.
\end{theorem}

\begin{proof}
\emph{Step 1 (item~(1)).} By Lemma~\ref{dp:lem:coefficients-unique}(2), with $j=\id_W$, $a_n=\beta_n\circ\mathrm{diag}$. So the diagonal of $(\beta_n)_{K'}$ is $(a_n)_{K'}$, by Proposition~\ref{dp:prop:law-extension}(3). Also $\norm{(\beta_n)_{K'}}=\norm{\beta_n}\le M_1R_1^{-n}$ by Theorem~\ref{dp:thm:base-change-multilinear}. So the series of $f_{K'}$ converges for $\pi(z)<R_1$ and equals $\sum_n(\beta_n)_{K'}(z,\dots,z)$.

\emph{Step 2 (item~(2)).} Let $\beta'_n$ be the $K'$-multilinear coefficients of $f_{K'}$. By Lemma~\ref{dp:lem:coefficients-unique}(2) over $K'$, the diagonal of $\beta'_n$ is $(a_n)_{K'}$. Put $\beta_n=\beta'_n\circ\iota_P^{\times n}$. This is $K$-multilinear with $\norm{\beta_n}\le\norm{\beta'_n}$ and diagonal $\iota_Wa_n$. So $\iota_Wf(x_0+h)=\sum_n\iota_Wa_n(h)=\sum_n\beta_n(h,\dots,h)$ wherever $f$ equals its series and the coefficients converge, that is, for $\norm h<\min(R_1,R)$.

For the converse, let $\beta_n\colon P^n\to\widehat{W_{K'}}$ be the coefficients of $\iota_W\circ f$. By Lemma~\ref{dp:lem:coefficients-unique}(2), with $j=\iota_W$, their diagonals are $\iota_Wa_n$. By Theorem~\ref{dp:thm:base-change-multilinear} with $Q=\widehat{W_{K'}}$, $\norm{\beta_n^{K'}}\le M_1R_1^{-n}$, and by Proposition~\ref{dp:prop:law-extension}(3) and~(2) the diagonal of $\beta_n^{K'}$ is $(a_n)_{K'}$. So the series of $f_{K'}$ converges for $\pi(z)<R_1$ and equals $\sum_n\beta_n^{K'}(z,\dots,z)$. Since $f_{K'}$ is defined by its series, no condition $R_1\le R$ is needed in this direction.

\emph{Step 3 (item~(3)).} By Theorem~\ref{dp:thm:base-change-isometric}(1), $\iota_W$ is an isometry, and the diagonals in Step~2 lie in $\iota_W(W)$.
\end{proof}

\begin{proposition}[descent of analyticity]\label{dp:prop:an-descent}
Identify $\Kh$ with the closure of $K$ in $\widehat{K'}$. A \emph{bounded retraction} with constant $C$ is a $K$-linear map $\psi\colon K'\to\Kh$ with $\psi(1)=1$ and $\abs{\psi(\lambda)}\le C\abs\lambda$. By continuity, it is the same as such a $\Kh$-linear map on $\widehat{K'}$.
\begin{enumerate}
\item Let $W$ be complete, let $\psi$ be a bounded retraction with constant $C$, and let $f$ be DP at $x_0$ with radius $R$. If $f_{K'}$ is analytic over $K'$ at $x_0$ with radius $R_1$ and bound $M_1$, then $f$ is analytic over $K$ at $x_0$ with radius $\min(R_1,R)$ and bound $CM_1$. The same holds with maxima.
\item If $\Kh$ is spherically complete, for instance if $K$ is discretely valued, there is a bounded retraction with $C=1$, for every $K'$. If $[\widehat{K'}:\Kh]<\infty$, there is one with $C=\mathfrak o^{-1}$ for every $\mathfrak o<1$.
\item If $K$ is complete and not spherically complete, and $K'$ is spherically complete, then there is no bounded retraction.
\end{enumerate}
\end{proposition}

\begin{proof}
\emph{Step 1 (item~(1)).} As $W$ is complete, it is a normed $\Kh$-space. Define $\Psi_0\colon W\otimes_KK'\to W$ by $\Psi_0(w\otimes\lambda)=\psi(\lambda)w$. It is $K$-linear. A representation gives $\norm{\Psi_0(z)}\le C\pi(z)$, and $\norm{\Psi_0(z)}\le C\pi_{\max}(z)$ if $W$ is ultrametric. So $\Psi_0$ extends to a $K$-linear map $\Psi\colon\widehat{W_{K'}}\to W$ with $\norm\Psi\le C$, and $\Psi\circ\iota_W=\id_W$. Let $\beta_n$ be as in Step~2 of the proof of Theorem~\ref{dp:thm:an-base-change}. Then $\Psi\circ\beta_n\colon P^n\to W$ is $K$-multilinear of norm at most $CM_1R_1^{-n}$, and $f(x_0+h)=\Psi(\iota_Wf(x_0+h))=\sum_n\Psi\beta_n(h,\dots,h)$ for $\norm h<\min(R_1,R)$.

\emph{Step 2 (item~(2)).} For the first claim apply~\cite[Theorem~D.3]{I} over $\Kh$ to the space $\widehat{K'}$, the subspace $\Kh$ and the identity of $\Kh$, with constant $1$. If $K$ is discretely valued, $\Kh$ is spherically complete by~\cite[Lemma~D.1]{I}. The second claim is Lemma~\ref{dp:bc:orthogonal} over $\Kh$ with $F=\widehat{K'}$.

\emph{Step 3 (item~(3)).} Suppose that $\psi\colon K'\to K$ is a bounded retraction. Let $\mathrm{pr}\colon\linf(\N,K)\to\mathsf E=\linf(\N,K)/c_0(\N,K)$ be the quotient map, with the quotient norm, and let $\mathbf 1$ be the constant sequence. Sequences in $c_0(\N,K)$ tend to $0$, so $\norm{\mathrm{pr}(\mathbf 1)}=1$, and $c\,\mathrm{pr}(\mathbf 1)\mapsto c$ is a $K$-linear contraction from $K\mathrm{pr}(\mathbf 1)$ to $K'$. It extends to a $K$-linear contraction $\chi\colon\mathsf E\to K'$. By Zorn's lemma it suffices to add one vector $\mathsf x$ to the domain $\mathsf D$ of a contraction $\chi_0$. The closed balls of center $-\chi_0(d)$ and radius $\norm{d+\mathsf x}$, $d\in\mathsf D$, have positive radii, as $c_0(\N,K)$ is closed in $\linf(\N,K)$. They meet pairwise, as $\abs{\chi_0(d)-\chi_0(d')}\le\norm{d-d'}\le\max(\norm{d+\mathsf x},\norm{d'+\mathsf x})$. As $K'$ is spherically complete, they have a common point $\mu$, which gives the contraction $d+c\mathsf x\mapsto\chi_0(d)+c\mu$: for $c\ne0$, $\abs{\chi_0(d)+c\mu}=\abs c\,\abs{\chi_0(c^{-1}d)+\mu}\le\abs c\,\norm{c^{-1}d+\mathsf x}=\norm{d+c\mathsf x}$. Now $\psi\circ\chi\circ\mathrm{pr}$ is a bounded functional on $\linf(\N,K)$ that vanishes on $c_0(\N,K)$ and takes the value $1$ at $\mathbf 1$. This contradicts~\cite[Th\'eor\`eme~4.1]{VdP}, as used in the proof of~\cite[Lemma~F.2]{I}: over such $K$, with unit vectors $e_i$, every bounded functional $\mathsf l$ on $\linf(\N,K)$ is $\mathsf l(x)=\sum_ix_i\mathsf l(e_i)$.
\end{proof}

Which other extensions have a bounded retraction is not determined here. Neither is descent along an extension without one; Question~\ref{dp:q:base-change}(3),(4) are test cases.

\begin{theorem}[ambient analytic germs]\label{dp:thm:aa-base-change}
Let $f$ be DP at $x_0$ with radius $R$ and coefficient laws $a_n$.
\begin{enumerate}
\item If $f$ is AA at $x_0$ with radius $\rho$ and bound $M$, then $f_{\Kh}$ is AA over $\Kh$ with radius $\rho$ and bound $M$.
\item Let $K$ be complete and $P,W$ ultrametric. If $f$ is AA at $x_0$ with radius $\rho$ and bound $M$ through an ultrametric space, then $f_{K'}$, with maxima, is AA over $K'$ with radius $\rho$ and bound $M$ through an ultrametric space, for every $K'$.
\item If $f_{K'}$ is AA over $K'$ with radius $\rho'$ and bound $M'$, then $f$ is AA over $K$ with radius $\min(\rho',R)$ and bound $M'$. So if $f$ is not AA over $K$, then $f_{K'}$ is not AA, and so not analytic, over any $K'$.
\item Let $K$ be complete, $P,W$ ultrametric, and $\rho,M>0$. Consider: \textup{(a)} $f$ is AA at $x_0$ with radius $\rho$ and bound $M$ through an ultrametric space; \textup{(b)} $\operatorname{amb}_{\max}(a_n)\le M\rho^{-n}$ for all $n$; \textup{(c)} $f_{K'}$, with maxima, is DP at $x_0$ with radius $\rho$ and bound $M$ for some densely valued $K'$; \textup{(d)} the same for every densely valued $K'$. Then \textup{(b)}, \textup{(c)} and \textup{(d)} are equivalent; \textup{(a)} implies \textup{(b)}; and \textup{(b)} implies \textup{(a)} with radius $\min(\rho,R)$.
\item If $K$ is densely valued, complete or not, and $P,W$ are ultrametric, and $f$ is DP at $x_0$ with radius $R$ and bound $M$, then $f$ is AA through an ultrametric space with the same radius and bound.
\end{enumerate}
\end{theorem}

\begin{proof}
Let $j\colon W\to Z$ and $\beta_n$ be the ambient data where they occur. By~\cite[Proposition~14.9]{I}, $ja_n=\beta_n\circ\mathrm{diag}$.

\emph{Step 1 (item~(1)).} By Theorem~\ref{dp:thm:base-change-isometric}(2), $\widehat{P_{\Kh}}=\widehat P$ and $\widehat{W_{\Kh}}=\widehat W$. As $\norm{a_n}\le\norm{\beta_n}$ and $\mathfrak e(P,\Kh)=1$, the germ $f_{\Kh}$ is given by its series on the ball of radius $\rho$ in $\widehat P$. Let $\hat\jmath\colon\widehat W\to\widehat Z$ and $\hat\beta_n\colon\widehat P^n\to\widehat Z$ be the extensions by continuity. They are $\Kh$-linear, respectively $\Kh$-multilinear, and $\hat\jmath$ is an isometry with $\norm{\hat\beta_n}=\norm{\beta_n}$. By density, the diagonal of $\hat\beta_n$ is $\hat\jmath\circ(a_n)_{\Kh}$.

\emph{Step 2 (item~(2)).} By~\eqref{dp:bc:amb-rep}, $\operatorname{amb}_{\max}(a_n)\le M\rho^{-n}$. By Theorem~\ref{dp:thm:base-change-amb}(2), $\operatorname{amb}_{\max}((a_n)_{K'})\le M\rho^{-n}$. So $\norm{(a_n)_{K'}}_{\pi_{\max}}\le M\rho^{-n}$. Apply Lemma~\ref{dp:bc:envelope} over $K'$ to the laws $(a_n)_{K'}$ between $P^{\max}_{K'}$ and $W^{\max}_{K'}$, and then complete the space and the maps, as in Step~1. The expansion holds on $\pi_{\max}(z)<\rho$.

\emph{Step 3 (item~(3)).} Let $j',\beta'_n$ be the ambient data of $f_{K'}$. By Lemma~\ref{dp:lem:coefficients-unique}(2) over $K'$, $j'(a_n)_{K'}=\beta'_n\circ\mathrm{diag}$. Then $j'\circ\iota_W$ is a $K$-linear isometry, and $\beta'_n\circ\iota_P^{\times n}$ is $K$-multilinear of norm at most $M'\rho'^{-n}$, with diagonal $j'\iota_Wa_n$. The expansion holds where $f$ equals its series, that is, for $\norm h<\min(\rho',R)$.

\emph{Step 4 (item~(4)).} For (a)$\Rightarrow$(b) use~\eqref{dp:bc:amb-rep}. If (b) holds, Lemma~\ref{dp:bc:envelope} gives ambient data with $\norm{\beta_n}\le M\rho^{-n}$ and $\beta_n\circ\mathrm{diag}=ja_n$. So $jf(x_0+h)=\sum_n\beta_n(h,\dots,h)$ for $\norm h<\min(\rho,R)$, which is (a) with that radius. For a densely valued $K'$, $\norm{(a_n)_{K'}}_{\pi_{\max}}=\operatorname{amb}_{\max}(a_n)$ by Theorem~\ref{dp:thm:base-change-amb}(3). This gives (b)$\Leftrightarrow$(c)$\Leftrightarrow$(d), provided a densely valued $K'$ exists. If $K$ is densely valued, take $K'=K$. If $K$ is discretely valued, take $K'=K(X)$ with the Gauss absolute value $\abs{\sum_ic_iX^i}=\max_i\abs{c_i}\mathsf r^i$, for $\mathsf r>0$ with $\log\mathsf r/\log\abs\vartheta$ irrational. It is multiplicative: for $\mathsf g,\mathsf g'\in K[X]$, if $i_0,i'_0$ are the least indices attaining the maxima for $\mathsf g,\mathsf g'$, then the coefficient of $X^{i_0+i'_0}$ in $\mathsf g\mathsf g'$ has a unique term of largest weight. So it extends to $K(X)$ by $\abs{\mathsf g/\mathsf g'}=\abs{\mathsf g}/\abs{\mathsf g'}$, and its value group $\abs\vartheta^{\Z}\mathsf r^{\Z}$ is dense in $(0,\infty)$.

\emph{Step 5 (item~(5)).} By Theorem~\ref{dp:thm:base-change-amb}(3), $\operatorname{amb}_{\max}(a_n)=\norm{a_n}\le MR^{-n}$. Apply Lemma~\ref{dp:bc:envelope} with $\rho=R$. The expansion holds for $\norm h<R$.
\end{proof}

\begin{example}[the loss of radius is real]\label{dp:ex:base-radius}
Let $K,K'$ be as in Example~\ref{dp:ex:base-law}, and work with maxima. Let $P=c_0(\N,K)$, with $\N$ partitioned into consecutive blocks $J_m$ of size $2m$, and let $f(h)=\sum_{m\ge1}b_m(h|_{J_m})$ on $B(0,1)$, with $b_m$ as in Example~\ref{dp:ex:base-law} in the coordinates of $J_m$. Then $\rho_{\mathrm{DP}}(f,0)=1$, while $\rho_{\mathrm{DP}}(f_{K'},0)=p^{-1/(p+1)}$. Also $f$ is analytic at $0$ with radius $p^{-1/(p+1)}$, and with no larger radius.

The coefficient laws are $a_{m(p+1)}(h)=b_m(h|_{J_m})$, of norm $1$, and $a_n=0$ for the other $n$. So $\rho_{\mathrm{DP}}(f,0)=1$. Each $a_{m(p+1)}$ is the diagonal of the form $\Phi_m$ in the coordinates of $J_m$, of norm at most $p^m$. So $\norm{(a_{m(p+1)})_{K'}}_{\pi_{\max}}\le p^m$, by Theorem~\ref{dp:thm:base-change-multilinear} and Proposition~\ref{dp:prop:law-extension}(3). The vector of Example~\ref{dp:ex:base-law} placed in $J_m$ gives equality. So $\limsup_n\norm{(a_n)_{K'}}_{\pi_{\max}}^{1/n}=p^{1/(p+1)}$. Put $R_0=p^{-1/(p+1)}<1$. The forms satisfy $\norm{\Phi_m}\le R_0^{-m(p+1)}$, so $f$ is analytic at $0$ with radius $R_0$ and bound $1$. If $f$ is analytic at $0$ with radius $R_1$ and bound $M_1$, then Theorem~\ref{dp:thm:an-base-change}(1) gives $\norm{(a_n)_{K'}}_{\pi_{\max}}\le M_1R_1^{-n}$, so $R_1\le R_0$.
\end{example}
%% ---- end sec-10-base-change.tex
%% ---- begin sec-11-agreement.tex
\section{Agreement with analyticity}\label{dp:sec:agreement}

Analytic germs are DP, with the same radius and bound (Theorem~\ref{dp:thm:calculus}(3); the diagonal of a multilinear map $\beta$ has norm at most $\norm\beta$). This section asks when a DP germ is analytic. The answer depends on the setting, so we keep the settings apart. In a fixed degree $n$, the condition $n!\ne0$ in $K$ is exactly the threshold on all normed spaces and on all Banach spaces (Theorem~\ref{dp:thm:degree-threshold}(1)). Between ultrametric spaces it is sufficient over every $K$ and necessary over complete, not spherically complete $K$; over discretely valued $K$ every law with Banach values is analytic (Theorem~\ref{dp:thm:degree-threshold}(2)). Theorem~\ref{dp:thm:field-level} compares this with the functor $\alpha_k$ of~\cite[n\textsuperscript{o}~7.8.1]{Bou}, field by field.

\subsection{Lifts}\label{dp:ssec:lifts}

Let $a\colon P\to W$ be a law of degree $n$. A \emph{lift} of $a$ is a bounded $n$\nobreakdash-linear map $\beta\colon P^n\to W$ with $\beta(h,\dots,h)=a(h)$ for all $h$. As in Definition~\ref{dp:def:strong-norm}(1), $\operatorname{lift}(a)$ is the infimum of $\norm\beta$ over the lifts $\beta$ of $a$, and $\infty$ if there are none. Always $\norm a\le\operatorname{lift}(a)$, and $\operatorname{lift}(\lambda a)=\abs\lambda\operatorname{lift}(a)$ for $\lambda\in K^\times$. For $n\ge1$ the \emph{lifting constant} of $P$ is
\[
 \mathfrak b_n(P)=\sup\{\operatorname{lift}(a)/\norm a\}\in[1,\infty],
\]
over all normed $W$ and all nonzero laws $a\colon P\to W$ of degree $n$, and $\mathfrak b_n(P)=1$ if there are none. Let $\mathfrak b^{\rm u}_n(P)$ be the same with $W$ ultrametric. Put $\mathfrak B_P=\sup_{n\ge1}\mathfrak b_n(P)^{1/n}$ and $\mathfrak B^{\rm u}_P=\sup_{n\ge1}\mathfrak b^{\rm u}_n(P)^{1/n}$.

\begin{lemma}[one point, one lift, everywhere]\label{dp:lem:lift-criterion}
For a law $a\colon P\to W$ of degree $n$ the following are equivalent.
\begin{enumerate}
\item[(i)] $a$ is analytic at some point of $P$.
\item[(ii)] $\operatorname{lift}(a)<\infty$.
\item[(iii)] $a$ has at every point of $P$ a finite analytic expansion, valid on all of $P$.
\end{enumerate}
Moreover $\mathfrak a_n(P)\le\mathfrak b_n(P)$ for $n\ge1$.
\end{lemma}

\begin{proof}
\emph{Step 1 ((i) implies (ii)).} Let $a(x_0+h)=\sum_l\beta_l(h,\dots,h)$ near $h=0$, with geometric bounds. By Theorem~\ref{dp:thm:mixed}(2), $a(x_0+h)=\sum_{i\le n}a^{[(n-i,i)]}(x_0,h)$ is the DP expansion of $a$ at $x_0$, and its term of degree $n$ is $a^{[(0,n)]}(x_0,h)=a(h)$. By Lemma~\ref{dp:lem:coefficients-unique}(2) with $j=\id_W$, $\beta_n(h,\dots,h)=a(h)$. So $\beta_n$ is a lift.

\emph{Step 2 ((ii) implies (iii), (iii) implies (i)).} Let $\beta$ be a lift and $x_0\in P$. Expanding $\beta(x_0+h,\dots,x_0+h)$ by multilinearity and grouping the $\binom ni$ terms with $h$ in $i$ slots gives bounded $i$\nobreakdash-linear maps $\beta_i$ with $a(x_0+h)=\sum_{i\le n}\beta_i(h,\dots,h)$ for all $h$. The last implication is trivial.

\emph{Step 3 (the ambience constant).} We may assume $\mathfrak b_n(P)<\infty$ and $\mathfrak a_n(P)>0$. Then the universal law $\mathsf v_n\colon P\to\Gamma^n_{\mathrm s}(P)$ of Theorem~\ref{dp:thm:normed-roby}(3) is nonzero, since $\norm\cdot_\Gamma$ does not vanish on $\Gamma^n(P)$, and of norm at most $1$. Let $\varepsilon>0$ and let $\beta$ be a lift of $\mathsf v_n$ with $\norm\beta\le\mathfrak b_n(P)+\varepsilon$. The linear map $\mathsf L$ on $P^{\otimes n}$ with $\mathsf L(h_1\otimes\dots\otimes h_n)=\beta(h_1,\dots,h_n)$ has $\norm{\mathsf L(\mathsf g)}\le\norm\beta\pi(\mathsf g)$ for $\mathsf g\in P^{\otimes n}$. On $\Gamma^n(P)$ it is the quotient map, as the two agree on powers. So $\norm{\mathsf g}_\Gamma\le\norm\beta\pi(\mathsf g)$ for $\mathsf g\in\Gamma^n(P)$, and $\mathfrak a_n(P)\le\mathfrak b_n(P)+\varepsilon$ for every $\varepsilon>0$.
\end{proof}

\begin{theorem}[agreement criterion]\label{dp:thm:agreement-criterion}
The following are equivalent.
\begin{enumerate}
\item For every normed $W$, every open $U\subseteq P$ and every $f\colon U\to W$: if $f$ is DP at $x_0\in U$, then $f$ is analytic at $x_0$.
\item $\mathfrak B_P<\infty$.
\end{enumerate}
The same holds with $W$ ultrametric in~\textup{(1)} and $\mathfrak B^{\rm u}_P$ in~\textup{(2)}. If $\mathfrak B_P<\infty$ and $f$ is DP at $x_0$ with radius $R$ and bound $M$, then $f$ is analytic at $x_0$ with radius $R/\mathfrak B_P$ and bound $2M$.
\end{theorem}

\begin{proof}
\emph{Step 1 ((2) implies (1), with the constants).} Let $a_n$ be the coefficient laws. For $n\ge1$ and $a_n\ne0$, choose a lift $\beta_n$ with $\norm{\beta_n}\le2\operatorname{lift}(a_n)\le2\mathfrak B_P^nMR^{-n}$. For $a_n=0$ take $\beta_n=0$, and put $\beta_0=a_0$. As $\mathfrak B_P\ge1$, the series $\sum_n\beta_n(h,\dots,h)$ equals $f(x_0+h)$ for $\norm h<R/\mathfrak B_P$.

\emph{Step 2 (a sum of spaces).} For the converse let $W_l$, $l\ge1$, be normed spaces, ultrametric in the second statement, and let $W$ be their $c_0$-sum: the sequences $(w_l)_l$ with $w_l\in W_l$ and $\norm{w_l}\to0$, with the supremum norm. It is ultrametric if all $W_l$ are. Let $\iota_l\colon W_l\to W$ and $\operatorname{pr}_l\colon W\to W_l$ be the inclusions and projections. They have norm at most $1$.

\emph{Step 3 (some $\mathfrak b_n(P)=\infty$).} In the second statement read $\mathfrak b^{\rm u}_n$ for $\mathfrak b_n$ here and in Step~4. Put $\mathsf z_l=\varpi^l$. The ratio $\operatorname{lift}/\norm\cdot$ is unbounded and invariant under scalars. So there are laws $\mathsf a_l\colon P\to W_l$ of degree $n$ with $\abs\varpi<\norm{\mathsf a_l}\le1$ and $\operatorname{lift}(\mathsf a_l)\ge l/\abs{\mathsf z_l}$. Put $\mathsf A(h)=(\mathsf z_l\mathsf a_l(h))_l$. For $h_1,\dots,h_r\in P$ we have $\mathsf A(\sum_is_ih_i)=\sum_{\abs\gamma=n}s^\gamma(\mathsf z_l\mathsf a_l^{[\gamma]}(h_1,\dots,h_r))_l$. The coefficient sequences lie in $W$, because $\norm{\mathsf a_l^{[\gamma]}(h_1,\dots,h_r)}\le(8/\abs\varpi)^n(\sum_i\norm{h_i})^n$ by Theorem~\ref{dp:thm:mixed}(1) and $\mathsf z_l\to0$. Also $\norm{\mathsf A(h)}\le\abs\varpi\norm h^n$. So $\mathsf A$ is a law of degree $n$, and it is DP on $P$. A lift $\beta$ of $\mathsf A$ gives the lifts $\operatorname{pr}_l\beta$ of $\mathsf z_l\mathsf a_l$, so $\norm\beta\ge\abs{\mathsf z_l}\operatorname{lift}(\mathsf a_l)\ge l$ for all $l$. Hence $\mathsf A$ has no lift, and by Lemma~\ref{dp:lem:lift-criterion} it is analytic at no point.

\emph{Step 4 (all $\mathfrak b_n(P)$ finite).} If the supremum is infinite while all terms are finite, then there are $n_1<n_2<\dots$ with $\mathfrak b_{n_l}(P)>(2l)^{n_l}$. Choose laws $\mathsf a_l\colon P\to W_l$ of degree $n_l$ with $\abs\varpi<\norm{\mathsf a_l}\le1$ and $\operatorname{lift}(\mathsf a_l)>(2l)^{n_l}\abs\varpi$. For $\norm h<1$ put $f(h)=(\mathsf a_l(h))_l$. It lies in $W$, as $\norm{\mathsf a_l(h)}\le\norm h^{n_l}\to0$, and $f=\sum_l\iota_l\mathsf a_l$ converges in $W$. So $f$ is DP at $0$ with radius $1$ and bound $1$. Suppose $f(h)=\sum_n\beta_n(h,\dots,h)$ for $\norm h<R'$ with $\norm{\beta_n}\le M'R'^{-n}$. By Lemma~\ref{dp:lem:coefficients-unique}(2), $\operatorname{pr}_l\beta_{n_l}$ is a lift of $\mathsf a_l$. So $(2l)^{n_l}\abs\varpi<M'R'^{-n_l}$, that is, $(2lR')^{n_l}<M'/\abs\varpi$. For $l>1/R'$ the left side exceeds $2^{n_l}$, and $n_l\ge l$. This fails for large $l$, so $f$ is not analytic at $0$.
\end{proof}

We recall from Section~\ref{dp:sec:laws} the \emph{scaling loss} $\mathfrak s(P)\in[\abs\varpi,1]$: the supremum of the $\mathsf c\in(0,1]$ such that every nonzero $h\in P$ has a multiple $\lambda h$, $\lambda\in K$, with $\mathsf c\le\norm{\lambda h}\le1$. Every $\mathsf c<\mathfrak s(P)$ has this property. If $\operatorname{char}K=0$, put $\mathfrak n_K=\inf_{n\ge1}\abs{n!}^{1/n}$. Then $\mathfrak n_K>0$. Indeed the integers $n$ with $\abs n<1$ form a prime ideal of $\Z$, by the ultrametric inequality. If it is $0$, then $\abs{n!}=1$ for all $n$. If it is $p\Z$, then $\abs{n!}=\abs p^{\ord_p(n!)}$ with $\ord_p(n!)\le n/(p-1)$, so $\mathfrak n_K\ge\abs p^{1/(p-1)}$.

\begin{theorem}[characteristic $0$ on ultrametric sources]\label{dp:thm:agreement-char0}
Let $P$ be ultrametric.
\begin{enumerate}
\item Let $n\ge1$ with $n!\ne0$ in $K$, and let $a\colon P\to W$ be a law of degree $n$. Then $\beta_a=a^{[1^n]}/n!$ is a lift of $a$, and
\begin{equation}\label{dp:agr:eq:char0}
 \norm{\beta_a}\le\abs{n!}^{-1}\bigl(2/\mathfrak s(P)\bigr)^n\norm a,\quad\text{and}\quad\norm{\beta_a}\le\abs{n!}^{-1}\mathfrak s(P)^{-n}\norm a\text{ if $W$ is ultrametric.}
\end{equation}
\item Let $\operatorname{char}K=0$. A map that is DP at $x_0$ with radius $R$ and bound $M$ is analytic at $x_0$ with radius $\mathfrak n_K\mathfrak s(P)R/2$ and bound $M$, and with radius $\mathfrak n_K\mathfrak s(P)R$ if $W$ is ultrametric. If moreover the residue characteristic is $0$, $W$ is ultrametric and $\mathfrak s(P)=1$, the two radii are equal.
\end{enumerate}
\end{theorem}

\begin{proof}
\emph{Step 1 (the lift).} The map $a^{[1^n]}$ is $n$\nobreakdash-linear by Lemma~\ref{dp:lem:elementary}(4), and $a^{[1^n]}(h,\dots,h)=n!\,a(h)$ by Lemma~\ref{dp:lem:hasse-law}(5).

\emph{Step 2 (the bound).} By Lemma~\ref{dp:lem:polarization}, $\norm{\beta_a(h_1,\dots,h_n)}\le\abs{n!}^{-1}2^n\norm a\max_i\norm{h_i}^n$, without $2^n$ if $W$ is ultrametric. Let $h_1,\dots,h_n$ be nonzero and $\mathsf c<\mathfrak s(P)$. Choose $\lambda_i$ with $\mathsf c\le\norm{\lambda_ih_i}\le1$. By multilinearity, $\norm{\beta_a(h_1,\dots,h_n)}=\prod_i\abs{\lambda_i}^{-1}\norm{\beta_a(\lambda_1h_1,\dots,\lambda_nh_n)}\le\abs{n!}^{-1}2^n\norm a\prod_i(\norm{h_i}/\mathsf c)$. Let $\mathsf c\to\mathfrak s(P)$.

\emph{Step 3 (germs).} Let $a_n$ be the coefficient laws and $\beta_n=\beta_{a_n}$ for $n\ge1$, $\beta_0=a_0$. As $\abs{n!}^{-1}\le\mathfrak n_K^{-n}$, \eqref{dp:agr:eq:char0} gives the radii, with bound $M$. In residue characteristic $0$ we have $\abs{n!}=1$.
\end{proof}

Compare \cite[Theorem~14.39(3) and Proposition~14.40(2)]{I}, with a larger constant.

\begin{theorem}[finite-dimensional sources]\label{dp:thm:agreement-finite}
\leavevmode
\begin{enumerate}
\item Let $T\colon K^r\to P$ be a bounded linear bijection with bounded inverse, $K^r$ with the maximum norm, and $P\ne\{0\}$. In every characteristic, $\mathfrak b_n(P)\le2^{r-1}(16r\norm T\norm{T^{-1}}/\abs\varpi)^n$. A map from an open subset of $P$ to a normed $W$ that is DP at $x_0$ with radius $R$ and bound $M$ is analytic there with radius $\abs\varpi R/(16r\norm T\norm{T^{-1}})$ and bound $2^{r-1}M$. If $K$ is complete, every finite-dimensional normed space has such a $T$~\cite[Lemma~A.3]{I}.
\item Over every $K$, let $P\ne\{0\}$ be finite-dimensional and $W$ Banach. There are an integer $\mathsf r\ge1$ and a real $C_P\ge1$, depending only on $P$ and $\varpi$, such that every law $a\colon P\to W$ of degree $n\ge1$ has a lift of norm at most $2^{\mathsf r}C_P^n\norm a$. A map from an open subset of $P$ to $W$ that is DP at $x_0$ with radius $R$ and bound $M$ is analytic there with radius $R/C_P$ and bound $2^{\mathsf r}M$.
\item Completeness of $W$ cannot be dropped in~\textup{(2)}. Let $K=\F_p(t)$ with the $t$-adic absolute value, let $k\ge p$, and let $(E_k,G)$ and $(D_k,G)$ be the pairs of finite-dimensional ultrametric spaces of~\cite[Theorem~H.4]{I}. The law $A^k_{E_k,D_k;G}$ of degree $k$, and the joint action $B$ of Theorem~\ref{dp:thm:calculus}(4) on $\cL(E_k,D_k)\times\cL(G,G)$, of degree $k+1$, are DP everywhere and analytic at no point. Their sources and targets are finite-dimensional and ultrametric.
\end{enumerate}
\end{theorem}

\begin{proof}
\emph{Step 1 (item (1)).} The bound on $\mathfrak b_n(P)$ is Step~2 of the proof of Theorem~\ref{dp:thm:calculus}, which constructs for each law $a$ a lift $\beta_a$ of norm at most $2^{r-1}(16r\norm T\norm{T^{-1}}/\abs\varpi)^n\norm a$. If $\norm{a_n}\le MR^{-n}$, then $\norm{\beta_{a_n}}\le2^{r-1}MR'^{-n}$ for $R'=\abs\varpi R/(16r\norm T\norm{T^{-1}})$, and $R'\le R$.

\emph{Step 2 (item (2): the completion).} The space $W$ is a normed space over $\Kh$, by continuity. By \cite[Lemma~Q.13(1),(2)]{I}, $a$ extends to a law $\hat a\colon\widehat P\to W$ of degree $n$ over $\Kh$ with $\norm{\hat a}\le\norm a$, and $\widehat P$ has finite dimension $\mathsf r\ge1$ over $\Kh$. Choose a $\Kh$\nobreakdash-basis $\mathsf l_1,\dots,\mathsf l_{\mathsf r}$ of $\widehat P$ with $\norm{\mathsf l_c}\le1$. By \cite[Lemma~A.3]{I} over $\Kh$, its coordinate forms $\mathsf l_c^*$ are bounded. Put $\mathsf m_P=\max_c\norm{\mathsf l_c^*}\ge1$ and $C_P=16\mathsf r\mathsf m_P/\abs\varpi$.

\emph{Step 3 (item (2): the lift).} By Theorem~\ref{dp:thm:mixed}(1) over $\Kh$, $\norm{\hat a^{[\gamma]}(\mathsf l_1,\dots,\mathsf l_{\mathsf r})}\le(8\mathsf r/\abs\varpi)^n\norm a$ for $\abs\gamma=n$. Put
\[
 \hat\beta(h_1,\dots,h_n)=\sum_{\mathsf w}\mathsf l^*_{\mathsf w_1}(h_1)\cdots\mathsf l^*_{\mathsf w_n}(h_n)\,\hat a^{[\gamma(\mathsf w)]}(\mathsf l_1,\dots,\mathsf l_{\mathsf r}),
\]
over the nondecreasing words $\mathsf w$ in $\{1,\dots,\mathsf r\}$ of length $n$, where $\gamma(\mathsf w)$ counts the letters of $\mathsf w$. Each monomial comes from exactly one such word, so $\hat\beta$ has diagonal $\hat a$. There are at most $2^{n+\mathsf r-1}$ words, so $\norm{\hat\beta}\le2^{\mathsf r}C_P^n\norm a$. The restriction of $\hat\beta$ to $P^n$ is a lift of $a$. The statement on germs follows as in Step~1.

\emph{Step 4 (item (3)).} Both maps are laws, DP on their whole domains, by Theorem~\ref{dp:thm:calculus}(4),(3),(1). The proof of \cite[Theorem~H.4]{I} fixes the morphism of $G$ to be $\id_G$, so it shows that $A^k_{E_k,D_k;G}$ is not analytic. As $A^k_{E_k,D_k;G}(u)=B(u,\id_G)$, the map $B$ is not analytic either. By Lemma~\ref{dp:lem:lift-criterion} neither is analytic at any point. The spaces $E_k,D_k,G$ are finite-dimensional, and operator norms into ultrametric spaces are ultrametric, so the sources and targets are as claimed. The sources are nonzero. So, by item~(2), the targets are not complete.
\end{proof}

\begin{theorem}[fixed degree]\label{dp:thm:degree-threshold}
Let $n\ge1$.
\begin{enumerate}
\item The following are equivalent: \textup{(a)} $n!\ne0$ in $K$, that is, $\operatorname{char}K=0$ or $\operatorname{char}K>n$; \textup{(b)} for all normed $P,W$, every law $P\to W$ of degree $n$ is analytic; \textup{(c)} for all Banach $P,W$, every law $P\to W$ of degree $n$ is analytic. Under~\textup{(a)}, $\operatorname{lift}(a)\le\abs{n!}^{-1}(8n/\abs\varpi^2)^n\norm a$ for every $P$, and $\operatorname{lift}(a)\le\abs{n!}^{-1}(2/\mathfrak s(P))^n\norm a$ for ultrametric $P$; here $\abs{n!}=1$ if $\operatorname{char}K>n$.
\item Between ultrametric spaces the condition is as follows. Under~\textup{(a)}, every law of degree $n$ is analytic, by~\textup{(1)}. Let $\operatorname{char}K=p\le n$.
\begin{enumerate}
\item If $K$ is complete and not spherically complete, there are ultrametric Banach spaces $P,W$ and a law $P\to W$ of degree $n$ that is analytic at no point.
\item If $K$ is discretely valued, every law of degree $n$ from an ultrametric $P$ to an ultrametric Banach space $W$ is analytic.
\item If $K$ is densely valued and either spherically complete or incomplete, every law of degree $n$ from an ultrametric $P$ to a space $W$ with an equivalent ultrametric spherically complete norm is analytic. For other ultrametric Banach $W$, if $\Kh$ is spherically complete, we do not decide this in general (Question~\ref{dp:q:residual}); if $K$ is incomplete and $\Kh$ is not spherically complete, some law of degree $n$ between ultrametric Banach spaces is analytic at no point (the paragraph on incomplete fields after that question). Let $K$ be spherically complete and $E,D,F$ ultrametric Banach. The law $A^n_{E,D;F}$ has a lift if $F$ has an equivalent ultrametric spherically complete norm, or if $E$ or $D$ is a bounded linear retract of some $c_0(J,K)$, for instance of countable type. In general we do not decide this for $A^n_{E,D;F}$ either.
\end{enumerate}
Let $W$ be ultrametric and incomplete. If $P$ is as in Theorem~\ref{dp:thm:agreement-finite}(1), or $P$ is the space $c_{00}(J)$ of Section~\ref{dp:ssec:agree-ultrametric}, every law of degree $n$ from $P$ to $W$ is analytic. Over $K=\F_p(t)$, some law of degree $n$ between finite-dimensional ultrametric spaces, with incomplete target, is analytic at no point. Otherwise we do not decide this in general (Question~\ref{dp:q:residual}).
\end{enumerate}
\end{theorem}

\begin{proof}
\emph{Step 1 ((a) implies (b), with the bounds).} As in Step~1 of the proof of Theorem~\ref{dp:thm:agreement-char0}, which does not use that $P$ is ultrametric, $\beta_a=a^{[1^n]}/n!$ is $n$\nobreakdash-linear with diagonal $a$. For nonzero $h_1,\dots,h_n$ choose $\lambda_i=\varpi^{l_i}$ with $\abs\varpi/n<\norm{\lambda_ih_i}\le1/n$. Then $\sum_i\norm{\lambda_ih_i}\le1$. So Theorem~\ref{dp:thm:mixed}(1) gives $\norm{a^{[1^n]}(\lambda_1h_1,\dots,\lambda_nh_n)}\le(8/\abs\varpi)^n\norm a$, and undoing the scaling gives $\norm{a^{[1^n]}(h_1,\dots,h_n)}\le(8n/\abs\varpi^2)^n\norm a\prod_i\norm{h_i}$. For ultrametric $P$ use~\eqref{dp:agr:eq:char0}. If $\operatorname{char}K=p>n$, then $n!$ is a nonzero element of $\F_p$, a root of unity, so $\abs{n!}=1$. A law with a lift is analytic by Lemma~\ref{dp:lem:lift-criterion}. Item~(b) implies item~(c) trivially.

\emph{Step 2 ((c) implies (a)).} Let $\operatorname{char}K=p\le n$. By \cite[Theorem~7.2(1)]{I} with $k=n$, there are Banach spaces $E,F$ such that $A^n_{E,E;F}$ is analytic at no point. It is a law of degree $n$ (Theorem~\ref{dp:thm:calculus}(4),(3)) between Banach spaces.

\emph{Step 3 (item (2)).} For (2)(a), take $A^n_{E,D;K}$ of Proposition~\ref{dp:prop:strict}(1)(c). Items (2)(b) and (2)(c) are Theorem~\ref{dp:thm:sph-agreement}(2),(1). A spherically complete $K$ is complete. So, for the laws $A^n_{E,D;F}$, \cite[Theorem~5.2]{I} and the sentence after it give a lift if $F$ has an equivalent ultrametric spherically complete norm, and \cite[Theorem~5.4]{I} gives one if $E$ or $D$ is a retract of some $c_0(J,K)$ or of countable type. In general the question is the one in \cite[Section~12]{I}. For $W$ incomplete, the sentences of item~(2) are Theorems~\ref{dp:thm:agreement-finite}(1), \ref{dp:thm:c0-agreement}(1) and~\ref{dp:thm:agreement-finite}(3).
\end{proof}

\subsection{Between ultrametric spaces}\label{dp:ssec:agree-ultrametric}

A normed space $P$ is of \emph{countable type} if it has a dense subspace of countable dimension. Let $c_{00}(J)$ be the finitely supported vectors of $c_0(J)=c_0(J,K)$, with the supremum norm and unit vectors $e_{\mathsf j}$, $\mathsf j\in J$.

\begin{lemma}[bases]\label{dp:lem:basis}
Let $K$ be complete and $P$ an infinite-dimensional ultrametric Banach space of countable type. Let $0<\mathfrak o\le1$, with $\mathfrak o<1$ unless $K$ is discretely valued. Put $\mathfrak o'=\mathfrak c_K$ if $K$ is discretely valued and $\mathfrak o'=\mathfrak o$ otherwise. Then $P$ has a sequence $\mathsf l_1,\mathsf l_2,\dots$ with $\mathfrak o'<\norm{\mathsf l_i}\le1$ and
\begin{equation}\label{dp:agr:eq:orthogonal}
 \Bigl\lVert\sum_i\lambda_i\mathsf l_i\Bigr\rVert\ge\mathfrak o\max_i\abs{\lambda_i}\norm{\mathsf l_i}
\end{equation}
for finitely supported $(\lambda_i)_i$, such that $\Xi(\lambda)=\sum_i\lambda_i\mathsf l_i$ is a linear bijection $c_0(\N,K)\to P$ with $\norm\Xi\le1$ and $\norm{\Xi^{-1}}\le(\mathfrak o\mathfrak o')^{-1}$.
\end{lemma}

\begin{proof}
\emph{Step 1 (a dense sequence).} Let $\mathsf D$ be a dense subspace of $P$ of countable dimension. A finite-dimensional subspace of $P$ is complete~\cite[Lemma~A.3]{I}, hence closed, and so not dense. So $\mathsf D$ is infinite-dimensional. Let $\mathsf x_1,\mathsf x_2,\dots$ be a basis of $\mathsf D$ and $P_m=\spn(\mathsf x_1,\dots,\mathsf x_m)$.

\emph{Step 2 (induction).} Choose $\mathsf t_m\in(0,1]$ with $\prod_m\mathsf t_m\ge\mathfrak o$: $\mathsf t_m=1$ if $\mathfrak o=1$, and $\mathsf t_m=\mathfrak o^{2^{-m}}$ otherwise. Put $\tau_1=1$ and $\tau_{m+1}=\mathsf t_m\tau_m$, so $\tau_m\ge\mathfrak o$. Suppose that $\mathsf l_1,\dots,\mathsf l_m$ is a basis of $P_m$ with $\norm{\sum_{i\le m}\lambda_i\mathsf l_i}\ge\tau_m\max_{i\le m}\abs{\lambda_i}\norm{\mathsf l_i}$; for $m=1$ take $\mathsf l_1=\mathsf x_1$. As $P_m$ is closed, $\mathsf d=\operatorname{dist}(\mathsf x_{m+1},P_m)>0$. Choose $y\in P_m$ with $\norm{\mathsf x_{m+1}-y}\le\mathsf t_m^{-1}\mathsf d$. This is possible for $\mathsf t_m<1$ because $\mathsf d>0$. For $\mathsf t_m=1$, $K$ is discretely valued and the distance is attained, by Step~4 of the proof of \cite[Lemma~Q.13]{I}, applied in $P_{m+1}$ with the subspace $P_m$. Put $\mathsf l_{m+1}=\mathsf x_{m+1}-y$. Then $\operatorname{dist}(\mathsf l_{m+1},P_m)=\mathsf d\ge\mathsf t_m\norm{\mathsf l_{m+1}}$. For $z=z'+\lambda\mathsf l_{m+1}$ with $z'\in P_m$, we get $\norm z\ge\mathsf t_m\abs\lambda\norm{\mathsf l_{m+1}}$, and $\norm{z'}\le\max(\norm z,\abs\lambda\norm{\mathsf l_{m+1}})\le\mathsf t_m^{-1}\norm z$. With the induction hypothesis for $z'$ and $\tau_m\le1$, this gives the inequality for $m+1$.

\emph{Step 3 (scaling).} Scaling each $\mathsf l_i$ keeps the inequality, and it gives~\eqref{dp:agr:eq:orthogonal}, as $\tau_m\ge\mathfrak o$. Scale so that $\mathfrak o'<\norm{\mathsf l_i}\le1$: by powers of $\vartheta$ if $K$ is discretely valued. Otherwise $\abs{K^\times}$ is dense in $(0,\infty)$, as in Step~6 of the proof of \cite[Lemma~Q.13]{I}, so it meets $(\mathfrak o/\norm{\mathsf l_i},1/\norm{\mathsf l_i}]$.

\emph{Step 4 (the bijection).} For finitely supported $\lambda$ we have $\mathfrak o\mathfrak o'\norm\lambda\le\mathfrak o\max_i\abs{\lambda_i}\norm{\mathsf l_i}\le\norm{\Xi\lambda}\le\max_i\abs{\lambda_i}\norm{\mathsf l_i}\le\norm\lambda$. For $\lambda\in c_0(\N,K)$ the terms $\lambda_i\mathsf l_i$ tend to $0$, so the series converges in the complete ultrametric $P$, and the bounds persist. So $\Xi$ is injective with closed range, as $c_0(\N,K)$ is complete. The range contains every $P_m$, hence is dense, hence is $P$.
\end{proof}

\begin{theorem}[$c_0$-type and countable-type sources]\label{dp:thm:c0-agreement}
Let $W$ be ultrametric.
\begin{enumerate}
\item Let $n\ge1$, let $P=c_{00}(J)$, and fix a total order on $J$. Every law $a\colon P\to W$ of degree $n$ is the diagonal of the $n$\nobreakdash-linear map
\[
 \beta(x^1,\dots,x^n)=\sum_{\mathsf w}v_{\mathsf w}\,x^1_{\mathsf w_1}\cdots x^n_{\mathsf w_n},
\]
over the nondecreasing words $\mathsf w$ in $J$ of length $n$, a finite sum on $c_{00}(J)$, where $v_{\mathsf w}=a^{[\gamma]}(e_{\mathsf j_1},\dots,e_{\mathsf j_r})$ for the distinct letters $\mathsf j_1<\dots<\mathsf j_r$ of $\mathsf w$ and their multiplicities $\gamma$. Moreover $\norm\beta\le\abs\varpi^{-n}\norm a$.
\item The same holds for $P=c_0(J)$ if $W$ is complete; the sum converges. In both cases, a map from an open subset of $P$ to $W$ that is DP at $x_0$ with radius $R$ and bound $M$ is analytic there with radius $\abs\varpi R$ and bound $M$.
\item Let $K$ be complete, $P$ an ultrametric Banach space of countable type and $W$ ultrametric Banach. For $\mathfrak o,\mathfrak o'$ as in Lemma~\ref{dp:lem:basis}, a germ on $P$ that is DP at $x_0$ with radius $R$ and bound $M$ is analytic there with radius $\abs\varpi\mathfrak o\mathfrak o'R$ and bound $M$.
\end{enumerate}
\end{theorem}

\begin{proof}
\emph{Step 1 (item (1)).} The map $\beta$ is $n$\nobreakdash-linear. For $h$ supported in a finite set $\{\mathsf j_1,\dots,\mathsf j_r\}$ we have $a(h)=\sum_\gamma h^\gamma a^{[\gamma]}(e_{\mathsf j_1},\dots,e_{\mathsf j_r})$, and each monomial comes from exactly one nondecreasing word. So $\beta(h,\dots,h)=a(h)$. The restriction of $a$ to $\spn(e_{\mathsf j_1},\dots,e_{\mathsf j_r})$, isometric to $c_0([r],K)$, has norm at most $\norm a$. So \cite[Lemma~Q.2(4)]{I} gives $\norm{v_{\mathsf w}}\le\abs\varpi^{-n}\norm a$. In the ultrametric $W$ the number of words is not counted, so $\norm{\beta(x^1,\dots,x^n)}\le\abs\varpi^{-n}\norm a\prod_c\norm{x^c}$.

\emph{Step 2 (item (2)).} For $x^c\in c_0(J)$ and $\varepsilon>0$, only finitely many words have $\prod_c\abs{x^c_{\mathsf w_c}}>\varepsilon$. So the family is summable in the complete ultrametric $W$, with the same bound. We have $\beta(h,\dots,h)=a(h)$ on the dense subspace $c_{00}(J)$, and both sides are continuous, $a$ by Theorem~\ref{dp:thm:mixed}(2). For the statement on germs, let $a_n$ be the coefficient laws and $\beta_n$ the map above with diagonal $a_n$, for $n\ge1$, and put $\beta_0=a_0$. Then $\norm{\beta_n}\le M(\abs\varpi R)^{-n}$.

\emph{Step 3 (item (3)).} We may assume $P\ne\{0\}$. If $P$ is infinite-dimensional, let $\Xi\colon c_0(\N,K)\to P$ be as in Lemma~\ref{dp:lem:basis}. If $\dim P=m$ is finite, let $\Xi\colon c_0([m],K)\to P$ be the map $\mathsf C$ of \cite[Lemma~Q.13(3)]{I} (with $\mathfrak o=1$ if $K$ is discretely valued, which only improves the bound). In both cases $\norm\Xi\le1$ and $\norm{\Xi^{-1}}\le(\mathfrak o\mathfrak o')^{-1}$. Let $a_n$ be the coefficient laws at $x_0$. The laws $a_n\circ\Xi$ have norm at most $MR^{-n}$ (Lemma~\ref{dp:lem:elementary}(2)). By item (2), for $n\ge1$, $a_n\circ\Xi$ is the diagonal of a map $\beta'_n$ of norm at most $\abs\varpi^{-n}MR^{-n}$. Then $\beta_n=\beta'_n\circ(\Xi^{-1})^{\times n}$ has diagonal $a_n$ and norm at most $M(\abs\varpi\mathfrak o\mathfrak o'R)^{-n}$. Put $\beta_0=a_0$.
\end{proof}

A norm $\norm\cdot'$ on $W$ is \emph{equivalent with constant $C_W$} if $\norm w\le\norm w'\le C_W\norm w$ for all $w\in W$. Every equivalent norm has this form after scaling by a positive real number, and scaling keeps a norm ultrametric and spherically complete.

\begin{theorem}[spherically complete targets, discretely valued fields]\label{dp:thm:sph-agreement}
Let $P$ be ultrametric, $U\subseteq P$ open and $f\colon U\to W$.
\begin{enumerate}
\item Let $W$ have an equivalent norm with constant $C_W$ that is ultrametric and spherically complete. Every law $a\colon P\to W$ of degree $n$ has a lift of norm at most $C_W(\abs\varpi\mathfrak c_K)^{-n}\norm a$. The map $f$ is DP at $x_0$ if and only if it is analytic at $x_0$. A DP expansion with radius $R$ and bound $M$ gives an analytic one with radius $\abs\varpi\mathfrak c_KR$ and bound $C_WM$.
\item Let $K$ be discretely valued and $W$ Banach with an equivalent ultrametric norm with constant $C_W$. Every law $a$ of degree $n$ from $P$ to $W$ has a lift of norm at most $\mathfrak c_K^{-1}C_W(\abs\varpi\mathfrak c_K)^{-n}\norm a$. The same equivalence holds, with radius $\abs\varpi\mathfrak c_KR$ and bound $\mathfrak c_K^{-1}C_WM$.
\end{enumerate}
\end{theorem}

\begin{proof}
A law into $W$ is a law into $(W,\norm\cdot')$, of norm at most $C_W$ times its old norm, and a lift for $\norm\cdot'$ is a lift for $\norm\cdot$ of no larger norm. So \cite[Proposition~14.40(1)]{I}, applied to $(W,\norm\cdot')$, gives both statements on laws and the passage from DP to analytic. In~(2), $(W,\norm\cdot')$ is an ultrametric Banach space, and we use the last sentence of that proposition. The converse is Theorem~\ref{dp:thm:calculus}(3).
\end{proof}

For incomplete ultrametric $W$ over discretely valued $K$, Theorem~\ref{dp:thm:sph-agreement}(2), applied to $\widehat W$, gives a lift only into $\widehat W$. A lift into $W$ need not exist over $\F_p(t)$ (Theorem~\ref{dp:thm:agreement-finite}(3)). For other $K$, whether it exists is part of Question~\ref{dp:q:residual}.

\begin{lemma}[operator spaces]\label{dp:lem:operator-sph}
Let $X$ be ultrametric and $Y$ ultrametric and spherically complete. Then $\cL(X,Y)$ with the operator norm is ultrametric and spherically complete.
\end{lemma}

\begin{proof}
The operator norm into an ultrametric space is ultrametric. Consider a nonempty family of pairwise meeting closed balls in $\cL(X,Y)$, with centers $u_\theta$ and radii $\varrho_\theta$. Then $\norm{u_\theta-u_{\theta'}}\le\max(\varrho_\theta,\varrho_{\theta'})$. Step~1 of the proof of \cite[Theorem~5.2]{I}, with $\{0\}$ as the subspace on which the map is given, yields a linear $u\colon X\to Y$ with $\norm{ux-u_\theta x}\le\varrho_\theta\norm x$ for all $x$ and $\theta$. So $u$ is bounded, and $u$ lies in every ball.
\end{proof}

\begin{theorem}[where a DP germ that is not analytic can occur]\label{dp:thm:conservativity}
Let $K$ be complete, $P$ ultrametric, $W$ an ultrametric Banach space, $U\subseteq P$ open and $f\colon U\to W$ DP at $x_0\in U$. Let $\widehat P$ be the completion of $P$.
\begin{enumerate}
\item The map $f$ is analytic at $x_0$ in each of the following cases: \textup{(a)} $\operatorname{char}K=0$; \textup{(b)} $K$ is discretely valued; \textup{(c)} $W$ has an equivalent ultrametric spherically complete norm; \textup{(d)} $\widehat P$ is linearly homeomorphic to $c_0(J,K)$ for some set $J$. Case~\textup{(d)} includes every $P$ of finite dimension and every $P$ of countable type.
\item So if $f$ is not analytic at $x_0$, then all of the following hold: $\operatorname{char}K=p>0$; $K$ is densely valued; $\widehat P$ is not linearly homeomorphic to any $c_0(J,K)$; and $W$ has no equivalent ultrametric spherically complete norm.
\item Such germs exist if $K$ has characteristic $p$ and is not spherically complete: $A^k_{E,D;K}$ for $k\ge p$ and the spaces of \cite[Theorem~7.2(2)]{I}.
\item Items \textup{(1)}--\textup{(3)} also hold with AA in place of DP.
\item Here $P$ and $W$ need not be ultrametric. Let $P$ be finite-dimensional and $W$ any normed space. A map from an open subset of $P$ to $W$ is DP at a point if and only if it is analytic there. So DP manifolds modeled on $K^m$ are the analytic manifolds, with the same maps, and DP vector bundles with normed fibers over such a base are the analytic ones.
\end{enumerate}
\end{theorem}

\begin{proof}
\emph{Step 1 (cases (a), (b), (c)).} These are Theorem~\ref{dp:thm:agreement-char0}(2) and Theorem~\ref{dp:thm:sph-agreement}(2),(1).

\emph{Step 2 (case (d)).} We may assume $P\ne\{0\}$. Let $T\colon c_0(J,K)\to\widehat P$ be a bounded linear bijection with bounded inverse, and let $f$ have radius $R$, bound $M$ and coefficient laws $a_n$. Here $\Kh=K$ and $\widehat W=W$. By \cite[Lemma~Q.13(1)]{I}, for $n\ge1$ the law $a_n$ extends to a law $\hat a_n\colon\widehat P\to W$ with $\norm{\hat a_n}\le\norm{a_n}$. By Theorem~\ref{dp:thm:c0-agreement}(2), $\hat a_n\circ T$ is the diagonal of a map $\beta'_n$ of norm at most $\abs\varpi^{-n}\norm T^n\norm{a_n}$. The restriction of $\beta'_n\circ(T^{-1})^{\times n}$ to $P^n$ is a lift of $a_n$ of norm at most $M(\abs\varpi R/(\norm T\norm{T^{-1}}))^{-n}$. With $\beta_0=a_0$, this makes $f$ analytic at $x_0$ with bound $M$. If $P$ has finite dimension, so has $\widehat P$, and \cite[Lemma~A.3]{I} gives $T$ with $J$ finite. If $P$ is of countable type, so is $\widehat P$, and Lemma~\ref{dp:lem:basis} gives $T$ with $J=\N$ if $\widehat P$ is infinite-dimensional.

\emph{Step 3 (items (2), (3), (4)).} Item (2) is the contrapositive of (1). For (3), see Proposition~\ref{dp:prop:strict}(1)(c): the source and the target are ultrametric Banach spaces, and the law is AA, hence DP, at every point and analytic at no point. For (4), AA germs are DP by \cite[Proposition~14.9(1)]{I}, and the germ of (3) is AA.

\emph{Step 4 (item (5)).} We may assume $P\ne\{0\}$. As $K$ is complete, $P$ has a bijection $T$ as in Theorem~\ref{dp:thm:agreement-finite}(1), so a DP germ on $P$ is analytic. The converse is Theorem~\ref{dp:thm:calculus}(3). Transitions, maps in charts and bundle transition families have finite-dimensional sources. So being DP and being analytic are the same for them.
\end{proof}

Item~(2) gives necessary conditions only. It does not say that a given pair $(P,W)$ with these properties carries such a germ. Item~(3) leaves one class of fields open: whether such germs exist over a densely valued, spherically complete $K$ of characteristic $p$ (Question~\ref{dp:q:residual}).

\subsection{General norms, and the comparison by fields}\label{dp:ssec:agree-general}

\begin{theorem}[general norms]\label{dp:thm:general-norms-agreement}
Let $P,W$ be normed. Always analytic germs are AA, and AA germs are DP.
\begin{enumerate}
\item \textup{(Agreement.)} DP, AA and analytic coincide at every point in each of these cases: \textup{(a)} $P$ is finite-dimensional, and $K$ is complete or $W$ is Banach; \textup{(b)} $\operatorname{char}K=0$ and $P$ is ultrametric. For $P$ finite-dimensional, DP and AA coincide for every $K$ and $W$. For $n!\ne0$ in $K$, every law of degree $n$ is analytic.
\item \textup{(AA $=$ DP, analytic $\ne$ AA, every characteristic.)} On $P=\ell^1(J)$ every DP germ with values in a normed space is AA. The germ $\mathsf f$ of Example~\ref{dp:ex:l1}, on $\ell^1(\N,K)$ with values in $\Kh$, is AA at $0$ and not analytic at $0$.
\item \textup{(Ultrametric sources.)} Let $P$ be ultrametric.
\begin{enumerate}
\item If $W$ is ultrametric, then DP $=$ AA. For $K$ complete and $W$ Banach, DP against analytic is Theorem~\ref{dp:thm:conservativity}. Now let $\operatorname{char}K=p$, with $K$ or $W$ incomplete. Then we have Theorems~\ref{dp:thm:agreement-finite}(1) and~\ref{dp:thm:c0-agreement}(1); for $W$ Banach, Theorems~\ref{dp:thm:agreement-finite}(2) and~\ref{dp:thm:c0-agreement}(2); for $W$ as there, Theorem~\ref{dp:thm:sph-agreement}(1); for $K$ discretely valued and $W$ Banach, Theorem~\ref{dp:thm:sph-agreement}(2); over $\F_p(t)$, the failure in Theorem~\ref{dp:thm:agreement-finite}(3); and for $K$ incomplete with $\Kh$ not spherically complete, the failure proved after Question~\ref{dp:q:residual}. We do not decide the rest.
\item If $W$ is not ultrametric, AA and analytic can differ. Let $\operatorname{char}K=p$ and $k\ge p$. Under condition~\textup{(L)} of \cite[Section~7.1]{I}, which makes $K$ discretely valued, or if $p=k=2$, the power map $c_0\to\hGamma^k_\pi c_0$, with $c_0=c_0(\N,K)$, is AA everywhere and analytic nowhere (Proposition~\ref{dp:prop:strict}(1)(b)). Its target is not ultrametric.
\item Every map from an open subset of an ultrametric space to a normed space that is DP at a point is AA there if and only if $\sup_n\mathfrak a_n(c_0(\N,K))^{1/n}<\infty$. This holds for $\operatorname{char}K=0$ \cite[Theorem~14.39(3)]{I}. For $\operatorname{char}K=p$ we do not decide it (Question~\ref{dp:q:c0}).
\end{enumerate}
\item \textup{(AA $\ne$ DP, characteristic $p$.)} Over every $K$ of characteristic $p$ there are DP germs that are AA at no point (Proposition~\ref{dp:prop:strict}(2)). A law that is not ambient has a non-ultrametric target, by \cite[Theorem~14.39(4)]{I}. The sources of the laws of Proposition~\ref{dp:prop:strict}(2) contain $\ell^1$-sums.
\item \textup{(Characteristic $0$, general sources.)} Whether DP $=$ AA for every normed $P$ is not decided here \cite[Question~14.42(1)]{I}. It holds for ultrametric $P$, for $\ell^1(J)$ and in finite dimension.
\end{enumerate}
\end{theorem}

\begin{proof}
\emph{Step 1 (items (1), (2), (3)(a), (4), (5)).} Item (1) is Theorems~\ref{dp:thm:agreement-finite}, \ref{dp:thm:agreement-char0}(2) and~\ref{dp:thm:degree-threshold}(1). For (2), $\mathfrak a_n(\ell^1(J))\le C_0^n$ \cite[Theorem~14.39(2)]{I}, and \cite[Theorem~14.38(1)]{I} applies. Item (3)(a) is \cite[Theorem~14.39(1)]{I} with the cited theorems. Items (4) and (5) are the cited results, with (1) and (2). The case of DP and AA in finite dimension is Step~2.

\emph{Step 2 (DP and AA for $P$ finite-dimensional).} We may assume $P\ne\{0\}$. Let $f\colon U\to W$ be DP at $x_0$ with radius $R$, bound $M$ and coefficient laws $a_n$, and let $j\colon W\to\widehat W$ be the inclusion. It is a linear isometry into a Banach space. By Theorem~\ref{dp:thm:agreement-finite}(2), for $n\ge1$ the law $j\circ a_n$ is the diagonal of a bounded $n$\nobreakdash-linear map $\beta_n\colon P^n\to\widehat W$ with $\norm{\beta_n}\le2^{\mathsf r}C_P^nMR^{-n}$. Put $\beta_0=j\circ a_0$. The diagonals lie in $j(W)$, and $jf(x_0+h)=\sum_n\beta_n(h,\dots,h)$ for $\norm h<R/C_P$. So $f$ is AA at $x_0$. The converse holds for every $P$, by \cite[Proposition~14.9(1)]{I}.

\emph{Step 3 (item (3)(b)).} The source $c_0$ is ultrametric. The map is AA everywhere and analytic nowhere by Proposition~\ref{dp:prop:strict}(1)(b). It remains to see that the target is not ultrametric. Under~(L), $K$ is discretely valued~\cite[Section~7.1]{I}, and the target is Banach. So an ultrametric target would make the map analytic, by Theorem~\ref{dp:thm:sph-agreement}(2). Let $p=k=2$. Let $\mathsf g=e_1\otimes e_1+e_2\otimes e_2$, and let $\mathsf g=\sum_lx_l\otimes y_l$ be a representation in $c_0\otimes c_0$. Projecting to the first two coordinates gives $\id_{K^2}=\sum_l\bar x_l\bar y_l^{\mathsf T}$, where $\bar x_l,\bar y_l\in K^2$ are the projections of $x_l,y_l$. Taking determinants, the Cauchy--Binet formula gives $1=\sum_{l<l'}\det(\bar x_l,\bar x_{l'})\det(\bar y_l,\bar y_{l'})$. By the ultrametric inequality in $K$, some $l\ne l'$ has $\abs{\det(\bar x_l,\bar x_{l'})\det(\bar y_l,\bar y_{l'})}\ge1$. So $\norm{x_l}\norm{y_l}\norm{x_{l'}}\norm{y_{l'}}\ge1$, and hence $\norm{x_l}\norm{y_l}+\norm{x_{l'}}\norm{y_{l'}}\ge2$. Thus $\pi(\mathsf g)\ge2$, while $\pi(e_i\otimes e_i)\le1$ for $i=1,2$. So $\pi$ does not satisfy the ultrametric inequality on $\Gamma^2(c_0)$.

\emph{Step 4 (item (3)(c)).} Let $\mathfrak q=\sup_n\mathfrak a_n(c_0(\N,K))^{1/n}<\infty$, and let $P$ be ultrametric. By Proposition~\ref{dp:prop:ultrametric}(3), $\mathfrak a_n(P,W)\le\mathfrak a_n(P)\le(\mathfrak c_K^{-1}\max(\mathfrak q,1))^n$. So \cite[Theorem~14.38(1)]{I} applies. Conversely, let $\mathfrak q=\infty$ and $P=c_0(\N,K)$. Let $W_n$ be the completion of $\Gamma^n_{\mathrm s}(P)$, and $W$ the $c_0$-sum of the Banach spaces $W_n$, as in Step~2 of the proof of Theorem~\ref{dp:thm:agreement-criterion}; it is Banach. The law $\iota_n\circ\mathsf v_n$ has norm at most $1$, and its linearization is $t\mapsto\iota_n[t]$, so $\mathfrak a_n(P,W)\ge\mathfrak a_n(P)$. Thus $(P,W)$ is not tame, and \cite[Proposition~14.37]{I} gives a germ into $W$ that is DP at a point and not AA there.
\end{proof}

The next table sums this up. An entry $=$ or $\ne$ means for all spaces of the row, or for some spaces of the row.

\begin{center}
\small
\begin{tabular}{p{4.6cm}p{2.9cm}p{6.1cm}}
\hline
Setting & DP against AA & AA against analytic\\
\hline
$P$ finite-dimensional, $K$ complete or $W$ Banach & $=$ & $=$\\
$P$ finite-dimensional, $K$ and $W$ incomplete & $=$ & $=$ for $P$ as in Theorem~\ref{dp:thm:agreement-finite}(1); $\ne$ for some spaces if $K=\F_p(t)$; not decided in general\\
$\operatorname{char}K=0$, $P$ ultrametric & $=$ & $=$\\
$\operatorname{char}K=0$, $P=\ell^1(\N,K)$ & $=$ & $\ne$\\
$\operatorname{char}K=0$, $P$ general & not decided & $\ne$\\
$\operatorname{char}K=p$, $P,W$ ultrametric, $K$ complete, $W$ Banach & $=$ & $=$ in the cases of Theorem~\ref{dp:thm:conservativity}(1); $\ne$ for some spaces if $K$ is not spherically complete; not decided in general if $K$ is densely valued and spherically complete\\
$\operatorname{char}K=p$, $P,W$ ultrametric, $K$ complete, $W$ incomplete & $=$ & $=$ for $P$ finite-dimensional and for $P=c_{00}(J)$; otherwise not decided\\
$\operatorname{char}K=p$, $P,W$ ultrametric, $K$ densely valued, incomplete & $=$ & $=$ for $W$ as in Theorem~\ref{dp:thm:sph-agreement}(1), for $P$ as in Theorem~\ref{dp:thm:agreement-finite}(1), for finite-dimensional $P$ with $W$ Banach, for $P=c_{00}(J)$, and for $P=c_0(J)$ with $W$ Banach; $\ne$ for some spaces if $\Kh$ is not spherically complete; otherwise not decided\\
$\operatorname{char}K=p$, $P,W$ ultrametric, $K$ discretely valued, incomplete & $=$ & $=$ for $W$ Banach, for $P$ as in Theorem~\ref{dp:thm:agreement-finite}(1), and for $P=c_{00}(J)$; $\ne$ for some finite-dimensional spaces with $W$ incomplete if $K=\F_p(t)$; otherwise not decided\\
$\operatorname{char}K=p$, $P$ ultrametric, $W$ not & not decided; $=$ for all $P$, $W$ if and only if the supremum in item~(3)(c) is finite & $\ne$ under (L) or for $p=2$\\
$\operatorname{char}K=p$, general & $\ne$ & $\ne$\\
\hline
\end{tabular}
\end{center}

So a DP germ that is not analytic needs characteristic $p$ or a non-ultrametric source, by item~(1)(b). With an ultrametric source and a non-ultrametric target it occurs over discretely valued $K$ under~(L), although Theorem~\ref{dp:thm:sph-agreement}(2) excludes it for ultrametric Banach targets. Between ultrametric spaces over the discretely valued field $\F_p(t)$ it occurs with an incomplete target (Theorem~\ref{dp:thm:agreement-finite}(3)). The DP germs that are not AA, among those listed here, have characteristic $p$ and non-ultrametric norms. Between ultrametric spaces there are none.

We now compare with the functor $\alpha_k$ of~\cite[n\textsuperscript{o}~7.8.1]{Bou}, where our $k$ is written $d$, and for which $\alpha_k(\id_F,u)=A^k_{E,D;F}(u)$~\cite[Section~10.2]{I}. Let $K$ be complete. Consider two statements.
\begin{itemize}
\item[$(\mathrm D_K)$] Every map between open subsets of ultrametric Banach spaces over $K$ that is DP at a point is analytic there.
\item[$(\mathrm B_K)$] For all ultrametric Banach spaces $E,D,F,F'$ and every $k\ge1$, the maps $A^k_{E,D;F}$ and the joint action $B$ on $\cL(E,D)\times\cL(F,F')$, restricted to alternating maps, are analytic everywhere.
\end{itemize}
The statement $(\mathrm B_K)$ is what \cite[paragraph \emph{Bourbaki's question}, and the proof of Corollary~10.5]{I} means by saying that $\alpha_k$ is, on ultrametric Banach spaces, a vector functor of the analytic class in every degree.

\begin{theorem}[comparison by fields]\label{dp:thm:field-level}
Let $K$ be complete.
\begin{enumerate}
\item $(\mathrm D_K)$ implies $(\mathrm B_K)$.
\item Both hold if $\operatorname{char}K=0$ or if $K$ is discretely valued.
\item Both fail if $\operatorname{char}K=p>0$ and $K$ is not spherically complete.
\end{enumerate}
\end{theorem}

\begin{proof}
\emph{Step 1 (item (1)).} By Theorem~\ref{dp:thm:calculus}(4),(1), $A^k_{E,D;F}$ and $B$ are laws, DP on their whole domains. Their sources are $\cL(E,D)$ and $\cL(E,D)\times\cL(F,F')$, and their target is $\cL(\Alt^k(D;F),\Alt^k(E;F'))$. These spaces are ultrametric, since operator norms into ultrametric spaces are, and they are complete, since $D$ and $F'$ are complete and $\Alt^k(E;F')$ is closed in $\Mult^k(E;F')$. So $(\mathrm D_K)$ applies at every point.

\emph{Step 2 (items (2), (3)).} If $\operatorname{char}K=0$, Theorem~\ref{dp:thm:agreement-char0}(2) gives $(\mathrm D_K)$. If $K$ is discretely valued, Theorem~\ref{dp:thm:sph-agreement}(2) gives $(\mathrm D_K)$. Then $(\mathrm B_K)$ holds by item (1). If $\operatorname{char}K=p$ and $K$ is not spherically complete, then for $k\ge p$ the map $A^k_{E,D;K}$ of \cite[Theorem~7.2(2)]{I} is analytic at no point. So $(\mathrm B_K)$ fails, and $(\mathrm D_K)$ fails by item (1).
\end{proof}

A complete discretely valued field is spherically complete~\cite[Lemma~D.1]{I}. So the one remaining case is $\operatorname{char}K=p$ with $K$ densely valued and spherically complete. There neither statement is decided here (\cite[Section~12]{I} and Question~\ref{dp:q:residual}). Theorem~\ref{dp:thm:field-level} is one implication, together with a comparison on the fields decided here: wherever $(\mathrm B_K)$ is decided here, $(\mathrm D_K)$ is decided with the same answer. Whether $(\mathrm B_K)$ implies $(\mathrm D_K)$ for such $K$ is open (Question~\ref{dp:q:residual}). A negative answer to $(\mathrm D_K)$ over a densely valued spherically complete $K$ would not decide $(\mathrm B_K)$; a failure of $(\mathrm B_K)$ there would give a failure of $(\mathrm D_K)$.

\begin{remark}[pairs of spaces]\label{dp:rem:field-level-spaces}
Theorem~\ref{dp:thm:field-level} is about fields. On a field where both statements fail, for many pairs of spaces every DP germ is analytic, by Theorems~\ref{dp:thm:c0-agreement} and~\ref{dp:thm:sph-agreement}. Those conditions concern the source and the target of a germ. The results of \cite{I} on $\alpha_k$ concern the fibers $E,D,F$, and the germ $A^k_{E,D;F}$ has source $\cL(E,D)$ and target $\cL(\Alt^k(D;F),\Alt^k(E;F))$.

\emph{Spherically complete values.} Let $E,D$ be ultrametric and $F$ ultrametric and spherically complete. By \cite[Theorem~5.2]{I}, $\Alt^k(E;F)$ and $\Alt^k(D;F)$ are ultrametric and spherically complete. So the target of $A^k_{E,D;F}$ is spherically complete, by Lemma~\ref{dp:lem:operator-sph}. The source $\cL(E,D)$ is ultrametric, as $D$ is. Hence Theorem~\ref{dp:thm:sph-agreement}(1) makes $A^k_{E,D;F}$ the diagonal of a bounded $k$\nobreakdash-linear map. For ultrametric $E,D$, this recovers the last conclusion of \cite[Theorem~5.2]{I}, that $A^k_{E,D;F}$ is a continuous polynomial, from its first assertion.

\emph{Countable type.} Here the results are only parallel. \cite[Corollary~10.5]{I} assumes $E$ or $D$ of countable type, and Theorem~\ref{dp:thm:c0-agreement}(3) assumes $P$ of countable type, and $\cL(E,D)$ need not be of countable type. Indeed, let $K$ be complete of characteristic $p$ and not spherically complete, and let $k\ge p$. For the spaces of \cite[Theorem~7.2(2)]{I} it is not, by Theorem~\ref{dp:thm:c0-agreement}(3), since $A^k_{E,D;K}$ is DP and not analytic. Neither result is deduced from the other here.
\end{remark}
%% ---- end sec-11-agreement.tex
%% ---- begin sec-12-functors.tex
\section{Vector functors and Bourbaki's \texorpdfstring{n\textsuperscript{o}~7.8}{no. 7.8}}\label{dp:sec:functors}

Bourbaki defines vector functors of class $C^m$ by a regularity condition on their maps on hom spaces~\cite[n\textsuperscript{o}~7.6.1]{Bou}. Bourbaki defines the alternating functor $\alpha_d$ when $K$ has characteristic $0$ or the bundles have finite rank, and states that it is not known whether this restriction can be removed~\cite[n\textsuperscript{o}~7.8]{Bou}. This section puts DP in place of $C^m$. We write $k$ for Bourbaki's $d$, so the functor is $\alpha_k$. Theorem~\ref{dp:thm:seven-eight-dp} collects the results.

\subsection{Polynomial functors}

Let $\mathsf{Vec}_K$ be the category of normed spaces and bounded linear maps, $\mathsf{Ban}_K$ its full subcategory of Banach spaces, and $\cC_{\rm um}$ the category of ultrametric Banach spaces, a full subcategory of $\mathsf{Ban}_K$. In this section $\cC$ is one of these three. Let $I=I_+\sqcup I_-$ be a finite set. For families $\mathbf V=(V_i)_{i\in I}$ and $\mathbf V'=(V'_i)_{i\in I}$ of objects of $\cC$, put
\[
 \operatorname{Hom}(\mathbf V,\mathbf V')=\prod_{i\in I_+}\cL(V_i,V_i')\times\prod_{i\in I_-}\cL(V_i',V_i),
\]
with the maximum norm. Composition is entrywise, in the reverse order at the entries in $I_-$. It is bilinear, associative, unital and of norm at most $1$. A \emph{vector functor of type $I$ on $\cC$} assigns to each $\mathbf V$ a normed space $\mathfrak F(\mathbf V)$, a Banach space if $\cC\ne\mathsf{Vec}_K$, and to each $\mathbf f\in\operatorname{Hom}(\mathbf V,\mathbf V')$ an operator $\mathfrak F(\mathbf f)\in\cL(\mathfrak F(\mathbf V),\mathfrak F(\mathbf V'))$, with $\mathfrak F(\id)=\id$ and $\mathfrak F(\mathbf f'\circ\mathbf f)=\mathfrak F(\mathbf f')\mathfrak F(\mathbf f)$. Its \emph{hom maps} are the maps $\mathfrak F_{\mathbf V,\mathbf V'}\colon\operatorname{Hom}(\mathbf V,\mathbf V')\to\cL(\mathfrak F(\mathbf V),\mathfrak F(\mathbf V'))$.

\begin{definition}\label{dp:def:poly-functor}
Let $\mathfrak F$ be a vector functor of type $I$ on $\cC$, and let $n\ge0$.
\begin{enumerate}
\item $\mathfrak F$ is a \emph{bounded strict polynomial functor, homogeneous of degree $n$}, if every hom map is a law of degree $n$. It is \emph{contractive} if these laws have norm at most $1$.
\item $\mathfrak F$ is \emph{multi-homogeneous of multidegree} $\mathbf d=(d_i)_{i\in I}\in\N^I$ if it is homogeneous of degree $\abs{\mathbf d}=\sum_id_i$ and $\mathfrak F((\lambda_if_i)_i)=\prod_i\lambda_i^{d_i}\,\mathfrak F(\mathbf f)$ for all scalars $\lambda_i$, with $\lambda^0=1$.
\item $\mathfrak F$ is a bounded strict polynomial functor \emph{of degree at most $n$} if every hom map is a finite sum of laws of degrees at most $n$.
\item $\mathfrak F$ is a \emph{DP vector functor} if every hom map is DP on its whole domain, and an \emph{analytic vector functor} if every hom map is analytic at every point of its domain.
\end{enumerate}
\end{definition}

By Theorem~\ref{dp:thm:calculus}(3), analytic vector functors are DP vector functors.

\emph{Examples.} Let $k\ge1$ and let $I_+=\{0\}$, $I_-=\{1\}$. The functors $\Mult^k(V_1;V_0)$ and $\alpha_k(V_1;V_0)=\Alt^k(V_1;V_0)$ act by $\alpha_k(f_0,f_1)(\beta)=f_0\circ\beta\circ(f_1\times\dots\times f_1)$, so $A^k_{E,D;F}(u)=\alpha_k(\id_F,u)$. The hom map of $\Mult^k$ is the diagonal of a bounded $(k+1)$\nobreakdash-linear map of norm at most $1$. So $\Mult^k$ is analytic and contractive, of multidegree $(1,k)$, and $\alpha_k$ is a closed subfunctor. In one covariant variable, $\hot^k_\pi$ is the completed projective tensor power. Its hom map is the diagonal of $(u_1,\dots,u_k)\mapsto u_1\otimes\dots\otimes u_k$, so it is analytic and contractive. The functor $\hLambda^k_\pi$ is its quotient by the kernel of the natural map $\widehat\wedge^k_E\colon\hot^k_\pi E\to\hLambda^k_\pi E$, a closed subfunctor, with the quotient norm~\cite[Lemma~3.7, Definition~3.8(1)]{I}.

\begin{lemma}[two witnesses]\label{dp:fn:lem:witnesses}
Let $p=\operatorname{char}K>0$ and $k\ge p$.
\begin{enumerate}
\item For the Banach spaces $E,F$ of~\cite[Theorem~7.2(1)]{I}, the maps $A^k_{E,E;F}$ and $u\mapsto\hLambda^k_\pi u$ are analytic at no point of $\cL(E,E)$.
\item If $K$ is complete and not spherically complete, then $A^k_{E,D;K}$ is analytic at no point of $\cL(E,D)$, for the ultrametric Banach spaces $E,D$ of~\cite[Theorem~7.2(2)]{I}.
\end{enumerate}
\end{lemma}

\begin{proof}
\emph{Step 1 (the maps $A^k$).} These statements are~\cite[Theorem~7.2]{I}.

\emph{Step 2 (the map $u\mapsto\hLambda^k_\pi u$).} Restriction to the dense subspace $\Lambda^k_\pi E$ is an isometric isomorphism of $\cL(\hLambda^k_\pi E,\hLambda^k_\pi E)$ onto $\cL(\Lambda^k_\pi E,\hLambda^k_\pi E)$, and it sends $\hLambda^k_\pi u$ to $\Lambda^k_\pi u$. If $u\mapsto\hLambda^k_\pi u$ were analytic at a point, then by~\cite[Lemma~3.4]{I} the pair $(E,E)$ would be good for Banach values in degree $k$, which means that $A^k_{E,E;F}$ is analytic for every Banach $F$~\cite[Definition~3.3]{I}. This contradicts~\cite[Theorem~7.2(1)]{I}.
\end{proof}

Recall the normed space $\Gamma^n_{\mathrm s}(P)$, its completion $\hGamma^n(P)$ and the universal law $\mathsf v_n$ of Theorem~\ref{dp:thm:normed-roby}.

\begin{lemma}[what represents laws]\label{dp:lem:functor-rep}
\leavevmode
\begin{enumerate}
\item \textup{(Coefficients.)} Let $W$ be a vector space and $W_0\subseteq W$ a linear subspace. Let $a\colon P\to W$ have values in $W_0$, and suppose that, for all $h_1,\dots,h_r\in P$, the map $s\mapsto a(\sum_cs_ch_c)$ is a homogeneous polynomial of degree $n$ with coefficients in $W$. Then the coefficients lie in $W_0$.
\item \textup{(Representing functor.)} For $n\ge0$, the assignment $\Gamma^n_{\mathrm s}(u)[t]=[u^{\otimes n}t]$ makes $E\mapsto\Gamma^n_{\mathrm s}(E)$ a contractive bounded strict polynomial functor, homogeneous of degree $n$, on $\mathsf{Vec}_K$. Extension to completions makes $\hGamma^n$ one on $\mathsf{Vec}_K$ and on $\mathsf{Ban}_K$. By Theorem~\ref{dp:thm:normed-roby}(3), $\cP_n(P;W)=\cL(\Gamma^n_{\mathrm s}(P),W)$ isometrically.
\item \textup{(The divided-power category.)} Let $\cC^I_n$ be the category whose objects are the families $\mathbf V$ of objects of $\cC$ and whose hom spaces are $\Gamma^n_{\mathrm s}(\operatorname{Hom}(\mathbf V,\mathbf V'))$, with the composition induced by $[\mathbf f'^{\otimes n}]\circ[\mathbf f^{\otimes n}]=[(\mathbf f'\circ\mathbf f)^{\otimes n}]$. This composition is bilinear of norm at most $1$. A homogeneous bounded strict polynomial functor $\mathfrak F$ of type $I$ and degree $n$ on $\cC$ is the same as a functor $\Phi$ from $\cC^I_n$ to normed spaces, Banach if $\cC\ne\mathsf{Vec}_K$, that is bounded linear on each hom space, through $\mathfrak F(\mathbf f)=\Phi([\mathbf f^{\otimes n}])$. Then $\Phi=L_{\mathfrak F_{\mathbf V,\mathbf V'}}$ on $\Gamma^n_{\mathrm s}(\operatorname{Hom}(\mathbf V,\mathbf V'))$, with the same norm as the law $\mathfrak F_{\mathbf V,\mathbf V'}$.
\end{enumerate}
\end{lemma}

\begin{proof}
\emph{Step 1 (item (1)).} As $K$ is infinite, interpolation in each variable writes each coefficient as a finite linear combination of values of the polynomial. These values lie in $W_0$.

\emph{Step 2 (item (2)).} A representation $t=\sum_l\lambda_lh_l^{\otimes n}$ gives $u^{\otimes n}t=\sum_l\lambda_l(uh_l)^{\otimes n}$, so $\norm{u^{\otimes n}t}_\Gamma\le\norm u^n\norm t_\Gamma$. For fixed $t$ and $u_1,\dots,u_r$, by~\eqref{dp:eq:orbit}, $s\mapsto(\sum_cs_cu_c)^{\otimes n}t$ is a homogeneous polynomial with coefficients in $\Gamma^n(D)$. The coefficients are linear in $[t]$, by uniqueness. Item (1), with $W$ the linear maps $\Gamma^n_{\mathrm s}(E)\to\Gamma^n_{\mathrm s}(D)$ and $W_0$ the bounded ones, shows that the hom map is a law of norm at most $1$. Extension to completions is a linear isometry of operator spaces.

\emph{Step 3 (item (3)).} On $\operatorname{Hom}(\mathbf V,\mathbf V')^{\otimes n}$ composition is defined factor by factor, as the composition of $\operatorname{Hom}$ is bilinear. On representations, $(\sum_c\lambda_c(\mathbf f'_c)^{\otimes n})\circ(\sum_l\mu_l\mathbf f_l^{\otimes n})=\sum_{c,l}\lambda_c\mu_l(\mathbf f'_c\circ\mathbf f_l)^{\otimes n}$, which has cost at most $\sum_{c,l}\abs{\lambda_c\mu_l}\norm{\mathbf f'_c}^n\norm{\mathbf f_l}^n$. So the norm is at most $1$, and the composition passes to the quotients by the null vectors. Associativity and units hold on powers. If $\mathfrak F$ is given, $\Phi=L_{\mathfrak F_{\mathbf V,\mathbf V'}}$ is multiplicative, as both sides are bilinear and agree on powers by functoriality. Conversely, if $\Phi$ is given, $\mathfrak F$ is functorial, and its hom maps are laws by Theorem~\ref{dp:thm:normed-roby}(2). The norms agree by Theorem~\ref{dp:thm:normed-roby}(3).
\end{proof}

\begin{proposition}[closure]\label{dp:prop:functor-closure}
Let $\mathfrak F,\mathfrak F'$ be bounded strict polynomial functors of type $I$ on $\cC$, homogeneous of degrees $n,n'$.
\begin{enumerate}
\item If $n=n'$, then $\mathfrak F\oplus\mathfrak F'$, with the maximum norm, is one of degree $n$.
\item If $\cC=\mathsf{Vec}_K$, then $\mathbf V\mapsto\mathfrak F(\mathbf V)\otimes_\pi\mathfrak F'(\mathbf V)$, separated, is one of degree $n+n'$. In all three cases its completion is one of degree $n+n'$ on $\cC$. Multidegrees add.
\item If $\mathfrak F''$ is one of type $\{1\}$, one covariant variable, homogeneous of degree $n''$, on one of the three categories, which contains the values of $\mathfrak F$ and is not $\mathsf{Vec}_K$ if $\cC\ne\mathsf{Vec}_K$, then $\mathfrak F''\circ\mathfrak F$ is one of degree $nn''$.
\item If $\mathfrak F_0(\mathbf V)\subseteq\mathfrak F(\mathbf V)$ are closed subspaces with $\mathfrak F(\mathbf f)\mathfrak F_0(\mathbf V)\subseteq\mathfrak F_0(\mathbf V')$ for all $\mathbf f$, then the closed subfunctor $\mathfrak F_0$, with the restricted norms, is one of degree $n$.
\item In the situation of~(4), the quotient $\mathfrak F/\mathfrak F_0$, with the quotient norms, is one of degree $n$.
\item In one covariant variable, the dual $\mathfrak F^\#(E)=\mathfrak F(E^*)^*$, $\mathfrak F^\#(u)=\mathfrak F(u^*)^*$, is one of degree $n$ on $\mathsf{Vec}_K$. If $K$ is complete, its values are Banach spaces, and it is one on $\mathsf{Ban}_K$ and on $\cC_{\rm um}$.
\end{enumerate}
If the given functors are contractive, so is the result of each of the operations \textup{(1)--(6)}. If $\mathfrak F,\mathfrak F',\mathfrak F''$ are DP vector functors, with no condition on degrees, the operations \textup{(1)--(6)} give DP vector functors.
\end{proposition}

\begin{proof}
\emph{Step 1 (items (1)--(3)).} Item (1) holds componentwise. With the maximum norm, its law has norm at most the larger of the two norms. For item (2), $s\mapsto\mathfrak F(\sum_cs_c\mathbf f_c)\otimes\mathfrak F'(\sum_cs_c\mathbf f_c)$ is a product of homogeneous polynomials with operator coefficients, and $\norm{T\otimes T'}\le\norm T\norm{T'}$ for $\pi$. So the law has norm at most the product of the two norms, and extension to the completions does not change this bound. For item (3), write $\mathfrak F(\sum_cs_c\mathbf f_c)=\sum_{\abs\gamma=n}s^\gamma T_\gamma$ with bounded operators $T_\gamma$. The law property of $\mathfrak F''$ on the vectors $T_\gamma$ makes $\mathfrak F''(\sum_\gamma z_\gamma T_\gamma)$ a homogeneous polynomial of degree $n''$ in the variables $z_\gamma$. Substitute $z_\gamma=s^\gamma$. The bound is the composite of the bounds of the two hom maps, so it is at most $1$ for contractive functors.

\emph{Step 2 (items (4) and (5)).} By Lemma~\ref{dp:lem:functor-rep}(1), the coefficients of the hom map of $\mathfrak F$ preserve $\mathfrak F_0$. So they induce operators on $\mathfrak F_0$ and on the quotients. Restriction does not increase operator norms. The quotient maps are metric surjections, so passing to quotients does not increase them either. So the laws of items (4) and (5) have norm at most that of the law of $\mathfrak F$.

\emph{Step 3 (item (6)).} The maps $u\mapsto u^*$ and $T\mapsto T^*$ are linear of norm at most $1$, so the hom map is a law by Lemma~\ref{dp:lem:elementary}(2). The same maps give the norm bound. Dual spaces over a complete field are complete. They are ultrametric, as $K$ is.

\emph{Step 4 (DP vector functors).} A pair of DP maps is DP, and $(T,T')\mapsto T\oplus T'$ is bounded linear. The map $(T,T')\mapsto T\otimes T'$ is the diagonal of a bounded bilinear map, so it is DP by Theorem~\ref{dp:thm:calculus}(3),(1), and composites of DP maps are DP by Theorem~\ref{dp:thm:calculus}(2). This gives items (1), (2), (3), (6). For item (4), $T\mapsto T|_{\mathfrak F_0(\mathbf V)}$ carries the hom map of $\mathfrak F$ to the hom map of $\mathfrak F_0$ followed by the inclusion of $\cL(\mathfrak F_0(\mathbf V),\mathfrak F_0(\mathbf V'))$ in $\cL(\mathfrak F_0(\mathbf V),\mathfrak F(\mathbf V'))$, a linear isometry with closed range. Theorem~\ref{dp:thm:target-descent}(2) applies. For item (5), with $q,q'$ the quotient maps at $\mathbf V,\mathbf V'$, the map $T\mapsto q'T$ carries the hom map of $\mathfrak F$ to the hom map of the quotient followed by $S\mapsto S\circ q$. This is a linear isometry into $\cL(\mathfrak F(\mathbf V),\mathfrak F(\mathbf V')/\mathfrak F_0(\mathbf V'))$, and its range, the operators vanishing on $\mathfrak F_0(\mathbf V)$, is closed. Theorem~\ref{dp:thm:target-descent}(2) applies again.
\end{proof}

\subsection{Functors on bundles}

Let $\mathcal M$ be a DP manifold and let $\mathbf M=(\mathsf M_i)_{i\in I}$ be DP vector bundles over $\mathcal M$, with fibers $G_i$ in $\cC$; put $\mathbf G=(G_i)_i$. Let $\mathsf t^i_x\colon G_i\to(\mathsf M_i)_x$ for $x\in U$ and $\mathsf t'^i_x\colon G_i\to(\mathsf M_i)_x$ for $x\in U'$ be local trivializations of all the $\mathsf M_i$. As in Corollary~\ref{dp:cor:bundles}, fibers carry norms up to equivalence, and $\mathfrak F(\mathbf M_x)$ is understood up to canonical isomorphism.

Let $\psi_x\in\operatorname{Hom}(\mathbf M_x,\mathbf G)$ have entries $(\mathsf t^i_x)^{-1}$ for $i\in I_+$ and $\mathsf t^i_x$ for $i\in I_-$, and define $\psi'_x$ in the same way. Then $\mathbf t(x)=\psi'_x\circ\psi_x^{-1}\in\operatorname{Hom}(\mathbf G,\mathbf G)$ has entries $(\mathsf t'^i_x)^{-1}\mathsf t^i_x$ for $i\in I_+$ and $(\mathsf t^i_x)^{-1}\mathsf t'^i_x$ for $i\in I_-$. These are transition families of the $\mathsf M_i$, so $x\mapsto\mathbf t(x)$ is DP on $U\cap U'$.

For a DP map $\mathsf w\colon\mathcal M\to\mathcal N$ and a family $\mathbf M'=(\mathsf M'_i)_i$ of DP vector bundles over $\mathcal N$ with fibers in $\cC$, a \emph{$\mathsf w$-morphism} $\mathsf M_i\to\mathsf M_i'$, respectively a \emph{$\mathsf w$-comorphism} $\mathsf M_i'\to\mathsf M_i$, is a family of linear maps $(\mathsf M_i)_x\to(\mathsf M_i')_{\mathsf w(x)}$, respectively $(\mathsf M_i')_{\mathsf w(x)}\to(\mathsf M_i)_x$, whose local operator-valued expressions are DP.

\begin{theorem}[polynomial functors act on DP bundles]\label{dp:thm:functors-dp}
\leavevmode
\begin{enumerate}
\item A bounded strict polynomial functor of finite degree on $\cC$ is a DP vector functor, with continuous hom maps.
\end{enumerate}
Let $\mathfrak F$ be a DP vector functor of type $I$ on $\cC$.
\begin{enumerate}
\setcounter{enumi}{1}
\item \textup{(n\textsuperscript{o}~7.6.2.)} There is exactly one structure of DP vector bundle on $\mathfrak F(\mathbf M)=\bigsqcup_x\mathfrak F(\mathbf M_x)$, with fiber $\mathfrak F(\mathbf G)$, for which every such family of trivializations gives the trivialization with inverse $\mathfrak F(\psi_x)$ over $U$. Its transition families are $x\mapsto\mathfrak F(\mathbf t(x))$.
\item \textup{(n\textsuperscript{o}~7.6.3.)} Let $\mathbf u=(u_i)_i$ consist of $\mathsf w$-morphisms $\mathsf M_i\to\mathsf M'_i$ for $i\in I_+$ and of $\mathsf w$-comorphisms $\mathsf M'_i\to\mathsf M_i$ for $i\in I_-$. There is exactly one bundle map $\mathfrak F(\mathbf u)\colon\mathfrak F(\mathbf M)\to\mathfrak F(\mathbf M')$ over $\mathsf w$ that is $\mathfrak F(\mathbf u_x)$ on the fiber at each $x$, and it is a $\mathsf w$-morphism. Also $\mathfrak F(\mathsf w^*\mathbf M')$ and $\mathsf w^*\mathfrak F(\mathbf M')$ are canonically isomorphic.
\item \textup{(n\textsuperscript{o}~7.6.4.)} Let $\mathfrak F_1,\dots,\mathfrak F_N$ be DP vector functors of type $I$ on $\cC$, and let $\theta$ be a family of bounded $N$\nobreakdash-linear maps $\theta_{\mathbf V}\colon\mathfrak F_1(\mathbf V)\times\dots\times\mathfrak F_N(\mathbf V)\to\mathfrak F(\mathbf V)$ with $\mathfrak F(\mathbf f)\circ\theta_{\mathbf V}=\theta_{\mathbf V'}\circ(\mathfrak F_1(\mathbf f)\times\dots\times\mathfrak F_N(\mathbf f))$. There is exactly one fiberwise $N$\nobreakdash-linear map $\theta_{\mathbf M}$ over $\id_{\mathcal M}$ equal to $\theta_{\mathbf M_x}$ on the fiber at each $x$. Its local operator-valued expressions are DP, and $\mathfrak F(\mathbf u)\circ\theta_{\mathbf M}=\theta_{\mathbf M'}\circ(\mathfrak F_1(\mathbf u)\times\dots\times\mathfrak F_N(\mathbf u))$.
\end{enumerate}
\end{theorem}

\begin{proof}
\emph{Step 1 (item (1)).} Each hom map is a finite sum of laws, which are continuous and DP on their domains by Theorem~\ref{dp:thm:calculus}(1). A finite sum of maps DP at a point is DP there, by adding the expansions.

\emph{Step 2 (item (2)).} By functoriality, $\mathfrak F(\psi'_x)\mathfrak F(\psi_x)^{-1}=\mathfrak F(\mathbf t(x))$, and this is DP in $x$ by Theorem~\ref{dp:thm:calculus}(2). These families give a DP atlas, as in Corollary~\ref{dp:cor:bundles}. An atlas determines the structure, which gives uniqueness.

\emph{Step 3 (item (3)).} A map on fibers is unique if it exists. In trivializations, the local expression of $\mathfrak F(\mathbf u)$ is $x\mapsto\mathfrak F(\mathbf u_{\rm loc}(x))$ with $\mathbf u_{\rm loc}$ DP, hence it is DP by Theorem~\ref{dp:thm:calculus}(2). The transitions of $\mathsf w^*\mathsf M_i'$ are those of $\mathsf M_i'$ composed with $\mathsf w$, so $\mathfrak F(\mathsf w^*\mathbf M')$ and $\mathsf w^*\mathfrak F(\mathbf M')$ have the same transition families.

\emph{Step 4 (item (4)).} In trivializations, $\theta_{\mathbf M}$ is the constant $\theta_{\mathbf G}$, which is DP. The definition does not depend on the trivializations, because $\mathfrak F(\mathbf t(x))\circ\theta_{\mathbf G}=\theta_{\mathbf G}\circ(\mathfrak F_1(\mathbf t(x))\times\dots\times\mathfrak F_N(\mathbf t(x)))$. The last identity holds on each fiber by the condition on $\theta$.
\end{proof}

\subsection{Analytic vector functors}

\begin{proposition}[closure of the analytic class]\label{dp:prop:analytic-closure}
\leavevmode
\begin{enumerate}
\item Analytic vector functors are closed under finite direct sums, the tensor product of Proposition~\ref{dp:prop:functor-closure}\textup{(2)}, the composition of Proposition~\ref{dp:prop:functor-closure}\textup{(3)}, the duals of Proposition~\ref{dp:prop:functor-closure}\textup{(6)} and bounded natural retracts: if $\mathfrak F'$ is analytic and $\theta_{\rm in}\colon\mathfrak F\to\mathfrak F'$, $\theta_{\rm out}\colon\mathfrak F'\to\mathfrak F$ are natural with bounded components and $\theta_{\rm out}\theta_{\rm in}=\id$, then $\mathfrak F$ is analytic.
\item Let $\operatorname{char}K=p>0$ and $k\ge p$. They are not closed under closed subfunctors or under quotients. If $K$ is complete and not spherically complete, then on $\cC_{\rm um}$, with ultrametric values, $\Mult^k$ is analytic and its closed subfunctor $\alpha_k$ is not. Over every $K$ of characteristic $p$, on $\mathsf{Vec}_K$ and on $\mathsf{Ban}_K$, $\Mult^k$ is analytic and $\alpha_k$ is not, and $\hot^k_\pi$ is analytic and its quotient $\hLambda^k_\pi$ is not.
\end{enumerate}
\end{proposition}

\begin{proof}
\emph{Step 1 (item (1)).} A pair of analytic maps is analytic. The map $(T,T')\mapsto T\otimes T'$ is the diagonal of a bounded bilinear map. Analytic germs compose, by~\cite[proof of Proposition~14.26, Step~1]{I}. That step assumes a complete final target. Apply it in the completion of the final target: the multilinear coefficients of the composite are finite sums of maps with values in the target itself, and the composite has values there, so the expansion holds in the target. This gives sums, tensor products, composition and duals. For a retract, naturality gives $\mathfrak F(\mathbf f)=\theta_{\rm out}\mathfrak F'(\mathbf f)\theta_{\rm in}$, a bounded linear image of $\mathfrak F'(\mathbf f)$.

\emph{Step 2 (item (2)).} If $\alpha_k$ were analytic, its composite with the bounded affine map $u\mapsto(\id_F,u)$, which is $A^k_{E,D;F}$, would be analytic. Lemma~\ref{dp:fn:lem:witnesses} excludes this, for $F=K$ in the first case. The quotient is Lemma~\ref{dp:fn:lem:witnesses}(1).
\end{proof}

By Proposition~\ref{dp:prop:functor-closure} the DP vector functors are closed under finite direct sums, tensor products, composition, duals, closed subfunctors and quotients. They are also closed under bounded natural retracts, as $\mathfrak F(\mathbf f)=\theta_{\rm out}\mathfrak F'(\mathbf f)\theta_{\rm in}$ is a bounded linear image of $\mathfrak F'(\mathbf f)$.

\begin{theorem}[multidegree threshold]\label{dp:thm:functor-threshold}
Let $\mathbf d=(d_i)_{i\in I}\in\N^I$ and $n=\abs{\mathbf d}$.
\begin{enumerate}
\item Let $\prod_id_i!\ne0$ in $K$, and let $\mathfrak F$ be a bounded strict polynomial functor of multidegree $\mathbf d$ on $\cC$. Let $a=\mathfrak F_{\mathbf V,\mathbf V'}$, and let $\iota_i$ be the inclusion of the $i$th entry into $\operatorname{Hom}(\mathbf V,\mathbf V')$. Then $a$ is the diagonal of the bounded $n$\nobreakdash-linear map
\begin{equation}\label{dp:fn:eq:lift}
 Q\bigl((f_{i,c})_{i\in I,\,1\le c\le d_i}\bigr)=\Bigl(\prod_id_i!\Bigr)^{-1}a^{[1^n]}\bigl((\iota_if_{i,c})_{i,c}\bigr),\qquad
 \norm Q\le\frac{(C_1n)^n\norm a}{\abs{\prod_id_i!}}\quad(\text{with }0^0=1).
\end{equation}
So $\mathfrak F$ is an analytic vector functor.
\item If $\prod_id_i!=0$ in $K$, there is a contractive bounded strict polynomial functor of multidegree $\mathbf d$, on $\mathsf{Vec}_K$ and on $\mathsf{Ban}_K$, whose hom map on some $\operatorname{Hom}(\mathbf V,\mathbf V)$, with all $V_i$ Banach, is analytic at no point.
\end{enumerate}
So all bounded strict polynomial functors of multidegree $\mathbf d$, on $\mathsf{Vec}_K$ or on $\mathsf{Ban}_K$, are analytic if and only if $\prod_id_i!\ne0$ in $K$. For $\alpha_k$, of multidegree $(1,k)$, the condition is $k!\ne0$.
\end{theorem}

In one variable the algebra of~\eqref{dp:fn:eq:lift} is~\cite[Th.~II.1]{Ro}; the bound is the estimate we need.

\begin{proof}
\emph{Step 1 (item (1)).} If $n=0$, then $a$ is constant. Let $n\ge1$. By Lemma~\ref{dp:lem:elementary}(4), $a^{[1^n]}$ is $n$\nobreakdash-linear, and~\eqref{dp:laws:scaled} with all exponents equal to $1$ gives the bound. On the diagonal $f_{i,c}=f_i$, the value $a^{[1^n]}((\iota_if_i)_{i,c})$ is the coefficient of $\prod_{i,c}s_{i,c}$ in $a((\sum_cs_{i,c}f_i)_i)=\prod_i(\sum_cs_{i,c})^{d_i}a(\mathbf f)$, which is $\prod_id_i!\,a(\mathbf f)$. Precompose $Q$ with the projection of the slot $(i,c)$ of $\operatorname{Hom}(\mathbf V,\mathbf V')^n$ onto the $i$th entry, which has norm at most $1$. This gives a bounded $n$\nobreakdash-linear map on $\operatorname{Hom}(\mathbf V,\mathbf V')^n$, of norm at most $\norm Q$, with diagonal $a$. Lemma~\ref{dp:lem:lift-criterion} gives analyticity at every point.

\emph{Step 2 (the functor for item (2)).} Let $p=\operatorname{char}K$ and choose $i_0$ with $d_{i_0}\ge p$. Let $E,F$ be the Banach spaces of~\cite[Theorem~7.2(1)]{I} for $k=d_{i_0}$. Put $\mathfrak H(V)=\hLambda^{d_{i_0}}_\pi V$ if $i_0\in I_+$, and $\mathfrak H(V)=\Alt^{d_{i_0}}(V;F)$ if $i_0\in I_-$, with $\mathfrak H(u)(\zeta)=\zeta\circ(u\times\dots\times u)$. For $j\ne i_0$ with $d_j\ge1$, put $\rho_j(V)=\hot^{d_j}_\pi V$ if $j\in I_+$ and $\rho_j(V)=\Mult^{d_j}(V;\Kh)$ if $j\in I_-$, with the action $\zeta\mapsto\zeta\circ(u\times\dots\times u)$. If $d_j=0$, put $\rho_j(V)=\Kh$, with the identity as the action on maps. Each is contractive, by the examples above and Theorem~\ref{dp:thm:calculus}(4). Composed with the projection of $\operatorname{Hom}$ onto its entry, each factor has degree $d_j$ in its own variable and $0$ in the others. By Proposition~\ref{dp:prop:functor-closure}(2),
\[
 \mathfrak F_{\mathbf d}(\mathbf V)=\mathfrak H(V_{i_0})\mathbin{\widehat\otimes_\pi}\rho_{j_1}(V_{j_1})\mathbin{\widehat\otimes_\pi}\dots\mathbin{\widehat\otimes_\pi}\rho_{j_{\abs I-1}}(V_{j_{\abs I-1}})
\]
is contractive of multidegree $\mathbf d$, where $j_1,\dots,j_{\abs I-1}$ are the elements of $I$ other than $i_0$, and the product is bracketed from the right, so that the first factor is $\mathfrak H(V_{i_0})$ and the second is the completed tensor product of the others. If $I=\{i_0\}$, then $\mathfrak F_{\mathbf d}=\mathfrak H$.

\emph{Step 3 (recovering $\mathfrak H$).} Let $V_{i_0}=E$ and $V_j=\Kh$ for $j\ne i_0$, a Banach space over $K$. If $I=\{i_0\}$, Step~4 holds with $\Psi=\id$. Let $I\ne\{i_0\}$. Let $X=\mathfrak H(E)$, a Banach space, and let $Y$ be the completed tensor product of the $\rho_j(\Kh)$, bracketed in the same way, so that $\mathfrak F_{\mathbf d}(\mathbf V)=X\mathbin{\widehat\otimes_\pi}Y$. A Banach space over $K$ is a $\Kh$-space, with $\norm{\lambda\xi}=\abs\lambda\norm\xi$ for $\lambda\in\Kh$. Multiplication of the entries of $\hot^{d_j}_\pi\Kh$, evaluation of $\Mult^{d_j}(\Kh;\Kh)$ at $(1,\dots,1)$, and the identity of $\Kh$ have norm at most $1$, by the triangle inequality and multiplicativity in $\Kh$. Their product is a linear map $\mathfrak m\colon Y\to\Kh$ of norm at most $1$. It takes the value $1$ at a vector $y\in Y$ of norm at most $1$: the tensor of the vectors $1^{\otimes d_j}$, of the multiplication in $\Mult^{d_j}(\Kh;\Kh)$, and of $1$. Then $\xi\otimes y'\mapsto\mathfrak m(y')\xi$ extends to a contraction $\mathsf r\colon X\mathbin{\widehat\otimes_\pi}Y\to X$. The map $\Psi(T_1)(\xi)=\mathsf r(T_1(\xi\otimes y))$ is linear of norm at most $1$ from $\cL(X\mathbin{\widehat\otimes_\pi}Y,X\mathbin{\widehat\otimes_\pi}Y)$ to $\cL(X,X)$, and $\Psi(T\mathbin{\widehat\otimes}\id_Y)=T$.

\emph{Step 4 (item (2)).} Let $\mathbf f_u\in\operatorname{Hom}(\mathbf V,\mathbf V)$ have entry $u$ at $i_0$ and $\id$ at the other entries. Then $\Psi(\mathfrak F_{\mathbf d}(\mathbf f_u))=\mathfrak H(u)$, and $u\mapsto\mathbf f_u$ is bounded affine. If the hom map of $\mathfrak F_{\mathbf d}$ were analytic at one point, it would be analytic everywhere by Lemma~\ref{dp:lem:lift-criterion}. Then $u\mapsto\mathfrak H(u)$ would be analytic on $\cL(E,E)$. This map is $\hLambda^{d_{i_0}}_\pi u$ if $i_0\in I_+$ and $A^{d_{i_0}}_{E,E;F}(u)$ if $i_0\in I_-$, which contradicts Lemma~\ref{dp:fn:lem:witnesses}(1).
\end{proof}

\begin{proposition}[the representing functor is not ambient]\label{dp:prop:gamma-not-ambient}
Let $n\ge1$, and let $\mathfrak G$ be $\Gamma^n_{\mathrm s}$ or $\hGamma^n$.
\begin{enumerate}
\item $\mathfrak G$ is a DP vector functor on $\mathsf{Vec}_K$, and $\hGamma^n$ is one on $\mathsf{Ban}_K$.
\item Let $p=\operatorname{char}K>0$, $n\ge p$, and let $D=K\oplus\mathcal E$ be the space of Proposition~\ref{dp:prop:degree-p}\textup{(2)}. The hom map of $\mathfrak G$ on $\cL(K,D)$ is not ambient of degree $n$. It is AA at no point, so it is analytic at no point. If $K$ is complete, $K$ and $D$ are Banach spaces, so for $\mathfrak G=\hGamma^n$ this happens on $\mathsf{Ban}_K$ too.
\item In the situation of~\textup{(2)}, $\mathfrak G$ is not a subquotient of an analytic vector functor: there do not exist an analytic vector functor $\widetilde{\mathfrak F}$, closed subfunctors $\mathfrak F_2\subseteq\mathfrak F_1\subseteq\widetilde{\mathfrak F}$, and natural bounded isomorphisms from $\mathfrak G$ to $\mathfrak F_1/\mathfrak F_2$ with bounded inverses, on a category containing $K$ and $D$.
\end{enumerate}
\end{proposition}

\begin{proof}
\emph{Step 1 (item (1)).} This is Lemma~\ref{dp:lem:functor-rep}(2) and Theorem~\ref{dp:thm:functors-dp}(1).

\emph{Step 2 (item (2)).} The multiplication of $K^{\otimes n}$ gives $\abs{\sum_l\lambda_lh_l^n}\le\sum_l\abs{\lambda_l}\abs{h_l}^n$, so $\norm{\lambda\,1^{\otimes n}}_\Gamma=\abs\lambda$ for $\lambda\in K$. Thus $\Gamma^n_{\mathrm s}(K)=K$ and $\hGamma^n(K)=\Kh$, with $[1^{\otimes n}]$ corresponding to $1$. So evaluation at $[1^{\otimes n}]$ is an isometric isomorphism of $\cL(\mathfrak G(K),\mathfrak G(D))$ onto $\mathfrak G(D)$; for $\mathfrak G=\hGamma^n$ this uses that a bounded $K$-linear map $T$ from $\Kh$ to a Banach space has $T(\lambda)=\lambda T(1)$ for $\lambda\in\Kh$. It sends $\mathfrak G(u)$ to $\mathsf v_n(u(1))$, by the formula of Lemma~\ref{dp:lem:functor-rep}(2), read in $\hGamma^n(D)$ if $\mathfrak G=\hGamma^n$. The inclusion of $\Gamma^n_{\mathrm s}(D)$ in $\hGamma^n(D)$ is isometric. Evaluation at $1$ is an isometric isomorphism of $\cL(K,D)$ onto $D$. So the ambient constant of the hom map equals that of $\mathsf v_n$, which is $\sup\{\norm t_\Gamma/\pi(t):t\in\Gamma^n(D)\}=\mathfrak a_n(D)$. This is infinite by Proposition~\ref{dp:prop:all-norms}(2). So the hom map is not ambient of degree $n$. If it were AA at a point, then by~\cite[Proposition~14.9(2)]{I} it would be ambient of degree $n$.

\emph{Step 3 (item (3)).} Suppose such data exist. Let $q\colon\mathfrak F_1(K)\to\mathfrak F_1(K)/\mathfrak F_2(K)$ and $q'\colon\widetilde{\mathfrak F}(D)\to\widetilde{\mathfrak F}(D)/\mathfrak F_2(D)$ be the quotient maps. The map $q$ is a metric surjection, and $\mathfrak F_1(D)/\mathfrak F_2(D)$ is a normed subspace of $\widetilde{\mathfrak F}(D)/\mathfrak F_2(D)$, as the two quotient norms are the same infimum. So $T\mapsto T\circ q$ is a linear isometry of $\cL(\mathfrak F_1(K)/\mathfrak F_2(K),\mathfrak F_1(D)/\mathfrak F_2(D))$ into $\cL(\mathfrak F_1(K),\widetilde{\mathfrak F}(D)/\mathfrak F_2(D))$. It sends the hom map of $\mathfrak F_1/\mathfrak F_2$ to $u\mapsto q'\circ\widetilde{\mathfrak F}(u)|_{\mathfrak F_1(K)}$, a bounded linear image of an analytic map. The hom map of $\mathfrak F_1/\mathfrak F_2$ is that of $\mathfrak G$ followed by a bounded linear bijection $\mathsf c$ with bounded inverse, so it is a law. By~\cite[Proposition~14.9(2)]{I}, applied to the isometry $T\mapsto T\circ q$, it is ambient. As $L_{\mathsf c\circ a}=\mathsf c\circ L_a$ for a law $a$, the hom map of $\mathfrak G$ has finite ambient constant too, against item~(2).
\end{proof}

\begin{theorem}[ultrametric spaces and values]\label{dp:thm:functors-um}
Let $K$ be complete and let $\mathfrak F$ be a vector functor of type $I$ on $\cC_{\rm um}$.
\begin{enumerate}
\item If every value of $\mathfrak F$ has an equivalent norm that is ultrametric and spherically complete, then $\mathfrak F$ is a DP vector functor if and only if it is an analytic vector functor.
\item If $K$ is discretely valued and every value of $\mathfrak F$ has an equivalent ultrametric norm, then $\mathfrak F$ is a DP vector functor if and only if it is an analytic vector functor.
\end{enumerate}
\end{theorem}

\begin{proof}
\emph{Step 1 (both items).} Analytic maps are DP by Theorem~\ref{dp:thm:calculus}(3). Conversely, $\operatorname{Hom}(\mathbf V,\mathbf V')$ is ultrametric, since an operator norm into an ultrametric space is ultrametric.

\emph{Step 2 (item (1)).} Give $\mathfrak F(\mathbf V)$ and $\mathfrak F(\mathbf V')$ the equivalent norms. The resulting operator norm on $\cL(\mathfrak F(\mathbf V),\mathfrak F(\mathbf V'))$ is equivalent to the given one, and it is ultrametric and spherically complete by Lemma~\ref{dp:lem:operator-sph}. Apply Theorem~\ref{dp:thm:sph-agreement}(1) at every point.

\emph{Step 3 (item (2)).} Give $\mathfrak F(\mathbf V)$ and $\mathfrak F(\mathbf V')$ the equivalent ultrametric norms. The resulting operator norm is ultrametric, and the operator space is Banach, as the values are. So Theorem~\ref{dp:thm:sph-agreement}(2) applies.
\end{proof}

\subsection{The statements of \texorpdfstring{n\textsuperscript{o}~7.8}{no. 7.8}}

Bourbaki also says that $\Alt^k(\mathsf M_1;\mathsf M_0)$ is a vector subbundle of $\Mult^k(\mathsf M_1;\mathsf M_0)$~\cite[n\textsuperscript{o}~7.8.1]{Bou}, which asks for closed subspaces with a topological complement in suitable charts~\cite[n\textsuperscript{o}~7.5.1]{Bou}. We use the same definition for DP vector bundles, with their local trivializations as charts. We say that $\Alt^k(E;F)$ is \emph{complemented} if there is a bounded linear projection $\operatorname{Pr}$ of $\Mult^k(E;F)$ onto $\Alt^k(E;F)$.

\begin{proposition}[the subbundle sentence]\label{dp:prop:subbundle}
Let $k\ge1$ and let $E,F$ be normed spaces.
\begin{enumerate}
\item If $\Alt^k(E;F)$ is complemented, then $A^k_{E,D;F}$ is the diagonal of $(u_1,\dots,u_k)\mapsto(\zeta\mapsto\operatorname{Pr}(\zeta\circ(u_1\times\dots\times u_k)))$ for every normed $D$~\cite[Proposition~5.1 and its proof]{I}.
\item Let $\operatorname{char}K=p$ and $k\ge p$. Then $\Alt^k(E;F)$ is not complemented for the Banach spaces $E,F$ of~\cite[Theorem~7.2(1)]{I}. If moreover $K$ is complete and not spherically complete, it is not complemented for the ultrametric Banach space $E$ of~\cite[Theorem~7.2(2)]{I} and $F=K$. Over every $K$, $\Alt^k(E;F)$ is complemented for all Banach $E,F$ if and only if $k!\ne0$ in $K$.
\item It is complemented in each of the following cases.
\begin{enumerate}
\item $k!\ne0$ in $K$.
\item $E$ has finite coordinates.
\item $K$ is arbitrary, $E$ is finite-dimensional and $F$ is Banach.
\item $E$ is arbitrary and $F$ has an equivalent norm that is ultrametric and spherically complete.
\item $E$ is arbitrary, $K$ is discretely valued and $F$ is Banach with an equivalent ultrametric norm.
\end{enumerate}
\item Let $\mathsf M_1,\mathsf M_0$ be DP vector bundles over a nonempty DP manifold, with fibers $G_1,G_0$. Then $\Alt^k(\mathsf M_1;\mathsf M_0)$ is a vector subbundle of $\Mult^k(\mathsf M_1;\mathsf M_0)$ if and only if $\Alt^k(G_1;G_0)$ is complemented.
\end{enumerate}
\end{proposition}

\begin{proof}
\emph{Step 1 (item (2)).} By item (1), a projection would make $A^k_{E,E;F}$, respectively $A^k_{E,D;K}$, the diagonal of a bounded $k$\nobreakdash-linear map, hence analytic, against Lemma~\ref{dp:fn:lem:witnesses}. If $k!\ne0$, item (3)(a) applies; if $k!=0$, then $\operatorname{char}K=p\le k$ and the first case applies.

\emph{Step 2 (item (3)(a),(b)).} For (a), take $\operatorname{Pr}(\beta)=\frac1{k!}\sum_{\nu\in S_k}\sgn(\nu)\,\beta\circ\nu$. For (b), let $\mathsf T\colon K^r\to E$ be a bounded linear bijection with bounded inverse, let $e_c$ be the images of the unit vectors, and let $e^*_c$ be the coordinate forms. Put
\[
 \operatorname{Pr}(\beta)(v_1,\dots,v_k)=\sum_{c_1<\dots<c_k}\det\bigl(e^*_{c_{\mathsf b}}(v_{\mathsf a})\bigr)_{\mathsf a,\mathsf b}\;\beta(e_{c_1},\dots,e_{c_k}).
\]
Two equal $v_{\mathsf a}$ give two equal rows, so $\operatorname{Pr}(\beta)$ is alternating. For alternating $\beta$, expanding by multilinearity gives $\operatorname{Pr}(\beta)=\beta$: terms with a repeated index vanish, and permuting the arguments changes the sign, in every characteristic.

\emph{Step 3 (item (3)(c)).} Let $\widehat E$ be the completion of $E$. The $\Kh$-span of $E$ in $\widehat E$ is finite-dimensional over $\Kh$, hence complete by~\cite[Lemma~A.3]{I}, hence equal to $\widehat E$. So $\widehat E$ has finite coordinates over $\Kh$, by the same lemma. By the completion principle of~\cite[text before Lemma~H.5]{I}, restriction is an isometric bijection from the bounded $\Kh$-multilinear maps $\widehat E^k\to F$ onto $\Mult^k(E;F)$, and by continuity it maps alternating maps onto alternating maps. Transport the projection of~(b).

\emph{Step 4 (item (3)(d),(e)).} For (d),~\cite[Theorem~5.2]{I} gives a contractive retraction for the equivalent norm; transport it to the given norm. For (e), rounding the equivalent norm as in~\cite[proof of Corollary~5.3]{I} gives an equivalent norm as in (d).

\emph{Step 5 (item (4)).} A chart of $\Mult^k(\mathsf M_1;\mathsf M_0)$ induced by trivializations of $\mathsf M_1,\mathsf M_0$ preserves alternating maps. So it maps the subset $\Alt^k(\mathsf M_1;\mathsf M_0)$ onto $U\times\Alt^k(G_1;G_0)$, and a complement gives the subbundle property. Conversely, if a chart at a point of the nonempty base maps this subset over $U$ onto $U\times\mathsf F_0$ with $\mathsf F_0$ complemented, then at a point of $U$ the fiber of $\Alt^k$ is complemented in the fiber of $\Mult^k$. Trivializations of $\mathsf M_1,\mathsf M_0$ identify these fibers with $\Alt^k(G_1;G_0)$ and $\Mult^k(G_1;G_0)$.
\end{proof}

\begin{theorem}[n\textsuperscript{o}~7.8 and DP]\label{dp:thm:seven-eight-dp}
Let $k\ge1$.
\begin{enumerate}
\item \textup{(DP, every characteristic, every $K$.)} On $\mathsf{Vec}_K$, on $\mathsf{Ban}_K$ and on $\cC_{\rm um}$, $\alpha_k$ is a DP vector functor, as is every bounded strict polynomial functor of finite degree, in particular $\hot^k_\pi$, $\hLambda^k_\pi$ and $\hGamma^k$. Theorem~\ref{dp:thm:functors-dp} holds for them. For DP vector bundles $\mathsf M_1,\mathsf M_0$ over a DP manifold $\mathcal M$, $\Alt^k(\mathsf M_1;\mathsf M_0)$ is a DP vector bundle and its inclusion into $\Mult^k(\mathsf M_1;\mathsf M_0)$ is a morphism. Let $k'\ge1$, and let $\mathsf M,\mathsf N,\mathsf N',\mathsf N''$ be DP vector bundles over $\mathcal M$ with fibers $G_1,G,G',G''$. Let $\varphi\colon\mathsf N\times_{\mathcal M}\mathsf N'\to\mathsf N''$ be a pairing whose local expression $x\mapsto\varphi_x$, with values in the bounded bilinear maps $G\times G'\to G''$, is DP. For $\omega$ in a fiber of $\Alt^k(\mathsf M;\mathsf N)$, $\omega'$ in the fiber of $\Alt^{k'}(\mathsf M;\mathsf N')$ at the same point, and $v_1,\dots,v_{k+k'}$ in the fiber of $\mathsf M$ there, the exterior product of~\cite[n\textsuperscript{o}~7.8.2]{Bou} is
\begin{equation}\label{dp:fn:eq:exterior}
 (\omega\wedge_\varphi\omega')(v_1,\dots,v_{k+k'})=\sum_\nu\sgn(\nu)\,\varphi\bigl(\omega(v_{\nu(1)},\dots,v_{\nu(k)}),\omega'(v_{\nu(k+1)},\dots,v_{\nu(k+k')})\bigr),
\end{equation}
summed over the permutations $\nu$ increasing on $\{1,\dots,k\}$ and on $\{k+1,\dots,k+k'\}$. It is a fiberwise bilinear map from $\Alt^k(\mathsf M;\mathsf N)\times_{\mathcal M}\Alt^{k'}(\mathsf M;\mathsf N')$ to $\Alt^{k+k'}(\mathsf M;\mathsf N'')$ with DP local expression.
\item \textup{(Analytic, all normed or all Banach spaces, every $K$.)} The functors $\alpha_k$ and $\hLambda^k_\pi$ are analytic if and only if $k!\ne0$ in $K$, that is, $\operatorname{char}K=0$ or $\operatorname{char}K>k$. The functor $\hot^k_\pi$ is analytic for every $k$. On $\mathsf{Vec}_K$, the same threshold holds for $\Gamma^k_{\mathrm s}$ and $\hGamma^k$. Theorem~\ref{dp:thm:functor-threshold} gives the threshold for all functors of a given multidegree.
\item \textup{(Ultrametric spaces, $K$ complete.)}
\begin{enumerate}
\item If $K$ is discretely valued, every DP vector functor on $\cC_{\rm um}$ with ultrametric values is analytic, in particular $\alpha_k$ in every degree. So in characteristic $p$ the condition $k!\ne0$ is sufficient and not necessary here.
\item If $\operatorname{char}K=p$, $K$ is not spherically complete and $k\ge p$, then $\alpha_k$ is DP and not analytic on $\cC_{\rm um}$. If $\operatorname{char}K=0$, every $\alpha_k$ is analytic.
\item If $\operatorname{char}K=p$, $K$ is densely valued and spherically complete and $k\ge p$, we do not decide whether $\alpha_k$ is analytic on $\cC_{\rm um}$ (\cite[Section~12]{I}; see also Question~\ref{dp:q:residual}). Every vector functor on $\cC_{\rm um}$ whose values have equivalent ultrametric spherically complete norms is DP if and only if it is analytic. For ultrametric Banach spaces $E,D,F,F'$ with $E$ or $D$ of countable type, the hom map of $\alpha_k$ on $\cL(F,F')\times\cL(E,D)$ is the diagonal of a bounded $(k+1)$\nobreakdash-linear map, so it is analytic, in every degree~\cite[Corollary~10.5]{I}.
\end{enumerate}
\item \textup{(Finite rank.)} Let $\cC_{\rm fd}$ be the full subcategory of finite-dimensional spaces in $\mathsf{Vec}_K$, and read Definition~\ref{dp:def:poly-functor} on it. Over every $K$, a vector functor on $\cC_{\rm fd}$ with Banach values is a DP vector functor if and only if it is an analytic vector functor. Over complete $K$ all finite-dimensional normed spaces are Banach, so $\alpha_k$ is analytic on $\cC_{\rm fd}$ in every characteristic. Over $K=\F_p(t)$ with the $t$-adic absolute value, for $k\ge p$, on finite-dimensional ultrametric spaces, where the values are finite-dimensional and incomplete, it can fail to be analytic.
\item \textup{(The subbundle sentence.)} For DP vector bundles $\mathsf M_1,\mathsf M_0$ over a nonempty DP manifold, with fibers $G_1,G_0$, the bundle $\Alt^k(\mathsf M_1;\mathsf M_0)$ is a vector subbundle of $\Mult^k(\mathsf M_1;\mathsf M_0)$ if and only if $\Alt^k(G_1;G_0)$ is complemented. This holds in the cases of Proposition~\ref{dp:prop:subbundle}\textup{(3)}, including $G_1$ finite-dimensional and $G_0$ Banach over every $K$. For the spaces of Proposition~\ref{dp:prop:subbundle}\textup{(2)}, $\Alt^k(\mathsf M_1;\mathsf M_0)$ is a DP vector bundle and not a vector subbundle.
\end{enumerate}
\end{theorem}

\begin{proof}
\emph{Step 1 (item (1), bundles).} By the examples above, Lemma~\ref{dp:lem:functor-rep}(2) and Proposition~\ref{dp:prop:functor-closure}, these functors are bounded strict polynomial functors on $\mathsf{Vec}_K$ and on $\mathsf{Ban}_K$, and on the full subcategory $\cC_{\rm um}$ by restriction. So Theorem~\ref{dp:thm:functors-dp}(1) applies. Both $\Mult^k$ and $\alpha_k$ are DP vector functors, and the inclusion is a natural bounded map, so Theorem~\ref{dp:thm:functors-dp}(4) with $N=1$ gives the morphism.

\emph{Step 2 (item (1), exterior product).} For a bounded bilinear map $\varphi_0\colon G\times G'\to G''$, formula~\eqref{dp:fn:eq:exterior} defines a bounded bilinear map $\mathsf P_{\varphi_0}$ from $\Alt^k(G_1;G)\times\Alt^{k'}(G_1;G')$ to $\Mult^{k+k'}(G_1;G'')$. Its values are alternating in every characteristic, as in~\cite[proof of Theorem~10.7, Step~2]{I}: terms with two equal arguments in one block vanish, and the others cancel in pairs. The map $\varphi_0\mapsto\mathsf P_{\varphi_0}$ is linear, of norm at most $\binom{k+k'}k$. In trivializations, the exterior product is $x\mapsto\mathsf P_{\varphi_x}$, a bounded linear image of a DP family.

\emph{Step 3 (item (2)).} For $\alpha_k$ this is~\cite[Corollary~7.3]{I}. For $\hLambda^k_\pi$, $\Gamma^k_{\mathrm s}$ and $\hGamma^k$ with $k!\ne0$, use Theorem~\ref{dp:thm:functor-threshold}(1). With $k!=0$, use Lemma~\ref{dp:fn:lem:witnesses}(1) and Proposition~\ref{dp:prop:gamma-not-ambient}(2), since $\operatorname{char}K\le k$. The examples above give $\hot^k_\pi$.

\emph{Step 4 (item (3)).} Item (a) is Theorem~\ref{dp:thm:functors-um}(2), as $\Alt^k(V_1;V_0)$ is an ultrametric Banach space when $V_0$ is. Item (b) is Proposition~\ref{dp:prop:analytic-closure}(2) with item (1), and Theorem~\ref{dp:thm:functor-threshold}(1) for characteristic $0$. Item (c) is Theorem~\ref{dp:thm:functors-um}(1) and the citation.

\emph{Step 5 (item (4)).} The hom spaces are finite-dimensional and the targets are Banach, so Theorem~\ref{dp:thm:agreement-finite}(2) applies at every point of every hom space (on the hom space $\{0\}$ every map is constant). Analytic maps are DP. Over complete $K$, the values of $\alpha_k$ on $\cC_{\rm fd}$ are finite-dimensional, so Banach by~\cite[Lemma~A.3]{I}, and $\alpha_k$ is DP by item (1). The last sentence is Theorem~\ref{dp:thm:agreement-finite}(3).

\emph{Step 6 (item (5)).} This is Proposition~\ref{dp:prop:subbundle}(4),(3),(2) with item (1).
\end{proof}
%% ---- end sec-12-functors.tex
%% ---- begin sec-13-jets.tex
\section{Jets}\label{dp:sec:jets}

The higher-order data of a DP map are its Hasse derivatives, by Theorem~\ref{dp:thm:quasi-analytic}(1) and Proposition~\ref{dp:prop:frechet-hasse}. This section assembles them into jet bundles over DP manifolds, in every characteristic and for all normed models. It then asks when the jet bundle of an analytic manifold is analytic. The answer is governed by a bounded linear section of the diagonal $\Mult^k(P;W)\to\cP_k(P;W)$.

\subsection{The composition rule as a law}\label{dp:ssec:chain-law}

In this section $P,P',W,Y$ are normed spaces and $r\ge1$. Put
\[
 \mathbb J_r(P;W)=\prod_{n=0}^r\cP_n(P;W),\qquad \mathbb J^-_r(P;W)=\prod_{n=1}^r\cP_n(P;W),
\]
with the maximum norm; here $\cP_0(P;W)=W$. For $a\in\mathbb J_r(P;W)$, $a^-\in\mathbb J^-_r(P;W)$ drops the component of degree $0$. For $\alpha\in\N^r$ we write $\deg\alpha=\sum_nn\alpha_n$.

Let $\alpha\in\N^r$ with $\deg\alpha=d\ge1$ and $\abs\alpha=m$. For $b\in\cP_m(W;Y)$ and $\ell=(\ell_n)_n\in\mathbb J^-_r(P;W)$, let $\mathcal F_\alpha(b,\ell)$ be the map $h\mapsto b^{[\alpha]}(\ell_1(h),\dots,\ell_r(h))$. Define
\[
 \Theta_r\colon\mathbb J_r(W;Y)\times\mathbb J^-_r(P;W)\to\mathbb J_r(P;Y)
\]
by
\[
 \Theta_r(b,\ell)_0=b_0,\qquad\Theta_r(b,\ell)_d=\sum_{\deg\alpha=d}\mathcal F_\alpha(b_{\abs\alpha},\ell)\quad(1\le d\le r).
\]

\begin{lemma}[the composition rule is a law]\label{dp:lem:chain-law}
\leavevmode
\begin{enumerate}
\item \textup{(Refinement.)} Let $b\in\cP_m(W;Y)$, $\alpha\in\N^q$ with $\abs\alpha=m$, let $I_1,\dots,I_q$ be finite sets, and let $w_{i,l}\in W$, $s_{i,l}\in K$ for $1\le i\le q$, $l\in I_i$. Then
\begin{equation}\label{dp:13:eq:refine}
 b^{[\alpha]}\Bigl(\sum_{l\in I_1}s_{1,l}w_{1,l},\dots,\sum_{l\in I_q}s_{q,l}w_{q,l}\Bigr)=\sum_\delta\prod_{i,l}s_{i,l}^{\delta_{i,l}}\;b^{[\delta]}\bigl((w_{i,l})_{i,l}\bigr),
\end{equation}
the sum over the arrays $\delta=(\delta_{i,l})$ of integers $\delta_{i,l}\ge0$ with $\sum_l\delta_{i,l}=\alpha_i$ for every $i$.
\item If moreover $\ell_i\in\cP_{d_i}(P;W)$ with $d_i\ge0$, then $h\mapsto b^{[\alpha]}(\ell_1(h),\dots,\ell_q(h))$ lies in $\cP_d(P;Y)$ with $d=\sum_i\alpha_id_i$, and its norm is at most $C_0^m\frac{m!}{\alpha!}\norm b\prod_i\norm{\ell_i}^{\alpha_i}$.
\item Let $\alpha\in\N^r$ with $\deg\alpha=d\ge1$, $\abs\alpha=m$. The map $\mathcal F_\alpha$ is a law of degree $1+m$ from $\cP_m(W;Y)\times\mathbb J^-_r(P;W)$ to $\cP_d(P;Y)$, of norm at most $C_0^mm!/\alpha!$, linear in $b$. The map $\ell\mapsto\mathcal F_\alpha({\cdot},\ell)$ is a law of degree $m$ from $\mathbb J^-_r(P;W)$ to $\cL(\cP_m(W;Y),\cP_d(P;Y))$, with the same bound.
\item $\Theta_r$ is DP on its whole domain, and linear in $b$.
\item \textup{(Jet chain rule.)} If $f$ is DP at $x_0$ with coefficient laws $a_n$, and $g$ is DP at $f(x_0)$ with coefficient laws $b_m$, then the Hasse jet of order $r$ of $g\circ f$ at $x_0$, that is, $(D^{(n)}(g\circ f)(x_0))_{0\le n\le r}$, is $\Theta_r((b_m)_{m\le r},(a_n)_{1\le n\le r})$.
\item \textup{(Associativity.)} For $b\in\mathbb J_r(W;Y)$, $a\in\mathbb J_r(P';W)$ and $\ell\in\mathbb J^-_r(P;P')$,
\[
 \Theta_r\bigl(b,\Theta_r(a,\ell)^-\bigr)=\Theta_r\bigl(\Theta_r(b,a^-),\ell\bigr).
\]
\end{enumerate}
\end{lemma}

\begin{proof}
\emph{Step 1 (item~(1)).} Put $\mathsf v_i=\sum_ls_{i,l}w_{i,l}$ and expand $b(\sum_i\lambda_i\mathsf v_i)$, $\lambda\in K^q$, in two ways. In the vectors $\mathsf v_i$ it is $\sum_{\abs{\alpha'}=m}\lambda^{\alpha'}b^{[\alpha']}(\mathsf v_1,\dots,\mathsf v_q)$. In the vectors $w_{i,l}$, with the scalars $\lambda_is_{i,l}$, it is $\sum_\delta\prod_{i,l}(\lambda_is_{i,l})^{\delta_{i,l}}b^{[\delta]}((w_{i,l})_{i,l})$. As $K$ is infinite, the coefficients of $\prod_i\lambda_i^{\alpha_i}$ agree.

\emph{Step 2 (item~(2)).} For $h_1,\dots,h_{l'}\in P$ and $s\in K^{l'}$ we have the expansion $\ell_i(\sum_\nu s_\nu h_\nu)=\sum_{\abs{\alpha'}=d_i}s^{\alpha'}\ell_i^{[\alpha']}(h_1,\dots,h_{l'})$. Apply~\eqref{dp:13:eq:refine} with the index $l$ running over these $\alpha'$, $w_{i,\alpha'}=\ell_i^{[\alpha']}(h_1,\dots,h_{l'})$ and $s_{i,\alpha'}=s^{\alpha'}$. The result is a polynomial in $s$ with coefficients in $Y$, each monomial of degree $\sum_{i,\alpha'}\abs{\alpha'}\delta_{i,\alpha'}=\sum_id_i\alpha_i=d$. The bound is~\eqref{dp:laws:scaled}.

\emph{Step 3 (item~(3)).} Item~(2) with $d_n=n$ gives the values of $\mathcal F_\alpha$ and $\norm{\mathcal F_\alpha(b,\ell)}\le C_0^m\frac{m!}{\alpha!}\norm b\prod_n\norm{\ell_n}^{\alpha_n}$. Let $(b^{(\nu)},\ell^{(\nu)})$, $1\le\nu\le l'$, be points and $s\in K^{l'}$. By linearity of $b\mapsto b^{[\alpha]}$ and~\eqref{dp:13:eq:refine},
\[
 \mathcal F_\alpha\Bigl(\sum_\nu s_\nu b^{(\nu)},\sum_\nu s_\nu\ell^{(\nu)}\Bigr)(h)=\sum_\nu\sum_\delta s_\nu\prod_{n,i}s_i^{\delta_{n,i}}\,(b^{(\nu)})^{[\delta]}\bigl(\ell^{(1)}_1(h),\dots,\ell^{(l')}_r(h)\bigr).
\]
By item~(2), each map of $h$ on the right lies in $\cP_d(P;Y)$, of degree $\sum_{n,i}n\delta_{n,i}=\sum_nn\alpha_n=d$. So we have a polynomial in $s$ with coefficients in $\cP_d(P;Y)$, homogeneous of degree $1+m$. For the operator statement, take $b^{(\nu)}=b$ fixed and vary only $\ell$: the coefficient of each monomial is the operator $b\mapsto[h\mapsto b^{[\delta]}(\ell^{(1)}_1(h),\dots)]$, bounded by item~(2).

\emph{Step 4 (items~(4), (5)).} $\Theta_r$ is a finite sum of laws, composed with bounded linear projections and inclusions. Laws are DP on their whole domain (Theorem~\ref{dp:thm:calculus}(1)), and maps into a product, sums and bounded linear images of DP maps are DP (Theorem~\ref{dp:thm:calculus}(2) and Section~\ref{dp:ssec:inverse}). Item~(5) is Proposition~\ref{dp:prop:chain-rule}.

\emph{Step 5 (item~(6)).} Let $\mathsf e(h)=\sum_{1\le n\le r}\ell_n(h)$ on $P$, $f(x)=\sum_{n\le r}a_n(x)$ on $P'$, and $g(w)=\sum_{m\le r}b_m(w-a_0)$ on $W$. These are DP everywhere, and their Hasse jets of order $r$ at $0$, $0$, $a_0$ are $\ell$, $a$, $b$. By item~(5) applied to $(g\circ f)\circ\mathsf e=g\circ(f\circ\mathsf e)$ at $0$, both sides are the Hasse jet of order $r$ of the same map.
\end{proof}

\subsection{DP jet bundles}\label{dp:ssec:jets}

\begin{definition}\label{dp:def:jets}
Let $\mathcal M,\mathcal N$ be DP manifolds (Definition~\ref{dp:def:manifold}(1)) modeled on $P$ and $W$. For DP maps $\mathsf F,\mathsf G$ defined on neighborhoods of $\mathsf m\in\mathcal M$ with values in $\mathcal N$, write $\mathsf F\sim^r_{\mathsf m}\mathsf G$ if $\mathsf F(\mathsf m)=\mathsf G(\mathsf m)$ and, for one pair of charts $\varphi$ at $\mathsf m$ and $\chi$ at $\mathsf F(\mathsf m)$,
\[
 D^{(n)}(\chi\circ\mathsf F\circ\varphi^{-1})(\varphi(\mathsf m))=D^{(n)}(\chi\circ\mathsf G\circ\varphi^{-1})(\varphi(\mathsf m))\qquad(1\le n\le r).
\]
The class $j^r_{\mathsf m}\mathsf F$ is the \emph{$r$-jet} of $\mathsf F$ at $\mathsf m$. The set of all $r$-jets is $J^r(\mathcal M,\mathcal N)$, with the \emph{source} $j^r_{\mathsf m}\mathsf F\mapsto\mathsf m$ and the \emph{target} $j^r_{\mathsf m}\mathsf F\mapsto\mathsf F(\mathsf m)$. Put $J^0(\mathcal M,\mathcal N)=\mathcal M\times\mathcal N$. For a normed space $W$, read as a DP manifold with the single chart $\id_W$, put $J^r(\mathcal M;W)=J^r(\mathcal M,W)$.
\end{definition}

The relation does not depend on the charts. For other charts $\varphi',\chi'$ put $\psi=\varphi'\circ\varphi^{-1}$ and $\xi=\chi'\circ\chi^{-1}$. Near $\varphi'(\mathsf m)$ we have $\chi'\circ\mathsf F\circ\varphi'^{-1}=\xi\circ(\chi\circ\mathsf F\circ\varphi^{-1})\circ\psi^{-1}$. By Lemma~\ref{dp:lem:chain-law}(5) applied twice, its Hasse jet of order $r$ is a function of that of $\chi\circ\mathsf F\circ\varphi^{-1}$ at $\varphi(\mathsf m)$ alone.

For a chart change $\psi$ on $\mathcal M$ and a point $\tilde x$ of its range put $\ell^\psi(\tilde x)=(D^{(n)}(\psi^{-1})(\tilde x))_{1\le n\le r}\in\mathbb J^-_r(P;P)$. Let $\mathsf T(\ell)\colon\mathbb J_r(P;W)\to\mathbb J_r(P;W)$, for $\ell\in\mathbb J^-_r(P;P)$, be the operator $\mathsf T(\ell)a=\Theta_r(a,\ell)$, and put $\mathsf T^\psi(x)=\mathsf T(\ell^\psi(\psi(x)))$. This is the \emph{operator family} of $\psi$.

\begin{theorem}[local form]\label{dp:thm:jet-local}
In the notation of Definition~\ref{dp:def:jets}, let $U_\varphi,U_\chi$ be the chart domains, and let $\Omega_{\varphi,\chi}\subseteq J^r(\mathcal M,\mathcal N)$ be the set of jets with source in $U_\varphi$ and target in $U_\chi$. The map
\[
 \Xi_{\varphi,\chi}\colon\Omega_{\varphi,\chi}\to\varphi(U_\varphi)\times\chi(U_\chi)\times\mathbb J^-_r(P;W),
\]
\[
 j^r_{\mathsf m}\mathsf F\mapsto\Bigl(\varphi(\mathsf m),\chi(\mathsf F(\mathsf m)),\bigl(D^{(n)}(\chi\circ\mathsf F\circ\varphi^{-1})(\varphi(\mathsf m))\bigr)_{1\le n\le r}\Bigr),
\]
is a bijection onto this open subset of $P\times W\times\mathbb J^-_r(P;W)$. For a second pair $(\varphi',\chi')$, the transition $\Xi_{\varphi',\chi'}\circ\Xi_{\varphi,\chi}^{-1}$ is defined on $\varphi(U_\varphi\cap U_{\varphi'})\times\chi(U_\chi\cap U_{\chi'})\times\mathbb J^-_r(P;W)$ and equals
\begin{equation}\label{dp:13:eq:transition}
 (x,y,a)\longmapsto\Bigl(\psi(x),\ \xi(y),\ \Theta_r\bigl(\mathsf b(y),\ \Theta_r\bigl((y,a),\ \ell^\psi(\psi(x))\bigr)^-\bigr)^-\Bigr),
\end{equation}
where $\mathsf b(y)=(D^{(m)}\xi(y))_{0\le m\le r}$ and $(y,a)=(y,a_1,\dots,a_r)$. It is DP, with DP inverse.
\end{theorem}

\begin{proof}
\emph{Step 1 (bijection).} Injectivity is the definition of $\sim^r_{\mathsf m}$. Let $(x,y,a)$ be in the product. The map $\hat{\mathsf G}(\tilde x)=y+\sum_{n=1}^ra_n(\tilde x-x)$ is DP everywhere and continuous, with Hasse jet $(y,a_1,\dots,a_r)$ at $x$. On an open neighborhood $O$ of $x$ in $\varphi(U_\varphi)$ with $\hat{\mathsf G}(O)\subseteq\chi(U_\chi)$, the map $\mathsf F=\chi^{-1}\circ\hat{\mathsf G}\circ\varphi$ is DP on $\varphi^{-1}(O)$, and $\Xi_{\varphi,\chi}(j^r_{\mathsf m}\mathsf F)=(x,y,a)$ for $\mathsf m=\varphi^{-1}(x)$.

\emph{Step 2 (the formula).} Let $\hat{\mathsf F}=\chi\circ\mathsf F\circ\varphi^{-1}$ near $x=\varphi(\mathsf m)$, with Hasse jet $(y,a)$ at $x$. Near $\psi(x)$ we have $\chi'\circ\mathsf F\circ\varphi'^{-1}=\xi\circ\hat{\mathsf F}\circ\psi^{-1}$. By Lemma~\ref{dp:lem:chain-law}(5), the jet of $\hat{\mathsf F}\circ\psi^{-1}$ at $\psi(x)$ is $\Theta_r((y,a),\ell^\psi(\psi(x)))$. Its value is $y$, so the jet of $\xi\circ(\hat{\mathsf F}\circ\psi^{-1})$ is the outer $\Theta_r$ in~\eqref{dp:13:eq:transition}. The target point is $\xi(y)$ and the source point is $\psi(x)$.

\emph{Step 3 (DP).} The maps $\psi,\xi,\psi^{-1}$ are DP on open sets. By Theorem~\ref{dp:thm:differentiable}(2), $\tilde x\mapsto D^{(n)}(\psi^{-1})(\tilde x)$ and $y\mapsto D^{(m)}\xi(y)$ are DP on open sets, with values in $\cP_n(P;P)$ and $\cP_m(W;W)$. So $x\mapsto\ell^\psi(\psi(x))$ and $y\mapsto\mathsf b(y)$ are DP. $\Theta_r$ is DP by Lemma~\ref{dp:lem:chain-law}(4). The map~\eqref{dp:13:eq:transition} is obtained from these by pairs, composition and projections, so it is DP on its open domain. Its inverse is the transition in the other direction, DP by the same argument. DP maps are continuous, so the transition is a homeomorphism between open sets.
\end{proof}

\begin{theorem}[DP jet bundles]\label{dp:thm:dp-jets}
Let $\mathcal M,\mathcal N$ be DP manifolds with normed models $P$ and $W$, in any characteristic, with no completeness and no ultrametric assumption.
\begin{enumerate}
\item There is a unique topology and maximal DP atlas on $J^r(\mathcal M,\mathcal N)$ for which every $\Xi_{\varphi,\chi}$ is a chart. With it, $J^r(\mathcal M,\mathcal N)$ is a DP manifold modeled on $P\times\mathbb J_r(P;W)$.
\item The forgetful maps $J^r\to J^{r-1}\to\dots\to J^0=\mathcal M\times\mathcal N$ are DP. In the charts of item~\textup{(1)} they are coordinate projections.
\item \textup{(Prolongation.)} For a DP map $\mathsf F\colon\mathcal M\to\mathcal N$, the map $j^r\mathsf F\colon\mathsf m\mapsto j^r_{\mathsf m}\mathsf F$ is a DP section of the source.
\item \textup{(Functoriality.)} For a DP map $\mathsf G\colon\mathcal N\to\mathcal N'$, the map $j^r_{\mathsf m}\mathsf F\mapsto j^r_{\mathsf m}(\mathsf G\circ\mathsf F)$ is DP. For a DP diffeomorphism $H\colon\mathcal M'\to\mathcal M$, the map $j^r_{\mathsf m}\mathsf F\mapsto j^r_{H^{-1}(\mathsf m)}(\mathsf F\circ H)$ from $J^r(\mathcal M,\mathcal N)$ to $J^r(\mathcal M',\mathcal N)$ is DP. Both are compatible with composition and identities.
\item The source $J^r(\mathcal M;W)\to\mathcal M$ is a DP vector bundle (Definition~\ref{dp:def:manifold}(2)) with fiber $\mathbb J_r(P;W)$, whose transition for a chart change $\psi$ is $(x,a)\mapsto(\psi(x),\mathsf T^\psi(x)a)$. The forgetful maps are DP bundle morphisms, and $j^r\mathsf F$ is a DP section for every DP map $\mathsf F\colon\mathcal M\to W$.
\end{enumerate}
\end{theorem}

\begin{proof}
\emph{Step 1 (item~(1)).} The sets $\Omega_{\varphi,\chi}$ cover $J^r(\mathcal M,\mathcal N)$. By Theorem~\ref{dp:thm:jet-local}, the charts are bijections onto open sets, and the images of the overlaps are open, with transitions DP homeomorphisms between them. So there is a unique topology making each $\Omega_{\varphi,\chi}$ open and each $\Xi_{\varphi,\chi}$ a homeomorphism onto its image, and the charts form a DP atlas. Two jets with different pairs (source, target) are separated by preimages of disjoint open sets under the continuous map to the Hausdorff space $\mathcal M\times\mathcal N$. Two jets with the same pair lie in one $\Omega_{\varphi,\chi}$, which is homeomorphic to a subset of a normed space.

\emph{Step 2 (items~(2), (3)).} In charts the forgetful map is bounded linear: it is $(x,y,a_1,\dots,a_r)\mapsto(x,y,a_1,\dots,a_{r-1})$. For $\hat{\mathsf F}=\chi\circ\mathsf F\circ\varphi^{-1}$, the map $j^r\mathsf F$ reads $x\mapsto(x,\hat{\mathsf F}(x),(D^{(n)}\hat{\mathsf F}(x))_{n\le r})$, which is DP by Theorem~\ref{dp:thm:differentiable}(2).

\emph{Step 3 (item~(4)).} For $\hat{\mathsf G}=\chi''\circ\mathsf G\circ\chi^{-1}$ and $\hat H=\varphi\circ H\circ\varphi_1^{-1}$ the two maps read
\begin{gather*}
 (x,y,a)\mapsto\bigl(x,\hat{\mathsf G}(y),\Theta_r((D^{(m)}\hat{\mathsf G}(y))_{m\le r},a)^-\bigr),\\
 (x,y,a)\mapsto\bigl(\hat H^{-1}(x),y,\Theta_r((y,a),(D^{(n)}\hat H(\hat H^{-1}(x)))_{1\le n\le r})^-\bigr).
\end{gather*}
Both are DP as in Step~3 of the proof of Theorem~\ref{dp:thm:jet-local}. Compatibility with composition and identities holds on representatives.

\emph{Step 4 (item~(5)).} For $\mathcal N=W$ we have $\xi=\id_W$ and $\mathsf b(y)=(y,\id_W,0,\dots,0)$. The only $\alpha$ with $\abs\alpha=1$ and $\deg\alpha=d$ is the unit vector at $d$, so the outer $\Theta_r$ in~\eqref{dp:13:eq:transition} is the identity on jets. So the transition is $(x,a)\mapsto(\psi(x),\mathsf T^\psi(x)a)$, for $a=(a_0,\dots,a_r)$, linear in $a$. By Lemma~\ref{dp:lem:chain-law}(3), $\ell\mapsto\mathsf T(\ell)$ is a finite sum of laws into $\cL(\mathbb J_r(P;W),\mathbb J_r(P;W))$, each $\mathcal F_\alpha({\cdot},\ell)$ placed in its block by bounded linear maps of norm $1$, plus the constant identity on the block of degree $0$. So it is DP, and $\mathsf T^\psi$ is DP as a composite. By Lemma~\ref{dp:lem:chain-law}(5),(6), the family $\mathsf T^\psi$ satisfies the cocycle identity, and it is invertible with inverse the family of the reversed chart change. In these charts the forgetful maps are the coordinate projections $\mathbb J_r\to\mathbb J_{r-1}$, and $j^r\mathsf F$ is DP by Step~2.
\end{proof}

\subsection{Jet bundles of analytic manifolds}\label{dp:ssec:analytic-jets}

An \emph{analytic manifold} is a Hausdorff space with an atlas modeled on a normed space whose transitions are analytic at every point of their domains. By Theorem~\ref{dp:thm:calculus}(3) it is a DP manifold. Below, charts of an analytic manifold are those of its analytic atlas. Analytic maps compose. Let $f$ be analytic at $x_0$ with multilinear coefficients $\beta_n$, $\norm{\beta_n}\le MR^{-n}$, and $g$ analytic at $f(x_0)$ with multilinear coefficients $\beta'_m$, $\norm{\beta'_m}\le M'R'^{-m}$. Substituting gives in degree $d\ge1$ the finite sum of the $\beta'_m(\beta_{n_1},\dots,\beta_{n_m})$ over $m\ge1$ and $n_1+\dots+n_m=d$ with all $n_i\ge1$. It is a bounded $d$-linear map into the target of $g$, of norm at most $M'R^{-d}(M/R')(1+M/R')^{d-1}$. For small $\norm h$ the family is absolutely summable in the completion of that target, and its sum lies in the target.

For $k\ge1$ let $\Delta_k\colon\Mult^k(P;W)\to\cP_k(P;W)$, $\Delta_k\beta(h)=\beta(h,\dots,h)$, a bounded linear map of norm at most $1$. We say that \emph{$(\mathrm{Sec}_k)$ holds for $(P,W)$} if $\Delta_k$ has a bounded linear right inverse $\mathsf s_k$, a \emph{section}. Read $\Mult^0(P;W)=W$, $\Delta_0=\id_W$ and $\mathsf s_0=\id_W$, so that $(\mathrm{Sec}_0)$ is trivial. Put $Z_k=\cP_k(P;W)$. The law
\[
 \ev_k\colon P\to\cL(Z_k,W),\qquad\ev_k(h)(a)=a(h),
\]
has degree $k$ and norm at most $1$: the coefficients of $s\mapsto\ev_k(\sum_is_ih_i)$ are the operators $a\mapsto a^{[\alpha]}(h_1,\dots,h_q)$, bounded by Theorem~\ref{dp:thm:mixed}. A section is one linear choice of lifts for all laws at once; item~(5) of Theorem~\ref{dp:thm:analytic-jets} uses this linearity.

\begin{lemma}[coefficients through a lift]\label{dp:13:lem:lifts}
\leavevmode
\begin{enumerate}
\item If $b=\Delta_m\beta$ for a bounded $m$-linear $\beta\colon W^m\to Y$, then for $\alpha\in\N^q$ with $\abs\alpha=m$,
\begin{equation}\label{dp:13:eq:lift-coeff}
 b^{[\alpha]}(w_1,\dots,w_q)=\sum_\sigma\beta\bigl(w_{\sigma(1)},\dots,w_{\sigma(m)}\bigr),
\end{equation}
the sum over the $m!/\alpha!$ maps $\sigma\colon[m]\to[q]$ with $\abs{\sigma^{-1}(i)}=\alpha_i$ for all $i$.
\item For $\alpha\in\N^r$ with $\deg\alpha=d\ge1$ and $\abs\alpha=m$, the formula
\[
 \mathsf N_\alpha\bigl(\beta,\ell^{(1)},\dots,\ell^{(m)}\bigr)(h)=\sum_\sigma\beta\bigl(\ell^{(1)}_{\sigma(1)}(h),\dots,\ell^{(m)}_{\sigma(m)}(h)\bigr),\qquad\ell^{(i)}\in\mathbb J^-_r(P;W),
\]
defines a bounded $(m+1)$-linear map from $\Mult^m(W;Y)\times\mathbb J^-_r(P;W)^m$ to $\cP_d(P;Y)$, and $\mathsf N_\alpha(\beta,\ell,\dots,\ell)=\mathcal F_\alpha(\Delta_m\beta,\ell)$.
\item Let $g\colon U\to Y$, $U\subseteq W$ open, be analytic at every point, and analytic at $w_0$ with radius $R$, bound $M$ and multilinear coefficients $\beta_m$. For $n\le m$ and $S\subseteq[m]$ with $\abs S=n$, let $\beta_m^S(h_1,\dots,h_{m-n};h'_1,\dots,h'_n)$ be $\beta_m$ with $h'_1,\dots,h'_n$ in the slots of $S$ and $h_1,\dots,h_{m-n}$ in the others, each in order, and put $\beta_{m,n}=\sum_S\beta_m^S$. Then $\norm{\beta_{m,n}}\le M(2/R)^m$, and for $\norm h<R/2$,
\begin{equation}\label{dp:13:eq:hasse-analytic}
 D^{(n)}g(w_0+h)=\Delta_n\mathsf H_n(w_0+h),\qquad\mathsf H_n(w_0+h)=\sum_{m\ge n}\beta_{m,n}(h,\dots,h)\ \text{ in }\Mult^n(W;\widehat Y).
\end{equation}
So $w\mapsto D^{(n)}g(w)$ is analytic on $U$ into $\cP_n(W;Y)$, and $\mathsf H_n$ is analytic into $\Mult^n(W;\widehat Y)$ with coefficients in $\Mult^n(W;Y)$.
\end{enumerate}
\end{lemma}

\begin{proof}
\emph{Step 1 (items~(1), (2)).} Expand $\beta(\sum_is_iw_i,\dots,\sum_is_iw_i)$ by multilinearity and compare the coefficients of $s^\alpha$, as $K$ is infinite. For item~(2), let $\ell_1,\dots,\ell_m$ be laws of degrees $d_1,\dots,d_m$ into $W$. Substituting their expansions shows that $h\mapsto\beta(\ell_1(h),\dots,\ell_m(h))$ is a law of degree $\sum_id_i$ and norm at most $\norm\beta\prod_i\norm{\ell_i}$. For each $\sigma$ the degree is $\sum_i\sigma(i)=\deg\alpha=d$. So $\mathsf N_\alpha$ has norm at most $m!/\alpha!$. The last formula is~\eqref{dp:13:eq:lift-coeff}.

\emph{Step 2 (item~(3)).} There are $\binom mn\le2^m$ sets $S$, which gives the bound. For $\norm h+\norm{h'}<R/2$, expanding $\beta_m(h+h',\dots,h+h')$ gives a family with $\sum_{m,n}\binom mnMR^{-m}\norm h^{m-n}\norm{h'}^n=\sum_mMR^{-m}(\norm h+\norm{h'})^m<\infty$, absolutely summable in $\widehat Y$. Grouping by the degree in $h'$ gives $g(w_0+h+h')=\sum_n\mathsf H_n(w_0+h)(h',\dots,h')$. As $g$ is DP at $w_0+h$, uniqueness of law expansions in $\widehat Y$ (Lemma~\ref{dp:lem:elementary}(5)) gives~\eqref{dp:13:eq:hasse-analytic}. Applying $\Delta_n$ to the series gives an expansion of $D^{(n)}g$ with coefficients in $\cP_n(W;Y)$.
\end{proof}

\begin{definition}[the shear manifold]\label{dp:13:def:shear}
For a normed space $V$, let $X_V=V\times V\times K$, with the two global charts $\id$ and $\psi_V(v,z,c)=(v,z+cv,c)$. The chart change $\psi_V$ and its inverse $(v,z,c)\mapsto(v,z-cv,c)$ are polynomial and analytic. This is an analytic manifold, the \emph{shear manifold} of $V$. Below, the operator family and the transition of the shear manifold mean those of $\psi=\psi_V$ or of $\psi=\psi_V^{-1}$, with $\xi=\id_W$ for the transition; each statement holds for either one.
\end{definition}

\begin{theorem}[jet bundles of analytic manifolds]\label{dp:thm:analytic-jets}
Let $k\ge1$.
\begin{enumerate}
\item \textup{(Reformulations.)} The following are equivalent: \textup{(a)} $(\mathrm{Sec}_k)$ holds for $(P,W)$; \textup{(b)} $\ev_k$ is the diagonal of a bounded $k$-linear map into $\cL(Z_k,W)$; \textup{(c)} $\ev_k$ is analytic at one point. Then every law in $Z_k$ has a lift of norm at most $\norm{\mathsf s_k}$ times its norm. $(\mathrm{Sec}_k)$ passes to bounded linear retracts of $P$.
\item \textup{(Sufficiency.)} If $(\mathrm{Sec}_k)$ holds for $(P,W)$ for $1\le k\le r$, then for all analytic manifolds $\mathcal M,\mathcal N$ with models $P,W$, the transitions~\eqref{dp:13:eq:transition} of $J^r(\mathcal M,\mathcal N)$ and the operator families $\mathsf T^\psi$ of $J^r(\mathcal M;W)$ are analytic.
\item \textup{(Necessity.)} If, for the shear manifold of $V$ and $\psi=\psi_V^{\pm1}$, the operator family $\mathsf T^\psi$ of $J^r(X_V;W)$ or the transition~\eqref{dp:13:eq:transition} with $\mathcal N=W$ is analytic at one point, then $(\mathrm{Sec}_k)$ holds for $(V,W)$ for $1\le k\le r$.
\item \textup{(Where sections exist.)} $(\mathrm{Sec}_k)$ holds for $(P,W)$ in each of these cases:
\begin{enumerate}
\item $k!\ne0$ in $K$, all normed $P,W$, with $\mathsf s_k(a)=a^{[1^k]}/k!$;
\item $P$ finite-dimensional, with $K$ complete or $W$ Banach;
\item $P=\ell^1(J)$, $W$ Banach;
\item $P$ ultrametric, $W$ with an equivalent ultrametric spherically complete norm; in particular $W=K$ with $K$ spherically complete;
\item $K$ discretely valued, $P$ ultrametric, $W$ Banach with an equivalent ultrametric norm;
\item $K$ complete, $P$ ultrametric, $W$ ultrametric Banach, $\widehat P$ linearly homeomorphic to some $c_0(J,K)$, which includes $\widehat P$ of countable type.
\end{enumerate}
\item \textup{(What a section gives.)} If $(\mathrm{Sec}_k)$ holds for $(\cL(E,D),F)$, then $A^k_{E,D;F}$ is the diagonal of a bounded $k$-linear map into $\cL(\Alt^k(D;F),\Alt^k(E;F))$.
\item \textup{(Where sections fail.)} Let $\operatorname{char}K=p>0$ and $k\ge p$.
\begin{enumerate}
\item Over every such $K$, with $E,F$ of \cite[Theorem~7.2(1)]{I}, $(\mathrm{Sec}_k)$ fails for $(\cL(E,E),F)$.
\item If $K$ is complete and not spherically complete, with $E,D$ of \cite[Theorem~7.2(2)]{I}, $(\mathrm{Sec}_k)$ fails for $(\cL(E,D),K)$, with scalar values and all spaces ultrametric Banach.
\item It fails for every pair $(P,W)$ with a law of degree $k$ that has no lift. Examples: $A^k_{E,E;F}$ as in~\textup{(a)}, on $\cL(E,E)$ with $W=\cL(\Alt^k(E;F),\Alt^k(E;F))$; $A^k_{E,D;K}$ as in~\textup{(b)}, for $K$ complete and not spherically complete; the power map $c_0\to\hGamma^k_\pi c_0$ under the conditions of Proposition~\ref{dp:prop:strict}\textup{(1)(b)}; and, with $k=p$, the law $\gamma_{\mathcal E}(h)=h^{\otimes p}$ on the space $\mathcal E$ of Proposition~\ref{dp:prop:degree-p}\textup{(2)}, with $W=\hGamma^p_\Gamma(\mathcal E)$, the completion of $\Gamma^p(\mathcal E)$ for $\norm\cdot_\Gamma$.
\end{enumerate}
So in each case, for $V$ the source space and $r\ge k$, the operator families of $J^r(X_V;W)$ are analytic at no point; they are DP by Theorem~\ref{dp:thm:dp-jets}.
\end{enumerate}
\end{theorem}

Items (1)--(3), (4)(a)--(c) and (5) assume no ultrametric norm, over every $K$. Ultrametric norms enter only in (4)(d)--(f), through the agreement theorems of Section~\ref{dp:sec:agreement}, and they place (6)(b) between ultrametric spaces.

\begin{proof}
\emph{Step 1 (item~(1)).} The formula $\mathsf E(x_1,\dots,x_k)(a)=\mathsf s_k(a)(x_1,\dots,x_k)$ matches sections with lifts $\mathsf E$ of $\ev_k$. The diagonal of $\mathsf s_k(a)$ is $a$ for all $a$ exactly when the diagonal of $\mathsf E$ is $\ev_k$, and the bounds agree. So (a) and (b) are equivalent. A diagonal of a bounded $k$-linear map is analytic at every point~\cite[Proposition~2.5]{I}. Conversely, let $\ev_k(h_0+h)=\sum_i\beta_i(h,\dots,h)$ near $0$. For fixed $h$ and small $\lambda\in K$, the left side at $\lambda h$ is a polynomial in $\lambda$ with coefficient of $\lambda^k$ equal to $\ev_k(h)$. By~\cite[Lemma~A.1]{I}, $\beta_k(h,\dots,h)=\ev_k(h)$. For retracts, let $\mathsf t\colon P'\to P$ and $\mathsf r\colon P\to P'$ be bounded linear with $\mathsf r\circ\mathsf t=\id$. Then $a\mapsto\mathsf s_k(a\circ\mathsf r)\circ\mathsf t^{\times k}$ is a section for $(P',W)$.

\emph{Step 2 (item~(2), the inner composition).} As $\psi^{-1}$ is analytic, $x\mapsto\ell^\psi(\psi(x))$ is analytic, by Lemma~\ref{dp:13:lem:lifts}(3) and composition. By Lemma~\ref{dp:13:lem:lifts}(2) for the spaces $P,W$ in place of $W,Y$, with $\beta=\mathsf s_m(a_m)$, the map $(a_m,\ell^{(1)},\dots,\ell^{(m)})\mapsto\mathsf N_\alpha(\mathsf s_m(a_m),\ell^{(1)},\dots,\ell^{(m)})$ is bounded $(m+1)$-linear, and its value at $\ell^{(i)}=\ell$ is $\mathcal F_\alpha(a_m,\ell)$. So $(x,y,a)\mapsto\Theta_r((y,a),\ell^\psi(\psi(x)))$ is analytic. Leaving $a_m$ as the input, the bounded $m$-linear map $(\ell^{(1)},\dots,\ell^{(m)})\mapsto\mathsf N_\alpha(\mathsf s_m({\cdot}),\ell^{(1)},\dots,\ell^{(m)})$ has diagonal $\ell\mapsto\mathcal F_\alpha({\cdot},\ell)$. So $\mathsf T^\psi$, with the constant identity on the block of degree~$0$, is analytic.

\emph{Step 3 (item~(2), the outer composition).} By Lemma~\ref{dp:13:lem:lifts}(3), $D^{(m)}\xi(y)=\Delta_m\mathsf H_m(y)$ with $\mathsf H_m$ analytic. By Lemma~\ref{dp:13:lem:lifts}(2), for the spaces $W,\widehat W$ in place of $W,Y$, we have $\mathcal F_\alpha(D^{(m)}\xi(y),\ell)=\mathsf N_\alpha(\mathsf H_m(y),\ell,\dots,\ell)$. No section is needed here. So the outer $\Theta_r$ in~\eqref{dp:13:eq:transition} is analytic in $(y,\ell)$, at first with values in $\cP_d(P;\widehat W)$. Each coefficient of the resulting series is a finite sum of bounded multilinear maps with values in $\cP_d(P;W)$, so the expansion lives in $\cP_d(P;W)$. The components $\psi$ and $\xi$ are analytic, and analytic maps compose.

\emph{Step 4 (item~(3), the operator family).} Let $\psi=\psi_V$. At $x=(v,z,c)$ the map $\psi^{-1}$ has $(\ell^\psi(\psi(x)))_1(v',z',c')=(v',z'-c'v-cv',c')$, since $\psi$ changes neither $v$ nor $c$. Let $\operatorname{pr}_2(v,z,c)=z$, and let $\mathsf i_k,\mathsf p_k$ be the inclusion of and the projection onto the block of degree $k$ of $\mathbb J_r(X_V;W)$. Define the bounded linear map $\mathsf C$ from $\cL(\mathbb J_r(X_V;W),\mathbb J_r(X_V;W))$ to $\cL(\cP_k(V;W),W)$ by
\[
 \mathsf C(\mathsf S)(a)=\bigl(\mathsf p_k\mathsf S\mathsf i_k(a\circ\operatorname{pr}_2)\bigr)(0,0,1).
\]
For a jet with only one component $a'\in\cP_k(X_V;W)$, the component of degree $k$ of $\mathsf T(\ell)a'$ comes only from $\alpha$ with $\abs\alpha=\deg\alpha=k$, that is, $\alpha=(k,0,\dots)$, and it is $a'\circ\ell_1$. So
\begin{equation}\label{dp:13:eq:compression}
 \mathsf C\bigl(\mathsf T^\psi(v,z,c)\bigr)(a)=a\bigl(\operatorname{pr}_2(0,-v,1)\bigr)=a(-v)=(-1)^k\ev_k(v)(a).
\end{equation}
If $\mathsf T^\psi$ is analytic at $(v_0,z_0,c_0)$, restriction to the slice $v\mapsto(v,z_0,c_0)$ and $\mathsf C$ make $\ev_k$ analytic at $v_0$, and item~(1) applies. For $\psi_V^{-1}$ the signs are $+$, and the same argument gives $\ev_k(v)$.

\emph{Step 5 (item~(3), the transition).} Let $\mathcal N=W$. By Theorem~\ref{dp:thm:dp-jets}(5), the block of degree $k$ of the transition is $\Phi(x,a)=\mathsf p_k\mathsf T^\psi(x)a$, linear in $a$. Let $\Phi(x_0+h,a^\circ+a'')=\sum_m\beta_m((h,a''),\dots,(h,a''))$ with $\norm{\beta_m}\le MR^{-m}$. By~\cite[Lemma~A.1]{I}, the coefficient of $s$ in $s\mapsto\Phi(x_0+h,a^\circ+sa'')$ gives
\[
 \mathsf p_k\mathsf T^\psi(x_0+h)a''=\sum_m\sum_{l\le m}\beta_m\bigl((h,0),\dots,(0,a''),\dots,(h,0)\bigr),
\]
with $(0,a'')$ in slot $l$. The right side is an operator-valued series in $h$ with coefficients of norm at most $mMR^{-m}$. So $x\mapsto\mathsf p_k\mathsf T^\psi(x)$ is analytic at $x_0$, and $\mathsf C$ uses only $\mathsf p_k\mathsf S$.

\emph{Step 6 (item~(4)).} (a): $a^{[1^k]}/k!$ is $k$-linear by Lemma~\ref{dp:lem:elementary}(4), bounded linearly in $a$ by~\eqref{dp:laws:scaled}, and has diagonal $a$ by Lemma~\ref{dp:lem:hasse-law}(5). (c): fix a total order on $J$, let $e_\theta$ be the unit vectors, and put
\[
 \mathsf s_k(a)(x^1,\dots,x^k)=\sum_{\mathsf w}a^{[\alpha(\mathsf w)]}(e_{\theta_1},\dots,e_{\theta_q})\,x^1_{\mathsf w_1}\cdots x^k_{\mathsf w_k},
\]
the sum over the nondecreasing words $\mathsf w$ of length $k$ in $J$, with distinct letters $\theta_1<\dots<\theta_q$ and multiplicities $\alpha(\mathsf w)$. By~\eqref{dp:laws:scaled} the coefficients have norm at most $C_0^kk!\norm a$, and $\sum_{\mathsf w}\prod_l\abs{x^l_{\mathsf w_l}}\le\prod_l\norm{x^l}$. So the sum converges in the Banach space $W$. Its diagonal is $a$ on finitely supported vectors, since words correspond to monomials, and both sides are continuous.

In the other cases $\ev_k$ is DP at $0$ by Theorem~\ref{dp:thm:calculus}(1), and we show that it is analytic there. Item~(1) then finishes. The target $\cL(Z_k,W)$ is Banach if $W$ is, and ultrametric if $W$ is. (b): Theorem~\ref{dp:thm:agreement-finite}(1) if $K$ is complete, and (2) if $W$ is Banach; for $P=\{0\}$ we have $Z_k=0$. (d): replace the norm of $W$ by the equivalent one, which does not affect $(\mathrm{Sec}_k)$. Then $Z_k$ is ultrametric, $\cL(Z_k,W)$ is ultrametric and spherically complete by Lemma~\ref{dp:lem:operator-sph}, and Theorem~\ref{dp:thm:sph-agreement}(1) applies. (e): replace the norm of $W$ as in (d). Then $\cL(Z_k,W)$ is an ultrametric Banach space, and Theorem~\ref{dp:thm:sph-agreement}(2) applies. (f): Theorem~\ref{dp:thm:conservativity}(1)(d).

\emph{Step 7 (item~(5)).} For $\omega\in\Alt^k(D;F)$ and $\mathsf v_1,\dots,\mathsf v_k\in E$ let $a_{\omega,\mathsf v}\in\cP_k(\cL(E,D);F)$ be $a_{\omega,\mathsf v}(u)=\omega(u\mathsf v_1,\dots,u\mathsf v_k)$, the diagonal of a bounded $k$-linear map, of norm at most $\norm\omega\prod_i\norm{\mathsf v_i}$. The map $(\omega,\mathsf v)\mapsto a_{\omega,\mathsf v}$ is linear in $\omega$, multilinear in $\mathsf v$, and zero if two $\mathsf v_i$ coincide. So
\[
 \mathsf A\colon\cL(\cP_k(\cL(E,D);F),F)\to\cL(\Alt^k(D;F),\Alt^k(E;F)),\qquad\mathsf A(\Upsilon)(\omega)(\mathsf v_1,\dots,\mathsf v_k)=\Upsilon(a_{\omega,\mathsf v}),
\]
is a well-defined bounded linear map of norm at most $1$: $\mathsf A(\Upsilon)(\omega)$ is alternating because $\Upsilon$ is linear. And $\mathsf A(\ev_k(u))(\omega)(\mathsf v_1,\dots,\mathsf v_k)=a_{\omega,\mathsf v}(u)=A^k_{E,D;F}(u)(\omega)(\mathsf v_1,\dots,\mathsf v_k)$. So $A^k_{E,D;F}=\mathsf A\circ\ev_k$. By item~(1), $\ev_k$ is the diagonal of a bounded $k$-linear map, and composing with $\mathsf A$ gives the claim. Each $a_{\omega,\mathsf v}$ has the lift $\omega(u_1\mathsf v_1,\dots,u_k\mathsf v_k)$, which is not alternating in $\mathsf v$. So the linearity of the section is what is used.

\emph{Step 8 (item~(6)).} (a), (b): by item~(5), a section would make $A^k$ a diagonal of a bounded $k$-linear map, hence analytic, against~\cite[Theorem~7.2]{I}. (c): the implication is item~(1). $A^k$ and the power map have no lift, since a lift would make them analytic, against~\cite[Theorem~7.2]{I} and Proposition~\ref{dp:prop:strict}(1)(b). A lift of $\gamma_{\mathcal E}$ would make it ambient of degree $p$. If $\gamma_{\mathcal E}$ were ambient of degree $p$, then so would be $(\lambda,h)\mapsto h^{\otimes p}$ on $K\oplus\mathcal E$, by precomposing a lift with the projection, against Proposition~\ref{dp:prop:degree-p}(2) with $n=p$. The last sentence is item~(3) and Theorem~\ref{dp:thm:dp-jets}.
\end{proof}

\begin{corollary}[the dichotomy]\label{dp:cor:jet-dichotomy}
Let $K$ be complete of characteristic $p>0$. For $r\ge p$ the following are equivalent: \textup{(i)} $K$ is spherically complete; \textup{(ii)} for every ultrametric normed space $P$ and every analytic manifold $\mathcal M$ modeled on $P$, the scalar jet bundle $J^r(\mathcal M;K)$ has analytic operator families and transitions. For $r<p$, \textup{(ii)} always holds. In all cases $J^r(\mathcal M;K)$ is a DP bundle.
\end{corollary}

\begin{proof}
(i)$\Rightarrow$(ii): Theorem~\ref{dp:thm:analytic-jets}(4)(d),(2). (ii)$\Rightarrow$(i): if $K$ is not spherically complete, let $V=\cL(E,D)$ be as in Theorem~\ref{dp:thm:analytic-jets}(6)(b) with $k=p$. Then $X_V$ is ultrametric, and by items (6)(b) and (3) its scalar jet bundle has operator families and transitions analytic at no point. For $r<p$: items (4)(a) and (2). The DP bundle is Theorem~\ref{dp:thm:dp-jets}(5).
\end{proof}

The statement of the corollary does not mention alternating forms. Its proof does, through Theorem~\ref{dp:thm:analytic-jets}(5). The proof produces no scalar law without a lift: the map that fails to be analytic has values in $\cL(\mathbb J_r(X_V;K),\mathbb J_r(X_V;K))$. The source $X_V$ and this target are ultrametric and the family is DP, so by~\cite[Theorem~14.39(1)]{I} the family is AA. So the corollary separates An from $\mathrm{AA}=\mathrm{DP}$, and not AA from DP. The separation from AA is Proposition~\ref{dp:prop:jets-not-ambient} below, outside ultrametric spaces.

Theorem~\ref{dp:thm:analytic-jets}(2) needs sections on the model, and item~(3) gives them on $V$ only. The next proposition closes this gap.

\begin{proposition}[sections and products]\label{dp:prop:sec-products}
Let $P_1,P_2,V$ be normed spaces and $k\ge1$.
\begin{enumerate}
\item If $(\mathrm{Sec}_i)$ holds for $(P_1,W)$ and $(\mathrm{Sec}_n)$ for $(P_2,W)$ whenever $i+n=k$, then $(\mathrm{Sec}_k)$ holds for $(P_1\times P_2,W)$.
\item $(\mathrm{Sec}_k)$ holds for $(K,W)$.
\item If $(\mathrm{Sec}_k)$ holds for $(P\times K,W)$, then $(\mathrm{Sec}_n)$ holds for $(P,W)$ for all $n\le k$.
\item The following are equivalent: \textup{(A)} $(\mathrm{Sec}_n)$ holds for $(V,W)$, $1\le n\le r$; \textup{(B)} $(\mathrm{Sec}_r)$ holds for $(X_V,W)$; \textup{(B$'$)} $(\mathrm{Sec}_n)$ holds for $(X_V,W)$, $1\le n\le r$; \textup{(C)} for all analytic manifolds $\mathcal M,\mathcal N$ with models $X_V,W$, the transitions of $J^r(\mathcal M,\mathcal N)$ and the operator families of $J^r(\mathcal M;W)$ are analytic; \textup{(D)} on the shear manifold of $V$, the operator family or the transition of $J^r(X_V;W)$ is analytic at one point.
\end{enumerate}
\end{proposition}

\begin{proof}
\emph{Step 1 (item~(1)).} For $a\in\cP_k(P_1\times P_2;W)$ write $a(x_1,x_2)=\sum_{i+n=k}a_{i,n}(x_1,x_2)$, with $a_{i,n}(x_1,x_2)=a^{[(i,n)]}((x_1,0),(0,x_2))$. By Theorem~\ref{dp:thm:mixed}(2), $\norm{a_{i,n}(x_1,x_2)}\le C_2^k\norm a\norm{x_1}^i\norm{x_2}^n$. By Lemma~\ref{dp:lem:chain-law}(2), $a_{i,n}({\cdot},x_2)\in\cP_i(P_1;W)$. By~\eqref{dp:13:eq:refine} and Lemma~\ref{dp:lem:chain-law}(2), $x_2\mapsto a_{i,n}({\cdot},x_2)$ is a law of degree $n$ into $\cP_i(P_1;W)$. Let $\mathsf s^{(1)}_i,\mathsf s^{(2)}_n$ be the sections. For $x^1,\dots,x^i\in P_1$ put $\mathsf c(x_2)=\mathsf s^{(1)}_i(a_{i,n}({\cdot},x_2))(x^1,\dots,x^i)$. Composing with the bounded linear map $\mathsf s^{(1)}_i({\cdot})(x^1,\dots,x^i)$ shows that $\mathsf c\in\cP_n(P_2;W)$, and $\mathsf c$ is linear in $a$ and multilinear in $(x^1,\dots,x^i)$. Then
\[
 \mathsf L_{i,n}(x^1,\dots,x^i;\mathsf y^1,\dots,\mathsf y^n)=\mathsf s^{(2)}_n(\mathsf c)(\mathsf y^1,\dots,\mathsf y^n)
\]
is bounded and multilinear in all arguments, because $\mathsf s^{(2)}_n$ is linear. At $x^l=x_1$, $\mathsf y^l=x_2$ its value is $\mathsf c(x_2)=a_{i,n}(x_1,x_2)$. With the projections $\operatorname{pr}_1,\operatorname{pr}_2$ of $P_1\times P_2$, the map
\[
 \mathsf s_k(a)(\mathsf z^1,\dots,\mathsf z^k)=\sum_{i+n=k}\mathsf L_{i,n}\bigl(\operatorname{pr}_1\mathsf z^1,\dots,\operatorname{pr}_1\mathsf z^i;\operatorname{pr}_2\mathsf z^{i+1},\dots,\operatorname{pr}_2\mathsf z^k\bigr)
\]
is a section.

\emph{Step 2 (items~(2), (3)).} (2): $\mathsf s_k(s\mapsto s^kw)(s_1,\dots,s_k)=s_1\cdots s_kw$. (3): for $b\in\cP_n(P;W)$ put $a(v,\lambda)=\lambda^{k-n}b(v)$, a law of degree $k$ with $\norm a\le\norm b$, linear in $b$, and let $\beta=\mathsf s_k(a)$. The coefficient of $\lambda^{k-n}$ in $\beta((v,\lambda),\dots,(v,\lambda))=\lambda^{k-n}b(v)$ is $\sum_S\beta^S(v,\dots,v)$, the sum over $S\subseteq[k]$ with $\abs S=k-n$, where $\beta^S$ has $(0,1)$ in the slots of $S$ and its arguments, in the form $(v,0)$, in the others. So $(v_1,\dots,v_n)\mapsto\sum_S\beta^S(v_1,\dots,v_n)$ is $n$-linear, bounded by $\binom kn\norm{\mathsf s_k}\norm b$, linear in $b$, with diagonal $b$.

\emph{Step 3 (item~(4)).} (A)$\Rightarrow$(B$'$): items (1), (2), first for $V\times K$ and then for $V\times(V\times K)$. (B$'$)$\Rightarrow$(B) is trivial. (B)$\Rightarrow$(A): $V\times K$ is a retract of $X_V$, so Theorem~\ref{dp:thm:analytic-jets}(1) and item~(3) apply. (B$'$)$\Rightarrow$(C) is Theorem~\ref{dp:thm:analytic-jets}(2). (C)$\Rightarrow$(D) is trivial. (D)$\Rightarrow$(A) is Theorem~\ref{dp:thm:analytic-jets}(3).
\end{proof}

So if $P$ is linearly homeomorphic to $P\times P\times K$, then the transitions of all $J^r(\mathcal M,\mathcal N)$ and the operator families of all $J^r(\mathcal M;W)$, with models $P,W$, are analytic exactly when $(\mathrm{Sec}_r)$ holds for $(P,W)$.

The law $\ev_k$ also decides ambience. Recall $\operatorname{amb}$ and $\mathfrak a_k(P,W)$ from Section~\ref{dp:sec:conv}, and $\cC_p$ from Definition~\ref{dp:def:cp}. For the next proposition let $p=\operatorname{char}K>0$ in item~(4), with $\mathcal E$ as in Proposition~\ref{dp:prop:degree-p}(2).

\begin{proposition}[the ambient form of the criterion]\label{dp:prop:jets-not-ambient}
\leavevmode
\begin{enumerate}
\item $\operatorname{amb}(\ev_k)=\mathfrak a_k(P,W)$. So $\ev_k$ is ambient if and only if $\mathfrak a_k(P,W)<\infty$, and $(\mathrm{Sec}_k)$ implies $\mathfrak a_k(P,W)\le\norm{\mathsf s_k}$.
\item \textup{(Sufficiency.)} If $\mathfrak a_m(P,W)<\infty$ for $1\le m\le r$, then for every analytic manifold $\mathcal M$ with model $P$, the operator families $\mathsf T^\psi$ of $J^r(\mathcal M;W)$ are AA at every point.
\item \textup{(Necessity.)} If on the shear manifold of $V$ the operator family of $J^r(X_V;W)$ is AA at one point, then $\mathfrak a_k(V,W)<\infty$ for $1\le k\le r$.
\item \textup{(Separation.)} Let $r\ge p$ and $W=\hGamma^p_\Gamma(\mathcal E)$. On the shear manifold of $\mathcal E$, the operator family of $J^r(X_{\mathcal E};W)$ is DP at every point, and its germ at no point is in $\cC_p$. So at no point is it AA, and at no point is it TA.
\end{enumerate}
\end{proposition}

So the single law $\ev_k$ sorts the three classes: it is analytic at a point if and only if $(\mathrm{Sec}_k)$ holds, it is AA at a point if and only if $\mathfrak a_k(P,W)<\infty$, and it is always DP. Whether $\mathfrak a_n(V,W)<\infty$ for $n\le r$ implies the same for $X_V$ is open, so we have no ambient form of Proposition~\ref{dp:prop:sec-products}(4). Item~(2) does not cover the transitions of $J^r(\mathcal M;W)$: they arise from the operator family by the bilinear evaluation, and an analytic map after an AA map need not be AA in characteristic $p$~\cite[Theorem~14.17]{I}.

\begin{proof}
\emph{Step 1 (item~(1)).} For $t\in\Gamma^k(P)$ and $a\in Z_k$ we have $L_{\ev_k}(t)(a)=L_a(t)$, since both sides are linear in $t$ and agree on powers. So
\[
 \operatorname{amb}(\ev_k)=\sup_t\sup_{a\ne0}\frac{\norm{L_a(t)}}{\norm a\,\pi(t)}=\sup_{a\ne0}\frac{\operatorname{amb}(a)}{\norm a}=\mathfrak a_k(P,W).
\]
By Theorem~\ref{dp:thm:analytic-jets}(1), a section gives a lift of $\ev_k$ of norm at most $\norm{\mathsf s_k}$. Its linearization on $P^{\otimes k}$ restricts to $L_{\ev_k}$, and it is bounded by $\norm{\mathsf s_k}$ for $\pi$.

\emph{Step 2 (item~(2)).} Let $1\le m\le r$ and $Z_m=\cP_m(P;W)$. By item~(1) there are a linear isometry $\iota\colon\cL(Z_m,W)\to Z$ and a bounded $m$-linear $\beta\colon P^m\to Z$ with $\iota\circ\ev_m=\Delta_m\beta$. Let $\alpha\in\N^r$ with $\abs\alpha=m$ and $\deg\alpha=d\ge1$. The map sending $\mathsf S'\in\cL(Z_m,\cP_d(P;W))$ to the law $h\mapsto\iota(a\mapsto\mathsf S'(a)(h))$ is a linear isometry $\mathsf J$ into $\cP_d(P;Z)$; the values are laws because the coefficients $a\mapsto\mathsf S'(a)^{[\alpha']}(h_1,\dots,h_{l'})$ are bounded in $a$, by Theorem~\ref{dp:thm:mixed}. We have $\mathcal F_\alpha(a,\ell)(h)=a^{[\alpha]}(\ell(h))=\ev_m^{[\alpha]}(\ell(h))(a)$, where $\ell(h)$ stands for $\ell_1(h),\dots,\ell_r(h)$. So~\eqref{dp:13:eq:lift-coeff} for $\iota\circ\ev_m$ gives $\mathsf J(\mathcal F_\alpha({\cdot},\ell))=\mathsf N_\alpha(\beta,\ell,\dots,\ell)$, for the spaces $P,Z$ in place of $W,Y$. By Lemma~\ref{dp:13:lem:lifts}(2) this is the diagonal of a bounded $m$-linear map of $\ell$. So $\ell\mapsto\mathcal F_\alpha({\cdot},\ell)$ is ambient of degree $m$. An ambient law is AA at every point: expanding its lift by multilinearity gives a finite expansion, whose diagonals are images of the values of its coefficient laws. So $\ell\mapsto\mathsf T(\ell)$ is a finite sum of AA maps, after bounded linear maps into blocks, plus a constant. By Lemma~\ref{dp:13:lem:lifts}(3) and composition, $x\mapsto\ell^\psi(\psi(x))$ is analytic. Now apply~\cite[Proposition~14.12(1)]{I}: linear postcomposition, analytic precomposition and sums.

\emph{Step 3 (item~(3)).} Restrict to the slice $v\mapsto(v,z_0,c_0)$ and apply $\mathsf C$ of Step~4 of the proof of Theorem~\ref{dp:thm:analytic-jets}, by~\cite[Proposition~14.12(1)]{I}. By~\eqref{dp:13:eq:compression}, and its form without the sign for $\psi_V^{-1}$, the map $v\mapsto\pm\ev_k(v)$ is AA at $v_0$. Its coefficient law of degree $k$ at $v_0$ is $\pm\ev_k$, the top part of the polynomial $s\mapsto\ev_k(v_0+sh)$. By~\cite[Proposition~14.9(1)]{I} this law is ambient, and item~(1) applies.

\emph{Step 4 (item~(4)).} The family is DP by Theorem~\ref{dp:thm:dp-jets}(5). By Step~8 of the proof of Theorem~\ref{dp:thm:analytic-jets}, $\gamma_{\mathcal E}\in\cP_p(\mathcal E;W)$ is not ambient, so $\mathfrak a_p(\mathcal E,W)=\infty$ and $\ev_p$ is not ambient, by item~(1). Let $\psi=\psi_{\mathcal E}^{\pm1}$, and suppose that the germ of $\mathsf T^\psi$ at $(v_0,z_0,c_0)$ is in $\cC_p$. The slice is an analytic map, so it is in $\cC_p$ by Theorem~\ref{dp:thm:cp}(1), and $\mathsf C$ is bounded linear. By Theorem~\ref{dp:thm:cp}(2) and~\eqref{dp:13:eq:compression}, with $k=p\le r$, the germ of $\pm\ev_p$ at $v_0$ is in $\cC_p$. Its coefficient law of degree $p$ is $\pm\ev_p$, so $\ev_p$ is ambient of degree $p$. This is a contradiction. AA germs are in $\cC_p$ by Theorem~\ref{dp:thm:cp}(1). The space $W$ is Banach, so $\cP_n(X_{\mathcal E};W)$ is complete by Theorem~\ref{dp:thm:normed-roby}(3), and $\mathsf T^\psi$ has values in the Banach space $\cL(\mathbb J_r(X_{\mathcal E};W),\mathbb J_r(X_{\mathcal E};W))$. So Theorem~\ref{dp:thm:chain} applies: $\mathrm{TA}\subseteq\cC_p$ among germs with Banach values.
\end{proof}
%% ---- end sec-13-jets.tex
%% ---- begin sec-14-strong.tex
\section{The strong norm}\label{dp:sec:strong}

The strong norm of Definition~\ref{dp:def:strong-norm} bounds all the coefficients of a law. This section studies it twice. As a norm on laws, it loses no radius on ultrametric spaces (Theorems~\ref{dp:thm:strong-inverse} and~\ref{dp:thm:strong-recenter}), and there it is the law norm over a densely valued extension (Proposition~\ref{dp:prop:strong-intrinsic}(5)). As a class of germs, the strongly DP germs compete with the DP germs, and on $\ell^1(\N,K)$ they fail: with a non-ultrametric target they do not contain the analytic germs (Example~\ref{dp:ex:an-not-sdp}), and they are not closed under composition or inverses (Theorem~\ref{dp:thm:strong-failure}). If the source has an equivalent ultrametric norm, the strongly DP germs are the DP germs.

The following expansion is used throughout. Let $a$ be a law of degree $n$, $\abs\alpha=n$, and let $x_c=\sum_l\lambda_{c,l}y_{c,l}$ be finite sums, for $1\le c\le r$. Expanding $a(\sum_cs_c\sum_l\lambda_{c,l}y_{c,l})$ in the vectors $y_{c,l}$ and comparing the coefficients of $s^\alpha$ gives
\begin{equation}\label{dp:s:eq:slots}
 a^{[\alpha]}(x_1,\dots,x_r)=\sum_\varepsilon\lambda^\varepsilon a^{[\varepsilon]}\bigl((y_{c,l})_{c,l}\bigr),
\end{equation}
the sum over the arrays $\varepsilon=(\varepsilon_{c,l})$ of integers $\ge0$ with $\sum_l\varepsilon_{c,l}=\alpha_c$ for each $c$, where $\lambda^\varepsilon=\prod_{c,l}\lambda_{c,l}^{\varepsilon_{c,l}}$.

\subsection{The hull}\label{dp:s:ssec-hull}

\begin{definition}\label{dp:def:hull}
For a normed space $P$ and $x\in P$ put
\[
 \mathsf u(x)=\inf\bigl\{\max_l\norm{x_l}:\ x=\sum_lx_l\text{ a finite sum}\bigr\}.
\]
This is an ultrametric seminorm with $\mathsf u\le\norm\cdot$: concatenate decompositions for $\mathsf u(x+y)\le\max(\mathsf u(x),\mathsf u(y))$. Let $P^{\mathsf u}=P/\ker\mathsf u$, normed by $\mathsf u$, and let $\operatorname{pr}\colon P\to P^{\mathsf u}$ be the quotient map, which is contractive. For a law $\tilde a$ on $P^{\mathsf u}$ we write $\norm{\tilde a}_{\mathsf u}$ for its law norm. In \cite[proof of Theorem~14.39(4)]{I} this seminorm is written $\nu$; here $\nu$ is a norm on the space $V$ of Section~\ref{dp:sec:nonambient}.
\end{definition}

On $\ell^1(J,K)$ the seminorm $\mathsf u$ is the supremum norm. Indeed, for $\xi\in J$ we have $\abs{x_\xi}\le\max_l\abs{(x_l)_\xi}\le\max_l\norm{x_l}$. Conversely, let $\eta>0$, list the support of $x$, and cut it greedily into blocks of $\ell^1$-mass at most $\sup_\xi\abs{x_\xi}+\eta$. Each block but the last has mass above $\eta$, so there are finitely many. The last block may have infinite support, and it still lies in $\ell^1(J,K)$. Hence the completion of $\ell^1(J,K)^{\mathsf u}$ is $c_0(J,\Kh)$.

\begin{proposition}[ultrametric targets see only the hull]\label{dp:prop:hull}
Let $P$ be normed and $W$ ultrametric.
\begin{enumerate}
\item A bounded $n$\nobreakdash-linear map $\beta\colon P^n\to W$ satisfies $\norm{\beta(x_1,\dots,x_n)}\le\norm\beta\prod_c\mathsf u(x_c)$. So it factors through $(P^{\mathsf u})^n$, with the same norm.
\item A law $a\colon P\to W$ of degree $n$ satisfies $\norm{a^{[\alpha]}(x_1,\dots,x_r)}\le\lvert\!\lvert\!\lvert a\rvert\!\rvert\!\rvert\prod_c\mathsf u(x_c)^{\alpha_c}$. So $a=\tilde a\circ\operatorname{pr}$ for a unique law $\tilde a$ on $P^{\mathsf u}$, with $\lvert\!\lvert\!\lvert \tilde a\rvert\!\rvert\!\rvert=\lvert\!\lvert\!\lvert a\rvert\!\rvert\!\rvert$ and
\[
 \norm{\tilde a}_{\mathsf u}\le\lvert\!\lvert\!\lvert a\rvert\!\rvert\!\rvert\le\theta_n(K)\mathfrak s(P^{\mathsf u})^{-n}\norm{\tilde a}_{\mathsf u}.
\]
Also $\operatorname{lift}(a)=\operatorname{lift}(\tilde a)$.
\end{enumerate}
\end{proposition}

\begin{proof}
\emph{Step 1 (item (1)).} Expand each slot over a decomposition $x_c=\sum_lx_{c,l}$. As $W$ is ultrametric, the norm of the sum is at most the maximum of the norms of the terms. Take infima, one slot at a time.

\emph{Step 2 (the bound of item (2)).} Let $x_c=\sum_lx_{c,l}$. By~\eqref{dp:s:eq:slots} with all $\lambda_{c,l}=1$, and~\eqref{dp:laws:strong}, each term has norm at most $\lvert\!\lvert\!\lvert a\rvert\!\rvert\!\rvert\prod_c\max_l\norm{x_{c,l}}^{\alpha_c}$. As $W$ is ultrametric, so does the sum. Take infima.

\emph{Step 3 (descent).} If $\mathsf u(y)=0$, then $a(x+y)=\sum_{i\le n}a^{[(n-i,i)]}(x,y)$, and by Step~2 the terms with $i\ge1$ vanish. So $\tilde a(\operatorname{pr}x)=a(x)$ is well defined, and $\tilde a(\sum_cs_c\operatorname{pr}x_c)=a(\sum_cs_cx_c)$ shows that it is a law with $\tilde a^{[\alpha]}(\operatorname{pr}x_1,\dots,\operatorname{pr}x_r)=a^{[\alpha]}(x_1,\dots,x_r)$. Step~2 gives $\lvert\!\lvert\!\lvert \tilde a\rvert\!\rvert\!\rvert\le\lvert\!\lvert\!\lvert a\rvert\!\rvert\!\rvert$, and Lemma~\ref{dp:lem:strong-formal}(1) with $T_2=\operatorname{pr}$ gives the reverse. Always $\norm{\tilde a}_{\mathsf u}\le\lvert\!\lvert\!\lvert \tilde a\rvert\!\rvert\!\rvert$, and the last inequality is Lemma~\ref{dp:lem:strong-coefficients}(2) on $P^{\mathsf u}$. By item (1), lifts of $a$ and of $\tilde a$ correspond, with the same norm.
\end{proof}

The hull can vanish. On $K$-valued step functions on $[0,1)$ with $\norm h=\int\abs h$, cutting $h$ into $\mathsf N$ pieces shows $\mathsf u\equiv0$. Then every law of degree $n\ge1$ into an ultrametric space vanishes, and DP germs into ultrametric spaces are constant near each point.

\begin{proposition}[the strong norm interpolates]\label{dp:prop:strong-interpolates}
\leavevmode
\begin{enumerate}
\item Let $W$ be ultrametric and $P$ any normed space. A map analytic at $x_0$ with radius $R$ and bound $M$ is sDP at $x_0$ with radius $R$ and bound $M$.
\item Let $\operatorname{char}K=0$, with $P,W$ arbitrary. A map sDP at $x_0$ with radius $R$ and bound $M$ is analytic at $x_0$ with radius $\mathfrak n_KR$ and bound $M$.
\item Let $P$ be ultrametric and $W$ arbitrary. A map DP at $x_0$ with radius $R$ and bound $M$ is sDP at $x_0$ with radius $R/C_1$ and bound $M$; with radius $\theta_K\mathfrak s(P)R$ if $W$ is also ultrametric. So on ultrametric $P$ the DP germs and the sDP germs are the same.
\item Let $K$ be complete and not spherically complete, $\operatorname{char}K=p\le k$, and $E,D$ as in \cite[Theorem~7.2(2)]{I}. Then $A^k_{E,D;K}$ is sDP at every point and analytic at no point. Over every $K$, the germ $\mathsf f$ of Example~\ref{dp:ex:l1} is DP at $0$ and not sDP at $0$.
\end{enumerate}
\end{proposition}

\begin{proof}
\emph{Step 1 (item (1)).} If $\beta_n$ expand the map, then $\beta_n$ has diagonal $a_n$ (Lemma~\ref{dp:lem:coefficients-unique}(2)). Apply Lemma~\ref{dp:lem:strong-formal}(5).

\emph{Step 2 (item (2)).} The constant $\mathfrak n_K\in(0,1]$ is that of Section~\ref{dp:sec:agreement}. For $n\ge1$ put $\beta_n=a_n^{[1^n]}/n!$. It is $n$\nobreakdash-linear (Lemma~\ref{dp:lem:elementary}(4)) and has diagonal $a_n$ (Lemma~\ref{dp:lem:hasse-law}(5)). By~\eqref{dp:laws:strong}, $\norm{\beta_n}\le\abs{n!}^{-1}\lvert\!\lvert\!\lvert a_n\rvert\!\rvert\!\rvert\le M(\mathfrak n_KR)^{-n}$. With $\beta_0=a_0$ these expand the map for $\norm h<\mathfrak n_KR\le R$.

\emph{Step 3 (item (3)).} The first radius is Lemma~\ref{dp:lem:strong-formal}(3). For the second, Lemma~\ref{dp:lem:strong-coefficients}(2) gives $\lvert\!\lvert\!\lvert a_n\rvert\!\rvert\!\rvert\le\theta_K^{1-n}\mathfrak s(P)^{-n}MR^{-n}\le M(\theta_K\mathfrak s(P)R)^{-n}$ for $n\ge1$, as $\theta_K\le1$. For $n=0$ use $\norm{a_0}\le M$. As $\theta_K\mathfrak s(P)\le1$, the expansion holds on the smaller ball. The converse direction is $\norm\cdot\le\lvert\!\lvert\!\lvert \cdot\rvert\!\rvert\!\rvert$.

\emph{Step 4 (item (4)).} By Step~3 of the proof of Theorem~\ref{dp:thm:calculus}, $A^k_{E,D;K}$ followed by an isometry is the diagonal of a $k$\nobreakdash-linear map of norm at most $1$ into $\cL(\Alt^k(D;K),\Mult^k(E;K))$, which is ultrametric. So $\lvert\!\lvert\!\lvert A^k_{E,D;K}\rvert\!\rvert\!\rvert\le1$ by Lemma~\ref{dp:lem:strong-formal}(1),(5), and a law is sDP at every point by Lemma~\ref{dp:lem:strong-formal}(3). It is analytic at no point by \cite[Theorem~7.2(2)]{I}. For $\mathsf f$, the law $a_n$ has coefficient $a_n^{[1^n]}(e_1,\dots,e_n)=\mathsf e_n$, so $\lvert\!\lvert\!\lvert a_n\rvert\!\rvert\!\rvert\ge n^n$. By Lemma~\ref{dp:lem:coefficients-unique}(1) the $a_n$ are the only possible laws of $\mathsf f$ at $0$, and $n^nR^n$ is unbounded for every $R>0$.
\end{proof}

So in characteristic $0$ the analytic and the sDP germs agree if $W$ or $P$ is ultrametric. They can differ otherwise.

\begin{example}[an analytic germ that is not strongly DP]\label{dp:ex:an-not-sdp}
Let $P=\ell^1(\N,K)$, over any $K$. Let $W$ be the $\ell^1$-sum of the spaces $\ell^1(\N^n,K)$, $n\ge1$, with inclusions $\iota_n$. Put $f(h)=\sum_{n\ge1}\iota_n(h^{\otimes n})$ for $\norm h<1$. Then $f$ is analytic at $0$ with radius $1$ and bound $1$, and it is not sDP at $0$. In characteristic $0$, by Proposition~\ref{dp:prop:strong-interpolates}(2), the sDP germs from $P$ to $W$ are strictly fewer than the analytic ones.
\end{example}

\begin{proof}[Verification]
The maps $(h^1,\dots,h^n)\mapsto\iota_n(h^1\otimes\dots\otimes h^n)$ have norm at most $1$ and expand $f$. The sum lies in $W$, as its coordinates lie in $K$ and $\sum_n\norm h^n<\infty$. The law of degree $n$ of $f$ is $a_n(h)=\iota_n(h^{\otimes n})$. The coefficient of $s_1\cdots s_n$ in $a_n(\sum_{\mathsf i\le n}s_{\mathsf i}e_{\mathsf i})$ is a sum of $n!$ distinct unit vectors, of $\ell^1$-norm $n!$ in every characteristic. So $\lvert\!\lvert\!\lvert a_n\rvert\!\rvert\!\rvert\ge n!$, and $n!R^n$ is unbounded for every $R>0$.
\end{proof}

\subsection{Ultrametric targets}\label{dp:s:ssec-targets}

With ultrametric Banach targets the class has a description with no reference to the strong norm. Recall the canonical cover $q_P$ (Definition~\ref{dp:def:cover}), and $\mathfrak c_K$ from Section~\ref{dp:sec:conv}.

\begin{theorem}[strongly DP germs into ultrametric spaces]\label{dp:thm:strong-um}
Let $W$ be an ultrametric Banach space, $P$ a normed space, $U\subseteq P$ open, $f\colon U\to W$ and $x_0\in U$. The following are equivalent.
\begin{enumerate}
\item $f$ is sDP at $x_0$.
\item $f\circ q_P$ is analytic at some preimage of $x_0$, equivalently at every preimage.
\item $f\circ Q$ is analytic at some preimage of $x_0$, for some bounded linear surjection $Q$ with a preimage constant.
\item $j\circ f$ is analytic at $x_0$ for some linear isometry $j$ of $W$ into an ultrametric normed space.
\item For some $R,M$, $f$ is DP at $x_0$ with radius $R$ and bound $M$, and its coefficient laws satisfy $\norm{\tilde a_n}_{\mathsf u}\le MR^{-n}$.
\end{enumerate}
Radius $R$ and bound $M$ in~\textup{(1)} give the same in~\textup{(2)} and~\textup{(5)}, and radius $\abs\varpi\mathfrak c_KR$ and bound $M$ in~\textup{(4)}. Radius $R$ and bound $M$ in~\textup{(3)}, with preimage constant $\mathsf c$, give radius $R/\mathsf c$ and bound $M$ in~\textup{(1)}. Radius $R$ and bound $M$ in~\textup{(4)} give the same in~\textup{(1)}. Radius $R$ and bound $M$ in~\textup{(5)} give radius $\theta_K\mathfrak s(P^{\mathsf u})R$ and bound $M$ in~\textup{(1)}.
\end{theorem}

\begin{proof}
\emph{Step 1 ((1) implies (2), (2) implies (3)).} Let $a_n$ be the laws of $f$ at $x_0$. The proof of Lemma~\ref{dp:lem:lift-cover} gives a lift of $a_n\circ q_P$ of norm at most $\lvert\!\lvert\!\lvert a_n\rvert\!\rvert\!\rvert\le MR^{-n}$. As $\norm{q_P}\le1$, these lifts expand $f\circ q_P$ on the ball of radius $R$ around every preimage. For (3), take $Q=q_P$.

\emph{Step 2 ((3) implies (1)).} The map $f\circ Q$ is DP with the same data (Theorem~\ref{dp:thm:calculus}(3)). By Proposition~\ref{dp:prop:source-descent}(1), $f$ is DP at $x_0$ with radius $R/\mathsf c$ and with laws $a_n$ such that $a_n\circ Q$ has a lift of norm at most $MR^{-n}$. Lemma~\ref{dp:lem:strong-formal}(5),(2) give $\lvert\!\lvert\!\lvert a_n\rvert\!\rvert\!\rvert\le\mathsf c^nMR^{-n}$.

\emph{Step 3 ((1) and (5)).} By Proposition~\ref{dp:prop:hull}(2), the laws descend. If (1) holds, then $\norm{\tilde a_n}_{\mathsf u}\le\lvert\!\lvert\!\lvert a_n\rvert\!\rvert\!\rvert\le MR^{-n}$, and $\norm{a_n}\le\lvert\!\lvert\!\lvert a_n\rvert\!\rvert\!\rvert$. If (5) holds, then for $n\ge1$,
\[
 \lvert\!\lvert\!\lvert a_n\rvert\!\rvert\!\rvert\le\theta_K^{1-n}\mathfrak s(P^{\mathsf u})^{-n}\norm{\tilde a_n}_{\mathsf u}\le M\bigl(\theta_K\mathfrak s(P^{\mathsf u})R\bigr)^{-n},
\]
and $\theta_K\mathfrak s(P^{\mathsf u})\le1$.

\emph{Step 4 ((5) implies (4)).} As $W$ is complete, $\tilde f=\sum_n\tilde a_n$ converges on the ball of radius $R$ of $P^{\mathsf u}$, and it is DP at $0$ with radius $R$ and bound $M$. As $P^{\mathsf u}$ and $W$ are ultrametric, \cite[Theorem~14.39(1)]{I} gives a linear isometry $j$ of $W$ into an ultrametric normed space $Z$ and $n$\nobreakdash-linear maps $\beta_n$ on $P^{\mathsf u}$, with $\norm{\beta_n}\le M(\abs\varpi\mathfrak c_KR)^{-n}$, that expand $j\circ\tilde f$. Since $\mathsf u\le\norm\cdot$, the maps $\beta_n\circ\operatorname{pr}^{\times n}$ expand $j\circ f$ at $x_0$.

\emph{Step 5 ((4) implies (1)).} Let $\beta_n\colon P^n\to Z$ expand $j\circ f$, with $\norm{\beta_n}\le MR^{-n}$ and $Z$ ultrametric. As $W$ is complete, Theorem~\ref{dp:thm:target-descent}(2) gives that $f$ is DP at $x_0$, with laws $a_n$ such that $j\circ a_n$ is the diagonal of $\beta_n$ (Lemma~\ref{dp:lem:coefficients-unique}(2)). By Lemma~\ref{dp:lem:strong-formal}(1),(5), $\lvert\!\lvert\!\lvert a_n\rvert\!\rvert\!\rvert=\lvert\!\lvert\!\lvert j\circ a_n\rvert\!\rvert\!\rvert\le\norm{\beta_n}$.
\end{proof}

\begin{corollary}[the two gaps]\label{dp:cor:two-gaps}
Let $W$ be an ultrametric Banach space and $P$ a normed space. Germs from $P$ to $W$ satisfy: analytic implies sDP, and sDP implies DP.
\begin{enumerate}
\item \textup{(The first gap.)} Let $f$ be sDP at $x_0$ with laws $a_n$. Then $f$ is analytic at $x_0$ if and only if $\sum_n\tilde a_n$ is analytic at $0$ on $P^{\mathsf u}$. So the first gap is the question between the ultrametric spaces $P^{\mathsf u}$ and $W$. It closes if $\operatorname{char}K=0$; if $K$ is discretely valued; if $W$ has an equivalent ultrametric spherically complete norm; and if $K$ is complete and the completion of $P^{\mathsf u}$ is linearly homeomorphic to $c_0(J,K)$. For $K$ complete, not spherically complete, of characteristic $p$, there are ultrametric $P$ and $W$ for which it does not close.
\item \textup{(The second gap.)} A DP germ is sDP if and only if the numbers $\norm{\tilde a_n}_{\mathsf u}$ have a geometric bound. By Lemma~\ref{dp:lem:strong-formal}(3), the excess $\lvert\!\lvert\!\lvert a\rvert\!\rvert\!\rvert/\norm a$ is at most $(C_1n)^n$. On $\ell^1(\N,K)$ the germ $\mathsf f$ has $\norm{a_n}<\abs\varpi^{-1}$ and $\lvert\!\lvert\!\lvert a_n\rvert\!\rvert\!\rvert\ge n^n$.
\item Let $\operatorname{char}K=0$. On every source, the analytic germs into ultrametric Banach spaces have source descent, and target descent along isometries between ultrametric Banach spaces, since they are the sDP germs.
\end{enumerate}
\end{corollary}

\begin{proof}
\emph{Step 1 (the inclusions).} Proposition~\ref{dp:prop:strong-interpolates}(1) and $\norm\cdot\le\lvert\!\lvert\!\lvert \cdot\rvert\!\rvert\!\rvert$.

\emph{Step 2 (item (1)).} If $\beta_n$ expand $f$, then $\beta_n$ has diagonal $a_n$, and Proposition~\ref{dp:prop:hull}(1) gives maps on $P^{\mathsf u}$ of the same norms with diagonals $\tilde a_n$. Conversely, compose an expansion on $P^{\mathsf u}$ with $\operatorname{pr}$, which is contractive. By Theorem~\ref{dp:thm:strong-um}(5), $\sum_n\tilde a_n$ is DP at $0$ on $P^{\mathsf u}$. The first three cases are Theorem~\ref{dp:thm:agreement-char0}(2) and Theorem~\ref{dp:thm:sph-agreement}(2),(1). For the fourth, let $K$ be complete. By \cite[Lemma~Q.13(1)]{I}, $\tilde a_n$ extends to the completion of $P^{\mathsf u}$ with the same bound, and Theorem~\ref{dp:thm:c0-agreement}(2), transported through the homeomorphism with its constant, gives a lift with a geometric bound. That the gap does not close is \cite[Proposition~14.40(3)]{I}, where $P$ is ultrametric, so that $P^{\mathsf u}=P$.

\emph{Step 3 (items (2) and (3)).} Item (2) is Theorem~\ref{dp:thm:strong-um}(5), Example~\ref{dp:ex:l1}(1) and Proposition~\ref{dp:prop:strong-interpolates}(4). For item (3), item (1) gives that the analytic germs are the sDP germs. If $f\circ Q$ is analytic, Theorem~\ref{dp:thm:strong-um}, (3) implies (1), shows that $f$ is sDP. If $j\circ f$ is analytic for an isometry $j$ into an ultrametric Banach space, (4) implies (1) shows the same.
\end{proof}

\subsection{Composition and inverse functions}\label{dp:s:ssec-composition}

Let $f$ map an open subset of $P$ into $W$ and $g$ an open subset of $W$ into $Y$, with $f(x_0+h)=\sum_na_n(h)$ for $\norm h<R$, $g(f(x_0)+w)=\sum_mb_m(w)$ for $\norm w<R'$, $\lvert\!\lvert\!\lvert a_n\rvert\!\rvert\!\rvert\le MR^{-n}$ for $n\ge1$, $\lvert\!\lvert\!\lvert b_m\rvert\!\rvert\!\rvert\le M'R'^{-m}$ and $M>0$. For a finitely supported sequence $\alpha=(\alpha_1,\alpha_2,\dots)$ put $\abs\alpha=\sum_n\alpha_n$ and $\deg\alpha=\sum_nn\alpha_n$. By \cite[Lemma~Q.3]{I}, the laws of $g\circ f$ at $x_0$ are
\[
 \Theta_d=\sum_{\deg\alpha=d}\phi_\alpha,\qquad\phi_\alpha(h)=b_{\abs\alpha}^{[\alpha]}(a_1(h),a_2(h),\dots).
\]
The algebra of composition is~\cite[VI~\S8, Prop.~VI.6, Cor.; Prop.~I.2; II~\S1]{Ro}; the bounds below are the estimates we need. Put $R_{\mathrm c}=R\min(1,R'/M)$. Put $R_{\mathrm e}=RR'/(M+R')$, and $R_{\mathrm e}=R_{\mathrm c}$ if $W$ is ultrametric.

\begin{theorem}[composition]\label{dp:thm:strong-composition}
\leavevmode
\begin{enumerate}
\item \textup{(Coefficients, all norms.)} For $h_1,\dots,h_r\in P$ and $\abs\delta=\deg\alpha$,
\begin{equation}\label{dp:s:eq:composite-coefficient}
 \phi_\alpha^{[\delta]}(h_1,\dots,h_r)=\sum_\varepsilon b_{\abs\alpha}^{[\varepsilon]}\bigl((a_n^{[\gamma]}(h_1,\dots,h_r))_{n,\gamma}\bigr),
\end{equation}
over the arrays $\varepsilon_{n,\gamma}\ge0$, with $\abs\gamma=n$, such that $\sum_\gamma\varepsilon_{n,\gamma}=\alpha_n$ and $\sum_{n,\gamma}\varepsilon_{n,\gamma}\gamma=\delta$. Each term has norm at most $\lvert\!\lvert\!\lvert b_{\abs\alpha}\rvert\!\rvert\!\rvert\prod_n\lvert\!\lvert\!\lvert a_n\rvert\!\rvert\!\rvert^{\alpha_n}\prod_c\norm{h_c}^{\delta_c}$.
\item \textup{($Y$ ultrametric.)} We have $\lvert\!\lvert\!\lvert \Theta_d\rvert\!\rvert\!\rvert\le M'R_{\mathrm c}^{-d}$ for $d\ge1$, and $g\circ f(x_0+h)=\sum_d\Theta_d(h)$ for $\norm h<R_{\mathrm e}$. So $g\circ f$ is sDP at $x_0$ with radius $R_{\mathrm e}$ and bound $M'$.
\item \textup{($W$ ultrametric, any $Y$.)} We have $\lvert\!\lvert\!\lvert \Theta_d\rvert\!\rvert\!\rvert\le M'(2C_1/R_{\mathrm c})^d$ for $d\ge1$. So $g\circ f$ is sDP at $x_0$.
\item \textup{($P$ ultrametric, any $W,Y$.)} The map $g\circ f$ is sDP at $x_0$.
\item \textup{($g$ a law, all norms.)} Let $g$ be a law of degree $k\ge1$ and $f(x_0)=0$. Then $\lvert\!\lvert\!\lvert \Theta_d\rvert\!\rvert\!\rvert\le\binom{d-1}{k-1}k^d\lvert\!\lvert\!\lvert g\rvert\!\rvert\!\rvert M^kR^{-d}$. So $g\circ f$ is sDP at $x_0$. If $f(x_0)\ne0$, then $g\circ f$ is still sDP at $x_0$. So sDP is stable under postcomposition with laws.
\item \textup{(Formal.)} Let $a\colon P\to W$ and $\bar a\colon P\to W'$ be laws of the same degree. The pairing $(a,\bar a)$, into $W\oplus W'$ with the maximum norm, has $\lvert\!\lvert\!\lvert (a,\bar a)\rvert\!\rvert\!\rvert=\max(\lvert\!\lvert\!\lvert a\rvert\!\rvert\!\rvert,\lvert\!\lvert\!\lvert \bar a\rvert\!\rvert\!\rvert)$. Constants are sDP, and so are the sums and pairings of sDP germs and their compositions with bounded linear maps on either side.
\end{enumerate}
So $g\circ f$ is sDP if one of $P,W,Y$ has an equivalent ultrametric norm, since a renorming with constant $\mathsf c$ multiplies strong norms of laws of degree $n$ by at most $\mathsf c^n$ on the source and by $\mathsf c$ on the target.
\end{theorem}

\begin{proof}
\emph{Step 1 (item (1)).} Put $v_n=\sum_{\abs\gamma=n}s^\gamma a_n^{[\gamma]}(h_1,\dots,h_r)$, so that $a_n(\sum_cs_ch_c)=v_n$. By~\eqref{dp:s:eq:slots} with the scalars $s^\gamma$, $b_{\abs\alpha}^{[\alpha]}(v_1,v_2,\dots)$ is the sum over the arrays with $\sum_\gamma\varepsilon_{n,\gamma}=\alpha_n$ of $s^{\sum\varepsilon_{n,\gamma}\gamma}b_{\abs\alpha}^{[\varepsilon]}((a_n^{[\gamma]}(h_1,\dots,h_r))_{n,\gamma})$. Read off the coefficient of $s^\delta$. The bound is~\eqref{dp:laws:strong}, twice.

\emph{Step 2 (item (2): bound).} Let $\deg\alpha=d$. As $Y$ is ultrametric, item (1) gives $\lvert\!\lvert\!\lvert \phi_\alpha\rvert\!\rvert\!\rvert\le M'R'^{-\abs\alpha}\prod_n(MR^{-n})^{\alpha_n}=M'(M/R')^{\abs\alpha}R^{-d}$. As $\abs\alpha\le d$, this is at most $M'R_{\mathrm c}^{-d}$. The sum over $\alpha$ is finite, and $Y$ is ultrametric.

\emph{Step 3 (item (2): expansion).} Let $\norm h<R_{\mathrm e}$ and $\mathsf x=\norm h/R<1$. Then $\norm{a_n(h)}\le M\mathsf x^n$, so $w=\sum_{n\ge1}a_n(h)$ has $\norm w\le M\mathsf x/(1-\mathsf x)<R'$, and $\norm w\le M\mathsf x<R'$ if $W$ is ultrametric. By~\eqref{dp:laws:strong}, $\norm{\phi_\alpha(h)}\le M'(M/R')^{\abs\alpha}\mathsf x^{\deg\alpha}$. As $R_{\mathrm e}\le R_{\mathrm c}$, we have $\mathsf x\max(1,M/R')<1$. So $\phi_\alpha(h)\to0$ along the finite-to-one map $\alpha\mapsto\deg\alpha$. Hence the family $(\phi_\alpha(h))_\alpha$ is summable in the ultrametric $\widehat Y$, and it may be grouped in any way. By Lemma~\ref{dp:lem:elementary}(3), $b_m(\sum_{n\le\overline n}a_n(h))$ is the sum of the $\phi_\alpha(h)$ with $\abs\alpha=m$ and support in $[\overline n]$. As $b_m$ is continuous, grouping by $\abs\alpha=m$ gives $b_m(w)$, and $\sum_mb_m(w)=g\circ f(x_0+h)$. Grouping by $\deg\alpha=d$ gives $\sum_d\Theta_d(h)$. Both sides lie in $Y$. Finally, $\lvert\!\lvert\!\lvert \Theta_d\rvert\!\rvert\!\rvert\le M'R_{\mathrm c}^{-d}\le M'R_{\mathrm e}^{-d}$.

\emph{Step 4 (item (3)).} By Proposition~\ref{dp:prop:hull}(2), $a_n=\tilde a_n\circ\operatorname{pr}$, so $\phi_\alpha=\tilde\phi_\alpha\circ\operatorname{pr}$ for a law $\tilde\phi_\alpha$ on $P^{\mathsf u}$. By~\eqref{dp:laws:strong} and Proposition~\ref{dp:prop:hull}(2), $\norm{\tilde\phi_\alpha}_{\mathsf u}\le\lvert\!\lvert\!\lvert b_{\abs\alpha}\rvert\!\rvert\!\rvert\prod_n\lvert\!\lvert\!\lvert a_n\rvert\!\rvert\!\rvert^{\alpha_n}\le M'R_{\mathrm c}^{-d}$. At most $2^{d-1}$ sequences $\alpha$ have degree $d$. On the ultrametric source $P^{\mathsf u}$, Lemma~\ref{dp:lem:strong-formal}(3) gives $\lvert\!\lvert\!\lvert \tilde\Theta_d\rvert\!\rvert\!\rvert\le C_1^d2^{d-1}M'R_{\mathrm c}^{-d}$, and $\lvert\!\lvert\!\lvert \Theta_d\rvert\!\rvert\!\rvert\le\lvert\!\lvert\!\lvert \tilde\Theta_d\rvert\!\rvert\!\rvert$ by Lemma~\ref{dp:lem:strong-formal}(1). By \cite[Lemma~Q.3]{I} the $\Theta_d$ expand $g\circ f$ near $x_0$.

\emph{Step 5 (item (4)).} By Proposition~\ref{dp:prop:strong-interpolates}(3) both germs are DP, $g\circ f$ is DP by Theorem~\ref{dp:thm:calculus}(2), and it is sDP by Proposition~\ref{dp:prop:strong-interpolates}(3).

\emph{Step 6 (item (5)).} Only $b_k=g$ occurs, so $\abs\alpha=k$, and at most $\binom{d-1}{k-1}$ sequences $\alpha$ have $\abs\alpha=k$ and degree $d$. Fix one. An array $\varepsilon$ of~\eqref{dp:s:eq:composite-coefficient} comes from at least one ordered decomposition of $\delta$ into $k$ vectors, by listing the $\gamma$ with their multiplicities. There are at most $\prod_ck^{\delta_c}=k^d$ such decompositions. So $\lvert\!\lvert\!\lvert \phi_\alpha\rvert\!\rvert\!\rvert\le k^d\lvert\!\lvert\!\lvert g\rvert\!\rvert\!\rvert\prod_n(MR^{-n})^{\alpha_n}=k^d\lvert\!\lvert\!\lvert g\rvert\!\rvert\!\rvert M^kR^{-d}$. By \cite[Lemma~Q.3]{I} the $\Theta_d$ expand $g\circ f$ near $x_0$. If $f(x_0)=w_0$, then $g(w_0+w)=\sum_{i\le k}g^{[(k-i,i)]}(w_0,w)$ is a finite sum of laws in $w$ (Lemma~\ref{dp:lem:strong-formal}(4)). The term $i=0$ is constant, and the first case applies to the others and to $f-w_0$. Conclude by item (6).

\emph{Step 7 (item (6)).} The coefficients of a pair are the pairs of coefficients. Sums use the triangle inequality, with the smaller radius. Bounded linear maps are Lemma~\ref{dp:lem:strong-formal}(1).
\end{proof}

The next theorem shows that composition fails for general norms. In~\eqref{dp:s:eq:composite-coefficient} the arrays are counted, and the norm of an $\ell^1$-target adds their terms. Let $\mathcal D$ be the set of unordered pairs $\{\mathsf i,\mathsf j\}$ with $\mathsf i\le\mathsf j$ in $\N$, with a total order. Let $\mathcal M_m$ be the set of multisets of size $m$ in $\mathcal D$, and $\mathcal M=\bigsqcup_{m\ge2}\mathcal M_m$. A sum of multisets adds multiplicities. Let $\mathsf X=\ell^1(\N\sqcup\mathcal D\sqcup\mathcal M,K)$, which is isometric to $\ell^1(\N,K)$, with coordinates $x=(x_\N,x_{\mathcal D},x_{\mathcal M})$. For $\mathsf m\in\mathcal M_m$ and $w\in\ell^1(\mathcal D,K)$ put $w^{\mathsf m}=\prod_{\varsigma\in\mathsf m}w_\varsigma$, with multiplicities. Define
\[
 \mathbf{Sq}(h)=\sum_{\mathsf i\le\mathsf j}h_{\mathsf i}h_{\mathsf j}e_{\{\mathsf i,\mathsf j\}},\qquad\mathbf b_m(w)=\sum_{\mathsf m\in\mathcal M_m}w^{\mathsf m}e_{\mathsf m},\qquad\mathbf g(w)=\sum_{m\ge2}\mathbf b_m(w)\quad(\norm w<1),
\]
\[
 F_1(x)=x+\mathbf{Sq}(x_\N),\qquad F_2(x)=x+\mathbf g(x_{\mathcal D}),\qquad\mathsf S(x)=x-\mathbf{Sq}(x_\N)-\mathbf g(x_{\mathcal D}).
\]

\begin{theorem}[failure on $\ell^1$]\label{dp:thm:strong-failure}
Let $K$ be arbitrary, of any characteristic.
\begin{enumerate}
\item We have $\lvert\!\lvert\!\lvert \mathbf{Sq}\rvert\!\rvert\!\rvert\le1$ and $\lvert\!\lvert\!\lvert \mathbf b_m\rvert\!\rvert\!\rvert\le1$. The maps $\mathbf{Sq},\mathbf g,F_1,F_2,\mathsf S$ are sDP at every point of their domains, with $\norm{x_{\mathcal D}}<1$ where $\mathbf g$ enters, and with values in $\mathsf X$, reading $\ell^1(\mathcal D,K)$ and $\ell^1(\mathcal M,K)$ inside it. All are analytic at $0$.
\item The map $\mathbf g\circ\mathbf{Sq}$ is analytic at $0$ and not sDP at $0$: its law of degree $2m$, $\mathbf t_m=\mathbf b_m\circ\mathbf{Sq}$, has $\lvert\!\lvert\!\lvert \mathbf t_m\rvert\!\rvert\!\rvert\ge(2m-1)!!$.
\item \textup{(Composition fails.)} $F_1,F_2$ are local diffeomorphisms at $0$, sDP with sDP inverses, and $F_2\circ F_1$ is not sDP at $0$.
\item \textup{(The inverse function theorem fails.)} $\mathsf S$ is sDP at $0$ with $D\mathsf S(0)=\id$. Its local inverse is $F_2\circ F_1$, which is analytic at $0$ and not sDP at $0$.
\end{enumerate}
\end{theorem}

\begin{proof}
\emph{Step 1 ($\mathbf{Sq}$).} The coefficient of $\mathbf{Sq}$ of index $(1,1)$ at $(y,y')$ has coordinates $y_{\mathsf i}y'_{\mathsf j}+y_{\mathsf j}y'_{\mathsf i}$ for $\mathsf i<\mathsf j$ and $2y_{\mathsf i}y'_{\mathsf i}$ for $\mathsf i=\mathsf j$. As $\abs2\le1$, its $\ell^1$-norm is at most $\sum_{\mathsf i,\mathsf j}\abs{y_{\mathsf i}}\abs{y'_{\mathsf j}}=\norm y\norm{y'}$. Also $\norm{\mathbf{Sq}(y)}\le\norm y^2$. A lift is $\beta_{\mathbf{Sq}}(y,y')=\sum_{\mathsf i\le\mathsf j}y_{\mathsf i}y'_{\mathsf j}e_{\{\mathsf i,\mathsf j\}}$, of norm at most $1$.

\emph{Step 2 ($\mathbf b_m$).} The coordinate at $\mathsf m$ of the coefficient of $s^\alpha$ in $\mathbf b_m(\sum_cs_cw_c)$ is a combination, with integer coefficients, of the products $\prod_cw_c^{\mathsf m_c}$ over the splittings $\mathsf m=\mathsf m_1+\dots+\mathsf m_r$ with $\abs{\mathsf m_c}=\alpha_c$. As integers have absolute value at most $1$, its absolute value is at most the sum of the $\prod_c\abs{w_c}^{\mathsf m_c}$. Each tuple $(\mathsf m_c)_c$ determines $\mathsf m$, so summing over $\mathsf m$ gives at most $\prod_c\sum_{\abs{\mathsf m_c}=\alpha_c}\abs{w_c}^{\mathsf m_c}\le\prod_c\norm{w_c}^{\alpha_c}$. So $\mathbf b_m$ is a law into $\ell^1(\mathcal M_m,K)$. A lift is
\[
 \beta_m(z^1,\dots,z^m)=\sum_{\varsigma_1\le\dots\le\varsigma_m}z^1_{\varsigma_1}\cdots z^m_{\varsigma_m}e_{\varsigma_1+\dots+\varsigma_m},
\]
over nondecreasing words in $\mathcal D$, where $\varsigma_1+\dots+\varsigma_m$ is the multiset with these elements. It has norm at most $1$.

\emph{Step 3 (item (1)).} For $\norm w<1$ the series $\mathbf g(w)$ has summable norms and coordinates in $K$, so it lies in $\ell^1(\mathcal M,K)$, also over incomplete $K$. Let $\norm{w_0}<1$. For $i\ge0$ put $\mathbf g'_i(y)=\sum_{m\ge\max(i,2)}\mathbf b_m^{[(m-i,i)]}(w_0,y)$. The summands lie in the disjoint blocks $\ell^1(\mathcal M_m,K)$, so each coordinate of the sum is one element of $K$, and the same holds for the coefficients of $\mathsf g_i$. So $\mathbf g'_i$ is a law into $\ell^1(\mathcal M,K)$, and by Lemma~\ref{dp:lem:strong-formal}(4), $\lvert\!\lvert\!\lvert \mathbf g'_i\rvert\!\rvert\!\rvert\le\sum_{m\ge i}\norm{w_0}^{m-i}=(1-\norm{w_0})^{-1}$. For $\norm y<1-\norm{w_0}$, comparing coordinates gives $\mathbf g(w_0+y)=\sum_i\mathbf g'_i(y)$. So $\mathbf g$ is sDP at $w_0$ with radius $1-\norm{w_0}$. Theorem~\ref{dp:thm:strong-composition}(6) gives the statements for $F_1,F_2,\mathsf S$, and the lifts of Steps~1,~2 give analyticity at $0$.

\emph{Step 4 (item (2)).} In $\mathbf t_m(\sum_{\mathsf i\le2m}s_{\mathsf i}e_{\mathsf i})$ the coordinate at $\mathsf m\in\mathcal M_m$ is $\prod_{\{\mathsf i,\mathsf j\}\in\mathsf m}s_{\mathsf i}s_{\mathsf j}$. It contains $s_1\cdots s_{2m}$, with coefficient $1$, exactly when $\mathsf m$ is a perfect matching of $[2m]$. So the coefficient of $s_1\cdots s_{2m}$ is a sum of $(2m-1)!!$ distinct unit vectors, and its $\ell^1$-norm is the real number $(2m-1)!!$, in every characteristic. By Lemma~\ref{dp:lem:coefficients-unique}(1), $\mathbf t_m$ is the law of degree $2m$ of $\mathbf g\circ\mathbf{Sq}$ at $0$, and $(2m-1)!!R^{2m}$ is unbounded for every $R>0$. The map $(x^1,\dots,x^{2m})\mapsto\beta_m(\beta_{\mathbf{Sq}}(x^1,x^2),\dots,\beta_{\mathbf{Sq}}(x^{2m-1},x^{2m}))$ is a lift of $\mathbf t_m$ of norm at most $1$, which gives analyticity.

\emph{Step 5 (item (3)).} We have $F_1^{-1}(x)=x-\mathbf{Sq}(x_\N)$ and $F_2^{-1}(x)=x-\mathbf g(x_{\mathcal D})$. The $\mathcal M$-component of $F_2\circ F_1(x)$ is $x_{\mathcal M}+\mathbf g(x_{\mathcal D}+\mathbf{Sq}(x_\N))$. Precomposing with $h\mapsto(h,0,0)$ and projecting to $\mathcal M$ are contractive and give $\mathbf g\circ\mathbf{Sq}$. Lemma~\ref{dp:lem:strong-formal}(1) and item (2) conclude.

\emph{Step 6 (item (4)).} By Step~5, $\mathsf S=F_1^{-1}\circ F_2^{-1}$, since both sides are $(x_\N,x_{\mathcal D}-\mathbf{Sq}(x_\N),x_{\mathcal M}-\mathbf g(x_{\mathcal D}))$. So the local inverse is $F_2\circ F_1$, which is not sDP at $0$ by item (3). As $\mathbf{Sq}$ and $\mathbf g$ have no terms of degree below $2$, $D\mathsf S(0)=\id$. For analyticity, $F_2\circ F_1(v)=v+\beta_{\mathbf{Sq}}(v_\N,v_\N)+\sum_m\beta_m(y,\dots,y)$ with $y=v_{\mathcal D}+\beta_{\mathbf{Sq}}(v_\N,v_\N)$. For each word in $\{1,2\}^m$ with entries of sum $N$, fill the slots of $\beta_m$ with $v_{\mathcal D}$ for the entry $1$ and with $\beta_{\mathbf{Sq}}(v_\N,v'_\N)$ for the entry $2$, taking $N$ arguments in order. At most $2^N$ words have sum $N$, so the part of degree $N$ is the diagonal of an $N$\nobreakdash-linear map of norm at most $2^N$. These maps expand $\sum_m\beta_m(y,\dots,y)$ for $\norm v<1/2$. Each coordinate of the sum over $N$ involves finitely many degrees, so the sum lies in $\ell^1(\mathcal M,K)$, over every $K$.
\end{proof}

Whether the same failure occurs on spaces with no equivalent ultrametric norm and no complemented copy of $\ell^1(\N,K)$ is open (Question~\ref{dp:q:strong}). On ultrametric Banach spaces there is an inverse function theorem with no loss. The formal inverse is~\cite[Th.~VI.3, Prop.~VI.8]{Ro}; the bound is the estimate we need.

\begin{theorem}[inverse functions on ultrametric Banach spaces]\label{dp:thm:strong-inverse}
Let $P$ be an ultrametric Banach space, and let $\psi(z)=z-\sum_{m\ge2}\Phi_m(z)$ for $\norm z<R$, with laws $\Phi_m$ of degree $m$ and $\lvert\!\lvert\!\lvert \Phi_m\rvert\!\rvert\!\rvert\le MR^{-m}$. Let $\rho_0=\min(R,R^2/M)$, with $\rho_0=R$ if $M=0$. Then for every $\rho\le\rho_0$, $\psi$ is an isometry of $B(0,\rho)$ onto itself, and its inverse is $\chi=\sum_{d\ge1}G_d$ with laws $G_d$ of degree $d$ and $\lvert\!\lvert\!\lvert G_d\rvert\!\rvert\!\rvert\le\rho_0^{1-d}$.
\end{theorem}

\begin{proof}
\emph{Step 1 (the formal inverse).} Put $G_1=\id$, and for $d\ge2$
\begin{equation}\label{dp:s:eq:inverse-recursion}
 G_d=\sum_{\abs\alpha\ge2,\ \deg\alpha=d}\Phi_{\abs\alpha}^{[\alpha]}(G_1,G_2,\dots).
\end{equation}
Only $G_n$ with $n<d$ enter. For $m\ge2$ we have $MR^{-m}\le\rho_0^{1-m}$, as $\rho_0\le R^2/M$ and $\rho_0\le R$. By Theorem~\ref{dp:thm:strong-composition}(1) and the ultrametric inequality, with induction, $\lvert\!\lvert\!\lvert G_d\rvert\!\rvert\!\rvert\le\max_mMR^{-m}\rho_0^{m-d}\le\rho_0^{1-d}$.

\emph{Step 2 (right inverse).} For $\norm z<\rho_0$ we have $\norm{G_d(z)}\le\norm z(\norm z/\rho_0)^{d-1}$. As $P$ is complete, $\chi$ converges on $B(0,\rho_0)$ and maps each $B(0,\rho)$ into itself. Apply Theorem~\ref{dp:thm:strong-composition}(2) with $g=\psi$, $f=\chi$, outer laws $b_1=\id$ and $b_m=-\Phi_m$ for $m\ge2$, and all four constants of that theorem equal to $\rho_0$. The bounds hold by Step~1. Then $R_{\mathrm e}=R_{\mathrm c}=\rho_0$, and by~\eqref{dp:s:eq:inverse-recursion} the laws of $\psi\circ\chi$ are $\Theta_1=\id$ and $\Theta_d=0$ for $d\ne1$. So $\psi\circ\chi=\id$ on $B(0,\rho_0)$.

\emph{Step 3 (isometry).} Let $z,z'\in B(0,\rho)$, $z\ne z'$, and $\rho'=\max(\norm{z'},\norm{z-z'})<\rho_0$. By~\eqref{dp:eq:coefficient-identity} and~\eqref{dp:laws:strong}, $\Phi_m(z)-\Phi_m(z')=\sum_{i\ge1}\Phi_m^{[(m-i,i)]}(z',z-z')$ has norm at most $MR^{-m}\rho'^{m-1}\norm{z-z'}\le(M\rho'/R^2)\norm{z-z'}$, and $M\rho'/R^2<1$. So $\norm{\psi(z)-\psi(z')}=\norm{z-z'}$. As $\psi(0)=0$, $\psi$ maps $B(0,\rho)$ into itself. With Step~2 it is onto, and $\chi$ is its inverse.
\end{proof}

\begin{proposition}[manifolds and the alternating bundle]\label{dp:prop:strong-bundle}
\leavevmode
\begin{enumerate}
\item \textup{(Transition law, all norms, every characteristic.)} Let $E,D,F$ be normed spaces and $k\ge1$. We have $\lvert\!\lvert\!\lvert A^k_{E,D;F}\rvert\!\rvert\!\rvert\le k!$, and $\lvert\!\lvert\!\lvert A^k_{E,D;F}\rvert\!\rvert\!\rvert\le1$ if $F$ is ultrametric. So $A^k_{E,D;F}$ is sDP at every point, and if $\mathsf{tr}$ is sDP into $\cL(E,D)$, then $A^k_{E,D;F}\circ\mathsf{tr}$ is sDP. The same holds for the maps $u\mapsto\Lambda^k_\pi u$ and $u\mapsto\widehat\Lambda^k_\pi u$ of Theorem~\ref{dp:thm:calculus}(4), with $k!$, and for the joint action $B$ there, with $(k+1)!$.
\item \textup{(Models with an equivalent ultrametric norm.)} On such a model the sDP germs on open sets are the DP germs, for every normed target. So an atlas has sDP transitions if and only if it has DP transitions, and Corollary~\ref{dp:cor:bundles} applies.
\item \textup{(Models with a complemented $\ell^1(\N)$.)} Let $P$ have a complemented subspace linearly homeomorphic to $\ell^1(\N,K)$, so that after an equivalent norm $P=\mathsf X\oplus P_1$ with the maximum norm. Over every $K$ the sDP local diffeomorphisms of $P$ do not form a pseudogroup: on a small ball around $0$, take the charts $\id$, $F_1\times\id$, $(F_2\circ F_1)\times\id$. The changes from each chart to the next are sDP, and the change from the first to the third is not. So sDP compatibility of charts is not transitive, and maximal atlases with sDP transitions are not defined.
\end{enumerate}
\end{proposition}

\begin{proof}
\emph{Step 1 (item (1)).} By Step~3 of the proof of Theorem~\ref{dp:thm:calculus}, $A^k_{E,D;F}$ followed by an isometry is the diagonal of a $k$\nobreakdash-linear map of norm at most $1$ into $\cL(\Alt^k(D;F),\Mult^k(E;F))$, which is ultrametric if $F$ is. Lemma~\ref{dp:lem:strong-formal}(1),(5) give the bounds. The same step gives the other maps. Laws are sDP at every point by Lemma~\ref{dp:lem:strong-formal}(3), and Theorem~\ref{dp:thm:strong-composition}(5) gives the composite.

\emph{Step 2 (item (2)).} Proposition~\ref{dp:prop:strong-interpolates}(3), after the renorming, which changes the strong norm of a law of degree $n$ by at most $\mathsf c^n$.

\emph{Step 3 (item (3)).} By Proposition~\ref{dp:prop:strong-intrinsic}(2), being sDP does not depend on the choice of equivalent norm. By Theorem~\ref{dp:thm:strong-composition}(6) and Theorem~\ref{dp:thm:strong-failure}(1), $F_1\times\id$ and $F_2\times\id$ are sDP near $0$. If $(F_2\circ F_1)\times\id$ were sDP at $0$, restriction to $\mathsf X$ and projection onto $\mathsf X$ would make $F_2\circ F_1$ sDP at $0$, against Theorem~\ref{dp:thm:strong-failure}(3).
\end{proof}

For other non-ultrametric models the question is open (Question~\ref{dp:q:strong}).

\subsection{Descents}\label{dp:s:ssec-descents}

We use the germs and the properties (C1), (C2), (C3) of Section~\ref{dp:ssec:minimal}. Two forms of (C2) are used. (C2)$_{\rm u}$ is (C2) restricted to isometries between ultrametric Banach spaces. (C2)$_\times$ is (C2) with $Z$ an arbitrary Banach space, which is (C2) itself, written so for contrast.

\begin{theorem}[what strong DP loses]\label{dp:thm:strong-descent}
In every characteristic:
\begin{enumerate}
\item The sDP germs with Banach values have \textup{(C2)} and \textup{(C3)}. With radius $R$ and bound $M$, target descent keeps both, and source descent along a surjection with preimage constant $\mathsf c$ gives radius $R/\mathsf c$ and bound $M$. They can fail \textup{(C1)} when $P$ and $W$ are both non-ultrametric (Example~\ref{dp:ex:an-not-sdp}).
\item Among germs with ultrametric Banach targets and all normed sources, the sDP germs have \textup{(C1)}, \textup{(C2)}$_{\rm u}$, \textup{(C3)}. They are the smallest class with \textup{(C1)} and \textup{(C3)}, and also the smallest class with \textup{(C1)} and \textup{(C2)}$_{\rm u}$. They are strictly fewer than the DP germs on $\ell^1(\N,K)$.
\end{enumerate}
\end{theorem}

\begin{proof}
\emph{Step 1 (item (1)).} Let $j\circ f$ be sDP with laws $\mathsf a_n$, for an isometry $j$ of Banach spaces. By Theorem~\ref{dp:thm:target-descent}(2), $f$ is DP with the same radius, with laws $a_n$ such that $j\circ a_n=\mathsf a_n$. By Lemma~\ref{dp:lem:strong-formal}(1), $\lvert\!\lvert\!\lvert a_n\rvert\!\rvert\!\rvert=\lvert\!\lvert\!\lvert \mathsf a_n\rvert\!\rvert\!\rvert$. Let $f\circ Q$ be sDP with laws $b_n$. By Proposition~\ref{dp:prop:source-descent}(1), $b_n=a_n\circ Q$, and Lemma~\ref{dp:lem:strong-formal}(2) gives $\lvert\!\lvert\!\lvert a_n\rvert\!\rvert\!\rvert\le\mathsf c^n\lvert\!\lvert\!\lvert b_n\rvert\!\rvert\!\rvert$. In Example~\ref{dp:ex:an-not-sdp}, pass to $\widehat W$ by Lemma~\ref{dp:lem:strong-formal}(1) if $K$ is not complete.

\emph{Step 2 (item (2)).} \textup{(C1)} is Proposition~\ref{dp:prop:strong-interpolates}(1). Let $\cC$ be a class with \textup{(C1)} and \textup{(C3)}, and let $f$ be sDP at $x_0$. By Theorem~\ref{dp:thm:strong-um}(2), $f\circ q_P$ is analytic, so it is in $\cC$, and so is $f$. Let $\cC$ have \textup{(C1)} and \textup{(C2)}$_{\rm u}$. By Theorem~\ref{dp:thm:strong-um}(4), $j\circ f$ is analytic for an isometry $j$ into an ultrametric normed space, which we may complete. So $j\circ f$, and then $f$, is in $\cC$. The germ $\mathsf f$ is DP and not sDP (Proposition~\ref{dp:prop:strong-interpolates}(4)).
\end{proof}

Theorem~\ref{dp:thm:minimal} does not apply to the class of item (1), as it fails (C1), nor to the class of item (2), as the enlargement in the proof of that theorem is not ultrametric. Table~\ref{dp:s:tab:descents} shows how the classes fail, for Banach values. Its rows for AA, TA and CA use Sections~\ref{dp:sec:chain} and~\ref{dp:sec:cover}.

\begin{table}[htbp]
\small
\caption{Properties of the classes, for Banach values, in every characteristic unless a cell says otherwise. A ``no'' means that the property fails for some spaces allowed by the cell.}\label{dp:s:tab:descents}
\begin{tabular}{@{}p{0.07\linewidth}p{0.06\linewidth}p{0.3\linewidth}p{0.07\linewidth}p{0.34\linewidth}@{}}
Class & (C1) & (C2)$_{\rm u}$ & (C2)$_\times$ & (C3)\\
\hline
An & yes & char.\ $0$: yes; char.\ $p$: yes for $K$ discrete or $W$ spherically complete up to equivalence; no for $K$ complete, not spherically complete; else open & no & char.\ $0$: yes if $P$ or $W$ is ultrametric, else open; char.\ $p$: no; for ultrametric $W$: yes for $K$ discrete or $W$ spherically complete up to equivalence, no for $K$ complete, not spherically complete, else open\\
AA & yes & yes & yes & no in char.\ $p$; char.\ $0$ open\\
TA & yes & yes & yes & no in char.\ $p$; char.\ $0$ open\\
sDP & no & yes & yes & yes\\
CA & yes & yes & no & yes\\
DP & yes & yes & yes & yes\\
\end{tabular}
\end{table}

\emph{Sources of the entries.} An: Corollary~\ref{dp:cor:two-gaps}, Theorem~\ref{dp:thm:strong-um}, Proposition~\ref{dp:prop:other-classes}(1), Theorem~\ref{dp:thm:ultrametric-minimal}(3), Theorem~\ref{dp:thm:forcing}(2), and the germ $\mathsf f$ of Example~\ref{dp:ex:l1} for (C2)$_\times$. AA: Proposition~\ref{dp:prop:other-classes}(2), Theorem~\ref{dp:thm:singles-out}(4),(5). TA: Theorem~\ref{dp:thm:ta-composition}(4), Theorem~\ref{dp:thm:target-descent}(2), Theorem~\ref{dp:thm:singles-out}(2),(4),(5). sDP: Theorem~\ref{dp:thm:strong-descent}. CA: Lemma~\ref{dp:lem:cover-flat}(3) for (C1) and (C3). For (C2)$_{\rm u}$, Lemma~\ref{dp:lem:cover-flat}(5), by which CA $=$ sDP for ultrametric $W$, and the row of sDP. For (C2)$_\times$, the germ $\mathsf f$: by Example~\ref{dp:ex:l1}(2) and Lemma~\ref{dp:lem:cover-flat}(4), $j\circ\mathsf f$ is CA for an isometry $j$, while $\mathsf f$ is not sDP (Proposition~\ref{dp:prop:strong-interpolates}(4)), so not CA by Lemma~\ref{dp:lem:cover-flat}(5). DP: Theorem~\ref{dp:thm:minimal}. For the axiom (S) of Definition~\ref{dp:def:s-axiom}: DP has it; the analytic, sDP and CA germs fail it on $\ell^1(\N,K)$ (Proposition~\ref{dp:prop:s-axiom}); AA and TA fail it in characteristic $p$ (Proposition~\ref{dp:prop:s-fails}); in characteristic $0$ they have it if and only if they equal DP (Proposition~\ref{dp:prop:s-char0}).

\subsection{The norm itself}\label{dp:s:ssec-norm}

On ultrametric spaces $\lvert\!\lvert\!\lvert \cdot\rvert\!\rvert\!\rvert$ is the law norm over a densely valued extension (item (5) below). Base change is that of Section~\ref{dp:sec:base-change}: $P_{K'}$, $W_{K'}$, $a_{K'}$ with $\pi$, and $P^{\max}_{K'}$, $W^{\max}_{K'}$ with $\pi_{\max}$.

\begin{proposition}[the strong norm is intrinsic]\label{dp:prop:strong-intrinsic}
Let $a\colon P\to W$ be a law of degree $n$.
\begin{enumerate}
\item The coefficients $a^{[\alpha]}$ are determined by $a$, as $K$ is infinite. So $\lvert\!\lvert\!\lvert a\rvert\!\rvert\!\rvert$ is determined by $a$, with no enlargement to choose.
\item An isometry on $W$ does not change $\lvert\!\lvert\!\lvert a\rvert\!\rvert\!\rvert$. A change of the norm of $P$ by a constant $\mathsf c$ changes it by at most the factor $\mathsf c^n$, and a change of the norm of $W$ by at most the factor $\mathsf c$. So sDP germs are invariant under linear homeomorphisms. Restriction to a subspace of $P$ does not increase it. For a quotient, see Lemma~\ref{dp:lem:strong-formal}(2).
\item If $W$ is ultrametric, then $\lvert\!\lvert\!\lvert a\rvert\!\rvert\!\rvert=\lvert\!\lvert\!\lvert \tilde a\rvert\!\rvert\!\rvert$ on $P^{\mathsf u}$.
\item \textup{(Base change, all norms.)} For every valued extension $K'/K$ we have $\lvert\!\lvert\!\lvert a_{K'}\rvert\!\rvert\!\rvert=\lvert\!\lvert\!\lvert a\rvert\!\rvert\!\rvert$. For ultrametric $P,W$ the same holds with $\pi_{\max}$. So if $f$ is sDP at $x_0$ with radius $R$, bound $M$ and laws $a_n$, then $\sum_n(a_n)_{K'}$, with values in the completion of $W_{K'}$, is sDP at $0$ with the same radius and bound; and the bounds $\lvert\!\lvert\!\lvert (a_n)_{K'}\rvert\!\rvert\!\rvert\le MR^{-n}$ descend to the $a_n$.
\item Let $P,W$ be ultrametric. For a densely valued extension $K'$, $\norm{a_{K'}}_{\pi_{\max}}=\lvert\!\lvert\!\lvert a\rvert\!\rvert\!\rvert$. If $K$ is complete, this number is $\operatorname{amb}_{\max}(a)$.
\end{enumerate}
\end{proposition}

\begin{proof}
\emph{Step 1 (items (1)--(3)).} Item (1) is the uniqueness of coefficients of a polynomial over an infinite field. Item (2) is Lemma~\ref{dp:lem:strong-formal}(1),(2). Item (3) is Proposition~\ref{dp:prop:hull}(2).

\emph{Step 2 (item (4)).} Let $x_c=\sum_ly_{c,l}\otimes\lambda_{c,l}$ in $P\otimes_KK'$. As in~\eqref{dp:s:eq:slots}, $a_{K'}^{[\alpha]}(x_1,\dots,x_r)=\sum_\varepsilon a^{[\varepsilon]}((y_{c,l})_{c,l})\otimes\lambda^\varepsilon$. By~\eqref{dp:laws:strong}, its projective seminorm is at most
\[
 \lvert\!\lvert\!\lvert a\rvert\!\rvert\!\rvert\sum_\varepsilon\prod_{c,l}\bigl(\abs{\lambda_{c,l}}\norm{y_{c,l}}\bigr)^{\varepsilon_{c,l}}\le\lvert\!\lvert\!\lvert a\rvert\!\rvert\!\rvert\prod_c\Bigl(\sum_l\abs{\lambda_{c,l}}\norm{y_{c,l}}\Bigr)^{\alpha_c},
\]
as multinomial coefficients are at least $1$. Take infima over the representations of the $x_c$. With $\pi_{\max}$, use maxima. For the reverse inequality, $P\to P_{K'}$ and $W\to W_{K'}$ are isometric (Theorem~\ref{dp:thm:base-change-isometric}), and $a_{K'}^{[\alpha]}(y_1\otimes1,\dots,y_r\otimes1)=a^{[\alpha]}(y_1,\dots,y_r)\otimes1$.

\emph{Step 3 (item (5)).} Lemma~\ref{dp:lem:strong-coefficients}(2) over $K'$, between $P^{\max}_{K'}$ and $W^{\max}_{K'}$, gives $\norm{a_{K'}}_{\pi_{\max}}=\lvert\!\lvert\!\lvert a_{K'}\rvert\!\rvert\!\rvert$, and item (4) gives $\lvert\!\lvert\!\lvert a\rvert\!\rvert\!\rvert$. For $K$ complete, Theorem~\ref{dp:thm:base-change-amb}(3) gives $\norm{a_{K'}}_{\pi_{\max}}=\operatorname{amb}_{\max}(a)$.
\end{proof}

The law norm can lose radius under base change (Example~\ref{dp:ex:base-radius}); by item (4) the strong norm does not.

\subsection{Radii}\label{dp:s:ssec-radii}

Let $W$ be complete, $U\subseteq P$ open, and let $f\colon U\to W$ be sDP at $x_0$ with laws $a_n$. Its \emph{strong radius} $R_{\rm s}(f,x_0)$ is the supremum of the $R>0$ with $B(x_0,R)\subseteq U$ such that $f(x_0+h)=\sum_na_n(h)$ for $\norm h<R$ and $\sup_n\lvert\!\lvert\!\lvert a_n\rvert\!\rvert\!\rvert R^n<\infty$.

\begin{theorem}[recentering in the strong norm]\label{dp:thm:strong-recenter}
Let $W$ be complete, $U\subseteq P$ open, and let $f\colon U\to W$ be sDP at $x_0$ with laws $a_n$. Let $x_1=x_0+x$ with $\varrho=\norm x<R_{\rm s}(f,x_0)$. Let $R>\varrho$ be such that $B(x_0,R)\subseteq U$, $f(x_0+h)=\sum_na_n(h)$ for $\norm h<R$ and $M_R=\sup_n\lvert\!\lvert\!\lvert a_n\rvert\!\rvert\!\rvert R^n<\infty$; every $R$ with $\varrho<R<R_{\rm s}(f,x_0)$ is of this kind. Put $a'_i=\sum_{n\ge i}a_n^{[(n-i,i)]}(x,\cdot\,)$.
\begin{enumerate}
\item The series converges in the strong norm, $a'_i$ is a law of degree $i$, and
\[
 \lvert\!\lvert\!\lvert a'_i\rvert\!\rvert\!\rvert\le\frac{M_RR^{-i}}{1-\varrho/R},\quad\text{and}\quad\lvert\!\lvert\!\lvert a'_i\rvert\!\rvert\!\rvert\le M_RR^{-i}\text{ if $W$ is ultrametric.}
\]
\item $f(x_1+y)=\sum_ia'_i(y)$ for $\norm y<R-\varrho$, and for $\norm y<R$ if $P$ is ultrametric.
\item So $f$ is sDP at $x_1$ and $R_{\rm s}(f,x_1)\ge R_{\rm s}(f,x_0)-\varrho$. If $P$ is ultrametric, then $R_{\rm s}(f,x_1)=R_{\rm s}(f,x_0)$.
\end{enumerate}
Without completeness of $W$, the same holds for $f$ as a map into $\widehat W$.
\end{theorem}

\begin{proof}
\emph{Step 1 (item (1)).} By~\eqref{dp:eq:coefficient-identity} and Lemma~\ref{dp:lem:strong-formal}(4), the term of degree $n$ has strong norm at most $\lvert\!\lvert\!\lvert a_n\rvert\!\rvert\!\rvert\varrho^{n-i}\le M_RR^{-i}(\varrho/R)^{n-i}$. Sum the geometric series, or take the maximum if $W$ is ultrametric. Convergence, and the law, are Lemma~\ref{dp:lem:strong-formal}(4).

\emph{Step 2 (item (2)).} The case $\norm y<R-\varrho$ is Lemma~\ref{dp:lem:strong-formal}(4). Let $P$ be ultrametric and $\norm y<R$. Then $\norm{x+y}<R$, and $\norm{a_n^{[(n-i,i)]}(x,y)}\le\lvert\!\lvert\!\lvert a_n\rvert\!\rvert\!\rvert\max(\varrho,\norm y)^n$. So the double series is bounded by $\sum_n(n+1)M_R(\max(\varrho,\norm y)/R)^n<\infty$. Summing over $i$ first gives $\sum_na_n(x+y)=f(x_1+y)$ by~\eqref{dp:eq:coefficient-identity}. Summing over $n$ first gives $\sum_ia'_i(y)$.

\emph{Step 3 (item (3)).} By item (1), $\sup_i\lvert\!\lvert\!\lvert a'_i\rvert\!\rvert\!\rvert R'^i<\infty$ for $R'<R$. With item (2), $R_{\rm s}(f,x_1)\ge R-\varrho$, and $\ge R$ if $P$ is ultrametric. Let $R\uparrow R_{\rm s}(f,x_0)$. If $P$ is ultrametric, then $x_0\in B(x_1,R_{\rm s}(f,x_1))$, and the same argument with $x_0,x_1$ exchanged gives equality.
\end{proof}

Equality in item (3) fails without an ultrametric $P$. Let $K=\mathbb Q_2$, $P=K^2$ with $\norm{(y_1,y_2)}=\abs{y_1}_2+\abs{y_2}_2$, and $f=0$ on $B(0,1)$, $f=1$ elsewhere. At $0$ all laws vanish and $R_{\rm s}(f,0)=1$. Let $x_1=(2,0)$, so $\norm{x_1}=1/2$. Then $f=0$ on $B(x_1,1/2)$, so all laws at $x_1$ vanish, while $f(x_1+(0,2))=1$ with $\norm{(0,2)}=1/2$. So $R_{\rm s}(f,x_1)=1/2$.

\begin{remark}[radii of the ambient and the divided-power class]\label{dp:rem:radius-contrast}
Let $W$ be Banach and $f$ be DP at $x_0$, with laws $a_n$. Put $\rho_{\rm AA}(f,x_0)=(\limsup_n\operatorname{amb}(a_n)^{1/n})^{-1}$. For ultrametric $P,W$, with $\operatorname{amb}_{\max}$ of Section~\ref{dp:sec:base-change}, put $\rho^{\max}_{\rm AA}(f,x_0)=(\limsup_n\operatorname{amb}_{\max}(a_n)^{1/n})^{-1}$. As $\norm a\le\operatorname{amb}(a)\le\operatorname{amb}_{\max}(a)$, we have
\[
 \rho^{\max}_{\rm AA}(f,x_0)\le\rho_{\rm AA}(f,x_0)\le\rho_{\rm DP}(f,x_0).
\]
So the DP radius at a point is never smaller than the ambient radii there. As in Definition~\ref{dp:def:radius}, in statements about another point $f$ is the sum of its law series on the ball in question. Below, $x_1=x_0+x$ and $\varrho=\norm x$.

\emph{The formula $R-\varrho$.} For $\varrho<\rho_{\rm AA}(f,x_0)$ we have $\rho_{\rm AA}(f,x_1)\ge\rho_{\rm AA}(f,x_0)-\varrho$, for all $K$ and norms \cite[Propositions~14.9(1), 14.12(2)]{I}. For $\varrho<\rho_{\rm DP}(f,x_0)$ the same holds for $\rho_{\rm DP}$ if $K$ is densely valued (Theorem~\ref{dp:thm:whole-ball}(2)). It fails over complete discretely valued $K$ with $\abs\vartheta<1/16$, for any residue field (Example~\ref{dp:ex:radius-drop}), so for local fields at the module normalization $\abs\vartheta=1/\mathsf q$ with $\mathsf q\ge17$. This inequality is not invariant under powers of the absolute value. The next one is.

\emph{Constancy on the ball.} Let $P,W$ be ultrametric. Then $\rho^{\max}_{\rm AA}(f,x_1)\ge\rho^{\max}_{\rm AA}(f,x_0)$ for $\varrho<\rho^{\max}_{\rm AA}(f,x_0)$. Indeed, let $\varrho<R<\rho^{\max}_{\rm AA}(f,x_0)$ and $M=\sup_n\operatorname{amb}_{\max}(a_n)R^n$. This is finite, as each $\operatorname{amb}_{\max}(a_n)$ is finite by \cite[Theorem~14.39(1), (14.9)]{I}. If $M=0$, then all $a_n$ vanish and the claim is clear. If $M>0$, Lemma~\ref{dp:bc:envelope} gives an isometry $j$ of $W$ into an ultrametric space $Z$, which we complete, and $n$\nobreakdash-linear $\beta_n$ with $\norm{\beta_n}\le MR^{-n}$ and diagonals $j\circ a_n$. Re-expand as in \cite[Proposition~14.12(2)]{I}, with the maximum in place of the sum, since the terms of the expansion of $\beta_n(x+h,\dots,x+h)$ have norm at most $MR^{-n}\max(\varrho,\norm h)^n$: the part $\beta'_i$ of degree $i$ has $\norm{\beta'_i}\le MR^{-i}$, and it converges for $\norm h<R$. Here $f$ is DP at $x_1$, with laws $a'_i$ (Theorem~\ref{dp:thm:whole-ball}(1)). By Lemma~\ref{dp:lem:coefficients-unique}(2) the diagonal of $\beta'_i$ is $j\circ a'_i$, so $\operatorname{amb}_{\max}(a'_i)\le MR^{-i}$ by~\eqref{dp:bc:amb-rep}. Then $\varrho<\rho^{\max}_{\rm AA}(f,x_1)$, and exchanging $x_0$ and $x_1$ gives equality. If $\abs{K^\times}$ is dense, or $\kappa$ is infinite and $\mathfrak s(P)=1$, then $\norm\cdot=\lvert\!\lvert\!\lvert \cdot\rvert\!\rvert\!\rvert$ on laws (Lemma~\ref{dp:lem:strong-coefficients}(2)), and Theorem~\ref{dp:thm:strong-recenter} gives the same for the DP radius. Over complete discretely valued $K$ the DP radius can drop inside its ball: with any residue field and $\mathfrak s(P)<1$ (Example~\ref{dp:ex:radius-drop}), and with finite residue field, $\mathfrak s(P)=1$ and radius at $x_0$ in $\abs{K^\times}$ (Example~\ref{dp:ex:radius-drop-b}). In general, on ultrametric spaces, a map DP at $x_0$ with radius $R_0$ is DP with radius $\theta_K\mathfrak s(P)R_0$ at every point of $B(x_0,\theta_K\mathfrak s(P)R_0)$, by Proposition~\ref{dp:prop:strong-interpolates}(3) and by Theorem~\ref{dp:thm:strong-recenter} with $R=\theta_K\mathfrak s(P)R_0$.
\end{remark}
%% ---- end sec-14-strong.tex
%% ---- begin sec-15-cover.tex
\section{Germs analytic after a cover}\label{dp:sec:cover}

This section studies a class of germs beside the DP germs: those that become analytic after the canonical cover $q_P$ of Definition~\ref{dp:def:cover}. We call them CA. The class has (C1) and (C3). For $K$ complete and Banach spaces it has composition, inverse functions and bundles. On ultrametric sources it is the class of DP germs. It does not contain the germ $\mathsf f$ of Example~\ref{dp:ex:l1}. So its theorems show what target descent into arbitrary Banach spaces adds to that list (Theorem~\ref{dp:thm:descents-independent}).

In this section $\mathcal C(N,l)$ is the set of $\vec n=(n_1,\dots,n_l)$ of positive integers with sum $N$, $\hat q_P$ is as in Definition~\ref{dp:def:cover}, and $\mathsf c=\abs\varpi^{-1}$ is the preimage constant of $q_P$ and $\hat q_P$. The conditions (C1), (C2), (C3) are those of Section~\ref{dp:ssec:minimal}.

\subsection{The class}\label{dp:ssec:ca-class}

\begin{definition}\label{dp:def:flat}
Let $P,W$ be normed. For $w\in W$ and an integer $\mathsf N\ge1$ put
\[
 \norm w_{(\mathsf N)}=\inf\Bigl\{\max_{i\le\mathsf N}\norm{w_i}:\ w=\sum_{i\le\mathsf N}w_i\Bigr\}.
\]
For $a\in\cP_n(P;W)$, $n\ge1$, put
\[
 \lvert\!\lvert\!\lvert a\rvert\!\rvert\!\rvert_\flat=\sup\frac{\norm{a^{[\alpha]}(h_1,\dots,h_l)}_{(n!/\alpha!)}}{\norm{h_1}^{\alpha_1}\cdots\norm{h_l}^{\alpha_l}},
\]
the supremum over $l\ge1$, nonzero $h_1,\dots,h_l\in P$ and $\abs\alpha=n$. For $n=0$ read $\lvert\!\lvert\!\lvert a\rvert\!\rvert\!\rvert_\flat=\norm{a(0)}$.
\end{definition}

Always $\norm w_{(\mathsf N)}\le\norm w$, and equality holds if $W$ is ultrametric. Also $\norm{\mu w}_{(\mathsf N)}=\abs\mu\norm w_{(\mathsf N)}$ for $\mu\in K$. For a positive integer $\mathsf m$, repeating each piece of a decomposition $\mathsf m$ times gives
\begin{equation}\label{dp:eq:ca-repeat}
 \norm{\mathsf m w}_{(\mathsf m\mathsf N)}\le\norm w_{(\mathsf N)}.
\end{equation}

\begin{definition}\label{dp:def:ca}
Let $U\subseteq P$ be open, $f\colon U\to W$ and $x_0\in U$. The map $f$ is \emph{analytic after a cover}, or \emph{CA}, at $x_0$ if it is DP at $x_0$, with coefficient laws $a_n$, and $\sup_n\operatorname{lift}(a_n\circ q_P)R^n<\infty$ for some $R>0$. It is CA on $U$ if it is CA at every point of $U$.
\end{definition}

Shrinking $R$, we can take $R$ to be a radius of $f$ at $x_0$. Recall from Section~\ref{dp:ssec:singles} the class $\mathfrak S$: the germs $f$ at $x_0$ such that $f\circ Q$ is analytic at a preimage of $x_0$, for some bounded linear surjection $Q$ with a preimage constant. Here we allow any normed target in it.

\begin{lemma}\label{dp:lem:cover-flat}
Let $P,W$ be normed and let $a\in\cP_n(P;W)$.
\begin{enumerate}
\item We have
\[
 \abs\varpi^n\lvert\!\lvert\!\lvert a\rvert\!\rvert\!\rvert_\flat\le\operatorname{lift}(a\circ q_P)\le\lvert\!\lvert\!\lvert a\rvert\!\rvert\!\rvert_\flat\le\min\bigl(\lvert\!\lvert\!\lvert a\rvert\!\rvert\!\rvert,\operatorname{lift}(a)\bigr),\qquad\norm a\le\lvert\!\lvert\!\lvert a\rvert\!\rvert\!\rvert_\flat.
\]
\item If $Q\colon P'\to P$ is a bounded linear surjection with preimage constant $\mathsf c'$, then $\operatorname{lift}(a\circ q_P)\le\mathsf c'^n\operatorname{lift}(a\circ Q)$.
\item For $f\colon U\to W$ and $x_0\in U$ the following are equivalent: $f$ is CA at $x_0$; $f\circ q_P$ is analytic at every preimage of $x_0$; $f\circ q_P$ is analytic at some preimage of $x_0$; the germ of $f$ at $x_0$ is in $\mathfrak S$. So, among germs with Banach values, CA is the smallest class with \textup{(C1)} and \textup{(C3)}.
\item Strongly DP germs are CA, and CA germs are DP. Analytic germs are CA. Every bounded law is CA at every point.
\item If $W$ is ultrametric, then CA germs are the strongly DP germs. If $P$ is ultrametric, then for every $W$ the CA germs, the strongly DP germs and the DP germs coincide.
\item Let $P=\ell^1_{00}(J)$ and $W$ be arbitrary, or $P=\ell^1(J)$ and $W$ Banach. Then $\operatorname{lift}(a\circ q_P)=\operatorname{lift}(a)$. So on these spaces a germ is CA if and only if it is analytic, with the same radius.
\end{enumerate}
\end{lemma}

\begin{proof}
\emph{Step 1 (a formula for the lifting norm).} Let $J$ be a set and $b\colon\ell^1_{00}(J)\to W$ a law of degree $n\ge1$. A bounded $n$-linear map $\beta$ on $\ell^1_{00}(J)$ is the same as an arbitrary bounded function $\mathsf w\mapsto v_{\mathsf w}=\beta(e_{\mathsf w_1},\dots,e_{\mathsf w_n})$ on words $\mathsf w\in J^n$, and $\norm\beta=\sup_{\mathsf w}\norm{v_{\mathsf w}}$. The \emph{content} of a word records how often each letter occurs. For finitely supported $h$,
\[
 \beta(h,\dots,h)=\sum_\alpha h^\alpha\sum_{\mathsf w\text{ of content }\alpha}v_{\mathsf w},\qquad b(h)=\sum_\alpha h^\alpha b^{[\alpha]}(e).
\]
As $K$ is infinite, $\beta$ has diagonal $b$ if and only if, for each content $\alpha$, the $n!/\alpha!$ vectors $v_{\mathsf w}$ of content $\alpha$ sum to $b^{[\alpha]}(e)$. Contents can be treated one at a time, so
\begin{equation}\label{dp:eq:ca-word}
 \operatorname{lift}(b)=\sup\bigl\{\norm{b^{[\alpha]}(e_{\theta_1},\dots,e_{\theta_l})}_{(n!/\alpha!)}:\ l\ge1,\ \theta_1,\dots,\theta_l\in J\text{ distinct},\ \alpha_c\ge1,\ \abs\alpha=n\bigr\}.
\end{equation}

\emph{Step 2 (item (1)).} The case $n=0$ is immediate. Let $n\ge1$. For $b=a\circ q_P$ and distinct $\xi_1,\dots,\xi_l\in I_P$ we have $b^{[\alpha]}(e_{\xi_1},\dots,e_{\xi_l})=a^{[\alpha]}(\xi_1,\dots,\xi_l)$. As $\norm{\xi_c}\le1$, formula~\eqref{dp:eq:ca-word} gives $\operatorname{lift}(a\circ q_P)\le\lvert\!\lvert\!\lvert a\rvert\!\rvert\!\rvert_\flat$.

For the left inequality, let $h_1,\dots,h_l\in P$ be nonzero and $\abs\alpha=n$. Arguments with $\alpha_c=0$ can be deleted, by Lemma~\ref{dp:lem:elementary}, without changing $n!/\alpha!$. Choose $\lambda_c\in K$ with $\abs\varpi<\abs{\lambda_c}\norm{h_c}\le1$, and put $\xi_c=\lambda_ch_c\in I_P$. Then $a^{[\alpha]}(\xi_1,\dots,\xi_l)=\lambda^\alpha a^{[\alpha]}(h_1,\dots,h_l)$ and $\abs{\lambda_c}^{-1}<\abs\varpi^{-1}\norm{h_c}$. If the $\xi_c$ are distinct, \eqref{dp:eq:ca-word} bounds $\norm{a^{[\alpha]}(\xi_1,\dots,\xi_l)}_{(n!/\alpha!)}$ by $\operatorname{lift}(a\circ q_P)$. Suppose $\xi_1=\xi_2$. Comparing coefficients in $a((s_1+s_2)\xi_1+\sum_{c\ge3}s_c\xi_c)$ gives $a^{[\alpha]}(\xi_1,\xi_1,\xi_3,\dots)=\binom{\alpha_1+\alpha_2}{\alpha_1}a^{[(\alpha_1+\alpha_2,\alpha_3,\dots)]}(\xi_1,\xi_3,\dots)$. Repeating, $a^{[\alpha]}(\xi_1,\dots,\xi_l)=\mathsf m\,a^{[\alpha']}(\xi')$ with $\xi'$ distinct and $\mathsf m=\alpha'!/\alpha!$ a positive integer. Since $n!/\alpha!=\mathsf m\,n!/\alpha'!$, \eqref{dp:eq:ca-repeat} gives the same bound. So $\norm{a^{[\alpha]}(h_1,\dots,h_l)}_{(n!/\alpha!)}\le\abs\varpi^{-n}\prod_c\norm{h_c}^{\alpha_c}\operatorname{lift}(a\circ q_P)$.

The inequality $\lvert\!\lvert\!\lvert a\rvert\!\rvert\!\rvert_\flat\le\lvert\!\lvert\!\lvert a\rvert\!\rvert\!\rvert$ follows from $\norm\cdot_{(\mathsf N)}\le\norm\cdot$. If $\beta$ is a lift of $a$, then $a^{[\alpha]}(h_1,\dots,h_l)$ is the sum of $\beta(h_{c_1},\dots,h_{c_n})$ over the $n!/\alpha!$ tuples in which each $c$ occurs $\alpha_c$ times, and each term has norm at most $\norm\beta\prod_c\norm{h_c}^{\alpha_c}$. So $\lvert\!\lvert\!\lvert a\rvert\!\rvert\!\rvert_\flat\le\operatorname{lift}(a)$. Taking $l=1$ gives $\norm a\le\lvert\!\lvert\!\lvert a\rvert\!\rvert\!\rvert_\flat$.

\emph{Step 3 (item (2)).} Choose $\tilde\xi\in P'$ with $Q\tilde\xi=\xi$ and $\norm{\tilde\xi}\le\mathsf c'\norm\xi\le\mathsf c'$ for each $\xi\in I_P$. The linear map $\mathsf t\colon\ell^1_{00}(I_P)\to P'$, $e_\xi\mapsto\tilde\xi$, has $\norm{\mathsf t}\le\mathsf c'$ and $Q\mathsf t=q_P$. If $\beta$ is a lift of $a\circ Q$, then $\beta\circ\mathsf t^{\times n}$ is a lift of $a\circ q_P$ of norm at most $\mathsf c'^n\norm\beta$.

\emph{Step 4 (item (3)).} Let $f$ be CA at $x_0$, with radius $R$ and $\operatorname{lift}(a_n\circ q_P)\le MR^{-n}$, and let $q_P(\tilde x_0)=x_0$. Choose lifts $\beta_n$ of $a_n\circ q_P$ with $\norm{\beta_n}\le(M+1)R^{-n}$. As $q_P$ is contractive, $f\circ q_P(\tilde x_0+\tilde h)=\sum_n\beta_n(\tilde h,\dots,\tilde h)$ for $\norm{\tilde h}<R$. So $f\circ q_P$ is analytic at $\tilde x_0$. The next two implications are trivial, with $Q=q_P$. Now let $f\circ Q$ be analytic at a preimage of $x_0$, with $Q$ of preimage constant $\mathsf c'$, radius $R$, bound $M$ and multilinear coefficients $\beta_n$. By Theorem~\ref{dp:thm:calculus}(3) and Proposition~\ref{dp:prop:source-descent}(1), $f$ is DP at $x_0$ with laws $a_n$ such that $a_n\circ Q$ are the laws of $f\circ Q$. By Lemma~\ref{dp:lem:coefficients-unique}(2), $\beta_n$ is a lift of $a_n\circ Q$. By item~(2), $\operatorname{lift}(a_n\circ q_P)\le\mathsf c'^nMR^{-n}$. So $f$ is CA at $x_0$. The last sentence is Theorem~\ref{dp:thm:singles-out}(3).

\emph{Step 5 (items (4) and (5)).} If $f$ is strongly DP with radius $R$ and bound $M$, Lemma~\ref{dp:lem:lift-cover} gives $\operatorname{lift}(a_n\circ q_P)\le MR^{-n}$. Every law is strongly DP at every point, by Lemma~\ref{dp:lem:strong-formal}(3). Analytic germs are in $\mathfrak S$, with $Q=\id$. If $W$ is ultrametric, $\lvert\!\lvert\!\lvert\cdot\rvert\!\rvert\!\rvert_\flat=\lvert\!\lvert\!\lvert\cdot\rvert\!\rvert\!\rvert$, so item~(1) gives $\lvert\!\lvert\!\lvert a_n\rvert\!\rvert\!\rvert\le\abs\varpi^{-n}\operatorname{lift}(a_n\circ q_P)$, and a CA germ with radius $R$ and $\operatorname{lift}(a_n\circ q_P)\le MR^{-n}$ is strongly DP with radius $\abs\varpi R$ and bound $M$. If $P$ is ultrametric, DP germs are strongly DP by Lemma~\ref{dp:lem:strong-formal}(3).

\emph{Step 6 (item (6)).} By item~(1), $\operatorname{lift}(a\circ q_P)\le\operatorname{lift}(a)$. For the converse, unit vectors have norm $1$, so they lie in $I_P$. The map $\mathsf t_0\colon\ell^1_{00}(J)\to\ell^1_{00}(I_P)$, $e_\theta\mapsto e_{e_\theta}$, is a linear isometry, and $q_P\mathsf t_0$ is the inclusion of $\ell^1_{00}(J)$ in $P$. If $\beta$ is a lift of $a\circ q_P$, then $\beta\circ\mathsf t_0^{\times n}$ has norm at most $\norm\beta$ and diagonal $a$ on $\ell^1_{00}(J)$. For $P=\ell^1_{00}(J)$ this is a lift of $a$. For $P=\ell^1(J)$ and $W$ Banach it extends to $P^n$ by density, with the same norm, and its diagonal is $a$, as $a$ is continuous (Theorem~\ref{dp:thm:calculus}(1)). For the last sentence, let $f$ be CA with radius $R$ and $\operatorname{lift}(a_n\circ q_P)\le MR^{-n}$. Then $\operatorname{lift}(a_n)\le MR^{-n}$, so there are lifts $\beta_n$ of $a_n$ with $\norm{\beta_n}\le(M+1)R^{-n}$, and $f(x_0+h)=\sum_n\beta_n(h,\dots,h)$ for $\norm h<R$. Conversely, let $f$ be analytic with radius $R$, bound $M$ and multilinear coefficients $\beta_n$. By Theorem~\ref{dp:thm:calculus}(3) and Lemma~\ref{dp:lem:coefficients-unique}(2), $f$ is DP with radius $R$ and $\beta_n$ is a lift of $a_n$, so $\operatorname{lift}(a_n\circ q_P)\le MR^{-n}$.
\end{proof}

For $P=\ell^1(J)$ with $J$ infinite and $W$ incomplete, we do not decide whether CA germs are analytic. The proof above gives lifts with values in $\widehat W$ only (Question~\ref{dp:q:cover}).

\begin{proposition}\label{dp:prop:ca-open}
Let $W$ be Banach. Let $f$ be CA at $x_0$, with radius $R$ and $\operatorname{lift}(a_n\circ q_P)\le MR^{-n}$. Then $f$ is CA at every point of $B(x_0,\abs\varpi R)$.
\end{proposition}

\begin{proof}
\emph{Step 1 (a preimage).} Choose lifts $\beta_n$ of $a_n\circ q_P$ with $\norm{\beta_n}\le(M+1)R^{-n}$, and $\tilde x_0$ with $q_P(\tilde x_0)=x_0$. Let $\norm{x-x_0}<\abs\varpi R$, and choose $\tilde h_1\in\ell^1_{00}(I_P)$ with $q_P\tilde h_1=x-x_0$ and $\norm{\tilde h_1}\le\abs\varpi^{-1}\norm{x-x_0}<R$.

\emph{Step 2 (grouping slots).} Let $\norm{\tilde h}<R-\norm{\tilde h_1}$. Then $f\circ q_P(\tilde x_0+\tilde h_1+\tilde h)=\sum_n\beta_n(\tilde h_1+\tilde h,\dots,\tilde h_1+\tilde h)$. Expand each term by multilinearity. For $d\ge0$ and $n\ge d$, let $\mathsf A_{n,d}$ be the $d$-linear map obtained by summing over the subsets $\sigma\subseteq\{1,\dots,n\}$ of size $d$, with its arguments placed in the slots $\sigma$ and $\tilde h_1$ in the other slots of $\beta_n$. Then $\norm{\mathsf A_{n,d}}\le\binom nd(M+1)R^{-n}\norm{\tilde h_1}^{n-d}$. With $\mathsf r=\norm{\tilde h_1}/R<1$,
\[
 \sum_{n\ge d}\binom nd\mathsf r^{n-d}=(1-\mathsf r)^{-d-1}.
\]
As $W$ is complete, $\mathsf A_d=\sum_{n\ge d}\mathsf A_{n,d}$ converges among bounded $d$-linear maps, with $\norm{\mathsf A_d}\le(M+1)R^{-d}(1-\mathsf r)^{-d-1}=\frac{M+1}{1-\mathsf r}(R-\norm{\tilde h_1})^{-d}$.

\emph{Step 3 (conclusion).} The double series is absolutely summable for $\norm{\tilde h}<R-\norm{\tilde h_1}$, since $\sum_n(M+1)((\norm{\tilde h_1}+\norm{\tilde h})/R)^n<\infty$. So $f\circ q_P(\tilde x_0+\tilde h_1+\tilde h)=\sum_d\mathsf A_d(\tilde h,\dots,\tilde h)$. Thus $f\circ q_P$ is analytic at the preimage $\tilde x_0+\tilde h_1$ of $x$. Apply Lemma~\ref{dp:lem:cover-flat}(3).
\end{proof}

\subsection{Composition and inverse functions}\label{dp:ssec:ca-calculus}

Pairs and bounded linear maps are formal. A pair of lifts is a lift of the pair, with the maximum of the norms, and $\operatorname{lift}(\Upsilon\circ b)\le\norm\Upsilon\operatorname{lift}(b)$ for a law $b$ and bounded linear $\Upsilon$. So pairs of CA germs into products with the maximum norm, and their bounded linear images, are CA.

\begin{theorem}[composition]\label{dp:thm:ca-composition}
Let $K$ be complete, of any characteristic, $P$ normed and $W,Y$ Banach. Let $f$ be CA at $x_0$, with laws $a_n$, and let $g$ map an open subset of $W$ to $Y$, CA at $f(x_0)$, with laws $b_m$. Then $g\circ f$ is CA at $x_0$. If $\operatorname{lift}(a_n\circ q_P)\le MR^{-n}$ for $n\ge1$ and $\operatorname{lift}(b_m\circ q_W)\le M'R'^{-m}$ for $m\ge0$, then for $N\ge1$,
\[
 \operatorname{lift}\bigl(D^{(N)}(g\circ f)(x_0)\circ q_P\bigr)\le M'\frac{\mathsf cM}{R'}\Bigl(1+\frac{\mathsf cM}{R'}\Bigr)^{N-1}R^{-N}.
\]
\end{theorem}

\begin{proof}
Fix $M_1>M$ and $M_1'>M'$. Let $\rho\le R$ and $\rho'\le R'$ be radii of $f$ at $x_0$ and of $g$ at $f(x_0)$, and let $q_P(\tilde x_0)=x_0$.

\emph{Step 1 (lifting through $q_W$).} Choose lifts $\beta_n$ of $a_n\circ q_P$, $n\ge1$, with $\norm{\beta_n}\le M_1R^{-n}$. Define a map $\mathsf l_W\colon W\to\ell^1_{00}(I_W)$, not linear, with $q_W\circ\mathsf l_W=\id_W$ and $\norm{\mathsf l_W(w)}\le\mathsf c\norm w$, as in the paragraph after Definition~\ref{dp:def:cover}. A multilinear map on $\ell^1_{00}(I_P)$ can be prescribed freely on words. So there is an $n$-linear $\tilde\beta_n\colon\ell^1_{00}(I_P)^n\to\ell^1_{00}(I_W)$ with $\tilde\beta_n(e_{\mathsf w_1},\dots,e_{\mathsf w_n})=\mathsf l_W(\beta_n(e_{\mathsf w_1},\dots,e_{\mathsf w_n}))$ on words $\mathsf w$. Then $\norm{\tilde\beta_n}\le\mathsf cM_1R^{-n}$ and $q_W\circ\tilde\beta_n=\beta_n$.

\emph{Step 2 (the inner map).} As $K$ is complete, $\ell^1(I_W)$ is a Banach space. So for $\norm{\tilde h}<\rho$ the series $\mathsf I(\tilde h)=\sum_{n\ge1}\tilde\beta_n(\tilde h,\dots,\tilde h)$ converges there, with $\sum_{n\ge1}\norm{\tilde\beta_n(\tilde h,\dots,\tilde h)}\le\mathsf V(\tilde h)$, where $\mathsf V(\tilde h)=\sum_{n\ge1}\mathsf cM_1(\norm{\tilde h}/R)^n$. Applying $\hat q_W$ gives $\hat q_W\mathsf I(\tilde h)=f(x_0+q_P\tilde h)-f(x_0)$.

\emph{Step 3 (the outer map).} Choose lifts $\beta'_m$ of $b_m\circ q_W$ with $\norm{\beta'_m}\le M_1'R'^{-m}$. As $Y$ is complete, they extend to $\ell^1(I_W)^m$ with the same norm. The diagonal of the extension is $b_m\circ\hat q_W$, as both are continuous and agree on the dense subspace. So $\sum_m\beta'_m(\tilde w,\dots,\tilde w)=g(f(x_0)+\hat q_W\tilde w)$ for $\norm{\tilde w}<\rho'$.

\emph{Step 4 (substitution).} For $\norm{\tilde h}<\rho$ with $\mathsf V(\tilde h)<\rho'$, substitute $\tilde w=\mathsf I(\tilde h)$. The resulting family is absolutely summable, since $\sum_m\norm{\beta'_m}\mathsf V(\tilde h)^m<\infty$. Grouping by total degree, $g\circ f\circ q_P(\tilde x_0+\tilde h)$ equals $b_0$ plus $\sum_{N\ge1}\Xi_N(\tilde h,\dots,\tilde h)$, where
\[
 \Xi_N=\sum_{l=1}^N\sum_{\vec n\in\mathcal C(N,l)}\beta'_l\circ(\tilde\beta_{n_1}\times\dots\times\tilde\beta_{n_l})
\]
is $N$-linear, each $\tilde\beta_{n_i}$ taking its own block of $n_i$ arguments. As $\mathcal C(N,l)$ has $\binom{N-1}{l-1}$ elements,
\[
 \norm{\Xi_N}\le\sum_{l=1}^N\binom{N-1}{l-1}M_1'R'^{-l}(\mathsf cM_1)^lR^{-N}=M_1'\frac{\mathsf cM_1}{R'}\Bigl(1+\frac{\mathsf cM_1}{R'}\Bigr)^{N-1}R^{-N}.
\]

\emph{Step 5 (conclusion).} The germ $g\circ f$ is DP at $x_0$ (Theorem~\ref{dp:thm:calculus}(2)), and $g\circ f\circ q_P$ is DP at $\tilde x_0$ with laws $D^{(N)}(g\circ f)(x_0)\circ q_P$ (Proposition~\ref{dp:prop:source-descent}(2)). By Lemma~\ref{dp:lem:coefficients-unique}(2) these are the diagonals of the $\Xi_N$. Letting $M_1\to M$ and $M_1'\to M'$ gives the bound. It is geometric in $N$, so $g\circ f$ is CA at $x_0$.
\end{proof}

\begin{theorem}[inverse and implicit functions]\label{dp:thm:ca-inverse}
Let $K$ be complete.
\begin{enumerate}
\item Let $P,W$ be Banach, let $f$ be CA at $x_0$, and let $\mathsf T_0=Df(x_0)$ be bijective with bounded inverse. Then the local inverse germ of $f$, which is DP by Theorem~\ref{dp:thm:inverse}, is CA at $f(x_0)$.
\item Let $P,Y,Y'$ be Banach. In Corollary~\ref{dp:cor:implicit}, if $\Phi$ is CA at $(x_0,y_0)$, then $\psi$ is CA at $x_0$.
\end{enumerate}
\end{theorem}

\begin{proof}
\emph{Step 1 (normal form).} Let $f$ be CA at $x_0$ with radius $R$ and laws $a_n$, and put $\tilde f(z)=\mathsf T_0^{-1}(f(x_0+z)-f(x_0))$. By Theorem~\ref{dp:thm:differentiable}(1), $a_1=\mathsf T_0$. So $\tilde f(z)=z-\sum_{m\ge2}\mathsf R_m(z)$ for $\norm z<R$, with $\mathsf R_m=-\mathsf T_0^{-1}\circ a_m$. Shrinking $R$, choose $M>0$ and lifts $\beta_m$ of $\mathsf R_m\circ q_P$ with $\norm{\beta_m}\le MR^{-m}$, $m\ge2$.

\emph{Step 2 (lifting to $\ell^1(I_P)$).} As in Step~1 of the proof of Theorem~\ref{dp:thm:ca-composition}, there are $m$-linear $\tilde\beta_m\colon\ell^1_{00}(I_P)^m\to\ell^1_{00}(I_P)$ with $q_P\circ\tilde\beta_m=\beta_m$ and $\norm{\tilde\beta_m}\le\mathsf cMR^{-m}$. As $K$ and $P$ are complete, $\tilde\beta_m$ and $\beta_m$ extend to $\ell^1(I_P)^m$, with values in $\ell^1(I_P)$ and $P$. By continuity, $\hat q_P\circ\tilde\beta_m=\beta_m$, and the diagonal of $\beta_m$ is $\mathsf R_m\circ\hat q_P$. Put $\mathsf F(u)=u-\sum_{m\ge2}\tilde\beta_m(u,\dots,u)$ for $u\in\ell^1(I_P)$, $\norm u<R$. Then $\hat q_P\circ\mathsf F=\tilde f\circ\hat q_P$ on this ball.

\emph{Step 3 (the inverse on the cover).} Define $N$-linear maps on $\ell^1(I_P)$ with values in $\ell^1(I_P)$: $\mathsf G_1=\id$ and, for $N\ge2$,
\[
 \mathsf G_N=\sum_{l=2}^N\sum_{\vec n\in\mathcal C(N,l)}\tilde\beta_l\circ(\mathsf G_{n_1}\times\dots\times\mathsf G_{n_l}),
\]
each $\mathsf G_{n_i}$ taking its own block of arguments. Then $\norm{\mathsf G_N}\le\mathsf y_N$, where $\mathsf y_1=1$ and
\[
 \mathsf y_N=\sum_{l=2}^N\mathsf cMR^{-l}\sum_{\vec n\in\mathcal C(N,l)}\prod_{i=1}^l\mathsf y_{n_i}\qquad(N\ge2).
\]
This is the sequence of Step~3 of the proof of~\cite[Lemma~Q.14]{I}, with $\mathsf cM$ for $\mathfrak b_{\mathsf T}$ and $1/R$ for $\varkappa$. Step~4 of that proof holds for all positive values of these constants. So there is $R_0>0$ with $\mathsf Y(t)=\sum_N\mathsf y_Nt^N<\infty$ for $0\le t<R_0$ and $\mathsf Y(t)\to0$ as $t\to0$. Shrink $R_0$ so that $\mathsf Y(t)<R$ on $[0,R_0)$. For $\norm u<R_0$ the series $\mathsf G(u)=\sum_N\mathsf G_N(u,\dots,u)$ converges, with $\norm{\mathsf G(u)}\le\mathsf Y(\norm u)$. Substituting and regrouping an absolutely summable family gives $\mathsf F(\mathsf G(u))=u$.

\emph{Step 4 (descent to $P$).} By Steps~2 and~3, $\tilde f(\hat q_P\mathsf G(u))=\hat q_P\mathsf F(\mathsf G(u))=\hat q_Pu$ for $\norm u<R_0$. By Theorem~\ref{dp:thm:inverse}, $\tilde f$ restricts to a bijection between neighborhoods of $0$, with inverse $\tilde g$. As $\norm{\hat q_P\mathsf G(u)}\le\mathsf Y(\norm u)$, for small $u$ we get $\hat q_P\mathsf G(u)=\tilde g(\hat q_Pu)$. So $\tilde g\circ\hat q_P$ has the expansion $\sum_N\hat q_P\circ\mathsf G_N$, with $\norm{\hat q_P\circ\mathsf G_N}\le\mathsf y_N$, geometric in $N$. Restricting to $\ell^1_{00}(I_P)$, $\tilde g\circ q_P$ is analytic at $0$.

\emph{Step 5 (item (1)).} The local inverse germ of $f$ is $w\mapsto x_0+\tilde g(\mathsf T_0^{-1}(w-f(x_0)))$. The map $Q=\mathsf T_0\circ q_P$ is a bounded linear surjection onto $W$ with preimage constant $\mathsf c\norm{\mathsf T_0^{-1}}$. Choose $\tilde w_0$ with $Q(\tilde w_0)=f(x_0)$. The local inverse germ evaluated at $Q(\tilde w_0+u)$ is $x_0+\tilde g(q_Pu)$. So the germ is in $\mathfrak S$, and it is CA by Lemma~\ref{dp:lem:cover-flat}(3).

\emph{Step 6 (item (2)).} The map $\mathsf H(x,y)=(x,\Phi(x,y))$ is CA at $(x_0,y_0)$, as a pair of CA germs. Its derivative has a bounded inverse, by Step~1 of the proof of Corollary~\ref{dp:cor:implicit}. By that proof, $\psi(x)$ is the second coordinate of $\mathsf H^{-1}(x,0)$. So $\psi$ is CA at $x_0$, by item~(1), Theorem~\ref{dp:thm:ca-composition}, and the fact that bounded linear maps are analytic.
\end{proof}

\emph{CA on open sets.} Let $K$ be complete. By Proposition~\ref{dp:prop:ca-open} and Theorems~\ref{dp:thm:ca-composition}, \ref{dp:thm:ca-inverse}, the local inverse of a map between Banach spaces that is CA on an open set, with derivative at the base point bijective with bounded inverse, is again CA on an open set. Maps that are CA on open sets, with Banach models and fibers, have the four properties (i)--(iv) of Section~\ref{dp:sec:conv}. Indeed, composition is Theorem~\ref{dp:thm:ca-composition}; tuples are pairs; evaluation and composition of operators are bounded bilinear, hence analytic; constants are analytic; and the definition is local. So Definition~\ref{dp:def:manifold}(1),(2), with DP replaced by CA on open sets and with Banach models and fibers, gives manifolds and vector bundles. The proof of Corollary~\ref{dp:cor:bundles} applies. As there, the replacement is also made in the definitions that \cite[Theorem~2.1]{I} uses: bundle morphisms are those of Definition~\ref{dp:def:manifold}(2), and a functor is in the class if all its maps on products of hom spaces are jointly CA on open sets. By Step~3 of that proof, the maps of the alternating bifunctor on Banach spaces and of the functor $\widehat\Lambda^k_\pi$ on hom spaces are bounded laws, here with Banach values, so they are CA at every point (Lemma~\ref{dp:lem:cover-flat}(4)). So in every characteristic the alternating and $\widehat\Lambda^k_\pi$ bundles of such a bundle are bundles of the same kind.

\subsection{Derivatives}\label{dp:ssec:ca-derivatives}

Tangent bundles and jets need the derivatives of a germ. In this subsection all targets are Banach, so each $\cP_d(P;W)$ is Banach (Theorem~\ref{dp:thm:normed-roby}(3)). If $P$ or $W$ is ultrametric, derivatives of CA maps are CA (Remark~\ref{dp:rem:cah-um}). In general we do not decide whether $Df$ is CA when $f$ is (Question~\ref{dp:q:cover}). So we build it into a definition.

\begin{definition}\label{dp:def:cah}
Let $U\subseteq P$ be open and $W$ Banach. A map $f\colon U\to W$ is \emph{CA$^{\rm H}$} on $U$ if it is DP on $U$ and, for every $d\ge0$, the map $x\mapsto D^{(d)}f(x)\in\cP_d(P;W)$ is CA at every point of $U$.
\end{definition}

For composition we use the map of Lemma~\ref{dp:lem:chain-law}, with $r=d$: for $\alpha\in\N^d$ with $\deg\alpha=d\ge1$, $\mathcal F_\alpha(b,(\ell_n)_{1\le n\le d})$ is the law $h\mapsto b^{[\alpha]}(\ell_1(h),\dots,\ell_d(h))$ of degree $d$. By Lemma~\ref{dp:lem:chain-law}(3), it is a bounded law in $(b,(\ell_n)_n)$ on $\cP_{\abs\alpha}(W;Y)\times\prod_{1\le n\le d}\cP_n(P;W)$. By Proposition~\ref{dp:prop:chain-rule},
\begin{equation}\label{dp:eq:ca-chain}
 D^{(d)}(g\circ f)(x)=\sum_{\deg\alpha=d}\mathcal F_\alpha\bigl(D^{(\abs\alpha)}g(f(x)),(D^{(n)}f(x))_{1\le n\le d}\bigr)\qquad(d\ge1).
\end{equation}

\begin{theorem}\label{dp:thm:cah}
Let all targets be Banach, and in items \textup{(1)--(4)} let $K$ be complete, of any characteristic.
\begin{enumerate}
\item CA$^{\rm H}$ contains the maps analytic on an open set and all bounded laws. It is closed under composition, pairing and bounded linear postcomposition. If $f$ is CA$^{\rm H}$ on $U$, then so is $D^{(d)}f\colon U\to\cP_d(P;W)$, for every $d$.
\item Let $P,W$ be Banach, and let $f$ be CA$^{\rm H}$ on $U$ with $Df(x_0)$ bijective with bounded inverse. Then the local inverse is CA$^{\rm H}$ on a neighborhood of $f(x_0)$. The map $\psi$ of Corollary~\ref{dp:cor:implicit} is CA$^{\rm H}$ near $x_0$ if $\Phi$ is CA$^{\rm H}$.
\item Theorem~2.1 of~\cite{I}, displayed in Section~\ref{dp:sec:conv}, holds with $C^m$ replaced by CA$^{\rm H}$ on open sets, for Banach models and fibers, with the replacement also made in the definitions that the theorem uses, as in Corollary~\ref{dp:cor:bundles}. So Definition~\ref{dp:def:manifold}, with DP replaced by CA$^{\rm H}$ on open sets and with Banach models and fibers, gives manifolds with tangent bundles, and the alternating and $\widehat\Lambda^k_\pi$ bundles of CA$^{\rm H}$ bundles are CA$^{\rm H}$ bundles.
\item For CA$^{\rm H}$ manifolds $\mathcal M,\mathcal N$ with Banach models, Theorem~\ref{dp:thm:dp-jets} holds with DP replaced by CA$^{\rm H}$. The jet space $J^r(\mathcal M,\mathcal N)$ is a CA$^{\rm H}$ manifold, with CA$^{\rm H}$ forgetful maps, prolongations and functorial maps.
\item Over every $K$, the germ of $\mathsf f$ of Example~\ref{dp:ex:l1} at $0$ is not CA, so it is not the germ of a CA$^{\rm H}$ map.
\end{enumerate}
\end{theorem}

\begin{proof}
\emph{Step 1 (item (1): analytic maps).} Let $f$ be analytic on $U$, and let $x_1\in U$, with radius $R$, bound $M$ and multilinear coefficients $\beta_n$. Then $f$ is DP on $U$. By~\eqref{dp:4b:eq:hasse-expansion}, $D^{(d)}f(x_1+h)=\sum_{n\ge d}D^{(d)}_{a_n}(h)$ near $0$, with $a_n$ the diagonal of $\beta_n$. Expanding $\beta_n(sh+s'z,\dots,sh+s'z)$ shows that $D^{(d)}_{a_n}(h)(z)$ is the sum, over the subsets $\sigma\subseteq\{1,\dots,n\}$ of size $d$, of $\beta_n$ with $z$ in the slots $\sigma$ and $h$ in the others. So $D^{(d)}_{a_n}$ is the diagonal of a bounded $(n-d)$-linear map into $\cP_d(P;W)$ of norm at most $\binom nd\norm{\beta_n}\le M(2/R)^n$. Hence $D^{(d)}f$ is analytic at $x_1$, so it is CA there.

\emph{Step 2 (item (1): laws).} A bounded law $a$ is DP on $P$. By Lemma~\ref{dp:lem:hasse-law}(1),(3) and Lemma~\ref{dp:lem:coefficients-unique}, $x\mapsto D^{(d)}a(x)=D^{(d)}_a(x)$ is a bounded law with values in $\cP_d(P;W)$, and it is $0$ if $d$ exceeds the degree of $a$. So it is CA at every point (Lemma~\ref{dp:lem:cover-flat}(4)).

\emph{Step 3 (item (1): composition).} Let $f$ be CA$^{\rm H}$ on $U$, $g$ CA$^{\rm H}$ on an open subset $U'$ of $W$, and $x\in f^{-1}(U')$. By Theorem~\ref{dp:thm:calculus}(2), $g\circ f$ is DP on the open set $f^{-1}(U')$. In~\eqref{dp:eq:ca-chain}, the map $D^{(\abs\alpha)}g\circ f$ is CA at $x$ by Theorem~\ref{dp:thm:ca-composition}, and each $D^{(n)}f$ is CA at $x$. The pair is CA, and $\mathcal F_\alpha$ is a bounded law, so it is CA at every point. Theorem~\ref{dp:thm:ca-composition} applies again, and finite sums are bounded linear images of tuples. For $d=0$, $g\circ f$ itself is CA by Theorem~\ref{dp:thm:ca-composition}.

\emph{Step 4 (item (1): the formal operations).} For a pair, $D^{(d)}(f_1,f_2)(x)$ is the image of $(D^{(d)}f_1(x),D^{(d)}f_2(x))$ under the isometry $\cP_d(P;W_1)\times\cP_d(P;W_2)\to\cP_d(P;W_1\times W_2)$. Bounded linear postcomposition is composition with an analytic map.

\emph{Step 5 (item (1): closure under $D^{(d)}$).} By Theorem~\ref{dp:thm:differentiable}(2), $D^{(d)}f$ is DP on $U$. Let $d'\ge0$. By~\eqref{dp:4b:eq:hasse-iterated-law}, $D^{(d')}(D^{(d)}f)(x)=\mathsf D_{d',d}(D^{(d'+d)}f(x))$, where $\mathsf D_{d',d}\colon\cP_{d'+d}(P;W)\to\cP_{d'}(P;\cP_d(P;W))$ sends $a$ to $h\mapsto a^{[(d',d)]}(h,\cdot\,)$. This map is linear, and it is bounded by Lemma~\ref{dp:lem:hasse-law}(1). So $D^{(d')}(D^{(d)}f)$ is CA at every point of $U$.

\emph{Step 6 (item (2)).} Theorem~\ref{dp:thm:inverse} gives open sets $V\ni x_0$ and $V'\ni f(x_0)$, with $f$ a bijection $V\to V'$ and inverse $\mathsf g$ DP on $V'$. For $w\in V'$, the chain rule shows that $Df(\mathsf g(w))$ is bijective with inverse $D\mathsf g(w)$. Fix $w\in V'$. Theorem~\ref{dp:thm:inverse} at $\mathsf g(w)$ gives a neighborhood $V_1$ of $\mathsf g(w)$ on which $f$ is injective, and a local inverse $\mathsf g_1$ with values in $V_1$, near $w$. As $\mathsf g$ is continuous, $\mathsf g(w')\in V_1$ for $w'$ near $w$, and $f(\mathsf g(w'))=w'$, so $\mathsf g=\mathsf g_1$ near $w$. As $f$ is CA at $\mathsf g(w)$, Theorem~\ref{dp:thm:ca-inverse}(1) gives that $\mathsf g$ is CA at $w$. So $\mathsf g$ is CA on $V'$. We show by induction on $d$ that $D^{(d)}\mathsf g$ is CA on $V'$. The case $d=0$ is done. On the open set of invertible operators in $\cL(P,W)$, the map $\operatorname{inv}\colon\mathsf T\mapsto\mathsf T^{-1}$ is analytic: near $\mathsf T_1$ it is $\mathsf T_1+\mathsf E\mapsto\sum_n\mathsf T_1^{-1}(-\mathsf E\mathsf T_1^{-1})^n$, a series of diagonals of bounded $n$-linear maps of norm at most $\norm{\mathsf T_1^{-1}}^{n+1}$. As $D\mathsf g=\operatorname{inv}\circ Df\circ\mathsf g$, the case $d=1$ follows from Theorem~\ref{dp:thm:ca-composition}. Let $d\ge2$. Apply~\eqref{dp:eq:ca-chain} to $f\circ\mathsf g=\id$. The only $\alpha$ with $\abs\alpha=1$ and $\deg\alpha=d$ is the unit vector at $d$, and its term is $Df(\mathsf g(w))\circ D^{(d)}\mathsf g(w)$. So
\[
 D^{(d)}\mathsf g(w)=-D\mathsf g(w)\circ\sum_{\abs\alpha\ge2,\ \deg\alpha=d}\mathcal F_\alpha\bigl(D^{(\abs\alpha)}f(\mathsf g(w)),(D^{(n)}\mathsf g(w))_{1\le n\le d}\bigr).
\]
In these terms $\alpha_d=0$, so the entry $n=d$ does not occur and may be replaced by $0$. Composition of operators with laws is bounded bilinear. By induction and Theorem~\ref{dp:thm:ca-composition}, $D^{(d)}\mathsf g$ is CA on $V'$. Implicit functions: as in Step~6 of the proof of Theorem~\ref{dp:thm:ca-inverse}, using item~(1).

\emph{Step 7 (item (3)).} By item~(1), CA$^{\rm H}$ on open sets has the properties (i)--(iv) of Section~\ref{dp:sec:conv}, as in the paragraph on CA on open sets above. The transitions of the tangent bundle are CA$^{\rm H}$ by item~(1) with $d=1$, under the isometry $\cP_1(P;P)=\cL(P,P)$. The maps of the functors on hom spaces are bounded laws, so they are CA$^{\rm H}$.

\emph{Step 8 (item (4)).} The proofs of Theorems~\ref{dp:thm:jet-local} and~\ref{dp:thm:dp-jets} and of Lemma~\ref{dp:lem:chain-law}(3),(4) use of the class only analytic maps (constants, translations, bounded linear maps), bounded laws, composition, pairing, finite sums, restriction, and closure under the maps $D^{(d)}$. Restriction is immediate from Definition~\ref{dp:def:cah}, finite sums are bounded linear images of tuples, and the others are in item~(1). The inverse of a chart change is itself a chart change, so no inverse function theorem enters. The model spaces $P\times W\times\prod_{1\le d\le r}\cP_d(P;W)$, and all the spaces of jets through which these proofs compose, are Banach.

\emph{Step 9 (item (5)).} On $P=\ell^1(\N,K)$ with the Banach target $\Kh$, Lemma~\ref{dp:lem:cover-flat}(6) shows that a CA germ is analytic. By Example~\ref{dp:ex:l1}(3), $\mathsf f$ is not analytic at $0$. As $D^{(0)}f=f$, a CA$^{\rm H}$ map is CA.
\end{proof}

\begin{remark}\label{dp:rem:cah-um}
Let $U\subseteq P$ be open and $W$ Banach. A map $f\colon U\to W$ that is strongly DP at every point of $U$ is CA$^{\rm H}$ on $U$. Indeed, it is DP on $U$. By Lemma~\ref{dp:lem:strong-formal}(4) and Lemma~\ref{dp:lem:coefficients-unique}, $D^{(d)}f$ is strongly DP at every point of $U$ as a map into $(\cP_d(P;W),\lvert\!\lvert\!\lvert\cdot\rvert\!\rvert\!\rvert)$. As $\norm\cdot\le\lvert\!\lvert\!\lvert\cdot\rvert\!\rvert\!\rvert$, it is strongly DP into $\cP_d(P;W)$ with the law norm, by Lemma~\ref{dp:lem:strong-formal}(1), hence CA by Lemma~\ref{dp:lem:cover-flat}(4). So, by Lemma~\ref{dp:lem:cover-flat}(5), if $P$ or $W$ is ultrametric, then a map that is CA on $U$ is CA$^{\rm H}$ on $U$.
\end{remark}

\begin{remark}\label{dp:rem:cah-types}
All germs of Types II, III, IVa and IVb of Definition~\ref{dp:def:types} are bounded laws, with values in a Banach space after completing the target if needed. So they are in CA$^{\rm H}$, by Step~2 of the proof of Theorem~\ref{dp:thm:cah}, which uses no completeness of $K$. The germ of Type I is excluded by Theorem~\ref{dp:thm:cah}(5).
\end{remark}

\subsection{What target descent adds}\label{dp:ssec:ca-descents}

Recall (C2)$_{\rm u}$ and (C2)$_\times$ from Section~\ref{dp:s:ssec-descents}. Consider one more property of a class $\cC$ of germs with Banach values that is closed under translation in the source. It is Definition~\ref{dp:def:s-axiom}.
\begin{itemize}
\item[(S)] Let $P$ be a normed space, $W$ a Banach space, and $b_n\colon P\to W$ laws of degree $n$, $n\ge0$, whose germs at $0$ are in $\cC$, with $\sup_n\norm{b_n}R^n<\infty$ for some $R>0$. Then the germ at $0$ of $h\mapsto\sum_nb_n(h)$, $\norm h<R$, is in $\cC$.
\end{itemize}

\begin{theorem}\label{dp:thm:descents-independent}
Consider germs with Banach values.
\begin{enumerate}
\item Let $K$ be complete. CA has \textup{(C1)}, \textup{(C2)}$_{\rm u}$, \textup{(C3)}, composition, pairing, inverse and implicit functions on Banach spaces, and it contains all bounded laws. It does not contain the germ of $\mathsf f$ of Example~\ref{dp:ex:l1} at $0$, and it fails \textup{(C2)}$_\times$. So, in every characteristic, DP is not forced by the properties of this list without \textup{(C2)}$_\times$.
\item Let $\operatorname{char}K=p$. TA (Definition~\ref{dp:def:ta}) has \textup{(C1)}, \textup{(C2)}$_\times$, composition, pairing, inverse and implicit functions on Banach spaces. It does not contain any germ of the map $\gamma^\sharp_p$ of Theorem~\ref{dp:thm:cp}(3), with values in $\hGamma^p_\Gamma(\mathcal E)$. So DP is not forced by the properties of this list without \textup{(C3)}.
\item Let $\operatorname{char}K=p$. AA and CA are incomparable, and so are TA and CA. The germ of $\mathsf f$ at $0$ lies in $\mathrm{AA}\subseteq\mathrm{TA}$ and not in CA. The germs of $\gamma^\sharp_p$, with values in $\hGamma^p_\Gamma(\mathcal E)$, lie in CA and in neither AA nor TA.
\item A class closed under translation in the source, with \textup{(C1)}, \textup{(C3)} and \textup{(S)}, contains every DP germ with Banach values, in every characteristic and with no target descent. CA fails \textup{(S)}, over every $K$.
\end{enumerate}
\end{theorem}

\begin{proof}
\emph{Step 1 (item (1)).} (C1), (C3) and the bounded laws are Lemma~\ref{dp:lem:cover-flat}(3),(4). Composition, pairing, inverse and implicit functions are Section~\ref{dp:ssec:ca-calculus}. For (C2)$_{\rm u}$, let $j\colon W\to Z$ be an isometry between ultrametric Banach spaces with $j\circ f$ CA. By Theorem~\ref{dp:thm:target-descent}(2), $f$ is DP with laws $a_n$, and the laws of $j\circ f$ are $j\circ a_n$. By Lemma~\ref{dp:lem:cover-flat}(5), $j\circ f$ is strongly DP. By Lemma~\ref{dp:lem:strong-formal}(1), $\lvert\!\lvert\!\lvert a_n\rvert\!\rvert\!\rvert=\lvert\!\lvert\!\lvert j\circ a_n\rvert\!\rvert\!\rvert$. So $f$ is strongly DP, hence CA. The germ of $\mathsf f$ is not CA by Theorem~\ref{dp:thm:cah}(5). By Example~\ref{dp:ex:l1}(2), $j\circ\mathsf f$ is analytic at $0$ for an isometry $j$ of $\Kh$ into a Banach space. So $j\circ\mathsf f$ is CA, and (C2)$_\times$ fails. As $\mathsf f$ is DP at $0$, CA is not DP.

\emph{Step 2 (item (2)).} TA contains AA, hence the analytic germs (Theorem~\ref{dp:thm:chain}). An isometry of Banach spaces has closed range, so Theorem~\ref{dp:thm:target-descent}(2) puts the laws of $j\circ f$ in $j(W)$, and Theorem~\ref{dp:thm:ta-composition}(2),(4) gives target descent, composition and pairing. Inverse and implicit functions are Theorem~\ref{dp:thm:ta-inverse}. By Theorem~\ref{dp:thm:cp}(3) no germ of $\gamma^\sharp_p$ is in $\cC_p$, and $\mathrm{TA}\subseteq\cC_p$ by Theorem~\ref{dp:thm:chain}. As $\gamma^\sharp_p$ is DP at every point, TA is not DP.

\emph{Step 3 (item (3)).} $\mathsf f$ is AA at $0$ by Example~\ref{dp:ex:l1}(2), and $\mathrm{AA}\subseteq\mathrm{TA}$. It is not CA by Theorem~\ref{dp:thm:cah}(5). The map $\gamma^\sharp_p$ is a bounded law, so it is CA at every point. It is not TA at any point by Step~2, and $\mathrm{AA}\subseteq\mathrm{TA}$.

\emph{Step 4 (item (4)).} Let $\cC$ have (C1), (C3), (S). By Lemma~\ref{dp:lem:cover-flat}(3),(4), $\cC$ contains every germ of a bounded law, so the germ at $0$ of each coefficient law of a DP germ. A DP germ at $x_0$ is, by translation, the germ at $0$ of $\sum_na_n$ with $\sup_n\norm{a_n}R^n<\infty$. So (S) applies. For the failure, CA is closed under translation in the source, as the laws and lifts do not change. Put $b_0=0$ and let $b_n=a_n$, $n\ge1$, be the laws of Example~\ref{dp:ex:l1}. These are bounded laws with $\norm{b_n}<\abs\varpi^{-1}$, so each of their germs at $0$ is CA. Their sum $\mathsf f$ is not CA at $0$ (Theorem~\ref{dp:thm:cah}(5)).
\end{proof}

On ultrametric sources the CA germs are the DP germs (Lemma~\ref{dp:lem:cover-flat}(5)), so the difference between the two classes lies on non-ultrametric sources. The germ $\mathsf f$ shows, in every characteristic, that AA is not contained in CA. In characteristic $p$, CA is not contained in AA (Theorem~\ref{dp:thm:descents-independent}(3)). In characteristic $0$ we do not decide whether CA is contained in AA. In characteristic $0$ we also do not decide, in general, whether the CA germs are the analytic germs when neither $P$ nor $W$ is ultrametric; they are for the spaces of Lemma~\ref{dp:lem:cover-flat}(6), and for $P$ with finite coordinates (Theorem~\ref{dp:thm:calculus}(3)). In every characteristic, whether CA$^{\rm H}$ has (C3) is open. See Question~\ref{dp:q:cover}.
%% ---- end sec-15-cover.tex
%% ---- begin sec-16-ledger.tex
\section{What forces each disagreement}\label{dp:sec:ledger}

This section lists the disagreements between DP germs and analytic germs that are stated in this paper and in~\cite{I}, and says what forces each one. It is a list of those disagreements only, and we do not claim that it is complete. In characteristic $p$, on ultrametric sources with non-ultrametric targets, whether DP and AA differ is open (Question~\ref{dp:q:c0}), and a germ that is DP and not AA there would be of none of the types below.

\subsection{The regimes}\label{dp:ssec:regimes}

In Table~\ref{dp:tab:ledger}, ``$=$'' means that DP germs and analytic germs coincide for all spaces of the row, ``$\ne$'' means that they differ for some spaces of the row, and ``open'' means that we do not decide this. The types are those of Definition~\ref{dp:def:types}. The rows marked ``see the text'' are treated in the paragraphs after the table.

\begin{table}[htbp]
\footnotesize
\caption{Regimes. Here ``um'' abbreviates ultrametric, ``sph'' abbreviates spherically complete, and ``discrete'' and ``dense'' refer to the value group of $K$.}\label{dp:tab:ledger}
\begin{tabular}{@{}p{0.05\linewidth}p{0.155\linewidth}p{0.135\linewidth}p{0.135\linewidth}p{0.195\linewidth}p{0.185\linewidth}@{}}
 & $K$ & $P$ & $W$ & DP against An & Where\\
\hline
L0 & $\mathbb R$, $\mathbb C$ & any & any & $=$ & see the text\\
L1 & any & finite-dim. & any if $K$ is complete; otherwise Banach & $=$ & Thm.~\ref{dp:thm:agreement-finite}\\
L2 & $\F_p(t)$; other incomplete $K$ not considered & finite-dim., um & finite-dim., incomplete & $\ne$, Type V & Thm.~\ref{dp:thm:agreement-finite}(3)\\
L3 & char.\ $0$ & um & any & $=$, radius factor $\mathfrak n_K\mathfrak s(P)/2$ & Thm.~\ref{dp:thm:agreement-char0}\\
L4 & char.\ $0$ & non-um, e.g.\ $\ell^1$ & um, e.g.\ $\Kh$ & $\ne$, Type I; sDP $=$ An & Ex.~\ref{dp:ex:l1}, Prop.~\ref{dp:prop:strong-interpolates}(1),(2), Cor.~\ref{dp:cor:two-gaps}\\
L5 & char.\ $0$ & non-um & non-um & $\ne$, Type I, for some; DP against AA open; An against CA open & see the text, Q.~\ref{dp:q:char0}, Q.~\ref{dp:q:cover}\\
\hline
L6 & char.\ $p$ & um & um, sph up to equivalence & $=$ & Thm.~\ref{dp:thm:sph-agreement}(1)\\
L7 & char.\ $p$, discrete, complete or not & um & um & $=$ for $W$ Banach; $W$ incomplete: $\ne$ over $\F_p(t)$ (L2), otherwise open & Thm.~\ref{dp:thm:sph-agreement}(2), Thm.~\ref{dp:thm:agreement-finite}(3), Q.~\ref{dp:q:residual}\\
L8 & char.\ $p$, complete & um; $\widehat P\cong c_0(J)$, or of countable type & um Banach & $=$ & Thm.~\ref{dp:thm:c0-agreement}, Thm.~\ref{dp:thm:conservativity}\\
L9 & char.\ $p$, complete, not sph & um & um, not sph & $\ne$, Type II & \cite[Theorem~7.2(2)]{I}, Prop.~\ref{dp:prop:strict}(1)(c)\\
L10 & char.\ $p$, complete, dense, sph & um, $\widehat P\not\cong c_0(J)$ & um, no equivalent sph norm & open & Q.~\ref{dp:q:residual}\\
L11 & char.\ $p$, incomplete, dense & um, infinite-dim. & um, no equivalent sph norm & $\ne$ for some spaces if $\Kh$ is not sph, by restriction of scalars; if $\Kh$ is sph and $W$ is Banach, L10 over $\Kh$; $W$ incomplete: open for general such spaces & Thm.~\ref{dp:thm:general-norms-agreement}(3)(a), Q.~\ref{dp:q:residual} and the text after it\\
L12 & char.\ $p$, discrete, or $p=2$ & um, e.g.\ $c_0$ & non-um & $\ne$, Type III & Prop.~\ref{dp:prop:strict}(1)(b),(d)\\
L13 & char.\ $p$ odd, $\Kh$ not sph & um & non-um, with the target of Type II as a bounded retract & $\ne$, Type II; for $K$ incomplete, after restriction of scalars & see the text, and the text after Q.~\ref{dp:q:residual}\\
L14 & char.\ $p$ odd, dense, $\Kh$ sph & um & non-um & open & Q.~\ref{dp:q:residual}, \cite[Remark~E.2]{I}\\
L15 & char.\ $p$ & non-um & um & $\ne$, Type I, if $P$ has a complemented copy of $\ell^1(\N)$; sDP against An is DP against An on $(P^{\mathsf u},W)$, for the hull $P^{\mathsf u}$ of Section~\ref{dp:sec:strong} & see the text, Cor.~\ref{dp:cor:two-gaps}\\
L16 & char.\ $p$ & non-um & non-um & $\ne$, Types I, III, IV; here also DP $\ne$ AA & Prop.~\ref{dp:prop:strict}(2), see the text\\
\end{tabular}
\end{table}

\emph{Row L0.} This row is outside the standing assumptions of this paper: over $K=\mathbb R$ or $\mathbb C$ the definitions are read verbatim. Let $a$ be a law of degree $n\ge1$. Comparing coefficients in $a(\sum_i\varepsilon_ix_i)$ gives $\sum_{\varepsilon\in\{\pm1\}^n}\varepsilon_1\cdots\varepsilon_na(\sum_i\varepsilon_ix_i)=2^na^{[1^n]}(x_1,\dots,x_n)$. As in Section~\ref{dp:sec:nonambient}, $a^{[1^n]}$ is $n$\nobreakdash-linear with $a^{[1^n]}(h,\dots,h)=n!\,a(h)$; that algebra does not use the absolute value. So $\norm{a^{[1^n]}(x_1,\dots,x_n)}\le\norm a(\sum_i\norm{x_i})^n$. Scaling to $\norm{x_i}=1/n$ shows that $a^{[1^n]}/n!$ is a lift of $a$ of norm at most $(n^n/n!)\norm a\le\mathrm e^n\norm a$. So a DP germ with radius $R$ and bound $M$ is analytic with radius $R/\mathrm e$ and bound $M$.

\emph{Rows with a retract.} Let $\iota\colon W'\to W$ and $\rho\colon W\to W'$ be bounded linear with $\rho\iota=\id$. If $f$ maps into $W'$, then $\iota\circ f$ is DP if $f$ is, and if $\iota\circ f$ is analytic at a point, then so is $f=\rho\circ\iota\circ f$. Let $k\ge p$, let $E,D$ be as in \cite[Theorem~7.2(2)]{I}, and put $W_0=\cL(\Alt^k(D;K),\Alt^k(E;K))$. The retract gives L13 from Type II, for instance with $W'=W_0$ and $W=W_0\oplus\ell^1(\N,K)$. It gives L5 from L4, with $f=\mathsf f$ and $W=\Kh\oplus\ell^1(\N,\Kh)$. The same applies on the source. Let $\iota'\colon P'\to P$ and $\rho'\colon P\to P'$ be bounded linear with $\rho'\iota'=\id$, and let $f$ be a germ at $0$ on $P'$ that is DP and not analytic. Then $f\circ\rho'$ is DP at $0$, and it is not analytic at $0$, as $f=f\circ\rho'\circ\iota'$. With $f=\mathsf f$ this gives L15. With both kinds of retract it gives the Types I and III in L16, from $\mathsf f$ and from the maps of Type III. For densely valued $K$ and odd $p$ only Type III (iii) is available.

In L13, containing that target isometrically is not enough. Let $\mathsf Y_K$ be the space $\cL(\Alt^k(D;K),\Mult^k(E;K))$. Then $\mathsf Y_K\oplus\ell^1(\N,K)$ contains $W_0$ isometrically, and $A^k_{E,D;K}$, as a map into it, is the diagonal of a bounded $k$\nobreakdash-linear map~\cite[Proposition~2.2]{I}, so it is analytic.

\emph{DP against AA across the table.} They agree on ultrametric $P,W$~\cite[Theorem~14.39(1)]{I} and on $P=\ell^1(J)$~\cite[Theorem~14.39(2)]{I}. They differ in L16, by Proposition~\ref{dp:prop:strict}(2). It is open in L5 (Question~\ref{dp:q:char0}), in L12--L14 (Question~\ref{dp:q:c0}(1), by Theorem~\ref{dp:thm:general-norms-agreement}(3)(c); laws of degree $2$ and of degree below $p$ are ambient, by Propositions~\ref{dp:prop:ultrametric}(1),(4) and~\ref{dp:prop:all-norms}(1)), and in L4 and L15 in general \cite[Question~14.42(4)]{I}.

\subsection{The types}\label{dp:ssec:types}

Recall (C1), (C2), (C3) from Section~\ref{dp:sec:descent}. As in Section~\ref{dp:ssec:ca-descents}, we write (C2)$_\times$ for (C2), along all linear isometries of Banach spaces, and (C2)$_{\rm u}$ for (C2) along linear isometries between ultrametric Banach spaces only.

Let $\mathsf f$ be the germ of Example~\ref{dp:ex:l1}, whose mechanism is described after that example. Let $V,\mu,\nu,R_n,H,H_0,G,c_0$ be as in Section~\ref{dp:sec:nonambient}, and let $\widehat V_\mu,\widehat V_\nu$ be the completions of $(V,\mu),(V,\nu)$. We write $(\lambda,x,w)$ for the points of $H=K\oplus c_0\oplus G$, so that $R_{p+1}(\lambda,x,w)=x^{\otimes p}\otimes w$. We write $f_{c_0},g_{c_0}$ for the maps $f_X,g_X$ of that section, with $X=c_0$. Conditions~(a) and~(b) are those of Proposition~\ref{dp:prop:nonambient}. Let $H^{\rm at},\widehat W^{\rm at},\widehat Y^{\rm at}$ be the spaces of \cite[Definition~14.16]{I} and the completions used after \cite[Theorem~14.17]{I}, with $X=X^{\rm at}$ there, and let $f^{\rm at},\hat g^{\rm at}$ be the maps $f_X,\hat g_X$ of that theorem and of the paragraph after it. Let $\mathcal E$ and $\hGamma^p_\Gamma(\mathcal E)$ be as in Section~\ref{dp:sec:chain}. For $\operatorname{char}K=p$ and $n\ge p$, put
\[
 \gamma^\sharp_n\colon K\oplus\mathcal E\to\hGamma^p_\Gamma(\mathcal E),\qquad\gamma^\sharp_n(\lambda,h)=\lambda^{n-p}h^{\otimes p},
\]
with $\lambda^0=1$, the laws of Proposition~\ref{dp:prop:degree-p}(2), with $\gamma^\sharp_p$ as in Theorem~\ref{dp:thm:cp}.

\begin{definition}[types]\label{dp:def:types}
We name the following germs and laws. Each is DP and not analytic.
\begin{itemize}
\item[] \textbf{Type I.} The germ $\mathsf f$ at $0$, on $\ell^1(\N,K)$, with values in $\Kh$, over every $K$. Its laws have norm less than $\abs\varpi^{-1}$, and a lift of its law of degree $n$ has norm at least $n^n$.
\item[] \textbf{Type II.} For $K$ complete, not spherically complete, of characteristic $p$, and $k\ge p$: the law $A^k_{E,D;K}$, for the ultrametric Banach spaces $E,D$ of \cite[Theorem~7.2(2)]{I}. It is ambient, with no lift.
\item[] \textbf{Type III.} For $\operatorname{char}K=p$: \textup{(i)} for $k\ge p$, the power map $c_0\to\hGamma^k_\pi c_0$, $h\mapsto h^{\otimes k}$, in the cases of Proposition~\ref{dp:prop:strict}(1)(b); \textup{(ii)} the power map $c_0\to\hGamma^p_\pi c_0$, $h\mapsto h^{\otimes p}$, if $K$ is discretely valued or $p=2$; \textup{(iii)} $A^k_{E,E;F}$, $k\ge p$, for the spaces $E,F$ of \cite[Theorem~7.2(1)]{I}. All are ambient, with no lift. In (i), (ii) the source is ultrametric, and the target has no equivalent ultrametric norm. In (iii), the space $F$ has no equivalent ultrametric norm, and over a general $K$ it is not asserted that $E$ is ultrametric \cite[Remark~E.2]{I}.
\item[] \textbf{Type IVa.} For $\operatorname{char}K=p$: \textup{(i)} the laws $R_n$, $n\ge p+1$, into $\widehat V_\mu$, under~(a) or~(b); \textup{(ii)} the laws $R_n$ into $\widehat V_\nu$, for $p=2$; \textup{(iii)} the law $\hat g^{\rm at}\circ f^{\rm at}$ on $H^{\rm at}$, into $\widehat Y^{\rm at}$, over every $K$. These are laws that are not ambient.
\item[] \textbf{Type IVb.} For $\operatorname{char}K=p$ and every $K$: the laws $\gamma^\sharp_n$, $n\ge p$. They are not ambient.
\item[] \textbf{Type V.} Over $K=\F_p(t)$ with the $t$-adic absolute value, and for $k\ge p$: the joint action $B$ on the morphisms of \cite[Theorem~H.4]{I}, with incomplete finite-dimensional values.
\end{itemize}
For $K$ discretely valued and $p$ odd, the laws $R_n$ into $\widehat V_\nu$ are not ambient, by Proposition~\ref{dp:prop:nonambient}(2), and we do not assign them a type (Question~\ref{dp:q:odd-c0}).
\end{definition}

We justify the assertions of Definition~\ref{dp:def:types}. The norm bound in Type I is Example~\ref{dp:ex:l1}, and the bound on lifts is Step~3 of its proof. The germ $\mathsf f$ is DP and not analytic by the same example. The maps of Types II and III are ambient and analytic at no point (Proposition~\ref{dp:prop:strict}(1)), so they have no lift (Lemma~\ref{dp:lem:lift-criterion}). In Type III (i), (ii), suppose that the target had an equivalent ultrametric norm. It would be an ultrametric Banach space, and by Theorem~\ref{dp:thm:c0-agreement}(2) the power map would be the diagonal of a bounded multilinear map, so it would be analytic. This is false. In (iii), the assertion on $F$ is \cite[Theorem~7.2(1)]{I}. The laws of Type IV are bounded, so they are DP (Theorem~\ref{dp:thm:calculus}(1)). They are not ambient, by Lemma~\ref{dp:lem:type-iv} and Proposition~\ref{dp:prop:degree-p}(2), so they are not analytic (Lemma~\ref{dp:lem:lift-criterion}). Type V is Theorem~\ref{dp:thm:agreement-finite}(3).

\begin{lemma}[Type IV]\label{dp:lem:type-iv}
Let $\operatorname{char}K=p$. A class means a class of germs with Banach values.
\begin{enumerate}
\item For every $K$ of characteristic $p$ and every $n\ge p+1$, $R_n$ into $\widehat V_\mu$ is in $\mathrm{AA}^\circ$ at every point. Every class that has \textup{(C1)}, \textup{(C2)}$_\times$, composition and pairing contains the germs of $R_n$ at every point. Under~\textup{(a)} or~\textup{(b)} it is not ambient of degree $n$.
\item If $\mathfrak a_p(c_0)<\infty$, then $R_n$ into $\widehat V_\nu$ is in $\mathrm{AA}^\circ$ at every point. This holds if $p=2$, and then $R_n$ into $\widehat V_\nu$ is not ambient of degree $n$. For odd $p$, whether $\mathfrak a_p(c_0)<\infty$ is open.
\item For every $K$ of characteristic $p$, $\hat g^{\rm at}\circ f^{\rm at}$ on $H^{\rm at}$ is in $\mathrm{AA}^\circ$ at every point, and it is a law of degree $p+1$ that is not ambient.
\item For every $n\ge p$, no germ of $\gamma^\sharp_n$ is in TA.
\end{enumerate}
So, as $\mathrm{AA}^\circ\subseteq\mathrm{TA}$ (Theorem~\ref{dp:thm:chain}), a law that is not ambient can lie in TA.
\end{lemma}

\begin{proof}
\emph{Step 1 (item (1): membership).} On $H$ let $\mathsf a'(\lambda,x,w)=x^{\otimes p}$, into the completion $\hGamma^p_\pi c_0$ of $\Gamma^p_\pi(c_0)$, and $\mathsf a''(\lambda,x,w)=w$, into $\widehat G$. Let $j$ be the isometry of $\hGamma^p_\pi c_0$ into the completion of $(c_0^{\otimes p},\pi)$. Then $j\circ\mathsf a'$ is the diagonal of the bounded $p$\nobreakdash-linear map $((\lambda_i,x_i,w_i))_{i\le p}\mapsto x_1\otimes\dots\otimes x_p$, so it is analytic at every point. And $\mathsf a''$ is bounded linear. So a class with (C1) and (C2)$_\times$ contains the germs of $\mathsf a'$ and $\mathsf a''$. By the definition of $\mu$, the map $(s,w)\mapsto s\otimes w$ from $\Gamma^p_\pi(c_0)\times G$ to $(V,\mu)$ has norm at most $1$. So it extends to a bounded bilinear map $\boxtimes$ from $\hGamma^p_\pi c_0\times\widehat G$ to $\widehat V_\mu$, which is analytic. We have $R_{p+1}=\boxtimes\circ(\mathsf a',\mathsf a'')$. For $n>p+1$, $R_n$ is the pair $(\lambda^{n-p-1},R_{p+1})$, with analytic first entry, followed by the bounded bilinear scalar multiplication $\Kh\times\widehat V_\mu\to\widehat V_\mu$. So pairing and composition put $R_n$ in the class. The class $\mathrm{AA}^\circ$ has these properties by definition.

\emph{Step 2 (item (1): not ambient).} Assume~(a) or~(b), and suppose that $R_n$ into $\widehat V_\mu$ is ambient of degree $n$: $j_1\circ R_n$ is the diagonal of a bounded $n$\nobreakdash-linear map $\beta$, for a linear isometry $j_1$ of $\widehat V_\mu$. For $(x,w)\in H_0$ we have $g_{c_0}\circ f_{c_0}(x,w)=R_n(1,x,w)$. Expanding $\beta((1,x,w),\dots,(1,x,w))$ by multilinearity shows that $j_1\circ g_{c_0}\circ f_{c_0}$ is analytic at $0$. By the one-point lemma \cite[Proposition~14.9(2)]{I}, the law $g_{c_0}\circ f_{c_0}$ into $\widehat V_\mu$ is ambient, so it is AA at $0$ \cite[Proposition~14.9(1)]{I}. This contradicts Theorem~\ref{dp:thm:c0-failure}(3).

\emph{Step 3 (item (2)).} As in Remark~\ref{dp:rem:nu-mu}, $\nu\le\mathfrak a_p(c_0)\mu$ on $V$. So if $\mathfrak a_p(c_0)<\infty$, the identity of $V$ extends to a bounded linear map $\widehat V_\mu\to\widehat V_\nu$. It is analytic, and composing with it keeps $\mathrm{AA}^\circ$. If $p=2$, then $\mathfrak a_2(c_0)\le\frac32\mathfrak c_K^{-2}$ by Proposition~\ref{dp:prop:ultrametric}(4), and non-ambience is Proposition~\ref{dp:prop:nonambient}(2) under~(b). The last sentence of item~(2) is part of \cite[Question~14.42(4)]{I}; see Question~\ref{dp:q:c0}.

\emph{Step 4 (item (3)).} By the paragraph after \cite[Theorem~14.17]{I}, $f^{\rm at}$ is AA into $\widehat W^{\rm at}$ at every point, and $\hat g^{\rm at}$ is the diagonal of a bounded bilinear map. So $\hat g^{\rm at}\circ f^{\rm at}$ is in $\mathrm{AA}^\circ$ at every point. By \cite[Theorem~14.17]{I} it is AA at no point. So it is not ambient, as an ambient law is AA at $0$ \cite[Proposition~14.9(1)]{I}.

\emph{Step 5 (item (4), $n=p$).} By Theorem~\ref{dp:thm:cp}(3), no germ of $\gamma^\sharp_p$ is in $\cC_p$, and $\mathrm{TA}\subseteq\cC_p$ by Theorem~\ref{dp:thm:chain}.

\emph{Step 6 (item (4), $n>p$).} Let $(\lambda_1,x_1)\in K\oplus\mathcal E$ with $\lambda_1\ne0$. Write the increment as $(\mathsf m,\mathsf h)$. The law of degree $p$ of $\gamma^\sharp_n$ at $(\lambda_1,x_1)$ is the part of degree $p$ in $(\mathsf m,\mathsf h)$ of $(\lambda_1+\mathsf m)^{n-p}(x_1+\mathsf h)^{\otimes p}$. At $\mathsf m=0$ it is $\mathsf h\mapsto\lambda_1^{n-p}\mathsf h^{\otimes p}$. If that law were ambient, then, restricting to $\mathcal E$, dividing by $\lambda_1^{n-p}$ and precomposing with the projection of $K\oplus\mathcal E$ onto $\mathcal E$, the law $\gamma^\sharp_p$ would be ambient, against Proposition~\ref{dp:prop:degree-p}(2). So the germ at $(\lambda_1,x_1)$ is not in $\cC_p$, hence not in TA. At a point $(0,x_1)$: by Theorem~\ref{dp:thm:chain}, a germ in TA is in $\cC_p^{\rm loc}$, that is, in $\cC_p$ at all nearby points. Points with $\lambda_1\ne0$ occur arbitrarily near. So the germ at $(0,x_1)$ is not in TA.
\end{proof}

\subsection{What forces them}\label{dp:ssec:forcing}

In this subsection a class means a class of germs with Banach values, as in Section~\ref{dp:sec:descent}, closed under bounded linear postcomposition. Type V has incomplete values and is treated separately. A class \emph{contains} a type if it contains the germs of its maps at every point of their domains, and at $0$ for Type I.

\begin{theorem}[forcing]\label{dp:thm:forcing}
\leavevmode
\begin{enumerate}
\item \textup{(Type I.)} Every class with \textup{(C1)} and \textup{(C2)}$_\times$ contains Type I. Let $K$ be complete. The properties \textup{(C1)}, \textup{(C3)}, composition, pairing, inverse and implicit functions, all bounded laws and alternating bundles do not force Type I, even jointly: the class CA, on open sets and for Banach models and fibers, has all of them and does not contain Type I. Among germs with ultrametric Banach targets, sDP has \textup{(C1)}, \textup{(C2)}$_{\rm u}$, \textup{(C3)}, composition and every $A^k_{E,D;F}$ with $F$ an ultrametric Banach space, and it does not contain Type I.
\item \textup{(Type II.)} Let $K$ be complete, not spherically complete, of characteristic $p$. Each of the following forces Type II: \textup{(C1)} with \textup{(C2)}$_{\rm u}$; \textup{(C1)} with \textup{(C3)}; containing the transition families $u\mapsto g(u)^*$ of the alternating bundle of the shear bundle \cite[Proposition~10.1, Corollary~10.2]{I}, with $F=K$.
\item \textup{(Type III.)} Let $\operatorname{char}K=p$. Each of \textup{(C1)} with \textup{(C2)}$_\times$ and \textup{(C1)} with \textup{(C3)} forces Type III. For Banach $E,F$, containing the families $u\mapsto g(u)^*$ with $D=E$ forces $A^k_{E,E;F}$.
\item \textup{(Type IV.)} Let $\operatorname{char}K=p$.
\begin{enumerate}
\item \textup{(C1)}, \textup{(C2)}$_\times$, composition and pairing force Type IVa.
\item \textup{(C1)} with \textup{(C3)} forces every bounded law, so Types IVa, IVb, and the laws $R_n$ into $\widehat V_\nu$.
\item TA has \textup{(C1)}, \textup{(C2)}$_\times$, composition, pairing, inverse and implicit functions and bundle lifting, and it does not contain Type IVb.
\item Let $r\ge p$. A class closed under restriction to affine subspaces that contains the operator family $\mathsf T^\psi$ of $J^r(X_{\mathcal E};\hGamma^p_\Gamma(\mathcal E))$, for $\psi=\psi_{\mathcal E}^{\pm1}$, on the shear manifold $X_{\mathcal E}$ of Proposition~\ref{dp:prop:jets-not-ambient}(4), contains the law $h\mapsto h^{\otimes p}$ on $\mathcal E$.
\end{enumerate}
\item \textup{(Type V.)} Among germs with arbitrary normed values, containing the joint actions forces Type V. Into the completion of its target, the map of Type V is analytic.
\end{enumerate}
\end{theorem}

\begin{proof}
\emph{Step 1 (item (1)).} By Example~\ref{dp:ex:l1}(2) there are a Banach space $Z$ and a linear isometry $j\colon\Kh\to Z$ with $j\circ\mathsf f$ analytic at $0$. Apply (C1) and (C2)$_\times$. The coefficient law $a_n$ of $\mathsf f$ of degree $n$ has $a_n^{[1^n]}(e_1,\dots,e_n)=\mathsf e_n$. So its strong norm is at least $n^n$, and by Lemma~\ref{dp:lem:coefficients-unique} $\mathsf f$ is not sDP at $0$. As $\Kh$ is ultrametric, $\mathsf f$ is not CA at $0$, by Lemma~\ref{dp:lem:cover-flat}(5). Now let $K$ be complete. By Theorem~\ref{dp:thm:descents-independent}(1), CA has (C1), (C3), composition, pairing, inverse and implicit functions and all bounded laws. CA is an open condition (Proposition~\ref{dp:prop:ca-open}), and on open sets it carries the alternating bundles, by the paragraph \emph{CA on open sets} of Section~\ref{dp:sec:cover}. The properties of sDP are Theorem~\ref{dp:thm:strong-descent}, Theorem~\ref{dp:thm:strong-composition} and Proposition~\ref{dp:prop:strong-bundle}.

\emph{Step 2 (every bounded law from \textup{(C1)}, \textup{(C3)}).} Let $a\colon P\to W$ be a bounded law with $W$ Banach, and let $x_0=q_P(\tilde x_0)$. By Lemma~\ref{dp:lem:laws-cover}, $a\circ q_P$ is analytic at $\tilde x_0$. So (C1) and (C3) put the germ of $a$ at $x_0$ in the class. This gives the (C1), (C3) clauses of items~(2) and~(3), and item~(4)(b).

\emph{Step 3 (the (C2) clauses of items (2), (3)).} Let $\mathsf Y_F=\cL(\Alt^k(D;F),\Mult^k(E;F))$. By \cite[Proposition~2.2]{I}, $A^k_{E,D;F}$ followed by the isometry $j_{E,F}\circ{-}$ into $\mathsf Y_F$ is the diagonal of a bounded $k$\nobreakdash-linear map, so it is analytic at every point. For item~(2), take $F=K$: then $\mathsf Y_K$ and $\cL(\Alt^k(D;K),\Alt^k(E;K))$ are ultrametric Banach spaces, and (C1), (C2)$_{\rm u}$ apply. For Type III (iii) apply (C1), (C2)$_\times$. For Type III (i), (ii), the power map into $\hGamma^k_\pi c_0$, followed by the isometry into the completion of $(c_0^{\otimes k},\pi)$, is the diagonal of $(x_1,\dots,x_k)\mapsto x_1\otimes\dots\otimes x_k$. Apply (C1), (C2)$_\times$.

\emph{Step 4 (the bundle clauses).} Let $\mathsf D_1=D\oplus E$, $g(u)(d,e)=(d+ue,e)$, and let $\mathsf J,\mathsf{Res}$ be the bounded linear maps of the proof of \cite[Proposition~10.1]{I}, with $\mathsf J\colon\Alt^k(D;F)\to\Alt^k(\mathsf D_1;F)$ and $\mathsf{Res}\colon\Alt^k(\mathsf D_1;F)\to\Alt^k(E;F)$. There $\mathsf D_1$ and $\mathsf{Res}$ are called $H$ and $R$. By \cite[(10.1)]{I}, $\mathsf{Res}\circ g(u)^*\circ\mathsf J=A^k_{E,D;F}(u)$. The map $\mathsf S\mapsto\mathsf{Res}\circ\mathsf S\circ\mathsf J$ is bounded linear. So a class closed under bounded linear postcomposition that contains the germs of $u\mapsto g(u)^*$ contains those of $A^k_{E,D;F}$.

\emph{Step 5 (item (4)(a),(c)).} By Definition~\ref{dp:def:ta}, $\mathrm{AA}^\circ$ is the smallest class with the properties of item~(a), and Type IVa lies in it, by Lemma~\ref{dp:lem:type-iv}(1),(2),(3). For item~(c), TA contains the analytic germs and has the listed properties by Theorems~\ref{dp:thm:ta-composition}, \ref{dp:thm:ta-inverse} and Corollary~\ref{dp:cor:ta-bundles}. It does not contain Type IVb, by Lemma~\ref{dp:lem:type-iv}(4).

\emph{Step 6 (item (4)(d)).} Let $Z_p=\cP_p(\mathcal E;\hGamma^p_\Gamma(\mathcal E))$. Fix $z_1\in\mathcal E$ and $\mathsf c_1\in K$. Restrict the operator family $\mathsf T^\psi$ to the slice $v\mapsto(v,z_1,\mathsf c_1)$ and apply the bounded linear map $\mathsf C$ of Step~4 of the proof of Theorem~\ref{dp:thm:analytic-jets}. By~\eqref{dp:13:eq:compression} this gives the map $v\mapsto\pm\ev_p(v)$ into $\cL(Z_p,\hGamma^p_\Gamma(\mathcal E))$, with $\ev_p(v)(a)=a(v)$. Evaluation at the law $h\mapsto h^{\otimes p}$ in $Z_p$, and multiplication by $\pm1$, are bounded linear. They give $v\mapsto v^{\otimes p}$.

\emph{Step 7 (item (5)).} The first assertion is the definition. The second is Theorem~\ref{dp:thm:agreement-finite}(2), as the source is finite-dimensional.
\end{proof}

Item~(4)(d) is about $\mathcal E$. To reach $\gamma^\sharp_p$ on $K\oplus\mathcal E$, one more precomposition, with the projection, is needed.

\subsection{Closure under series of laws}\label{dp:ssec:series}

Outside the list in Theorem~\ref{dp:thm:forcing}(1), and besides (C2)$_\times$, one more property forces Type I. We recall it from Section~\ref{dp:ssec:ca-descents}.

\begin{definition}\label{dp:def:s-axiom}
Let $\cC$ be a class of germs with Banach values, closed under translation in the source. It has \emph{\textup{(S)}} if the following holds. Let $P$ be a normed space, $W$ a Banach space, and $b_n\colon P\to W$ laws of degree $n$, $n\ge0$, whose germs at $0$ are in $\cC$, with $\sup_n\norm{b_n}R^n<\infty$ for some $R>0$. Then the germ at $0$ of $h\mapsto\sum_nb_n(h)$, $\norm h<R$, is in $\cC$.
\end{definition}

\begin{proposition}[series of laws]\label{dp:prop:s-axiom}
\leavevmode
\begin{enumerate}
\item A class with \textup{(C1)} and \textup{(S)} contains Type I, over every $K$, with no target descent.
\item The classes An, sDP and CA fail \textup{(S)} on $(\ell^1(\N,K),\Kh)$, over every $K$; so does $\mathrm{CA}^{\rm H}$ for $K$ complete.
\item DP has \textup{(S)}. If $\operatorname{char}K=0$ and $P$ is ultrametric, the analytic germs on $P$ have the property in \textup{(S)}, for every Banach $W$.
\item A class with \textup{(C1)}, \textup{(C3)} and \textup{(S)} contains every DP germ with Banach values, with no target descent.
\end{enumerate}
\end{proposition}

\begin{proof}
\emph{Step 1 (item (1)).} The law $a_n(h)=\mathsf e_nh_1\cdots h_n$ on $\ell^1(\N,K)$ is the diagonal of the bounded map $(x^1,\dots,x^n)\mapsto\mathsf e_nx^1_1\cdots x^n_n$, so its germ at $0$ is analytic. Also $\norm{a_n}<\abs\varpi^{-1}$ by Example~\ref{dp:ex:l1}(1). Apply (S) with $R=1$ and $a_0=0$.

\emph{Step 2 (item (2)).} Each $a_n$ is analytic, so its germ is in An. It is in sDP by Lemma~\ref{dp:lem:strong-formal}(5), as $\Kh$ is ultrametric. It is in CA by Lemma~\ref{dp:lem:cover-flat}(4), and in $\mathrm{CA}^{\rm H}$ by Theorem~\ref{dp:thm:cah}(1). The sum $\mathsf f$ is not analytic (Example~\ref{dp:ex:l1}(3)), and it is neither sDP nor CA, by Step~1 of the proof of Theorem~\ref{dp:thm:forcing}. As $\mathrm{CA}^{\rm H}\subseteq\mathrm{CA}$, it is not in $\mathrm{CA}^{\rm H}$.

\emph{Step 3 (item (3)).} Let $M=\sup_n\norm{b_n}R^n$. For $\norm h<R$ the sum of the $\norm{b_n(h)}$ is at most $M\sum_n(\norm h/R)^n$, which is finite, and $W$ is complete. So $\sum_nb_n$ is DP at $0$ with radius $R$ and bound $M$, by Definition~\ref{dp:def:laws}(3). If $\operatorname{char}K=0$ and $P$ is ultrametric, this germ is analytic, by Theorem~\ref{dp:thm:agreement-char0}(2).

\emph{Step 4 (item (4)).} By Step~2 of the proof of Theorem~\ref{dp:thm:forcing}, (C1) and (C3) put every bounded law in the class. A DP germ at $x_0$ is the translate of $\sum_n\mathsf a_n$, with coefficient laws $\mathsf a_n$ and $\sup_n\norm{\mathsf a_n}R^n<\infty$. Apply (S).
\end{proof}

\begin{proposition}[in characteristic $0$]\label{dp:prop:s-char0}
Let $\operatorname{char}K=0$, and let $\cC$ be a class of germs with Banach values, closed under translation in the source, with \textup{(C1)} and \textup{(S)}. Then $\cC$ contains every DP germ with Banach values. So a class with \textup{(C1)} and translations, inside DP, has \textup{(S)} if and only if it equals DP. In particular AA has \textup{(S)} if and only if $\mathrm{AA}=\mathrm{DP}$, and TA has \textup{(S)} if and only if $\mathrm{TA}=\mathrm{DP}$.
\end{proposition}

\begin{proof}
\emph{Step 1 (laws).} Let $a$ be a law of degree $n\ge1$. By Theorem~\ref{dp:thm:degree-threshold}(1), as $n!\ne0$, $a$ is the diagonal of the bounded $n$\nobreakdash-linear map $a^{[1^n]}/n!$. So the germ of $a$ at $0$ is analytic, hence in $\cC$. A law of degree $0$ is constant.

\emph{Step 2 (series).} A DP germ at $x_0$ is the translate of $\sum_n\mathsf a_n$, with coefficient laws $\mathsf a_n$ and $\sup_n\norm{\mathsf a_n}R^n<\infty$. Apply (S). The second sentence follows with Proposition~\ref{dp:prop:s-axiom}(3). AA and TA contain the analytic germs, lie in DP and are closed under translation.
\end{proof}

\begin{proposition}[in characteristic $p$]\label{dp:prop:s-fails}
Let $\operatorname{char}K=p$. There are laws $\mathsf b_n\colon K\oplus\mathcal E\to\hGamma^p_\Gamma(\mathcal E)$ of degree $n$, $n\ge p$, each the diagonal of a bounded $n$\nobreakdash-linear map, with $\norm{\mathsf b_n}\le1$, whose sum is not in TA at $0$. So AA and TA fail \textup{(S)}.
\end{proposition}

\begin{proof}
\emph{Step 1 (the laws).} Let $\mathcal E_m,T^{(m)}$ be as in Proposition~\ref{dp:prop:degree-p}, and $\operatorname{pr}_m\colon\mathcal E\to\mathcal E_m$ the contractive projection. By~\eqref{dp:eq:degree-p-ratio}, for each $n\ge p$ there is $m_n$ with $\norm{T^{(m_n)}}_\Gamma\ge n^n\pi(T^{(m_n)})$ in $\Gamma^p(\mathcal E_{m_n})$. Here $\pi(T^{(m_n)})>0$, as $T^{(m_n)}\ne0$ and $\pi$ is a norm on $\mathcal E_{m_n}^{\otimes p}$. Put $\mathsf b_n(\lambda,h)=\lambda^{n-p}(\operatorname{pr}_{m_n}h)^{\otimes p}$. Then $\norm{\mathsf b_n}\le1$.

\emph{Step 2 (diagonals).} The law $\mathsf b_n$ factors through the bounded linear map $\id_K\times\operatorname{pr}_{m_n}$ onto the finite-dimensional space $K\oplus\mathcal E_{m_n}$. On it, the law is the diagonal of a bounded $n$\nobreakdash-linear map into the Banach space $\hGamma^p_\Gamma(\mathcal E)$, by Theorem~\ref{dp:thm:agreement-finite}(2). So the germ of $\mathsf b_n$ at $0$ is analytic, hence AA and TA.

\emph{Step 3 (a coefficient).} Suppose that $\sum_n\mathsf b_n$ is in TA at $0$. By Theorem~\ref{dp:thm:chain} it is in $\cC_p^{\rm u}$: for some $R,M$ the linearizations $\mathsf L_n$ of its coefficient laws, which are the $\mathsf b_n$, satisfy~\eqref{dp:eq:cpu}. Let $\mathsf z_0=(1,0)$. For $h\in\mathcal E_{m_n}$, the coefficient of $\lambda^{n-p}$ in $\mathsf L_n((\lambda\mathsf z_0+h)^{\otimes n})=\mathsf b_n(\lambda,h)=\lambda^{n-p}h^{\otimes p}$ gives $\mathsf L_n(\mathsf z_0^{\otimes(n-p)}\cdot h^{\otimes p})=h^{\otimes p}$. As the powers span, $\mathsf L_n(\mathsf z_0^{\otimes(n-p)}\cdot s)=s$ for $s\in\Gamma^p(\mathcal E_{m_n})$.

\emph{Step 4 (the bound).} By~\eqref{dp:eq:cpu}, with $N$ there equal to our $n$, with $n$ there equal to $p$, and with $h=\mathsf z_0$, we get $\norm s_\Gamma\le MR^{-n}\pi(s)$. Here $\norm s_\Gamma$ is computed in $\Gamma^p(\mathcal E)$, and $\pi$ in $(K\oplus\mathcal E)^{\otimes p}$, which is at most $\pi$ in $\mathcal E_{m_n}^{\otimes p}$, as inclusions are contractive. Applying $\operatorname{pr}_{m_n}^{\otimes p}$ to a representation of $s$ in $\mathcal E$ shows that $\norm s_\Gamma$ computed in $\mathcal E$ is at least $\norm s_\Gamma$ computed in $\mathcal E_{m_n}$. Take $s=T^{(m_n)}$: $n^n\le MR^{-n}$ for all $n\ge p$, which is false.
\end{proof}

So, in characteristic $p$, (S) separates DP from every other class of Tables~\ref{dp:tab:compare} and~\ref{dp:tab:compare-b}. In characteristic $0$ it separates DP from each such class that differs from DP. The property (S) is close to the definition of DP.

\subsection{Summary}\label{dp:ssec:summary}

Table~\ref{dp:tab:types} says, for each type, which failure it answers and what forces it. In the row of Type I, the analytic germs on $(\ell^1,W)$ with $W$ ultrametric Banach and $\operatorname{char}K=0$ have (C2)$_{\rm u}$ and (C3) (Corollary~\ref{dp:cor:two-gaps}(3)), composition \cite[Proposition~14.26, proof, Step~1]{I}, every $A^k$ (Theorem~\ref{dp:thm:degree-threshold}(1)) and jet sections (Theorem~\ref{dp:thm:analytic-jets}(4)(a),(2)); they fail (C2)$_\times$ and (S).

\begin{table}[htbp]
\small
\caption{The types.}\label{dp:tab:types}
\begin{tabular}{@{}p{0.08\linewidth}p{0.46\linewidth}p{0.4\linewidth}@{}}
Type & Failure it answers & What forces it\\
\hline
I & (C2)$_\times$ and (S), on $(\ell^1,W)$, $W$ ultrametric Banach, $\operatorname{char}K=0$. By Theorem~\ref{dp:thm:forcing}(1) the other properties of that list do not force it & (C1) with (C2)$_\times$; (C1) with (S)\\
II & The functor $\alpha_k$ of~\cite[n\textsuperscript{o}~7.8.1]{Bou} (Theorem~\ref{dp:thm:seven-eight-dp}), its bundles \cite[Corollary~10.2]{I}, scalar jets (Corollary~\ref{dp:cor:jet-dichotomy}) & (C1) with (C2)$_{\rm u}$; (C1) with (C3); the alternating bundle of the shear bundle\\
III & The power maps of $c_0$ (Proposition~\ref{dp:prop:strict}(1)); $\Alt^k$ with values in $F$ of \cite[Theorem~7.2(1)]{I} & (C1) with (C2)$_\times$; (C1) with (C3); for (iii), the alternating bundle of the shear bundle\\
IVa & Composition for AA \cite[Theorem~14.17]{I}, Theorem~\ref{dp:thm:c0-failure} & (C1), (C2)$_\times$, composition, pairing; (C1) with (C3)\\
IVb & (C3) for AA, TA (Proposition~\ref{dp:prop:other-classes}(2), Theorem~\ref{dp:thm:singles-out}(5)); jets (Proposition~\ref{dp:prop:jets-not-ambient}) & (C1) with (C3); jets, for $r\ge p$, on $\mathcal E$. Not composition (Theorem~\ref{dp:thm:forcing}(4)(c))\\
V & Incomplete values; none after completion & Containing the joint actions\\
\end{tabular}
\end{table}

Tables~\ref{dp:tab:compare} and~\ref{dp:tab:compare-b} compare the classes. They hold for $\operatorname{char}K=p$, $K$ complete and Banach values, unless a row or a cell says otherwise. A dash means that the entry is not considered here. Their last row says that, between ultrametric spaces, AA, TA, sDP, CA and DP are one class of germs. There $\mathrm{CA}^{\rm H}$ is this class too, as a DP germ is DP on a ball (Theorem~\ref{dp:thm:whole-ball}), its derivatives $D^{(j)}$ are DP (Theorem~\ref{dp:thm:differentiable}), and DP germs on an ultrametric source are CA (Lemma~\ref{dp:lem:cover-flat}(5)).

Two cells of Table~\ref{dp:tab:compare-b} need one more remark, as CA$^{\rm H}$ germs are germs of maps on open sets. A map into a Banach space that is analytic at a point is analytic on a ball around it (Section~\ref{dp:ssec:reach}), so its germ is in $\mathrm{CA}^{\rm H}$ by Theorem~\ref{dp:thm:cah}(1). So $\mathrm{CA}^{\rm H}$ fails (C2)$_\times$: with $j$ as in Example~\ref{dp:ex:l1}(2), $j\circ\mathsf f$ is in $\mathrm{CA}^{\rm H}$ at $0$, and $\mathsf f$ is not, by Theorem~\ref{dp:thm:cah}(5). And if $\operatorname{char}K=0$ and $P$ or $W$ is ultrametric, then $\mathrm{CA}^{\rm H}=\mathrm{An}$. The inclusion $\mathrm{An}\subseteq\mathrm{CA}^{\rm H}$ is the second sentence of this paragraph. Conversely $\mathrm{CA}^{\rm H}\subseteq\mathrm{CA}\subseteq\mathrm{DP}$, for $P$ ultrametric $\mathrm{DP}=\mathrm{An}$ (Theorem~\ref{dp:thm:agreement-char0}(2)), and for $W$ ultrametric $\mathrm{CA}=\mathrm{sDP}=\mathrm{An}$ (Lemma~\ref{dp:lem:cover-flat}(5), Proposition~\ref{dp:prop:strong-interpolates}(1),(2)).

\begin{table}[htbp]
\small
\caption{Comparison, first part.}\label{dp:tab:compare}
\begin{tabular}{@{}p{0.15\linewidth}p{0.29\linewidth}p{0.23\linewidth}p{0.23\linewidth}@{}}
 & An & AA & TA\\
\hline
Composition & yes, \cite[Proposition~14.26, proof, Step~1]{I} & no, \cite[Theorem~14.17]{I} & yes, Thm.~\ref{dp:thm:ta-composition}\\
Inverse functions & --- & no, \cite[Corollary~14.19(1)]{I} & yes, Thm.~\ref{dp:thm:ta-inverse}\\
Charts: compatibility is transitive & yes, by composition, \cite[Proposition~14.26, proof, Step~1]{I} & no, \cite[Corollary~14.19(3)]{I} & yes, Cor.~\ref{dp:cor:ta-bundles}\\
Alternating bundle & no for $k\ge p$, \cite[Theorem~7.2, Corollary~10.2]{I} & sections depend on the trivialization, \cite[Corollary~14.19(4)]{I} & yes, Cor.~\ref{dp:cor:ta-bundles}\\
Jets, $r\ge p$ & no for some non-ultrametric $W$, Thm.~\ref{dp:thm:analytic-jets}(3),(6)(a); scalar, ultrametric models: no iff $K$ is not spherically complete, Cor.~\ref{dp:cor:jet-dichotomy} & no, Prop.~\ref{dp:prop:jets-not-ambient} & no, Prop.~\ref{dp:prop:jets-not-ambient}\\
(C2)$_{\rm u}$ & no if $K$ is complete and not spherically complete, Thm.~\ref{dp:thm:forcing}(2); yes if $K$ is discrete or $W$ is spherically complete up to equivalence, Thm.~\ref{dp:thm:strong-um}, Cor.~\ref{dp:cor:two-gaps}(1); else open & yes, Prop.~\ref{dp:prop:other-classes}(2) & yes, Thm.~\ref{dp:thm:ta-composition}\\
(C2)$_\times$ & no, Ex.~\ref{dp:ex:l1} & yes, Prop.~\ref{dp:prop:other-classes}(2) & yes, Thm.~\ref{dp:thm:ta-composition}\\
(C3) & no, Prop.~\ref{dp:prop:other-classes}(1) & no, Prop.~\ref{dp:prop:other-classes}(2) & no, Thm.~\ref{dp:thm:singles-out}(5)\\
(S) & no on $\ell^1$, Prop.~\ref{dp:prop:s-axiom}(2) & no, Prop.~\ref{dp:prop:s-fails}; char.\ $0$: iff $=$ DP, Prop.~\ref{dp:prop:s-char0} & no, Prop.~\ref{dp:prop:s-fails}; char.\ $0$: iff $=$ DP, Prop.~\ref{dp:prop:s-char0}\\
Contains An & & yes, \cite[Definition~14.7]{I} & yes, Thm.~\ref{dp:thm:chain}\\
All bounded laws & no, Prop.~\ref{dp:prop:strict}(1) & no, Prop.~\ref{dp:prop:strict}(2) & no, Lem.~\ref{dp:lem:type-iv}(4)\\
Excludes Type I & yes, Ex.~\ref{dp:ex:l1}(3) & no, Ex.~\ref{dp:ex:l1}(2) & no, Ex.~\ref{dp:ex:l1}(2), Thm.~\ref{dp:thm:chain}\\
Char.\ $0$: agrees with An & & $P$ ultrametric: yes, \cite[Proposition~14.40(2)]{I}; on $\ell^1$: no, Ex.~\ref{dp:ex:l1} & $P$ ultrametric: yes, Thm.~\ref{dp:thm:agreement-char0}; on $\ell^1$: no, Ex.~\ref{dp:ex:l1}\\
Between ultrametric spaces & Rows L6--L11 of Table~\ref{dp:tab:ledger} & $=$ DP, \cite[Theorem~14.39(1)]{I} & $=$ DP, Thm.~\ref{dp:thm:singles-out}(4)\\
\end{tabular}
\end{table}

\begin{table}[htbp]
\small
\caption{Comparison, second part.}\label{dp:tab:compare-b}
\begin{tabular}{@{}p{0.15\linewidth}p{0.3\linewidth}p{0.27\linewidth}p{0.18\linewidth}@{}}
 & sDP & $\mathrm{CA}^{\rm H}$ & DP\\
\hline
Composition & yes if one of $P,W,Y$ has an equivalent ultrametric norm, Thm.~\ref{dp:thm:strong-composition}; no on $\ell^1$, Thm.~\ref{dp:thm:strong-failure} & yes, Thm.~\ref{dp:thm:cah}(1); for CA, Thm.~\ref{dp:thm:ca-composition} & yes, Thm.~\ref{dp:thm:calculus}(2)\\
Inverse functions & no on $\ell^1$, Thm.~\ref{dp:thm:strong-failure}; yes on ultrametric Banach spaces, Thm.~\ref{dp:thm:strong-inverse} & yes, Thm.~\ref{dp:thm:cah}(2); for CA, Thm.~\ref{dp:thm:ca-inverse} & yes, Thm.~\ref{dp:thm:inverse}\\
Charts: compatibility is transitive & no on $\ell^1$, Prop.~\ref{dp:prop:strong-bundle} & yes, Thm.~\ref{dp:thm:cah}(1) & yes, Prop.~\ref{dp:prop:categories}\\
Alternating bundle & the law: yes; manifolds: on models with an equivalent ultrametric norm; not with a complemented $\ell^1(\N)$; else open, Prop.~\ref{dp:prop:strong-bundle} & yes, Thm.~\ref{dp:thm:cah}(3); also CA on open sets, Thm.~\ref{dp:thm:forcing}(1) & yes, Cor.~\ref{dp:cor:bundles}\\
Jets, $r\ge p$ & on models with an equivalent ultrametric norm, Prop.~\ref{dp:prop:strong-bundle}, Thm.~\ref{dp:thm:dp-jets}; not with a complemented $\ell^1(\N)$; else open & yes, Thm.~\ref{dp:thm:cah}(4) & yes, Thm.~\ref{dp:thm:dp-jets}\\
(C2)$_{\rm u}$ & yes, Thm.~\ref{dp:thm:strong-descent} & yes, Thm.~\ref{dp:thm:descents-independent}(1), Rem.~\ref{dp:rem:cah-um} & yes, Thm.~\ref{dp:thm:minimal}(1)\\
(C2)$_\times$ & yes, Thm.~\ref{dp:thm:strong-descent} & no, Ex.~\ref{dp:ex:l1}(2), Thm.~\ref{dp:thm:cah}(1),(5) & yes, Thm.~\ref{dp:thm:minimal}(1)\\
(C3) & yes, Thm.~\ref{dp:thm:strong-descent} & for CA: yes, Lem.~\ref{dp:lem:cover-flat}(3); for $\mathrm{CA}^{\rm H}$: open, Q.~\ref{dp:q:cover} & yes, Thm.~\ref{dp:thm:minimal}(1)\\
(S) & no on $\ell^1$, Prop.~\ref{dp:prop:s-axiom}(2) & no on $\ell^1$, Prop.~\ref{dp:prop:s-axiom}(2) & yes, Prop.~\ref{dp:prop:s-axiom}(3)\\
Contains An & if $P$ or $W$ is ultrametric, Prop.~\ref{dp:prop:strong-interpolates}; not on $(\ell^1,W)$ for some non-ultrametric $W$, Ex.~\ref{dp:ex:an-not-sdp} & yes, Thm.~\ref{dp:thm:cah}(1) & yes, Thm.~\ref{dp:thm:calculus}(3)\\
All bounded laws & yes, Lem.~\ref{dp:lem:strong-formal}(3) & yes, Thm.~\ref{dp:thm:cah}(1) & yes, Thm.~\ref{dp:thm:calculus}(1)\\
Excludes Type I & yes, Thm.~\ref{dp:thm:forcing}(1) & yes, Thm.~\ref{dp:thm:cah}(5) & no, Ex.~\ref{dp:ex:l1}(1)\\
Char.\ $0$: agrees with An & $P$ or $W$ ultrametric: yes, Prop.~\ref{dp:prop:strong-interpolates}; strictly smaller on $(\ell^1,W)$ for some non-ultrametric $W$, Ex.~\ref{dp:ex:an-not-sdp} & $P$ or $W$ ultrametric: yes, Thm.~\ref{dp:thm:cah}(1), Thm.~\ref{dp:thm:agreement-char0}(2), Lem.~\ref{dp:lem:cover-flat}(5), Prop.~\ref{dp:prop:strong-interpolates}; $\ell^1(J)$, $W$ Banach: yes, Lem.~\ref{dp:lem:cover-flat}(6); else open, Q.~\ref{dp:q:cover} & $P$ ultrametric: yes, Thm.~\ref{dp:thm:agreement-char0}; on $\ell^1$: no, Ex.~\ref{dp:ex:l1}\\
Between ultrametric spaces & $=$ DP, Lem.~\ref{dp:lem:strong-formal}(3) & $=$ DP, see the text & Rows L6--L11 of Table~\ref{dp:tab:ledger}\\
\end{tabular}
\end{table}

On ultrametric sources $\mathrm{DP}=\mathrm{sDP}=\mathrm{CA}$, so (C1) and (C3) with $\ell^1$ covers already give DP; on general sources DP can be larger than CA, by Type I, and that excess is what (C2)$_\times$ or (S) adds.
%% ---- end sec-16-ledger.tex
%% ---- begin sec-17-questions.tex
\section{Questions}\label{dp:sec:questions}

This section lists what this paper and~\cite{I} do not decide, with what they do give. We prove no case of the conjectures below. We write $c_0=c_0(\N,K)$.

\begin{question}\label{q:quest:polynomial-laws}
Let $K$ be nonarchimedean of characteristic $p>0$. Is every bounded polynomial law on $c_0(\N,K)$ ambient of its degree?
\end{question}

Bounded polynomial laws of degree below $p$ are ambient (Section~\ref{dp:sec:nonambient}), and for $p=2$ so are those of degree $2$ on ultrametric normed spaces (Proposition~\ref{dp:prop:ultrametric}(1),(4)). By Proposition~\ref{dp:prop:ultrametric}(1),(3), every bounded polynomial law of degree $n$ on $c_0(\N,K)$ is ambient of degree $n$ if and only if every bounded polynomial law of degree $n$ on every ultrametric normed space over $K$ is. As $K\oplus c_0(\N,K)$, with the maximum norm, is isometric to $c_0(\N,K)$, Proposition~\ref{dp:prop:ultrametric}(1),(2) shows that the degrees with a positive answer form an initial segment of the positive integers: a negative answer in some degree gives a negative answer in every higher degree, and a positive answer in infinitely many degrees gives one in all. Frobenius laws are ambient, and every bounded polynomial law of degree $p$ on a normed space $P$, with values in any normed space, is ambient of degree $p$ if and only if $P$ satisfies $(\mathrm N_P)$ with some constant (Theorem~\ref{dp:thm:frobenius}(2),(3)); so for $p\ge3$, in degree $p$ the question asks whether $c_0(\N,K)$ satisfies $(\mathrm N_{c_0(\N,K)})$ with some constant. Over every $K$ and in every degree $n\ge p$, Proposition~\ref{dp:prop:degree-p}(2) gives a law that is not ambient of degree $n$, on a space that is a Banach space if $K$ is complete; if $K$ is discretely valued or $p=2$, Proposition~\ref{dp:prop:nonambient}(2) gives such laws in every degree $n\ge p+1$ also on the space $H=K\oplus c_0(\N,K)\oplus G$ of Section~\ref{dp:sec:nonambient}. These spaces have no equivalent ultrametric norm: they contain $\ell^1$-sums of finite-dimensional quotients of weighted $\ell^1$-spaces, and both proofs obtain their lower bounds, for $\nu$ and for $\norm\cdot_\Gamma$, by adding up contributions of the coordinates of these $\ell^1$-spaces.

\begin{question}\label{dp:q:char0}
Let $\operatorname{char}K=0$. Is $\sup_{n\ge1}\mathfrak a_n(P)^{1/n}<\infty$ for every normed space $P$? Is $\mathrm{TA}=\mathrm{DP}$?
\end{question}

By \cite[Proposition~14.37]{I}, the first question is the second part of \cite[Question~14.42(1)]{I}: a positive answer for $P$ says that every DP germ from $P$ into a Banach space is AA. The supremum is finite for ultrametric $P$ and for $\ell^1(J)$ \cite[Theorem~14.39(2),(3)]{I}, and in finite dimension (Theorem~\ref{dp:thm:agreement-finite}(2)). For general $P$ the bound of Proposition~\ref{dp:prop:all-norms} is not geometric. By Proposition~\ref{dp:prop:s-char0}, AA has the property (S) of Definition~\ref{dp:def:s-axiom} if and only if $\mathrm{AA}=\mathrm{DP}$, and the same holds for TA. By Proposition~\ref{dp:prop:all-sources}, for $K$ complete, the first part of \cite[Question~14.42(1)]{I}, whether $\mathfrak k(P)<\infty$ for every normed $P$, is a question on finite-dimensional spaces.

\begin{question}\label{dp:q:c0}
\leavevmode
\begin{enumerate}
\item Let $\operatorname{char}K=p$. Is $\sup_{n\ge1}\mathfrak a_n(c_0)^{1/n}<\infty$?
\item Let $\operatorname{char}K=p$. Among germs with ultrametric sources and Banach targets, are the DP germs the smallest class with \textup{(C1), (C2), (C3)}?
\item Over every $K$, is $(P,W)$ tame for every normed $P$ and every ultrametric $W$?
\end{enumerate}
\end{question}

For $n<p$ and for $n=2$, $\mathfrak a_n(c_0)$ is finite, by Propositions~\ref{dp:prop:all-norms} and~\ref{dp:prop:ultrametric}(4). So Question~\ref{q:quest:polynomial-laws} is undecided only for $n\ge\max(3,p)$. Item~(1) concerns non-ultrametric targets only, as $(P,W)$ is tame for ultrametric $P,W$ \cite[Theorem~14.39(1)]{I}. Item~(3) is the first sentence of \cite[Question~14.42(4)]{I}; the other sentences there are item~(1) and Question~\ref{q:quest:polynomial-laws}. In each degree the constant in item~(3) is finite \cite[Theorem~14.39(4)]{I}.

Items~(1) and~(2) are equivalent.

\emph{Step 1 ((1) implies (2)).} Let $f\colon U\to W$ be DP at $x_0$, with $P$ ultrametric and $W$ Banach, and let $q^{\rm u}_P$ be the cover of Theorem~\ref{dp:thm:ultrametric-minimal}. Then $f\circ q^{\rm u}_P$ is DP (Proposition~\ref{dp:prop:source-descent}). By Proposition~\ref{dp:prop:ultrametric}(3), $\mathfrak a_n(c_{00}(I_P))\le\mathfrak c_K^{-n}\mathfrak a_n(c_0)$, which is geometric, so $f\circ q^{\rm u}_P$ is AA \cite[Theorem~14.38(1)]{I}. By \cite[Theorem~14.11]{I} there is a linear isometry $\mathsf j$ of $W$ into a Banach space with $\mathsf j\circ f\circ q^{\rm u}_P$ analytic. Then (C1), (C2), (C3) put $f$ in the class, and all sources involved are ultrametric. The DP germs have (C1), (C2), (C3), by Theorem~\ref{dp:thm:minimal}(1).

\emph{Step 2 (a class with the three properties).} Let $\mathfrak M$ be the class of germs $(P,U,f,x_0)$, with $P$ ultrametric and Banach values, for which there are an ultrametric $P'$, a bounded linear surjection $Q\colon P'\to P$ with a preimage constant, and $\tilde x_0\in Q^{-1}(x_0)$ with $f\circ Q$ AA at $\tilde x_0$. It has (C1), with $Q=\id$. It has (C2), as AA has (C2) (Proposition~\ref{dp:prop:other-classes}(2)). It has (C3), as a composite of two such surjections is again one.

\emph{Step 3 ((2) implies (1)).} Assume~(2). Then $\mathfrak M$ contains every DP germ on $c_{00}(\N)$ into a Banach space $W$. For such a germ $f$ at $x_0$, let $Q,\tilde x_0$ be as in Step~2, with preimage constant $\mathsf d$. Choose $\tilde e_i\in P'$ with $Q\tilde e_i=e_i$ and $\norm{\tilde e_i}\le\mathsf d$. As $P'$ is ultrametric, $\mathsf s(x)=\sum_ix_i\tilde e_i$ is a bounded linear map, and $Q\mathsf s=\id$. So $f(x)=(f\circ Q)(\tilde x_0+\mathsf s(x-x_0))$ near $x_0$, and $f$ is AA at $x_0$ \cite[Proposition~14.12(1)(b)]{I}. By \cite[Proposition~14.37]{I}, $(c_{00}(\N),W)$ is tame for every Banach $W$. Take for $W$ the $c_0$-sum of the completions of the spaces $\Gamma^n_{\mathrm s}(c_{00}(\N))$, as in Step~4 of the proof of Theorem~\ref{dp:thm:general-norms-agreement}. This gives $\sup_n\mathfrak a_n(c_{00}(\N))^{1/n}<\infty$. The inclusion of $c_0([m],K)$ into $c_{00}(\N)$ and the coordinate projection back are contractive, so $\mathfrak a_n(c_0([m],K))\le\mathfrak a_n(c_{00}(\N))$. Now Proposition~\ref{dp:prop:ultrametric}(3) gives~(1).

\begin{question}\label{dp:q:residual}
Let $\operatorname{char}K=p$.
\begin{enumerate}
\item Let $K$ be complete, densely valued and spherically complete. Does $(\mathrm D_K)$ hold, that is, is every DP germ between ultrametric Banach spaces analytic? Does $(\mathrm B_K)$ hold? Does $(\mathrm B_K)$ imply $(\mathrm D_K)$?
\item Let $p$ be odd, and let $K$ be densely valued and spherically complete. Let $P$ be ultrametric and $W$ a Banach space with no equivalent ultrametric norm. Is every DP germ from $P$ to $W$ analytic?
\item Let $K$ be discretely valued and complete, $n\ge p$, $P$ ultrametric and infinite-dimensional, and $W$ ultrametric and incomplete. Is every law of degree $n$ from $P$ to $W$ the diagonal of a bounded $n$\nobreakdash-linear map into $W$?
\end{enumerate}
\end{question}

These are the rows L10, L14 and L7 of Table~\ref{dp:tab:ledger}. In~(1), $(\mathrm D_K)$ and $(\mathrm B_K)$ are the statements given before Theorem~\ref{dp:thm:field-level}, which proves that $(\mathrm D_K)$ implies $(\mathrm B_K)$, and neither statement is decided for this $K$ (\cite[Section~12]{I}). Also in~(1), by Theorem~\ref{dp:thm:conservativity}, a germ that is DP and not analytic needs $\widehat P\not\cong c_0(J,K)$ and $W$ with no equivalent ultrametric spherically complete norm. In~(3), for $W$ Banach the answer is yes, by Theorem~\ref{dp:thm:sph-agreement}(2). For $P$ finite-dimensional it is yes, by Theorem~\ref{dp:thm:agreement-finite}(1), and for $P=c_{00}(J)$ it is yes, by Theorem~\ref{dp:thm:c0-agreement}(1). The row L5, characteristic $0$ with $P$ and $W$ not ultrametric, is in Questions~\ref{dp:q:char0} and~\ref{dp:q:cover}.

\emph{Incomplete fields, Banach targets.} Let $K$ be incomplete, $P$ ultrametric and $W$ a Banach space. Then $W$ and the completion $\widehat P$ are normed spaces over $\Kh$, and bounded $K$-multilinear maps between normed spaces over $\Kh$ are $\Kh$-multilinear, by continuity. So a germ between normed spaces over $\Kh$ with Banach target is analytic over $K$ if and only if it is analytic over $\Kh$. Coefficient laws from $P$ to $W$ extend to $\widehat P$ over $\Kh$ with the same bounds \cite[Lemma~Q.13(1),(2)]{I}, a DP germ extends through its law series, and a lift over $\Kh$ restricts to a lift over $K$. So every DP germ from an ultrametric space into an ultrametric Banach space over $K$ is analytic if and only if $(\mathrm D_{\Kh})$ holds. If $\Kh$ is spherically complete, this is~(1) over $\Kh$, and in the same way the case of $W$ Banach with no equivalent ultrametric norm is~(2) over $\Kh$. If $\Kh$ is not spherically complete, the answer is no. Let $k\ge p$, let $E,D$ be the spaces of \cite[Theorem~7.2(2)]{I} over $\Kh$, and read them over $K$. The map $A^k_{E,D;\Kh}$ is a law over $K$, so it is DP, and it is analytic at no point over $K$, by the above and \cite[Theorem~7.2(2)]{I}. Its source and target are ultrametric Banach spaces. Its target has no equivalent ultrametric spherically complete norm, by Theorem~\ref{dp:thm:sph-agreement}(1), and its source is infinite-dimensional, by Theorem~\ref{dp:thm:agreement-finite}(2). By the retract argument of Section~\ref{dp:ssec:regimes}, the same holds with the target $W_0\oplus\ell^1(\N,\Kh)$, where $W_0$ is the target above, and this has no equivalent ultrametric norm.

\emph{Incomplete fields, incomplete targets.} Over the incomplete field $\F_p(t)$ with the $t$-adic absolute value, the question in~(3), without the completeness of $K$ and the infinite dimension of $P$, has the answer no in every degree $n\ge p$ (Theorem~\ref{dp:thm:agreement-finite}(3), Lemma~\ref{dp:lem:lift-criterion}). For other incomplete discretely valued $K$ we do not decide it. Over incomplete densely valued $K$ with $W$ incomplete, whether every DP germ between ultrametric spaces is analytic is not decided here either, beyond the cases of Theorem~\ref{dp:thm:general-norms-agreement}(3)(a).

\begin{question}\label{dp:q:odd-c0}
Let $\operatorname{char}K=p$ be odd.
\begin{enumerate}
\item Let $K$ be densely valued. Is $\mathfrak l_p(c_0)=\infty$?
\item Is the law $R_{p+1}$ of Proposition~\ref{dp:prop:nonambient}, into the completion of $(V,\nu)$, in TA?
\end{enumerate}
\end{question}

Item~(1) is part of \cite[Question~14.42(2)]{I}. A positive answer gives the conclusions of Theorem~\ref{dp:thm:c0-failure} over this $K$, with $G$ replaced by a bad space for $X=c_0$, $n=p$, by \cite[Proposition~14.18]{I}.

Item~(2) is equivalent to $\mathfrak a_p(c_0)<\infty$. If $\mathfrak a_p(c_0)<\infty$, the answer is yes, by Lemma~\ref{dp:lem:type-iv}(2) and Theorem~\ref{dp:thm:chain}. Conversely, let the germ of $R_{p+1}$ at some point be in TA. By Theorem~\ref{dp:thm:chain} it is in $\cC_p^{\rm u}$. Its coefficient law of degree $p+1$ at that point is $R_{p+1}$, by Theorem~\ref{dp:thm:mixed}(2), so~\eqref{dp:eq:cpu} holds for the linearization $L$ of $R_{p+1}$, with some radius $R'$ and bound $M$. As $Q_1=0$, we have $N_1=0$ and $G_1=\ell^1(I_1,K)$. Let $w\in G_1$ be a unit vector of this space, and let $\phi$ be the coordinate form with $\phi(w)=1$, composed with the projection of $G$ onto $G_1$; it has norm $1$. Apply~\eqref{dp:eq:cpu} with $N=p+1$, $n=p$, with $(0,0,w)\in H$ as $h$, and with $t\in\Gamma^p(c_0)$. For $t=x^{\otimes p}$, comparing coefficients in $R_{p+1}(0,sx,s'w)=s^ps'\,x^{\otimes p}\otimes w$ gives $L((0,0,w)\cdot t)=t\otimes w$, and the powers span $\Gamma^p(c_0)$. So $\nu(t\otimes w)\le MR'^{-p-1}\pi(t)$, with $\pi$ the projective norm of $c_0^{\otimes p}$, which dominates that of $H^{\otimes p}$. Applying $\id\otimes\phi$ to a representation of $t\otimes w$ for $\nu$ gives $\norm t_\Gamma\le\nu(t\otimes w)$. Hence $\mathfrak a_p(c_0)\le MR'^{-p-1}$. Whether $\mathfrak a_p(c_0)<\infty$ for odd $p$ is Question~\ref{q:quest:polynomial-laws} in degree $p$.

\begin{question}\label{dp:q:ta}
Let $\operatorname{char}K=p$.
\begin{enumerate}
\item Is $\mathrm{AA}^\circ=\mathrm{TA}$? Does $\mathrm{AA}^\circ$ have the inverse function theorem?
\item Is there a constant $\mathsf A$, depending only on $p$, with $\mathfrak t_{p^2+1}\le\mathsf A\,\mathfrak t^{(2)}_{p^2+1}$ on $\Gamma^{p^2+1}(P)$ for every normed $P$?
\item Is every germ of $\cC_p^{\rm loc}$ in $\cC_p^{\rm u}$?
\end{enumerate}
\end{question}

By Theorem~\ref{dp:thm:composition-closure}, for $K$ complete and Banach targets, the first part of~(1) asks whether $\mathrm{TA}=\bigcup_\ell\mathrm{TA}_\ell$. The proof of Theorem~\ref{dp:thm:ta-inverse} does not bound the depth of the trees (Remark~\ref{dp:rem:ta-depth}). Degrees $N\le p^2$ are decided by Proposition~\ref{dp:prop:base-p}(3); the first undecided case is~(2).

\begin{conjecture}\label{dp:conj:ta}
Let $\operatorname{char}K=p$. Among germs with Banach values, $\mathrm{TA}\subsetneq\cC_p^{\rm u}$.
\end{conjecture}

The inclusion is Theorem~\ref{dp:thm:chain}.

\begin{question}\label{dp:q:strong}
\leavevmode
\begin{enumerate}
\item Which models $P$ with no equivalent ultrametric norm admit a pseudogroup of strongly DP local diffeomorphisms? Is failure tied to a complemented copy of $\ell^1(\N,K)$?
\item For which $P$ is $\sup_n\bigl(\sup_a\lvert\!\lvert\!\lvert a\rvert\!\rvert\!\rvert/\norm a\bigr)^{1/n}<\infty$, with $a$ over the nonzero laws of degree $n$ on $P$, with values in any normed space?
\end{enumerate}
\end{question}

Proposition~\ref{dp:prop:strong-bundle} treats models with an equivalent ultrametric norm and models with a complemented copy of $\ell^1(\N,K)$. A positive answer to~(2) for $P$ gives $\mathrm{DP}=\mathrm{sDP}$ on $P$ for all $W$. It holds for ultrametric $P$ (Lemma~\ref{dp:lem:strong-formal}(3)) and fails on $\ell^1(\N,K)$, by the coefficient laws $\mathsf e_nh_1\cdots h_n$ of Example~\ref{dp:ex:l1}.

\begin{question}\label{dp:q:cover}
\leavevmode
\begin{enumerate}
\item Let $\operatorname{char}K=0$. Is $\mathrm{CA}=\mathrm{An}$ for all $P,W$ with no equivalent ultrametric norm, with $P$ infinite-dimensional and not linearly homeomorphic to any $\ell^1(J)$ or $\ell^1_{00}(J)$?
\item Let $W$ be Banach. If $f$ is CA on an open set, is $Df$ CA? Is $\mathrm{CA}^{\rm H}=\mathrm{CA}$? Does $\mathrm{CA}^{\rm H}$ have \textup{(C3)}?
\item Let $\operatorname{char}K=0$. Is $\mathrm{CA}\subseteq\mathrm{AA}$?
\item Is there a largest class containing CA, with composition and inverse functions, that excludes the germ $\mathsf f$ of Example~\ref{dp:ex:l1}? Is $\mathrm{CA}^{\rm H}$ maximal for any list of properties?
\item On $\ell^1(J)$ with $J$ infinite and $W$ incomplete, is $\mathrm{CA}=\mathrm{An}$?
\end{enumerate}
\end{question}

By Lemma~\ref{dp:lem:cover-flat}(3), among germs with Banach values, CA is the smallest class with (C1), (C3). So, for $W$ Banach,~(1) asks whether An has (C3) there. A counterexample needs a surjection $Q$ with no bounded linear section $\mathsf s$, since $f=(f\circ Q)\circ\mathsf s$ for a section. In~(2), the answer on $Df$ is yes if $P$ or $W$ is ultrametric (Remark~\ref{dp:rem:cah-um}, Lemma~\ref{dp:lem:cover-flat}(5)). In~(3), the reverse inclusion fails, by $\mathsf f$, and in characteristic $p$ the two classes are incomparable (Theorem~\ref{dp:thm:descents-independent}). A class as in~(4) fails (S), by Proposition~\ref{dp:prop:s-axiom}. In~(5), for $W$ Banach the answer is yes (Lemma~\ref{dp:lem:cover-flat}(6)).

\begin{question}\label{dp:q:base-change}
In items~(1)--(4), let $K'$ be a valued extension of $K$.
\begin{enumerate}
\item Let $K$ be densely valued. Is $\mathfrak e_n(P,K')=1$ for all normed $P$ and all $n$? Is there a universal bound better than the $C_0^n$ of Theorem~\ref{dp:thm:base-change-law}?
\item Let $K$ be complete and $[K':K]=\infty$, and let $X_0\subseteq X$ be a linear isometry with $X$ not ultrametric. Is $(X_0)_{K'}\to X_{K'}$ bounded below?
\item Let $K$ be complete, not spherically complete, of characteristic $p$, and $K'$ spherically complete. Between ultrametric spaces, with maxima, is there a germ $f$, DP at its base point, that is not analytic over $K$ with $f_{K'}$ analytic over $K'$?
\item Let $K,K'$ be as in~\textup{(3)}, let $k\ge p$, and let $E,D$ be as in \cite[Theorem~7.2(2)]{I}. Is $(A^k_{E,D;K})_{K'}$ analytic over $K'$?
\item Let $K$ be complete, not spherically complete, of characteristic $p$, and $P$ ultrametric. Is every law $a\colon P\to K$ of degree $n$ the diagonal of a bounded $n$\nobreakdash-linear form $\beta$? Is there a constant $C$ with $\norm\beta\le C^n\norm a$?
\end{enumerate}
\end{question}

In~(1), with maxima the constant is $1$ over densely valued $K$ (Theorem~\ref{dp:thm:base-change-um}). A positive answer to~(2) is what the proof of Theorem~\ref{dp:thm:aa-base-change}(2) needs for general norms. For $W$ complete, analyticity descends along $K'$ if $\widehat{K'}$ has a bounded retraction onto $\Kh$, and no spherically complete $K'$ over a complete, not spherically complete $K$ has one (Proposition~\ref{dp:prop:an-descent}). Which other extensions have one is not determined here, and neither is descent along an extension with none, or for incomplete $W$. Items~(3), (4) are the test cases. By Theorem~\ref{dp:thm:an-base-change}(2), with maxima, (4) asks whether $A^k_{E,D;K}$ is analytic over $K$ as a map into the completion of $(W_{E,D})^{\max}_{K'}$, with $W_{E,D}=\cL(\Alt^k(D;K),\Alt^k(E;K))$. That germ is DP and AA through an ultrametric space, so its base change is too (Theorems~\ref{dp:thm:germ-base-change}, \ref{dp:thm:aa-base-change}(2)). In~(5), a counterexample needs $\widehat P$ not linearly homeomorphic to any $c_0(J,K)$ (Theorem~\ref{dp:thm:conservativity}(1)(d)). The operator-valued question is decided by \cite[Theorem~7.2(2)]{I}, and the scalar one is not.

\begin{question}\label{dp:q:jets}
\leavevmode
\begin{enumerate}
\item Which sequences of laws $(a_i)$ are given by a $C^\infty_{\mathrm{BGN}}$ map $f$ and a point $x$, through finite Taylor expansions $f(x+\tau v)=\sum_{i\le k}\tau^ia_i(v)+o(\abs\tau^k)$?
\item Let $K$ be densely valued and spherically complete, of characteristic $p$. Does $(\mathrm{Sec}_k)$ hold for all ultrametric $P$ and all ultrametric Banach $W$?
\item Over an incomplete $K$, and for $P$ not ultrametric with scalar values over a spherically complete $K$, when does $(\mathrm{Sec}_k)$ hold?
\item Let $\cG$ be a group in DP manifolds, modeled on $P$. Let $\operatorname{Dist}_r(\cG)$ be the dual of the jets of order $r$ of scalar DP germs at the identity. Is $\bigcup_r\operatorname{Dist}_r(\cG)$ an algebra under the convolution dual to the group law, through the truncated composition of Lemma~\ref{dp:lem:chain-law}? If $\operatorname{char}K=p$, does it see a $p$\nobreakdash-operation?
\item Does $\mathfrak a_n(V,W)<\infty$ for $1\le n\le r$ imply $\mathfrak a_n(V\times V\times K,W)<\infty$ for $1\le n\le r$?
\end{enumerate}
\end{question}

In~(1), for DP germs into a complete $W$ the answer is the geometric sequences (Theorems~\ref{dp:thm:bgn}(4), \ref{dp:thm:quasi-analytic}). For $P$ Banach, a positive answer to Question~\ref{dp:q:residual}(1) gives~(2), as $\ev_k$ is a law into an ultrametric Banach space, and $(\mathrm{Sec}_k)$ is the analyticity of $\ev_k$ at one point (Theorem~\ref{dp:thm:analytic-jets}(1)). In~(4), by Theorem~\ref{dp:thm:normed-roby} the space $\operatorname{Dist}_r(\cG)$ receives a canonical map from $\bigoplus_{j\le r}\Gamma^j_{\mathrm s}(P)$.

\begin{question}\label{dp:q:quotients}
\leavevmode
\begin{enumerate}
\item Let $P$ be ultrametric. For a map $f$ that is DP at $x_0$, is $\norm{f^{[k]}(z)}\le M\sigma^{-k}$ on a full box of $U^{[k]}$: point coordinates near $x_0$, direction coordinates of norm at most $1$, scalar coordinates small?
\item Let $P$ be two-dimensional. If $f$ is defined on a polydisc in $P$ and, along every affine coordinate line, is given by one power series converging on the whole disc, is $f$ analytic near the center?
\item Let $K$ be discretely valued with finite residue field. If $f$ is bounded on $B(x_0,R)$ and, on every slice of dimension at most $2$, is a series of homogeneous polynomials converging on the $R$-ball, is $f$ DP at $x_0$ with some radius?
\item Let $K$ be a local field and $W$ complete. What is the infimum of $\rho_{\mathrm{DP}}(f,x_1)$, over normed $P$ and over $f$ that are the sum of their law series on $B(x_0,R)$, with $\rho_{\mathrm{DP}}(f,x_0)=R$ and $\norm{x_1-x_0}=\varrho<R$? Is the behaviour of the radius in Remark~\ref{dp:rem:whole-ball} sharp as $\varrho\to R$?
\item Let $W$ be complete, and let $K$ be discretely valued with $\mathsf q\le16$ residue elements and $\abs\vartheta=1/\mathsf q$. If $f$ is the sum of its law series on $B(x_0,\rho_{\mathrm{DP}}(f,x_0))$ and $\varrho=\norm{x_1-x_0}<\rho_{\mathrm{DP}}(f,x_0)$, is $\rho_{\mathrm{DP}}(f,x_1)\ge\rho_{\mathrm{DP}}(f,x_0)-\varrho$?
\end{enumerate}
\end{question}

In~(1), for $W$ complete, a bound on a box that contains $\mathcal Z_k(\rho)$ implies DP (Proposition~\ref{dp:prop:div-bgn}(4)). On $\ell^1(\N,K)$ no such bound holds for the germ $\mathsf f$ of Example~\ref{dp:ex:l1} at $0$. Put $z_0=0$ and $z_i=(z_{i-1},(e_i,0,\dots,0),0)$, with $(e_i,0,\dots,0)\in P_{i-1}$. By Theorem~\ref{dp:thm:bgn}(1),(3) and induction, $\mathsf f^{[k]}(z_k)$ is the coefficient of $s_1\cdots s_k$ in $\mathsf f(\sum_is_ie_i)$. This is $\mathsf e_k$, of absolute value at least $k^k$.

In~(2), the map of Example~\ref{dp:ex:slice-needed}(1) is not separately analytic in this sense. In~(3), Theorem~\ref{dp:thm:banach-steinhaus} gives this for $P$ Banach when the homogeneous parts are bounded laws. Values do not bound coefficients over such $K$: for $\mathsf u(x)=(x^{\mathsf q}-x)/\vartheta$ the iterate $\mathsf u^{\circ l}$ is bounded by $1$ on the unit ball, has degree $\mathsf q^l$, and its linear coefficient has absolute value $\abs\vartheta^{-l}$. The uniform radius in~(3) cannot be dropped: the map $g_2$ of Example~\ref{dp:ex:slice-needed}(2) is locally bounded, analytic on every finite-dimensional affine subspace of $c_0$, and not DP at $0$. In~(5), Theorem~\ref{dp:thm:whole-ball}(2) gives the inequality over densely valued $K$, and Example~\ref{dp:ex:radius-drop} shows that it fails for $\mathsf q\ge17$.

\begin{question}\label{dp:q:functors}
Let $\operatorname{char}K=p$. For a normed space $P$ let $P^{(1)}=K\otimes_{\mathrm{Frob}}P$, with $c\otimes\lambda x=c\lambda^p\otimes x$, and let $\norm\cdot_{\mathrm s}$ be the largest seminorm with $\norm{c\otimes x}_{\mathrm s}\le\abs c\norm x^p$. Let $\widehat P^{(1)}$ be the separated completion and $\mathrm{Fr}_P\colon P\to\widehat P^{(1)}$, $x\mapsto1\otimes x$.
\begin{enumerate}
\item Is $\mathrm{Fr}_P$ the diagonal of a bounded $p$\nobreakdash-linear map, for every Banach space $P$?
\item Let $K$ be densely valued, and let $E$ be an ultrametric Banach space. Here $E_{K^{1/p}}$ and $\pi_{\max}$ are as in Section~\ref{dp:sec:base-change}. Is there a bounded map $\beta$ from $E^p$ to the completion of $(E_{K^{1/p}},\pi_{\max})$, additive in each variable, with $\beta(\dots,\lambda x_i,\dots)=\lambda^{1/p}\beta(\dots,x_i,\dots)$ and $\beta(x,\dots,x)=x\otimes1$?
\item Which bounded strict polynomial functors of degree $n\ge p$ are analytic on all Banach spaces? Is every such functor a natural retract of one built from tensor powers by $\oplus$, $\otimes_\pi$, composition and duality?
\item Let $K$ be complete. Is every DP vector functor on $\mathsf{Ban}_K$ obtained from the analytic vector functors and the completed divided powers $\hGamma^n$ by finite sums, $\otimes_\pi$, composition, closed subfunctors, quotients and isomorphism?
\end{enumerate}
\end{question}

In~(1), $\mathrm{Fr}_P$ is a Frobenius law, so it is ambient of degree $p$ (Theorem~\ref{dp:thm:frobenius}(2)). The twist $u\mapsto u^{(1)}=\id\otimes u$, with values in the completed twists, has $(\sum_ls_lu_l)^{(1)}=\sum_ls_l^pu_l^{(1)}$ and $\norm{u^{(1)}}\le\norm u^p$. So it is a contractive bounded strict polynomial functor of degree $p$ on $\mathsf{Ban}_K$, and a DP vector functor (Theorem~\ref{dp:thm:functors-dp}). Let $K$ be complete. Its hom map on $\cL(K,P)=P$ is $\mathrm{Fr}_P$, so by Lemma~\ref{dp:lem:lift-criterion} a negative answer to~(1) makes it not analytic. The answer is yes for $K$ complete, $[K:K^p]<\infty$ and $P=c_0(J,K)$. Indeed, the $K$-linear map $P^{(1)}\to P$, $c\otimes x\mapsto(cx_i^p)_{i\in J}$, has the inverse $\mathsf D(z)=\sum_\xi\xi\otimes(\mathsf r_\xi(z_i))_i$, where $\xi$ runs over a basis of $K$ over $K^p$ and $c=\sum_\xi\xi\,\mathsf r_\xi(c)^p$. It is bounded, as all $K^p$-norms on $K$ are equivalent, and $(h_1,\dots,h_p)\mapsto\mathsf D((h_{1,i}\cdots h_{p,i})_i)$ is a lift. A test case is $c_0$ with $[K:K^p]=\infty$.

In~(2), for $K$ complete and $E=c_0(J,K)$, the map $\beta(h_1,\dots,h_p)=((h_{1,i}\cdots h_{p,i})^{1/p})_i$ is such a map, into $c_0(J,K^{1/p})$. A test case is $E=\ell^\infty(\N,K)$ with $[K:K^p]=\infty$.

Let $K$ be complete. Item~(4) includes the divided powers because finite sums, closed subfunctors, quotients and isomorphism alone do not reach them. A closed subfunctor of a quotient $\mathfrak F_1/\mathfrak F_2$ is a quotient of a closed subfunctor, and a quotient of $\mathfrak F_1/\mathfrak F_2$ is a quotient of $\mathfrak F_1$. So finite sums, closed subfunctors, quotients and isomorphism do not leave the functors isomorphic to subquotients of analytic vector functors. For $n\ge p$, $\hGamma^n$ is a DP vector functor and not such a subquotient (Proposition~\ref{dp:prop:gamma-not-ambient}(3)). So the DP vector functors are not the smallest class containing the analytic ones with these closure properties. Whether $\otimes_\pi$ and composition reach the divided powers is not decided here.

Let $S^n_\pi E$ be the quotient of $(E^{\otimes n},\pi)$ by the closed span of the $y-\varsigma y$, with $\varsigma$ a permutation of the factors, and $\widehat S^n_\pi E$ its completion.

\begin{conjecture}\label{dp:conj:symmetric}
For every $K$ of odd characteristic $p$ and every $n\ge p$, $\widehat S^n_\pi$ is not an analytic vector functor on $\mathsf{Ban}_K$.
\end{conjecture}

For $n!\ne0$ it is analytic (Theorem~\ref{dp:thm:functor-threshold}).

\begin{question}\label{dp:q:forms}
Let $\operatorname{char}K=p$.
\begin{enumerate}
\item What replaces the constant rank theorem for DP maps?
\item Can two analytic Banach manifolds, or analytic bundles, that are not analytically isomorphic become DP isomorphic?
\end{enumerate}
\end{question}

In~(1), the classical statement fails: $x\mapsto x^p$ on $K$ has zero derivative and is not locally constant.

Let $K$ be complete of characteristic $p$, let $P$ be a Banach space, and let $\mathrm H^k$ be the $k$-th cohomology of the complex of scalar DP germs at $0$ on $P$. By Theorem~\ref{dp:thm:weight-splitting}, for $k\ge1$ every class has a representative with laws of total weight in $p\Z$. For DP functions $\theta,\psi_1,\dots,\psi_k$ put
\[
 \operatorname{Car}^{-1}(\theta\,d\psi_1\wedge\dots\wedge d\psi_k)=\theta^p\psi_1^{p-1}\cdots\psi_k^{p-1}\,d\psi_1\wedge\dots\wedge d\psi_k,
\]
a closed DP form.

\begin{conjecture}\label{dp:conj:cartier}
The map $\operatorname{Car}^{-1}$ extends additively to a well-defined map from finite sums of such products to $\mathrm H^k$.
\end{conjecture}

The image of this map does not span $\mathrm H^k$ over $K$ in general. Let $P=K$ with $[K:K^p]=\infty$, and let $k=1$. Here every $1$-form is closed, and scalar DP germs at $0$ are convergent power series in one variable $x$ (Theorem~\ref{dp:thm:divided-one}).

\emph{Step 1 (a coefficient).} For a power series $\psi$ and $a\ge0$, the coefficient of $x^{pa+p-1}\,dx$ in $\psi^{p-1}\,d\psi$ is the $p$-th power of the coefficient of $x^a$ in $\psi'$. Indeed, this coefficient depends on finitely many coefficients of $\psi$. Modulo exact forms, $\psi\mapsto\psi^{p-1}\,d\psi$ is additive, as the polynomial $((y+z)^p-y^p-z^p)/p$ has integer coefficients. Exact forms have no monomial $x^{pa+p-1}\,dx$. For $\psi=cx^i$ we get $ic^px^{ip-1}\,dx$, and $ic^p=(ic)^p$. So the coefficient of $x^{pa+p-1}\,dx$ in $\theta^p\psi^{p-1}\,d\psi$ is the $p$-th power of the coefficient of $x^a$ in $\theta\psi'$, and it lies in $K^p$.

\emph{Step 2 (the class).} Choose $\mathsf c_a\in K$, $a\ge0$, linearly independent over $K^p$, with $\abs{\mathsf c_a}\le1$; this is possible after scaling by elements of $K^p$. The form $\sum_a\mathsf c_ax^{pa+p-1}\,dx$ is DP at $0$ with radius $1$. Suppose that it equals $\sum_{l\le\mathsf L}\lambda_l\theta_l^p\psi_l^{p-1}\,d\psi_l$ plus an exact form near $0$, with $\lambda_l\in K$. By Step~1, every $\mathsf c_a$ lies in the $K^p$-span of $\lambda_1,\dots,\lambda_{\mathsf L}$. This is false, as that span is finite-dimensional over $K^p$.

The space $P=K$ satisfies $(\mathrm S^j_P)$ and $(\mathrm R^j_P)$ of \cite[Section~3]{I} for every $j\ge1$, as $\Lambda^j_\pi K=0$ for $j\ge2$. The form of Step~2 is a convergent series of values of $\operatorname{Car}^{-1}$ with coefficients in $K$.

\begin{question}\label{dp:q:cartier}
Let $P$ satisfy $(\mathrm S^j_P)$ and $(\mathrm R^j_P)$ of \cite[Section~3]{I} for every $j\ge1$. Is every class of $\mathrm H^k$, $k\ge1$, represented by the sum of a series, convergent on a ball, of forms $c\,\theta^p\psi_1^{p-1}\cdots\psi_k^{p-1}\,d\psi_1\wedge\dots\wedge d\psi_k$ with $c\in K$?
\end{question}

The hypothesis comes from forms with values in a Banach space. Under condition~(L) of \cite[Section~7.1]{I}, for $Y=c_0$, the closed constant form $\omega_Y$ with values in $\hLambda^p_\pi Y$ has no primitive near any point (\cite[Proposition~10.9]{I}, \cite[Theorem~7.7]{I}). This paper gives no scalar form of this kind.

\begin{question}\label{dp:q:largest}
Is there a list of properties, with no bound on slices or on divided differences, for which the DP germs are the largest class inside $C^\infty_{\mathrm{BGN}}$ with these properties? Two properties that such a list can contain are determination by the Hasse jet at a point, as in Theorem~\ref{dp:thm:quasi-analytic}, and closure under composition.
\end{question}

With a bound on slices in the list the answer is yes and says nothing, by Proposition~\ref{dp:prop:slices}; the same holds for planar divided differences, by Theorem~\ref{dp:thm:divided-planes}. Theorem~\ref{dp:thm:minimal} says that the DP germs are the smallest class with (C1), (C2), (C3).
%% ---- end sec-17-questions.tex
\appendix
%% ---- begin app-A-failure.tex
\section{Proofs for Section~\ref{dp:sec:nonambient}}\label{dp:app:proofs}

\subsection{Tensors as functions}\label{dp:ssec:tensor-notation}

For the proofs of Proposition~\ref{dp:prop:degree-p} and Lemma~\ref{dp:lem:bases} we use the following notation. It is the notation of \cite[Section~Q.5]{I}. For a field $\Bbbk$ and integers $m,n\ge1$, read a tensor $t\in(\Bbbk^m)^{\otimes n}$ as the function $t(i_1,\dots,i_n)$ on $[m]^n$, where $[m]=\{1,\dots,m\}$. For linear forms $\phi_1,\dots,\phi_n$ on $\Bbbk^m$ let $t(\phi_1,\dots,\phi_n)$ be the value of $\phi_1\otimes\dots\otimes\phi_n$ at $t$. Let $\Gamma^n(\Bbbk^m)$ be the space of symmetric tensors, also for finite $\Bbbk$, where it can be larger than the span of the powers.

The proof of Lemma~\ref{dp:lem:c0-lower} uses two more notions from the same place. The \emph{flattening rank} $\operatorname{frk}(t)$ is the rank of the $m\times m^{n-1}$ matrix of the entries $t(i_1,\dots,i_n)$, with rows indexed by $i_1$. A \emph{projection onto the symmetric tensors} is a linear map of $(\Bbbk^m)^{\otimes n}$ into $\Gamma^n(\Bbbk^m)$ that fixes $\Gamma^n(\Bbbk^m)$. Symmetric tensors with entries in a commutative $\Bbbk$\nobreakdash-algebra are combinations of symmetric tensors over $\Bbbk$. So the linear extension of such a map to these entries maps tensors with these entries to symmetric ones and fixes the symmetric ones.

\subsection{Non-ambient laws and composition on \texorpdfstring{$c_0\oplus G$}{c0+G}}\label{dp:af:c0}

In this subsection $p=\operatorname{char}K>0$, and the notation is that of Section~\ref{dp:sec:nonambient}. Let $\kappa$ be the residue field of $K$, so $\F_p\subseteq\kappa$. Let $\operatorname{pr}_m\colon c_0\to K^m$ and $\operatorname{pr}_{G_m}\colon G\to G_m$ be the coordinate projections. Both have norm at most $1$. We write $\mathsf X$ for an indeterminate, which is called $X$ in \cite[Appendix~K]{I}; here $X$ stands for the normed space of \cite[Definition~14.15]{I}.

\begin{proof}[Proof of Lemma~\ref{dp:lem:c0-lower}]
\emph{Step 0 (the field under (a)).} This step is used only under~(a). The completion $\Kh$ has characteristic $p$, so its valuation ring is a complete discrete valuation ring containing the prime field $\F_p$, with $\vartheta$ as a uniformizer and with residue field $\kappa$. By Cohen's structure theorem for complete regular local rings containing a field~\cite{Co}, this ring contains a coefficient field, that is, a field which the residue map sends isomorphically onto $\kappa$. We identify this field with $\kappa$; it contains $\F_p$. As the ring is complete and $\vartheta$ is a uniformizer, every element of the ring is uniquely a series $\sum_{i\ge0}\zeta_i\vartheta^i$ with $\zeta_i\in\kappa$. So the ring is isomorphic to the power series ring $\kappa[[\mathsf X]]$, by an isomorphism taking $\vartheta$ to $\mathsf X$. It extends to an isomorphism of $\Kh$ onto $\kappa((\mathsf X))$, which is isometric for the absolute value $\abs\xi=\abs\vartheta^{\ord\xi}$ on $\kappa((\mathsf X))$. All objects in the lemma are defined from $K$ as a normed field. So under~(a) we may assume that $K$ is a subfield of $\kappa((\mathsf X))$ with this absolute value. Write $\coeff_0$ for the constant coefficient of a Laurent series, applied entrywise to tensors.

\emph{Step 1 (projections for the family).} As $I_m$ contains the tuples of unit vectors, the span $E_0$ of the tensors $y_l$, $l\in I_m$, is $(K^m)^{\otimes p}$. As $K$ is infinite, $\Gamma^p(c_0)$ is the space of symmetric tensors in $c_0^{\otimes p}$, by Section~\ref{dp:sec:conv}. The inclusion of $(K^m)^{\otimes p}$ in $c_0^{\otimes p}$ is injective and commutes with permutations of the factors, so $E_0\cap\Gamma^p(c_0)=\Gamma^p(K^m)$. Let $\Psi\in\mathcal P(y^{(m)})$, in the sense of \cite[Definition~14.15]{I}, and put $\Psi_m=\operatorname{pr}_m^{\otimes p}\circ\Psi$. The map $\operatorname{pr}_m^{\otimes p}$ commutes with permutations of the factors, so it maps $\Gamma^p(c_0)$ into $\Gamma^p(K^m)$, and it is the identity on $(K^m)^{\otimes p}$. As $\Psi$ fixes $\Gamma^p(K^m)$, the map $\Psi_m$ is a projection of $(K^m)^{\otimes p}$ onto the symmetric tensors. As $\operatorname{pr}_m$ has norm at most $1$, the map $\operatorname{pr}_m^{\otimes p}$ does not increase the cost of a representation, so
\begin{equation}\label{dp:af:projection-cost}
 \pi_m(\Psi_my_l)\le\pi(\Psi y_l)\qquad(l\in I_m).
\end{equation}

\emph{Step 2 (item (1) under (a)).} Assume~(a), let $K\subseteq\kappa((\mathsf X))$ be as in Step~0, let $\Psi_K$ be a projection of $(K^m)^{\otimes p}$ onto the symmetric tensors, and put $\sigma_l=\Psi_K(y_l)$. By Section~\ref{dp:ssec:tensor-notation} with $\Bbbk=K$, the linear extension $\widetilde\Psi$ of $\Psi_K$ to entries in $\kappa((\mathsf X))$ is a projection of $(\kappa((\mathsf X))^m)^{\otimes p}$ onto $\Gamma^p(\kappa((\mathsf X))^m)$. For $\mathsf z\in(\kappa^m)^{\otimes p}$ put $\Psi_\kappa(\mathsf z)=\coeff_0\widetilde\Psi(\mathsf z)$. Then $\Psi_\kappa$ is $\kappa$\nobreakdash-linear with symmetric values and fixes the symmetric tensors over $\kappa$. So it is a projection of $(\kappa^m)^{\otimes p}$ onto $\Gamma^p(\kappa^m)$. As the $y_l$ have entries in $\F_p\subseteq K\cap\kappa$, we have $\widetilde\Psi(y_l)=\sigma_l$ and $\Psi_\kappa(y_l)=\coeff_0\sigma_l$.

Read $\sigma_l$ as a matrix with rows indexed by the first index, in the completion of $\linf([m],K)\otimes_\pi\linf([m]^{p-1},K)$. Its norm there is at most $\pi_m(\sigma_l)$, since an elementary tensor of $p-1$ vectors of $K^m$, read as a function on $[m]^{p-1}$, has as maximum norm the product of their maximum norms. By the opening paragraph of \cite[Appendix~K]{I}, \cite[Lemma~K.2(2)]{I} holds for every field of characteristic $p$ in place of the finite field there, so for $\kappa$. With $\abs\vartheta$ in place of $\varrho$ and the index sets $[m]$ and $[m]^{p-1}$, it gives $\operatorname{frk}\Psi_\kappa(y_l)\le M_{\abs\vartheta}\pi_m(\sigma_l)$. Summing over $l$,
\begin{equation}\label{dp:eq:rank-sum}
 \sum_{l\in I_m}\operatorname{frk}\Psi_\kappa(y_l)\le M_{\abs\vartheta}\sum_{l\in I_m}\pi_m(\sigma_l).
\end{equation}
The tensors $y_l$, $l\in I_m$, are the $y^1\otimes\dots\otimes y^p$ with $(y^1,\dots,y^p)\in(\F_p^m)^p$. So by \cite[Lemma~Q.6(4)]{I}, applied to $\Psi_\kappa$ with $\Bbbk=\kappa$, the left side of~\eqref{dp:eq:rank-sum} is at least $mp^{pm}/\bigl(2^p(1+\lfloor\log_pm\rfloor)\bigr)$. Hence $\sum_{l\in I_m}\pi_m(\sigma_l)\ge\mathsf b_m$.

\emph{Step 3 (item (1) under (b)).} Assume~(b), so $p=2$. For $S\subseteq[m]$ let $1_S\in\F_2^m$ have the entries $1$ at the elements of $S$ and $0$ elsewhere, so that $I_m$ consists of the pairs $l=(1_S,1_{S'})$, $S,S'\subseteq[m]$. Identify $(K^m)^{\otimes2}$ with $K^{m\times m}$ by $u\otimes v\mapsto uv^{\mathsf T}$. Then $\Gamma^2(K^m)$ is the space of symmetric matrices, and $y_l=1_S1_{S'}^{\mathsf T}$ for $l=(1_S,1_{S'})$. Let $\Psi_K$ be a projection onto the symmetric tensors, and put $\sigma_l=\Psi_K(y_l)$.

First, $\abs{\det\sigma_{l_0}}\ge1$ for some $l_0\in I_m$. If $m=1$, then $\Psi_K=\id$ and $\sigma_{l_0}=1$ for $l_0=(1_{\{1\}},1_{\{1\}})$. Let $m\ge2$, and let $\wedge\colon K^{m\times m}\to\Lambda^2K^m$ be the linear map of \cite[Theorem~7.20(1)]{I}, with $m$ in place of its $n$:
\[
 \wedge\mathsf Z=\sum_{i<i'}(\mathsf Z_{ii'}+\mathsf Z_{i'i})\,e_i\wedge e_{i'}.
\]
Here $e_1,\dots,e_m$ are the unit vectors of $K^m$. The map is onto, its kernel is the space of symmetric matrices, and $\wedge(uv^{\mathsf T})=u\wedge v$, as $-1=1$ in $K$. As $\Psi_K-\id$ vanishes on the symmetric matrices, it is $\chi\circ\wedge$ for a linear map $\chi\colon\Lambda^2K^m\to K^{m\times m}$. As $\Psi_K$ takes symmetric values, $\wedge\circ\chi\circ\wedge=-\wedge=\wedge$, and $\wedge$ is onto, so $\wedge\circ\chi=\id$. Moreover $\Psi_K(uv^{\mathsf T})=\chi(u\wedge v)+uv^{\mathsf T}$. So \cite[Proposition~N.12]{I}, applied over the field $K$ to $\chi$, with $m$ in place of its $n$, gives $\sum_{S,S'\subseteq[m]}\det\Psi_K(1_S1_{S'}^{\mathsf T})=1$. As $K$ is nonarchimedean, some term has absolute value at least $1$.

By \cite[Lemma~N.15(2)]{I}, $\pi_m(\sigma_{l_0})\ge m$. By \cite[Lemma~Q.6(3)]{I}, applied to the norm $\pi_m$ and to the map $(y^1,y^2)\mapsto\Psi_K(y^1\otimes y^2)$, the average of $\pi_m(\sigma_l)$ over the $4^m$ elements $l\in I_m$ is at least $\pi_m(\sigma_{l_0})/4$. So $\sum_{l\in I_m}\pi_m(\sigma_l)\ge m4^{m-1}=\mathsf b_m$.

\emph{Step 4 (items (2) and (3)).} Let $\norm{y_l}_{\rm e}=\norm{y_l^1}\cdots\norm{y_l^p}$, as in \cite[Definition~14.15]{I}. The factors of the $y_l$ have entries in $\F_p$, so they have norm at most $1$, and the tuple $l=(e_1,\dots,e_1)$ of unit vectors has $\norm{y_l}_{\rm e}=1$. So $0<\sum_{l\in I_m}\norm{y_l}_{\rm e}\le\abs{I_m}=p^{pm}$. Let $\Psi\in\mathcal P(y^{(m)})$. By~\eqref{dp:af:projection-cost} and item~(1) for the projection $\Psi_m$,
\begin{equation}\label{dp:af:family-sum}
 \sum_{l\in I_m}\pi(\Psi y_l)\ge\sum_{l\in I_m}\pi_m(\Psi_my_l)\ge\mathsf b_m.
\end{equation}
Hence $\mathfrak l(y^{(m)})\ge\mathsf b_m/p^{pm}$. This is $m/\bigl(2^pM_{\abs\vartheta}(1+\lfloor\log_pm\rfloor)\bigr)$ under~(a) and $m/4$ under~(b). Both tend to infinity with $m$, so $\mathfrak l_p(c_0)=\infty$.

For item~(3), recall from Section~\ref{dp:sec:nonambient} that $G(y^{(m)})=G_m$, with the same elements $\varepsilon_l$, so that the tensor $t(y^{(m)})=\sum_{l\in I_m}y_l\otimes\varepsilon_l$ of \cite[Definition~14.15]{I} is $t_m$. By \cite[Lemma~Q.7(1)]{I}, with $c_0$ and $p$ in place of its $X$ and $n$, and by~\eqref{dp:af:family-sum}, the norm $\mu$ of $t_m$ in $\Gamma^p_\pi(c_0)\otimes G_m$ is at least $\mathsf b_m$. It remains to pass from $G_m$ to $G$. Let $t_m=\sum_q\mathsf s_q\otimes w_q$ with $\mathsf s_q\in\Gamma^p(c_0)$ and $w_q\in G$. As $\operatorname{pr}_{G_m}$ fixes $G_m$, also $t_m=\sum_q\mathsf s_q\otimes\operatorname{pr}_{G_m}(w_q)$, and this representation in $\Gamma^p(c_0)\otimes G_m$ has no greater cost. So $\mu(t_m)\ge\mathsf b_m$.
\end{proof}

\begin{proof}[Proof of Theorem~\ref{dp:thm:c0-failure}]
We check the hypotheses of \cite[Proposition~14.18]{I} with $c_0$ and $p\ge2$ in place of its $X$ and $n$. The coordinate forms of $c_0$ are bounded, and their products are bounded for $\pi$ and separate the points of $c_0^{\otimes p}$, so $\pi$ is a norm on $c_0^{\otimes p}$. By Lemma~\ref{dp:lem:c0-lower}(2), $\mathfrak l_p(c_0)=\infty$ and $\sup_m\mathfrak l(y^{(m)})=\infty$. The space $G$ is the $\ell^1$-sum of the universal spaces $G(y^{(m)})=G_m$, $m\ge1$, so it is a bad space. With these spaces and this $G$, the objects $H_X$, $W_X$, $Y_X$, $f_X$, $g_X$ of \cite[Definition~14.15]{I} are $H_0$, $W_X$, $Y_X$, $f_X$, $g_X$ of Section~\ref{dp:sec:nonambient}. So \cite[Proposition~14.18]{I} gives items~(1)--(4). Its statement on $g_X\circ f_X$ includes the map into the completion of $Y_X$.
\end{proof}

\begin{proof}[Proof of Proposition~\ref{dp:prop:nonambient}]
\emph{Step 1 (item (1)).} By \cite[Lemma~Q.7(1)]{I}, with $c_0$ and $p$ in place of its $X$ and $n$, the tensor $t_m=t(y^{(m)})$ lies in $\Gamma^p(c_0)\otimes G_m$, and its projective norm in $c_0^{\otimes p}\otimes G_m$ is at most $\sum_{l\in I_m}\norm{y_l}_{\rm e}$, which is at most $p^{pm}$ by Step~4 of the proof of Lemma~\ref{dp:lem:c0-lower}. The inclusion of $G_m$ in $G$ is isometric. So $t_m\in\Gamma^p(c_0)\otimes G=V$ and $\pi(t_m)\le p^{pm}$. By Section~\ref{dp:sec:nonambient}, $\nu\ge\mu$ on $V$, and by Lemma~\ref{dp:lem:c0-lower}(3), $\mu(t_m)\ge\mathsf b_m$. This is item~(1).

\emph{Step 2 (item (2): a law).} Let $n\ge p+1$ and $r\ge1$, and let $(\lambda_i,x_i,g_i)\in H$ and $s_i\in K$ for $1\le i\le r$. By~\eqref{dp:eq:orbit},
\[
 R_n\Bigl(\sum_is_i(\lambda_i,x_i,g_i)\Bigr)=\Bigl(\sum_is_i\lambda_i\Bigr)^{n-p-1}\Bigl(\sum_{\abs\alpha=p}s^\alpha O_\alpha(x_1,\dots,x_r)\Bigr)\otimes\Bigl(\sum_is_ig_i\Bigr),
\]
a homogeneous polynomial of degree $n$ in $s$. Its coefficients lie in $\Gamma^p(c_0)\otimes G=V$, as $O_\alpha(x_1,\dots,x_r)\in\Gamma^p(c_0)$. Moreover $\nu(R_n(\lambda,x,g))\le\abs\lambda^{n-p-1}\norm x^p\norm g\le\norm{(\lambda,x,g)}^n$. So $R_n$ is a law of degree $n$ with $\norm{R_n}\le1$.

\emph{Step 3 (item (2): comparing coefficients).} Suppose $jR_n(h)=\beta(h,\dots,h)$ for $h\in H$, for a linear isometry $j$ of $V$ or of its completion into a normed space $Z$, and a bounded $n$\nobreakdash-linear map $\beta\colon H^n\to Z$. Let $\beta_\otimes\colon(H^{\otimes n},\pi)\to Z$ be the linear map with $\beta_\otimes(h_1\otimes\dots\otimes h_n)=\beta(h_1,\dots,h_n)$, of norm at most $\norm\beta$. Put $\mathbf 1=(1,0,0)\in H$ and $d_n=n!/(p!\,(n-p-1)!)$. For $t\in c_0^{\otimes p}\otimes G$, read in $H^{\otimes(p+1)}$, let $\iota(t)\in H^{\otimes n}$ be the sum of one term for each of the $d_n$ words $\mathsf w\in\{1,2,3\}^n$ with $p$ letters $1$ and one letter $2$. The term of $\mathsf w$ is the tensor obtained from $t\otimes\mathbf 1^{\otimes(n-p-1)}$ by moving the first $p$ factors, in order, to the positions of the letters $1$ of $\mathsf w$, the $(p+1)$-st factor to the position of the letter $2$, and the factors $\mathbf 1$ to the positions of the letters $3$. The inclusions of $c_0$ and $G$ in $H$ are isometric, permutations of the factors preserve the projective norm, and $\norm{\mathbf 1}=1$. So $\iota(t)$ has projective norm at most $d_n\pi(t)$ in $H^{\otimes n}$.

For $x\in c_0$, $g\in G$ and $s=(s_1,s_2,s_3)\in K^3$ put $h_s=(s_3,s_1x,s_2g)$. Then
\[
 \beta_\otimes(h_s^{\otimes n})=\beta(h_s,\dots,h_s)=jR_n(h_s)=s_1^ps_2s_3^{n-p-1}\,j(x^{\otimes p}\otimes g),
\]
and $\iota(x^{\otimes p}\otimes g)$ is the coefficient of $s_1^ps_2s_3^{n-p-1}$ in the polynomial $h_s^{\otimes n}$. As $K$ is infinite, comparing coefficients gives $\beta_\otimes(\iota(x^{\otimes p}\otimes g))=j(x^{\otimes p}\otimes g)$. So $\beta_\otimes\circ\iota=j$ on $V$.

\emph{Step 4 (item (2): the contradiction).} By Steps~1 and~3,
\[
 \nu(t_m)=\norm{j(t_m)}\le\norm\beta\,d_n\pi(t_m)\le d_n\norm\beta\,p^{pm}.
\]
With item~(1) this gives $m\le2^pd_nM_{\abs\vartheta}\norm\beta(1+\lfloor\log_pm\rfloor)$ under~(a) and $m\le4d_n\norm\beta$ under~(b), for every $m\ge1$. As $m/(1+\lfloor\log_pm\rfloor)$ is unbounded, this is impossible.
\end{proof}
%% ---- end app-A-failure.tex
%% ---- begin app-A-degree-p.tex
\subsection{Non-ambient laws over every \texorpdfstring{$K$}{K}}\label{dp:ssec:degree-p}

For the proof of Proposition~\ref{dp:prop:degree-p} we extend the notation of Section~\ref{dp:ssec:tensor-notation} from $\Bbbk^m$ to $\Bbbk^{\mathcal J}$, for a field $\Bbbk$ and a finite set $\mathcal J$. A tensor $t\in(\Bbbk^{\mathcal J})^{\otimes d}$ is read as a function on the words $w=(w_1,\dots,w_d)\in\mathcal J^d$. The space $\Gamma^d(\Bbbk^{\mathcal J})$ is the space of symmetric tensors, and $\Gamma^0(\Bbbk^{\mathcal J})=\Bbbk$. Let $\N^{\mathcal J}_d$ be the set of $\alpha\in\N^{\mathcal J}$ with $\abs\alpha=d$; for $\mathcal J=[m]$ it is the set $\N^m_d$ of Section~\ref{dp:sec:nonambient}. The \emph{content} of a word $w\in\mathcal J^d$ is the $\alpha\in\N^{\mathcal J}_d$ whose entry at each element of $\mathcal J$ is the number of letters of $w$ equal to it. Let $O_\alpha\in\Gamma^d(\Bbbk^{\mathcal J})$ be the sum, over the words $w$ of content $\alpha$, of the tensor product of the unit vectors of $\Bbbk^{\mathcal J}$ at the letters $w_1,\dots,w_d$. The $O_\alpha$, $\alpha\in\N^{\mathcal J}_d$, form a basis of $\Gamma^d(\Bbbk^{\mathcal J})$, and $t=\sum_\alpha t_\alpha O_\alpha$, where $t_\alpha$ is the value of $t$ at any word of content $\alpha$.

The \emph{product} of $t\in\Gamma^d(\Bbbk^{\mathcal J})$ and $t'\in\Gamma^{d'}(\Bbbk^{\mathcal J})$ is the sum $t\cdot t'$ of the $\binom{d+d'}d$ tensors obtained from $t\otimes t'$ by the permutations of the $d+d'$ factors that keep the order of the first $d$ factors and the order of the last $d'$. A scalar in $\Gamma^0(\Bbbk^{\mathcal J})$ multiplies. These permutations represent the cosets of the subgroup that permutes the first $d$ factors and the last $d'$ factors among themselves, and this subgroup fixes $t\otimes t'$. So $t\cdot t'$ is symmetric, and it lies in $\Gamma^{d+d'}(\Bbbk^{\mathcal J})$. For $\Bbbk=K$ this is the shuffle product of Section~\ref{dp:sec:conv}. Tensor powers of linear maps commute with permutations of the factors, so $q^{\otimes(d+d')}(t\cdot t')=q^{\otimes d}(t)\cdot q^{\otimes d'}(t')$ for every linear map $q$ of $\Bbbk^{\mathcal J}$ into some $\Bbbk^{\mathcal J'}$.

Let $\mathcal J=[m]$, $\alpha\in\N^m_d$ and $\alpha'\in\N^m_{d'}$. The value of $O_\alpha\cdot O_{\alpha'}$ at a word of content other than $\alpha+\alpha'$ is $0$. Its value at a word of content $\alpha+\alpha'$ is the number of ways to choose the $d$ of its letters that come from the first factor, which must contain $\alpha_i$ of the $\alpha_i+\alpha'_i$ letters equal to $i$, for each $i$. So
\begin{equation}\label{dp:eq:orbit-product}
 O_\alpha\cdot O_{\alpha'}=\prod_{i=1}^m\binom{\alpha_i+\alpha'_i}{\alpha_i}\,O_{\alpha+\alpha'},
\end{equation}
and the product is commutative and associative on the spaces $\Gamma^d(\Bbbk^m)$, $d\ge0$, as it is bilinear and $\binom{a+b}a\binom{a+b+c}{a+b}=(a+b+c)!/(a!\,b!\,c!)$ for integers $a,b,c\ge0$. For $t\in\Gamma^{d+d'}(\Bbbk^m)$ and $\gamma\in\N^m_d$ let $\partial^\gamma t\in\Gamma^{d'}(\Bbbk^m)$ have the coefficients $(\partial^\gamma t)_\alpha=t_{\alpha+\gamma}$, $\alpha\in\N^m_{d'}$. For $p=\operatorname{char}K>0$, $\Bbbk=\F_p$, $d=p-1$ and $d'=1$ this is the slice of Section~\ref{dp:sec:nonambient}. Then $\partial^\gamma\partial^{\gamma'}=\partial^{\gamma+\gamma'}$. Let $\alpha\in\N^m_{d+d'}$. Then $\partial^\gamma O_\alpha$ is $O_{\alpha-\gamma}$ if $\gamma\le\alpha$ entrywise and $0$ otherwise. By~\eqref{dp:eq:orbit-product}, $(\partial^\gamma O_\alpha)\cdot O_\gamma=\prod_i\binom{\alpha_i}{\gamma_i}O_\alpha$ for $\gamma\le\alpha$. Vandermonde's identity gives $\sum_\gamma\prod_i\binom{\alpha_i}{\gamma_i}=\binom{d+d'}d$, the sum over the $\gamma\in\N^m_d$ with $\gamma\le\alpha$. Hence, by linearity in $t$,
\begin{equation}\label{dp:eq:contraction-sum}
 \sum_{\gamma\in\N^m_d}(\partial^\gamma t)\cdot O_\gamma=\binom{d+d'}d\,t\qquad(t\in\Gamma^{d+d'}(\Bbbk^m)).
\end{equation}

\begin{lemma}[an identity for symmetric tensors]\label{dp:lem:linear}
Let $\Bbbk$ be a field of characteristic $p>0$ and $m\ge1$, let $\mathsf L$ be a matrix over $\Bbbk$ with rows and columns indexed by $\N^m_{p-1}$, and let $\mathsf L^\sharp$ be the linear map of $\Gamma^p(\Bbbk^m)$ with
\[
 \mathsf L^\sharp(t)=\sum_{\gamma,\gamma'\in\N^m_{p-1}}\mathsf L_{\gamma\gamma'}\,(\partial^\gamma t)\cdot O_{\gamma'}.
\]
\begin{enumerate}
\item We have $t=\sum_{\gamma\in\N^m_{p-1}}\bigl((\id-\mathsf L^\sharp)(\partial^\gamma t)\bigr)\cdot O_\gamma$ for every $t\in\Gamma^{2p-1}(\Bbbk^m)$.
\item We have $\rank(\id-\mathsf L^\sharp)\ge\binom{m+2p-2}{2p-1}\big/\binom{m+p-2}{p-1}\ge(p-1)!\,m^p/(2p-1)!$.
\end{enumerate}
\end{lemma}

\begin{proof}
\emph{Step 1 (the identity part).} Let $t\in\Gamma^{2p-1}(\Bbbk^m)$. By~\eqref{dp:eq:contraction-sum} with $d=p-1$ and $d'=p$, we have $\sum_\gamma(\partial^\gamma t)\cdot O_\gamma=\binom{2p-1}{p-1}t$, the sum over $\N^m_{p-1}$. The integer $\binom{2p-1}{p-1}(p-1)!$ is the product $(p+1)(p+2)\cdots(2p-1)$, which is congruent to $(p-1)!$ modulo $p$, and $(p-1)!$ is prime to $p$. So $\binom{2p-1}{p-1}=1$ in $\Bbbk$, and $\sum_\gamma(\partial^\gamma t)\cdot O_\gamma=t$.

\emph{Step 2 (the part of $\mathsf L^\sharp$).} For $\gamma\in\N^m_{p-1}$ we have $\mathsf L^\sharp(\partial^\gamma t)=\sum_{\gamma',\gamma''}\mathsf L_{\gamma'\gamma''}\,(\partial^{\gamma'}\partial^\gamma t)\cdot O_{\gamma''}$. The product is commutative and associative, and $\partial^{\gamma'}\partial^\gamma=\partial^\gamma\partial^{\gamma'}$. So
\[
 \sum_\gamma\mathsf L^\sharp(\partial^\gamma t)\cdot O_\gamma
 =\sum_{\gamma',\gamma''}\mathsf L_{\gamma'\gamma''}\,O_{\gamma''}\cdot
 \sum_\gamma\bigl(\partial^\gamma(\partial^{\gamma'}t)\bigr)\cdot O_\gamma,
\]
all sums over $\N^m_{p-1}$. By~\eqref{dp:eq:contraction-sum} with $d=p-1$ and $d'=1$, applied to $\partial^{\gamma'}t\in\Gamma^p(\Bbbk^m)$, the inner sum is $p\,\partial^{\gamma'}t=0$. With Step~1 this is item~(1).

\emph{Step 3 (item (2)).} By item~(1), the linear map $t\mapsto\bigl((\id-\mathsf L^\sharp)(\partial^\gamma t)\bigr)_{\gamma\in\N^m_{p-1}}$ from $\Gamma^{2p-1}(\Bbbk^m)$ into the product of $\abs{\N^m_{p-1}}$ copies of the image of $\id-\mathsf L^\sharp$ is injective. As $\dim\Gamma^d(\Bbbk^m)=\abs{\N^m_d}=\binom{m+d-1}d$, this gives $\binom{m+2p-2}{2p-1}\le\binom{m+p-2}{p-1}\rank(\id-\mathsf L^\sharp)$. The quotient of the two binomial coefficients is $(p-1)!/(2p-1)!$ times the product of the $p$ numbers $m+p-1,\dots,m+2p-2$, so it is at least $(p-1)!\,m^p/(2p-1)!$.
\end{proof}

\begin{proof}[Proof of Proposition~\ref{dp:prop:degree-p}]
Let $p=\operatorname{char}K>0$ and $m\ge1$, and let $\mathcal J_m$ be the disjoint union of $[m]$ and $\N^m_{p-1}$. For a field $\Bbbk$ containing $\F_p$ write $e_i$, $i\in[m]$, and $e_\gamma$, $\gamma\in\N^m_{p-1}$, for the unit vectors of $\Bbbk^{\mathcal J_m}$, and read $\Bbbk^m$ as the span of the $e_i$. For $T\in\Gamma^p(\F_p^m)$ let $q_T\colon\Bbbk^{\mathcal J_m}\to\Bbbk^m$ be the linear map with $q_T(e_i)=e_i$ and $q_T(e_\gamma)=\partial^\gamma T$. Its matrix has entries in $\F_p$. For a real number $\tau\ge1$ and $z=\sum_iz_ie_i+\sum_\gamma z_\gamma e_\gamma$ in $K^{\mathcal J_m}$ put $\norm z_\tau=\sum_i\abs{z_i}+\tau\sum_\gamma\abs{z_\gamma}$. By~\eqref{dp:eq:weighted-norm}, $\norm h_{T,\tau}=\inf\{\norm z_\tau:q_T(z)=h\}$.

Call the elements $\mathbf a$ of $\N^{\mathcal J_m}_p$ \emph{positions}. Let $\abs{\mathbf a}_\partial$ be the sum of the entries of $\mathbf a$ at the elements of $\N^m_{p-1}$, and call $\mathbf a$ \emph{nonlinear} if $\abs{\mathbf a}_\partial\ge2$. For a word $w\in\mathcal J_m^p$ let $\abs w_\partial$ be the number of its letters in $\N^m_{p-1}$. For $S\in\Gamma^p(\Bbbk^{\mathcal J_m})$ let $\operatorname{supp}S$ be the set of positions $\mathbf a$ with $S_{\mathbf a}\ne0$. Put $\mathsf r_m=(p-1)!\,m^p/(2p-1)!$. In Steps~1 to~3, $T\in\Gamma^p(\F_p^m)$ and $\tau\ge1$ are arbitrary.

\emph{Step 1 (norms).} If $h=q_T(z)$, the $i$-th entry of $h$ is $z_i+\sum_\gamma z_\gamma T_{\gamma^{+i}}$, with $\gamma^{+i}$ as in Section~\ref{dp:sec:nonambient}, of absolute value at most $\norm z_\tau$, as $\abs{T_\alpha}\le1$ and $\tau\ge1$. So no entry of $h\in\mathcal E_{T,\tau}$ has absolute value above $\norm h_{T,\tau}$. Thus the coordinate forms of $\mathcal E_{T,\tau}$ are bounded and separate points. By Section~\ref{dp:sec:conv}, $\pi$ is a norm on $\mathcal E_{T,\tau}^{\otimes p}$, $\norm\cdot_\Gamma$ is a norm on $\Gamma^p(\mathcal E_{T,\tau})$, and $\pi\le\norm\cdot_\Gamma$.

\emph{Step 2 (the projective norm).} In $\mathcal E_{T,\tau}$ we have $\norm{e_i}\le1$ and $\norm{\partial^\gamma T}\le\tau$, by~\eqref{dp:eq:weighted-norm}. For a word $(i_1,\dots,i_{p-1})\in[m]^{p-1}$ let $\gamma(i_1,\dots,i_{p-1})$ be its content. The value of $T$ at $(i_1,\dots,i_{p-1},i)$ is the $i$-th entry of $\partial^{\gamma(i_1,\dots,i_{p-1})}T$, so
\[
 T=\sum_{(i_1,\dots,i_{p-1})\in[m]^{p-1}}e_{i_1}\otimes\dots\otimes e_{i_{p-1}}\otimes\partial^{\gamma(i_1,\dots,i_{p-1})}T.
\]
This is a sum of $m^{p-1}$ elementary tensors whose factors have norms with product at most $\tau$. Hence $\pi(T)\le\tau m^{p-1}$ in $\mathcal E_{T,\tau}^{\otimes p}$.

\emph{Step 3 (a representation gives a symmetric tensor over $\F_p$).} The elements of $K$ of absolute value below $1$ form an $\F_p$\nobreakdash-linear subspace that does not contain $1$. Extend a basis of this subspace, together with $1$, to a basis of $K$ over $\F_p$. So there is an $\F_p$\nobreakdash-linear map $\rho\colon K\to\F_p$ with $\rho(1)=1$ that vanishes on it. Then $\rho$ fixes $\F_p$, and $\abs x\ge1$ for every $x\in K$ with $\rho(x)\ne0$.

Let $\delta>0$ and $T=\sum_lc_lh_l^{\otimes p}$ with $c_l\in K$, $h_l\in\mathcal E_{T,\tau}$ and $\sum_l\abs{c_l}\norm{h_l}_{T,\tau}^p\le\norm T_\Gamma+\delta$. Choose $\hat h_l\in K^{\mathcal J_m}$ with $q_T(\hat h_l)=h_l$ and $\norm{\hat h_l}_\tau\le(1+\delta)\norm{h_l}_{T,\tau}$, with $\hat h_l=0$ if $h_l=0$. The tensor $\widetilde S=\sum_lc_l\hat h_l^{\otimes p}$ is symmetric, and $q_T^{\otimes p}(\widetilde S)=T$. For $z\in K^{\mathcal J_m}$ the sum of the numbers $\tau^{\abs w_\partial}\abs{z^{\otimes p}(w)}$, $w\in\mathcal J_m^p$, is $\norm z_\tau^p$. So, by the triangle inequality,
\begin{equation}\label{dp:eq:weighted-lift}
 \sum_{w\in\mathcal J_m^p}\tau^{\abs w_\partial}\abs{\widetilde S(w)}\le\sum_l\abs{c_l}\norm{\hat h_l}_\tau^p\le(1+\delta)^p(\norm T_\Gamma+\delta).
\end{equation}
Let $S$ be the tensor over $\F_p$ with the values $S(w)=\rho(\widetilde S(w))$. It is symmetric, as $\widetilde S$ is, so $S\in\Gamma^p(\F_p^{\mathcal J_m})$. As the matrix of $q_T$ has entries in $\F_p$ and $\rho$ is $\F_p$\nobreakdash-linear, each value of $q_T^{\otimes p}(S)$ is the image under $\rho$ of the corresponding value of $q_T^{\otimes p}(\widetilde S)=T$. These values lie in $\F_p$, so $q_T^{\otimes p}(S)=T$. For $\mathbf a\in\operatorname{supp}S$ and a word $w$ of content $\mathbf a$ we have $\rho(\widetilde S(w))=S_{\mathbf a}\ne0$, so $\abs{\widetilde S(w)}\ge1$; and $\abs w_\partial=\abs{\mathbf a}_\partial$. Taking one such word for each $\mathbf a\in\operatorname{supp}S$, \eqref{dp:eq:weighted-lift} gives $\sum_{\mathbf a\in\operatorname{supp}S}\tau^{\abs{\mathbf a}_\partial}\le(1+\delta)^p(\norm T_\Gamma+\delta)$. As $\Gamma^p(\F_p^{\mathcal J_m})$ is finite, one $S$ occurs for arbitrarily small $\delta$. So
\begin{equation}\label{dp:eq:support-bound}
 \norm T_\Gamma\ge\sum_{\mathbf a\in\operatorname{supp}S}\tau^{\abs{\mathbf a}_\partial}\quad\text{for some }S\in\Gamma^p(\F_p^{\mathcal J_m})\text{ with }q_T^{\otimes p}(S)=T.
\end{equation}

\emph{Step 4 (supports with few nonlinear positions).} Let $S\in\Gamma^p(\F_p^{\mathcal J_m})$. Let $\mathcal K_S\subseteq\N^m_p$ be the set of the $\gamma^{+i}$, $i\in[m]$, for the $\gamma\in\N^m_{p-1}$ at which some nonlinear $\mathbf a\in\operatorname{supp}S$ has a nonzero entry. A position has nonzero entries at no more than $p$ elements of $\N^m_{p-1}$, so $\abs{\mathcal K_S}$ is at most $pm$ times the number of nonlinear positions in $\operatorname{supp}S$. We show that at most $p^{\abs{\mathcal K_S}+\binom{m+p-1}p-\mathsf r_m}$ tensors $T'\in\Gamma^p(\F_p^m)$ satisfy $q_{T'}^{\otimes p}(S)=T'$.

Write $S=S^{(0)}+S^{(1)}+S^{(2)}$, where $S^{(0)}$, $S^{(1)}$ and $S^{(2)}$ are supported in the positions $\mathbf a$ with $\abs{\mathbf a}_\partial=0$, $\abs{\mathbf a}_\partial=1$ and $\abs{\mathbf a}_\partial\ge2$. The positions with $\abs{\mathbf a}_\partial=0$ are the $\alpha\in\N^m_p$, extended by $0$, and $q_{T'}$ fixes the $e_i$. So $q_{T'}^{\otimes p}(S^{(0)})=\sum_{\alpha\in\N^m_p}S_\alpha O_\alpha$ does not depend on $T'$.

For $\gamma,\gamma'\in\N^m_{p-1}$ let $\mathbf a(\gamma,\gamma')$ be the position with the entry $1$ at $\gamma$, the entries of $\gamma'$ on $[m]$ and $0$ elsewhere. These are the positions with $\abs{\mathbf a}_\partial=1$. Also $O_{\mathbf a(\gamma,\gamma')}=e_\gamma\cdot O_{\gamma'}$, with $\gamma'$ extended by $0$, since each word of content $\mathbf a(\gamma,\gamma')$ arises from exactly one of the $p$ permutations that define this product, applied to exactly one tensor $e_\gamma\otimes e_{i_1}\otimes\dots\otimes e_{i_{p-1}}$ with $(i_1,\dots,i_{p-1})$ of content $\gamma'$. As the product commutes with $q_{T'}$, and $q_{T'}$ fixes the $e_i$,
\[
 q_{T'}^{\otimes p}(S^{(1)})=\sum_{\gamma,\gamma'\in\N^m_{p-1}}S_{\mathbf a(\gamma,\gamma')}\,(\partial^\gamma T')\cdot O_{\gamma'}=\mathsf L^\sharp(T'),
\]
with $\mathsf L_{\gamma\gamma'}=S_{\mathbf a(\gamma,\gamma')}$ in Lemma~\ref{dp:lem:linear}, for $\Bbbk=\F_p$.

Finally, $q_{T'}^{\otimes p}(S^{(2)})$ depends only on the vectors $q_{T'}(e_\gamma)=\partial^\gamma T'$ for the $\gamma$ at which some $\mathbf a\in\operatorname{supp}S^{(2)}$ has a nonzero entry. That is, it depends only on the coefficients $T'_\alpha$, $\alpha\in\mathcal K_S$. So if $q_{T'}^{\otimes p}(S)=T'$, then $(\id-\mathsf L^\sharp)(T')$ equals $q_{T'}^{\otimes p}(S^{(0)})+q_{T'}^{\otimes p}(S^{(2)})$, which depends only on $(T'_\alpha)_{\alpha\in\mathcal K_S}$. For each of the $p^{\abs{\mathcal K_S}}$ choices of these coefficients the solutions $T'$ lie in one coset of the kernel of $\id-\mathsf L^\sharp$. This kernel has at most $p^{\binom{m+p-1}p-\mathsf r_m}$ elements, by Lemma~\ref{dp:lem:linear}(2), as $\dim\Gamma^p(\F_p^m)=\binom{m+p-1}p$.

\emph{Step 5 (item (1)).} Put $\mathsf n_1=\mathsf r_m/(3p^2(1+\log_pm))$, $\mathsf n_2=\mathsf r_m/(3pm)$ and $\tau=\max\{1,(m/(p(1+\log_pm)))^{1/2}\}$. There are $\binom{\abs{\mathcal J_m}+p-1}p\le\abs{\mathcal J_m}^p$ positions, and $\abs{\mathcal J_m}=m+\binom{m+p-2}{p-1}\le m+m^{p-1}\le2m^{p-1}$. So $p$ times the number of positions is at most $p\,2^pm^{p(p-1)}\le p^{p^2}m^{p^2}=p^{p^2(1+\log_pm)}$, as $p\,2^p\le p^{p^2}$.

Every $S\in\Gamma^p(\F_p^{\mathcal J_m})$ that is nonzero at fewer than $\mathsf n_1$ positions is the sum of $\lceil\mathsf n_1\rceil-1$ tensors $O_{\mathbf a}$, $\mathbf a$ a position, each multiplied by an element of $\F_p$. By the bound on the number of positions above, and as $\lceil\mathsf n_1\rceil-1<\mathsf n_1$, there are at most $p^{\mathsf n_1p^2(1+\log_pm)}=p^{\mathsf r_m/3}$ of them. If moreover fewer than $\mathsf n_2$ of its positions are nonlinear, then $\abs{\mathcal K_S}\le pm\mathsf n_2=\mathsf r_m/3$, and by Step~4 at most $p^{\binom{m+p-1}p-2\mathsf r_m/3}$ tensors $T'$ satisfy $q_{T'}^{\otimes p}(S)=T'$. So at most $p^{\binom{m+p-1}p-\mathsf r_m/3}$ of the $p^{\binom{m+p-1}p}$ elements $T'$ of $\Gamma^p(\F_p^m)$ satisfy $q_{T'}^{\otimes p}(S)=T'$ for such an $S$.

As $\mathsf r_m>0$, these are not all of $\Gamma^p(\F_p^m)$. Choose $T\in\Gamma^p(\F_p^m)$ outside them. Then every $S$ with $q_T^{\otimes p}(S)=T$ is nonzero at no fewer than $\mathsf n_1$ positions, or at no fewer than $\mathsf n_2$ nonlinear positions. The weights $\tau^{\abs{\mathbf a}_\partial}$ in~\eqref{dp:eq:support-bound} are at least $1$, and at least $\tau^2$ for nonlinear $\mathbf a$. Also $\tau^2\mathsf n_2\ge\mathsf n_1$, as $\tau^2\ge m/(p(1+\log_pm))$. So $\norm T_\Gamma\ge\mathsf n_1>0$ in $\Gamma^p(\mathcal E_{T,\tau})$, and $T\ne0$, so $\pi(T)>0$ by Step~1. If $\tau>1$, Step~2 gives
\[
 \frac{\norm T_\Gamma}{\pi(T)}\ge\frac{\mathsf n_1}{\tau m^{p-1}}=\frac{(p-1)!}{3p^{3/2}(2p-1)!}\Bigl(\frac m{1+\log_pm}\Bigr)^{1/2}.
\]
If $\tau=1$, then $m\le p(1+\log_pm)$, so the right side of~\eqref{dp:eq:degree-p-ratio} is at most $(p-1)!\,\pi(T)/(3p(2p-1)!)<\pi(T)\le\norm T_\Gamma$, by Step~1.

\emph{Step 6 (item (2): the law).} Let $T^{(m)},\tau^{(m)},\mathcal E_m$, $m\ge1$, and $\mathcal E$ be as in item~(2), let $n\ge p$, and let $\mathsf g_n\colon K\oplus\mathcal E\to\Gamma^p(\mathcal E)$ be the map $\mathsf g_n(\lambda,h)=\lambda^{n-p}h^{\otimes p}$, with $\lambda^0=1$. Each norm of $\mathcal E_m$ lies between the maximum of the absolute values of the entries, by Step~1, and their sum. So if $K$ is complete, the spaces $\mathcal E_m$ are complete, and so is their $\ell^1$-sum $\mathcal E$. By Step~1 the coordinate forms of the spaces $\mathcal E_m$, composed with the projections of $\mathcal E$ onto them, are bounded, and they separate points. So $\pi$ is a norm on $\mathcal E^{\otimes p}$, and $\norm\cdot_\Gamma\ge\pi$ is a norm on $\Gamma^p(\mathcal E)$.

For finitely many vectors $(\lambda_l,h_l)$ of $K\oplus\mathcal E$ and scalars $s_l$, the value of $\mathsf g_n$ at $\sum_ls_l(\lambda_l,h_l)$ is $(\sum_ls_l\lambda_l)^{n-p}(\sum_ls_lh_l)^{\otimes p}$. By~\eqref{dp:eq:orbit} this is a homogeneous polynomial of degree $n$ in the $s_l$ with coefficients in $\Gamma^p(\mathcal E)$. Moreover $\norm{\lambda^{n-p}h^{\otimes p}}_\Gamma\le\abs\lambda^{n-p}\norm h^p\le\norm{(\lambda,h)}^n$. So $\mathsf g_n$ is a bounded polynomial law of degree $n$ and norm at most $1$.

\emph{Step 7 (item (2): the ratio on $\Gamma^p(\mathcal E)$).} The inclusion of $\mathcal E_m$ in $\mathcal E$ is isometric, and the projection of $\mathcal E$ onto $\mathcal E_m$ is a contractive left inverse. Let $t\in\Gamma^p(\mathcal E_m)$. Then $\pi(t)$ computed in $\mathcal E^{\otimes p}$ is at most $\pi(t)$ computed in $\mathcal E_m^{\otimes p}$. Also $\norm t_\Gamma$ computed in $\Gamma^p(\mathcal E)$ is at least $\norm t_\Gamma$ computed in $\Gamma^p(\mathcal E_m)$, since the $p$-th tensor power of the projection fixes $t$ and does not increase the cost $\sum_l\abs{c_l}\norm{h_l}^p$ of a representation. By item~(1), applied to $t=T^{(m)}\ne0$, and as $m/(1+\log_pm)\to\infty$ with $m$, the quotient $\norm\cdot_\Gamma/\pi$ is unbounded on the nonzero elements of $\Gamma^p(\mathcal E)$.

\emph{Step 8 (item (2): comparing coefficients).} Suppose that $j\mathsf g_n(\lambda,h)=\beta((\lambda,h),\dots,(\lambda,h))$ for all $(\lambda,h)\in K\oplus\mathcal E$, for a linear isometry $j$ of $\Gamma^p(\mathcal E)$, or of its completion, into a normed space $Z$ and a bounded $n$\nobreakdash-linear map $\beta\colon(K\oplus\mathcal E)^n\to Z$. Put $\mathbf 1=(1,0)$ and $h^+=(0,h)$ for $h\in\mathcal E$. For $h_1,\dots,h_p\in\mathcal E$ let $\beta'(h_1,\dots,h_p)$ be the sum, over the $\binom np$ sets of $p$ of the $n$ slots, of the value of $\beta$ at the $n$-tuple with $h_1^+,\dots,h_p^+$ in these slots, in this order, and $\mathbf 1$ in the others. Then $\beta'$ is $p$\nobreakdash-linear, and $\norm{\beta'}\le\binom np\norm\beta$, as $\norm{\mathbf 1}=1$ and $\norm{h^+}=\norm h$.

For $h\in\mathcal E$ and $s=(s_1,s_2)\in K^2$ we have $(s_2,s_1h)=s_2\mathbf 1+s_1h^+$. So $\beta((s_2,s_1h),\dots,(s_2,s_1h))$ is a homogeneous polynomial of degree $n$ in $s$ whose coefficient of $s_1^ps_2^{n-p}$ is $\beta'(h,\dots,h)$. On the other side, $j\mathsf g_n(s_2,s_1h)=s_1^ps_2^{n-p}j(h^{\otimes p})$. As $K$ is infinite, comparing coefficients gives $j(h^{\otimes p})=\beta'(h,\dots,h)$. Let $\beta'_\otimes\colon(\mathcal E^{\otimes p},\pi)\to Z$ be the linear map with $\beta'_\otimes(h_1\otimes\dots\otimes h_p)=\beta'(h_1,\dots,h_p)$, of norm $\norm{\beta'}$. As $\Gamma^p(\mathcal E)$ is spanned by powers, $j(t)=\beta'_\otimes(t)$ for $t\in\Gamma^p(\mathcal E)$. So $\norm t_\Gamma=\norm{j(t)}\le\binom np\norm\beta\,\pi(t)$ for $t\in\Gamma^p(\mathcal E)$, against Step~7. So $\mathsf g_n$ is not ambient of degree $n$, also as a map into the completion of $\Gamma^p(\mathcal E)$.
\end{proof}
%% ---- end app-A-degree-p.tex
%% ---- begin app-A-constants.tex
\subsection{Laws of degree \texorpdfstring{$p$}{p}}\label{dp:ac:frobenius}

In this subsection $p=\operatorname{char}K>0$, and $\Gamma^p(P)$, $\norm\cdot_\Gamma$, $\norm\cdot_{\mathrm F}$, $\pi$ and the hollow tensors are as in Section~\ref{dp:sec:nonambient}. For a subspace $P_0$ of $P$ we read $P_0^{\otimes p}$ as a subspace of $P^{\otimes p}$, since $K$ is a field. For $h_1,\dots,h_r\in P$ and $\alpha\in\N^r_p$, the orbit sum $O_\alpha(h_1,\dots,h_r)$ of Section~\ref{dp:sec:conv} is the sum of the tensors $h_{w_1}\otimes\dots\otimes h_{w_p}$ over the words $w\in[r]^p$ of content $\alpha$, in the sense of Section~\ref{dp:ssec:degree-p}, and it lies in $\Gamma^p(P)$. With the unit vectors of $K^r$ in place of $h_1,\dots,h_r$ it is the tensor $O_\alpha$ of Section~\ref{dp:ssec:degree-p}, for $\Bbbk=K$ and $\mathcal J=[r]$. We use~\eqref{dp:eq:orbit} with $n=p$.

We use \cite[Lemma~Q.1]{I} in every degree $n\ge0$. A law $a\colon P\to W$ of degree $n$ has a unique linear map $L_a\colon\Gamma^n(P)\to W$ with $L_a(h^{\otimes n})=a(h)$, and $\norm{L_a(t)}\le\norm a\norm t_\Gamma$ (item~(1)). Without the bound this is~\cite[Th.~IV.1, Prop.~IV.5]{Ro}. A map $a$ with $ja(h)=\beta(h,\dots,h)$ for all $h\in P$, for a linear isometry $j$ of $W$ into a normed space and a bounded $n$\nobreakdash-linear map $\beta$ on $P^n$, is such a law, and $\norm{L_a(t)}\le\norm\beta\pi(t)$ (item~(3)). If $\norm{L_a(t)}\le C\pi(t)$ for all $t\in\Gamma^n(P)$, then $a$ is ambient of degree $n$, with an $n$\nobreakdash-linear map of norm at most $C$ (item~(4)).

\begin{lemma}[linearization]\label{dp:lem:linearization}
\leavevmode
\begin{enumerate}
\item If $a\colon P\to W$ is a Frobenius law, then $L_a$ vanishes on the hollow tensors, and $\norm{L_a(t)}\le\norm a\norm t_{\mathrm F}$ for $t\in\Gamma^p(P)$.
\item If $P$ satisfies $(\mathrm F_P)$ with constant $C_{\mathrm F}$ and $(\mathrm N_P)$ with constant $C_{\mathrm N}$, then $\norm t_\Gamma\le\bigl(C_{\mathrm F}+(1+C_{\mathrm F})C_{\mathrm N}\bigr)\pi(t)$ for $t\in\Gamma^p(P)$.
\end{enumerate}
\end{lemma}

\begin{proof}
\emph{Step 1 (item (1)).} Let $h_1,\dots,h_p\in P$. By~\eqref{dp:eq:orbit}, for $s_1,\dots,s_p\in K$,
\[
 \sum_{\alpha\in\N^p_p}s_1^{\alpha_1}\cdots s_p^{\alpha_p}L_a\bigl(O_\alpha(h_1,\dots,h_p)\bigr)
 =a(s_1h_1+\dots+s_ph_p)=s_1^pa(h_1)+\dots+s_p^pa(h_p).
\]
As $K$ is infinite and $p\ge2$, comparing the coefficients of $s_1\cdots s_p$ gives $L_a(O_{(1,\dots,1)}(h_1,\dots,h_p))=0$. The words of content $(1,\dots,1)$ are the orderings of $[p]$, so $O_{(1,\dots,1)}(h_1,\dots,h_p)$ is one of the tensors that span the hollow tensors. So $L_a$ vanishes on the hollow tensors. If $t-t'$ is hollow, then $\norm{L_a(t)}=\norm{L_a(t')}\le\norm a\norm{t'}_\Gamma$; the infimum over $t'$ is the bound.

\emph{Step 2 (item (2)).} Let $t\in\Gamma^p(P)$ and $\delta>0$. By $(\mathrm F_P)$ there is $t'\in\Gamma^p(P)$ with $t-t'$ hollow and $\norm{t'}_\Gamma\le C_{\mathrm F}\pi(t)+\delta$. As $\pi\le\norm\cdot_\Gamma$, we get $\pi(t-t')\le\pi(t)+\norm{t'}_\Gamma\le(1+C_{\mathrm F})\pi(t)+\delta$. By $(\mathrm N_P)$, $\norm t_\Gamma\le\norm{t'}_\Gamma+C_{\mathrm N}\pi(t-t')\le C_{\mathrm F}\pi(t)+\delta+C_{\mathrm N}\bigl((1+C_{\mathrm F})\pi(t)+\delta\bigr)$. Let $\delta\to0$.
\end{proof}

\begin{lemma}[bases of representations]\label{dp:lem:bases}
Let $P_0$ be a finite-dimensional subspace of $P$ and $t\in\Gamma^p(P)\cap P_0^{\otimes p}$. For linear forms $\phi_1,\dots,\phi_p$ on $P_0$ let $t(\phi_1,\dots,\phi_p)$ be the value of $\phi_1\otimes\dots\otimes\phi_p$ at $t$, as in Section~\ref{dp:ssec:tensor-notation}.
\begin{enumerate}
\item Let $b_1,\dots,b_m$ be a basis of $P_0$, with dual basis $b_1^*,\dots,b_m^*$. Then $t-\sum_{i=1}^mt(b_i^*,\dots,b_i^*)\,b_i^{\otimes p}$ is hollow.
\item Let $t=\sum_{l=1}^rx_{l,1}\otimes\dots\otimes x_{l,p}$ with nonzero vectors $x_{l,c}$ that span $P_0$. There is a set $\mathcal B\subseteq[r]\times[p]$ such that the vectors $x_{l,c}$, $(l,c)\in\mathcal B$, form a basis of $P_0$ whose dual basis $(x_{l,c}^*)_{(l,c)\in\mathcal B}$ satisfies
\[
 \abs{t(x_{l,c}^*,\dots,x_{l,c}^*)}\,\norm{x_{l,c}}^p\le\norm{x_{l,1}}\cdots\norm{x_{l,p}}\qquad\bigl((l,c)\in\mathcal B\bigr).
\]
\end{enumerate}
\end{lemma}

The congruence in item~(1) belongs to the algebra of divided powers; compare \cite[Th.~IV.2, Prop.~III.1 and App.~I, Props.~1,~2]{Ro}. The bound in item~(2) is the estimate we need.

\begin{proof}
\emph{Step 1 (item (1)).} Let $\Pi_0$ be a linear projection of $P$ onto $P_0$. As $\Pi_0^{\otimes p}$ fixes $P_0^{\otimes p}$ and maps powers to powers, $t=\Pi_0^{\otimes p}(t)$ is a combination of powers $h^{\otimes p}$, $h\in P_0$. For a word $w\in[m]^p$ put $b_w=b_{w_1}\otimes\dots\otimes b_{w_p}$. These tensors form a basis of $P_0^{\otimes p}$, and $t=\sum_wt(b_{w_1}^*,\dots,b_{w_p}^*)\,b_w$. For a power $h^{\otimes p}$, $h\in P_0$, the coefficient $b_{w_1}^*(h)\cdots b_{w_p}^*(h)$ of $b_w$ depends only on the content $\alpha\in\N^m_p$ of $w$. Hence so does that of $t$, and $t$ is a combination of the tensors $O_\alpha(b_1,\dots,b_m)$.

Let $\alpha\in\N^m_p$. If every $\alpha_i<p$, then $\alpha!$ is invertible in $K$. For a word $w$ of content $\alpha$, the hollow tensor $O_{(1,\dots,1)}(b_{w_1},\dots,b_{w_p})$ is $\alpha!\,O_\alpha(b_1,\dots,b_m)$, as the reorderings of the letters of $w$ give each word of content $\alpha$ exactly $\alpha!$ times. So $O_\alpha(b_1,\dots,b_m)$ is hollow. Otherwise $\alpha_i=p$ for one $i$, the tensor $O_\alpha(b_1,\dots,b_m)$ is $b_i^{\otimes p}$, and its coefficient is $t(b_i^*,\dots,b_i^*)$. Summing over $\alpha$ gives item~(1).

\emph{Step 2 (item (2): the choice of $\mathcal B$).} Call $\mathcal B\subseteq[r]\times[p]$ a \emph{basis set} if the vectors $x_{l,c}$, $(l,c)\in\mathcal B$, form a basis of $P_0$. Basis sets exist, as these vectors span $P_0$. Fix a basis of $P_0$, and for a basis set $\mathcal B$ let $\det_{\mathcal B}$ be the determinant of the coordinates of the vectors $x_{l,c}$, $(l,c)\in\mathcal B$, in a fixed order of $[r]\times[p]$. Put $\upsilon_{l,c}=(\norm{x_{l,1}}\cdots\norm{x_{l,p}})^{1/p}/\norm{x_{l,c}}>0$, so that $\upsilon_{l,1}\cdots\upsilon_{l,p}=1$ for each $l$. Choose, among the finitely many basis sets, a basis set $\mathcal B$ with the largest value of $\abs{\det_{\mathcal B}}\prod_{(l,c)\in\mathcal B}\upsilon_{l,c}$. Then
\begin{equation}\label{dp:eq:exchange}
 \abs{x_{l,c}^*(x_{l',c'})}\,\upsilon_{l',c'}\le\upsilon_{l,c}\qquad
 \bigl((l,c)\in\mathcal B,\ (l',c')\in[r]\times[p]\bigr).
\end{equation}
This is clear if $(l',c')\in\mathcal B$ or $x_{l,c}^*(x_{l',c'})=0$. Otherwise the expansion of $x_{l',c'}$ in the basis has the nonzero coefficient $x_{l,c}^*(x_{l',c'})$ at $(l,c)$. So replacing $(l,c)$ by $(l',c')$ gives a basis set $\mathcal B'$, and $\det_{\mathcal B'}=\pm x_{l,c}^*(x_{l',c'})\det_{\mathcal B}$, since the determinant is multilinear and alternating. The choice of $\mathcal B$ gives $\abs{x_{l,c}^*(x_{l',c'})}\,\upsilon_{l',c'}/\upsilon_{l,c}\le1$.

\emph{Step 3 (item (2): the bound).} Let $(l,c)\in\mathcal B$. Then $t(x_{l,c}^*,\dots,x_{l,c}^*)=\sum_{l'=1}^r\prod_{c'=1}^px_{l,c}^*(x_{l',c'})$. By the ultrametric inequality in $K$ and~\eqref{dp:eq:exchange},
\[
 \abs{t(x_{l,c}^*,\dots,x_{l,c}^*)}\le\max_{l'}\prod_{c'=1}^p\frac{\upsilon_{l,c}}{\upsilon_{l',c'}}=\upsilon_{l,c}^p
 =\frac{\norm{x_{l,1}}\cdots\norm{x_{l,p}}}{\norm{x_{l,c}}^p}.\qedhere
\]
\end{proof}

\begin{proof}[Proof of Theorem~\ref{dp:thm:frobenius}]
\emph{Step 1 (item (1)).} Let $t\in\Gamma^p(P)$ and $\delta>0$; we may assume $t\ne0$. Choose a representation $t=\sum_{l=1}^rx_{l,1}\otimes\dots\otimes x_{l,p}$ with $\sum_l\norm{x_{l,1}}\cdots\norm{x_{l,p}}\le\pi(t)+\delta$. Dropping the terms with a zero factor, all $x_{l,c}$ are nonzero. Let $P_0$ be their span, so $\dim P_0\le\dim P$ and $t\in P_0^{\otimes p}$, and take $\mathcal B$ from Lemma~\ref{dp:lem:bases}(2). By Lemma~\ref{dp:lem:bases}(1), $t-t'$ is hollow for $t'=\sum_{(l,c)\in\mathcal B}t(x_{l,c}^*,\dots,x_{l,c}^*)\,x_{l,c}^{\otimes p}$. So
\begin{align*}
 \norm t_{\mathrm F}\le\norm{t'}_\Gamma
 &\le\sum_{(l,c)\in\mathcal B}\abs{t(x_{l,c}^*,\dots,x_{l,c}^*)}\,\norm{x_{l,c}}^p\\
 &\le\sum_{l=1}^r\abs{\mathcal B\cap(\{l\}\times[p])}\,\norm{x_{l,1}}\cdots\norm{x_{l,p}}
 \le\min(p,\dim P)\,(\pi(t)+\delta),
\end{align*}
since $\mathcal B$ has $\dim P_0\le\dim P$ elements. Let $\delta\to0$.

\emph{Step 2 (item (2)).} By Lemma~\ref{dp:lem:linearization}(1) and item~(1), $\norm{L_a(t)}\le\norm a\norm t_{\mathrm F}\le p\norm a\,\pi(t)$ for $t\in\Gamma^p(P)$. Apply \cite[Lemma~Q.1(4)]{I}.

\emph{Step 3 (item (3): the bound).} By item~(1) and Lemma~\ref{dp:lem:linearization}(2) with $C_{\mathrm F}=p$ and $C_{\mathrm N}=C$, $\norm t_\Gamma\le(p+(p+1)C)\pi(t)$ for $t\in\Gamma^p(P)$. So $\norm{L_a(t)}\le(p+(p+1)C)\norm a\,\pi(t)$, by \cite[Lemma~Q.1(1)]{I}. Apply \cite[Lemma~Q.1(4)]{I}.

\emph{Step 4 (item (3): the converse).} By Theorem~\ref{dp:thm:normed-roby}, $\mathsf v_p\colon P\to\Gamma^p_{\mathrm s}(P)$ is a law of degree $p$ whose linearization is the quotient map $t\mapsto[t]$. By hypothesis $j\mathsf v_p(h)=\beta(h,\dots,h)$ for a linear isometry $j$ of $\Gamma^p_{\mathrm s}(P)$ into a normed space and a bounded $p$\nobreakdash-linear map $\beta$ on $P^p$. By \cite[Lemma~Q.1(3)]{I}, $\norm t_\Gamma=\norm{[t]}\le\norm\beta\pi(t)$ for $t\in\Gamma^p(P)$, in particular for hollow $t$.
\end{proof}

\subsection{Properties of the ambience constants}\label{dp:ac:constants}

From here on $n\ge1$ and the characteristic of $K$ is arbitrary, with $\Gamma^n(P)$, $\norm\cdot_\Gamma$, $\pi$, $\mathfrak a_n(P)$, $[m]$, $c_0([m],K)$ and $\mathfrak c_K$ as in Sections~\ref{dp:sec:conv} and~\ref{dp:sec:nonambient}. By Theorem~\ref{dp:thm:normed-roby}, the map $\mathsf v_n\colon P\to\Gamma^n_{\mathrm s}(P)$, $\mathsf v_n(h)=[h^{\otimes n}]$, is a law of degree $n$ and norm at most $1$, and its linearization is the quotient map $t\mapsto[t]$, with $\norm{[t]}=\norm t_\Gamma$. For a subspace $P_0$ of $P$ we read $P_0^{\otimes n}$ as a subspace of $P^{\otimes n}$. When two spaces or two fields occur, we say in which space and over which field $\pi$, $\norm\cdot_\Gamma$ and $\mathfrak a_n$ are computed. The completion $\widehat P$ of $P$ is a normed space over $\Kh$, the scalar multiplication extending by continuity, and it is ultrametric if $P$ is.

\begin{lemma}[properties of the ambience constants]\label{dp:lem:constants}
Let $P'$ be a normed space.
\begin{enumerate}
\item If $j\mathsf v_n(h)=\beta(h,\dots,h)$ for all $h\in P$, for a linear isometry $j$ of $\Gamma^n_{\mathrm s}(P)$ into a normed space and a bounded $n$\nobreakdash-linear map $\beta$ on $P^n$, then $\mathfrak a_n(P)\le\norm\beta$.
\item If $\Xi$ is a bounded linear bijection of $P$ onto $P'$ with bounded inverse, then $\mathfrak a_n(P')\le(\norm\Xi\norm{\Xi^{-1}})^n\mathfrak a_n(P)$.
\item If $P'$ is a subspace of $P$, with the restricted norm, and some linear projection of $P$ onto $P'$ has norm at most $1$, then $\mathfrak a_n(P')\le\mathfrak a_n(P)$.
\item $\mathfrak a_n(P)$ is at most the supremum of $\mathfrak a_n(P_0)$ over the finite-dimensional subspaces $P_0$ of $P$, with the restricted norm.
\end{enumerate}
\end{lemma}

\begin{proof}
\emph{Step 1 (item (1)).} By \cite[Lemma~Q.1(3)]{I}, $\norm t_\Gamma=\norm{[t]}\le\norm\beta\pi(t)$ for $t\in\Gamma^n(P)$. So $\mathfrak a_n(P)\le\norm\beta$.

\emph{Step 2 (item (2)).} We may assume that $\mathfrak a_n(P)$ is finite. Let $t'\in\Gamma^n(P')$ and $t=(\Xi^{-1})^{\otimes n}(t')$. The maps $\Xi^{\otimes n}$ and $(\Xi^{-1})^{\otimes n}$ carry powers to powers and elementary tensors to elementary tensors, so $t\in\Gamma^n(P)$. They multiply the cost of a representation by at most $\norm\Xi^n$ and $\norm{\Xi^{-1}}^n$, respectively. Hence $\norm{t'}_\Gamma\le\norm\Xi^n\norm t_\Gamma\le\norm\Xi^n\mathfrak a_n(P)\pi(t)\le(\norm\Xi\norm{\Xi^{-1}})^n\mathfrak a_n(P)\pi(t')$.

\emph{Step 3 (item (3)).} We may assume that $\mathfrak a_n(P)$ is finite. Let $\Pi$ be such a projection and $t\in\Gamma^n(P')$. Then $\Pi^{\otimes n}$ fixes $t$ and carries a representation of $t$ in $\Gamma^n(P)$ to one in $\Gamma^n(P')$ of no larger cost, and representations of $t$ in $P'^{\otimes n}$ are representations in $P^{\otimes n}$. So $\norm t_\Gamma$ computed in $\Gamma^n(P')$ is at most $\norm t_\Gamma$ computed in $\Gamma^n(P)$, which is at most $\mathfrak a_n(P)\pi(t)$ with $\pi(t)$ computed in $P^{\otimes n}$, hence at most $\mathfrak a_n(P)\pi(t)$ with $\pi(t)$ computed in $P'^{\otimes n}$.

\emph{Step 4 (item (4)).} Let $t\in\Gamma^n(P)$ and $\delta>0$, and let $t=\sum_lx_{l,1}\otimes\dots\otimes x_{l,n}$ with $\sum_l\norm{x_{l,1}}\cdots\norm{x_{l,n}}\le\pi(t)+\delta$. Let $P_0$ be the span of the $x_{l,c}$. Let $\Pi_0$ be a linear projection of $P$ onto $P_0$. As $\Pi_0^{\otimes n}$ fixes $P_0^{\otimes n}$ and maps powers to powers, $t=\Pi_0^{\otimes n}(t)\in\Gamma^n(P_0)$. Then $\norm t_\Gamma$ computed in $\Gamma^n(P)$ is at most $\norm t_\Gamma$ computed in $\Gamma^n(P_0)$. As the representation lies in $P_0^{\otimes n}$, this is at most $\mathfrak a_n(P_0)(\pi(t)+\delta)$ with $\pi(t)$ computed in $P^{\otimes n}$, and so at most the supremum in~(4) times $\pi(t)+\delta$. Let $\delta\to0$.
\end{proof}

\begin{lemma}[orthogonal bases]\label{dp:lem:orthogonal}
Let $K$ be complete, and let $P$ be ultrametric, of finite dimension $m\ge1$. Then $\mathfrak a_n(P)\le\mathfrak c_K^{-n}\mathfrak a_n(c_0([m],K))$.
\end{lemma}

\begin{proof}
If $K$ is discretely valued, put $\mathfrak o=1$ and $\mathfrak o'=\mathfrak c_K$; otherwise let $0<\mathfrak o<1$ and put $\mathfrak o'=\mathfrak o$. By \cite[Lemma~Q.13(3)]{I} there is a linear bijection $\mathsf C\colon c_0([m],K)\to P$ with $\norm{\mathsf C}\le1$ and $\norm{\mathsf C^{-1}}\le(\mathfrak o\mathfrak o')^{-1}$. Lemma~\ref{dp:lem:constants}(2) gives $\mathfrak a_n(P)\le(\mathfrak o\mathfrak o')^{-n}\mathfrak a_n(c_0([m],K))$. This is the claim if $K$ is discretely valued; otherwise let $\mathfrak o\to1$.
\end{proof}

\begin{lemma}[completion]\label{dp:lem:completion}
\leavevmode
\begin{enumerate}
\item $\mathfrak a_n(P)$ is at most $\mathfrak a_n(\widehat P)$ computed over $\Kh$.
\item For $m\ge1$, $\mathfrak a_n(c_0([m],\Kh))$ computed over $\Kh$ is at most $\mathfrak a_n(c_0([m],K))$.
\end{enumerate}
\end{lemma}

\begin{proof}
\emph{Step 1 (item (1)).} Let $C=\mathfrak a_n(\widehat P)$, computed over $\Kh$, be finite. By \cite[Lemma~Q.13(1)]{I}, $\mathsf v_n$ extends to a law $\hat{\mathsf v}_n\colon\widehat P\to\hGamma^n(P)$ of degree $n$ over $\Kh$ and norm at most $1$. By \cite[Lemma~Q.1(1)]{I}, its linearization $\hat L$ over $\Kh$ satisfies $\norm{\hat L(t)}\le\norm t_\Gamma\le C\pi(t)$ on $\Gamma^n(\widehat P)$. So \cite[Lemma~Q.1(4)]{I} over $\Kh$ gives a $\Kh$\nobreakdash-linear isometry $\hat\jmath$ of $\hGamma^n(P)$ into a normed space $Z$ over $\Kh$ and a bounded $n$\nobreakdash-linear map $\hat\beta$ on $\widehat P^n$ over $\Kh$, of norm at most $C$, with $\hat\jmath\hat{\mathsf v}_n(h)=\hat\beta(h,\dots,h)$ for $h\in\widehat P$. The restrictions of $\hat\jmath$ to $\Gamma^n_{\mathrm s}(P)$ and of $\hat\beta$ to $P^n$ are a linear isometry into $Z$, read as a normed space over $K$, and a bounded $n$\nobreakdash-linear map of norm at most $C$, and $\hat\jmath\mathsf v_n(h)=\hat\beta(h,\dots,h)$ for $h\in P$. Lemma~\ref{dp:lem:constants}(1) gives $\mathfrak a_n(P)\le C$.

\emph{Step 2 (projective norms in $c_0([m],K)$).} Read $K^m\subseteq\Kh^m$, so that the tensors $e_{w_1}\otimes\dots\otimes e_{w_n}$, $w\in[m]^n$, are a basis of $(K^m)^{\otimes n}$ over $K$ and of $(\Kh^m)^{\otimes n}$ over $\Kh$. Call the coefficients of a tensor in this basis its coefficients. For $z\in(K^m)^{\otimes n}$ we show that $\pi(z)$ is the same over $K$ and over $\Kh$, the spaces being $c_0([m],K)$ and $c_0([m],\Kh)$. Over $\Kh$ it is at most its value over $K$. Over $K$, $\pi(z)$ is at most the sum of the absolute values of the coefficients of $z$.

Let $z=\sum_lx_{l,1}\otimes\dots\otimes x_{l,n}$ with $x_{l,c}\in\Kh^m$, let $\mathfrak r=\sum_l\norm{x_{l,1}}\cdots\norm{x_{l,n}}$, and let $0<\delta\le1$. Choose $x'_{l,c}\in K^m$ with $\norm{x'_{l,c}-x_{l,c}}\le\delta\norm{x_{l,c}}$; then $\norm{x'_{l,c}}\le\norm{x_{l,c}}$. Every coefficient of
\[
 x_{l,1}\otimes\dots\otimes x_{l,n}-x'_{l,1}\otimes\dots\otimes x'_{l,n}
 =\sum_{c=1}^nx'_{l,1}\otimes\dots\otimes x'_{l,c-1}\otimes(x_{l,c}-x'_{l,c})\otimes x_{l,c+1}\otimes\dots\otimes x_{l,n}
\]
has absolute value at most $\delta\norm{x_{l,1}}\cdots\norm{x_{l,n}}$. So $z'=\sum_lx'_{l,1}\otimes\dots\otimes x'_{l,n}$ lies in $(K^m)^{\otimes n}$, $\pi(z')\le\mathfrak r$ over $K$, and $z-z'$ has coefficients of absolute value at most $\delta\mathfrak r$. Over $K$ this gives $\pi(z)\le\mathfrak r+m^n\delta\mathfrak r$. Let $\delta\to0$ and take the infimum over the representations of $z$ over $\Kh$.

\emph{Step 3 (item (2)).} Let $\mathsf A=\mathfrak a_n(c_0([m],K))$ be finite, let $t\in\Gamma^n(c_0([m],\Kh))$, computed over $\Kh$, and let $\delta>0$. Write $O_\alpha=O_\alpha(e_1,\dots,e_m)$ for $\alpha\in\N^m_n$. By Section~\ref{dp:ssec:degree-p}, with $\Bbbk=\Kh$ and $\mathcal J=[m]$, we have $t=\sum_\alpha t_\alpha O_\alpha$ with $t_\alpha\in\Kh$, as $t$ is symmetric. By Section~\ref{dp:sec:conv} the $O_\alpha$ lie in $\Gamma^n(c_0([m],K))$, computed over $K$; let $\mathfrak s$ be the sum of their seminorms $\norm{O_\alpha}_\Gamma$ there. Choose $t'_\alpha\in K$ with $\abs{t_\alpha-t'_\alpha}\le\delta$ and put $t'=\sum_\alpha t'_\alpha O_\alpha$. Representations over $K$ are representations over $\Kh$, so, over $\Kh$, $\norm{t-t'}_\Gamma\le\delta\mathfrak s$, hence $\pi(t-t')\le\delta\mathfrak s$ and $\pi(t')\le\pi(t)+\delta\mathfrak s$. With Step~2,
\[
 \norm t_\Gamma\le\norm{t'}_\Gamma+\delta\mathfrak s\le\mathsf A\pi(t')+\delta\mathfrak s\le\mathsf A(\pi(t)+\delta\mathfrak s)+\delta\mathfrak s,
\]
where $\norm t_\Gamma$ and $\pi(t)$ are computed over $\Kh$, $\norm{t'}_\Gamma$ over $K$, which is at least its value over $\Kh$, and $\pi(t')$ over either field. Let $\delta\to0$.
\end{proof}

\subsection{Degree two}\label{dp:ac:degree-two}

\begin{lemma}[degree two]\label{dp:lem:degree-two}
For every $m\ge1$, $\mathfrak a_2(c_0([m],K))\le3/2$.
\end{lemma}

Without the condition $\norm a<\infty$, the laws $a$ of degree $2$ are the quadratic maps~\cite[Prop.~II.1]{Ro}.

\begin{proof}
Identify $(K^m)^{\otimes2}$ with $K^{m\times m}$ by $x\otimes x'\mapsto xx'^{\mathsf T}$, as in Step~3 of the proof of Lemma~\ref{dp:lem:c0-lower}. Then the elements of $\Gamma^2(c_0([m],K))$ are symmetric matrices, $\norm{\mathsf S}_\Gamma$ is the infimum of $\sum_l\abs{\lambda_l}\norm{h_l}^2$ over the representations $\mathsf S=\sum_l\lambda_lh_lh_l^{\mathsf T}$, and $\pi(\mathsf S)$ is the infimum of $\sum_l\norm{x_l}\norm{x'_l}$ over the representations $\mathsf S=\sum_lx_lx_l'^{\mathsf T}$. A matrix is \emph{integral} if its entries have absolute value at most $1$. For a matrix $\mathsf M$ and $k\ge1$ let $\mathfrak d_k(\mathsf M)$ be the largest absolute value of a $k\times k$ minor of $\mathsf M$, or $0$ if there is none, and put $\mathfrak d_0(\mathsf M)=1$. A determinant is a sum of products with one entry from each column. So, by the ultrametric inequality, the minors of an integral matrix have absolute value at most $1$, and, by the Cauchy--Binet formula, $\mathfrak d_k(\mathsf M)$ does not increase when $\mathsf M$ is multiplied on either side by an integral matrix.

\emph{Step 1 (a pivot).} Let $\mathsf S\in K^{m\times m}$ be symmetric and nonzero, with entries $\mathsf S_{ii'}$, and let $\mathfrak f$ be the largest absolute value of an entry. If $\abs{\mathsf S_{ii}}=\mathfrak f$ for some $i$, put $\mathcal V=\{i\}$; otherwise put $\mathcal V=\{i,i'\}$ for some $i\ne i'$ with $\abs{\mathsf S_{ii'}}=\mathfrak f$. Let $\mathsf S_{\mathcal V}$ be the submatrix of $\mathsf S$ with rows and columns in $\mathcal V$, and $\mathsf S_{[m]\mathcal V}$ the one with columns in $\mathcal V$. Then $\abs{\det\mathsf S_{\mathcal V}}=\mathfrak f^{\abs{\mathcal V}}$; in the second case $\det\mathsf S_{\mathcal V}=\mathsf S_{ii}\mathsf S_{i'i'}-\mathsf S_{ii'}^2$ with $\abs{\mathsf S_{ii}\mathsf S_{i'i'}}<\mathfrak f^2$. So $\mathsf S_{\mathcal V}^{-1}$, its adjugate divided by its determinant, has entries of absolute value at most $\mathfrak f^{-1}$, and $\mathsf G=\mathsf S_{[m]\mathcal V}\mathsf S_{\mathcal V}^{-1}$ is integral, with the identity matrix as its rows in $\mathcal V$. Put $\mathsf S'=\mathsf S-\mathsf G\mathsf S_{\mathcal V}\mathsf G^{\mathsf T}$. It is symmetric with entries of absolute value at most $\mathfrak f$. Its rows and columns in $\mathcal V$ vanish, as the rows in $\mathcal V$ of $\mathsf G\mathsf S_{\mathcal V}\mathsf G^{\mathsf T}$ form $\mathsf S_{\mathcal V}\mathsf G^{\mathsf T}=\mathsf S_{[m]\mathcal V}^{\mathsf T}$.

If $\mathcal V=\{i\}$, then $\mathsf G\mathsf S_{\mathcal V}\mathsf G^{\mathsf T}=\mathsf S_{ii}uu^{\mathsf T}$ for the column $u$ of $\mathsf G$. If $\mathcal V=\{i,i'\}$ and $u_1$, $u_2$ are the columns of $\mathsf G$ at $i$ and $i'$, then
\[
 \mathsf G\mathsf S_{\mathcal V}\mathsf G^{\mathsf T}=(\mathsf S_{ii}-\mathsf S_{ii'})u_1u_1^{\mathsf T}
 +(\mathsf S_{i'i'}-\mathsf S_{ii'})u_2u_2^{\mathsf T}+\mathsf S_{ii'}(u_1+u_2)(u_1+u_2)^{\mathsf T},
\]
three terms with coefficients of absolute value at most $\mathfrak f$ and vectors of norm at most $1$. In both cases
\begin{equation}\label{dp:eq:pivot-cost}
 \norm{\mathsf G\mathsf S_{\mathcal V}\mathsf G^{\mathsf T}}_\Gamma\le\tfrac32\abs{\mathcal V}\mathfrak f.
\end{equation}

Let $\mathsf U\in K^{m\times m}$ have the columns of $\mathsf G$ at the indices in $\mathcal V$ and the columns of the identity matrix $\id$ elsewhere, and let $\mathsf F\in K^{m\times m}$ agree with $\mathsf S$ at the positions in $\mathcal V\times\mathcal V$ and vanish elsewhere. The nonzero columns of $\mathsf U-\id$ lie at indices in $\mathcal V$, and its rows in $\mathcal V$ vanish. So $(\mathsf U-\id)^2=0$, $\mathsf U^{-1}=2\id-\mathsf U$ is integral, and $\mathsf U\mathsf S'=\mathsf S'=\mathsf S'\mathsf U^{\mathsf T}$, since the rows and columns of $\mathsf S'$ in $\mathcal V$ vanish. As $\mathsf U\mathsf F\mathsf U^{\mathsf T}=\mathsf G\mathsf S_{\mathcal V}\mathsf G^{\mathsf T}$, this gives $\mathsf S=\mathsf U(\mathsf F+\mathsf S')\mathsf U^{\mathsf T}$, so $\mathfrak d_k(\mathsf S)\ge\mathfrak d_k(\mathsf F+\mathsf S')$.

For $0\le k_1\le\min(k,\abs{\mathcal V})$, a $k_1\times k_1$ minor of $\mathsf S_{\mathcal V}$ and a $(k-k_1)\times(k-k_1)$ minor of $\mathsf S'$ with rows and columns outside $\mathcal V$ together form a $k\times k$ minor of $\mathsf F+\mathsf S'$ whose matrix is, after reordering its rows and columns, block diagonal. As the nonzero minors of $\mathsf S'$ have rows and columns outside $\mathcal V$,
\begin{equation}\label{dp:eq:minor-split}
 \mathfrak d_k(\mathsf S)\ge\mathfrak d_{k_1}(\mathsf S_{\mathcal V})\,\mathfrak d_{k-k_1}(\mathsf S').
\end{equation}
Here $\mathfrak d_1(\mathsf S_{\mathcal V})=\mathfrak f$, and $\mathfrak d_2(\mathsf S_{\mathcal V})=\mathfrak f^2$ if $\abs{\mathcal V}=2$.

\emph{Step 2 (the pivot sequence).} Apply Step~1 to $\mathsf S$, then to $\mathsf S'$ if it is nonzero, and so on. A row that vanishes stays zero, as the row of $\mathsf G$ at its index, and so that of $\mathsf G\mathsf S_{\mathcal V}\mathsf G^{\mathsf T}$, vanishes. The rows in $\mathcal V$, which contain an entry of absolute value $\mathfrak f$, vanish after the step. So after at most $m$ steps the matrix is zero, and $\mathsf S$ is the sum of the matrices $\mathsf G\mathsf S_{\mathcal V}\mathsf G^{\mathsf T}$ of the steps. List the numbers $\mathfrak f$ of the steps, each once if $\abs{\mathcal V}=1$ and twice if $\abs{\mathcal V}=2$, as $\mathfrak f_1,\dots,\mathfrak f_{\mathfrak n}$. As the entries of $\mathsf S'$ have absolute value at most $\mathfrak f$, we have $\mathfrak f_1\ge\dots\ge\mathfrak f_{\mathfrak n}>0$. By~\eqref{dp:eq:pivot-cost}, $\norm{\mathsf S}_\Gamma\le\frac32(\mathfrak f_1+\dots+\mathfrak f_{\mathfrak n})$. By~\eqref{dp:eq:minor-split}, the last sentence of Step~1 and induction on the number of steps, applied to $\mathsf S'$, whose sequence is $\mathfrak f_{\abs{\mathcal V}+1},\dots,\mathfrak f_{\mathfrak n}$,
\begin{equation}\label{dp:eq:minor-lower}
 \mathfrak d_k(\mathsf S)\ge\mathfrak f_1\cdots\mathfrak f_k\qquad(1\le k\le\mathfrak n):
\end{equation}
for $k\le\abs{\mathcal V}$ take $k_1=k$, and for $k>\abs{\mathcal V}$ take $k_1=\abs{\mathcal V}$.

\emph{Step 3 (products and sums).} Let $\mathfrak g_1,\dots,\mathfrak g_{\mathfrak n}>0$ with $\mathfrak f_1\cdots\mathfrak f_k\le\mathfrak g_1\cdots\mathfrak g_k$ for $1\le k\le\mathfrak n$. Then $\mathfrak f_1+\dots+\mathfrak f_{\mathfrak n}\le\mathfrak g_1+\dots+\mathfrak g_{\mathfrak n}$. Indeed, $\log(\mathfrak g_k/\mathfrak f_k)\le\mathfrak g_k/\mathfrak f_k-1$, so $\mathfrak g_k-\mathfrak f_k\ge\mathfrak f_k\log(\mathfrak g_k/\mathfrak f_k)$. Put $\mathfrak h_k=\sum_{i\le k}\log(\mathfrak g_i/\mathfrak f_i)\ge0$, $\mathfrak h_0=0$ and $\mathfrak f_{\mathfrak n+1}=0$. Summation by parts gives
\[
 \sum_{k=1}^{\mathfrak n}\mathfrak f_k\log(\mathfrak g_k/\mathfrak f_k)=\sum_{k=1}^{\mathfrak n}\mathfrak f_k(\mathfrak h_k-\mathfrak h_{k-1})
 =\sum_{k=1}^{\mathfrak n}(\mathfrak f_k-\mathfrak f_{k+1})\mathfrak h_k\ge0.
\]

\emph{Step 4 (the projective norm).} Let $\mathsf S=\sum_{l=1}^rx_lx_l'^{\mathsf T}$, and let $\mathfrak g_1\ge\mathfrak g_2\ge\dots\ge\mathfrak g_r$ be the numbers $\norm{x_l}\norm{x'_l}$ in nonincreasing order. By the Cauchy--Binet formula, a $k\times k$ minor of $\mathsf S$ is the sum, over the sets of $k$ indices $l$, of the products of the corresponding $k\times k$ minors of the matrices with the columns $x_l$ and $x'_l$, $l$ in the set. By the ultrametric inequality such a minor of the first matrix has absolute value at most the product of the $\norm{x_l}$, and likewise for the second. So $\mathfrak d_k(\mathsf S)\le\mathfrak g_1\cdots\mathfrak g_k$ if $k\le r$, and $\mathfrak d_k(\mathsf S)=0$ if $k>r$. By~\eqref{dp:eq:minor-lower}, $r\ge\mathfrak n$ and $\mathfrak f_1\cdots\mathfrak f_k\le\mathfrak g_1\cdots\mathfrak g_k$ for $k\le\mathfrak n$, so $\mathfrak g_1,\dots,\mathfrak g_{\mathfrak n}>0$. Step~3 gives $\mathfrak f_1+\dots+\mathfrak f_{\mathfrak n}\le\mathfrak g_1+\dots+\mathfrak g_{\mathfrak n}\le\sum_l\norm{x_l}\norm{x'_l}$. Taking the infimum over the representations, $\mathfrak f_1+\dots+\mathfrak f_{\mathfrak n}\le\pi(\mathsf S)$, and Step~2 gives $\norm{\mathsf S}_\Gamma\le\frac32\pi(\mathsf S)$, also for $\mathsf S=0$.
\end{proof}

\subsection{Proofs of the propositions on ambience constants}\label{dp:ac:proofs}

\begin{proof}[Proof of Proposition~\ref{dp:prop:ultrametric}]
\emph{Step 1 (item (1)).} Let $\mathfrak a_n(P)<\infty$ and let $a\colon P\to W$ be a law of degree $n$. By \cite[Lemma~Q.1(1)]{I}, $\norm{L_a(t)}\le\norm a\norm t_\Gamma\le\mathfrak a_n(P)\norm a\,\pi(t)$ for $t\in\Gamma^n(P)$. Apply \cite[Lemma~Q.1(4)]{I}. Conversely, if every law of degree $n$ on $P$ is ambient of degree $n$, so is $\mathsf v_n$, and Lemma~\ref{dp:lem:constants}(1) gives $\mathfrak a_n(P)<\infty$.

\emph{Step 2 (item (2)).} We may assume that $C=\mathfrak a_{n+1}(K\oplus P)$ is finite. The map $a_+(\lambda,h)=\lambda\mathsf v_n(h)$ from $K\oplus P$ to $\Gamma^n_{\mathrm s}(P)$ is a law of degree $n+1$ with $\norm{a_+}\le1$: for finitely many $(\lambda_l,h_l)\in K\oplus P$ and scalars $s_l$, we have $a_+(\sum_ls_l(\lambda_l,h_l))=(\sum_ls_l\lambda_l)\,\mathsf v_n(\sum_ls_lh_l)$, and $\norm{a_+(\lambda,h)}\le\abs\lambda\norm h^n\le\norm{(\lambda,h)}^{n+1}$. By item~(1), $ja_+(\lambda,h)=\beta((\lambda,h),\dots,(\lambda,h))$ for a linear isometry $j$ of $\Gamma^n_{\mathrm s}(P)$ into a normed space and a bounded $(n+1)$\nobreakdash-linear map $\beta$ on $(K\oplus P)^{n+1}$ with $\norm\beta\le C$.

Put $\mathbf 1=(1,0)$ and $h^+=(0,h)$ for $h\in P$. For $h_1,\dots,h_n\in P$ let $\beta'(h_1,\dots,h_n)$ be the sum, over the $n+1$ slots, of the value of $\beta$ at the $(n+1)$\nobreakdash-tuple with $\mathbf 1$ in this slot and $h_1^+,\dots,h_n^+$ in the others, in this order. Then $\beta'$ is $n$\nobreakdash-linear, and $\norm{\beta'}\le(n+1)\norm\beta$, as $\norm{\mathbf 1}=1$ and $\norm{h^+}=\norm h$. For $h\in P$ and $s\in K$, the value $\beta((1,sh),\dots,(1,sh))$, with $(1,sh)=\mathbf 1+sh^+$, is a polynomial in $s$ whose coefficient of $s^n$ is $\beta'(h,\dots,h)$, while $ja_+(1,sh)=s^nj\mathsf v_n(h)$. As $K$ is infinite, comparing coefficients gives $j\mathsf v_n(h)=\beta'(h,\dots,h)$. Lemma~\ref{dp:lem:constants}(1) gives $\mathfrak a_n(P)\le\norm{\beta'}\le(n+1)C$.

\emph{Step 3 (item (3)).} The second inequality follows from Lemma~\ref{dp:lem:constants}(3), applied to the coordinate projection of $c_0(\N,K)$ onto $c_0([m],K)$. For the first, by Lemma~\ref{dp:lem:constants}(4), and as $\mathfrak a_n(\{0\})=0$, it suffices to show $\mathfrak a_n(P_0)\le\mathfrak c_K^{-n}\sup_{m\ge1}\mathfrak a_n(c_0([m],K))$ for every finite-dimensional subspace $P_0\ne\{0\}$ of $P$. Its completion $\widehat P_0$ is an ultrametric normed space over $\Kh$. By \cite[Lemma~Q.13(2)]{I} its dimension $m'$ over $\Kh$ satisfies $1\le m'<\infty$, and $\mathfrak c_{\Kh}=\mathfrak c_K$. Lemma~\ref{dp:lem:completion}(1), Lemma~\ref{dp:lem:orthogonal} over $\Kh$ and Lemma~\ref{dp:lem:completion}(2) give
\[
 \mathfrak a_n(P_0)\le\mathfrak a_n(\widehat P_0)\le\mathfrak c_K^{-n}\,\mathfrak a_n(c_0([m'],\Kh))
 \le\mathfrak c_K^{-n}\,\mathfrak a_n(c_0([m'],K)),
\]
where the second and third constants are computed over $\Kh$. The ultrametric norm of $P$ enters items~(3) and~(4) only here: it makes $\widehat P_0$ ultrametric, as Lemma~\ref{dp:lem:orthogonal} requires.

\emph{Step 4 (item (4)).} Apply item~(3) with $n=2$ and Lemma~\ref{dp:lem:degree-two}.
\end{proof}

\begin{proof}[Proof of Proposition~\ref{dp:prop:all-norms}]
\emph{Step 1 (item (1)).} Let $n!\ne0$ in $K$, and let $t=\sum_lx_{l,1}\otimes\dots\otimes x_{l,n}\in\Gamma^n(P)$ with nonzero factors; dropping the terms with a zero factor does not raise the cost of a representation. As $t$ is symmetric, $n!\,t=\sum_l\sum_{\sigma\in S_n}x_{l,\sigma(1)}\otimes\dots\otimes x_{l,\sigma(n)}$.

Fix $l$, and choose $\lambda_c\in K$ with $\abs\varpi<\norm{\lambda_cx_{l,c}}\le1$, using powers of $\varpi$. Then $\abs{\lambda_c}^{-1}<\norm{x_{l,c}}/\abs\varpi$. Put $v_c=\lambda_cx_{l,c}$. By~\cite[(4.2)]{I},
\[
 \Bigl(\prod_{c=1}^n\lambda_c\Bigr)\sum_{\sigma\in S_n}x_{l,\sigma(1)}\otimes\dots\otimes x_{l,\sigma(n)}
 =\sum_{\mathcal I\subseteq[n]}(-1)^{n-\abs{\mathcal I}}\Bigl(\sum_{c\in\mathcal I}v_c\Bigr)^{\otimes n}.
\]
Here $\norm{\sum_{c\in\mathcal I}v_c}\le n$, and $\le1$ if $P$ is ultrametric. There are $2^n$ sets $\mathcal I$. So, by~\eqref{dp:eq:gamma-norm},
\[
 \Bigl\|\sum_{\sigma\in S_n}x_{l,\sigma(1)}\otimes\dots\otimes x_{l,\sigma(n)}\Bigr\|_\Gamma
 \le\Bigl(\frac{2n}{\abs\varpi}\Bigr)^n\norm{x_{l,1}}\cdots\norm{x_{l,n}},
\]
with $2/\abs\varpi$ in place of $2n/\abs\varpi$ if $P$ is ultrametric. Summing over $l$ and taking the infimum over the representations gives $\abs{n!}\norm t_\Gamma\le(2n/\abs\varpi)^n\pi(t)$, and likewise in the ultrametric case. As $n!\ne0$, this is the bound on $\mathfrak a_n(P)$.

\emph{Step 2 (item (2)).} Let $n\ge p=\operatorname{char}K>0$. By Proposition~\ref{dp:prop:degree-p}(2) there is a law of degree $n$ on $K\oplus\mathcal E$ that is not ambient of degree $n$. By Proposition~\ref{dp:prop:ultrametric}(1), $\mathfrak a_n(K\oplus\mathcal E)=\infty$.

\emph{Step 3 (the last assertion).} If $n!\ne0$ in $K$, the bound of item~(1) does not depend on $P$. If $n!=0$ in $K$, then $n\ge p=\operatorname{char}K>0$, and item~(2) applies.
\end{proof}
%% ---- end app-A-constants.tex
%% ---- begin bib.tex
\begin{thebibliography}{99}
\bibitem{BGN} W.~Bertram, H.~Gl\"ockner and K.-H.~Neeb,
\emph{Differential calculus, manifolds and Lie groups over arbitrary infinite fields},
arXiv:math/0303300 (2003); published as \emph{Differential calculus over general base fields and rings}, Expo.\ Math.\ \textbf{22} (2004), 213--282,\\
{\small\url{https://arxiv.org/abs/math/0303300}}.

\bibitem{Bou} N.~Bourbaki, \emph{Vari\'et\'es diff\'erentielles et analytiques. Fascicule de r\'esultats}, Hermann, Paris,
1967 (\S\S1--7) and 1971 (\S\S8--15); reprinted by Springer, 2007. We cite the reprint.

\bibitem{Co} I.~S.~Cohen, \emph{On the structure and ideal theory of complete local rings}, Trans.\ Amer.\ Math.\ Soc.\
\textbf{59} (1946), 54--106.

\bibitem[I]{I} J.~McCarthy, \emph{Alternating bundles and the splitting of exterior powers}, in preparation.

\bibitem{VdP} M.~van der Put, \emph{Espaces de Banach non archim\'ediens}, Bull.\ Soc.\ Math.\ France
\textbf{97} (1969), 309--320,\\
{\small\url{https://www.numdam.org/article/BSMF_1969__97__309_0.pdf}}.

\bibitem{Ro} N.~Roby, \emph{Lois polynomes et lois formelles en th\'eorie des modules}, Ann.\ Sci.\ \'Ecole Norm.\
Sup.\ (3) \textbf{80} (1963), 213--348.

\end{thebibliography}
%% ---- end bib.tex
\end{document}
